%!TEX program = lualatex
\documentclass[letterpaper, 11pt]{report}
\usepackage{fontspec}
\usepackage{fancyhdr}

%%% MARGINS %%%
% use this for normal letter documents
% \usepackage[
%     letterpaper,
%     margin=1in,
%     includehead
% ]{geometry}

%use this for bookletized pdfs
\usepackage[
    paperwidth=6.875in,
    paperheight=10.625in,
    margin=0.75in,     % Now your margins will be exactly what you type
    includehead
]{geometry}

%%% FONT %%%
\setmainfont{Liberation Serif}

%%% SPEAKER \& TALK VARIABLES %%%
\newcommand{\speaker}{}
\newcommand{\talk}{}
\newcommand{\talkdatestr}{}
\newcommand{\setspeaker}[1]{\renewcommand{\speaker}{#1}}
\newcommand{\settalk}[1]{\renewcommand{\talk}{#1}}
\newcommand{\settalkdate}[1]{\renewcommand{\talkdatestr}{#1}}

%%% HEADER \& FOOTER SETUP %%%
\pagestyle{fancy}
\fancyhf{}

\fancyhead[L]{\textbf{\talk}}
\fancyhead[R]{\textbf{\speaker}}

\renewcommand{\headrulewidth}{0.4pt}

%%% CUSTOM HEADER CLASS FOR TALK TITLE %%%
\newcommand{\talkheading}[2]{%
    \vspace{1em}
    \noindent{\huge\bfseries #1}
    \par\nopagebreak\vspace{0.25em}
    \noindent{\normalsize\bfseries #2\ifx\talkdatestr\empty\else,\ \talkdatestr\fi}
    \par\nopagebreak\vspace{0.5em}
}

%%% IN-TALK HEADERS
\newcommand{\intalkheader}[1]{
    \vspace{0.75em}
    \noindent{\normalsize\bfseries #1}
    \par\nopagebreak % Prevents a break immediately after the text
    \vspace{0.25em}
    \nopagebreak     % Prevents a break after the whitespace
}

\begin{document}
\newpage
\settalk{Strive for Excellence}
\setspeaker{Dallin H. Oaks}
\settalkdate{October 1971}
\talkheading{Strive for Excellence}{Dallin H. Oaks}

My beloved brethren who hold the priesthood of God: Over the past several months I have had occasion to ponder the problems and to reflect upon the ideals of education in The Church of Jesus Christ of Latter-day Saints. As I have studied this subject, my thoughts have been drawn to some lines I read at BYU some twenty years ago. The first of these lines you will think strange as an illustration of the subject of education in the Church. They were written by Thomas Hobbes, the seventeenth century English political philosopher, in his greatest work, The Leviathan.

In describing the nature of man, Hobbes wrote that “the life of man [is] solitary, poor, nasty, brutish, and short.” This is a classic example of the philosophies of man. I am grateful that my education exposed me to that thought and others like it, because my familiarity with these thoughts has helped me to understand the world and its peoples and its problems.

But most of all, I am grateful that my educational program was such that at the time I was exposed to this view of man, I was also being taught these lines:

“Adam fell that men might be; and men are, that they might have joy.” (2 Ne. 2:25.)

“Remember the worth of souls is great in the sight of God.” (D\&C 18:10.)

The worlds were created by the Only Begotten of the Father, “and the inhabitants thereof are begotten sons and daughters unto God.” (D\&C 76:24.)

“For a wise and glorious purpose thou hast placed me here on earth. …” (“O My Father,” Hymns, No. 138.)

“… they who keep their second estate shall have glory added upon their heads for ever and ever.” (Abr. 3:26.)

“Wherefore, as it is written, they are gods, even the sons of God—Wherefore, all things are theirs. …” (D\&C 76:58–59.)

My personal experience converts me to the wisdom of the educational philosophy that joins spiritual with secular instruction. At Brigham Young University and in the other institutions of the Church Educational System, we are concerned with teaching the fundaments of spiritual and secular knowledge and with bringing those teachings into harmony in the lives of men and women in order to prepare them for a balanced and full life of service to God and fellowman.

From this philosophy I distill four thoughts that I offer for the special attention of the young men of the priesthood:

I am grateful for my membership in the Church. I am proud to hold the priesthood of God. No worldly honor or position can add much to the dignity, beauty, and power of the priesthood. I am thankful to my Heavenly Father for the testimony I have of the truth of the gospel. I have measured its requirements by reason and found them satisfying. I have put its precepts into practice and felt their good effects in my life. I have seen the gospel work good in the lives of others. I have observed miraculous things. But these signs follow them that believe. I know that the gospel is true because my Father in heaven has answered my prayers and borne witness to me by the power of the Holy Ghost. I am devoted to the gospel of Jesus Christ. I am loyal to the chosen servants of the Lord, whom I sustain with all my heart.

I bear this testimony to you and ask your prayers and the blessings of our Father in heaven upon every teacher and worker in the Church Educational System, that we may meet our responsibilities to him and to his children.

\newpage
\settalk{Why Do We Serve?}
\setspeaker{Dallin H. Oaks}
\settalkdate{October 1984}
\talkheading{Why Do We Serve?}{Dallin H. Oaks}

My dear brothers and sisters, because it was not appropriate for me to commence my Church service until I had concluded my judicial duties in state government, I did not speak at the April conference where I was sustained. Consequently, this semiannual conference is my first opportunity to speak to the general membership of the Church, to express acceptance of my calling to the Council of the Twelve.

I am thrilled with this calling. Having been “called of God, by prophecy, and by the laying on of hands by those who are in authority” (A of F 1:5), I have gladly forsaken my professional activities to spend the rest of my days in the service of the Lord. I will devote my whole heart, might, mind, and strength to the great trusts placed in me, especially to the responsibilities of a special witness of the name of Jesus Christ in all the world.

Many men and women were called to Church service last April. Eight men were called as General Authorities. Six women were called to the presidencies of the Relief Society and Young Women. Over two hundred men were called to serve as bishops, and over seventeen hundred men and women were called as full-time missionaries. In that same month, tens of thousands of others were called as officers and teachers and other workers in the many Church organizations throughout the world. Those called in April joined millions of others already serving in similar capacities in the restored Church.

As I contemplated my own calling and the callings of millions of others already in service, I was led to consider this question: Why do we serve?

Service is an imperative for those who worship Jesus Christ. To followers who were vying for prominent positions in his kingdom, the Savior taught, “Whosoever will be chief among you, let him be your servant.” (Matt. 20:27.) On a later occasion, he spoke of ministering to the needs of the hungry, the naked, the sick, and the imprisoned. He concluded that teaching with these words: “Verily I say unto you, Inasmuch as ye have done it unto one of the least of these my brethren, ye have done it unto me.” (Matt. 25:40.)

In latter-day revelation the Lord has commanded that we “succor the weak, lift up the hands which hang down, and strengthen the feeble knees.” (D\&C 81:5.) In another section of the Doctrine and Covenants, he instructed us to be “anxiously engaged in a good cause, and do many things of [our] own free will, and bring to pass much righteousness.” (D\&C 58:27.) Holders of the Melchizedek Priesthood receive it upon a covenant to use its powers in the service of others. Indeed, service is a covenant obligation of all members of the Church of Jesus Christ.

Whether our service is to our fellowmen or to God, it is the same. (See Mosiah 2:17.) If we love him, we should keep his commandments and feed his sheep. (See John 21:16–17.)

When we think of service, we usually think of the acts of our hands. But the scriptures teach that the Lord looks to our thoughts as well as to our acts. One of God’s earliest commandments to Israel was that they should love him and “serve him with all your heart and with all your soul.” (Deut. 11:13.) When the prophet Samuel was sent to Bethlehem to choose and anoint one of the sons of Jesse as a new king for Israel, the Lord told him to reject the first son, though he was a man of fine appearance. The Lord explained, “Look not on his countenance, or on the height of his stature; because I have refused him: for the Lord seeth not as man seeth; for man looketh on the outward appearance, but the Lord looketh on the heart.” (1 Sam. 16:7.)

We are familiar with the proverb which states that as a man “thinketh in his heart, so is he.” (Prov. 23:7.) We also read in Proverbs: “All the ways of a man are clean in his own eyes; but the Lord weigheth the spirits.” (Prov. 16:2.)

Latter-day revelation declares that the Lord requires not only the acts of the children of men, but “the Lord requireth the heart and a willing mind.” (D\&C 64:34.)

Numerous scriptures teach that our Heavenly Father knows our thoughts and the intents of our heart. (See D\&C 6:16; Mosiah 24:12; Alma 18:32.) The prophet Moroni taught that if our works are to be credited for good, they must be done for the right reasons. If a man “offereth a gift, or prayeth unto God, except he shall do it with real intent it profiteth him nothing.

“For behold, it is not counted unto him for righteousness.” (Moro. 7:6–7.)

Similarly, the prophet Alma taught that if we have hardened our hearts against the word of God, we will “not dare to look up to our God” at the final judgment because “all our works will condemn us; … and our thoughts will also condemn us.” (Alma 12:14.)

These scriptures make clear that in order to purify our service in the Church and to our fellowmen, it is necessary to consider not only how we serve, but also why we serve.

People serve one another for different reasons, and some reasons are better than others. Perhaps none of us serves in every capacity all the time for only a single reason. Since we are imperfect beings, most of us probably serve for a combination of reasons, and the combinations may be different from time to time as we grow spiritually. But we should all strive to serve for the reasons that are highest and best.

What are some of the reasons for service? By way of illustration, and without pretending to be exhaustive, I will suggest six reasons. I will discuss these in ascending order from the lesser to the greater reasons for service.

Some may serve for hope of earthly reward. Such a man or woman might serve in Church positions or in private acts of mercy in an effort to achieve prominence or cultivate contacts that would increase income or aid in acquiring wealth. Others might serve in order to obtain worldly honors, prominence, or power.

The scriptures have a word for gospel service “for the sake of riches and honor”; it is “priestcraft.” (Alma 1:16.) Nephi said, “Priestcrafts are that men preach and set themselves up for a light unto the world, that they may get gain and praise of the world; but they seek not the welfare of Zion.” (2 Ne. 26:29.) In these latter days, we are commanded to “seek to bring forth and establish the cause of Zion.” (D\&C 6:6.) Unfortunately, not all who accomplish works under that heading are really intending to build up Zion or strengthen the faith of the people of God. Other motives can be at work.

Service that is ostensibly unselfish but is really for the sake of riches or honor surely comes within the Savior’s condemnation of those who “outwardly appear righteous unto men, but within … are full of hypocrisy and iniquity.” (Matt. 23:28.) Such service earns no gospel reward.

“I would that ye should do alms unto the poor,” the Savior declared, “but take heed that ye do not your alms before men to be seen of them; otherwise ye have no reward of your Father who is in heaven.” (3 Ne. 13:1; see also Matt. 6:1–2.) The Savior continued:

“Therefore, when ye shall do your alms do not sound a trumpet before you, as will hypocrites do in the synagogues and in the streets, that they may have glory of men. Verily I say unto you, they have their reward.” (3 Ne. 13:2; see also Matt. 6:2.)

In contrast, those who serve quietly, even “in secret,” qualify for the Savior’s promise that “thy Father, who seeth in secret, shall reward thee openly.” (3 Ne. 13:18; see also Matt. 6:4.)

Another reason for service—probably more worthy than the first, but still in the category of service in search of earthly reward—is that motivated by a personal desire to obtain good companionship. We surely have good associations in our Church service, but is that why we serve?

I once knew a person who was active in Church service until a socially prominent friend and fellow worker moved away. When the friend moved from the ward, this person ceased to serve. In this case, a Church worker was willing to serve only when the fellow workers were acceptable.

Persons who serve only to obtain good companionship are more selective in choosing their friends than the Master was in choosing his servants or associates. Jesus called most of his servants from those in humble circumstances. And he associated with sinners. He answered critics of such association by saying, “They that are whole need not a physician; but they that are sick. I came not to call the righteous, but sinners to repentance.” (Luke 5:31–32.)

The first section of the Doctrine and Covenants, which speaks of people in the last days, gives a description that seems to include those who serve for hope of earthly reward of one sort or another: “They seek not the Lord to establish his righteousness, but every man walketh in his own way, and after the image of his own god, whose image is in the likeness of the world, and whose substance is that of an idol.” (D\&C 1:16.)

These first two reasons for service are selfish and self-centered and unworthy of Saints. As the Apostle Paul said, we that are strong enough to bear the infirmities of the weak should not do so “to please ourselves.” (Rom. 15:1.) Reasons aimed at earthly rewards are distinctly lesser in character and reward than the other reasons I will discuss.

Some may serve out of fear of punishment. The scriptures abound with descriptions of the miserable state of those who fail to follow the commandments of God. Thus, King Benjamin taught his people that the soul of the unrepentant transgressor would be filled with “a lively sense of his own guilt, which doth cause him to shrink from the presence of the Lord, and doth fill his breast with guilt, and pain, and anguish, which is like an unquenchable fire, whose flame ascendeth up forever and ever.” (Mosiah 2:38.) Such descriptions surely offer sufficient incentive for keeping the commandment of service. But service out of fear of punishment is a lesser motive at best.

Other persons may serve out of a sense of duty or out of loyalty to friends or family or traditions. These are those I would call the good soldiers, who instinctively do what they are asked without question and sometimes without giving much thought to the reasons for their service. Such persons fill the ranks of voluntary organizations everywhere, and they do much good. We have all benefited by the good works of such persons. Those who serve out of a sense of duty or loyalty to various wholesome causes are the good and honorable men and women of the earth.

Service of the character I have just described is worthy of praise and will surely qualify for blessings, especially if it is done willingly and joyfully. As the Apostle Paul wrote in his second letter to the Corinthians:

“But this I say, He which soweth sparingly shall reap also sparingly; and he which soweth bountifully shall reap also bountifully.

“Every man according as he purposeth in his heart, so let him give; not grudgingly, or of necessity: for God loveth a cheerful giver.” (2 Cor. 9:6–7.)

“It is obeying God willingly that is accepted,” an anonymous writer has said. “The Lord hates that which is forced—it is rather a tax than an offering.”

Although those who serve out of fear of punishment or out of a sense of duty undoubtedly qualify for the blessings of heaven, there are still higher reasons for service.

One such higher reason for service is the hope of an eternal reward. This hope—the expectation of enjoying the fruits of our labors—is one of the most powerful sources of motivation. As a reason for service, it necessarily involves faith in God and in the fulfillment of his prophecies. The scriptures are rich in promises of eternal rewards. For example, in a revelation given through the Prophet Joseph Smith in June 1829, the Lord said: “If you keep my commandments and endure to the end you shall have eternal life, which gift is the greatest of all the gifts of God.” (D\&C 14:7.)

The last motive I will discuss is, in my opinion, the highest reason of all. In its relationship to service, it is what the scriptures call “a more excellent way.” (1 Cor. 12:31.)

“Charity is the pure love of Christ.” (Moro. 7:47.) The Book of Mormon teaches us that this virtue is “the greatest of all.” (Moro. 7:46.) The Apostle Paul affirmed and illustrated that truth in his great teaching about the reasons for service:

“Though I speak with the tongues of men and of angels, and have not charity, I am become as sounding brass, or a tinkling cymbal. …

“And though I bestow all my goods to feed the poor, … and have not charity, it profiteth me nothing.” (1 Cor. 13:1–3.)

We know from these inspired words that even the most extreme acts of service—such as giving all of our goods to feed the poor—profit us nothing unless our service is motivated by the pure love of Christ.

If our service is to be most efficacious, it must be accomplished for the love of God and the love of his children. The Savior applied that principle in the Sermon on the Mount, in which he commanded us to love our enemies, bless them that curse us, do good to them that hate us, and pray for them that despitefully use us and persecute us. (See Matt. 5:44.) He explained the purpose of that commandment as follows:

“For if ye love them which love you, what reward have ye? do not even the publicans the same?

“And if ye salute your brethren only, what do ye more than others? do not even the publicans so?” (Matt. 5:46–47.)

This principle—that our service should be for the love of God and the love of fellowmen rather than for personal advantage or any other lesser motive—is admittedly a high standard. The Savior must have seen it so, since he joined his commandment for selfless and complete love directly with the ideal of perfection. The very next verse of the Sermon on the Mount contains this great commandment: “Be ye therefore perfect, even as your Father which is in heaven is perfect.” (Matt. 5:48.)

This principle of service is reaffirmed in the fourth section of the Doctrine and Covenants:

“Therefore, O ye that embark in the service of God, see that ye serve him with all your heart, might, mind and strength, that ye may stand blameless before God at the last day.” (D\&C 4:2.)

We learn from this command that it is not enough to serve God with all of our might and strength. He who looks into our hearts and knows our minds demands more than this. In order to stand blameless before God at the last day, we must also serve him with all our heart and mind.

Service with all of our heart and mind is a high challenge for all of us. Such service must be free of selfish ambition. It must be motivated only by the pure love of Christ.

If we have difficulty with the command that we serve for love, a Book of Mormon teaching can help us. After describing the importance of charity, the prophet Moroni counseled:

“Wherefore, my beloved brethren, pray unto the Father with all the energy of heart, that ye may be filled with this love, which he hath bestowed upon all who are true followers of his Son, Jesus Christ.” (Moro. 7:48.)

The service of persons filled with that love will meet the high test expressed in the Twenty-fourth Psalm:

“Who shall ascend into the hill of the Lord? or who shall stand in his holy place?

“He that hath clean hands, and a pure heart.” (Ps. 24:3–4.)

I know that God expects us to work to purify our hearts and our thoughts so that we may serve one another for the highest and best reason, the pure love of Christ.

Most of all, I know that God lives, and I know that his Only Begotten Son, Jesus Christ, died for our sins and is our Savior. And I know that God has restored the fulness of the gospel through the Prophet Joseph Smith in these latter days. In the name of Jesus Christ, amen.

\newpage
\settalk{Taking upon Us the Name of Jesus Christ}
\setspeaker{Dallin H. Oaks}
\settalkdate{April 1985}
\talkheading{Taking upon Us the Name of Jesus Christ}{Dallin H. Oaks}

On this Easter Sunday we rejoice in the resurrection of our Lord and Savior, Jesus Christ. This is a holy day for all of Christianity. Countless Christians attend worship services on this day to partake of the sacrament of the Lord’s Supper, which many call communion. (See 1 Cor. 10:16.)

Members of The Church of Jesus Christ of Latter-day Saints are commanded to partake of the sacrament each week. (See D\&C 59:9, 12.) In doing so, they witness unto God the Eternal Father, as stated in the prayer on the bread, that they are “willing to take upon them the name of thy Son, and always remember him and keep his commandments which he has given them.” (D\&C 20:77; Moro. 4:3.) We should ponder these sacred covenants during the sacrament service.

On this Easter Sunday it is appropriate to reflect on what it means to partake of the sacrament. I will focus on the first of these solemn “witnesses” to God the Eternal Father: that we are willing to take upon us the name of his Son. What does this mean?

Our witness that we are willing to take upon us the name of Jesus Christ has several different meanings. Some of these meanings are obvious, and well within the understanding of our children. Others are only evident to those who have searched the scriptures and pondered the wonders of eternal life.

One of the obvious meanings renews a promise we made when we were baptized. Following the scriptural pattern, persons who are baptized “witness before the Church that they have truly repented of all their sins, and are willing to take upon them the name of Jesus Christ, having a determination to serve him to the end.” (D\&C 20:37; see also 2 Ne. 31:13; Moro. 6:3.) When we partake of the sacrament, we renew this covenant and all the other covenants we made in the waters of baptism. (See Joseph Fielding Smith, Doctrines of Salvation, comp. Bruce R. McConkie, 3 vols., Salt Lake City: Bookcraft, 1954–56, 2:341, 346.)

As a second obvious meaning, we take upon us our Savior’s name when we become members of The Church of Jesus Christ of Latter-day Saints. By his commandment, this church bears his name. (See D\&C 115:4; 3 Ne. 27:7–8.) Every member, young and old, is a member of the “household of God.” (Eph. 2:19.) As true believers in Christ, as Christians, we have gladly taken his name upon us. (See Alma 46:15.) As King Benjamin taught his people, “Because of the covenant which ye have made ye shall be called the children of Christ, his sons, and his daughters; for behold, this day he hath spiritually begotten you.” (Mosiah 5:7; see also Alma 5:14; 36:23–26.)

We also take upon us the name of Jesus Christ whenever we publicly proclaim our belief in him. Each of us has many opportunities to proclaim our belief to friends and neighbors, fellow workers, and casual acquaintances. As the Apostle Peter taught the Saints of his day, we should “sanctify the Lord God in [our] hearts: and be ready always to give an answer to every man that asketh [us] a reason of the hope that is in [us].” (1 Pet. 3:15.) In this, we keep the modern commandment: “Take upon you the name of Christ, and speak the truth in soberness.” (D\&C 18:21.)

A third meaning appeals to the understanding of those mature enough to know that a follower of Christ is obligated to serve him. Many scriptural references to the name of the Lord seem to be references to the work of his kingdom. Thus, when Peter and the other Apostles were beaten, they rejoiced “that they were counted worthy to suffer shame for his name.” (Acts 5:41.) Paul wrote certain members who had ministered to the Saints that the Lord would not forget the labor of love they had “shewed toward his name.” (Heb. 6:10.) According to this meaning, by witnessing our willingness to take upon us the name of Jesus Christ, we signify our willingness to do the work of his kingdom.

In these three relatively obvious meanings, we see that we take upon us the name of Christ when we are baptized in his name, when we belong to his Church and profess our belief in him, and when we do the work of his kingdom.

There are other meanings as well, deeper meanings that the more mature members of the Church should understand and ponder as he or she partakes of the sacrament.

It is significant that when we partake of the sacrament we do not witness that we take upon us the name of Jesus Christ. We witness that we are willing to do so. (See D\&C 20:77.) The fact that we only witness to our willingness suggests that something else must happen before we actually take that sacred name upon us in the most important sense.

What future event or events could this covenant contemplate? The scriptures suggest two sacred possibilities, one concerning the authority of God, especially as exercised in the temples, and the other—closely related—concerning exaltation in the celestial kingdom.

The name of God is sacred. The Lord’s Prayer begins with the words, “Our Father which art in heaven, Hallowed be thy name.” (Matt. 6:9.) From Sinai came the commandment, “Thou shalt not take the name of the Lord thy God in vain.” (Ex. 20:7, Deut. 5:11.) Latter-day revelation equates this with using the name of God without authority. “Let all men beware how they take my name in their lips,” the Lord declares in a modern revelation, for “many there be who … use the name of the Lord, and use it in vain, having not authority.” (D\&C 63:61–62.)

Consistent with these references, many scriptures that refer to “the name of Jesus Christ” are obviously references to the authority of the Savior. This was surely the meaning conveyed when the seventy reported to Jesus that “even the devils are subject unto us through thy name.” (Luke 10:17.) The Doctrine and Covenants employs this same meaning when it describes the Twelve Apostles of this dispensation as “they who shall desire to take upon them my name with full purpose of heart.” (D\&C 18:27.) The Twelve are later designated as “special witnesses of the name of Christ in all the world,” and as those who “officiate in the name of the Lord, under the direction of the Presidency of the Church.” (D\&C 107:23, 33.)

By way of further illustration, the Old Testament contains scores of references to the name of the Lord in a context where it clearly means the authority of the Lord. Most of these references have to do with the temple.

When the children of Israel were still on the other side of the Jordan, the Lord told them that when they entered the promised land there should be a place where the Lord their God would “cause his name to dwell.” (Deut. 12:11; see also Deut. 14:23–24; 16:6.) Time after time in succeeding revelations, the Lord and his servants referred to the future temple as a house for “the name” of the Lord God of Israel. (See 1 Kgs. 3:2; 5:5; 8:16–20, 29, 44, 48; 1 Chr. 22:8–10, 19; 29:16; 2 Chr. 2:4; 6:5–10, 20, 34, 38.) After the temple was dedicated, the Lord appeared to Solomon and told him that He had hallowed the temple “to put my name there for ever.” (1 Kgs. 9:3; 2 Chr. 7:16.)

Similarly, in modern revelations the Lord refers to temples as houses built “unto my holy name.” (D\&C 124:39; 105:33; 109:2–5.) In the inspired dedicatory prayer of the Kirtland Temple, the Prophet Joseph Smith asked the Lord for a blessing upon “thy people upon whom thy name shall be put in this house.” (D\&C 109:26.)

All of these references to ancient and modern temples as houses for “the name” of the Lord obviously involve something far more significant than a mere inscription of his sacred name on the structure. The scriptures speak of the Lord’s putting his name in a temple because he gives authority for his name to be used in the sacred ordinances of that house. That is the meaning of the Prophet’s reference to the Lord’s putting his name upon his people in that holy house. (See D\&C 109:26.)

Willingness to take upon us the name of Jesus Christ can therefore be understood as willingness to take upon us the authority of Jesus Christ. According to this meaning, by partaking of the sacrament we witness our willingness to participate in the sacred ordinances of the temple and to receive the highest blessings available through the name and by the authority of the Savior when he chooses to confer them upon us.

Another future event we may anticipate when we witness our willingness to take that sacred name upon us concerns our relationship to our Savior and the incomprehensible blessings available to those who will be called by his name at the last day.

King Benjamin told his people, “There shall be no other name given nor any other way nor means whereby salvation can come unto the children of men, only in and through the name of Christ, the Lord Omnipotent.” (Mosiah 3:17; see also 2 Ne. 31:21.) Peter proclaimed “the name of Jesus Christ of Nazareth” to the leaders of the Jews, declaring that “there is none other name under heaven given among men, whereby we must be saved.” (Acts 4:10, 12; see also D\&C 18:21.)

The scriptures proclaim that the Savior’s atoning sacrifice was for those who “believe on his name.” Alma taught that Jesus Christ, the Son, the Only Begotten of the Father, would come “to take away the sins of the world, yea, the sins of every man who steadfastly believeth on his name.” (Alma 5:48; 9:27; 11:40; Hel. 14:2.) In the words of King Benjamin, “Whosoever doeth this shall be found at the right hand of God, for he shall know the name by which he is called; for he shall be called by the name of Christ.” (Mosiah 5:9.)

Thus, those who exercise faith in the sacred name of Jesus Christ and repent of their sins and enter into his covenant and keep his commandments (see Mosiah 5:8) can lay claim on the atoning sacrifice of Jesus Christ. Those who do so will be called by his name at the last day.

When the Savior taught the Nephites following his resurrection, he referred to the scriptural statement that “ye must take upon you the name of Christ.” He explained, “For by this name shall ye be called at the last day; And whoso taketh upon him my name, and endureth to the end, the same shall be saved at the last day.” (3 Ne. 27:5–6.) That same teaching is repeated in a modern revelation, which adds the caution that “if they know not the name by which they are called, they cannot have place in the kingdom of my Father.” (D\&C 18:25; see also Alma 5:38.)

The Book of Mormon explains the significance of being called by the name of Jesus Christ. When the Savior showed his spirit body to the brother of Jared, he introduced himself as the Father and the Son, declaring that through his redeeming sacrifice all mankind who believed on his name should have life eternal through him, “and they shall become my sons and my daughters.” (Ether 3:14.) Abinadi said of those who believed in the Lord and looked to him for a remission of their sins “that these are his seed, or they are heirs of the kingdom of God.” (Mosiah 15:11.) He continued this explanation as follows:

“For these are they whose sins he has borne; these are they for whom he has died, to redeem them from their transgressions. And now, are they not his seed?” (Mosiah 15:12.)

Speaking through the prophet Alma, the Lord explained the significance of this relationship: “For behold, in my name are they called; and if they know me they shall come forth, and shall have a place eternally at my right hand.” (Mosiah 26:24.)

In these great scriptures from the Book of Mormon, we learn that those who are qualified by faith and repentance and compliance with the laws and ordinances of the gospel will have their sins borne by the Lord Jesus Christ. In spiritual and figurative terms they will become the sons and daughters of Christ, heirs to his kingdom. These are they who will be called by his name in the last day.

According to this meaning, when we witness our willingness to take upon us the name of Jesus Christ, we are signifying our commitment to do all that we can to achieve eternal life in the kingdom of our Father. We are expressing our candidacy—our determination to strive for—exaltation in the celestial kingdom.

Those who are found worthy to take upon them the name of Jesus Christ at the last day are described in the great revelations recorded in the ninety-third and seventy-sixth sections of the Doctrine and Covenants. Here the Savior revealed to Joseph Smith that in due time, if we keep the commandments of God, we can receive the “fulness” of the Father. (D\&C 93:19–20.) Here the Savior bears record that “all those who are begotten through me are partakers of the glory of the [Father], and are the church of the Firstborn.” (D\&C 93:22.) “They are they into whose hands the Father has given all things. … Wherefore, as it is written, they are gods” who “shall dwell in the presence of God and his Christ forever and ever.” (D\&C 76:55, 58, 62.) “And this is life eternal, that they might know thee the only true God, and Jesus Christ, whom thou hast sent.” (John 17:3; see also D\&C 88:4–5.) This is the ultimate significance of taking upon us the name of Jesus Christ.

When the priest offers the scriptural prayer on the bread at the sacrament table, he prays that all who partake may “witness” unto God, the Eternal Father, “that they are willing to take upon them the name of thy Son.” (D\&C 20:77; Moro. 4:3.) This witness has several different meanings.

It causes us to renew the covenant we made in the waters of baptism to take upon us the name of Jesus Christ and serve him to the end. We also take upon us his name as we publicly profess our belief in him, as we fulfill our obligations as members of his Church, and as we do the work of his kingdom.

But there is something beyond these familiar meanings, because what we witness is not that we take upon us his name but that we are willing to do so. In this sense, our witness relates to some future event or status whose attainment is not self-assumed, but depends on the authority or initiative of the Savior himself.

Scriptural references to the name of Jesus Christ often signify the authority of Jesus Christ. In that sense, our willingness to take upon us his name signifies our willingness to take upon us the authority of Jesus Christ in the sacred ordinances of the temple, and to receive the highest blessings available through his authority when he chooses to confer them upon us.

Finally, our willingness to take upon us the name of Jesus Christ affirms our commitment to do all that we can to be counted among those whom he will choose to stand at his right hand and be called by his name at the last day. In this sacred sense, our witness that we are willing to take upon us the name of Jesus Christ constitutes our declaration of candidacy for exaltation in the celestial kingdom. Exaltation is eternal life, “the greatest of all the gifts of God.” (D\&C 14:7.)

That is what we should ponder as we partake of the sacred emblems of the sacrament. As we do so, we glory in the mission of the risen Lord, who lived and taught and suffered and died and rose again that all mankind might have immortality and eternal life. Of this I testify in the sacred name of the Lord Jesus Christ, whose witness I am, amen.

\newpage
\settalk{Spirituality}
\setspeaker{Dallin H. Oaks}
\settalkdate{October 1985}
\talkheading{Spirituality}{Dallin H. Oaks}

As faithful members of The Church of Jesus Christ of Latter-day Saints, we have a distinctive way of looking at life. We view our experiences in terms of eternity. As we draw farther from worldliness, we feel closer to our Father in Heaven and more able to be guided by his Spirit. We call this quality of life spirituality.

To the faithful, spirituality is a lens through which we view life and a gauge by which we evaluate it. The Apostle Paul expressed this thought in two of his letters:

“We look not at the things which are seen, but at the things which are not seen: for the things which are seen are temporal; but the things which are not seen are eternal.” (2 Cor. 4:18.)

“For they that are after the flesh do mind the things of the flesh; but they that are after the Spirit the things of the Spirit.

“For to be carnally minded is death; but to be spiritually minded is life and peace.” (Rom. 8:5–6.)

To be spiritually minded is to view and evaluate our experiences in terms of the enlarged perspective of eternity.

Each of us has a personal lens through which we view the world. Our lens gives its special tint to all we see. It can suppress some features and emphasize others. It can also reveal things otherwise invisible. Through the lens of spirituality, we can know “the things of God” by “the Spirit of God.” (1 Cor. 2:11.) As the Apostle Paul taught, such things are “foolishness” to the “natural man.” He cannot see them “because they are spiritually discerned.” (See 1 Cor. 2:14.)

How we interpret our experiences is also a function of our degree of spirituality. Some interpret mortality solely in terms of worldly accomplishments and possessions. In contrast, we who have a testimony of the gospel of Jesus Christ should interpret our experiences in terms of our knowledge of the purpose of life, the mission of our Savior, and the eternal destiny of the children of God.

Spirituality is not a function of occupation or calling. A scientist may be more spiritual than a theologian; a teacher may be more spiritual than an officer. Spirituality is determined by personal outlook and priorities. It is evident in our words and actions. Elder John Taylor showed his spirituality in these words, uttered as he reported on his mission to Europe in 1852:

“Some people have said to me, sometimes, Are you not afraid to cross over the seas, and deserts, where there are wolves and bears, and other ferocious animals? … Are you not afraid that you will drop by the way, and leave your body on the desert track, or beneath the ocean’s wave? No. Who cares anything about it? What of it, if we should happen to drop by the way? … These things don’t trouble me, but I have felt to rejoice all the day long, that God has revealed the principle of eternal life, that I am put in possession of that truth, and that I am counted worthy to engage in the work of the Lord.” (Journal of Discourses, 1:17.)

The scriptures contain great illustrations of spirituality as it relates to everyday living. One of these, recorded in the tenth chapter of Luke, tells how the Savior came to a particular village:

“And a certain woman named Martha received him into her house.

“And she had a sister called Mary, which also sat at Jesus’ feet, and heard his word.

“But Martha was cumbered about much serving, and came to him, and said, Lord, dost thou not care that my sister hath left me to serve alone? bid her therefore that she help me.

“And Jesus answered and said unto her, Martha, Martha, thou art careful and troubled about many things:

“But one thing is needful: and Mary hath chosen that good part, which shall not be taken away from her.” (Luke 10:38–42.)

This scripture reminds every Martha, male and female, that we should not be so occupied with what is routine and temporal that we fail to cherish those opportunities that are unique and spiritual.

The contrast between the spiritual and the temporal is also illustrated by the twins Esau and Jacob and their different attitudes toward their birthright. The firstborn, Esau, “despised his birthright.” (Gen. 25:34.) Jacob, the second twin, desired it. Jacob valued the spiritual, while Esau sought the things of this world. When he was hungry, Esau sold his birthright for a mess of pottage. “Behold,” he explained, “I am at the point to die: and what profit shall this birthright do to me?” (Gen. 25:32.) Many Esaus have given up something of eternal value in order to satisfy a momentary hunger for the things of the world.

The Roman soldiers of Pilate provided an unforgettable illustration of the different perspectives of the carnal mind and the spiritual mind. During a tragic afternoon on Calvary, a handful of soldiers waited at the foot of a cross. The most important event in all eternity was taking place on the cross above their heads. Oblivious to that fact, they occupied themselves casting lots to divide the earthly property of the dying Son of God. (See Matt. 27:35; Luke 23:34; John 19:24.) Their example reminds each of us that we should not be casting our lots for the things of the world while the things of eternity, including our families and the work of the Lord, suffer for our attention.

Here is an example of a spiritual and a temporal evaluation of an everyday experience. In a BYU devotional several years ago, Elder Loren C. Dunn described how his father, a busy stake president in Tooele, gave his two young sons the responsibility of raising cows on the family farm. He gave the boys large latitude in what they could do, and they made some mistakes. These were observed by an alert neighbor, who complained to their father about what the young cow-raisers were doing. “Jim, you don’t understand,” President Dunn replied. “You see, I’m raising boys, not cows.” (“Our Spiritual Heritage,” in Brigham Young University 1981–82 Fireside and Devotional Speeches, Provo: Brigham Young University Press, 1982, p. 138.) What a marvelous insight! What an example for parents who are inclined to view and evaluate their children’s performance solely in temporal terms.

What we see around us depends on what we seek in life. The Spanish conquerors took irreplaceable objects of art from the craftsmen of the New World and melted them down into gold bullion. The enemies of the young prophet, Joseph Smith, hounded him to get possession of the golden plates from which he was to translate the Book of Mormon. They sought the golden plates to get money, not a message. The temporal value of the plates had a price; their spiritual value was priceless.

Elder John A. Widtsoe taught that “there is a spiritual meaning of all human acts and earthly events. … It is the business of man to find the spiritual meaning of earthly things. … No man is quite so happy … as he who backs all his labors by such a spiritual interpretation and understanding of the acts of his life.” (In Conference Report, Apr. 1922, pp. 96–97.)

The Latter-day Saint men and women who settled these valleys of the mountains acted upon that principle. Judged in terms of the values and aspirations of the world, some pioneer enterprises were failures. The iron mission did not succeed in making significant quantities of iron. The cotton mission did not give the Utah Territory self-sufficiency in cotton production. Efforts to manufacture sugar did not achieve material success for forty years. The Perpetual Immigration Fund did not perpetuate itself because many immigrants were unable to pay their debt to it.

But, when measured against the eternal values of loyalty, cooperation, and consecration, some of the most conspicuous worldly failures are seen as the pioneer enterprisers’ greatest triumphs. Whatever their financial outcome, these enterprises called forth the sacrifices that molded pioneers into Saints and prepared Saints for exaltation. Unto God, “all things … are spiritual.” (D\&C 29:34.)

In another great event in Mormon history, several hundred men marched from Ohio to give military relief to the persecuted Saints in Zion—western Missouri. But when the men of Zion’s Camp approached their intended destination, the Prophet Joseph Smith disbanded them. According to its ostensible purpose, the expedition was a failure. But most of the men who were to lead the Church for the next half-century, including those who would take the Saints across the plains and colonize the Intermountain West, came to know the Prophet Joseph and received their formative leadership training in the march of Zion’s Camp. As Elder Orson F. Whitney said of Zion’s Camp:

“The redemption of Zion is more than the purchase or recovery of lands, the building of cities, or even the founding of nations. It is the conquest of the heart, the subjugation of the soul, the sanctifying of the flesh, the purifying and ennobling of the passions.” (The Life of Heber C. Kimball, 2d ed., Salt Lake City: Stevens \& Wallis, 1945, p. 65.)

The first of the Ten Commandments—“Thou shalt have no other gods before me” (Ex. 20:3)—epitomizes the nature of spirituality. A spiritual person has no priorities ahead of God. A person who seeks or serves other objectives, such as power or prominence, is not spiritual.

The primacy of the spiritual over the temporal is evident in the teachings of the Savior’s three senior Apostles. Peter taught:

“All flesh is as grass, and all the glory of man as the flower of grass. The grass withereth, and the flower thereof falleth away:

“But the word of the Lord endureth for ever.” (1 Pet. 1:24–25.)

James asked: “Know ye not that the friendship of the world is enmity with God? whosoever therefore will be a friend of the world is the enemy of God.” (James 4:4.)

And the Apostle John wrote:

“Love not the world, neither the things that are in the world. If any man love the world, the love of the Father is not in him.

“For all that is in the world, the lust of the flesh, and the lust of the eyes, and the pride of life, is not of the Father, but is of the world.

“And the world passeth away, and the lust thereof: but he that doeth the will of God abideth for ever.” (1 Jn. 2:15–17.)

Materialism, which gives priority to material needs and objects, is obviously the opposite of spirituality. The Savior taught that we should not lay up “treasures upon earth, where moth and rust doth corrupt, and where thieves break through and steal.” (Matt. 6:19.) We should lay up treasures in heaven: “For where your treasure is, there will your heart be also.” (Matt. 6:21.)

Like other Book of Mormon prophets, Samuel the Lamanite warned the Nephites that they were cursed because of their riches, “Because ye have set your hearts upon them, and have not hearkened unto the words of him who gave them unto you.” (Hel. 13:21; see also Hel. 6:17; 7:21.)

The Apostle Paul counseled young Timothy, “Charge them that are rich in this world, that they be not highminded, nor trust in uncertain riches, but in the living God, who giveth us richly all things to enjoy.” (1 Tim. 6:17.)

There is nothing inherently evil about money. The Good Samaritan used the same coinage to serve his fellowman that Judas used to betray the Master. It is “the love of money [which] is the root of all evil.” (1 Tim. 6:10; italics added.) The critical difference is the degree of spirituality we exercise in viewing, evaluating, and managing the things of this world and our experiences in it.

If allowed to become an object of worship or priority, money can make us selfish and prideful, “puffed up in the vain things of the world.” (Alma 5:37.) In contrast, if used for fulfilling our legal obligations and for paying our tithes and offerings, money can demonstrate integrity and develop unselfishness. The spiritually enlightened use of property can help prepare us for the higher law of a celestial glory.

The qualities of spirituality we have been able to embody in our lives are often evident in the way we react to death or other apparent tragedies or misfortunes. As faithful Latter-day Saints, we can bear the death of loved ones because we have faith in the resurrection and the eternal nature of family ties. We can repent and rise above our mistakes and inadequacies because we know that our Savior “suffered these things for all, that they might not suffer if they would repent.” (D\&C 19:16.)

Seen with the perspective of eternity, a temporal setback can be an opportunity to develop soul power of eternal significance. Strength is forged in adversity. Faith is developed in a setting where we cannot see what lies ahead.

Lehi promised his son Jacob that God “shall consecrate thine afflictions for thy gain.” (2 Ne. 2:2.) In the midst of the Missouri persecutions, the Lord assured the Saints that “all things wherewith you have been afflicted shall work together for your good.” (D\&C 98:3.) Those who can look upon their afflictions in this manner have spirituality.

How do we achieve spirituality? How do we attain that degree of holiness where we can have the constant companionship of the Holy Ghost? How do we come to view and evaluate the things of this world with the perspective of eternity?

We seek spirituality through faith, repentance, and baptism; through forgiveness of one another; through fasting and prayer; through righteous desires and pure thoughts and actions. We seek spirituality through service to our fellowmen; through worship; through feasting on the word of God, in the scriptures and the teachings of the living prophets. We attain spirituality through making and keeping covenants, through conscientiously trying to keep all the commandments of God. Spirituality is not acquired suddenly. It is the consequence of a succession of right choices. It is the harvest of a righteous life.

Through the lens of spirituality, we see all the commandments of God as invitations to blessings. Obedience and sacrifice, loyalty and love, fidelity and family all appear in eternal perspective. The words of the Savior, given to the world in the Prophet Joseph Smith’s inspired translation of the Bible, have renewed significance:

“And whosoever will lose his life in this world, for my sake, shall find it in the world to come.

“Therefore, forsake the world, and save your souls; for what is a man profited, if he shall gain the whole world, and lose his own soul? Or what shall a man give in exchange for his soul?” (JST, Matt. 16:28–29.)

The fruits of spirituality were revealed to the Prophet Joseph Smith in the eighty-eighth section of the Doctrine and Covenants:

“And if your eye be single to my glory, your whole bodies shall be filled with light, and there shall be no darkness in you; and that body which is filled with light comprehendeth all things.

“Therefore, sanctify yourselves that your minds become single to God.” (D\&C 88:67–68.)

I testify that we are the children of God the Eternal Father. Through the atoning sacrifice of his Only Begotten Son, our Savior Jesus Christ, he has given us the means by which we may be cleansed of our sins. Through his prophets he has given us the eternal perspective of spirituality.

May we strive to attain that level of spirituality where we sanctify ourselves that our minds become single to God. (See D\&C 88:68.) As we do so, we will enjoy his promised blessings, including the blessing of eternal life, “the greatest of all the gifts of God.” (D\&C 14:7.) Of this I testify in the name of Jesus Christ, amen.

\newpage
\settalk{Reverent and Clean}
\setspeaker{Dallin H. Oaks}
\settalkdate{April 1986}
\talkheading{Reverent and Clean}{Dallin H. Oaks}

Recently our family was viewing what was supposed to be a wholesome movie on videotape. Suddenly, one of the actors used a vulgar expression. Embarrassed, we began to smooth this over for our ten-year-old daughter. She quickly assured us that we needn’t worry because she heard worse than that every day from the boys and girls at her school.

I am sure most LDS parents have had similar experiences. The nature and extent of profanity and vulgarity in our society is a measure of its deterioration.

I cannot remember when I first heard profane and vulgar expressions in common use around me. I suppose it was from adults in the barnyard or the barracks. Today, our young people hear such expressions from boys and girls in their grade schools, from actors on stage and in the movies, from popular novels, and even from public officials and sports heroes. Television and videotapes bring profanity and vulgarity into our homes.

For many in our day, the profane has become commonplace and the vulgar has become acceptable. Surely this is one fulfillment of the Book of Mormon prophecy that in the last days “there shall be great pollutions upon the face of the earth.” (Morm. 8:31.)

The people of God have always been commanded to abstain from language that is profane or vulgar. Latter-day Saints should understand why.

The names of God the Father and his Son, Jesus Christ, are sacred. The prophet Isaiah taught that the Lord will not suffer these names to be dishonored—“polluted” as the scriptures say. (See 1 Ne. 20:11; Isa. 48:11.)

In the third of the Ten Commandments, the Lord commanded ancient Israel: “Thou shalt not take the name of the Lord thy God in vain; for the Lord will not hold him guiltless that taketh his name in vain.” (Ex. 20:7.) This same commandment was repeated to the Book of Mormon people by the prophet Abinadi (see Mosiah 13:15) and to each of us through modern prophets. (See D\&C 136:21.)

The Doctrine and Covenants gives this example:

“Let all men beware how they take my name in their lips—

“For behold, verily I say, that many there be who are under this condemnation, who use the name of the Lord, and use it in vain, having not authority.” (D\&C 63:61–62.)

This scripture shows that we take the name of the Lord in vain when we use his name without authority. This obviously occurs when the sacred names of God the Father and his Son, Jesus Christ, are used in what is called profanity: in hateful cursings, in angry denunciations, or as marks of punctuation in common discourse.

The names of the Father and the Son are used with authority when we reverently teach and testify of them, when we pray, and when we perform the sacred ordinances of the priesthood.

There are no more sacred or significant words in all of our language than the names of God the Father and his Son, Jesus Christ.

As we read in the Book of Mormon, after the Savior appeared to the people on this continent he taught them that they must take upon them the name of Christ:

“For by this name shall ye be called at the last day;

“And whoso taketh upon him my name, and endureth to the end, the same shall be saved at the last day.” (3 Ne. 27:5–6.)

He has instructed his followers to call the Church in his name. (See 3 Ne. 27:7–8; D\&C 115:4.) In our time this is The Church of Jesus Christ of Latter-day Saints.

The Savior taught that we should begin our prayers by saying, “Our Father who art in heaven, hallowed be thy name.” (3 Ne. 13:9; see also Luke 11:2.) In the Book of Mormon, the risen Lord gave these further instructions:

“Therefore ye must always pray unto the Father in my name;

“And whatsoever ye shall ask the Father in my name, which is right, believing that ye shall receive, behold it shall be given unto you.

“Pray in your families unto the Father, always in my name, that your wives and your children may be blessed.” (3 Ne. 18:19–21; see also 3 Ne. 27:7; John 14:13; 15:16.)

The scriptures are replete with declarations that the name of Jesus Christ is “the only name which shall be given under heaven, whereby salvation shall come unto the children of men.” (Moses 6:52; see also Acts 4:12; 2 Ne. 25:20; 31:21; Mosiah 3:17.)

The Bible has hundreds of references to the name of God, a sacred word which usually refers to God the Father, or Elohim. (See Gen.; John 3:16.) The ancient prophets also knew and revered the name of Jehovah, the Holy One of Israel, Jesus Christ, whom the Bible usually refers to as the Lord. (See JST, Ex. 6:3; Abr. 1:16; 2:8; Ether 3:1–28; Isa. 43:3.)

These names were so sacred that the children of Israel were repeatedly commanded not to profane the holy name of their God. (See Lev. 18:21; 19:12; 20:3; 21:6.) One who blasphemed the name of the Lord was commanded to be put to death by stoning. (See Lev. 24:16.)

Cataloging the sins of his countrymen, the prophet Ezekiel said, “Her priests have … profaned mine holy things: they have put no difference between the holy and the profane … and I am profaned among them.” (Ezek. 22:26; see also 36:20–23.)

Throughout the ages, the Lord has directed that “whatsoever ye shall do, ye shall do it in my name.” (3 Ne. 27:7.) God the Father commanded that Adam and Eve and all of their descendants should be baptized “in the name of mine Only Begotten Son, who is full of grace and truth, which is Jesus Christ.” (Moses 6:52.)

At the conclusion of his ministry, the risen Lord identified signs that would follow those who believed (see Mark 16:17–18):

“In my name they shall do many wonderful works;

“In my name they shall cast out devils;

“In my name they shall heal the sick;

“In my name they shall open the eyes of the blind, and unstop the ears of the deaf.” (D\&C 84:66–69.)

When Peter healed the lame beggar, he spoke these words: “Such as I have give I thee: In the name of Jesus Christ of Nazareth rise up and walk.” (Acts 3:6.)

When the names of God the Father and his Son, Jesus Christ, are used with reverence and authority, they invoke a power beyond what mortal man can comprehend.

It should be obvious to every believer that these mighty names—by which miracles are wrought, by which the world was formed, through which man was created, and by which we can be saved—are holy and must be treated with the utmost reverence. As we read in modern revelation, “Remember that that which cometh from above is sacred, and must be spoken with care, and by constraint of the Spirit.” (D\&C 63:64.)

So it is that the Holy Priesthood after the Order of the Son of God is called the Melchizedek Priesthood “out of respect or reverence to the name of the Supreme Being, to avoid the too frequent repetition of his name.” (D\&C 107:3–4.)

The desire and work of Satan is to mislead and corrupt. He seeks to frustrate the gospel plan by which God has provided the opportunity of eternal life for His children.

Satan seeks to discredit the sacred names of God the Father and his Son, Jesus Christ, the names through which their work is done. He succeeds in a measure whenever he is able to influence any man or woman, boy or girl, to make holy names common and to associate them with coarse thoughts and evil acts. Those who use sacred names in vain are, by that act, promoters of Satan’s purposes.

Profanity is profoundly offensive to those who worship the God whose name is desecrated. We all remember how a prophet reacted from a hospital bed when an operating room attendant stumbled and cursed in his presence. Even half-conscious, Elder [Spencer W.] Kimball “recoiled and implored: ‘Please! Please! That is my Lord whose names you revile.’” (Improvement Era, May 1953, p. 320.)

The words we speak are important. The Savior taught that men will be held to account for “every idle word” in the day of judgment. “For by thy words thou shalt be justified, and by thy words thou shalt be condemned.” (Matt. 12:36–37.) He also said, “That which cometh out of the mouth, this defileth a man.” (Matt. 15:11.)

Truly, as the Apostle James taught, “The tongue is a fire, … an unruly evil” that can defile the whole body. (James 3:6, 8.)

Profanity also takes its toll on the one who uses it. As we read in Proverbs, “A wholesome tongue is a tree of life: but perverseness therein is a breach in the spirit.” (Prov. 15:4.) The Spirit of the Lord, the Holy Ghost, testifies of God the Father and his Son, Jesus Christ. (See 2 Ne. 31:18.) When those names are dishonored, that Spirit, which “doth not dwell in unholy temples” (Hel. 4:24), is offended and withdraws. For this reason, those who profane the name of God inevitably relinquish the companionship of his Spirit.

As the Apostle Paul taught Timothy, in order to be “approved unto God,” we must “shun profane and vain babblings: for they will increase unto more ungodliness.” (2 Tim. 2:15–16.) Profanity leads to more ungodliness because the Spirit of the Lord withdraws and the profane are left without guidance.

Vulgar and crude expressions are also offensive to the Spirit of the Lord.

The Apostle James taught that followers of Christ should be “slow to speak, slow to wrath,” and should “lay apart all filthiness.” (James 1:19, 21.) In the Bible, filthiness is a term associated with sexual sin and with lewd language. (See Ezek. 16:36; 24:13; Eph. 5:3–4.) Thus, Paul was surely condemning vulgarity when he wrote the Colossians, “Also put off all these; anger, wrath, malice, blasphemy, filthy communication out of your mouth.” (Col. 3:8.)

These biblical condemnations of vulgarity are needed in our day.

Indecent and vulgar expressions pollute the air around us. Relations that are sacred between husband and wife are branded with coarse expressions that degrade what is intimate in marriage and make commonplace what is forbidden outside it. Moral sins that should be unspeakable are in the common vernacular. Human conduct plunging downward from the merely immodest to the utterly revolting is written on the walls and shouted in the streets. Twentieth-century men and women of sensitivity can easily understand how Lot, a fugitive from the actions and speech of Sodom and Gomorrah, could have been “vexed with the filthy conversation of the wicked.” (2 Pet. 2:7.)

How soberly we must regard the Book of Mormon teachings that “there cannot any unclean thing enter into the kingdom of God; wherefore there must needs be a place of filthiness prepared for that which is filthy.” (1 Ne. 15:34; see also Alma 7:21.)

Profane and vulgar expressions are public evidence of a speaker’s ignorance, inadequacy, or immaturity.

A speaker who profanes must be ignorant or indifferent to God’s stern command that his name must be treated with reverence and not used in vain.

A speaker who mouths profanity or vulgarity to punctuate or emphasize speech confesses inadequacy in his or her own language skills. Properly used, modern languages require no such artificial boosters.

A speaker who employs profanity or vulgarity to catch someone’s attention with shock effect engages in a babyish device that is inexcusable as juvenile or adult behavior. Such language is morally bankrupt. It also progressively self-defeating, since shock diminishes with familiarity and the user can only maintain its effect by escalating its excess.

Members of the Church, young or old, should never allow profane or vulgar words to pass their lips. The language we use projects the images of our hearts, and our hearts should be pure. As the Savior taught:

“Out of the abundance of the heart the mouth speaketh.

“A good man out of the good treasure of the heart bringeth forth good things: and an evil man out of the evil treasure bringeth forth evil things.” (Matt. 12:34–35.)

The Book of Mormon teaches us that when we are brought before the judgment bar of God “our words will condemn us … and our thoughts will also condemn us.” (Alma 12:14.) Let us recognize profanity and vulgarity for what they are. They are sins that separate us from God and cripple our spiritual defenses by causing the Holy Ghost to withdraw from us.

We should abstain and we should teach our children to abstain from all such expressions.

We can also encourage our associates to do likewise. Where we have the courage to make a friendly request, like Elder Kimball, we will often receive a respectful and cooperative reply. Our married daughter who lives in Illinois had such an experience. As she took her turn carpooling the twelve-year-olds home from the soccer game, her noisy passengers filled the air with profanity. Firmly, but with good humor, she told the boys, “In our family we only use that name when we worship, so we ask you, please don’t say that name disrespectfully in our car.” The boys immediately complied, and, what is even more surprising, most of them still remembered the next time it was her turn to drive.

We obviously cannot control all that goes on in our presence. Modern revelation suggests one alternative for those who would be clean: “Go ye out from among the wicked. Save yourselves.” (D\&C 38:42.) Sometimes we can remove ourselves from language that is profane or vulgar. If this is not possible, we can at least register an objection so that others cannot conclude that our silence means approval or acquiescence.

Our thirteenth article of faith commits us to seek after things that are “virtuous, lovely, or of good report or praiseworthy.” The language of Latter-day Saints should be reverent and clean. We understand the eternal requirement of cleanliness, and we understand the sacred significance of the names of the Father and the Son.

I testify of God the Father and his Son, Jesus Christ, and pray that we may be more faithful in honoring their holy names. In the name of Jesus Christ, amen.

\newpage
\settalk{“Brother’s Keeper”}
\setspeaker{Dallin H. Oaks}
\settalkdate{October 1986}
\talkheading{“Brother’s Keeper”}{Dallin H. Oaks}

One of the consequences of mortality is the necessity of earning our daily bread (see Gen. 3:19; Moses 4:25). We do so as employees, as business people, and as investors. In all of our earning activities, we have the challenge of dealing fairly and considerately with others.

Our duty is clear. The Savior gave us the Golden Rule: “All things whatsoever ye would that men should do to you, do ye even so to them” (Matt. 7:12).

Satan’s position is the opposite. He sponsors self-interest, raw and unrefined by any other consideration. One of his most effective tools is the temptation to take unfair advantage in order to get gain. It has been so from the beginning.

Cain set the pattern of the world. Cain coveted the flocks of his brother Abel, and Satan showed him how to obtain them (see JST, Gen. 5:14, 23; Moses 5:29, 38). Satan taught Cain that a man could get worldly wealth by committing some evil against its owner (see JST, Gen. 5:16; Moses 5:31).

Cain killed Abel. The scriptures say that he did so “for the sake of getting gain” (Moses 5:50), the flocks of his brother (JST, Gen. 5:18; Moses 5:33). Seeing this, the Lord asked Cain, “Where is Abel thy brother?” Cain first attempted to cover his sin with a lie: “I know not.” Then he added a rationalization: “Am I my brother’s keeper?” (Gen. 4:9; Moses 5:34).

Are we our brothers’ keepers? In other words, are we responsible to look after the well-being of our neighbors as we seek to earn our daily bread? The Savior’s Golden Rule says we are. Satan says we are not.

Tempted of Satan, some have followed the example of Cain. They covet property and then sin to obtain it. The sin may be murder, robbery, or theft. It may be fraud or deception. It may even be some clever but legal manipulation of facts or influence to take unfair advantage of another. Always the excuse is the same: “Am I my brother’s keeper?”

Those who follow the example of Cain fulfill a Book of Mormon prophecy. Seeing our day, Nephi prophesied that many would say, “Lie a little, take the advantage of one because of his words, dig a pit for thy neighbor; there is no harm in this” (2 Ne. 28:8).

We live in a world where many look on the marketplace as a ruthless arena where the buyer must beware, where no one is obligated to do more than the law requires, and where fraud isn’t fraud unless you can prove it in court.

Members of the Church of Jesus Christ have a higher standard. President Harold B. Lee said, “The standard … in the Church must be visibly higher than the standard … in the world” (Ye Are the Light of the World, Salt Lake City: Deseret Book Co., 1974, p. 13). We are commanded to live the Golden Rule.

Despite that high standard, some who profess to be Christians seek to earn their living by systematically victimizing their neighbors.

Some seize wealth by trafficking in illegal drugs or pornography. Traders in these products enrich themselves by transactions that ruin the bodies, minds, or morals of their customers.

Other criminals live by stealing. And not all stealing is at gunpoint or by dark of night. Some theft is by deception, where the thief manipulates the confidence of his victim.

The white-collar cousin of stealing is fraud, which gets its gain by lying about an essential fact in a transaction.

Scheming promoters with glib tongues and ingratiating manners deceive their neighbors into investments the promoters know to be more speculative than they dare reveal.

Difficulties of proof make fraud a hard crime to enforce. But the inadequacies of the laws of man provide no license for transgression under the laws of God. Though their method of thievery may be immune from correction in this life, sophisticated thieves in white shirts and ties will ultimately be seen and punished for what they are. He who presides over that Eternal Tribunal knows our secret acts, and he is “a discerner of the thoughts and intents of the heart” (Heb. 4:12; D\&C 33:1).

Most of us can be relatively comfortable when a message on the Golden Rule in the workplace uses examples like illegal drugs and theft by deception. What follows is more challenging. And it should be. We cannot expect to be comfortable if we measure our conduct against the Savior’s command, “I would that ye should be perfect even as I” (3 Ne. 12:48). To follow in the footsteps of the only perfect person who ever lived, we must expect to stretch our souls.

Followers of Christ have the moral responsibility of earning their livings and conducting their financial transactions in ways that are consistent with the principles of the gospel and the teachings of the Savior. Members of The Church of Jesus Christ of Latter-day Saints should not be involved in employment or other activities upon which they cannot conscientiously ask the blessings of the Lord.

Persons who prosecute frivolous lawsuits do not measure up to this high standard. Groundless litigation rewards some plaintiffs handsomely, but it injures everyone else by raising the price of products and services.

An employee who receives the compensation agreed upon but does not perform the service agreed upon earns part of his living by injuring others.

So does an employer who is unfair to his employees. An idealistic young professional wrote Church headquarters about the plight of migrant farm workers. He had observed treatment that was probably illegal and certainly unchristian. When I read his letter, I thought of the positive example of Jesse Knight, the great benefactor of Brigham Young Academy. At a time when most mine owners exploited their workers, this Christian employer paid his miners something extra so they could earn their living in six days’ labor and rest on the Sabbath. He did not require them to patronize a company store. He built his workers a building for recreation, worship, and schooling. And Brother Knight would not permit the superintendent to question his workers about their religion or politics (see Jesse William Knight, The Jesse Knight Family, Salt Lake City: Deseret News Press, 1940, pp. 43–44; and Gary Fuller Reese, “Uncle Jesse,” master’s thesis, Brigham Young University, 1961, pp. 26–28).

Of course, we understand that what an employer can pay his employees is limited by what his business can obtain for its products or services in a competitive marketplace. Contracts also impose limits on legitimate economic expectations.

Christian standards should also apply to those who earn a living by selling or advertising products in the marketplace.

The marketplace for products and services has many potential buyers who are vulnerable because they are poorly informed or excessively trusting. For example, a friend told me of a young student couple who didn’t have enough money for rent, groceries, and tuition but were persuaded to sign up for an expensive self-improvement course. Can a seller ever justify obtaining personal profit by persuading someone to assume a financial burden he cannot wisely bear in order to acquire something he does not really need? The Prophet Joseph Smith taught that Latter-day Saints should deal justly with their neighbors and mercifully with the poor (see History of the Church, 5:401).

To cite another kind of example, an owner who keeps his business open on Sunday prevents his employees from attending worship services and being with their families on the Sabbath. Modern-day prophets have encouraged us not to shop on Sunday (see, for example, Spencer W. Kimball, in Conference Report, Oct. 1974, p. 6; or Ensign, Nov. 1974, p. 6). Those of us who shop on the Sabbath cannot escape responsibility for encouraging businesses to remain open on that day. Essential services must be provided, but most Sabbath transactions could be avoided if merchants and customers were determined to avoid doing business on the Lord’s day.

Last year the Deseret News carried an article about a Salt Lake City pharmacist who stopped selling cigarettes in his drugstore. He explained, “It’s just incompatible for a profession dedicated to saving people’s lives to sell a product that does nothing but kill” (“S.L. Pharmacist Using Backhoe to Snuff out Cigarette Stock,” 20 Dec. 1985, p. B1). That merchant was more concerned about his customers’ welfare than his personal profits.

Sister Oaks called my attention to a similar example in the world of advertising. The magazine Women’s Sports and Fitness does not accept cigarette ads, thus foregoing much-needed revenue. A woman columnist and physician, Dr. Joan Ullyot, praised this policy and contrasted it to the practice of another organization:

“I am dismayed that a prominent women’s sport, tennis, continues to take support from a cigarette company. Surely the top women in this sport, none of whom smoke, have the [courage] to say no to this hypocrisy and stop lending their names and prestige to sanction and glamorize a lethal product. Any role model in sport who accepts support or sponsorship from a company whose products destroy health and fitness should take a hard look at what she is, by association, endorsing” (Women’s Sports and Fitness, Sept. 1986, p. 12).

Wouldn’t it be wonderful if this same attitude of looking after the interests of others governed Latter-day Saints who are making a profit from the sale or promotion of alcoholic beverages? Consider the terrible effects of alcohol. Alcohol-related accidents are the leading cause of death of those under twenty-five. The physical, social, and financial effects of alcohol ruin marriages and family life. By dulling inhibitions, alcohol leads to untold numbers of crimes and moral transgressions. Alcohol is the number one addictive drug in our day.

The consumption of alcohol is increasing among youth. Targeting young audiences, advertisers portray beer and wine as joyful, socially desirable, and harmless. Producers are promoting new types of alcoholic beverages as competitors in the huge soft-drink market. Grocery and convenience stores and gas stations stock alcoholic beverages side by side with soda pop. Can Christians who are involved in this commerce be indifferent to the physical and moral effects of the alcohol from which they are making their profits?

Other examples could be given, but these few are sufficient to illustrate the principle that the Golden Rule applies to our earning activities. We are our brother’s keeper, even in the marketplace.

I am aware that this is a high standard which cannot be met overnight. But it is important to recognize our responsibility and begin to work toward it. And we should do so joyfully. The gospel is the good news. Commandments lead to blessings. The Prophet Joseph Smith told our first missionaries that when preaching we should “warn in compassion.” We “have no right … to scare mankind to repentance,” he said. We should preach the gospel as “glad tidings of great joy unto all people” (History of the Church, 1:280).

We should also remember that the principle that the Golden Rule governs our earning activities is difficult to apply in practice. We should not consider employees responsible for policies they regret but cannot control. A decision that is made by the owner of a market should not inflict feelings of guilt on a conscientious but powerless Christian who runs the checkout stand. Similarly, a part-owner does not have freedom to impose his standards on business policies if he has partners who do not share his moral concerns. An incorporated business may be controlled by stockholders who have no concern for the destructive human effects of a profitable product or policy.

We live in a complex society, where even the simplest principle can be exquisitely difficult to apply. I admire investors who are determined not to obtain income or investment profits from transactions that add to the sum total of sin and misery in the world. But they will have difficulty finding investments that meet this high standard. Good things are often packaged with bad, so decisions usually involve balancing. In a world of corporate diversification, we are likely to find that a business dealing in beverages sells milk in one division and alcohol in another. Just when we think that our investments are entirely unspotted from the world, we may find that our life insurance is partially funded by investments we wish to avoid. Or our savings may be deposited in a bank that is lending to ventures we could not approve. Such complexities make it difficult to prescribe firm rules.

We must rely on teaching correct principles, which each member should personally apply to govern his or her own circumstances. To that end, each of us should give thoughtful and prayerful consideration to whether we are looking after the well-being of our neighbors in the way we earn our daily bread.

The motive of Cain is at the headwaters of wickedness. Cain’s sin was murder, but his motive was personal gain. That motive has produced all manner of wickedness, including murder, thievery, and fraud. That motive is also at work in the legal but immoral practices of those who get gain by preying on the weaknesses or ignorance of their neighbors. Always such activities involve Cain’s ancient rationalization: “Am I my brother’s keeper?”

In contrast, the Savior taught us to “love [our] enemies, bless them that curse [us], do good to them that hate [us], and pray for them who despitefully use [us] and persecute [us]” (3 Ne. 12:44). When we have that duty toward our enemies, we cannot allow ourselves to do less for our partners, our customers, our employees, and others with whom we deal in the marketplace.

What a beautiful and happy world this would be if all of us would strive to live these principles to the fullest. Our efforts and influence would affect millions. Examples improve society more than sermons. Most people would rather see a sermon than hear one.

In those brilliant generations that followed the appearance of the resurrected Christ in the New World, “there were no contentions and disputations among [the people], and every man did deal justly one with another” (4 Ne. 1:2). Fourth Nephi records: “Surely there could not be a happier people among all the people who had been created by the hand of God” (4 Ne. 1:16). We should be striving to regain that condition. As modern revelation declares: “Zion must increase in beauty, and in holiness” (D\&C 82:14). One of the ways prescribed to achieve that increase is “every man seeking the interest of his neighbor, and doing all things with an eye single to the glory of God” (D\&C 82:19).

May God bless us to live the Golden Rule in our earning activities. As we seek to be our brother’s keeper, we will be attempting to follow in the footsteps of the Master. I testify of Jesus Christ, our Savior, whose blood has atoned for repented sins and whose resurrection has broken the bands of death for all. The fulness of the gospel was restored through the Prophet Joseph Smith. His successor, President Ezra Taft Benson, holds the keys of the everlasting gospel in our day. In the name of Jesus Christ, amen.

\newpage
\settalk{Priesthood Blessings}
\setspeaker{Dallin H. Oaks}
\settalkdate{April 1987}
\talkheading{Priesthood Blessings}{Dallin H. Oaks}

In the spring of 1866, during what is called the Black Hawk War, our pioneers were struggling to beat back deadly Indian attacks on many settlements in southern Utah. Two of President Heber C. Kimball’s sons were called into military service for a three-month expedition against the Indians. Before they left he gave them a priesthood blessing. Apparently concerned that his sons might shed the blood of their Lamanite brothers, he first counseled them about the great promises God has made to this branch of the house of Israel. He then blessed his sons and promised them they would not see a single Indian on their campaign. His sons, full of fight and eager to smell gunpowder, were disappointed at this promise, but the blessing was fulfilled. When they returned three months later, they reported:

“We … rode hundreds of miles, following the tracks of different bands of hostile Indians, and were close upon them a great many times. They were attacking settlements all around us, killing the settlers and driving off stock.” But the company did not see a single Indian (Orson F. Whitney, Life of Heber C. Kimball, an Apostle, 2d ed., Salt Lake City: Stevens and Wallis, 1945, p. 429).

In a priesthood blessing a servant of the Lord exercises the priesthood, as moved upon by the Holy Ghost, to call upon the powers of heaven for the benefit of the person being blessed. Such blessings are conferred by holders of the Melchizedek Priesthood, which has the keys of all the spiritual blessings of the Church (see D\&C 107:18, 67).

There are many kinds of priesthood blessings. As I give various examples, please remember that priesthood blessings are available for all who need them, but they are only given on request.

Blessings for the healing of the sick are preceded by anointing with oil, as the scriptures direct (see James 5:14–15; Mark 6:13; D\&C 24:13–14; 42:43–48; 66:9). Patriarchal blessings are conferred by an ordained patriarch.

Persons desiring guidance in an important decision can receive a priesthood blessing. Persons who need extra spiritual power to overcome a personal challenge can receive a blessing. Expectant mothers can be blessed before they give birth. Many LDS families remember a sacred occasion where a worthy father gave a priesthood blessing to a son or daughter who was about to be married. Priesthood blessings are often requested from fathers before children leave home for other purposes, such as school, service in the military, or a long trip.

Newly called missionaries often request a father’s blessing before they depart. I have a friend who is blind. He remembers how his father blessed him that despite his physical disability he would be able to complete the mission, be successful in his calling, and develop a great love for the people. I am a witness to the fulfillment of that blessing in the life of a wonderful Latter-day Saint.

Blessings given in circumstances such as I have just described are sometimes called blessings of comfort or counsel. They are usually given by fathers or husbands or other elders in the family. They can be recorded and kept in family records for the personal spiritual guidance of the persons blessed.

Over ten years ago a teenage boy requested a blessing from President Ezra Taft Benson. Even though the boy’s father was not an active elder, President Benson asked, “How would you like to talk to him at an opportune time and ask him if he would be willing to give you a father’s blessing?” Though doubtful, the young man agreed to try. He later reported:

“Brother Benson, that’s the sweetest thing that has happened in our family. … He gave me one of the most beautiful blessings you could ever ask for. … When he got through there was a bond of appreciation and gratitude and love between us that we have never had in our home” (in Conference Report, Oct. 1977, pp. 45–46; or Ensign, Nov. 1977, p. 32).

Priesthood blessings are also given in connection with a priesthood ordination or with the setting apart of a man or woman for a calling in the Church. These are probably the most frequent occasions for priesthood blessings.

Many of us have requested a priesthood blessing when we were about to embark upon a new responsibility in our employment. I received such a blessing many years ago and felt its immediate comfort and long-term guidance.

In setting Dr. Russell M. Nelson apart as a stake president, a General Authority blessed him with power to meet the extremely demanding time requirements of his profession as a heart surgeon. Elder Nelson has described how that blessing was fulfilled by significant reductions in the risk of certain heart operations and in the time required for postoperative care. Eight years later, the man who had blessed him became his patient. Elder Spencer W. Kimball was scheduled for a complex heart operation. Presidents Harold B. Lee and N. Eldon Tanner blessed Dr. Nelson “that the operation would be performed without error, that all would go well, and that [he] need not fear for [his] own inadequacies, for [he] had been raised up by the Lord to perform this operation” (Ensign, May 1984, p. 88). That blessing was realized. A little over a year later, his fully recovered and vigorous patient became President of the Church and gave leadership through events and growth that will never be forgotten.

What is the significance of a priesthood blessing? Think of a young man preparing to leave home to seek his fortune in the world. If his father gave him a compass, he might use this worldly tool to help him find his way. If his father gave him money, he could use this to give him power over worldly things. A priesthood blessing is a conferral of power over spiritual things. Though it cannot be touched or weighed, it is of great significance in helping us overcome obstacles on the path to eternal life.

Remember how the Savior intervened to assure that the little children could come unto him. And then “he took them up in his arms, put his hands upon them, and blessed them” (Mark 10:16). When the risen Lord visited the people on this continent, “he took their little children, one by one, and blessed them, and prayed unto the Father for them” (3 Ne. 17:21).

It is a very sacred responsibility for a Melchizedek Priesthood holder to speak for the Lord in giving a priesthood blessing. As the Lord has told us in modern revelation, “My word … shall all be fulfilled, whether by mine own voice or by the voice of my servants, it is the same” (D\&C 1:38). If a servant of the Lord speaks as he is moved upon by the Holy Ghost, his words are “the will of the Lord, … the mind of the Lord, … the word of the Lord, … [and] the voice of the Lord” (D\&C 68:4).

But if the words of a blessing only represent the priesthood holder’s own desires and opinions, uninspired by the Holy Ghost, then the blessing is conditioned on whether it represents the will of the Lord.

Worthy Melchizedek Priesthood holders can give blessings to their posterity. The scriptures record many such blessings, including Adam’s (see D\&C 107:53–57), Isaac’s (see Gen. 27:28–29, 39–40; 28:3–4; Heb. 11:20), Jacob’s (see Gen. 48:9–22; 49; Heb. 11:21), and Lehi’s (see 2 Ne. 1:28–32; 4).

When Joseph Smith, Sr., was dying, his children gathered to receive his final blessing. After first blessing his wife, Father Smith began with Hyrum, his eldest, and gave each child what he called a “dying blessing” (see Lucy Mack Smith, History of Joseph Smith, ed. Preston Nibley, Salt Lake City: Bookcraft, 1956, pp. 308–13; Pearson H. Corbett, Hyrum Smith, Patriarch, Salt Lake City: Deseret Book Co., 1963, pp. 240–41).

In modern revelation, parents who are members of the Church are commanded to bring their children “before the church,” where the elders “are to lay their hands upon them in the name of Jesus Christ, and bless them in his name” (D\&C 20:70). This is why parents bring babies to a sacrament meeting, where an elder—usually the father—gives them a name and a blessing.

If any of the young men in this priesthood meeting has thought he has never received a priesthood blessing, I hope he has now realized that he has already received at least two and perhaps more.

Priesthood blessings are not limited to those blessings spoken as hands are laid on the head of one person. Blessings are sometimes pronounced on groups of people. The prophet Moses blessed all the children of Israel before his death (see Deut. 33:1). The Prophet Joseph Smith “pronounced a blessing upon the sisters” working on the Kirtland Temple. He also blessed “the congregation” (History of the Church, 2:399). As recently as last April conference President Benson blessed the Latter-day Saints and “good people everywhere … with increased power to do good and to resist evil,” and “with increased understanding of the Book of Mormon” (in Conference Report, Apr. 1986, p. 100; or Ensign, May 1986, p. 78).

Priesthood blessings are also pronounced on places. Nations are blessed and dedicated for the preaching of the gospel. Temples and houses of worship are dedicated to the Lord by a priesthood blessing. Other buildings may be dedicated when they are used in the service of the Lord. “Church members may dedicate their homes … as sacred edifices where the Holy Spirit can reside” (General Handbook of Instructions [1985], p. 11-5). Missionaries and other priesthood holders can leave a priesthood blessing upon homes where they have been received (see D\&C 75:19; Alma 10:7–11). Young men, within a short time you may be asked to give such a blessing. I hope you are preparing yourselves spiritually.

In the time that remains, I will mention some other examples of priesthood blessings.

About a hundred years ago, Sarah Young Vance qualified as a midwife. Before she began serving the women of Arizona, a priesthood leader blessed her that she would “always do only what was right and what was best for the welfare of her patients.” Over a period of forty-five years Sarah delivered approximately fifteen hundred babies without the loss of a single mother or child. “Whenever I came up against a difficult problem,” she recalled, “something always seemed to inspire me and somehow I would know what was the right thing to do” (L. J. Arrington and S. A. Madsen, Sunbonnet Sisters: True Stories of Mormon Women and Frontier Life, Salt Lake City: Bookcraft, 1984, p. 105).

In 1864, Joseph A. Young was called on a special mission to transact Church business in the East. His father, President Brigham Young, blessed him to go and return in safety. As he was returning, he was involved in a severe train wreck. “The whole train was smashed,” he reported, “including the car I was in to within one seat of where I sat, [but] I escaped without a scratch” (Letters of Brigham Young to His Sons, ed. Dean C. Jessee, Salt Lake City: Deseret Book Co., 1974, p. 4).

As a boy, I was inspired by a story of courage in Nauvoo, which involved my grandfather’s uncle. In the spring of 1844, some men were plotting against the Prophet Joseph Smith. One of the leaders, William Law, held a secret meeting at his home in Nauvoo. Among those invited were nineteen-year-old Dennison Lott Harris and his friend, Robert Scott. Dennison’s father, Emer Harris, who is my second great-grandfather, was also invited. He sought counsel from the Prophet Joseph Smith, who told him not to attend the meeting but to have the young men attend. The Prophet instructed them to pay close attention and report what was said.

The spokesmen at this first meeting denounced Joseph Smith as a fallen prophet and stated their determination to destroy him. When the Prophet heard this, he asked the young men to attend the second meeting. They did so, and reported the plotting.

A third meeting was to be held a week later. Again the Prophet asked them to attend, but he told them this would be their last meeting. “Be careful to remain silent and not to make any covenants or promises with them,” he counseled. He also cautioned them on the great danger of their mission. Although he thought it unlikely, it was possible they would be killed. Then, the Prophet Joseph Smith blessed Dennison and Robert by the power of the priesthood, promising them that if their lives were taken, their reward would be great.

In the strength of this priesthood blessing, they attended the third meeting and listened to the murderous plans. Then, when each person was required to take an oath to join the plot and keep it secret, they bravely refused. After everyone else had sworn secrecy, the whole group turned on Dennison and Robert, threatening to kill them unless they took the oath also. Because any refusal threatened the secrecy of their plans, about half of the plotters proposed to kill these two immediately. Knives were drawn, and angry men began to force them down into a basement to kill them.

Other plotters shouted to wait. Parents probably knew where they were. If they didn’t return, an alarm would be sounded and a search could reveal the boys’ deaths and the secret plans. During a long argument, two lives hung in the balance. Finally, the group decided to threaten to kill the young men if they ever revealed anything that had occurred and then to release them. This was done. Despite this threat, and because they had followed the Prophet’s counsel not to make any promises to the conspirators, Dennison and Robert promptly reported everything to the Prophet Joseph Smith.

For their own protection, the Prophet had these courageous young men promise him that they would never reveal this experience, not even to their fathers, for at least twenty years. A few months later, the Prophet Joseph Smith was murdered.

Many years passed. The members of the Church settled in the West. While Dennison L. Harris was serving as bishop of the Monroe Ward in southern Utah, he met a member of the First Presidency at a Church meeting in Ephraim. There, on Sunday, 15 May 1881, thirty-seven years after the Prophet Joseph Smith had sealed his lips to protect him against mob vengeance, Dennison Harris recited this experience to President Joseph F. Smith (see Verbal Statement of Bishop Dennison L. Harris, 15 May 1881, MS 2725, Historical Department, The Church of Jesus Christ of Latter-day Saints, Salt Lake City; the account was later published in the Contributor, Apr. 1884, pp. 251–60). Dennison Harris’s posterity includes many notable Latter-day Saints, including Franklin S. Harris, long-time president of Brigham Young University.

As I speak of priesthood blessings, I have a flood of memories: I remember my sons and daughters asking for blessings to help them through the most stressful experiences of their lives. I rejoice as I recall inspired promises and the strengthened faith that came when they were fulfilled. I feel pride in the faith of a new generation when I think of a son, apprehensive about a professional examination and unable to reach his faraway father, seeking a priesthood blessing from the most accessible priesthood holder in his family, the husband of his sister. I remember a confused young convert to the Church seeking a blessing to help him change the self-destructive pattern of his life. He received a blessing so unusual I was astonished when I heard the words I spoke.

Brethren, young and old, do not be hesitant to ask for a priesthood blessing when you are in need of spiritual power. Fathers and other elders, cherish and magnify the privilege of blessing your children and the other children of our Heavenly Father. Be prepared to give priesthood blessings under the influence of the Holy Ghost whenever you are requested in sincerity and faith.

This is the true church of our Savior. I testify of the saving mission of Jesus Christ. We are bearers of his priesthood. God bless us to exercise that priesthood under his direction, for the blessing of his children. In the name of Jesus Christ, amen.

\newpage
\settalk{“The Light and Life of the World”}
\setspeaker{Dallin H. Oaks}
\settalkdate{October 1987}
\talkheading{“The Light and Life of the World”}{Dallin H. Oaks}

My dear brothers and sisters, I rejoice with you in the privilege of coming together on this beautiful Sabbath day to worship our Father in Heaven and his Son Jesus Christ and to be instructed by their servants.

The Book of Mormon tells of the resurrected Lord visiting some of the people of the Americas. Clothed in a white robe, he descended out of heaven. Standing in the midst of a multitude, he stretched forth his hand and said:

“Behold, I am Jesus Christ, whom the prophets testified shall come into the world.

“And behold, I am the light and the life of the world” (3 Ne. 11:10–11).

He has repeated this declaration in many modern revelations (see D\&C 12:9; 39:2; 45:7). In harmony with his words, we solemnly affirm that Jesus Christ, the Only Begotten Son of God the Eternal Father, is the light and life of the world.

Jesus Christ is the light and life of the world because all things were made by him. Under the direction and according to the plan of God the Father, Jesus Christ is the Creator, the source of the light and life of all things. Through modern revelation we have the testimony of John, who bore record that Jesus Christ is “the light and the Redeemer of the world, the Spirit of truth, who came into the world, because the world was made by him, and in him was the life of men and the light of men.

“The worlds were made by him; men were made by him; all things were made by him, and through him, and of him” (D\&C 93:9–10).

\intalkheader{The Light of the World}

Jesus Christ is the light of the world because he is the source of the light which “proceedeth forth from the presence of God to fill the immensity of space” (D\&C 88:12). His light is “the true light that lighteth every man that cometh into the world” (D\&C 93:2; see also D\&C 84:46). The scriptures call this universal light “the light of truth” (D\&C 88:6), “the light of Christ” (D\&C 88:7; Moro. 7:18), and the “Spirit of Christ” (Moro. 7:16). This is the light that quickens our understanding (see D\&C 88:11). It is “the light by which [we] may judge” (Moro. 7:18). It “is given to every man, that he may know good from evil” (Moro. 7:16).

Jesus Christ is also the light of the world because his example and his teachings illuminate the path we should walk to return to the presence of our Father in Heaven. Before Jesus was born, Zacharias prophesied that the Lord God of Israel would visit his people “to give light to them that sit in darkness and in the shadow of death, to guide [their] feet into the way of peace” (Luke 1:79).

During his ministry Jesus taught, “Behold I am the light; I have set an example for you” (3 Ne. 18:16). Later, he told his Apostles, “Hold up your light that it may shine unto the world,” adding, “Behold, I am the light which ye shall hold up—that which ye have seen me do” (3 Ne. 18:24). He taught the Nephite multitude, “Ye know the things that ye must do in my church; for the works which ye have seen me do that shall ye also do” (3 Ne. 27:21).

The Savior emphasized the close relationship between his light and his commandments when he taught the Nephites, “Behold I am the law, and the light” (3 Ne. 15:9).

The Psalmist expressed that relationship: “Thy word is a lamp unto my feet, and a light unto my path” (Ps. 119:105).

As the Lord led Lehi and his people out of Jerusalem, he said, “I will also be your light in the wilderness; and I will prepare the way before you, if it so be that ye shall keep my commandments” (1 Ne. 17:13).

As we keep the Lord’s commandments, we see his light ever brighter on our path and we realize the fulfillment of Isaiah’s promise, “And the Lord shall guide thee continually” (Isa. 58:11).

Jesus Christ is also the light of the world because his power persuades us to do good. The prophet Mormon taught: “All things which are good cometh of God; …

“Wherefore, every thing which inviteth and enticeth to do good, and to love God, and to serve him, is inspired of God” (Moro. 7:12–13). Mormon’s words anticipate what the Lord later told Moroni while he was compiling the Book of Mormon: “He that believeth these things which I have spoken … shall know that these things are true; for it persuadeth men to do good.

“And whatsoever thing persuadeth men to do good is of me; for good cometh of none save it be of me. … I am the light, and the life, and the truth of the world” (Ether 4:11–12; see also D\&C 11:12).

And so we see that Jesus Christ is the light of the world because he is the source of the light that quickens our understanding, because his teachings and his example illuminate our path, and because his power persuades us to do good.

\intalkheader{The Life of the World}

Jesus Christ is the life of the world because of his unique position in what the scriptures call “the great and eternal plan of deliverance from death” (2 Ne. 11:5).

Jesus taught: “I am the door: by me if any man enter in, he shall be saved. …

“I am come that they might have life, and that they might have it more abundantly” (John 10:9–10).

Later, Jesus explained to his Apostles, “I am the way, the truth, and the life: no man cometh unto the Father, but by me” (John 14:6).

We come to the Father through the life-giving mission of the Son in two ways. In each of these ways, Jesus Christ is the life of the world, our Savior and our Redeemer.

Through the power and example of the infinite atonement of Jesus Christ, all mankind will be resurrected (see 2 Ne. 9:7, 12). Our mortal life came into being because of his creative act. Our immortal life has now been assured because the Resurrected Lord has redeemed us from death. According to the plan of the Father, the Son was “the firstborn from the dead” (Col. 1:18). “As in Adam all die, even so in Christ shall all be made alive” (1 Cor. 15:22).

Jesus Christ is also the life of the world because he has atoned for the sins of the world. By yielding to temptation, Adam and Eve were “cut off from the presence of the Lord” (Hel. 14:16). In the scriptures this separation is called spiritual death (see Hel. 14:16; D\&C 29:41).

The atonement of our Savior overcame this spiritual death. The scriptures say, “The Son of God hath atoned for original guilt” (Moses 6:54). As Paul taught the Saints in Rome: “Therefore as by the offence of one judgment came upon all men to condemnation; even so by the righteousness of one the free gift came upon all men unto justification of life” (Rom. 5:18). As a result of this atonement, “men will be punished for their own sins, and not for Adam’s transgression” (A of F 1:2).

Our Savior has redeemed us from the sin of Adam, but what about the effects of our own sins? Since “all have sinned” (Rom. 3:23), we are all spiritually dead. Again, our only hope for life is our Savior, who, the prophet Lehi taught, “offereth himself a sacrifice for sin, to answer the ends of the law” (2 Ne. 2:7).

In order to lay claim upon our Savior’s life-giving triumph over the spiritual death we suffer because of our own sins, we must follow the conditions he has prescribed. As he has told us in modern revelation, “I, God, have suffered these things for all, that they might not suffer if they would repent;

“But if they would not repent they must suffer even as I” (D\&C 19:16–17).

Our third article of faith describes the Savior’s conditions in these words: “We believe that through the Atonement of Christ, all mankind may be saved, by obedience to the laws and ordinances of the Gospel.”

In the words of our Savior, recorded in the Book of Mormon as he taught the people on this continent, “And whosoever will hearken unto my words and repenteth and is baptized, the same shall be saved” (3 Ne. 23:5).

In summary, the Lord Jesus Christ, our Savior and our Redeemer, is the life of the world because his resurrection and his atonement save us from both physical and spiritual death. Jacob rejoiced in this gift of life: “O how great the goodness of our God, who prepareth a way for our escape from the grasp of this awful monster; yea, that monster, death and hell, which I call the death of the body, and also the death of the spirit” (2 Ne. 9:10).

I wish that everyone could understand our belief and hear our testimony that Jesus Christ, our Savior and our Redeemer, is the light and life of the world.

\intalkheader{Our Savior and Our Redeemer}

Some who profess to be followers of Christ insist that members of The Church of Jesus Christ of Latter-day Saints are not Christians. Indeed, there are those who make their living attacking our church and its doctrines. I wish all of them could have the experience I shared recently.

A friend who was making his first visit to Salt Lake City, called on me in my office. He is a well-educated man and a devout and sincere Christian. Although we have not discussed this with each other, we both know that some leaders of his denomination have taught that members of our church are not Christians.

After a short discussion on a matter of common interest, I told my friend I had something I would like him to see. We walked over to Temple Square and into the North Visitors’ Center. We viewed the pictures of Bible and Book of Mormon Apostles and prophets. Then we turned our steps up the inclined walkway to the second level. Here Thorvaldsen’s great statue of the risen Christ dominates a setting suggestive of the immensity of space and the grandeur of the creations of God.

As we emerged and beheld this majestic likeness of the Christus, arms outstretched and hands showing the wounds of his crucifixion, my friend drew a sharp breath. We stood quietly for a few minutes, enjoying a reverent communion of worshipful thoughts about our Savior. Then, without further conversation, we made our way down to the street level. On the way we walked past the small diorama showing the Prophet Joseph Smith kneeling in the Sacred Grove.

As we left Temple Square and took our leave of one another, my friend took me by the hand. “Thank you for showing me that,” he said. “Now I understand something about your faith that I have never understood before.” I hope that every person who has ever had doubts about whether we are Christians can achieve that same understanding.

We love the Lord Jesus Christ. He is the Messiah, our Savior and our Redeemer. His is the only name by which we can be saved (see Mosiah 3:17, Mosiah 5:8; D\&C 18:23). We seek to serve him. We belong to his church, The Church of Jesus Christ of Latter-day Saints. Our missionaries and members testify of Jesus Christ in many nations of the world. As the prophet Nephi wrote in the Book of Mormon, “We talk of Christ, we rejoice in Christ, we preach of Christ, we prophesy of Christ, and we write according to our prophecies, that our children may know to what source they may look for a remission of their sins” (2 Ne. 25:26).

As we state in our first article of faith, “We believe in God, the Eternal Father, and in His Son, Jesus Christ, and in the Holy Ghost.” God the Father, the great Elohim, the Almighty God, is the Father of our spirits, the framer of heaven and earth, and the author of the plan of our salvation (see Moses 1:31–33; 2:1–2; D\&C 20:17–21). Jesus Christ is his Only Begotten Son, Jehovah, the Holy One and God of Israel, the Messiah, “the God of the whole earth” (3 Ne. 11:14). As the Book of Mormon declares, “Salvation was, and is, and is to come, in and through the atoning blood of Christ, the Lord Omnipotent” (Mosiah 3:18; see also Moses 6:52, 59). The scriptures proclaim and we reverently affirm that Jesus Christ is the light and life of the world.

What does this knowledge mean to Latter-day Saints? (We call ourselves “Saints” because this is the scriptural term for those who have sought to make their lives holy by entering into covenants to follow Christ.)

Our Savior is the light of the world. We should live so that we can be enlightened by his Spirit, and so that we may hear and heed the ratifying seal of the Holy Ghost, which testifies of the Father and the Son (see D\&C 20:26). We should study the principles of his gospel and receive its ordinances. We should keep his commandments, including his two great commandments to love God and to love and serve our neighbors (see Matt. 22:36–40). We should be faithful to the covenants we have made in the name of Jesus Christ.

Our Savior is also the life of the world. We should give thanks for his absolute gift of immortality. We should receive the ordinances and keep the covenants necessary to receive his conditional gift of life eternal, the greatest of all the gifts of God (see D\&C 14:7).

In short, Latter-day Saints invite each other and all men and women everywhere to “come unto Christ.” As a prophet has told us in the Book of Mormon: “I would that ye should come unto Christ, who is the Holy One of Israel, and partake of his salvation, and the power of his redemption. Yea, come unto him, and offer your whole souls as an offering unto him, and continue in fasting and praying, and endure to the end; and as the Lord liveth ye will be saved” (Omni 1:26).

May God bless all of us to come unto Christ. I testify that he is our Savior and our Redeemer, the light and the life of the world, in the name of Jesus Christ, amen.

\newpage
\settalk{Always Remember Him}
\setspeaker{Dallin H. Oaks}
\settalkdate{April 1988}
\talkheading{Always Remember Him}{Dallin H. Oaks}

In April 1830 the Lord commanded the members of his newly restored Church to “meet together often to partake of bread and wine in the remembrance of the Lord Jesus” (D\&C 20:75). This was the same instruction he gave when he introduced this ordinance nearly 2,000 years ago. Luke writes:

“And he took bread, and gave thanks, and brake it, and gave unto them, saying, This is my body which is given for you: this do in remembrance of me” (Luke 22:19).

When we partake of the sacrament, we witness unto God the Eternal Father that we “do always remember” his Son (see D\&C 20:77, 79; 3 Ne. 18:7, 11). Each Sabbath day millions of Latter-day Saints make this promise. What does it mean to “always remember” our Savior?

To remember means to keep in memory. In the scriptures, it often means to keep a person in memory, together with associated emotions like love, loyalty, or gratitude. The stronger the emotion, the more vivid and influential the memory. Here are some examples:

By now you must surely realize that I have given these three examples because the reasons why I will always remember these persons are related to the reasons why we should always remember Jesus Christ: He is our Creator, our Redeemer, and our Teacher.

\intalkheader{Our Creator, Our Redeemer, Our Teacher}

Under the direction and according to the plan of God the Father, his Son Jehovah “created the heavens and the earth, and all things that in them are” (3 Ne. 9:15). He gave us life in the beginning of this world, and through the power of his resurrection he will give each of us life again after we have died in mortality. Jesus Christ is the life of the world.

He is our Redeemer. According to the Father’s plan, he provided the atoning sacrifice that can rescue us from the extremity of spiritual death. As a free-will offering, the Only Begotten Son of God came to earth and shed his blood for the remission of our sins (see D\&C 27:2).

Our Creator and our Redeemer is also our teacher. He taught us how to live. He gave us commandments, and if we follow them, we will receive blessings and happiness in this world and eternal life in the world to come.

And so we see that He whom we should always remember is He who gave us mortal life, He who showed us the way to a happy life, and He who redeems us so we can have immortality and eternal life.

If we keep our covenant that we will always remember him, we can always have his Spirit to be with us (see D\&C 20:77, 79). That Spirit will testify of him, and it will guide us into truth.

His teachings and his example will guide and strengthen us in the way we should live. The effect was described in the words of the once popular song, “Try to remember, and if you remember, then follow” (“Try to Remember,” words by Tom Jones).

I will now refer to some of these teachings we should remember and follow.

\intalkheader{Serve as Called}

Follow is the word the Savior used when he called his helpers to the ministry. As he was walking by the Sea of Galilee, he saw two fishermen, Simon Peter and his brother Andrew, at work in their vocation. “And he saith unto them, Follow me, and I will make you fishers of men” (Matt. 4:19). “And straightway they forsook their nets, and followed him” (Mark 1:18).

Here the Savior established a pattern for those he calls to do his work. Acting through his servants, for he has said that “by mine own voice or by the voice of my servants, it is the same” (D\&C 1:38), he calls us to take time from our daily activities to follow him and serve our fellowmen. Even the greatest among us should be the servant of all (see Mark 10:43–44). Those who always remember him will straightway assume and faithfully fulfill the responsibilities to which they are called by his servants.

\intalkheader{Forgive Others}

Among the things we should remember about the Savior is that there are things we should forget about our fellowmen—the wrongs they have done us. “Lord,” the Apostle Peter asked the Master, “how oft shall my brother sin against me, and I forgive him? till seven times?” (Matt. 18:21). In response, Jesus taught the parable of the unforgiving servant. This man owed a large debt to his king. When he begged for mercy, the king was moved with compassion and forgave the debt. But when a fellow servant owed him a debt, this man took his debtor by the throat and cast him into prison until he should pay it. When the unforgiving servant was brought to judgment, the king said:

“Shouldest not thou also have had compassion on thy fellowservant, even as I had pity on thee?

“And his lord was wroth, and delivered him to the tormentors, till he should pay all that was due unto him.

“So likewise shall my heavenly Father do also unto you,” Jesus concluded (Matt. 18:33–35; see also Matt. 6:14–15; 3 Ne. 13:14–15).

As the Lord has told us in modern revelation, “He that forgiveth not his brother his trespasses standeth condemned before the Lord; for there remaineth in him the greater sin” (D\&C 64:9). If we always remember our Savior, we will forgive and forget grievances against those who have wronged us.

\intalkheader{Receive Ordinances}

At the beginning of his ministry, Jesus sought out John the Baptist, who was preaching the baptism of repentance for the remission of sins (see Mark 1:4).

“Then cometh Jesus from Galilee to Jordan unto John, to be baptized of him.

“But John forbade him, saying, I have need to be baptized of thee, and comest thou to me?

“And Jesus answering said unto him, Suffer it to be so now: for thus it becometh us to fulfil all righteousness” (Matt. 3:13–15).

Those who seek to follow the Savior will understand the importance of the ordinance of baptism. The Lamb without Blemish saw fit to submit himself to baptism by one holding the authority of the priesthood in order to “fulfil all righteousness.” How much more each of us has need of the cleansing and saving power of this ordinance and the other ordinances of the gospel.

As we always remember him, we should strive to assure that we and our family members and, indeed, all the sons and daughters of God everywhere follow our Savior into the waters of baptism. This reminds each of us of our duties to proclaim the gospel, perfect the Saints, and redeem the dead.

\intalkheader{Endure Afflictions}

Remembering the Savior can also help us understand and endure the inevitable afflictions of this life. The Savior taught:

“Blessed are ye, when men shall revile you, and persecute you, and shall say all manner of evil against you falsely, for my sake.

“Rejoice, and be exceeding glad: for great is your reward in heaven: for so persecuted they the prophets which were before you” (Matt. 5:11–12).

\intalkheader{Minister to the Sick and the Afflicted}

When the Risen Lord appeared to the people on this continent, he taught them and called leaders and gave them the authority of his priesthood. Next he healed the sick, the lame, the blind, and all others who were afflicted in any manner. Then “he commanded that their little children should be brought” (3 Ne. 17:11). And he “blessed them, and prayed unto the Father for them” (v. 21).

As I remember this inspiring example, I also remember visits and letters I have had from persons caring for loved ones who are sick or who are afflicted with the infirmities of old age. I also remember loved ones grieving over little children with life-shortening or crippling physical or emotional disabilities. How their hearts ache for their little ones! How they need our love and support! I also remember the words, “Inasmuch as ye have done it unto one of the least of these my brethren, ye have done it unto me” (Matt. 25:40). Here our Savior gives an assurance of blessings for those who carry such burdens and a challenge for others who can lend them support.

\intalkheader{Love Our Neighbors}

We should always remember how the Savior taught us to love and do good to one another. Loving and serving one another can solve so many problems!

I recently received a letter from a sister in another country. She wrote about the plight of single adult members of the Church. “Where do I fit in?” she asked. She longed to join in church social activities, but she said they were always designed for couples. She felt herself the “odd one out,” forced by circumstances rather than choice to forego these wholesome associations “rather than risk breaking up even numbers.”

She wrote of the trauma of being single, especially when this resulted from a companion’s desertion, divorce, or death. When she was a married woman, she said, “I never once gave much thought to the plight of the single sisters, except experiencing a kind of helpless pity for them.” Now in that circumstance herself, she felt that the married sisters of her acquaintance tended to shun the sisters who were single. She asked me what could be done to help the single adult members of the Church with what she described as their “feelings of rejection, nonacceptance, and noncaring by their fellow Church members.” Judging from the letters we receive, I believe there are many thousands of single adult members, our brothers and sisters, with similar feelings.

Our Savior gave us the parable of the good shepherd who left the multitude and went out in search of a single sheep who was lost (see Luke 15:3–6). Doesn’t that same principle require couples who enjoy loving companionship to go out of their way to include in their social circles brothers and sisters who have been deprived of that companionship? “Try to remember, and if you remember, then follow.”

A few years ago I spoke by assignment to a chamber of commerce group in Salt Lake City. During a question-answer period, I listened to a fine woman who was not of our faith. She spoke movingly of the pain her children had experienced when they were shunned by LDS youth in school and social activities. More recently, a Utah convert to the Church has written of his concern at the way some non-LDS adults with good basic values come to Utah with high expectations for a life among good neighbors and then, as he wrote, “find themselves excluded at best and ostracized at worst.”

Of course, there will be differences in the personal standards and social activities of faithful Latter-day Saints and members of other groups. But these differences are no excuse for ostracism, arrogance, or unkindness by LDS people. As my convert friend wrote, “I personally believe that Satan is as active among the Saints in turning them away from their neighbors as he is in turning disaffected persons against the Church.”

As we covenant that we will always remember our Savior, we must not forget Jehovah’s command to Israel:

“But the stranger that dwelleth with you shall be unto you as one born among you, and thou shalt love him as thyself” (Lev. 19:34; see Ex. 22:21; Deut. 10:19).

We should always remember how Jesus commanded us to love our neighbor as ourselves. He illustrated that great teaching with the example of the Good Samaritan, who crossed the social barriers of his day to perform acts of kindness and mercy. Then the Savior said, “Go, and do thou likewise” (Luke 10:37).

A decade ago President Spencer W. Kimball said, “Let us fellowship the students from all nations as they come to our land, so that we, above all other people, treat them as brothers and sisters in true friendship, whether or not they are interested in the gospel” (Regional Representatives’ seminar, 29 Sept. 1978).

That prophetic instruction should guide our relationships with all of our neighbors.

\intalkheader{Much Given, Much Required}

As we remember our Lord and Savior, we should contemplate the great blessings we have as members of The Church of Jesus Christ of Latter-day Saints. We have been taught by the Lord Jesus Christ. We have been led by his prophets. We have received the sealing ordinances of his gospel. He has blessed us bounteously.

As we remember all of this, we should also remember the divine caution: “For of him unto whom much is given much is required” (D\&C 82:3; see also Luke 12:48). That eternal principle of law and justice is a measure of what God expects of us.

May we always remember, as we covenant to do, is my humble prayer in the name of Jesus Christ, amen.

\newpage
\settalk{“What Think Ye of Christ?”}
\setspeaker{Dallin H. Oaks}
\settalkdate{October 1988}
\talkheading{“What Think Ye of Christ?”}{Dallin H. Oaks}

“What think ye of Christ?” (Matt. 22:42.) That question is as penetrating today as when Jesus used it to confound the Pharisees almost two thousand years ago. Like a sword, sharp and powerful, it uncovers what is hidden, divides truth from error, and goes to the heart of religious belief.

Here are some answers being given today.

Some praise Jesus Christ as the greatest teacher who ever lived, but deny that he is Messiah, Savior, or Redeemer. Some prominent theologians teach that our secularized world needs “a new concept of God, stripped of the … supernatural.” They believe that “not even a suffering God can help to solve the pain and tragedy of modern man.” (John A. Hardon, Christianity in the Twentieth Century, Garden City, N.Y.: Doubleday and Co., 1971, pp. 356, 359.)

A bishop in one Christian denomination has declared that “Jesus was in every sense a human being, just as we are.” (“One Clergyman’s Views on the ‘Death of God,’” U.S. News \& World Report, 18 Apr. 1966, p. 57).

Under the influence of such teachings, the religion of many is like the creed of the humanists, who declare that “no deity will save us; we must save ourselves.” (The Encyclopedia of American Religions: Religious Creeds, 1st ed., ed. J. Gordon Melton, Detroit: Gale Research Co., 1973, p. 641).

Another church that claims roots in “Christianity” maintains that Jesus’ crucifixion was not the fulfillment of his mission, but evidence of its failure. They teach that he did not cleanse men of original sin, but that another messiah must come to complete our salvation and establish the kingdom of heaven on earth. (Holy Spirit Association for the Unification of World Christianity, Outline of the Principle, Level 4, 1980, pp. 79–83, 238–39, 247–48, 252, 298–99.)

Many years ago a young Latter-day Saint enrolled in a midwestern university and applied for a scholarship only available to Christians. Both the applicant and the university officials were unsure whether a Mormon was eligible. After consulting a panel of theologians, they concluded that this Mormon was a Christian.

When I first heard of that event over thirty years ago, I was shocked that anyone, especially a member of our church, would entertain any doubt that we are Christians. I have come to a better understanding of that confusion. I think we sometimes thoughtlessly give others cause to wonder. How does this happen?

For many years I was a teacher of law. A frequent teaching method in that discipline is to concentrate classroom instruction on the difficult questions—the obscure and debatable matters that lie at the fringes of learning. Some law teachers believe that the simple general rules that answer most legal questions are so obvious that students can learn them by independent study. As a result, these teachers devote little time to teaching the basics.

I believe some of us sometimes do the same thing in gospel teaching. We neglect to teach and testify to some simple, basic truths of paramount importance. This omission permits some members and nonmembers to get wrong ideas about our faith and belief.

What do members of The Church of Jesus Christ of Latter-day Saints think of Christ?

Jesus Christ is the Only Begotten Son of God the Eternal Father. He is our Creator. He is our Teacher. He is our Savior. His atonement paid for the sin of Adam and won victory over death, assuring resurrection and immortality for all men.

He is all of these, but he is more. Jesus Christ is the Savior, whose atoning sacrifice opens the door for us to be cleansed of our personal sins so that we can be readmitted to the presence of God. He is our Redeemer.

The Messiah’s atoning sacrifice is the central message of the prophets of all ages. It was prefigured by the animal sacrifices prescribed by the law of Moses, whose whole meaning, one prophet explained, “point[ed] to that great and last sacrifice [of] … the Son of God, yea, infinite and eternal.” (Alma 34:14.) The Atonement was promised and predicted by the Old Testament prophets. Isaiah declared:

“He was wounded for our transgressions, he was bruised for our iniquities: … and with his stripes we are healed.

“All we like sheep have gone astray; we have turned every one to his own way; and the Lord hath laid on him the iniquity of us all.

“He was … brought as a lamb to the slaughter. …

“He was cut off out of the land of the living: for the transgression of my people was he stricken. …

“He bare the sin of many, and made intercession for the transgressors.” (Isa. 53:5–8, 12.)

At the beginning of the Savior’s ministry, John the Baptist exclaimed, “Behold the Lamb of God, which taketh away the sin of the world.” (John 1:29.)

At the end of his ministry, as Jesus blessed the cup and gave it to his disciples, he said, “For this is my blood of the new testament, which is shed for many for the remission of sins.” (Matt. 26:28.) As Latter-day Saints partake of the sacrament of the Lord’s Supper, we drink water in remembrance of his blood, which was shed for us. (See D\&C 20:79.)

The writers of the New Testament teach that our Savior’s suffering and his blood atoned for our sins.

The Apostle Paul told the Corinthians that the first principle of the gospel he preached to them was “how that Christ died for our sins according to the scriptures.” (1 Cor. 15:3.) And to the Colossians he wrote, “We have redemption through his blood, even the forgiveness of sins.” (Col. 1:14; see also Heb. 2:17; 10:10.)

Peter described how Christ “bare our sins in his own body on the tree, that we, being dead to sins, should live unto righteousness: by whose stripes ye were healed.” (1 Pet. 2:24.)

John wrote that “the blood of Jesus Christ … cleanseth us from all sin.” (1 Jn. 1:7; see also 2:2; 3:5; 4:10.)

We revere the Bible. And so we and our fellow believers in Christ sing these words from that inspiring hymn “How Great Thou Art”:

\begin{verse}
And when I think that God, his Son not sparing,\\
Sent him to die, I scarce can take it in,\\
That on the cross my burden gladly bearing\\
He bled and died to take away my sin.
\end{verse}

(Hymns, 1985, no. 86.)

Although the Bible’s explanation of atonement for individual sins should be unmistakable, that doctrine has been misunderstood by many who have only the Bible to explain it.

Modern prophets declare that the Book of Mormon contains the fulness of the everlasting gospel in greater clarity than any other scripture (see D\&C 20:8–9; 27:5). In a day when many are challenging the divinity of Jesus Christ or doubting the reality of his atonement and resurrection, the message of that second witness, the Book of Mormon, is needed more urgently than ever.

President Ezra Taft Benson has reminded us again and again that the Book of Mormon “was written for our day” and that it “is the keystone in our witness of Jesus Christ.” (In Conference Report, Oct. 1986, pp. 4, 5; or Ensign, Nov. 1986, pp. 5, 6.) I believe that the reason our Heavenly Father has had his prophet direct us into a more intensive study of the Book of Mormon is that this generation needs its message more than any of its forebears. As President Benson has said, the Book of Mormon “provides the most complete explanation of the doctrine of the Atonement,” and “its testimony of the Master is clear, undiluted, and full of power.” (Ensign, Nov. 1986, p. 5.)

In contrast, what is called “liberal theology” teaches that Jesus Christ is important not because he atoned for our sins, but only because he taught us the way to come to God by perfecting ourselves. In this theology, human beings can be reconciled to God entirely through their own righteousness. (See O. Kendall White, Jr., Mormon Neo-Orthodoxy: A Crisis Theology, Salt Lake City: Signature Books, 1987, pp. 43–44.)

Another group—secular rather than religious—believes that Jesus was not God, that man is God, and that you can create your own destiny through the powers of your own mind. (See “Age-old Fear of New Age Concerns,” Insight, 11 July 1988, p. 55.)

Are Latter-day Saints susceptible to such heresies? The Apostle Paul wrote that we should “work out [our] own salvation with fear and trembling.” (Philip. 2:12.) Could that familiar expression mean that the sum total of our own righteousness will win us salvation and exaltation? Could some of us believe that our heavenly parentage and our divine destiny allow us to pass through mortality and attain eternal life solely on our own merits?

On the basis of what I have heard, I believe that some of us, some of the time, say things that can create that impression. We can forget that keeping the commandments, which is necessary, is not sufficient. As Nephi said, we must labor diligently to persuade everyone “to believe in Christ, and to be reconciled to God; for we know that it is by grace that we are saved, after all we can do.” (2 Ne. 25:23.)

In his famous poem “Invictus,” William Ernest Henley hurled man’s challenge against Fate. With head “bloody, but unbowed,” determined man is unconquerable. The last verse reads:

\begin{verse}
It matters not how strait the gate,\\
How charged with punishments the scroll,\\
I am the master of my fate.\\
I am the captain of my soul.
\end{verse}

(In Out of the Best Books, 5 vols., ed. Bruce B. Clark and Robert K. Thomas, Salt Lake City: Deseret Book Co., 1968, 4:93.)

Writing a half-century later, Elder Orson F. Whitney replied with these lines:

\begin{verse}
Art thou in truth? Then what of him\\
Who bought thee with his blood?\\
Who plunged into devouring seas\\
And snatched thee from the flood?
\end{verse}

(Improvement Era, May 1926, p. 611.)

Man unquestionably has impressive powers and can bring to pass great things by tireless efforts and indomitable will. But after all our obedience and good works, we cannot be saved from the effect of our sins without the grace extended by the atonement of Jesus Christ.

The Book of Mormon puts us right. It teaches that “salvation doth not come by the law alone” (Mosiah 13:28); that is, salvation does not come by keeping the commandments alone. “By the law no flesh is justified.” (2 Ne. 2:5.) Even those who serve God with their whole souls are unprofitable servants. (See Mosiah 2:21.) Man cannot earn his own salvation.

The Book of Mormon teaches, “Since man had fallen he could not merit anything of himself.” (Alma 22:14.) “There can be nothing which is short of an infinite atonement which will suffice for the sins of the world.” (Alma 34:12; see also 2 Ne. 9:7; Alma 34:8–16.) “Wherefore, redemption cometh in and through the Holy Messiah; … he offereth himself a sacrifice for sin, to answer the ends of the law.” (2 Ne. 2:6–7.) And so we “preach of Christ … that our children may know to what source they may look for a remission of their sins.” (2 Ne. 25:26.)

In the Book of Mormon the Savior explains the gospel, including the Atonement and its relationship to repentance, baptism, works of righteousness, and the ultimate judgment:

“My Father sent me that I might be lifted up upon the cross, … that I might draw all men unto me, … that they may be judged according to their works. And … whoso repenteth and is baptized in my name shall be filled; and if he endureth to the end, behold, him will I hold guiltless before my Father at that day when I shall stand to judge the world.” (3 Ne. 27:14–16.)

In that same teaching the Savior restates these principles in a way that emphasizes our everlasting reliance on the Atonement worked out by the shedding of his blood:

“And no unclean thing can enter into [the Father’s] kingdom; therefore nothing entereth into his rest save it be those who have washed their garments in my blood, because of their faith, and the repentance of all their sins, and their faithfulness unto the end.” (3 Ne. 27:19.)

Joseph Smith stated this same relationship in our third article of faith: “We believe that through the Atonement of Christ, all mankind may be saved, by obedience to the laws and ordinances of the Gospel.” [A of F 1:3]

Why is Christ the only way? How was it possible for him to take upon himself the sins of all mankind? Why was it necessary for his blood to be shed? And how can our soiled and sinful selves be cleansed by his blood?

These are mysteries I do not understand. To me, as to President John Taylor, the miracle of the atonement of Jesus Christ is “incomprehensible and inexplicable.” (See The Mediation and Atonement of Our Lord and Savior Jesus Christ, Salt Lake City: Deseret News Co., 1882, pp. 148–49.) But the Holy Ghost has given me a witness of its truthfulness, and I rejoice that I can spend my life in proclaiming it.

I testify with the ancient and modern prophets that there is no other name and no other way under heaven by which man can be saved except by Jesus Christ. (See Acts 4:10, 12; 2 Ne. 25:20; Alma 38:9; D\&C 18:23.)

I witness with the prophet Lehi that “there is no flesh that can dwell in the presence of God, save it be through the merits, and mercy, and grace of the Holy Messiah.” (2 Ne. 2:8.)

I testify with the prophet Alma that no man can be saved except he is cleansed from all stain, through the blood of Jesus Christ. (See Alma 5:21.) As he explains, “repentance could not come unto men except there were a punishment” (Alma 42:16), and “therefore God himself atoneth for the sins of the world, to bring about the plan of mercy, to appease the demands of justice” (Alma 42:15).

I witness with the prophets of the Book of Mormon that the Messiah, the Holy One of Israel, suffered “according to the flesh” (Alma 7:13), the pains, the infirmities, and the griefs and sorrows of every living creature in the family of Adam. (See 2 Ne. 9:21; Alma 7:12–13; Mosiah 14:4; D\&C 18:11.)

I testify that when the Savior suffered and died for all men, all men became subject unto him (see 2 Ne. 9:5) and to his commandment that all must repent and be baptized in his name, having faith in him, “or they cannot be saved in the kingdom of God.” (2 Ne. 9:23; see also Alma 11:40; John 3:5; 8:24.)

Speaking through the Prophet Joseph Smith in our dispensation, the Savior said:

“I am … Christ the Lord, … the Redeemer of the world.

“I [have] accomplished and finished the will of him whose I am, even the Father, concerning me—having done this that I might subdue all things unto myself—

“Retaining all power, even to … judging every man according to his works and the deeds which he hath done.

“And surely every man must repent or suffer, for I, God, am endless. …

“Wherefore, I command you to repent. …

“For behold, I, God, have suffered these things for all, that they might not suffer if they would repent;

“But if they would not repent they must suffer even as I.” (D\&C 19:1–4, 13, 16–17.)

What think we of Christ? As members of The Church of Jesus Christ of Latter-day Saints, we testify with the Book of Mormon prophet-king Benjamin that “there shall be no other name given nor any other way nor means whereby salvation can come unto the children of men, only in and through the name of Christ, the Lord Omnipotent,

“For behold, … salvation was, and is, and is to come, in and through the atoning blood of Christ.” (Mosiah 3:17–18.)

And as we repent of our sins and seek to keep his commandments and our covenants, we cry out, as Benjamin’s people cried out, “O have mercy, and apply the atoning blood of Christ that we may receive forgiveness of our sins” (Mosiah 4:2).

In all of this, we remember and rely on the Lord’s sure word: “Keep my commandments in all things. And, if you keep my commandments and endure to the end you shall have eternal life, which gift is the greatest of all the gifts of God.” (D\&C 14:6–7.) In the name of Jesus Christ, amen.

\newpage
\settalk{Alternate Voices}
\setspeaker{Dallin H. Oaks}
\settalkdate{April 1989}
\talkheading{Alternate Voices}{Dallin H. Oaks}

Last summer, at a Pioneer parade in Wyoming, I saw a young colt separated from its mother. The lost youngster whinnied and trotted about, listening to a chorus of voices as it sought the voice that would guide it back to the side of the one it loved.

At other times I have seen lambs lost in a moving herd of sheep. A great chorus of voices rises from the herd, but each lamb listens for the one voice that can guide it. The Savior used this ageless example in the allegory of the Good Shepherd. “The sheep hear his voice: … and the sheep follow him: for they know his voice. And a stranger will they not follow, … for they know not the voice of strangers.” (John 10:3–5.)

From among the chorus of voices we hear in mortality, we must recognize the voice of the Good Shepherd, who calls us to follow him toward our heavenly home.

As Paul said to the Corinthians, “There are … so many kinds of voices in the world, and none of them is without signification.” (1 Cor. 14:10.)

Some voices speak of the things of the world, providing the useful information we need to make our way in mortality. I will make no further reference to these voices. My remarks will refer to those voices that speak of God, of his commandments, and of the doctrines, ordinances, and practices of his church. Some of those who speak on these subjects have been called and given divine authority to do so. Others, whom I choose to call alternate voices, speak on these subjects without calling or authority.

In the five years since I was called as a General Authority, I have seen many instances where Church leaders and members have been troubled by things said by these alternate voices. I am convinced that some members are confused about the Church’s relationship to the alternate voices. As a result, members can be misled in their personal choices, and the work of the Lord can suffer.

Some alternate voices are those of well-motivated men and women who are merely trying to serve their brothers and sisters and further the cause of Zion. Their efforts fit within the Lord’s teaching that his servants should not have to be commanded in all things, but “should be anxiously engaged in a good cause, and do many things of their own free will, and bring to pass much righteousness.” (D\&C 58:27.)

Other alternate voices are pursuing selfish personal interests, such as property, pride, prominence, or power. Other voices are the bleatings of lost souls who cannot hear the voice of the Shepherd and trot about trying to find their way without his guidance. Some of these voices call out guidance for others—the lost leading the lost.

Some alternate voices are of those whose avowed or secret object is to deceive and devour the flock. The Good Shepherd warned, “Beware of false prophets, which come to you in sheep’s clothing, but inwardly they are ravening wolves.” (Matt. 7:15; see also 3 Ne. 14:15.) In both the Bible and the Book of Mormon the Savior charged his shepherds to watch over and protect the flock from such wolves. (See Acts 20:28–29; Alma 5:59.)

There have always been alternate voices whose purpose or effect is to deceive. Their existence is part of the Plan. The prophet Lehi taught that there “must needs be … an opposition in all things.” (2 Ne. 2:11; italics added.) And there have always been other alternate voices whose purpose or effect is unselfish and wholesome.

In most instances, alternate voices are heard in the same kinds of communications the Church uses to perform its mission. The Church has magazines and other official publications, a newspaper supplement, letters from Church leaders, general conferences, and regular meetings and conferences in local units. Similarly, alternate voices are heard in magazines, journals, and newspapers and at lectures, symposia, and conferences.

The Church of Jesus Christ of Latter-day Saints does not attempt to isolate its members from alternate voices. Its approach, as counseled by the Prophet Joseph Smith, is to teach correct principles and then leave its members to govern themselves by personal choices.

Of course, the Church does have a responsibility to point out what is the voice of the Church and what is not. This is especially necessary when some alternate voice, deliberately or inadvertently, communicates a message in a way that implies Church sponsorship or acquiescence.

For the same reason, the Church does approve or disapprove those publications that are to be published or used in the official activities of the Church, general or local. For example, we have procedures to ensure approved content for materials published in the name of the Church or used for instruction in its classes. These procedures can be somewhat slow and cumbersome, but they have an important benefit. They provide a spiritual quality control that allows members to rely on the truth of what is said. Members who listen to the voice of the Church need not be on guard against being misled. They have no such assurance for what they hear from alternate voices.

Local Church leaders also have a responsibility to review the content of what is taught in classes or presented in worship services, as well as the spiritual qualifications of those they use as teachers or speakers. Leaders must do all they can to avoid expressed or implied Church endorsement for teachings that are not orthodox or for teachers who will use their Church position or prominence to promote something other than gospel truth.

Church leaders are sometimes invited to state the Church’s position at a debate or symposium about some doctrine, ordinance, or practice of the Church. This kind of presentation gives an audience the benefit of whatever illumination results from the adversarial clash of opposing viewpoints. Representatives of a business organization, a political party, or a social action group might welcome such an invitation. But the Church is directed to avoid disputation and contention. Moreover, if a representative of the Church participated in such an event, this could have the unwanted effect of encouraging Church members to look to the sponsors of alternate voices to bring them information on the positions of the Church.

Members of the Church are free to participate or to listen to any alternate voices they choose, but Church leaders should avoid official involvement, directly or indirectly.

There are disadvantages to official nonparticipation in events where Church doctrines, ordinances, or practices are discussed. In some instances, the overall presentation will be decidedly inaccurate or unfair because the position of the Church and the knowledge of its leaders are not presented. In other instances, a volunteer will step forward to present what he or she considers to be the Church’s position. Sometimes these volunteers are well-informed and capable, and they contribute to a balanced presentation. Sometimes they are not, and their contribution makes matters worse. When attacked by error, truth is better served by silence than by a bad argument.

In any case, volunteers do not speak for the Church. As long as Church leaders feel they should not participate in an event where the Church or its doctrines are discussed, the overall presentation will be incomplete and unbalanced. In such circumstances, no one should think that the Church’s silence constitutes an admission of facts asserted in that setting.

Individual members of the Church may also confront difficult questions when they are invited to participate. The question is more complicated when the invitation does not relate to a publication or a lecture on a single subject, but to a group of articles, a series of publications, or a conference or symposium with a large number of subjects. One article or one issue of a publication or one session of a conference may be edifying and uplifting, something a faithful Latter-day Saint would wish to support or enjoy. But another article or another session may be destructive, something a faithful Latter-day Saint would not wish to support or promote.

Some of life’s most complicated decisions involve mixtures of good and evil. To what extent can one seek the benefit of something good one desires when this can only be done by simultaneously promoting something bad one opposes? That is a personal decision, but it needs to be made with a sophisticated view of the entire circumstance and with a prayer for heavenly guidance.

There are surely limits at which every faithful Latter-day Saint would draw the line. For example, in my view a person who has made covenants in the holy temple would not make his or her influence available to support or promote a source that publishes or discusses the temple ceremonies, even if other parts of the publication or program are unobjectionable. I would not want my support or my name used to further a public discussion of things I have covenanted to hold sacred.

As Latter-day Saints consider their personal relationship to various alternate voices, they will be helped by considering the ways we acquire knowledge, especially knowledge of sacred things.

In modern revelation the Lord has told us to “seek learning … by study and also by faith.” (D\&C 109:7.)

We seek learning by studying the accumulated wisdom of various disciplines and by using the powers of reasoning placed in us by our Creator.

We should also seek learning by faith in God, the giver of revelation. I believe that many of the great discoveries and achievements in science and the arts have resulted from a God-given revelation. Seekers who have paid the price in perspiration have been magnified by inspiration.

The acquisition of knowledge by revelation is an extra bonus to seekers in the sciences and the arts, but it is the fundamental method for those who seek to know God and the doctrines of his gospel. In this area of knowledge, scholarship and reason are insufficient.

A seeker of truth about God must rely on revelation. I believe this is what the Book of Mormon prophet meant when he said, “To be learned is good if they hearken unto the counsels of God.” (2 Ne. 9:29.) It is surely what the Savior taught when he said, “Blessed art thou, Simon Bar-jona: for flesh and blood hath not revealed it unto thee, but my Father which is in heaven.” (Matt. 16:17.)

The way to revelation is righteousness. Marveling at the Master’s teachings, his enemies asked:

“How knoweth this man letters, having never learned?

“Jesus answered them, and said, My doctrine is not mine, but his that sent me.

“If any man will do his will, he shall know of the doctrine, whether it be of God, or whether I speak of myself.” (John 7:15–17.)

The Book of Mormon teaches that those who diligently seek shall have “the mysteries of God … unfolded unto them, by the power of the Holy Ghost.” (1 Ne. 10:19; see also 1 Cor. 2:4–16; Alma 18:35; D\&C 121:26.) The prophet Jacob declared the impossibility of uninspired man’s understanding God: “No man knoweth of his ways save it be revealed unto him; wherefore, brethren, despise not the revelations of God.” (Jacob 4:8.)

The Lord’s prescribed methods of acquiring sacred knowledge are very different from the methods used by those who acquire learning exclusively by study. For example, a frequent technique of scholarship is debate or adversarial discussion, a method with which I have had considerable personal experience. But the Lord has instructed us in ancient and modern scriptures that we should not contend over the points of his doctrine. (See 3 Ne. 11:28–30; D\&C 10:63.) Those who teach the gospel are instructed not to preach with “wrath” or “strife” (D\&C 60:14; see also 2 Tim. 2:23–25), but in “mildness and in meekness” (D\&C 38:41), “reviling not against revilers” (D\&C 19:30). Similarly, techniques devised for adversary debate or to search out differences and work out compromises are not effective in acquiring gospel knowledge. Gospel truths and testimony are received from the Holy Ghost through reverent personal study and quiet contemplation.

In the scriptures, the Lord has specified how we learn by faith. We must be humble, cultivate faith, repent of our sins, serve our fellowmen, and keep the commandments of God. (See Ether 12:27; D\&C 1:28; 12:8; 50:28; 63:23; 136:32–33.) As the Book of Mormon says, “Yea, he that repenteth and exerciseth faith, and bringeth forth good works, and prayeth continually without ceasing—unto such it is given to know the mysteries of God.” (Alma 26:22.)

I have seen some persons attempt to understand or undertake to criticize the gospel or the Church by the method of reason alone, unaccompanied by the use or recognition of revelation. When reason is adopted as the only—or even the principal—method of judging the gospel, the outcome is predetermined. One cannot find God or understand his doctrines and ordinances by closing the door on the means He has prescribed for receiving the truths of his gospel. That is why gospel truths have been corrupted and gospel ordinances have been lost when left to the interpretation and sponsorship of scholars who lack the authority and reject the revelations of God.

That is what the Savior told his professional critics, as recorded in the eleventh chapter of Luke. He was confronted by a group who had hypocritically built monuments to the prophets their predecessors had murdered, while personally rejecting the living prophets God was sending them. (See Luke 11:47–49.) In what I understand to be a condemnation of their rejection of revelation, the Savior pronounced woe upon these worldly professionals: “For ye have taken away the key of knowledge: ye entered not in yourselves, and them that were entering in ye hindered.” (Luke 11:52.)

The early leaders of the restored church had to learn that same truth. In several revelations the Lord rebuked Joseph Smith, David Whitmer, and others for not having their minds on the things of God, for yielding to “the persuasions of men” (D\&C 3:6; 5:21), and for being “persuaded by those whom I have not commanded” (D\&C 30:2).

The correct relationship between study and faith in the receipt of sacred knowledge is illustrated in Oliver Cowdery’s attempt to translate ancient records. He failed because he “took no thought,” but only asked God. (D\&C 9:7.) The Lord told him he should have “stud[ied] it out in [his] mind” and then asked if it was right. (D\&C 9:8.) Only then would the Lord reveal whether the translation was correct or not. And only on receiving that revelation could the text be written, because “you cannot write that which is sacred save it be given you from me.” (D\&C 9:9.) In the acquisition of sacred knowledge, scholarship and reason are not alternatives to revelation. They are a means to an end, and the end is revelation from God.

God has promised that if we ask him, we will “receive revelation upon revelation, knowledge upon knowledge, that [we may] know the mysteries and peaceable things—that which bringeth joy, that which bringeth life eternal.” (D\&C 42:61.)

In our day we are experiencing an explosion of knowledge about the world and its people. But the people of the world are not experiencing a comparable expansion of knowledge about God and his plan for his children. On that subject, what the world needs is not more scholarship and technology but more righteousness and revelation.

I long for the day prophesied by Isaiah when “the earth shall be full of the knowledge of the Lord.” (Isa. 11:9; 2 Ne. 21:9.) In an inspired utterance, the Prophet Joseph Smith described the Lord’s “pouring down knowledge from heaven upon the heads of the Latter-day Saints.” (D\&C 121:33.) This will not happen for those whose “hearts are set so much upon the things of this world, and aspire to the honors of men.” (V. 35.) Those who fail to learn and use “principles of righteousness” (V. 36) will be left to themselves to kick against those in authority, “to persecute the saints, and to fight against God” (V. 38). In contrast, the Lord makes this great promise to the faithful:

“The doctrine of the priesthood shall distil upon thy soul as the dews from heaven.

“The Holy Ghost shall be thy constant companion, and thy scepter an unchanging scepter of righteousness and truth; and thy dominion shall be an everlasting dominion, and without compulsory means it shall flow unto thee forever and ever.” (D\&C 121:45–46.)

I testify of these things in the name of Jesus Christ, amen.

\newpage
\settalk{Modern Pioneers}
\setspeaker{Dallin H. Oaks}
\settalkdate{October 1989}
\talkheading{Modern Pioneers}{Dallin H. Oaks}

The days of the pioneers are not past. There are modern pioneers whose achievements are an inspiration to all of us.

In a message about the pioneers who crossed the plains over a century ago, President J. Reuben Clark spoke words that apply to pioneers in every age. In his description of “Them of the Last Wagon,” President Clark paid tribute to the rank and file, “those great souls, … in name unknown, unremembered, unhonored in the pages of history, but lovingly revered round the hearthstones of their children and their children’s children.” (J. Reuben Clark: Selected Papers on Religion, Education, and Youth, ed. David H. Yarn, Jr., Provo: Brigham Young University Press, 1984, pp. 67–68; see also Improvement Era, Nov. 1947, pp. 704–5, 747–48.)

In every great cause there are leaders and followers. In the wagon trains, the leaders were “out in front where the air was clear and clean and where they had unbroken vision of the blue vault of heaven.” (Clark, p. 69.) But, as President Clark observed, “Back in the last wagon, not always could they see the brethren way out in front and the blue heaven was often shut out from their sight by heavy, dense clouds of the dust of the earth. Yet day after day, they of the last wagon pressed forward, worn and tired, footsore, sometimes almost disheartened, borne up by their faith that God loved them, that the Restored Gospel was true, and that the Lord led and directed the brethren out in front.” (Clark, p. 69.)

The purposes of God were accomplished by the unswerving loyalty and backbreaking work of the faithful tens of thousands who pushed on, as President Clark said, “with little praise, with not too much encouragement, and never with adulation.” (Clark, pp. 69–70.)

“And thousands upon thousands of these … measured to their humble calling and to their destiny as fully as Brother Brigham and the others measured to theirs, and God will so reward them. They were pioneers in word and thought and act and faith, even as were they of more exalted station. … God keep their memories ever fresh among us … to help us meet our duties even as they met theirs.” (Clark, pp. 73–74.)

President Clark’s words of tribute also apply to the membership of The Church of Jesus Christ of Latter-day Saints in our day. In every nation, in every worthy occupation and activity, members of this church face hardships, overcome obstacles, and follow the servants of the Lord Jesus Christ as valiantly as the pioneers of any age. They pay their tithes and offerings. They serve as missionaries or as Church Service volunteers, or they support others who do so. Like the noble young mothers who postpone the pursuit of their personal goals in order to provide the needs of their children, they sacrifice immediate pleasures to keep commitments that are eternal. They accept callings and, in the service of others, they willingly give their time and sometimes their lives.

They do as the Savior taught: They deny themselves; they take up their crosses daily; they follow Him. (See Luke 9:23.) These are those the Savior likened to the seed that fell on good ground: “in an honest and good heart, having heard the word, [they] keep it, and bring forth fruit with patience.” (Luke 8:15.)

The fruits of the gospel issue from every honest and good heart, without regard to past origins or current positions in the Church. As President Clark declared, “There is no aristocracy of birth in this Church; it belongs equally to the highest and the lowliest.” (Clark, p. 73.)

I will give some illustrations of modern pioneers. My examples are not necessarily the most notable, but I believe they are typical of the rank-and-file Latter-day Saints who are the heart and the hands of this great latter-day work.

Our older couple missionaries, now numbering over 2,600 throughout the world, provide an unequaled example of Christian service. Who could calculate the contribution these couples are making in furthering the mission of the Church? They preach the gospel, strengthen leaders and members in struggling branches, serve in temples and visitors’ centers, and in countless other ways accomplish the essential work of the kingdom, both the important and the routine.

In a missionary meeting in a remote corner of the world, Sister Oaks and I listened as a devoted brother said, “I never thought I could teach the gospel. I only thought I could fish. But now that I am here, I get so wrapped up in telling people about the gospel!”

A few minutes later, another devoted missionary, his wife, said, “I feel so sorry for those who have nothing to worry about and occupy them except how many steps to the swimming pool or the golf course!”

Time after time, the pioneers President Clark praised left their homes, loaded their wagons, and moved to new hardships at the direction of their prophet. In our day, many couples go on mission after mission. One dear veteran described her family’s reaction: “Our children say, ‘We hope you’ll come by and at least have dinner with us before you go on another mission.’”

Every day other thousands set aside personal preferences and give devoted service as teachers and leaders, as temple workers, in name extraction, and in so many other ways.

The Apostle Paul described the followers of Christ as “rejoicing in hope; patient in tribulation.” (Rom. 12:12.) We are tested for those qualities in different ways at different times.

A few weeks ago, some members of my family visited the Winter Quarters cemetery at Florence, Nebraska. There they saw Avard T. Fairbanks’s marvelous statue of the pioneer parents looking down at the body of their baby, soon to be left in its grave at the side of the trail. Those pioneers received some of their toughest tests at graveside. Some modern pioneers receive their tests at bedside. One sister wrote:

“My mother cared for her mother until [Grandma] was ninety-eight. My dad now has Alzheimer’s disease, and my mother patiently cares for him. … The amazing part of this is the attitude of my mother. She always thought she would travel after she retired. She has always kept a beautiful home, loving to entertain others. She maintains her home as best she can, but has had to put aside many things that bring her joy. The amazing part is the joy my mother radiates. Her attitude is so beautiful. She finds real joy in the simple things of life. She is the pillar of strength to the whole family as she uplifts us all with her positive attitude.”

There are hidden heroines and heroes among the Latter-day Saints—“those of the last wagon” whose fidelity to duty and devotion to righteousness go unnoticed by anyone except the One whose notice really matters.

Others, including those who have been called to prominent positions, are more noticeable, but surely no more noble. I am one of these. At a public occasion a mother introduced me to her teenage son. “Do you know who this is?” she asked him.

“Sure,” the boy replied. “He’s one of those guys who hangs on the wall at seminary.”

Prominent position—“hanging on the wall at seminary”—does not put anyone on a fast-track to exaltation. The criteria for that ultimate goal is the same for every person—leader or follower, prominent or obscure: Have we received the ordinances of salvation and kept our covenants? A member of the Church in Great Britain said it best. He had served as stake president. As that period of prominence came to an end, he told Elder Boyd K. Packer why it did not bother him to be released: “I served because I am under covenant. And I can keep my covenants quite as well as a home teacher as I can serving as stake president.” (In Conference Report, Apr. 1987, p. 26; or Ensign, May 1987, p. 24.)

Numberless officers, teachers, advisers, and clerks keep their covenants in that same way. Their service is almost invisible, except to Him who sees all things and promises all who do good that they shall “in nowise lose their reward.” (D\&C 58:28; see also Matt. 10:42.)

The pioneers who crossed the plains took their directions from the trails blazed by their leaders. For safety, those pioneers traveled in groups. Then, as now, a pioneer who got separated from the company and off the marked trail walked a lonely and dangerous path until he could rejoin the group. So it is today. A letter said it this way:

“One and a half years ago I was excommunicated. I was guilty of great hypocrisy and deception before God in matters of infidelity. This Saturday I am going to be baptized and receive the gift of the Holy Ghost. As the day approaches, my gratitude deepens for the Lord’s mercy extended to me, allowing me to repent and experience the mighty change in my heart. It grieves me to know of the great contribution I made to the Lord’s suffering in Gethsemane, but I glory in the proposition that I, as a result of that suffering, might turn my life and make His purposes my purposes.”

This writer expressed gratitude for “the Lord’s repentance process,” which would now “allow me to become the father, son, and priesthood-bearer that I always appeared to be. The feeling of finally being an honest, truly honest, man is indescribable.”

One of the best qualities in any of the sons and daughters of God, whatever their circumstance, is a determination to become better. Since we all have a need to improve, we should always be willing to recognize goodness and encourage improvement in everyone.

One of the most Godlike expressions of the human soul is the act of forgiveness. Everyone is wronged at some point by someone, and many suffer serious wrongs. Christians everywhere stand in awe of those pioneers who have climbed that steep slope to the spiritual summit attained by those who have heeded the Savior’s command to forgive all men. (See Matt. 6:14–15; D\&C 64:9–10.) Forgiveness is mortality’s mirror image of the mercy of God.

A sister wrote me about her feelings toward a relative who had abused her as a child, leaving her with a painful physical condition. In her words, “I have to live with the pain and try to function around it.” She wrote, “At times I [felt] angry and wonder[ed] why I had to suffer the abuse in the first place and why must I continue to pay a price now.”

One day, as she listened to a talk in church, her heart was touched. The Spirit bore witness that she should forgive the man who had wronged her and that she could do so with the help of our Lord Jesus Christ. Her letter explained: “The price for that sin has already been paid by Him in Gethsemane. I have no right to hold on to it and demand justice, so I gladly hand it back to Him and rejoice in His love and mercy.”

Her letter described the result of her decision: “My heart is so full of joy, peace, and gratitude and love! Isn’t His work glorious? How I do love Him! Words cannot express my feelings.”

Like this sister who forgave, many modern Saints do their pioneering on the frontiers of their own attitudes and emotions. The proverb says, “He that ruleth his spirit [is better] than he that taketh a city.” (Prov. 16:32.) Modern Saints know that one who subdues his own spirit is just as much a pioneer as one who conquers a continent.

The path of modern pioneers is not easy. Burdens carried in the heart can be just as heavy as those pulled in a handcart. And just as some early pioneers struggled for the benefit of others, so some modern pioneers carry burdens imposed by the transgressions or thoughtlessness of others.

Another letter came from a woman who had been divorced. Although she said that the ten years that followed her divorce were a time of trial, heartache, struggle, and loneliness, she described that experience as “a blessing”—“a refining process.” She expressed gratitude “for what I now have. It has brought me so close to my Heavenly Father and particularly to the Lord Jesus Christ. It is a feeling that I’m not sure can be expressed in words. I literally came before the Lord with a broken heart and a contrite spirit. No physical pain I have ever experienced has been as painful as the emotional pain I have felt. But each time I feel it, it draws me so close to the Lord because I think of all He suffered, and it makes me so grateful. I love Him with all my heart and soul for His sacrifice and for all He represents.”

Many of our members are struggling valiantly to try to do it all. They support themselves and provide for their families. They strive to carry out the responsibilities of their church callings. They spend many hours transporting their children to numberless church and school activities. They try to be generous with money and time for worthy causes in the community. They strive to improve themselves. They hope, after all of this, to have some little time left for togetherness and recreation.

One sister wrote, “We are having great difficulty [just] trying to cope.” Many could say the same. Yet they do cope. They carry on without complaint, even when they have just cause for complaint. And even when they fall short, the Lord blesses them for their righteous desires (see Mosiah 4:24–25), for, as King Benjamin taught, “it is not requisite that a man should run faster than he has strength” (Mosiah 4:27).

How grateful we are for the service and example of these faithful members! Like all my Brethren among the General Authorities, I look to the rank-and-file members of this church for my models of faithfulness and nobility. When I visit a conference and mingle with the Saints, I always receive more than I give. I agree with the sentiment voiced by President Gordon B. Hinckley. After describing the faithful Saints he had met at a conference, he added, “We have the responsibility of leading them, when, in fact, we can learn so much from them.”

Our faith and resolve are strengthened by the spiritual achievements and service of ordinary Latter-day Saints. There are thousands of such inspirational examples, but they are rarely published except on the pages of the Church News and the Church magazines—the Ensign, New Era, and Friend. I encourage everyone to have these unique publications in their home.

Some of the unsung heroes and heroines of our day are the faithful home teachers and visiting teachers who feed the Master’s sheep. When the Apostle Paul likened the Church to a body, he referred to such less-visible members as the hands and the feet, saying that upon these we should “bestow more abundant honour.” (1 Cor. 12:23.)

An LDS girl whose two parents took no part in Church activities later wrote this recollection to an elder who had been her home teacher:

“You were the bright hope in my often difficult life. There is no greater call than a home teacher. You loved and showed respect for my parents. You honored them and at the same time supported me. You were there! … As I look back now, I realize you and the truth you offered were my life-support.

“Behind the doors were years of pain, tears, and fear. You were able to come into our home and chase them away, if only for a short time. No one else could do that.”

In our day, as in the days of earlier pioneers, those in the lead wagons set the direction and signal onward, but it is the faithful men and women in the wagons which follow that provide the momentum and motive power for this great work.

As modern pioneers press forward, they suffer hardships and make sacrifices. But they are sustained by an assurance given by the Lord Himself. These words, first spoken to the struggling Saints in Ohio, apply also to the faithful of our day:

“Verily I say unto you my friends, fear not, let your hearts be comforted; yea, rejoice evermore, and in everything give thanks;

“Waiting patiently on the Lord, for your prayers have entered into the ears of the Lord of Sabaoth, and are recorded with this seal and testament—the Lord hath sworn and decreed that they shall be granted.

“Therefore, he giveth this promise unto you, with an immutable covenant that they shall be fulfilled; and all things wherewith you have been afflicted shall work together for your good, and to my name’s glory, saith the Lord.” (D\&C 98:1–3.)

This is His work. We are His children. He loves us—one and all. Of this I testify in the name of Jesus Christ, amen.

\newpage
\settalk{World Peace}
\setspeaker{Dallin H. Oaks}
\settalkdate{April 1990}
\talkheading{World Peace}{Dallin H. Oaks}

Some years ago, an acquaintance of mine who was moving to Washington, D.C., went to the district offices to take the driver’s license examination. He had to fill out a form that asked for his business address and his occupation. He had just been appointed a justice of the United States Supreme Court, so he used that as his business address. In the blank marked “occupation” he wrote the word justice. The person at the counter examined this answer, frowned, and said, “Justice? Justice! Well, I guess that’s all right. Last week a fellow wrote peace.”

Each of us should pursue the occupation of “peace.” But what is peace, and how do we seek it?

Many think of peace as the absence of war. Everyone wants that kind of peace. Songs celebrate it, and bumper stickers proclaim it.

Many good people promote peace by opposing war. They advocate laws or treaties to abolish war, to require disarmament, or to reduce armed forces.

Those methods may reduce the likelihood or the costs of war. But opposition to war cannot ensure peace, because peace is more than the absence of war.

For over fifty years, I have heard the leaders of this Church preach that peace can only come through the gospel of Jesus Christ. I am coming to understand why.

The peace the gospel brings is not just the absence of war. It is the opposite of war. Gospel peace is the opposite of any conflict, armed or unarmed. It is the opposite of national or ethnic hostilities, of civil or family strife.

In the midst of World War I, President Joseph F. Smith declared:

“For years it has been held that peace comes only by preparation for war; the present conflict should prove that peace comes only by preparing for peace, through training the people in righteousness and justice, and selecting rulers who respect the righteous will of the people. …

“There is only one thing that can bring peace into the world. It is the adoption of the gospel of Jesus Christ, rightly understood, obeyed and practiced by rulers and people alike” (Improvement Era, Sept. 1914, pp. 1074–75).

A generation later, during the savage hostilities of World War II, President David O. McKay declared:

“Peace will come and be maintained only through the triumph of the principles of peace, and by the consequent subjection of the enemies of peace, which are hatred, envy, ill-gotten gain, the exercise of unrighteous dominion of men. Yielding to these evils brings misery to the individual, unhappiness to the home, war among nations” (Gospel Ideals [Salt Lake City: Deseret Book Co., 1953], p. 280).

Such has been the message of the prophets in all ages. Referring to the first families of the earth, Moses wrote, “And in those days Satan had great dominion among men, and raged in their hearts; and from thenceforth came wars and bloodshed” (Moses 6:15).

In his own day, Moses gave the Lord’s promise to the children of Israel: “If ye walk in my statutes, and keep my commandments, … I will give peace in the land, … neither shall the sword go through your land” (Lev. 26:3, 6).

Throughout the Book of Mormon, the Lord declares, “Inasmuch as ye shall keep my commandments ye shall prosper in the land” (2 Ne. 1:20).

As we seek to understand the causes of wars, persecutions, and civil strife, we can see that they are almost always rooted in wickedness.

The mass-murders of the twentieth century are among the bloodiest crimes ever committed against humanity. We can hardly comprehend the magnitude of the Nazi holocaust murders of over five million Jews in Europe, Stalin’s purges and labor camps that killed five to ten million in the Soviet Union, and the two to three million noncombatants who were killed or who died of hunger in the Biafran War (see Isidor Walliman and Michael N. Dobkowski, eds., Genocide and the Modern Age [New York: Greenwood Press, 1987], p. 46; The Nation, 6 Mar. 1989, p. 294, 7/14 Aug. 1989, p. 154).

All of these slaughters, and others like them, were rooted in the ancient wickedness Satan taught—that a man could murder to get gain (see Moses 5:30–31). The mass-murderers of this century killed to acquire property and to secure power over others.

Through the prophet Moses, the Lord God of Israel commanded:

“Thou shalt not kill.

“Thou shalt not commit adultery.

“Thou shalt not steal.

“Thou shalt not bear false witness. …

“Thou shalt not covet” (Ex. 20:13–17).

Obedience to these commandments, which are the bedrock moral foundation for all Christians and Jews, would have prevented the greatest tragedies of this century.

We still live in a time of turmoil. There are wars between some nations, armed conflicts within others, and violent controversies in most. People are killed every day in some places, and hatred is practiced in many more. Peace is a victim everywhere.

If only we could heed the call of the Lord God of Israel, “Come unto me all ye ends of the earth” (2 Ne. 26:25). As the Book of Mormon teaches, he has created all flesh, “and the one being is as precious in his sight as the other” (Jacob 2:21). He has given salvation “free for all men” (2 Ne. 26:27), and “all men are privileged the one like unto the other, and none are forbidden” (v. 28).

“And he inviteth [all men] to come unto him and partake of his goodness; and he denieth none that come unto him, black and white, bond and free, male and female; and he remembereth the heathen; and all are alike unto God” (2 Ne. 26:33).

The blessings of the gospel are universal, and so is the formula for peace: keep the commandments of God. War and conflict are the result of wickedness; peace is the product of righteousness.

During the past year we have seen revolutionary changes in the governments of many nations. We are gratified that in most nations these changes have been accomplished without war or bloodshed. Nevertheless, we are far from securing peace in these nations or in any others throughout the world.

Many take comfort from the Old Testament prophecy that nations will “beat their swords into plowshares, and their spears into pruninghooks” (Micah 4:3). But this prophecy only applies to that time of peace which follows the time when the God of Jacob “will teach us of his ways, and we will walk in his paths” (Micah 4:2).

For now, we have wars and conflicts, and everywhere they are rooted in violations of the commandments of God.

The leaders of some nations have systematically murdered their opposition.

Persons in power in some nations have stolen public and private property so they could live in luxury. At the same time, they have neglected the most basic needs of the hungry and homeless among their people.

Some private citizens have promoted poverty by stealing, by corrupting public officials, and by oppressing the poor and defenseless.

Just across the borders of some nations are the wretched camps of refugees whose suffering circumstances are also traceable to man’s inability to keep the commandments of God.

The moral climate in some nations is reminiscent of the prophet Ezekiel’s description of “the bloody city” of Jerusalem:

“Her princes in the midst thereof are like wolves ravening the prey, to shed blood, and to destroy souls, to get dishonest gain. …

“The people of the land have used oppression, and exercised robbery, and have vexed the poor and needy” (Ezek. 22:27, 29).

Democracy does not ensure peace. When a nation is governed according to the voice of its people, its actions will mirror the righteousness or wickedness of its people.

We cannot have peace among nations without achieving general righteousness among the people who comprise them. Elder John A. Widtsoe said:

“The only way to build a peaceful community is to build men and women who are lovers and makers of peace. Each individual, by that doctrine of Christ and His Church, holds in his own hands the peace of the world.

“That makes me responsible for the peace of the world, and makes you individually responsible for the peace of the world. The responsibility cannot be shifted to someone else. It cannot be placed upon the shoulders of Congress or Parliament, or any other organization of men with governing authority” (in Conference Report, Oct. 1943, p. 113).

If citizens do not have a basic goodness to govern their actions toward one another, we can never achieve peace in the world. One nation’s greed, hatred, or desire for power over another is simply a reflection of the greeds, hatreds, and selfish desires of individuals within that nation.

Conversely, each citizen furthers the cause of world peace when he or she keeps the commandments of God and lives at peace with family and neighbors. Such citizens are living the prayer expressed in the words of a popular song, “Let there be peace on earth, and let it begin with me” (Sy Miller and Jill Jackson, “Let There Be Peace on Earth”).

The Savior and his Apostles had no program for world peace other than individual righteousness. They mounted no opposition to the rule of Rome or to the regime of its local tyrants. They preached individual righteousness and taught that the children of God should love their enemies (see Matt. 5:44) and “live peaceably with all men” (Rom. 12:18).

Recent history reminds us that people who continue to hate one another after a war will have another war, whereas the victor and vanquished who forgive one another will share peace and prosperity.

Our Church members demonstrated the healing and pacifying power of love in their shipment of food and clothing to relieve the suffering of the German Saints just after World War II. U.S. President Harry S. Truman was amazed when President George Albert Smith told him the supplies would not be sold. “You don’t mean you are going to give it to them?” he exclaimed.

President Smith replied simply, “They are our brothers and sisters and are in distress” (in Edward L. Kimball and Andrew E. Kimball, Jr., Spencer W. Kimball [Salt Lake City: Bookcraft, 1977], p. 222).

A few months later, Elder Ezra Taft Benson saw a German member in tears as he ran his fingers through a container of cracked wheat and whispered, “Brother Benson, it is hard for me to believe that people who have never seen us could do so much for us” (in Sheri L. Dew, Ezra Taft Benson [Salt Lake City: Deseret Book Co., 1987], p. 219).

What can one person do to promote world peace? The answer is simple: keep God’s commandments and serve his children.

A bishop who seeks to heal a troubled marriage or resolve a personal controversy is working for peace. So is a victim of abuse who is conscientiously working on the long process of forgiving the transgressor.

Young men and women contribute to peace when they forgo the temporary pleasure of self-gratifying activities and involve themselves in service projects and other acts of kindness.

The most powerful workers for peace may be faithful mothers and fathers. Some of the most terrible crimes committed against humanity are the acts of persons who have been scarred and twisted by the sins of others—often their own parents or others who cared for them. Parents who lovingly care for their own children or shelter foster children and raise them in righteousness are working for peace. So are parents who teach their children in the way King Benjamin counseled, to forgo conflicts and “to love one another, and to serve one another” (Mosiah 4:15).

Persons who seek to reduce human suffering and persons who work to promote understanding among different peoples are also important workers for peace.

A personal act of kindness or reconciliation also has an impact for peace. Lincoln’s biographer described such an act. A Union officer applied to his commander-in-chief for permission to leave his regiment to attend to the burial of his wife. Lincoln gruffly refused. Another battle was imminent, and every officer was needed. The next morning President Lincoln reconsidered and granted the request. He went to the room of the grieving man, took his hand, and said:

“My dear Colonel, I was a brute last night. I have no excuse to offer. I was weary to the last extent; but I had no right to treat a man with rudeness who had offered his life for his country, much more a man who came to me in great affliction. I have had a regretful night, and come now to beg your forgiveness” (Carl Sandburg, Abraham Lincoln, The War Years, 4 vols. [New York: Harcourt, Brace, and Co., 1939], 1:514).

Our missionaries, young men and women and older couples, are workers for world peace. So are the faithful souls who support them.

Like the church that sends them forth, our missionaries have no political agenda and no specific program for disarmament or reduction of forces. They circulate no petitions, advocate no legislation, support no candidates. They are the Lord’s servants, and his program for world peace depends on righteousness, not rhetoric. His methods involve repentance and reformation, not placards and picketing.

By preaching righteousness, our missionaries seek to treat the causes of war. They preach repentance from personal corruption, greed, and oppression because only by individual reformation can we overcome corruption and oppression by groups or nations. By inviting all to repent and come unto Christ, our missionaries are working for peace in this world by changing the hearts and behavior of individual men and women.

In The Church of Jesus Christ of Latter-day Saints, we follow the formula prescribed by the prophet-king Benjamin. He taught that those who receive a remission of their sins through the atonement of Christ are filled with the love of God and the knowledge of that which is just and true. That kind of person “will not have a mind to injure one another, but to live peaceably” with all people (Mosiah 4:13).

That is our method, and salvation and peace for all mankind is our goal.

Jesus Christ is our Savior. He has taught us the way to live. If we follow him and have goodwill toward all men, we can have peace on earth.

May God bless all of us in that great effort, I pray, in the name of Jesus Christ, amen.

\newpage
\settalk{Witnesses of Christ}
\setspeaker{Dallin H. Oaks}
\settalkdate{October 1990}
\talkheading{Witnesses of Christ}{Dallin H. Oaks}

A few months ago, I received a letter from a Church member who posed an unusual question: “Do I have a right to bear testimony of the Savior? Or is that the sole prerogative of the Twelve?” In response, I will share some thoughts on why every member of this Church should bear witness and testimony of Jesus Christ.

In the beginning, God commanded Adam, “Thou shalt do all that thou doest in the name of the Son, and thou shalt repent and call upon God in the name of the Son forevermore” (Moses 5:8). Then the Holy Ghost, “which beareth record of the Father and the Son,” came upon Adam and Eve, and they “blessed the name of God, and they made all things known unto their sons and their daughters” (Moses 5:9, 12).

Later, Enoch described how God had taught Adam that all must repent and be baptized in the name of Jesus Christ, whose atoning sacrifice made possible the forgiveness of sins, and that they must teach these things to their children (see Moses 6:52–59).

And so our first parents established the pattern, receiving a testimony from the Holy Ghost and then bearing witness of the Father and the Son to those around them.

The prophet Nephi described the ordinance of baptism as an occasion when persons would witness unto the Father that they were willing to take upon them the name of Christ (see 2 Ne. 31:13). Similarly, the Lord has specified that those who desire to be baptized in this dispensation should “come forth with broken hearts and contrite spirits, and witness before the church that they … are willing to take upon them the name of Jesus Christ” (D\&C 20:37; see also Moro. 6:3). We renew that promise when we partake of the sacrament (see D\&C 20:77; Moro. 4:3).

We also witness of Christ by our membership in the Church that bears his name (see 3 Ne. 27:7; D\&C 115:4).

We are commanded to pray unto the Father in the name of his Son, Jesus Christ (see 3 Ne. 18:19, 21, 23; see also Moses 5:8), and to do “all things … in the name of Christ” (D\&C 46:31).

If we follow these commandments, we serve as witnesses of Jesus Christ through our baptism, our membership in his Church, our partaking of the sacrament, and our prayers and other actions in his name.

But our duty to be witnesses of Jesus Christ requires more than this, and I fear that some of us fall short. Latter-day Saints can become so preoccupied with our own agendas that we can forget to witness and testify of Christ.

I quote from a recent letter I received from a member in the United States. He described what he heard in his fast and testimony meeting:

“I sat and listened to seventeen testimonies and never heard Jesus mentioned or referred to in any way. I thought I might be in [some other denomination], but I supposed not because there were no references to God, either. …

“The following Sunday, I again attended church. I sat through a priesthood lesson, a Gospel Doctrine lesson, and seven sacrament meeting speakers and never once heard the name of Jesus or any reference to him.”

Perhaps that description is exaggerated. Surely, it is exceptional. I quote it because it provides a vivid reminder for all of us.

In answer to the question, “What are the fundamental principles of your religion?” the Prophet Joseph Smith said, “The fundamental principles of our religion are the testimony of the Apostles and Prophets, concerning Jesus Christ, that He died, was buried, and rose again the third day, and ascended into heaven; and all other things which pertain to our religion are only appendages to it” (Teachings of the Prophet Joseph Smith, sel. Joseph Fielding Smith [Salt Lake City: Deseret Book Co., 1938], p. 121).

When Alma spoke to a group of prospective members at the Waters of Mormon, he instructed them on the duties of those who were “desirous to come into the fold of God, and to be called his people” (Mosiah 18:8). One of those duties was “to stand as witnesses of God at all times and in all things, and in all places that ye may be in, even until death” (Mosiah 18:9).

How do members become witnesses? The original Apostles were eyewitnesses to the ministry and resurrection of the Savior (see Acts 10:39–41). He told them, “Ye shall be witnesses unto me both in Jerusalem, and in all Judaea, and in Samaria, and unto the uttermost part of the earth” (Acts 1:8; see also Acts 10:42–43). However, he cautioned them that their witnessing would be after they had received the Holy Ghost (see Acts 1:8; see also Luke 24:49).

An eyewitness was not enough. Even the witness and testimony of the original Apostles had to be rooted in the testimony of the Holy Ghost. A prophet has told us that the witness of the Holy Ghost makes an impression on our soul that is more significant than “a visitation of an angel” (Joseph Fielding Smith, Doctrines of Salvation, comp. Bruce R. McConkie, 3 vols. [Salt Lake City: Bookcraft, 1954–56], 1:44). And the Bible shows that when we testify on the basis of this witness, the Holy Ghost testifies to those who hear our words (see Acts 2; 10:44–47).

When Peter and the other Apostles were brought before the civil authorities, he testified that Jesus Christ was “a Prince and a Saviour, for to give repentance to Israel, and forgiveness of sins” (Acts 5:31). Then Peter added, “And we are his witnesses of these things; and so is also the Holy Ghost, whom God hath given to them that obey him” (v. 32). The mission of the Holy Ghost is to witness of the Father and the Son (see 2 Ne. 31:18; 3 Ne. 28:11; D\&C 20:27). Consequently, everyone who has received the witness of the Holy Ghost has a duty to share that testimony with others.

Apostles have the calling and ordination to be special witnesses of the name of Christ in all the world (see D\&C 107:23), but the duty to witness and testify of Christ at all times and in all places applies to every member of the Church who has received the testimony of the Holy Ghost.

The book of Luke records two examples of this. In obedience to the law of Moses, Joseph and Mary brought the infant Jesus to the temple at Jerusalem after forty days, to present him to the Lord. There, two aged and spiritual temple workers received a witness of his identity and testified of him. Simeon, who had known by revelation from the Holy Ghost that he should not taste of death until he had seen the Messiah, took the infant in his arms and testified to his divine mission (see Luke 2:25–35). Anna, whom the scripture called “a prophetess” (Luke 2:36), recognized the Messiah “and spake of him to all them that looked for redemption in Jerusalem” (Luke 2:38).

Anna and Simeon were eyewitnesses to the infant, but, just like the Apostles, their knowledge of his divine mission came through the witness of the Holy Ghost. “The testimony of Jesus is the spirit of prophecy” (Rev. 19:10). Therefore, we can properly say that when each received this witness, Simeon was a prophet and Anna was a prophetess. Each then fulfilled the prophetic duty to testify to those around them. As Peter said, “To [Christ] give all the prophets witness” (Acts 10:43). This was what Moses meant when he expressed the wish “that all the Lord’s people were prophets, and that the Lord would put his spirit upon them!” (Num. 11:29).

The scriptures describe other occasions when ordinary members—men and women—bore witness of Christ. The Book of Mormon tells of King Lamoni and his queen, who testified of their Redeemer (see Alma 19). The Bible describes the witness of the Holy Ghost coming upon the kinsmen and friends of Cornelius, who were then heard to “magnify God” (Acts 10:24, 46).

Our scriptural duty to witness of the Savior and to testify of his divine Sonship has been affirmed by the prophets in our own day.

We are told that the commandments are given and the gospel is proclaimed that every person “might speak in the name of God the Lord, even the Savior of the world” (D\&C 1:20).

Spiritual gifts come by the power of the Holy Ghost, that all the faithful may be benefited. One of these gifts is “to know that Jesus Christ is the Son of God, and that he was crucified for the sins of the world” (D\&C 46:13). Those who receive that gift have the duty to testify of it. We know this because immediately after describing the gift of knowing that Jesus Christ is the Son of God, the Lord says, “To others it is given to believe on their words, that they also might have eternal life if they continue faithful” (D\&C 46:14; see also 3 Ne. 19:28). Those who have the gift to know must give their witness so that those who have the gift to believe on their words can enjoy the benefit of that gift.

Speaking to some of the earliest missionaries of this dispensation, the Lord said: “But with some I am not well pleased, for they will not open their mouths, but they hide the talent which I have given unto them, because of the fear of man. Wo unto such, for mine anger is kindled against them” (D\&C 60:2).

In contrast, the Lord gave this great promise to those who were valiant in bearing testimony: “For I will forgive you of your sins with this commandment—that you remain steadfast … in bearing testimony to all the world of those things which are communicated unto you” (D\&C 84:61).

This caution and promise were directed specifically to missionaries, but other scriptures suggest that they apply to other members as well.

In his vision of the spirits of the dead, President Joseph F. Smith described “the spirits of the just” as those “who had been faithful in the testimony of Jesus while they lived in mortality” (D\&C 138:12).

In contrast, in his vision of the three degrees of glory, the Prophet Joseph Smith described those souls who go to the terrestrial kingdom as the “honorable men of the earth,” who were “not valiant in the testimony of Jesus” (D\&C 76:75, 79).

What does it mean to be “valiant in the testimony of Jesus”? Surely this includes keeping his commandments and serving him. But wouldn’t it also include bearing witness of Jesus Christ, our Savior and our Redeemer, to believers and nonbelievers alike? As the Apostle Peter taught the Saints of his day, we should “sanctify the Lord God in [our] hearts: and be ready always to give an answer to every man that asketh [us] a reason of the hope that is in [us]” (1 Pet. 3:15).

All of us need to be valiant in the testimony of Jesus. As believers in Christ, we affirm the truth of Peter’s testimony in the name of Jesus Christ of Nazareth that “there is none other name under heaven given among men, whereby we must be saved” (Acts 4:12; see also D\&C 109:4). We know from modern revelation that we can come unto the Father only in his name (see D\&C 93:19). As the Book of Mormon teaches, salvation is “in and through the atoning blood of Christ, the Lord Omnipotent” (Mosiah 3:18; see also Moses 6:52, 59).

To those who are devoted to the Lord Jesus Christ, I say there has never been a greater need for us to profess our faith, privately and publicly.

When the gospel was first restored, the pulpits of this land were aflame with the testimony of Jesus, the divine Son of God and Savior of the world. True, the fulness of his doctrine and the power of his priesthood had been lost from the earth, but there were many good and honorable men and women who were valiant in their own testimonies of Jesus. Our earliest missionaries concentrated their message on the Restoration—the calling of the Prophet Joseph Smith and the restoring of priesthood authority—since they could assume that most of those they taught had a fundamental belief in Jesus Christ as our Savior.

Today, our missionaries cannot make that assumption. There are still many God-fearing people who testify to the divinity of Jesus Christ. But there are many more—even in the formal ranks of Christianity—who doubt his existence or deny his divinity. As I see the deterioration in religious faith that has happened in my own lifetime, I am convinced that we who are members of his Church need to be increasingly valiant in our testimony of Jesus.

Speaking almost twenty years ago, President Harold B. Lee said: “Fifty years ago or more, when I was a missionary, our greatest responsibility was to defend the great truth that the Prophet Joseph Smith was divinely called and inspired and that the Book of Mormon was indeed the word of God. But even at that time there were the unmistakable evidences that there was coming into the religious world actually a question about the Bible and about the divine calling of the Master, himself. Now, fifty years later, our greatest responsibility and anxiety is to defend the divine mission of our Lord and Master, Jesus Christ, for all about us, even among those who claim to be professors of the Christian faith, are those not willing to stand squarely in defense of the great truth that our Lord and Master, Jesus Christ, was indeed the Son of God” (address delivered at LDS Student Association Fireside, Utah State University, 10 Oct. 1971).

Our knowledge of the literal divinity, resurrection, and atonement of Jesus Christ is more certain and more distinctive with each passing year. That is one reason the Lord inspired his prophet, Ezra Taft Benson, to have us reemphasize our study and testimony of the Book of Mormon, whose mission is “the convincing of the Jew and Gentile that Jesus is the Christ, the Eternal God” (Book of Mormon title page).

We live in a time when too many who purport to be Christians have a cause that comes ahead of Christ. For example, a national magazine recently reported an innovation by a new bishop of a prominent Christian church. Their ministers have always consecrated the emblems of the flesh and blood of Jesus Christ in the name of the “Father, Son, and Holy Ghost.” However, in an effort to use what are called “nonsexist words,” this new bishop has begun to consecrate the eucharist in the name of the “Creator, Redeemer, and Sustainer” (“Fretful Murmur in the Cathedral,” Insight, 24 Apr. 1989, p. 47). Such trendy and expedient tampering with the Christian faith is illustrative of the extent to which some are unwilling to witness of Jesus Christ, the Son of God.

Such deliberate deviations are not likely to be made by faithful Latter-day Saints. However, we need to be on guard against careless omissions and oversights in our personal testimonies, in our formal instruction, and in our worship and funeral services.

In addition, each of us has many opportunities to proclaim our belief to friends and neighbors, fellow workers, and casual acquaintances. I hope we will take these opportunities to express our love for our Savior, our witness of his divine mission, and our determination to serve him.

If we do all of this, we can say, like the Apostle Paul, “I am not ashamed of the gospel of Christ: for it is the power of God unto salvation to every one that believeth” (Rom. 1:16).

And, we can say, like the prophet Nephi, “We talk of Christ, we rejoice in Christ, we preach of Christ, we prophesy of Christ, … that our children may know to what source they may look for a remission of their sins” (2 Ne. 25:26).

I testify of Jesus Christ, the Lord God of Israel, the light and life of the world, as I affirm the truth of his gospel, in the name of Jesus Christ, amen.

\newpage
\settalk{“Honour Thy Father and Thy Mother”}
\setspeaker{Dallin H. Oaks}
\settalkdate{April 1991}
\talkheading{“Honour Thy Father and Thy Mother”}{Dallin H. Oaks}

Since our conference last October, many have ridden roller coasters of emotion from the war in the Persian Gulf. Many Latter-day Saints had their lives changed by that conflict. In the military theater of operations, we had over 140 Latter-day Saint groups providing leadership, worship, and fellowship for members of the armed forces. At home, families were separated and in stress. We pay tribute to the Church leaders and members who shouldered extra burdens in looking after the families of our service people. They are still doing so. Stake, ward, quorum, and Relief Society leaders acted and are still acting in the best traditions of brotherly and sisterly service.

During this crisis, our hearts went out to those who were oppressed and in jeopardy on both sides of the conflict. Week after week, in the leading councils of the Church, in Church meetings everywhere, in our homes, and in public and private gatherings, we prayed for the well-being of those in uniform. We also prayed that the war would be short and that the numbers of dead and wounded would be as small as possible.

Our prayers were answered, and in this national period of thanksgiving designated by presidential proclamation, we join millions of religious people everywhere in prayers of thanks to a merciful Father in Heaven. We offer love and sympathy to the families of those who lost their lives. And we continue to pray that the leaders who preside over the peacekeeping process and the care and repatriation of prisoners and refugees will be wise and considerate and successful in binding up the wounds of war.

Thousands of years ago, on a mountain across the Arabian peninsula from the recent conflict, the Lord God of Israel gave his people ten commandments. The fifth commandment that the Lord gave through the prophet Moses was:

“Honour thy father and thy mother: that thy days may be long upon the land which the Lord thy God giveth thee” (Ex. 20:12).

This morning I wish to speak about honoring our parents and the aged among us.

The commandment to honor our parents has strands that run through the entire fabric of the gospel. It is inherent in our relationship to God our Father. It embraces the divine destiny of the children of God. This commandment relates to the government of the family, which is patterned after the government of heaven.

The commandment to honor our parents echoes the sacred spirit of family relationships in which—at their best—we have sublime expressions of heavenly love and care for one another. We sense the importance of these relationships when we realize that our greatest expressions of joy or pain in mortality come from the members of our families.

Other manifestations of this commandment include the bearing and care of children, the preparation of family histories, and efforts to see that the ordinances of eternity are performed for our departed ancestors.

The Savior reemphasized the importance of the fifth commandment during his ministry. He reminded the scribes and Pharisees that we are commanded to honor our father and our mother and that God had directed that whoever cursed father or mother should be put to death (see Lev. 20:9; Deut. 21:18–21; Matt. 15:4; Mark 7:10). In this day, failing to honor our parents is not a capital crime in any country of which I am aware. However, the divine direction to honor our father and our mother has never been revoked (see Mosiah 13:20; Matt. 19:19; Luke 18:20).

Like many scriptures, this commandment has multiple meanings.

To young people, honoring parents is appropriately understood to focus on obedience, respect, and emulation of righteous parents. The Apostle Paul illuminated that focus when he taught, “Children, obey your parents in all things [I believe he meant all righteous things]: for this is well pleasing unto the Lord” (Col. 3:20).

President Spencer W. Kimball combined the ideas of obedience and emulation in these words:

“If we truly honor [our parents], we will seek to emulate their best characteristics and to fulfill their highest aspirations for us. No gift purchased from a store can begin to match in value to parents some simple, sincere words of appreciation. Nothing we could give them would be more prized than righteous living for each youngster” (The Teachings of Spencer W. Kimball, ed. Edward L. Kimball [Salt Lake City: Bookcraft, 1982], p. 348).

Young people, if you honor your parents, you will love them, respect them, confide in them, be considerate of them, express appreciation for them, and demonstrate all of these things by following their counsel in righteousness and by obeying the commandments of God.

To persons whose parents are dead, honoring parents is likely to involve thoughts of family reunions, family histories, temple work, and commitment to the great causes in which departed parents spent their lives.

Middle-aged persons are likely to think of the commandment to honor our fathers and our mothers in terms of caring for aged parents. In a message given a year and a half ago, President Ezra Taft Benson encouraged families “to give their elderly parents and grandparents the love, care, and attention they deserve.” He said:

“Remember, parents and grandparents are our responsibility, and we are to care for them to the very best of our ability. When the elderly have no families to care for them, priesthood and Relief Society leaders should make every effort to meet their needs in the same loving way” (in Conference Report, Oct. 1989, p. 6; or Ensign, Nov. 1989, p. 6).

Recent years have seen great increases in the numbers and percentage of older people in our population. A recent study estimated that in another ten years, one-seventh of the population of the United States, about 35 million people, will be at least 65 years old. At that time, about five million citizens will be age 85 or older (see “Consumer Issues and the Elderly,” Deseret News, 7 May 1990, p. C1).

From time to time, Church leaders hear of grown children who seem to be good Latter-day Saints but are negligent or even maliciously indifferent in caring for their aged parents. Some have encouraged parents to distribute their property and then have put them away in institutions, sometimes with inadequate care and sometimes without regular visits and expressions of love from their children. I believe this was the kind of circumstance the Lord’s spokesman, the prophet Isaiah, thundered against when he commanded, “Hide not thyself from thine own flesh” (Isa. 58:7).

The best way to care for the aged is to preserve their independence as long as possible. President Benson explained:

“Even when parents become elderly, we ought to honor them by allowing them freedom of choice and the opportunity for independence as long as possible. Let us not take away from them choices which they can still make. Some parents are able to live and care for themselves well into their advancing years and would prefer to do so. Where they can, let them.

“If they become less able to live independently, then family, Church, and community resources may be needed to help them. When the elderly become unable to care for themselves, even with supplemental aid, care can be provided in the home of a family member when possible. Church and community resources may also be needed in this situation” (in Conference Report, Oct. 1989, pp. 6–7; or Ensign, Nov. 1989, p. 7).

When aged parents who are not able to live alone are invited to live with their children, this keeps them in the family circle and allows them to continue their close ties with all members of the family. When a parent lives with one child, the other children should make arrangements to share the burdens and blessings of this arrangement.

When it is not possible for parents to be cared for in the homes of their children, so that some type of institutional care is obtained, their children should remember that institutional care will generally focus on physical needs. Members of the family should make regular visits and contacts to provide the spiritual and emotional sustenance and the love that must continue in the family relationship for mortal life and throughout all eternity.

In some nations where our members reside, the obligation to care for aged parents is more keenly felt and more faithfully observed than in the United States. I saw this in Asia. But the care of aged parents is still a strongly felt obligation among most Americans. Six out of ten older persons questioned in a recent national survey had weekly personal visits with their children, and three-fourths of them talked on the telephone with their children at least weekly. Two-thirds of those surveyed expect to take care of their elderly parents. (See Deseret News, 7 May 1990, p. C1).

Latter-day Saints have a good record of caring for their aged parents and for older citizens generally. I have seen wonderful examples of this in my own family and among my LDS friends and associates. Many of our General Authorities and their companions have been exemplary in caring for their aged parents.

When I was a young boy in a small Utah town, I remember seeing my grandmother overseeing the provision of food, favors, transportation, and entertainment for a large group of elderly in the community. As a counselor in the stake Relief Society presidency, she was making preparations for “Old Folks Day.”

Most of you have never heard of Old Folks Day. It was a unique Utah Mormon institution. It began in 1875, when Charles R. Savage, the pioneer photographer, persuaded Presiding Bishop Edward Hunter to declare a day for honoring what we now call senior citizens. The first Old Folks Day transported guests by rail to an outing at Saltair, west of Salt Lake City. A monument honoring that celebration and its founder stands on the southeast corner of Temple Square.

The annual Old Folks Day celebrations were held in nearly every community in Utah. Travel, refreshments, and entertainment were given to all citizens seventy years of age and older. Although this holiday was conducted by the leaders and members of this Church, it was stipulated that “there are none to be excluded because of their religion, and the oldest guest present is the special guest of the occasion whether they be white or black or whatever the complexion of their religious belief” (quoted in Joseph Heinerman, “The Old Folks Day: A Unique Utah Tradition,” Utah Historical Quarterly, Spring 1985, p. 158).

The committee directing these celebrations was dissolved in 1970, and the responsibility for honoring those who had come to be called senior citizens was passed to the stake presidents of the Church. Since that time we have had further increases in the number of senior citizens in our midst, but perhaps not significant increases in the amount of honor accorded them. Fortunately, the advances in medical science that have produced increased longevity have also increased our senior citizens’ effective participation in church, community, business, and social events. But the need for honor, especially for our fathers and our mothers, is undiminished.

The fifth commandment is often referred to as the first commandment with a promise: “Honour thy father and thy mother: that thy days may be long upon the land which the Lord thy God giveth thee” (Ex. 20:12). I have wondered about the relationship between the commandment and the promise. How could honoring our parents increase our longevity?

During almost forty years of marriage, I have observed something that provides at least a partial explanation of how this promise is fulfilled.

In the early days of our marriage, I spent many happy hours in the home of my wife’s parents, Charles and True Dixon. There I met June’s maternal grandmother, Adelaide White Call. Then a widow about eighty-five years of age, Grandma Call was a survivor of what older people called “the exodus.” She had been among those valiant Latter-day Saints expelled from northern Mexico in 1912. Now her sons and daughters were living throughout the United States. In her later years, they helped her locate in Utah County, near June’s parents.

During my visits, I saw the gentleness and love and concern with which the Call children and their companions looked after this older parent. They visited her frequently. My wife’s mother looked in on her every day and often had her in their home. They made her part of every occasion in which she desired to participate, and they gave her every consideration and respect. They cared for her every need when she was ill. Surely, I said to myself, these Call children honor their mother.

It has been about forty years since I saw that honor given. Now I see its effects. I see June and her brother and sisters honoring their mother as they saw their mother honoring her own mother. Fortunately, True Dixon is blessed with good health and vigor and has no present need for the kind of care her mother required. Still, her children are attentive. There are frequent visits and phone calls and invitations that include her in all the family activities she desires. I believe her days will be longer upon the land because of the attentiveness and companionship of her children, who learned the way to honor a parent by seeing how their own mother honored hers.

I am grateful for this example and for this principle, especially when I anticipate the effect of having our daughters and sons observe how their mother honors her mother. I am sure that when the time comes, my own companion’s days will be lengthened upon the land because of the care her children will give to her because of the example she has set for them. A worthy example repeats itself from generation to generation. Truly, righteousness is a beacon and a worthy act is its own reward. As the Lord said, “He who doeth the works of righteousness shall receive his reward, even peace in this world, and eternal life in the world to come” (D\&C 59:23).

In time to come, each of us will be judged by the Lord God of Israel, who commanded us to honor our fathers and our mothers. I pray that each of us will conduct ourselves toward our parents in such a way that we will be guiltless before God at that day. In the name of Jesus Christ, amen.

\newpage
\settalk{Joy and Mercy}
\setspeaker{Dallin H. Oaks}
\settalkdate{October 1991}
\talkheading{Joy and Mercy}{Dallin H. Oaks}

One of the greatest of all God’s revelations is Father Lehi’s teaching that “men are, that they might have joy” (2 Ne. 2:25). Joy is more than happiness. Joy is the ultimate sensation of well-being. It comes from being complete and in harmony with our Creator and his eternal laws.

The opposite of joy is misery. Misery is more than unhappiness, sorrow, or suffering. Misery is the ultimate state of disharmony with God and his laws.

Joy and misery are eternal emotions whose ultimate extent we are not likely to experience in mortality. In this life we have some mortal simulations, which we call happiness or pleasure and unhappiness or pain. In the midst of these emotions is suffering. Some suffering comes from our own sins or those of others, but much suffering is simply an inevitable part of the mortal condition, like an accidental injury.

An example of our two emotional extremes occurred two years ago. As part of an outing, a group of Latter-day Saint Boy Scouts entered an abandoned mine in the mountains not far from here. Somehow, young Joshua Dennis was separated from the group and became lost in the mine. Anyone who has ever lost something valuable will remember that terrible feeling. The pain is most extreme when we have lost a loved one. Joshua’s family and friends feared their loss might be permanent.

Search efforts were organized. For days, many good people dropped everything they were doing to search for the one who was lost. Many shared the pain of his loss. Then, miraculously, he was found. Prayers were answered, and the mercy of a loving Heavenly Father was manifest in the happiness of family and friends reunited with the one who was lost. The pain of loss turned to the overwhelming joy of reunion. (See "Hidden Treasure," Ensign, Aug. 1991, pp. 30–35.)

Few experiences illustrate mortal feelings of joy better than the recovery of something precious we fear we have lost. The scriptures illustrate this experience. In the parable of the lost sheep, the shepherd invited his friends to rejoice with him, for he had “found [his] sheep which was lost” (Luke 15:6). “Likewise,” Jesus explained, “joy shall be in heaven over one sinner that repenteth” (Luke 15:7). In another parable, a father rejoiced in the return of a prodigal son, explaining, “For this thy brother was dead, and is alive again; and was lost, and is found” (Luke 15:32).

These experiences are symbolic of our eternal journey. The Fall has separated us from our heavenly home. We must choose which way we will go. Satan, who was separated from God and permanently lost, would like to see our separation become permanent also. Father Lehi taught that Satan’s purpose is to make men miserable. “Because he had fallen from heaven, and had become miserable forever, [Satan] sought also the misery of all mankind” (2 Ne. 2:18; see also 2 Ne. 2:27). Those who yield to his enticings are on the way to the same miserable destiny as he. Shut out from the presence of God, they will be in “a state of misery and endless torment” (Mosiah 3:25; see also 2 Ne. 2:5). As the Lord said about the wicked, “misery shall be their doom” (Moses 7:37).

Much of the misery encouraged by Satan comes from losses. Satan experienced that kind of misery when he lost his first estate. Now he tries to inflict similar losses on those who have proceeded to mortality, the second estate. Satan encourages a loss of virtue, a loss of integrity, a loss of reputation, a loss of ideals, a loss of wholesome associations, and even a loss of life.

In contrast, our Heavenly Father created us to resist and to overcome such losses, to be whole, to have joy. He wants us to return to him, and he has provided a way for that reunion to be achieved. No wonder we say that his gospel is “a voice of mercy from heaven; … a voice of gladness for the living and the dead; glad tidings of great joy” (D\&C 128:19).

Our Creator wants us to be happy in this life. The Prophet Joseph Smith taught that “happiness is the object and design of our existence” (Teachings of the Prophet Joseph Smith, sel. Joseph Fielding Smith [Salt Lake City: Deseret Book Co., 1938], p. 255). The things of the earth were created for our happiness. Modern revelation tells us that “all things which come of the earth … are made for the benefit and the use of man, both to please the eye and to gladden the heart” (D\&C 59:18). Even on the Sabbath, a day of worship, the Lord expects us to have “a glad heart and a cheerful countenance” (D\&C 59:15). A prophet has called the gospel plan “the great plan of happiness” (Alma 42:8).

How do we seek happiness or joy? King Benjamin invited his people to consider “the blessed and happy state of those that keep the commandments of God. For behold,” he said, “they are blessed in all things, both temporal and spiritual; and if they hold out faithful to the end they are received into heaven … [to] dwell with God in a state of never-ending happiness (Mosiah 2:41).

Father Lehi explained that if Adam and Eve had not gone through the process we call the Fall, they would have remained forever in their initial state of innocence, “having no joy, for they knew no misery; doing no good, for they knew no sin” (2 Ne. 2:23).

Our first parents recognized this principle. When the Holy Ghost fell upon them and bore record of the Father and the Son, Adam blessed the name of the Lord, declaring that “because of my transgression my eyes are opened, and in this life I shall have joy” (Moses 5:10). With inspired insight, Eve explained the purpose of life and the source of joy:

“Were it not for our transgression we never should have had seed, and never should have known good and evil, and the joy of our redemption, and the eternal life which God giveth unto all the obedient” (Moses 5:11).

When I think of happiness or joy in this life, I begin with some experiences that are simple and basic. I see the expression on the face of a one-year-old taking those first steps. I remember a two-year-old immersed in a soft ice-cream cone. I think of a child loving a puppy or a kitten.

If the more mature have not dulled their physical or spiritual sensitivities by excess or disuse, they can also experience joy in what is simple and basic—in flowers and other growing things, in a sunrise or sunset or other beauties of nature, in wholesome companionship.

Another source of happiness and mortal joy is the accomplishment of worthy goals, simple things like physical exercise or more complex goals like the completion of an arduous task.

Other goals have eternal significance. Their completion produces joy in this life and the promise of eternal joy in the world to come. A few weeks ago Elder Rex D. Pinegar and I saw an example of this as we visited the beautifully renovated temple in Cardston, Alberta. In the brides’ room stood a lovely young woman in her wedding dress. She was there because she had made righteous choices. The look in her eyes and the expression on her face were a perfect embodiment of joy.

But despite all we can do, we cannot have a fulness of joy in this world or through our own efforts (see D\&C 101:36). Only in Christ can our joy be full. This is why the angel proclaimed: “I bring you good tidings of great joy, which shall be to all people.

“For unto you is born this day … a Saviour, which is Christ the Lord” (Luke 2:10–11).

We are able to have a fulness of joy only when spirit and body are inseparably connected in the glorious resurrection to celestial glory (see D\&C 93:33; 76:50–70). That joy, of course, comes only through the mercy of the Holy Messiah, whose resurrection broke the bands of death and whose atonement unlocks the reservoir of mercy by which we can be cleansed of our sins and come into the presence of God to receive the fulness of the Father.

We joyously proclaim that “there is no flesh that can dwell in the presence of God, save it be through the merits, and mercy, and grace of the Holy Messiah” (2 Ne. 2:8). God’s mercy is the only source of the ultimate and eternal joy, which restores every loss, dries every tear, and erases every pain. Eternal joy transcends all suffering. In this life and in the life to come, that joy comes about through the Resurrection and the remission of sins.

Remembering his early life, Alma told how his sins caused him to be “tormented with the pains of hell” (Alma 36:13). He said he was “racked, even with the pains of a damned soul” (Alma 36:16). At length he remembered his father’s teachings about the Savior, who would atone for the sins of the world. He cried out, “O Jesus, thou Son of God, have mercy on me” (Alma 36:18).

The miracle of forgiveness was wrought in his life, and the bitter pain of sin was replaced by the sweet joy of redemption. In his words, “Oh, what joy, and what marvelous light I did behold; yea, my soul was filled with joy as exceeding as was my pain!” (Alma 36:20).

Alma learned the eternal truth that the pain and misery that come from sin can only be erased by repentance. Physical pain ends with death. Spiritual pain, or misery, is everlasting, unless we repent.

The Book of Mormon tells of an entire people who received a remission of their sins and “were filled with joy” and “peace of conscience” (Mosiah 4:3). King Benjamin reminded them that they had “received a remission of [their] sins, which causeth such exceedingly great joy in [their] souls” (Mosiah 4:11).

The joy that follows the remission of sins comes from the Spirit of the Lord (see Mosiah 4:3, 20). It is a fulfillment of the Lord’s promise that “I will impart unto you of my Spirit, … which shall fill your soul with joy” (D\&C 11:13). As the Apostle Paul taught, “The fruit of the Spirit is love, joy, peace” (Gal. 5:22). It comes in the same way to everyone—to rich and poor, to the prominent and the obscure. In conferring his greatest gift of mercy through the Atonement, God is not a respecter of persons.

In contrast, those who yield to the enticing of Satan may, as the scripture says, “enjoy the pleasures of sin for a season” (Heb. 11:25), but that kind of pleasure can never lead to lasting happiness or eternal joy. The spirit and influence of Satan bring feelings of confusion, contention, darkness, disorder, anger, hate, and misery.

Persons who pursue Satan’s way are certain to have Satan’s misery. Unless they repent they will “remain with the father of lies, in misery, like unto himself” (2 Ne. 9:9). As Alma said to his wayward son, “Wickedness never was happiness” (Alma 41:10).

Our personal experiences and almost every newscast and newspaper remind us of the unhappiness and pain suffered on this earth. Some of this is traceable to sin. Many letters sent to Church headquarters describe the pain people inflict upon one another, often within those family relationships that should be the source of life’s greatest joy.

We have seen the pain of parents when their children stray from the path of truth. We have seen the pain that comes when a wife abandons her family and flies off to seek pleasure in forbidden paths, like a moth to a flame. We have seen the pain that comes in those more numerous instances when a husband abandons his wife and children in his self-destructive search for pleasure.

Other unhappiness results from the lust for power and wealth. A few weeks ago some of my Brethren and I visited a country that had been ruled for decades by an oppressive dictatorship, recently overthrown. We saw the conditions produced by rulers who had gratified their lusts at the expense of their people. Their palaces stood in shameful contrast to the workers’ apartment houses built without indoor plumbing. Everywhere we saw the ugly effects of the neglect of public services. Truly, the Proverbs say, “When the wicked … rule, the people mourn” (Prov. 29:2).

The consequences of wickedness are far-reaching. They continue from generation to generation. The pitiful victims of physical and sexual and emotional abuse are more susceptible to Satan’s manipulations. They are more likely to perpetuate these ugly transgressions within their own family relationships. Like highly contagious physical ailments, the spiritual diseases of lust and greed and corruption spread their evil effects, and Satan rejoices in the unhappiness and pain they cause in each succeeding generation.

Brothers and sisters, old and young, I plead with each of you to remember that wickedness never was happiness and that sin leads to misery. Young people, do not seek happiness in the glittering but shallow things of the world. We cannot achieve lasting happiness by pursuing the wrong things. Someone once said, “You can never get enough of what you don’t need because what you don’t need won’t satisfy you.”

Young and old, turn your eyes and your hearts away from the deceptive messages of the media. There is no happiness in alcohol or drugs, only enslavement. There is no happiness in violence, only pain and sorrow. There is no happiness in sexual relations and physical familiarities outside the bonds of marriage, only degradation and increased momentum along the way to spiritual death.

There is no lasting happiness in what we possess. Happiness and joy come from what a person is, not from what he or she possesses or appears to be. Youth, hold fast to your standards. Study and use that saving pamphlet For the Strength of Youth.

Righteousness fosters righteousness. The effects of righteous examples are felt for generations to come. Good parenting produces youth who make good parents. Just as many of us have been strengthened by the noble examples of our pioneering ancestors in many lands, so the righteous choices and sacrifices of our day can bless our families and our friends and our nations for all the years to come.

We appreciate the work of our Church News and our Church magazines, which share worthy and wholesome examples for the strengthening of all. We are saddened by the negative examples highlighted in the media and in many other public communications. Examples are powerful. We should use them to build up, not to tear down.

I testify to you of the lasting happiness and ultimate joy of those who exercise faith in God and keep his commandments. I urge each of you to seek the joy that comes from keeping the commandments of God and exerting a righteous influence for the benefit of those we love.

I testify to the truth of God’s promise that the faithful “shall enter into the joy of [the] Lord, and shall inherit eternal life” (D\&C 51:19). I testify of God the Eternal Father and his Son Jesus Christ, the authors of the great plan of happiness and mercy. I pray that we may do those things that are required for happiness and peace in this life and for joy and eternal life in the world to come, in the name of Jesus Christ, amen.

\newpage
\settalk{The Relief Society and the Church}
\setspeaker{Dallin H. Oaks}
\settalkdate{April 1992}
\talkheading{The Relief Society and the Church}{Dallin H. Oaks}

This year we are celebrating the 150th anniversary of the Relief Society, organized in Nauvoo, Illinois, on March 17, 1842. Last month’s anniversary program was carried by satellite to most continents of the world. Books are being published to review the history and celebrate the worldwide sisterhood of Relief Society. Ward and stake Relief Societies are celebrating through service in their local communities. Far-reaching efforts to promote literacy will be formally announced later this year.

We are grateful for the effective leadership of President Elaine L. Jack and her counselors and board, who are directing this celebration, and for the earlier leaders and workers whose accomplishments we celebrate.

The Relief Society has great significance for every member of the Church. All of us have been blessed through the example and service of its members.

I am the beneficiary of at least four different generations of Relief Society service: my grandmother, my mother, my wife, and our daughters.

The most vivid memories of my childhood include my grandmother all dressed up to leave the farm and drive into town, resolute and cheerful in her Relief Society service. My mother’s leadership in the Relief Society of one of the BYU stakes was influential in the lives of hundreds of young women being prepared for a lifetime of service in family, church, and community. I have met these women in many of my visits throughout the Church.

In Chicago, our children and I were schooled in Christian love and service by a mother and wife working in her calling as ward Relief Society president. Later, at BYU, we rejoiced as our daughters were called to leadership and service in the Relief Societies of their BYU branches. The entire family enjoys benefits and blessings through Relief Society service.

From its beginning, the Relief Society has led out in charitable work. At the first meeting, President Emma Smith said, “Each member should be ambitious to do good” (minutes of the Female Relief Society of Nauvoo, 17 Mar. 1842, p. 13; see also “Ambitious to Do Good,” Ensign, Mar. 1992, p. 4; quotations are taken from original documents, hereafter referred to as minutes). The minutes of those initial meetings are filled with accounts of how the sisters obtained work opportunities for the needy, took in the homeless, and made donations to help those in need of food, shelter, and schooling.

A decade after the departure from Nauvoo, sisters trained in the principles of the Relief Society were still leading in efforts to provide for those in need. In a session of conference, President Brigham Young announced that the Saints in two handcart companies were stranded by early snows and suffering in the mountains of Wyoming. He called for immediate help to rescue them, and before they left the Tabernacle many sisters had begun to gather clothing to send to the Saints in the mountains (see Kenneth W. Godfrey et al., Women’s Voices: An Untold History of the Latter-day Saints, 1830–1900 [Salt Lake City: Deseret Book Co., 1982], p. 269).

In the initial meetings of Relief Society, the Prophet Joseph Smith taught that the society “is not only to relieve the poor, but to save souls” (minutes, 9 June 1842, p. 63; History of the Church, 5:25). A later First Presidency explained: “One of the purposes of the organization of the Relief Society was that a system might be inaugurated by which study of religious subjects, or Church doctrine and government, might be pursued by women. The administration of charity under the direction of the Bishopric … was to be part of their active work. But this was not intended to absorb their activities to the exclusion of the development of faith, and the advancement of women in literary, social and domestic activities of life” (in James R. Clark, comp., Messages of the First Presidency of The Church of Jesus Christ of Latter-day Saints, 6 vols. [Salt Lake City: Bookcraft, 1965–75], 5:217).

“To save souls opens the whole field of human activity and development,” Elder John A. Widtsoe later declared. “Relief of poverty, relief of illness; relief of doubt, relief of ignorance—relief of all that hinders the joy and progress of woman. What a magnificent commission!” (John A. Widtsoe, Evidences and Reconciliations [Salt Lake City: Bookcraft, 1987], p. 308).

That commission included teaching. In a revelation given in 1830, the Lord told Emma Smith that the Prophet would authorize her “to expound scriptures, and to exhort the church, according as it shall be given thee by my Spirit” (D\&C 25:7). When she was later selected to lead the Relief Society, her prophet husband referred to this revelation that she would “expound the scriptures to all” and “teach the female part of the community.” He declared “that not she alone, but others, may attain to the same blessings” (minutes, 17 Mar. 1842, p. 8).

Succeeding Presidents of the Church have reemphasized this important duty to teach, and the leaders and teachers of the Relief Society have fulfilled this responsibility with great distinction.

The Relief Society was organized upon the initiative of the women of Nauvoo. Desiring to organize a society to promote sisterhood and to accomplish benevolent works, a group of women asked Eliza R. Snow to draft a constitution and bylaws. When Joseph Smith learned of this, he asked that the sisters be called together so that he could provide “something better for them than a written Constitution.” One sister recalled his saying, “I will organize the women under the priesthood after the pattern of the priesthood” (Sarah M. Kimball, “Auto-Biography,” Woman’s Exponent, 1 Sept. 1883, p. 51).

We are fortunate to have careful minutes of the first two years’ meetings of the Relief Society. From these minutes we know the substance of the Prophet Joseph Smith’s instructions to the new organization and its members. This anniversary is an appropriate time to recall and reemphasize these prophetic directions.

In his first formal instruction to the newly founded organization, the Prophet said he was “deeply interested that [the Relief Society] might be built up to the Most High in an acceptable manner.” He taught that “when instructed, we must obey that voice … that the blessings of heaven may rest down upon us. All must act in concert, or nothing can be done [that the Society] should move according to the ancient Priesthood” (minutes, 30 Mar. 1842, p. 22; History of the Church, 4:570).

The Prophet’s counsel apparently sought to give this new organization the benefit of an early revelation in which the Lord instructed the newly organized First Presidency “how you may act before me, that it may turn to you for your salvation.

“I, the Lord, am bound when ye do what I say; but when ye do not what I say, ye have no promise” (D\&C 82:9–10). The Relief Society’s promised blessings were dependent upon its leaders and members functioning within the limits the Lord had set.

The next time he met with the Relief Society, Joseph Smith “exhorted the sisters always to concentrate their faith and prayers for, and place confidence in those whom God has appointed to honor, whom God has placed at the head to lead” (minutes, 28 Apr. 1842, p. 37; History of the Church, 4:604–5). This counsel, of course, furthered the direction in the earlier revelation on priesthood, which declared that all “authorities or offices in the church are appendages” to the Melchizedek Priesthood and that this priesthood “holds the right of presidency, and has power and authority over all the offices in the church in all ages of the world” (D\&C 107:5, 8). Consequently, the Relief Society and the auxiliaries organized later have always functioned and have thrived under the direction of the presiding authorities of the priesthood.

At this same meeting, the Prophet spoke the words that President Gordon B. Hinckley recently characterized as “a charter … of the Relief Society of The Church of Jesus Christ of Latter-day Saints” (“Ambitious to Do Good,” Ensign, Mar. 1992, p. 4):

“This Society is to get instruction thro’ the order which God has established—thro’ the medium of those appointed to lead” (minutes, 28 Apr. 1842).

Here the Prophet declared that the Relief Society was to receive instruction and direction from the priesthood leaders who presided over their activities. Like the quorums of priesthood holders in the Church, the Relief Society was to be self-governing, but it was not to be an independent organization. It was an integral part of the Church, not a separate church for women.

The Prophet continued, “I now turn the key to you in the name of God and this Society shall rejoice and knowledge and intelligence shall flow down from this time—this is the beginning of better days to this Society” (minutes, 28 Apr. 1842, p. 40).

When he “turn[ed] the key,” the Prophet Joseph Smith made the Relief Society an official part of the Church and kingdom of God. This opened to women new opportunities for receiving knowledge and intelligence from on high, such as through the temple ordinances that were soon to be instituted. Similarly, as the Prophet promised them in connection with their charitable service, “If you live up to your privileges, the angels cannot be restrained from being your associates” (minutes, 28 Apr. 1842, p. 38; History of the Church, 4:605).

President Joseph Fielding Smith taught that the Prophet’s action opened to women the possibility of exercising “some measure of divine authority, particularly in the direction of government and instruction in behalf of the women of the Church” (“Relief Society—An Aid to the Priesthood,” Relief Society Magazine, Jan. 1965, p. 5). President Smith explained: “While the sisters have not been given the Priesthood, … that does not mean that the Lord has not given unto them authority. Authority and Priesthood are two different things. A person may have authority given to him, or a sister to her, to do certain things in the Church that are binding and absolutely necessary for our salvation, such as the work that our sisters do in the House of the Lord” (“The Relief Society Organized by Revelation,” Relief Society Magazine, Jan. 1959, p. 4).

President Smith’s teaching on authority explains what the Prophet Joseph Smith meant when he said that he organized the Relief Society “under the priesthood after the pattern of the priesthood.” The authority to be exercised by the officers and teachers of the Relief Society, as with the other auxiliary organizations, was the authority that would flow to them through their organizational connection with The Church of Jesus Christ of Latter-day Saints and through their individual setting apart under the hands of the priesthood leaders by whom they were called.

No priesthood keys were delivered to the Relief Society. Keys are conferred on individuals, not organizations. The same is true of priesthood authority and of the related authority exercised under priesthood direction. Organizations may channel the exercise of such authority, but they do not embody it. Thus, the priesthood keys were delivered to the members of the First Presidency and the Quorum of Twelve Apostles, not to any organizations (see Topical Guide, “Priesthood, Keys of”).

Under the priesthood authority of the bishop, the president of a ward Relief Society presides over and directs the activities of the Relief Society in the ward. A stake Relief Society president presides and exercises authority over the function to which she has been called. The same is true for the other auxiliaries. Similarly, women called as missionaries are set apart to go forth with authority to teach the everlasting gospel, and women called to work in a temple are given authority for the sacred functions to which they have been called. All function under the direction of the priesthood leader who has been given the priesthood keys to direct those who labor in his area of responsibility.

The Prophet Joseph Smith told the early sisters that he had something better for them than a written constitution. Being organized under priesthood authority, they were to reject worldly concepts of power and seek the power that flows down from heaven for those functions and to those individuals who are using their time and talents in the Lord’s way.

In considering the Prophet’s instructions to the first Relief Society, we should remember that in those earliest days in Church history more revelation was to come. Thus, when he spoke to the sisters about the appropriateness of their laying on hands to bless one another, the Prophet cautioned “that the time had not been before that these things could be in their proper order—that the Church is not now organized in its proper order, and cannot be until the Temple is completed” (minutes, 28 Apr. 1842, p. 36). During the century that followed, as temples became accessible to most members, “proper order” required that these and other sacred practices be confined within those temples.

I will conclude by offering some counsel on the responsibilities of fathers and mothers and priesthood leaders, with special emphasis on matters of interest to the Relief Society.

President Harold B. Lee repeatedly told men that “the greatest of the Lord’s work you … will ever do … will be within the walls of your own home” (in Conference Report, Apr. 1973, p. 130; or Ensign, July 1973, p. 98). That direction also applies to women, and it should engage the best teaching efforts of the Relief Society. We cannot overstate the supreme importance of the task our Father in Heaven has assigned to the mothers, who are the teachers and workers and standard-setters in the homes of the Latter-day Saints. The mothers in those homes give the impressionable sons and daughters of God their earliest and most formative orientation for their mortal journey toward eternal life.

Brethren, we know that the priesthood is the power of God delegated to men to act for the blessing and salvation of all mankind. While we sometimes refer to priesthood holders as “the priesthood,” we must never forget that the priesthood is not owned by or embodied in those who hold it. It is held in a sacred trust to be used for the benefit of men, women, and children alike. Elder John A. Widtsoe said, “Men have no greater claim than women upon the blessings that issue from the Priesthood and accompany its possession” (Priesthood and Church Government [Salt Lake City: Deseret Book Co., 1938], p. 83). For example, our young women should have just as many opportunities for blessings from priesthood leaders as our young men.

Some leaders at various levels of the Church have neglected to apply these basic principles. Some have failed to have the regular consultation with auxiliary leaders that is specified in our Church handbooks of instruction. President Spencer W. Kimball taught the governing principle to the priesthood leaders of the Church when he said: “Our sisters do not wish to be indulged or to be treated condescendingly; they desire to be respected and revered as our sisters and our equals. I mention all these things, my brethren, not because the doctrines or the teachings of the Church regarding women are in any doubt, but because in some situations our behavior is of doubtful quality” (in Conference Report, Oct. 1979, p. 72; or Ensign, Nov. 1979, p. 49).

Priesthood leaders are directed to work in close harmony and partnership with the leaders of our auxiliaries: “As auxiliary leaders work with priesthood leaders to accomplish the mission of the Church, the Lord’s earthly kingdom will prosper and individual lives will be blessed” (Melchizedek Priesthood Leadership Handbook [1990], p. 2).

Only by unity can we follow the way of the Lord, who said, “Be one; and if ye are not one ye are not mine” (D\&C 38:27).

One of the great functions of Relief Society is to provide sisterhood for women, just as priesthood quorums provide brotherhood for men. But all should remember that neither sisterhood nor brotherhood is an end in itself. Each is a means of individual spiritual growth and cooperative service. The ultimate and highest expression of womanhood and manhood is in the new and everlasting covenant of marriage between a man and a woman. Only this relationship culminates in exaltation. As the Apostle Paul taught, “Neither is the man without the woman, neither the woman without the man, in the Lord” (1 Cor. 11:11). Thus, the common objective of brotherhood in our priesthood quorums and sisterhood in our Relief Societies is to bring men and women together in the sacred marriage and family relationships that lead toward eternal life, “the greatest of all the gifts of God” (D\&C 14:7).

We give thanks for the Savior, who made this great goal attainable, for His priesthood authority that administers the essential ordinances, and for the great men and women whose lives are an inspiring legacy of godly service. In the name of Jesus Christ, amen.

\newpage
\settalk{Bible Stories and Personal Protection}
\setspeaker{Dallin H. Oaks}
\settalkdate{October 1992}
\talkheading{Bible Stories and Personal Protection}{Dallin H. Oaks}

My dear brethren, this is an important occasion, when holders of the holy priesthood all over the world gather for instruction and inspiration.

Like many of the older men in this gathering, I have sons and grandsons listening in various locations. We want this meeting to be valuable and interesting to the young men of the priesthood. I am directing my talk primarily to them.

When I was a boy, I spent most of my evenings reading books. One of my favorites was Hurlbut’s Story of the Bible. Published by a Protestant minister to help teach Bible truths to young people, this book tells 168 stories from the Bible.

I loved these stories and read them many times. I will share some of my favorites and comment on their teachings and their impact on my life.

I begin with a story I thought I understood as a boy but did not begin to understand until later.

The Lord spoke to Abraham and told him to take his son Isaac and go to the top of a mountain in the land of Moriah “and offer him there for a burnt offering” (Gen. 22:2).

The first time I read this story I didn’t know what a burnt offering was. But I lived on a farm with animals and mountains nearby, so I could easily understand the rest of the story.

Abraham got up early in the morning and saddled one of his animals, and they started out. I thought that Isaac must have felt privileged to be with his father on such a trip.

On the third day, Abraham and Isaac climbed the mountain to worship. Like most young men, Isaac was curious. He saw the fire and the wood and the knife they carried, “but,” he asked his father, “where is the lamb for a burnt offering?” (Gen. 22:7). I did not realize until I had sons of my own how much pain Abraham must have felt when he answered simply, “My son, God will provide” (Gen. 22:8).

When they came to the prescribed place, Abraham built an altar and laid wood upon it. Then, the Bible says, “Abraham … bound Isaac his son, and laid him on the altar upon the wood” (Gen. 22:9). What did Isaac think when Abraham did such a strange thing? The Bible mentions no struggle or objection. Isaac’s silence can only be explained in terms of his trust in and obedience to his father.

And then the Bible says, “Abraham stretched forth his hand, and took the knife to slay his son” (Gen. 22:10).

As you know, Abraham had passed his test, and the Lord saved young Isaac. “Lay not thine hand upon the lad,” an angel commanded Abraham (Gen. 22:12). A ram whose horns were caught in a thicket became the offering, instead of Isaac.

As a young man, I saw mostly the adventure in that story, though I was surely impressed with Isaac’s obedience. When I was older, I learned that the experience of Abraham and Isaac was what the scriptures call a type, which is a likeness or reminder of something else. The Book of Mormon prophet Jacob said that the command for Abraham to sacrifice his son Isaac was “a similitude of God and his Only Begotten Son” (Jacob 4:5).

This story also shows the goodness of God in protecting Isaac and in providing a substitute so he would not have to die. Because of our sins and our mortality, we, like Isaac, are condemned to death. When all other hope is gone, our Father in Heaven provides the Lamb of God, and we are saved by his sacrifice.

The Apostle Paul taught that the scriptures are “given by inspiration of God” and are “profitable … for instruction in righteousness” (2 Tim. 3:16). We obtain instruction in righteousness from the experiences recorded in the scriptures. They provide what we might call case studies of the results of keeping or breaking the commandments of God.

One example of special importance to young people involved young Joseph, who was sold into Egypt. Though only a slave, Joseph’s abilities were so impressive to his master that he was put in charge of all that his master had in the house and in the field (see Gen. 39:4–6). Then, in that position of prominence and power, Joseph met a test.

His master’s wife tempted him to commit adultery with her. Joseph rejected her advances, telling her he would not betray the trust of her husband or the even greater trust he would violate by sinning against God in doing what Joseph called “this great wickedness” (Gen. 39:9). He rejected her again and again. Then one day, when no one else was in the house, she seized hold of his clothing. In a marvelously vivid description, the scripture says, “And he left his garment in her hand, and fled, and got him out” (Gen. 39:12).

What a persuasive instruction in righteousness! The same teaching was given in this modern revelation: “And go ye out from among the wicked. Save yourselves. Be ye clean that bear the vessels of the Lord” (D\&C 38:42). Those words command all of us to follow the example of Joseph.

A common element in many of my best-loved Bible stories was the way the Lord protected his righteous and faithful sons. When I was young, that was my favorite part of the story of Joseph.

We all remember how the jealous older sons of Jacob plotted to kill their favored younger brother. After seizing Joseph and throwing him into a pit, they decided to sell him into slavery instead. Even as they told their father that Joseph had been killed by wild beasts, the traders who had purchased him on the plains of Canaan were leading him down into Egypt and slavery. (See Gen. 37.)

In Egypt, Joseph was unjustly sent to prison. But even there he excelled, and the Lord blessed him. In time he came forth to interpret Pharaoh’s dream, and he was made ruler of all Egypt. In that powerful position, he became the instrument to save his people from famine and to love and forgive the brothers who had wronged him. (See Gen. 40–45).

As a young boy, I was thrilled with Joseph’s adventures and impressed with how the Lord had rescued him from the perils of murder, slavery, and prison. The first time I read the Book of Mormon, I found the statement that “Joseph … who was sold into Egypt … was preserved by the hand of the Lord” (1 Ne. 5:14). In later readings in the scriptures, I found that this kind of protection is available to all. For example, the Bible states that “the Lord preserveth the faithful” (Ps. 31:23) and that God “is a shield unto them that put their trust in him” (Prov. 30:5).

Another favorite example of God’s protecting care is the shepherd boy David. David had a firm faith in the God of Israel, and that faith gave him great courage.

When the armies of the Philistines were gathered to battle against the Israelites, the mighty Goliath came forward and hurled his challenge to individual combat. King Saul and all Israel “were dismayed, and greatly afraid” (1 Sam. 17:11). Day after day he renewed his challenge, but no one would face him.

When young David came to the camp of Israel to deliver provisions, he heard Goliath’s roar. In surprise David asked, “Who is this uncircumcised Philistine, that he should defy the armies of the living God?” (1 Sam. 17:26). David asked if he could fight the man. The king refused, saying, “Thou art but a youth” (1 Sam. 17:33). David replied with courage and faith: “The Lord that delivered me out of the paw of the lion, … he will deliver me out of the hand of this Philistine” (1 Sam. 17:37).

As David went onto the field of battle, Goliath mocked him for his youth, cursed him by his gods, and shouted that he would feed his flesh to the birds and beasts of the field (see 1 Sam. 17:42–44).

David’s reply is one of the great expressions of faith and courage in all our literature. It thrilled me as a boy, and it still thrills me.

“Thou comest to me with a sword, and with a spear, and with a shield: but I come to thee in the name of the Lord of hosts, the God of the armies of Israel, whom thou hast defied.

“This day will the Lord deliver thee into mine hand; and I will smite thee, and take thine head from thee; and I will give the carcases of the host of the Philistines this day unto the fowls of the air, and to the wild beasts of the earth; that all the earth may know that there is a God in Israel.

“And all this assembly shall know that the Lord saveth not with sword and spear: for the battle is the Lord’s, and he will give you into our hands” (1 Sam. 17:45–47).

You all know what happened next. David stunned the Philistine with a sling-stone and cut off his head with his own sword. Frightened by the fall of their champion, the Philistines fled. Shouting in triumph, the armies of Israel pursued them and won a great victory.

Countless young people have been inspired by this marvelous instruction in righteousness. At times all of us must stand against those who mock and revile. Some of us, sometime, will face some earthly power as mighty as Goliath. When that happens, we should emulate the courage of David, who was mighty because he had faith and he went forth in a righteous cause in the name of the Lord of Hosts.

Our missionaries also seem weak and defenseless, powerless against the armaments of the adversary and those who serve him. But the Lord has promised them that he “will be their shield” (D\&C 35:14), and that promise is fulfilled every day in many places around the world.

The shield the Lord gives to the faithful also protects us against our own harmful impulses. The revelation that commands modern Saints to refrain from alcohol, tobacco, hot drinks, and other harmful things promises the faithful that “the destroying angel shall pass by them, as the children of Israel, and not slay them” (D\&C 89:21).

Another story of protection involved a prophet and his young servant. Because Elisha had helped the kingdom of Israel repel the Syrians, they sent a great army with horses and chariots to capture the prophet. When Elisha’s young servant saw the armies surrounding their city, he cried out in fear, but Elisha reassured him:

“Fear not: for they that be with us are more than they that be with them.

“And Elisha prayed, and said, Lord, I pray thee, open his eyes, that he may see. And the Lord opened the eyes of the young man; and he saw: and, behold, the mountain was full of horses and chariots of fire round about Elisha” (2 Kgs. 6:16–17).

The Lord intervened to confuse and blind the Syrians, and they were taken prisoners by the armies of Israel.

When I read this wonderful story as a boy, I always identified with the young servant of Elisha. I thought, If I am ever surrounded by the forces of evil while I am in the Lord’s service, I hope the Lord will open my eyes and give me faith to understand that when we are in the work of the Lord, those who are with us are always more powerful than those who oppose us.

Bible stories such as these do not mean that the servants of God are delivered from all hardship or that they are always saved from death. Some believers lose their lives in persecutions, and some suffer great hardships as a result of their faith. But the protection promised to the faithful servants of God is a reality today as it was in Bible times.

All over the world, faithful Latter-day Saints are protected from the powers of the evil one and his servants until they have finished their missions in mortality. For some the mortal mission is brief, as with some valiant young men who have lost their lives in missionary service. But for most of us the mortal journey is long, and we continue our course with the protection of guardian angels.

During my life I have had many experiences of being guided in what I should do and in being protected from injury and also from evil. The Lord’s protecting care has shielded me from the evil acts of others and has also protected me from surrendering to my own worst impulses. I enjoyed that protection one warm summer night on the streets of Chicago. I have never shared this experience in public. I do so now because it is a persuasive illustration of my subject.

My wife, June, had attended a ward officers’ meeting. When I came to drive her home, she was accompanied by a sister we would take home on our way. She lived in the nearby Woodlawn area, which was the territory of a gang called the Blackstone Rangers.

I parked at the curb outside this sister’s apartment house and accompanied her into the lobby and up the stairs to her door. June remained in the car on 61st Street. She locked all of the doors, and I left the keys in the ignition in case she needed to drive away. We had lived on the south side of Chicago for quite a few years and were accustomed to such precautions.

Back in the lobby, and before stepping out into the street, I looked carefully in each direction. By the light of a nearby streetlight, I could see that the street was deserted except for three young men walking by. I waited until they were out of sight and then walked quickly toward our car.

As I came to the driver’s side and paused for June to unlock the door, I saw one of these young men running back toward me. He had something in his right hand, and I knew what it would be. There was no time to get into the car and drive away before he came within range.

Fortunately, as June leaned across to open the door, she glanced through the back window and saw this fellow coming around the end of the car with a gun in his hand. Wisely, she did not unlock the door. For the next two or three minutes, which seemed like an eternity, she was a horrified spectator to an event happening at her eye level, just outside the driver’s window.

The young man pushed the gun against my stomach and said, “Give me your money.” I took the wallet out of my pocket and showed him it was empty. I wasn’t even wearing a watch I could offer him because my watchband had broken earlier that day. I offered him some coins I had in my pocket, but he growled a rejection.

“Give me your car keys,” he demanded. “They are in the car,” I told him. “Tell her to open the car,” he replied. For a moment I considered the new possibilities that would present, and then refused. He was furious. He jabbed me in the stomach with his gun and said, “Do it, or I’ll kill you.”

Although this event happened twenty-two years ago, I remember it as clearly as if it were yesterday. I read somewhere that nothing concentrates the mind as wonderfully as having someone stand in front of you with a deadly weapon and tell you he intends to kill you.

When I refused, the young robber repeated his demands, this time emphasizing them with an angrier tone and more motion with his gun. I remember thinking that he probably wouldn’t shoot me on purpose, but if he wasn’t careful in the way he kept jabbing that gun into my stomach, he might shoot me by mistake. His gun looked like a cheap one, and I was nervous about its firing mechanism.

“Give me your money.” “I don’t have any.” “Give me your car keys.” “They’re in the car.” “Tell her to open the car.” “I won’t do it.” “I’ll kill you if you don’t.” “I won’t do it.”

Inside the car June couldn’t hear the conversation, but she could see the action with the gun. She agonized over what she should do. Should she unlock the door? Should she honk the horn? Should she drive away? Everything she considered seemed to have the possibility of making matters worse, so she just waited and prayed. Then a peaceful feeling came over her. She felt it would be all right.

Then, for the first time, I saw the possibility of help. From behind the robber, a city bus approached. It stopped about twenty feet away. A passenger stepped off and scurried away. The driver looked directly at me, but I could see that he was not going to offer any assistance.

While this was happening behind the young robber, out of his view, he became nervous and distracted. His gun wavered from my stomach until its barrel pointed slightly to my left. My arm was already partly raised, and with a quick motion I could seize the gun and struggle with him without the likelihood of being shot. I was taller and heavier than this young man and at that time of my life was somewhat athletic. I had no doubt that I could prevail in a quick wrestling match if I could get his gun out of the contest.

Just as I was about to make my move, I had a unique experience. I did not see anything or hear anything, but I knew something. I knew what would happen if I grabbed that gun. We would struggle, and I would turn the gun into that young man’s chest. It would fire, and he would die. I also understood that I must not have the blood of that young man on my conscience for the rest of my life.

I relaxed, and as the bus pulled away I followed an impulse to put my right hand on his shoulder and give him a lecture. June and I had some teenage children at that time, and giving lectures came naturally.

“Look here,” I said. “This isn’t right. What you’re doing just isn’t right. The next car might be a policeman, and you could get killed or sent to jail for this.”

With the gun back in my stomach, the young robber replied to my lecture by going through his demands for the third time. But this time his voice was subdued. When he offered the final threat to kill me, he didn’t sound persuasive. When I refused again, he hesitated for a moment and then stuck the gun in his pocket and ran away. June unlocked the door, and we drove off, uttering a prayer of thanks. We had experienced the kind of miraculous protection illustrated in the Bible stories I had read as a boy.

I have often pondered the significance of that event in relation to the responsibilities that came later in my life. Less than a year after that August night, I was chosen as president of Brigham Young University. Almost fourteen years after that experience, I received my present calling.

I am grateful that the Lord gave me the vision and strength to refrain from trusting in the arm of flesh and to put my trust in the protecting care of our Heavenly Father. I am grateful for the Book of Mormon promise to us of the last days that “the righteous need not fear,” for the Lord “will preserve the righteous by his power” (1 Ne. 22:17). I am grateful for the protection promised to those who have kept their covenants and qualified for the blessings promised in sacred places.

These and all promises to the faithful children of God are made by the voice and power of the Lord God of Israel. I testify of that God, our Savior Jesus Christ, whose resurrection and atonement have assured immortality and provided the opportunity and direction toward eternal life. In the name of Jesus Christ, amen.

\newpage
\settalk{The Language of Prayer}
\setspeaker{Dallin H. Oaks}
\settalkdate{April 1993}
\talkheading{The Language of Prayer}{Dallin H. Oaks}

When I was young, I learned that great respect was owed to those who held the office of bishop. As a sign of that respect, we always addressed our bishop as Bishop Christensen or Bishop Calder or Brother Jones. We never called our bishop Mr. or by his first name, as we did in speaking to others. With the bishop, we always used an honored title.

When I was seventeen, I joined the Utah National Guard. There I learned that a soldier must use certain words in speaking to an officer. I saw this as another mark of respect for authority. I also observed that this special language served as a way of reminding both the soldier and the officer of the responsibilities of their positions. I later understood that same reasoning as explaining why full-time missionaries should always be called by the dignified titles of elder or sister, or the equivalent in other languages.

In my legal training I became familiar with the formal language lawyers use to address judges during court proceedings. After graduation I worked for a year as a law clerk to the chief justice of the United States. We always used the formal title of his office, Chief Justice. Similarly, communications to our most senior government leaders should be addressed in a particular way, such as Mr. President, Your Excellency, or Your Majesty. The use of titles signifies respect for office and authority.

The words we use in speaking to someone can identify the nature of our relationship to that person. They can also remind speaker and listener of the responsibilities they owe one another in that relationship. The form of address can also serve as a mark of respect or affection.

So it is with the language of prayer. The Church of Jesus Christ of Latter-day Saints teaches its members to use special language in addressing prayers to our Father in Heaven.

When we go to worship in a temple or a church, we put aside our working clothes and dress ourselves in something better. This change of clothing is a mark of respect. Similarly, when we address our Heavenly Father, we should put aside our working words and clothe our prayers in special language of reverence and respect. In offering prayers in the English language, members of our Church do not address our Heavenly Father with the same words we use in speaking to a fellow worker, to an employee or employer, or to a merchant in the marketplace. We use special words that have been sanctified by use in inspired communications, words that have been recommended to us and modeled for us by those we sustain as prophets and inspired teachers.

The special language of prayer follows different forms in different languages, but the principle is always the same. We should address prayers to our Heavenly Father in words which speakers of that language associate with love and respect and reverence and closeness. The application of this principle will, of course, vary according to the nature of a particular language, including the forms that were used when the scriptures were translated into that language. Some languages have intimate or familiar pronouns and verbs used only in addressing family and very close friends. Other languages have honorific forms of address that signify great respect, such as words used only when speaking to a king or other person of high rank. Both of these kinds of special words are appropriately used in offering prayers in other languages because they communicate the desired feelings of love, respect, reverence, or closeness.

Modern English has no special verbs or pronouns that are intimate, familiar, or honorific. When we address prayers to our Heavenly Father in English, our only available alternatives are the common words of speech like you and your or the dignified but uncommon words like thee, thou, and thy, which were used in the King James Version of the Bible almost five hundred years ago. Latter-day Saints, of course, prefer the latter. In our prayers we use language that is dignified and different, even archaic.

The men whom we sustain as prophets, seers, and revelators have consistently taught and urged English-speaking members of our Church to phrase their petitions to the Almighty in the special language of prayer. President Spencer W. Kimball said, “In all our prayers, it is well to use the pronouns thee, thou, thy, and thine instead of you, your, and yours inasmuch as they have come to indicate respect” (Faith Precedes the Miracle [Salt Lake City: Deseret Book Co., 1972], p. 201). Numerous other Church leaders have given the same counsel (see Stephen L Richards, in Conference Report, Oct. 1951, p. 175; Bruce R. McConkie, “Why the Lord Ordained Prayer,” Ensign, Jan. 1976, p. 12; and L. Tom Perry, in Conference Report, Oct. 1983, pp. 14–15; or Ensign, Nov. 1983, p. 13).

Perhaps some who are listening to this sermon in English are already saying, “But this is unfamiliar and difficult. Why should we have to use words that have not been in common use in the English language for hundreds of years? If we require a special language of prayer in English, we will discourage the saying of prayers by little children, by new members, and by others who are just learning to pray.”

Brothers and sisters, the special language of prayer is much more than an artifact of the translation of the scriptures into English. Its use serves an important, current purpose. We know this because of modern revelations and because of the teachings and examples of modern prophets. The way we pray is important.

The English words thee, thou, thy, and thine occur throughout the prayers the prophets of the Lord have revealed for use in our day.

A revelation given in 1830, the year the Church was organized, directs that the elder or priest who administers the sacrament “shall kneel … and call upon the Father in solemn prayer, saying:

“O God, the Eternal Father, we ask thee in the name of thy Son, Jesus Christ” (D\&C 20:76–77, 79).

The prayer offered at the dedication of the Kirtland Temple in 1836 is another model that illustrates the language of prayer used by the Prophet Joseph Smith:

“And now, Holy Father, we ask thee to assist us, thy people, with thy grace, in calling our solemn assembly, …

“That thy glory may rest down upon thy people, and upon this thy house, which we now dedicate to thee, that it may be sanctified and consecrated to be holy, and that thy holy presence may be continually in this house” (D\&C 109:10, 12).

This prophetic model of the language of prayer has been faithfully followed in all of the sacred petitions by which the prophets have dedicated temples to the Lord. Exactly one hundred years ago this week, at a spot not far from where I stand, President Wilford Woodruff began the dedicatory prayer of the Salt Lake Temple with these words:

“Our Father in heaven, thou who hast created the heavens and the earth, and all things that are therein; thou most glorious One, … we, thy children, come this day before thee, and in this house which we have built to thy most holy name, humbly plead the atoning blood of thine Only Begotten Son, that our sins may be remembered no more against us forever, but that our prayers may ascend unto thee and have free access to thy throne, that we may be heard in thy holy habitation” (Deseret Semi-Weekly News, 7 Apr. 1893, p. 2; see also Gordon B. Hinckley, “The Salt Lake Temple,” Ensign, Mar. 1993, p. 2).

When the Prophet Joseph Smith was imprisoned in the jail at Liberty, Missouri, he wrote an inspired prayer, which we now read in the 121st section of the Doctrine and Covenants. Note the special language the Prophet used in addressing our Father in Heaven:

“O God, where art thou? And where is the pavilion that covereth thy hiding place? …

“Remember thy suffering saints, O our God; and thy servants will rejoice in thy name forever” (D\&C 121:1, 6).

Other prayers offered by the Prophet Joseph Smith also use the special, formal language of prayer (see The Personal Writings of Joseph Smith, ed. Dean C. Jessee [Salt Lake City: Deseret Book Co., 1984], pp. 283–84, 536–37).

To cite more recent examples, we are all aware that the prayers offered at these general conferences of the Church always use the special language of prayer we have learned from the examples of modern prophets and teachers.

We are also guided by the special language we read in the prayers recorded in the King James Translation of the Bible and in the Book of Mormon.

We have scriptural record of three beautiful translated prayers the Savior offered during his earthly ministry. They are models for all of us. Notable in each of these prayers are the words thee, thou, thy, and thine instead of you, your, and yours.

In teaching his disciples what we call the Lord’s Prayer, the Savior said, “After this manner therefore pray ye: Our Father which art in heaven, Hallowed be thy name” (Matt. 6:9; see also 3 Ne. 13:9).

In his great intercessory prayer, uttered on the night before his crucifixion, the Savior used these words:

“Father, the hour is come; glorify thy Son, that thy Son also may glorify thee. …

“And this is life eternal, that they might know thee the only true God, and Jesus Christ, whom thou hast sent” (John 17:1, 3).

The Book of Mormon records this prayer the Savior offered during his visit to the righteous remnant of Israel on the American continent following his resurrection:

“Father, I thank thee that thou hast given the Holy Ghost unto these whom I have chosen. …

“Father, I pray thee that thou wilt give the Holy Ghost unto all them that shall believe in their words” (3 Ne. 19:20–21).

The special language of prayer that Latter-day Saints use in English has sometimes been explained by reference to the history of the English language. It has been suggested that thee, thou, thy, and thine are simply holdovers from forms of address once used to signify respect for persons of higher rank. But more careful scholarship shows that the words we now use in the language of prayer were once commonly used by persons of rank in addressing persons of inferior position. These same English words were also used in communications between persons in an intimate relationship. There are many instances where usages of English words have changed over the centuries. But the history of English usage is not the point.

Scholarship can contradict mortal explanations, but it cannot rescind divine commands or inspired counsel. In our day the English words thee, thou, thy, and thine are suitable for the language of prayer, not because of how they were used anciently but because they are currently obsolete in common English discourse. Being unused in everyday communications, they are now available as a distinctive form of address in English, appropriate to symbolize respect, closeness, and reverence for the one being addressed.

I hope this renewal of counsel that we use special language in our prayers will not be misunderstood.

Literary excellence is not our desire. We do not advocate flowery and wordy prayers. We do not wish to be among those who “pray to be heard of men, and to be praised for their wisdom” (Alma 38:13). We wish to follow the Savior’s teaching, “When ye pray, use not vain repetitions, as the heathen do: for they think that they shall be heard for their much speaking” (Matt. 6:7; see also 3 Ne. 13:7). Our prayers should be simple, direct, and sincere.

We should also remember that our position on special prayer language in English is based on modern revelations and the teachings and examples of modern prophets. It is not part of the teachings known and accepted by our brothers and sisters of other Christian and Jewish faiths. When leaders or members of other churches or synagogues phrase their prayers in the familiar forms of you or your, this does not signify a lack of reverence or respect in their belief and practice but only a preference for the more modern language. Significantly, this modern language is frequently the language used in the scriptural translations with which they are most familiar.

We are especially anxious that our position on special language in prayers in English not cause some to be reluctant to pray in our Church meetings or in other settings where their prayers are heard. We have particular concern for converts and others who have not yet had experience in using these words.

I am sure that our Heavenly Father, who loves all of his children, hears and answers all prayers, however phrased. If he is offended in connection with prayers, it is likely to be by their absence, not their phraseology.

When one of our daughters was about three years old, she did something that always delighted her parents. When we called her name, she would usually answer by saying, “Here me is.” This childish reply was among the sweetest things her parents heard. But when she was grown, we expected her to use appropriate language when she spoke, and of course she did. As the Apostle Paul said, “When I was a child, I spake as a child, I understood as a child, I thought as a child: but when I became a man, I put away childish things” (1 Cor. 13:11).

The same is true of prayer. Our earliest efforts will be heard with joy by our Heavenly Father, however they are phrased. They will be heard in the same way by loving members of our Church. But as we gain experience as members of The Church of Jesus Christ of Latter-day Saints, we need to become more mature in all of our efforts, including our prayers.

Men and women who wish to show respect will take the time to learn the special language of prayer. Persons spend many hours mastering communication skills in other mediums, such as poetry or prose, vocal or instrumental music, and even the language of access to computers. My brothers and sisters, the manner of addressing our Heavenly Father in prayer is at least as important as these.

It requires a little time for adults to learn how to use the language of prayer. But it is not really very difficult. In fact, we are more than 75 percent of the way in English prayers when we simply delete you and your and substitute thee and thy (see Don E. Norton, Jr., “The Language of Formal Prayer,” Ensign, Jan. 1976, pp. 44–47). The special language of prayer is even easier in most other languages.

Modern revelation commands parents to “teach their children to pray” (D\&C 68:28). This requires parents to learn and pray with the special language of prayer. We learn our native language simply by listening to those who speak it. This is also true of the language with which we address our Heavenly Father. The language of prayer is easier and sweeter to learn than any other tongue. We should give our children the privilege of learning this language by listening to their parents use it in the various prayers offered daily in our homes.

The Prophet Joseph Smith said, “It is a great thing to inquire at the hands of God, or to come into His presence” (Teachings of the Prophet Joseph Smith, sel. Joseph Fielding Smith [Salt Lake City: Deseret Book Co., 1938], p. 22). The special language of prayer reminds us of the greatness of that privilege. I pray that all of us will be more sensitive to the importance of using this reverent and loving language as we offer our public and private prayers.

I testify that this is the Church of Jesus Christ, which our Savior has restored in these latter days with the authority and duty to preach his gospel and his commandments to every nation, kindred, tongue, and people. In the name of Jesus Christ, amen.

\newpage
\settalk{“The Great Plan of Happiness”}
\setspeaker{Dallin H. Oaks}
\settalkdate{October 1993}
\talkheading{“The Great Plan of Happiness”}{Dallin H. Oaks}

Questions like, Where did we come from? Why are we here? and Where are we going? are answered in the gospel of Jesus Christ. Prophets have called it the plan of salvation and “the great plan of happiness” (Alma 42:8). Through inspiration we can understand this road map of eternity and use it to guide our path in mortality.

The gospel teaches us that we are the spirit children of heavenly parents. Before our mortal birth we had “a pre-existent, spiritual personality, as the sons and daughters of the Eternal Father” (statement of the First Presidency, Improvement Era, Mar. 1912, p. 417; see also Jer. 1:5). We were placed here on earth to progress toward our destiny of eternal life. These truths give us a unique perspective and different values to guide our decisions from those who doubt the existence of God and believe that life is the result of random processes.

Our understanding of life begins with a council in heaven. There the spirit children of God were taught his eternal plan for their destiny. We had progressed as far as we could without a physical body and an experience in mortality. To realize a fulness of joy, we had to prove our willingness to keep the commandments of God in a circumstance where we had no memory of what preceded our mortal birth.

In the course of mortality, we would become subject to death, and we would be soiled by sin. To reclaim us from death and sin, our Heavenly Father’s plan provided us a Savior, whose atonement would redeem all from death and pay the price necessary for all to be cleansed from sin on the conditions he prescribed (see 2 Ne. 9:19–24).

Satan had his own plan. He proposed to save all the spirit children of God, assuring that result by removing their power to choose and thus eliminating the possibility of sin. When Satan’s plan was rejected, he and the spirits who followed him opposed the Father’s plan and were cast out.

All of the myriads of mortals who have been born on this earth chose the Father’s plan and fought for it. Many of us also made covenants with the Father concerning what we would do in mortality. In ways that have not been revealed, our actions in the spirit world influence us in mortality.

Although Satan and his followers have lost their opportunity to have a physical body, they are permitted to use their spirit powers to try to frustrate God’s plan. This provides the opposition necessary to test how mortals will use their freedom to choose. Satan’s most strenuous opposition is directed at whatever is most important to the Father’s plan. Satan seeks to discredit the Savior and divine authority, to nullify the effects of the Atonement, to counterfeit revelation, to lead people away from the truth, to contradict individual accountability, to confuse gender, to undermine marriage, and to discourage childbearing (especially by parents who will raise children in righteousness).

Maleness and femaleness, marriage, and the bearing and nurturing of children are all essential to the great plan of happiness. Modern revelation makes clear that what we call gender was part of our existence prior to our birth. God declares that he created “male and female” (D\&C 20:18; Moses 2:27; Gen. 1:27). Elder James E. Talmage explained: “The distinction between male and female is no condition peculiar to the relatively brief period of mortal life; it was an essential characteristic of our pre-existent condition” (Millennial Star, 24 Aug. 1922, p. 539).

To the first man and woman on earth, the Lord said, “Be fruitful, and multiply” (Moses 2:28; Gen. 1:28; see also Abr. 4:28). This commandment was first in sequence and first in importance. It was essential that God’s spirit children have mortal birth and an opportunity to progress toward eternal life. Consequently, all things related to procreation are prime targets for the adversary’s efforts to thwart the plan of God.

When Adam and Eve received the first commandment, they were in a transitional state, no longer in the spirit world but with physical bodies not yet subject to death and not yet capable of procreation. They could not fulfill the Father’s first commandment without transgressing the barrier between the bliss of the Garden of Eden and the terrible trials and wonderful opportunities of mortal life.

For reasons that have not been revealed, this transition, or “fall,” could not happen without a transgression—an exercise of moral agency amounting to a willful breaking of a law (see Moses 6:59). This would be a planned offense, a formality to serve an eternal purpose. The Prophet Lehi explained that “if Adam had not transgressed he would not have fallen” but would have remained in the same state in which he was created (2 Ne. 2:22).

“And they would have had no children; wherefore they would have remained in a state of innocence, having no joy, for they knew no misery; doing no good, for they knew no sin” (2 Ne. 2:23).

But the Fall was planned, Lehi concludes, because “all things have been done in the wisdom of him who knoweth all things” (2 Ne. 2:24).

It was Eve who first transgressed the limits of Eden in order to initiate the conditions of mortality. Her act, whatever its nature, was formally a transgression but eternally a glorious necessity to open the doorway toward eternal life. Adam showed his wisdom by doing the same. And thus Eve and “Adam fell that men might be” (2 Ne. 2:25).

Some Christians condemn Eve for her act, concluding that she and her daughters are somehow flawed by it. Not the Latter-day Saints! Informed by revelation, we celebrate Eve’s act and honor her wisdom and courage in the great episode called the Fall (see Bruce R. McConkie, “Eve and the Fall,” in Woman [Salt Lake City: Deseret Book Co., 1979], pp. 67–68). Joseph Smith taught that it was not a “sin” because God had decreed it (see The Words of Joseph Smith, ed. Andrew F. Ehat and Lyndon W. Cook [Provo: Religious Studies Center, Brigham Young University, 1980], p. 63). Brigham Young declared, “We should never blame Mother Eve, not the least” (in Journal of Discourses, 13:145). Elder Joseph Fielding Smith said: “I never speak of the part Eve took in this fall as a sin, nor do I accuse Adam of a sin. … This was a transgression of the law, but not a sin … for it was something that Adam and Eve had to do!” (Doctrines of Salvation, comp. Bruce R. McConkie, 3 vols. [Salt Lake City: Bookcraft, 1954–56], 1:114–15).

This suggested contrast between a sin and a transgression reminds us of the careful wording in the second article of faith: “We believe that men will be punished for their own sins, and not for Adam’s transgression” (emphasis added). It also echoes a familiar distinction in the law. Some acts, like murder, are crimes because they are inherently wrong. Other acts, like operating without a license, are crimes only because they are legally prohibited. Under these distinctions, the act that produced the Fall was not a sin—inherently wrong—but a transgression—wrong because it was formally prohibited. These words are not always used to denote something different, but this distinction seems meaningful in the circumstances of the Fall.

Modern revelation shows that our first parents understood the necessity of the Fall. Adam declared, “Blessed be the name of God, for because of my transgression my eyes are opened, and in this life I shall have joy, and again in the flesh I shall see God” (Moses 5:10).

Note the different perspective and the special wisdom of Eve, who focused on the purpose and effect of the great plan of happiness: “Were it not for our transgression we never should have had seed, and never should have known good and evil, and the joy of our redemption, and the eternal life which God giveth unto all the obedient” (Moses 5:11). In his vision of the redemption of the dead, President Joseph F. Smith saw “the great and mighty ones” assembled to meet the Son of God, and among them was “our glorious Mother Eve” (D\&C 138:38–39).

When we understand the plan of salvation, we also understand the purpose and effect of the commandments God has given his children. He teaches us correct principles and invites us to govern ourselves. We do this by the choices we make in mortality.

We live in a day when there are many political, legal, and social pressures for changes that confuse gender and homogenize the differences between men and women. Our eternal perspective sets us against changes that alter those separate duties and privileges of men and women that are essential to accomplish the great plan of happiness. We do not oppose all changes in the treatment of men and women, since some changes in laws or customs simply correct old wrongs that were never grounded in eternal principles.

The power to create mortal life is the most exalted power God has given his children. Its use was mandated in the first commandment, but another important commandment was given to forbid its misuse. The emphasis we place on the law of chastity is explained by our understanding of the purpose of our procreative powers in the accomplishment of God’s plan.

The expression of our procreative powers is pleasing to God, but he has commanded that this be confined within the relationship of marriage. President Spencer W. Kimball taught that “in the context of lawful marriage, the intimacy of sexual relations is right and divinely approved. There is nothing unholy or degrading about sexuality in itself, for by that means men and women join in a process of creation and in an expression of love” (The Teachings of Spencer W. Kimball, ed. Edward L. Kimball [Salt Lake City: Bookcraft, 1982], p. 311).

Outside the bonds of marriage, all uses of the procreative power are to one degree or another a sinful degrading and perversion of the most divine attribute of men and women. The Book of Mormon teaches that unchastity is “most abominable above all sins save it be the shedding of innocent blood or denying the Holy Ghost” (Alma 39:5). In our own day, the First Presidency of the Church has declared the doctrine of this Church “that sexual sin—the illicit sexual relations of men and women—stands, in its enormity, next to murder” (in James R. Clark, comp., Messages of the First Presidency of The Church of Jesus Christ of Latter-day Saints, 6 vols. [Salt Lake City: Bookcraft, 1965–75], 6:176). Some who do not know the plan of salvation behave like promiscuous animals, but Latter-day Saints—especially those who are under sacred covenants—have no such latitude. We are solemnly responsible to God for the destruction or misuse of the creative powers he has placed within us.

The ultimate act of destruction is to take a life. That is why abortion is such a serious sin. Our attitude toward abortion is not based on revealed knowledge of when mortal life begins for legal purposes. It is fixed by our knowledge that according to an eternal plan, all of the spirit children of God must come to this earth for a glorious purpose, and that individual identity began long before conception and will continue for all the eternities to come. We rely on the prophets of God, who have told us that while there may be “rare” exceptions, “the practice of elective abortion is fundamentally contrary to the Lord’s injunction, ‘Thou shalt not … kill, nor do anything like unto it’ (D\&C 59:6)” (1991 Supplement to the 1989 General Handbook of Instructions, p. 1).

Our knowledge of the great plan of happiness also gives us a unique perspective on the subject of marriage and the bearing of children. In this we also run counter to some strong current forces in custom, law, and economics.

Marriage is disdained by an increasing number of couples, and many who marry choose to forgo children or place severe limits on their number. In recent years strong economic pressures in many nations have altered the traditional assumption of a single breadwinner per family. Increases in the number of working mothers of young children inevitably signal a reduced commitment of parental time to nurturing the young. The effect of these reductions is evident in the rising numbers of abortions, divorces, child neglect, and juvenile crime.

We are taught that marriage is necessary for the accomplishment of God’s plan, to provide the approved setting for mortal birth, and to prepare family members for eternal life. “Marriage is ordained of God unto man,” the Lord said, “that the earth might answer the end of its creation; and that it might be filled with the measure of man, according to his creation before the world was made” (D\&C 49:15–17).

Our concept of marriage is motivated by revealed truth, not by worldly sociology. The Apostle Paul taught, “Neither is the man without the woman, neither the woman without the man, in the Lord” (1 Cor. 11:11). President Spencer W. Kimball explained, “Without proper and successful marriage, one will never be exalted” (Marriage and Divorce [Salt Lake City: Deseret Book Co., 1976], p. 24).

According to custom, men are expected to take the initiative in seeking marriage. That is why President Joseph F. Smith directed his prophetic pressure at men. He said, “No man who is marriageable is fully living his religion who remains unmarried” (Gospel Doctrine [Salt Lake City: Deseret Book Co., 1939], p. 275). We hear of some worthy LDS men in their thirties who are busy accumulating property and enjoying freedom from family responsibilities without any sense of urgency about marriage. Beware, brethren. You are deficient in a sacred duty.

Knowledge of the great plan of happiness also gives Latter-day Saints a distinctive attitude toward the bearing and nurturing of children.

In some times and places, children have been regarded as no more than laborers in a family economic enterprise or as insurers of support for their parents. Though repelled by these repressions, some persons in our day have no compunctions against similar attitudes that subordinate the welfare of a spirit child of God to the comfort or convenience of parents.

The Savior taught that we should not lay up treasures on earth but should lay up treasures in heaven (see Matt. 6:19–21). In light of the ultimate purpose of the great plan of happiness, I believe that the ultimate treasures on earth and in heaven are our children and our posterity.

President Kimball said, “It is an act of extreme selfishness for a married couple to refuse to have children when they are able to do so” (in Conference Report, Apr. 1979, p. 6; or Ensign, May 1979, p. 6). When married couples postpone childbearing until after they have satisfied their material goals, the mere passage of time assures that they seriously reduce their potential to participate in furthering our Heavenly Father’s plan for all of his spirit children. Faithful Latter-day Saints cannot afford to look upon children as an interference with what the world calls “self-fulfillment.” Our covenants with God and the ultimate purpose of life are tied up in those little ones who reach for our time, our love, and our sacrifices.

How many children should a couple have? All they can care for! Of course, to care for children means more than simply giving them life. Children must be loved, nurtured, taught, fed, clothed, housed, and well started in their capacities to be good parents themselves. Exercising faith in God’s promises to bless them when they are keeping his commandments, many LDS parents have large families. Others seek but are not blessed with children or with the number of children they desire. In a matter as intimate as this, we should not judge one another.

President Gordon B. Hinckley gave this inspired counsel to an audience of young Latter-day Saints:

“I like to think of the positive side of the equation, of the meaning and sanctity of life, of the purpose of this estate in our eternal journey, of the need for the experiences of mortal life under the great plan of God our Father, of the joy that is to be found only where there are children in the home, of the blessings that come of good posterity. When I think of these values and see them taught and observed, then I am willing to leave the question of numbers to the man and the woman and the Lord” (“If I Were You, What Would I Do?” Brigham Young University 1983–84 Fireside and Devotional Speeches [Provo: University Publications, 1984], p. 11).

Some who are listening to this message are probably saying, “But what about me?” We know that many worthy and wonderful Latter-day Saints currently lack the ideal opportunities and essential requirements for their progress. Singleness, childlessness, death, and divorce frustrate ideals and postpone the fulfillment of promised blessings. In addition, some women who desire to be full-time mothers and homemakers have been literally compelled to enter the full-time workforce. But these frustrations are only temporary. The Lord has promised that in the eternities no blessing will be denied his sons and daughters who keep the commandments, are true to their covenants, and desire what is right.

Many of the most important deprivations of mortality will be set right in the Millennium, which is the time for fulfilling all that is incomplete in the great plan of happiness for all of our Father’s worthy children. We know that will be true of temple ordinances. I believe it will also be true of family relationships and experiences.

I pray that we will not let the challenges and temporary diversions of mortality cause us to forget our covenants and lose sight of our eternal destiny. We who know God’s plan for his children, we who have covenanted to participate, have a clear responsibility. We must desire to do what is right, and we must do all that we can in our own circumstances in mortality.

In all of this, we should remember King Benjamin’s caution to “see that all these things are done in wisdom and order; for it is not requisite that a man should run faster than he has strength” (Mosiah 4:27). I think of that inspired teaching whenever I feel inadequate, frustrated, or depressed.

When we have done all that we are able, we can rely on God’s promised mercy. We have a Savior, who has taken upon him not just the sins, but also “the pains and the sicknesses of his people … that he may know according to the flesh how to succor his people according to their infirmities” (Alma 7:11–12). He is our Savior, and when we have done all that we can, he will make up the difference, in his own way and in his own time. Of that I testify in the name of Jesus Christ, amen.

\newpage
\settalk{“Faith in the Lord Jesus Christ”}
\setspeaker{Dallin H. Oaks}
\settalkdate{April 1994}
\talkheading{“Faith in the Lord Jesus Christ”}{Dallin H. Oaks}

My beloved young sisters, I have been inspired by the prayers, music, and words in this marvelous meeting. I feel that each young woman who is listening has been strengthened in her resolve to become what President Janette C. Hales has challenged her to become: a righteous, problem-solving woman of faith.

These wonderful women who are the general presidency of the Young Women of the Lord’s church have told us how this can be accomplished—how we can seek and obtain and increase faith in the Lord Jesus Christ. Sister Pearce gave us inspiring examples of men and women who exercised faith and trust in our Heavenly Father and his Son, Jesus Christ, by believing that they are in charge of this world, that they know us and love us, and that they have a plan for us. Sister Pinegar taught us that we can and should seek and choose to believe in our Savior and his love.

These teachings and these teachers are true. I feel challenged by the responsibility of concluding a meeting on this most fundamental subject.

The first principle of the gospel is not “faith.” The first principle of the gospel is “Faith in the Lord Jesus Christ” (A of F 1:4). I wish to speak to you young women about that supremely important truth.

Faith does not exist by itself. Faith requires an object. It must be faith in something or someone.

In that respect, faith is like love. Love cannot exist without an object.

A personal experience illustrates that point. Sister Oaks and I are the parents of six, including four daughters. Our youngest daughter is still in her teens. As parents, we have learned a lot about teenage girls. I remember when one of our teenage daughters announced that she was in love with eight boys. She produced a list of their names. I made silent note of the fact that she had never even dated some of these boys, and one of them she had never even met. Within a few weeks she dropped several names off her list and added others. When I asked her how she could fall in love and out of love with so many boys so quickly, she wisely admitted, “I guess I’m not in love with those boys. I’m just in love with love.” Your parents and grandparents will remember the words of an old song, “Falling in love with love is falling for make believe” (Lorenz Hart, “Falling in Love with Love,” The Boys from Syracuse, n.p.: Chappell \& Co., 1938).

Love is meaningless unless it is directed toward something or someone. We love our parents. We love our brothers and sisters. We love the Lord.

Faith is the same. If we think we have faith, we should ask, faith in whom or faith in what? For some, faith is nothing more than faith in themselves. That is only self-confidence or self-centeredness. Others have faith in faith, which is something like relying on the power of positive thinking or betting on the proposition that we can get what we want by manipulating the powers within us.

The first principle of the gospel is faith in the Lord Jesus Christ. Without this faith, the prophet Mormon said, we “are not fit to be numbered among the people of his church” (Moro. 7:39).

The scriptures teach us that faith comes by hearing the word of God (see Rom. 10:17). That word, which comes to us by scripture, by prophetic teaching, and by personal revelation, teaches us that we are children of God, the Eternal Father. It teaches us about the identity and mission of Jesus Christ, his Only Begotten Son, our Savior and Redeemer. Founded on our knowledge of those things, faith in the Lord Jesus Christ is a conviction and trust that God knows us and loves us and will hear our prayers and answer them with what is best for us.

In fact, God will do more than what is best for us. He will do what is best for us and for all of our Heavenly Father’s children. The conviction that the Lord knows more than we do and that he will answer our prayers in the way that is best for us and for all of his other children is a vital ingredient of faith in the Lord Jesus Christ. This important reality is beautifully described in an experience recorded in Elder John H. Groberg’s recent book, In the Eye of the Storm. He describes a lesson he learned as a young missionary traveling on a sailboat in the Tongan islands.

“We would always pray for protection, success, and good seas and wind to take us to our destination. Once I asked the Lord to bless us with a good tail wind so we could get to Foa quickly. As we got under way, one of the older men said, ‘Elder Groberg, you need to modify your prayers a little.’

“‘How’s that?’ I replied.

“‘You asked the Lord for a tail wind to take us rapidly to Foa. If you pray for a tail wind to Foa, what about the people who are trying to come from Foa to Pangai? They are good people, and you are praying against them. Just pray for a good wind, not a tail wind.’

“That taught me something important. Sometimes we pray for things that will benefit us but may hurt others. We may pray for a particular type of weather, or to preserve someone’s life, when that answer to our prayer may hurt someone else. That’s why we must always pray in faith, because we can’t have true, God-given faith in something that is not according to His will. If it’s according to His will, all parties will benefit. I learned to pray for a good wind and the ability to get there safely, not necessarily a tail wind” (Salt Lake City: Bookcraft, 1993, p. 175).

Faith must include trust. I am glad that each member of the presidency stressed that fact in her talk. When we have faith in the Lord Jesus Christ, we must have trust in him. We must trust him enough that we are content to accept his will, knowing that he knows what is best for us.

The kind of faith that includes trust in the Lord stands in contrast to many imitations. Some people trust no one but themselves. Some put their highest trust in a friend or another family member, perhaps because they feel that person is more righteous or more wise than they. But that is not the Lord’s way. He told us to put our faith and our trust in the Lord Jesus Christ.

The Savior gave us the model for that kind of faith and trust. Remember how he prayed to the Father in the agony of Gethsemane? This was the culminating event of his life, the climactic fulfillment of his mission as our Savior. The gospel of Luke, as corrected in the inspired translation of the Prophet Joseph Smith, describes how he knelt down and prayed: “Father, if thou be willing, remove this cup from me; nevertheless, not my will, but thine be done” (JST, Luke 22:42 [The Holy Scriptures: Inspired Version]).

Here we see the Savior’s absolute faith and trust in the Father. “Nevertheless,” he said, “not my will, but thine be done.” The Father’s answer was to deny the plea of his Only Begotten Son. The Atonement had to be worked out by that lamb without blemish. But though the Son’s request was denied, his prayer was answered. The scripture records: “And there appeared an angel unto him from heaven, strengthening him” (JST, Luke 22:43 [The Holy Scriptures: Inspired Version]).

Strengthened from heaven to do the will of the Father, the Savior fulfilled his mission. “And being in an agony, he prayed more earnestly; and he sweat as it were great drops of blood falling down to the ground” (JST, Luke 22:44).

When we try to develop faith in the Lord Jesus Christ rather than merely cultivating faith as an abstract principle of power, we understand the meaning of the Savior’s words: “If ye will have faith in me ye shall have power to do whatsoever thing is expedient in me” (Moro. 7:33).

Similarly, the Savior taught the Nephites that they must always pray to the Father in his name, adding: “And whatsoever ye shall ask the Father in my name, which is right, believing that ye shall receive, behold it shall be given unto you” (3 Ne. 18:20).

Here the Savior reminds us that faith, no matter how strong it is, cannot produce a result contrary to the will of him whose power it is. The exercise of faith in the Lord Jesus Christ is always subject to the order of heaven, to the goodness and will and wisdom and timing of the Lord. That is why we cannot have true faith in the Lord without also having complete trust in the Lord’s will and in the Lord’s timing. When we have that kind of faith and trust in the Lord, we have true security in our lives. President Spencer W. Kimball said, “Security is not born of inexhaustible wealth but of unquenchable faith” (The Teachings of Spencer W. Kimball, ed. Edward L. Kimball, Salt Lake City: Bookcraft, 1982, pp. 72–73).

I read of a young woman who exercised that kind of faith and trust. For many months her mother had been seriously ill. Finally, the faithful father called the children to her bedside and told them to say good-bye to their mother because she was dying. The twelve-year-old daughter protested:

“Papa, I do not want my mamma to die. I have been with her in the hospital … for six months; time and time again … you have administered to her, and she has been relieved of her pain and quietly gone to sleep. I want you to lay hands upon my mamma and heal her.”

The father, who was Elder Heber J. Grant, told the children that he felt in his heart that their mother’s time had arrived. The children left, and he knelt by his wife’s bedside. Later he recalled his prayer: “I told the Lord I acknowledged his hand in life [and] in death. … But I told the Lord that I lacked the strength to have my wife die and to have it affect the faith of my little children.” He pleaded with the Lord to give his daughter “a knowledge that it was his mind and his will that her mamma should die.”

Within an hour the mother died. When Elder Grant called the children back into her room and told them, his little six-year-old boy began to weep bitterly. The twelve-year-old sister took him in her arms and said: “Do not weep, Heber; since we went out of this room, the voice of the Lord from heaven has said to me, In the death of your mamma the will of the Lord shall be done” (Bryant S. Hinckley, Heber J. Grant: Highlights in the Life of a Great Leader, Salt Lake City: Deseret Book Co., 1951, pp. 243–44).

When we have the kind of faith and trust exhibited by that young woman, we have the strength to sustain us in every important event in our lives. President Spencer W. Kimball said that we need what he called “reservoirs of faith” to stand firm and strong against all the temptations and adversities of life (Spencer W. Kimball, Faith Precedes the Miracle, Salt Lake City: Deseret Book Co., 1972, pp. 110–11).

My beloved young sisters, each of you needs to build a reservoir of faith so you can draw upon it when someone you love or respect betrays you, when some scientific discovery seems to cast doubt on a gospel principle, or when someone makes light of sacred things, such as the name of God or the sacred ceremonies of the temple. You need to draw on your reservoir of faith when you are weak or when someone else calls on you to strengthen them. You also need to draw on your reservoir of faith when some requirement of Church membership or service interferes with your personal preferences.

You need the strength that comes from faith and trust in the Lord Jesus Christ if you are to fulfill your duty “to stand as witnesses of God at all times and in all things, and in all places” (Mosiah 18:9). In times of trial you need the comfort offered in the holy scriptures, which assure you that when you have the shield of faith you will “be able to quench all the fiery darts of the wicked” (D\&C 27:17).

Faith in the Lord Jesus Christ prepares you for whatever life brings. This kind of faith prepares you to deal with life’s opportunities—to take advantage of those that are received and to persist through the disappointments of those that are lost.

Most importantly, faith in the Lord Jesus Christ opens the door of salvation and exaltation: “For no [one] can be saved, according to the words of Christ, save they shall have faith in his name” (Moro. 7:38).

I testify that these things are true. I invoke the blessings of Almighty God upon you, my faithful young sisters, as you seek to develop and exercise your faith and trust in the Lord Jesus Christ and as you seek to serve him and keep his commandments, in the name of Jesus Christ, amen.

\newpage
\settalk{Tithing}
\setspeaker{Dallin H. Oaks}
\settalkdate{April 1994}
\talkheading{Tithing}{Dallin H. Oaks}

When the risen Lord appeared to the faithful on this continent, he taught them the commandments the prophet Malachi had already given to other children of Israel. The Lord commanded that they should record these words (see 3 Ne. 24:1).

“Will a man rob God? Yet ye have robbed me. But ye say: Wherein have we robbed thee? In tithes and offerings.

“Ye are cursed with a curse, for ye have robbed me, even this whole nation.

“Bring ye all the tithes into the storehouse, that there may be meat in my house; and prove me now herewith, saith the Lord of Hosts, if I will not open you the windows of heaven, and pour you out a blessing that there shall not be room enough to receive it” (3 Ne. 24:8–10; see also Mal. 3:8–10).

After the Savior quoted these words, “he expounded them unto the multitude” and said, “These scriptures, which ye had not with you, the Father commanded that I should give unto you; for it was wisdom in him that they should be given unto future generations” (3 Ne. 26:1–2).

Here we see that the law of tithing is not a remote Old Testament practice but a commandment directly from the Savior to the people of our day. The Lord reaffirmed that law in modern revelation, commanding his people to pay “one-tenth of all their interest annually” and declaring that “this shall be a standing law unto them forever” (D\&C 119:4).

No prophet of the Lord in modern times has preached the law of tithing more fervently than Heber J. Grant. As an Apostle and later as President of the Church, he frequently called upon the Saints to pay an honest tithe and made firm promises to those who would do so.

In a general conference in 1912, Elder Heber J. Grant declared:

“I bear witness—and I know that the witness I bear is true—that [to] the men and the women who have been absolutely honest with God, who have paid their tithing, … God has given … wisdom whereby they have been able to utilize the remaining nine-tenths, and it has been of greater value to them, and they have accomplished more with it than they would if they had not been honest with the Lord” (in Conference Report, Apr. 1912, p. 30).

In 1929, President Heber J. Grant said:

“I appeal to the Latter-day Saints to be honest with the Lord and I promise them that peace, prosperity and financial success will attend those who are honest with our Heavenly Father. … When we set our hearts upon the things of this world and fail to be strictly honest with the Lord we do not grow in the light and power and strength of the gospel as we otherwise would do” (in Conference Report, Oct. 1929, pp. 4–5).

During the Great Depression, President Grant continued to remind the Saints that the payment of tithing would open the windows of heaven for blessings needed by the faithful. In that stressful period, some of our bishops observed that members who paid their tithing were able to support their families more effectively than those who did not. The tithe payers tended to keep their employment, enjoy good health, and be free from the most devastating effects of economic and spiritual depression (see “Tithing Is a Builder of Faith,” Church News, 9 Dec. 1961, p. 16). Countless tithe-paying Latter-day Saints can testify to similar blessings today.

I am grateful to President Grant and other prophets for teaching the principle of tithing to my parents and to them for teaching it to me. My attitude toward the law of tithing was set in place by the example and words of my mother, illustrated in a conversation I remember from my youth.

During World War II, my widowed mother supported her three young children on a schoolteacher’s salary that was meager. When I became conscious that we went without some desirable things because we didn’t have enough money, I asked my mother why she paid so much of her salary as tithing. I have never forgotten her explanation: “Dallin, there might be some people who can get along without paying tithing, but we can’t. The Lord has chosen to take your father and leave me to raise you children. I cannot do that without the blessings of the Lord, and I obtain those blessings by paying an honest tithing. When I pay my tithing, I have the Lord’s promise that he will bless us, and we must have those blessings if we are to get along.”

Years later I read President Joseph F. Smith’s memory of a similar testimony and teaching by his widowed mother. In the April 1900 conference, President Smith shared this memory from his childhood:

“My mother was a widow, with a large family to provide for. One spring when we opened our potato pits she had her boys get a load of the best potatoes, and she took them to the tithing office; potatoes were scarce that season. I was a little boy at the time, and drove the team. When we drove up to the steps of the tithing office, ready to unload the potatoes, one of the clerks came out and said to my mother, ‘Widow Smith, it’s a shame that you should have to pay tithing.’ … He chided my mother for paying her tithing, called her anything but wise or prudent; and said there were others who were strong and able to work that were supported from the tithing office. My mother turned upon him and said: ‘William, you ought to be ashamed of yourself. Would you deny me a blessing? If I did not pay my tithing, I should expect the Lord to withhold His blessings from me. I pay my tithing, not only because it is a law of God, but because I expect a blessing by doing it. By keeping this and other laws, I expect to prosper and to be able to provide for my family’” (in Conference Report, Apr. 1900, p. 48).

Some people say, “I can’t afford to pay tithing.” Those who place their faith in the Lord’s promises say, “I can’t afford not to pay tithing.”

Some time ago I was speaking to a meeting of Church leaders in a country outside of North America. As I spoke about tithing, I found myself saying something I had not intended to say. I told them the Lord was grieved that only a small fraction of the members in their nations relied on the Lord’s promises and paid a full tithing. I warned that the Lord would withhold material and spiritual blessings when his covenant children were not keeping this vital commandment.

I hope those leaders taught that principle to the members of the stakes and districts in their countries. The law of tithing and the promise of blessings to those who live it apply to the people of the Lord in every nation. I hope our members will qualify for the blessings of the Lord by paying a full tithing.

Tithing is a commandment with a promise. The words of Malachi, reaffirmed by the Savior, promise those who bring their tithes into the storehouse that the Lord will open “the windows of heaven, and pour [them] out a blessing, that there shall not be room enough to receive it.” The promised blessings are temporal and spiritual. The Lord promises to “rebuke the devourer,” and he also promises tithe payers that “all nations shall call you blessed, for ye shall be a delightsome land” (3 Ne. 24:10–12; see also Mal. 3:10–12).

I believe these are promises to the nations in which we reside. When the people of God withheld their tithes and offerings, Malachi condemned “this whole nation” (Mal. 3:9). Similarly, I believe that when many citizens of a nation are faithful in the payment of tithes, they summon the blessings of heaven upon their entire nation. The Bible teaches that “righteousness exalteth a nation” (Prov. 14:34) and “a little leaven leaveneth the whole lump” (Gal. 5:9; see also Matt. 13:33).

The payment of tithing also brings the individual tithe payer unique spiritual blessings. Tithe paying is evidence that we accept the law of sacrifice. It also prepares us for the law of consecration and the other higher laws of the celestial kingdom. The Lectures on Faith, prepared by the early leaders of the restored Church, part the curtain on that subject when they say:

“Let us here observe that a religion that does not require the sacrifice of all things never has power sufficient to produce the faith necessary unto life and salvation; for, from the first existence of man, the faith necessary unto the enjoyment of life and salvation never could be obtained without the sacrifice of all earthly things” (Lectures on Faith, 6:7).

We should not think that the payment and blessings of tithing are unique to members of The Church of Jesus Christ of Latter-day Saints. Tithe paying is commanded in the Bible. Abraham paid tithes to Melchizedek (see Gen. 14:20). Jacob covenanted to “give the tenth” unto God (Gen. 28:22). After the children of Israel were brought out of Egypt, the prophet Moses commanded that they should give a tenth to the Lord (see Lev. 27:30–34).

The Savior reaffirmed that teaching when the Pharisees asked him whether it was lawful to pay taxes. The Savior replied with this command: “Render therefore unto Caesar the things which are Caesar’s; and unto God the things that are God’s” (Matt. 22:21).

A few years ago the New York Times carried a feature article on a dozen highly paid professional athletes who were giving a fixed share (usually 10 percent) of their income to their church (see George Vecsey, “As They Look Past Their Riches, Athletes Are Turning to Religion,” New York Times, 29 Apr. 1991, pp. A1, C7). None of the featured athletes was LDS. If the names of our tithe-paying LDS professional athletes had been added to the list, it would have been much longer.

There are accounts of good Christian businessmen who promised to give the Lord a share of their profits and then attributed their business success to the fact that the Lord was their partner (see Betty Munson, “His Two Strips of Wheat,” Guideposts, Dec. 1991, pp. 24–27; William G. Shepherd, “Men Who Tithe,” Improvement Era, June 1928, pp. 633–45). BYU president Ernest L. Wilkinson, who often spoke of the blessings he had received from paying his tithing, quoted this statement from a non-Mormon businessman:

“We would not lend a neighbor money with which to run his business without interest. Neither would we expect him to lend us money without paying interest. I found I was using God’s money and the business talents He had given me without paying Him interest. That’s all I’ve done in tithing—just met my interest obligations!” (“The Principle and Practice of Paying Tithing,” Brigham Young University Bulletin, 10 Dec. 1957, pp. 10–11).

In the Lord’s commandment to the people of this day, tithing is “‘one-tenth of all their interest annually,’ which is understood to mean income.” The First Presidency has said, “No one is justified in making any other statement than this” (First Presidency letter, 19 Mar. 1970, quoted in the General Handbook of Instructions [1989], p. 9-1; see also D\&C 119).

We pay tithing, as the Savior taught, by bringing the tithes “into the storehouse” (3 Ne. 24:10; Mal. 3:10). We do this by paying our tithing to our bishop or branch president. We do not pay tithing by contributing to our favorite charities. The contributions we should make to charities come from our own funds, not from the tithes we are commanded to pay to the storehouse of the Lord.

The Lord has directed by revelation that the expenditure of his tithes will be directed by his servants, the First Presidency, the Quorum of the Twelve, and the Presiding Bishopric (see D\&C 120). Those funds are spent to build and maintain temples and houses of worship, to conduct our worldwide missionary work, to translate and publish scriptures, to provide resources to redeem the dead, to fund religious education, and to support other Church purposes selected by the designated servants of the Lord.

In earlier times, tithing was paid in kind—a tenth of the herdsman’s increase, a tenth of the farmer’s produce. I am sorry that our modern cash economy deprives parents of the wonderful teaching opportunities presented by the payment of tithing in kind. In a recent book, Tongan Saints: Legacy of Faith, the author quotes a Tongan bishop’s memories of one such example:

“Grandpa Vanisi’s spirituality inspired an awe in me as a child. I remember following him daily to his plantation. He would always point out to me the very best of his taro, bananas, or yams and say: ‘These will be for our tithing.’ His greatest care was given to these ‘chosen’ ones. During the harvest, I was often the one assigned to take our load of tithing to the branch president. I remember sitting on the family horse. Grandfather would lift onto its back a sack of fine taro which I balanced in front of me. Then with a very serious look in his eyes, he said to me, ‘Simi, be very careful because this is our tithing.’ From my grandfather I learned early in life that you give only your best to the Lord” (in Eric B. Shumway, trans. and ed., Tongan Saints: Legacy of Faith [Laie, Hawaii: The Institute for Polynesian Studies, 1991], pp. 79–80).

I had a similar experience as a young boy on my grandparents’ farm. They taught me about tithing with examples of one egg or one bushel of peaches out of ten. Years later I used those same kinds of examples to try to teach the principles of tithing to our own children.

Parents are always looking for better ways to teach, and the results of their efforts are sometimes unexpected. Attempting to teach tithing to our young son, I explained the principle of a tenth and how it would apply to the eggs gathered in a chicken farm and the young calves or horses born in a breeding herd. When I finished what I was sure was a clear explanation, I wanted to test whether our seven-year-old had understood. I asked him to imagine that he was a farmer with a harvest of eggs and young animals. I supplied the figures and then asked our little boy what he would give to the bishop as tithing. He thought deeply for a moment and then said, “I would give him a very old horse.”

We obviously had some further conversations on the principle of tithing, and I am proud of the way he and his brother and sisters learned and practiced that principle. But I have often thought of that little boy’s words as I have observed how some adult Church members relate to the law of tithing. I think we still have some whose attitude and performance consist of giving the bishop something like “a very old horse.”

The payment of tithing is a test of priorities. The Savior taught that reality when he gave this parable:

“The ground of a certain rich man brought forth plentifully:

“And he thought within himself, saying, What shall I do, because I have no room where to bestow my fruits?

“And he said, This will I do: I will pull down my barns, and build greater; and there will I bestow all my fruits and my goods.

“And I will say to my soul, Soul, thou hast much goods laid up for many years; take thine ease, eat, drink, and be merry.

“But God said unto him, Thou fool, this night thy soul shall be required of thee: then whose shall those things be, which thou hast provided?

“So is he that layeth up treasure for himself, and is not rich toward God” (Luke 12:16–21).

A modern illustration of that principle is suggested in the apocryphal story of two men standing before the casket of a wealthy friend. Asked one, “How much property did he leave?” Replied the other, “He left all of it.”

President Lorenzo Snow taught that “the law of tithing is one of the most important ever revealed to man” (quoted in Le Roi C. Snow, “The Lord’s Way Out of Bondage,” Improvement Era, July 1938, p. 442). Faithful adherence to this law opens the windows of heaven for blessings temporal and spiritual. As a lifelong recipient of those blessings, I testify to the goodness of our God and his bounteous blessings to his children.

I pray that each member of this church will qualify for the blessings promised and bestowed on those who bring all their tithes into the storehouse. In the name of Jesus Christ, amen.

\newpage
\settalk{Worship through Music}
\setspeaker{Dallin H. Oaks}
\settalkdate{October 1994}
\talkheading{Worship through Music}{Dallin H. Oaks}

President Hunter, we have been thrilled by your inspired message. We express our love to you. We also congratulate the newly called and sustained General Authorities and general officers of the Church.

Our hearts have united with the Mormon Youth Chorus’s spirited singing of Charles Wesley’s inspired words, “Rejoice, the Lord is King! Your Lord and King adore!” (Hymns, 1985, no. 66). With the events of this solemn assembly, we have also felt the overwhelming gratitude expressed in our beloved hymn “We Thank Thee, O God, for a Prophet” (Hymns, 1985, no. 19). We have rejoiced in the privilege of sustaining President Howard W. Hunter as President of the Church and Presidents Hinckley and Monson as his Counselors. In this worldwide assembly, we have pledged our prayers and best efforts to support the men whom the Lord has called to lead his church. I testify that what we have done has been recorded in the heavens and that each of us will be accountable to God for the way we respond to the leadership we have sustained in this solemn and sacred way.

Last spring I made my first visit to Brasília, Brazil. Over three thousand Saints gathered for a regional conference. The printed program listed the musical numbers, but the Portuguese words meant nothing to me. But when their beautiful choir began to sing, the music crossed all barriers of language and spoke to my soul:

\begin{verse}
The morning breaks, the shadows flee;\\
Lo, Zion’s standard is unfurled! …\\
The dawning of a brighter day\\
Majestic rises on the world.
\end{verse}

[“The Morning Breaks,” Hymns, 1985, no. 1]

Through the miracle of sacred music, the Spirit of the Lord descended upon us, and we were made ready for gospel instruction and worship.

The First Presidency has said:

“Inspirational music is an essential part of our church meetings. The hymns invite the Spirit of the Lord, create a feeling of reverence, unify us as members, and provide a way for us to offer praises to the Lord.

“Some of the greatest sermons are preached by the singing of hymns. Hymns move us to repentance and good works, build testimony and faith, comfort the weary, console the mourning, and inspire us to endure to the end” (Hymns, 1985, p. ix).

The singing of hymns is one of the best ways to put ourselves in tune with the Spirit of the Lord. I wonder if we are making enough use of this heaven-sent resource in our meetings, in our classes, and in our homes.

Last July I visited the Church’s Polynesian Cultural Center in Hawaii. Before the evening show of dancing and music from various island cultures, I went backstage to thank the performers. I arrived during those frantic moments before the show began. Scores of performers were hurrying through the last-minute tasks required to coordinate their efforts in a fast-moving performance. I wondered how the director would bring this turmoil to order in preparation for my brief remarks.

It happened as if by miracle. On signal, one strong voice began, and the strains of “We Thank Thee, O God, for a Prophet” quickly swelled into a beautiful chorus as the uniquely talented young people brought their thoughts into harmony with the Lord.

We had a similar experience in our family. Last spring some of our children and fourteen of our grandchildren had a family outing in the mountains. One of our activities was a meeting to share experiences and testimonies. We gathered at the appointed time, but the little people were only gathered in body. The large spirits in those little bodies were clamoring for more of the exciting outdoor activities they had been enjoying. The cabin where we met was too small to contain them, and it seemed as if a dozen restless children and their outcries were ricocheting off the walls in every direction. Grandparents will appreciate the apprehension I felt at trying to sponsor something serious in that setting.

Suddenly the instinctive wisdom of young mothers rescued our efforts. Two mothers began to sing a song familiar to the children. Others joined in, and within a few minutes the mood had changed and all spirits were subdued and receptive to spiritual things. I offered a silent prayer of thanks for hymns and for mothers who know how to use them!

The singing of hymns is one of the best ways to learn the doctrine of the restored gospel. Elder Stephen D. Nadauld captured this unique strength in some lines he wrote and shared in a General Authority meeting:

\begin{verse}
If I would teach with power\\
The doctrine and the plan,\\
I’d wish for gentle music\\
To prepare the soul of man.
\end{verse}

The scriptures contain many affirmations that hymn singing is a glorious way to worship. Before the Savior and his Apostles left the upper room where they had the sublime experience of the Last Supper, they sang a hymn. After their hymn, the Savior led them to the Mount of Olives (see Matt. 26:30).

The Apostle Paul advised the Colossians that they should be “teaching and admonishing one another in psalms and hymns and spiritual songs, singing with grace in your hearts to the Lord” (Col. 3:16; see also Alma 26:8).

Modern revelation reaffirms the importance of sacred music. In one of the earliest revelations given through the Prophet Joseph Smith, the Lord appointed Emma Smith “to make a selection of sacred hymns, as it shall be given thee, which is pleasing unto me, to be had in my church.

“For my soul delighteth in the song of the heart; yea, the song of the righteous is a prayer unto me, and it shall be answered with a blessing upon their heads” (D\&C 25:11–12).

In a revelation given through another prophet a generation later, the Lord commanded his people to “praise the Lord with singing, [and] with music” (D\&C 136:28).

This direction to praise the Lord with singing is not limited to large meetings. When the Lord’s Apostles meet in modern times, the singing of hymns is still part of their meetings. The weekly meetings of the First Presidency and Quorum of the Twelve Apostles in the Salt Lake Temple always begin with a hymn. Elder Russell M. Nelson plays the organ accompaniment. Members of the First Presidency, who conduct these meetings, rotate the privilege of selecting the opening song. Most of us record the date each hymn is sung. According to my records, the opening song most frequently sung during the decade of my participation has been “I Need Thee Every Hour” (Hymns, 1985, no. 98). Picture the spiritual impact of a handful of the Lord’s servants singing that song before praying for his guidance in fulfilling their mighty responsibilities.

The veil is very thin in the temples, especially when we join in worshipping through music. At temple dedications I have seen more tears of joy elicited by music than by the spoken word. I have read accounts of angelic choirs joining in these hymns of praise, and I think I have experienced this on several occasions. In dedicatory sessions featuring beautiful and well-trained choirs of about thirty voices, there are times when I have heard what seemed to be ten times thirty voices praising God with a quality and intensity of feeling that can be experienced but not explained. Some who are listening today will know what I mean.

Sacred music has a unique capacity to communicate our feelings of love for the Lord. This kind of communication is a wonderful aid to our worship. Many have difficulty expressing worshipful feelings in words, but all can join in communicating such feelings through the inspired words of our hymns.

When a congregation worships through singing, all present should participate. Here I share another experience. I had finished a special assignment on a Sunday morning in Salt Lake City and desired to attend a sacrament meeting. I stopped at a convenient ward meetinghouse and slipped unnoticed into the overflow area just as the congregation was beginning to sing these sacred words of the sacrament song:

\begin{verse}
’Tis sweet to sing the matchless love\\
Of Him who left his home above\\
And came to earth—oh, wondrous plan—\\
To suffer, bleed, and die for man!
\end{verse}

[“’Tis Sweet to Sing the Matchless Love,” Hymns, 1985, no. 177]

My heart swelled as we sang this worshipful hymn and contemplated renewing our covenants by partaking of the sacrament. Our voices raised the concluding strains:

\begin{verse}
For Jesus died on Calvary,\\
That all thru him might ransomed be.\\
Then sing hosannas to his name;\\
Let heav’n and earth his love proclaim.
\end{verse}

As we sang these words, I glanced around at members of the congregation and was stunned to observe that about a third of them were not singing. How could this be? Were those who did not even mouth the words suggesting that for them it was not “sweet to sing the matchless love” or to “sing hosannas to his name”? What are we saying, what are we thinking, when we fail to join in singing in our worship services?

I believe some of us in North America are getting neglectful in our worship, including the singing of hymns. I have observed that the Saints elsewhere are more diligent in doing this. We in the center stakes of Zion should renew our fervent participation in the singing of our hymns.

There are a few conventions all of us should observe as we worship through music. As we sing we should think about the messages of the words. Our hymns contain matchless doctrinal sermons, surpassed only by the scriptures in their truth and poetic impact.

We depend on our choristers and organists to lead us at the prescribed pace. Too slow or too fast can detract from a worshipful mood.

We should be careful what music we use in settings where we desire to contribute to worship. Many musical numbers good for other wholesome settings are not appropriate for Church meetings.

Our hymns have been chosen because they have been proven effective to invite the Spirit of the Lord. A daughter who plays the violin described that reality. “I love to play classical music,” she said, “but when I play our hymns, I can just feel the Spirit of the Lord in my practice room.”

Soloists should remember that music in our worship services is not for demonstration but for worship. Vocal or instrumental numbers should be chosen to facilitate worship, not to provide performance opportunity for artists, no matter how accomplished.

Our sacred music prepares us to be taught the truths of the gospel. This is why we are selective in the kinds of music and the kinds of instruments we use in our worship services. This is why we encourage our choirs to use the hymnbook as their basic resource. We can make selective use of other music that is in harmony with the spirit of our hymns, such as Charles Gounod’s marvelous “O Divine Redeemer,” sung at the funeral of President Ezra Taft Benson. But a hymnbook’s hymn is often the most inspiring and appropriate musical selection for a choir, a vocalist, or an instrumentalist (see Michael F. Moody, “Conversation: Learning and Singing Hymns,” Ensign, Aug. 1994, p. 79).

Sacred music can help us even where there is no formal performance. For example, when temptation comes, we can neutralize its effect by humming or repeating the words of a favorite hymn (see Boyd K. Packer, in Conference Report, Oct. 1973, pp. 21–26; or Ensign, Jan. 1974, pp. 25–28).

Our hymns can work their miraculous effect even when the chorus of voices is few and even when hardly a sound can be heard. I felt this a few months ago as I participated in a musical performance that was unique in my church experience. I had been invited to speak at the Great Basin LDS Deaf Conference, hosted by the Salt Lake Valley (Deaf) Ward of the Salt Lake Park Stake. Over three hundred deaf brothers and sisters were in attendance. The members of the stake presidency and I were almost the only adults in the congregation who could hear and who attempted to sing audibly. The rest of that large assembly sang with their hands. Hardly a lip moved, and hardly a sound was heard except the organ and four faint voices from the stand. In the audience, all hands moved in unison with the leader as the audience signed, “The Spirit of God like a fire is burning!” (Hymns, 1985, no. 2). As we sang together, the Spirit of the Lord descended upon us, and we were made ready for prayer. Our sacred music is a powerful preparation for prayer and gospel teaching.

We need to make more use of our hymns to put us in tune with the Spirit of the Lord, to unify us, and to help us teach and learn our doctrine. We need to make better use of our hymns in missionary teaching, in gospel classes, in quorum meetings, in home evenings, and in home teaching visits. Music is an effective way to worship our Heavenly Father and his Son, Jesus Christ. We should use hymns when we need spiritual strength and inspiration.

We who have “felt to sing the song of redeeming love” (Alma 5:26) need to keep singing that we may draw ever closer to him who has inspired sacred music and commanded that it be used to worship him. May we be diligent in doing so is my humble prayer, which I offer with a testimony of the truthfulness of the gospel of Jesus Christ and of the divine calling of those we have sustained today. In the name of Jesus Christ, amen.

\newpage
\settalk{Apostasy and Restoration}
\setspeaker{Dallin H. Oaks}
\settalkdate{April 1995}
\talkheading{Apostasy and Restoration}{Dallin H. Oaks}

The Church of Jesus Christ of Latter-day Saints has many beliefs in common with other Christian churches. But we have differences, and those differences explain why we send missionaries to other Christians, why we build temples in addition to churches, and why our beliefs bring us such happiness and strength to deal with the challenges of life and death. I wish to speak about some of the important additions our doctrines make to the Christian faith. My subject is apostasy and restoration.

Last year searchers discovered a Roman fort and city in the Sinai close to the Suez Canal. Though once a major city, its location had been covered by desert sands and its existence had been forgotten for hundreds of years (see “Remains of Roman Fortress Emerge from Sinai Desert,” Deseret News, 6 Oct. 1994, p. A20). Discoveries like this contradict the common assumption that knowledge increases with the passage of time. In fact, on some matters the general knowledge of mankind regresses as some important truths are distorted or ignored and eventually forgotten. For example, the American Indians were in many respects more successful at living in harmony with nature than our modern society. Similarly, modern artists and craftsmen have been unable to recapture some of the superior techniques and materials of the past, like the varnish on a Stradivarius violin.

We would be wiser if we could restore the knowledge of some important things that have been distorted, ignored, or forgotten. This also applies to religious knowledge. It explains the need for the gospel restoration we proclaim.

When Joseph Smith was asked to explain the major tenets of our faith, he wrote what we now call the Articles of Faith. The first article states, “We believe in God, the Eternal Father, and in His Son, Jesus Christ, and in the Holy Ghost.” The Prophet later declared that “the simple and first principles of the gospel” include knowing “for a certainty the character of God” (“Conference Minutes,” Times and Seasons, 15 Aug. 1844, p. 614). We must begin with the truth about God and our relationship to him. Everything else follows from that.

In common with the rest of Christianity, we believe in a Godhead of Father, Son, and Holy Ghost. However, we testify that these three members of the Godhead are three separate and distinct beings. We also testify that God the Father is not just a spirit but is a glorified person with a tangible body, as is his resurrected Son, Jesus Christ.

When first communicated to mankind by prophets, the teachings we now have in the Bible were “plain and pure, and most precious and easy” to understand (1 Ne. 14:23). Even in the transmitted and translated version we have today, the Bible language confirms that God the Father and his resurrected Son, Jesus Christ, are tangible, separate beings. To cite only two of many such teachings, the Bible declares that man was created in the image of God, and it describes three separate members of the Godhead manifested at the baptism of Jesus (see Gen. 1:27; Matt. 3:13–17).

In contrast, many Christians reject the idea of a tangible, personal God and a Godhead of three separate beings. They believe that God is a spirit and that the Godhead is only one God. In our view, these concepts are evidence of the falling away we call the Great Apostasy.

We maintain that the concepts identified by such nonscriptural terms as “the incomprehensible mystery of God” and “the mystery of the Holy Trinity” are attributable to the ideas of Greek philosophy. These philosophical concepts transformed Christianity in the first few centuries following the deaths of the Apostles. For example, philosophers then maintained that physical matter was evil and that God was a spirit without feelings or passions. Persons of this persuasion, including learned men who became influential converts to Christianity, had a hard time accepting the simple teachings of early Christianity: an Only Begotten Son who said he was in the express image of his Father in Heaven and who taught his followers to be one as he and his Father were one, and a Messiah who died on a cross and later appeared to his followers as a resurrected being with flesh and bones.

The collision between the speculative world of Greek philosophy and the simple, literal faith and practice of the earliest Christians produced sharp contentions that threatened to widen political divisions in the fragmenting Roman empire. This led Emperor Constantine to convene the first churchwide council in a.d. 325. The action of this council of Nicaea remains the most important single event after the death of the Apostles in formulating the modern Christian concept of deity. The Nicene Creed erased the idea of the separate being of Father and Son by defining God the Son as being of “one substance with the Father.”

Other councils followed, and from their decisions and the writings of churchmen and philosophers there came a synthesis of Greek philosophy and Christian doctrine in which the orthodox Christians of that day lost the fulness of truth about the nature of God and the Godhead. The consequences persist in the various creeds of Christianity, which declare a Godhead of only one being and which describe that single being or God as “incomprehensible” and “without body, parts, or passions.” One of the distinguishing features of the doctrine of The Church of Jesus Christ of Latter-day Saints is its rejection of all of these postbiblical creeds (see Stephen E. Robinson, Are Mormons Christians? [Salt Lake City: Bookcraft, 1991]; Encyclopedia of Mormonism, ed. Daniel H. Ludlow, 5 vols. [New York: Macmillan Publishing Co., 1992], 1:56–58, 1:393–404, 2:548–53).

In the process of what we call the Apostasy, the tangible, personal God described in the Old and New Testaments was replaced by the abstract, incomprehensible deity defined by compromise with the speculative principles of Greek philosophy. The received language of the Bible remained, but the so-called hidden meanings of scriptural words were now explained in the vocabulary of a philosophy alien to their origins. In the language of that philosophy, God the Father ceased to be a Father in any but an allegorical sense. He ceased to exist as a comprehensible and compassionate being. And the separate identity of his Only Begotten Son was swallowed up in a philosophical abstraction that attempted to define a common substance and an incomprehensible relationship.

These descriptions of a religious philosophy are surely undiplomatic, but I hasten to add that Latter-day Saints do not apply such criticism to the men and women who profess these beliefs. We believe that most religious leaders and followers are sincere believers who love God and understand and serve him to the best of their abilities. We are indebted to the men and women who kept the light of faith and learning alive through the centuries to the present day. We have only to contrast the lesser light that exists among peoples unfamiliar with the names of God and Jesus Christ to realize the great contribution made by Christian teachers through the ages. We honor them as servants of God.

Then came the First Vision. An unschooled boy, seeking knowledge from the ultimate source, saw two personages of indescribable brightness and glory and heard one of them say, while pointing to the other, “This is My Beloved Son. Hear Him!” (JS—H 1:17). The divine teaching in that vision began the restoration of the fulness of the gospel of Jesus Christ. God the Son told the boy prophet that all the “creeds” of the churches of that day “were an abomination in his sight” (JS—H 1:19). We affirm that this divine declaration was a condemnation of the creeds, not of the faithful seekers who believed in them. Joseph Smith’s first vision showed that the prevailing concepts of the nature of God and the Godhead were untrue and could not lead their adherents to the destiny God desired for them.

After a subsequent outpouring of modern scripture and revelation, this modern prophet declared, “The Father has a body of flesh and bones as tangible as man’s; the Son also; but the Holy Ghost has not a body of flesh and bones, but is a personage of Spirit” (D\&C 130:22).

This belief does not mean that we claim sufficient spiritual maturity to comprehend God. Nor do we equate our imperfect mortal bodies to his immortal, glorified being. But we can comprehend the fundamentals he has revealed about himself and the other members of the Godhead. And that knowledge is essential to our understanding of the purpose of mortal life and of our eternal destiny as resurrected beings after mortal life.

In the theology of the restored church of Jesus Christ, the purpose of mortal life is to prepare us to realize our destiny as sons and daughters of God—to become like him. Joseph Smith and Brigham Young both taught that “no man … can know himself unless he knows God, and he can not know God unless he knows himself” (in Journal of Discourses, 16:75; see also The Words of Joseph Smith, ed. Andrew F. Ehat and Lyndon W. Cook [Provo: Religious Studies Center, Brigham Young University, 1980], p. 340). The Bible describes mortals as “the children of God” and as “heirs of God, and joint-heirs with Christ” (Rom. 8:16–17). It also declares that “we suffer with him, that we may be also glorified together” (Rom. 8:17) and that “when he shall appear, we shall be like him” (1 Jn. 3:2). We take these Bible teachings literally. We believe that the purpose of mortal life is to acquire a physical body and, through the atonement of Jesus Christ and by obedience to the laws and ordinances of the gospel, to qualify for the glorified, resurrected celestial state that is called exaltation or eternal life.

Like other Christians, we believe in a heaven or paradise and a hell following mortal life, but to us that two-part division of the righteous and the wicked is merely temporary while the spirits of the dead await their resurrections and final judgments. The destinations that follow the final judgments are much more diverse. Our restored knowledge of the separateness of the three members of the Godhead provides a key to help us understand the diversities of resurrected glory.

In their final judgment, the children of God will be assigned to a kingdom of glory for which their obedience has qualified them. In his letters to the Corinthians, the Apostle Paul described these places. He told of a vision in which he was “caught up to the third heaven” and “heard unspeakable words, which it is not lawful for a man to utter” (2 Cor. 12:2, 4). Speaking of the resurrection of the dead, he described “celestial bodies,” “bodies terrestrial” (1 Cor. 15:40), and “bodies telestial” (JST, 1 Cor. 15:40), each pertaining to a different degree of glory. He likened these different glories to the sun, to the moon, and to different stars (see 1 Cor. 15:41).

We learn from modern revelation that these three different degrees of glory have a special relationship to the three different members of the Godhead.

The lowest degree is the telestial domain of those who “received not the gospel, neither the testimony of Jesus, neither the prophets” (D\&C 76:101) and who have had to suffer for their wickedness. But even this degree has a glory that “surpasses all understanding” (D\&C 76:89). Its occupants receive the Holy Spirit and the administering of angels, for even those who have been wicked will ultimately be “heirs of [this degree of] salvation” (D\&C 76:88).

The next higher degree of glory, the terrestrial, “excels in all things the glory of the telestial, even in glory, and in power, and in might, and in dominion” (D\&C 76:91). The terrestrial is the abode of those who were the “honorable men of the earth” (D\&C 76:75). Its most distinguishing feature is that those who qualify for terrestrial glory “receive of the presence of the Son” (D\&C 76:77). Concepts familiar to all Christians might liken this higher kingdom to heaven because it has the presence of the Son.

In contrast to traditional Christianity, we join with Paul in affirming the existence of a third or higher heaven. Modern revelation describes it as the celestial kingdom—the abode of those “whose bodies are celestial, whose glory is that of the sun, even the glory of God” (D\&C 76:70). Those who qualify for this kingdom of glory “shall dwell in the presence of God and his Christ forever and ever” (D\&C 76:62). Those who have met the highest requirements for this kingdom, including faithfulness to covenants made in a temple of God and marriage for eternity, will be exalted to the godlike state referred to as the “fulness” of the Father or eternal life (D\&C 76:56, 94; see also D\&C 131; 132:19–20). (This destiny of eternal life or God’s life should be familiar to all who have studied the ancient Christian doctrine of and belief in deification or apotheosis.) For us, eternal life is not a mystical union with an incomprehensible spirit-god. Eternal life is family life with a loving Father in Heaven and with our progenitors and our posterity.

The theology of the restored gospel of Jesus Christ is comprehensive, universal, merciful, and true. Following the necessary experience of mortal life, all sons and daughters of God will ultimately be resurrected and go to a kingdom of glory. The righteous—regardless of current religious denomination or belief—will ultimately go to a kingdom of glory more wonderful than any of us can comprehend. Even the wicked, or almost all of them, will ultimately go to a marvelous—though lesser—kingdom of glory. All of that will occur because of God’s love for his children and because of the atonement and resurrection of Jesus Christ, “who glorifies the Father, and saves all the works of his hands” (D\&C 76:43).

The purpose of The Church of Jesus Christ of Latter-day Saints is to help all of the children of God understand their potential and achieve their highest destiny. This church exists to provide the sons and daughters of God with the means of entrance into and exaltation in the celestial kingdom. This is a family-centered church in doctrine and practices. Our understanding of the nature and purpose of God the Eternal Father explains our destiny and our relationship in his eternal family. Our theology begins with heavenly parents. Our highest aspiration is to be like them. Under the merciful plan of the Father, all of this is possible through the atonement of the Only Begotten of the Father, our Lord and Savior, Jesus Christ. As earthly parents we participate in the gospel plan by providing mortal bodies for the spirit children of God. The fulness of eternal salvation is a family matter.

It is the reality of these glorious possibilities that causes us to proclaim our message of restored Christianity to all people, even to good practicing Christians with other beliefs. This is why we build temples. This is the faith that gives us strength and joy to confront the challenges of mortal life. We offer these truths and opportunities to all people and testify to their truthfulness in the name of Jesus Christ, amen.

\newpage
\settalk{Powerful Ideas}
\setspeaker{Dallin H. Oaks}
\settalkdate{October 1995}
\talkheading{Powerful Ideas}{Dallin H. Oaks}

Last summer I attended the funeral of an elect lady. One speaker described three of her great qualities: loyalty, obedience, and faith. As he elaborated on her life, I thought how appropriate it was to speak of such powerful qualities in a funeral tribute. A life is not a trivial thing, and its passing should not be memorialized with trivial things. A funeral service is a time to speak of powerful ideas—ideas that can appropriately stand beside the importance of life, ideas that are powerful in their influence on those who remain behind.

As I enjoyed the spirit of this inspiring funeral, my thoughts were directed toward the application of this principle in other settings. Parents should also teach powerful ideas. So should home teachers, visiting teachers, and the teachers in various classes. The Savior warned that we will be judged for “every idle word that [we] shall speak” (Matt. 12:36). Modern revelation commands us to cease from “light speeches” and “light-mindedness” (D\&C 88:121) and to cast away “idle thoughts” and “excess of laughter” (D\&C 88:69). There are plenty of other spokesmen for trivial things. Latter-day Saints should be constantly concerned with teaching and emphasizing those great and powerful eternal truths that will help us find our way back to the presence of our Heavenly Father.

About thirty years ago, some scholars authored a book on general education—the body of knowledge expected of all educated persons. Its title, The Knowledge Most Worth Having, is a good reminder of the fact that knowledge is not of equal value (Wayne C. Booth, ed. [Chicago: University of Chicago Press, 1967]). Some knowledge is more important than others. That principle also applies to what we call spiritual knowledge.

Consider the power of the idea taught in our beloved song “I Am a Child of God” (Hymns, 1985, no. 301), sung so impressively by the choir at the beginning of this session. Here is the answer to one of life’s great questions, “Who am I?” I am a child of God with a spirit lineage to heavenly parents. That parentage defines our eternal potential. That powerful idea is a potent antidepressant. It can strengthen each of us to make righteous choices and to seek the best that is within us. Establish in the mind of a young person the powerful idea that he or she is a child of God, and you have given self-respect and motivation to move against the problems of life.

When we understand our relationship to God, we also understand our relationship to one another. All men and women on this earth are the offspring of God, spirit brothers and sisters. What a powerful idea! No wonder God’s Only Begotten Son commanded us to love one another. If only we could do so! What a different world it would be if brotherly and sisterly love and unselfish assistance could transcend all boundaries of nation, creed, and color. Such love would not erase all differences of opinion and action, but it would encourage each of us to focus our opposition on actions rather than actors.

The eternal truth that our Heavenly Father loves all His children is an immensely powerful idea. It is especially powerful when children can visualize it through the love and sacrifice of their earthly parents. Love is the most powerful force in the world. Arthur Henry King has said: “Love is not just an ecstasy, not just an intense feeling. It is a driving force. It is something that carries us through our life of joyful duty” (The Abundance of the Heart [Salt Lake City: Bookcraft, 1986], p. 84).

We all have our own examples of the power of love. More than twenty-five years ago I recorded some memories I had of my father, who died before I was eight years old. What I wrote then illustrates the power of love in the life of a boy:

“The strongest impression I have of my relationship with my father I cannot document with any event or any words I can recall. It is a feeling. Based on words and actions long since lost to mind, this feeling persists with all the clarity of perfect faith. He loved me and he was proud of me. … That is the kind of memory a boy can treasure, and also a man” (“Memories of My Father,” 15 Oct. 1967).

Another powerful idea we should teach one another is that mortal life has a purpose and that mortal death is not the end but only a transition to the next phase of an existence that is immortal. President Brigham Young taught that “our existence here is for the sole purpose of exaltation and restoration to the presence of our Father and God” (Discourses of Brigham Young, sel. John A. Widtsoe [Salt Lake City: Deseret Book, 1941], p. 37). The idea of eternal progress is one of the most powerful ideas in our theology. It gives us hope when we falter and challenge when we soar. Surely this is one of the great “solemnities of eternity” that we are commanded to let “rest upon [our] minds” (D\&C 43:34).

Another idea that is powerful to lift us from discouragement is that the work of The Church of Jesus Christ of Latter-day Saints, “to bring to pass the … eternal life of man” (Moses 1:39), is an eternal work. Not all problems are overcome and not all needed relationships are fixed in mortality. The work of salvation goes on beyond the veil of death, and we should not be too apprehensive about incompleteness within the limits of mortality.

A powerful idea with immediate practical application is the reality that we can pray to our Heavenly Father and He will hear our prayers and help us in the way that is best for us. Most of us have experienced the terrible empty feeling that comes from being separated from those who love us. If we remember that we can pray and be heard and helped, we can always withstand that feeling of emptiness. We can always be in touch with a powerful friend who loves us and helps us in His own time and in His own way.

Thousands of experiences show that we can pray and have our prayers answered. Some of the choicest involve young children. In the biography of President Spencer W. Kimball we read:

“Again and again Spencer watched his parents take their problems to the Lord. One day when Spencer was five and out doing his chores, little one-year-old Fannie wandered from the house and was lost. No one could find her. Clare, sixteen, said, ‘Ma, if we pray, the Lord will direct us to Fannie.’ So the mother and children prayed. Immediately after the prayer Gordon walked to the very spot where Fannie was fast asleep in a large box behind the chicken coop. ‘We thanked our Heavenly Father over and over,’ Olive recorded in her journal” (in Edward L. Kimball and Andrew E. Kimball, Jr., Spencer W. Kimball [Salt Lake City: Bookcraft, 1977], p. 31).

Every follower of Jesus Christ knows that the most powerful ideas of the Christian faith are the resurrection and the atonement of Jesus Christ. Because of Him we can be forgiven of our sins and we will live again. Those powerful ideas have been explained in countless sermons from this pulpit and a million others. They are well known but not well applied in the lives of most of us.

Our model is not the latest popular hero of sports or entertainment, not our accumulated property or prestige, and not the expensive toys and diversions that encourage us to concentrate on what is temporary and forget what is eternal. Our model—our first priority—is Jesus Christ. We must testify of Him and teach one another how we can apply His teachings and His example in our lives.

Brigham Young gave us some practical advice on how to do this. “The difference between God and the Devil,” he said, “is that God creates and organizes, while the whole study of the Devil is to destroy” (Discourses of Brigham Young, p. 69). In that contrast we have an important example of the reality of “opposition in all things” (2 Ne. 2:11).

Remember, our Savior, Jesus Christ, always builds us up and never tears us down. We should apply the power of that example in the ways we use our time, including our recreation and diversions. Consider the themes of the books, magazines, movies, television, and music we make popular by our patronage. Do the purposes and actions portrayed in our chosen entertainment build up or tear down the children of God? During my lifetime I have seen a strong trend to displace what builds up and dignifies the children of God with portrayals and performances that are depressing, demeaning, and destructive.

The powerful idea in this example is that whatever builds people up serves the cause of the Master, and whatever tears people down serves the cause of the adversary. We support one cause or the other every day by our patronage. This should remind us of our responsibility and motivate us toward fulfilling it in a way that would be pleasing to Him whose suffering offers us hope and whose example should give us direction.

We should always put the Savior first. The first commandment Jehovah gave to the children of Israel was, “Thou shalt have no other gods before me” (Ex. 20:3). This seems like a simple idea, but in practice many find it difficult.

It is surprisingly easy to take what should be our first devotion and subordinate it to other priorities. Fifty years ago, the Christian philosopher C. S. Lewis illustrated that tendency with an example that is distressingly applicable in our own day. In his book The Screwtape Letters, a senior devil explains how to corrupt Christians and frustrate the work of Jesus Christ. One letter explains how any “extreme devotion” can lead Christians away from the Lord and the practice of Christianity. Lewis gives two examples, extreme patriotism or extreme pacifism, and explains how either “extreme devotion” can corrupt its adherent:

“Let him begin by treating the Patriotism or the Pacifism as a part of his religion. Then let him, under the influence of partisan spirit, come to regard it as the most important part. Then quietly and gradually nurse him on to the stage at which the religion becomes merely part of the ‘Cause,’ in which Christianity is valued chiefly because of the excellent arguments it can produce in favour of the British war effort or of pacifism. … Once you have made the World an end, and faith a means, you have almost won your man, and it makes very little difference what kind of worldly end he is pursuing” (C. S. Lewis, The Screwtape Letters, rev. ed. [New York: Macmillan, 1982], p. 35).

We can readily see that tendency in our own time, with many causes that are good in themselves but become spiritually corrupting when they assume priorities ahead of Him who commanded, “Thou shalt have no other gods before me.” Jesus Christ and His work come first. Anything that would use Him or His kingdom or His church as a means to an end serves the cause of the adversary.

Two other powerful ideas were given voice by a noble young woman who survived a terrible experience. Virginia Reed was a survivor of the tragic Donner-Reed party, which made one of the earliest wagon treks into California. If this wagon train had followed the established Oregon Trail from Fort Bridger (Wyoming) northwest to Fort Hall (Idaho) and then southwest toward California, they would have reached their destination in safety. Instead, they were misled by a promoter. Lansford W. Hastings persuaded them they could save significant distance and time by following his so-called Hastings Cutoff. The Donner-Reed party left the proven trail at Fort Bridger and struggled southwest. They blazed a trail through the rugged Wasatch Mountains and then south of the Great Salt Lake and westward over the soggy surface of the salt flats in furnace heat.

The delays and incredible energies expended on this unproven route cost the Donner-Reed party an extra month in reaching the Sierra Nevada Mountains. As they hastened up the eastern slope trying to beat the first snows, they were caught in a tragic winter storm only one day short of the summit and a downhill passage into California. Marooned for the winter, half their group perished from starvation and cold.

After months in the mountains and incredible hardships of hunger and terror, thirteen-year-old Virginia Reed reached California and sent a letter to her cousin in the Midwest. After recounting her experiences and the terrible sufferings of their party, she concluded with this wise advice: “Never take no cutofs and hury along as fast as you can” (Virginia E. B. Reed letter to her cousin Mary Gillespie, 16 May 1847, quoted in J. Roderic Korns and Dale L. Morgan, eds., West from Fort Bridger [Logan, Utah: Utah State University Press, 1994], p. 238).

That is powerful and true advice, especially for teenagers. Young people are surrounded by many beckoning paths and many persuasive promoters who offer advice and cutoffs as substitutes for the proven way. “Try out this detour” and “tarry here for a while” are familiar proposals on the journey of life. My young friends, remember Virginia Reed’s advice: “Never take no cutofs and hury along as fast as you can.”

I conclude with an example from the life of the Apostle Paul. During his ministry he was exposed to ample light-mindedness, idle thoughts, and trivial things. In Athens he observed that “all the Athenians and strangers which were there [in his immediate vicinity] spent their time in nothing else, but … to tell, or to hear some new thing” (Acts 17:21). Paul’s determination to focus on powerful ideas is evident in one of his letters to the Saints in Corinth. He had not come “with excellency of speech or of wisdom,” he reminded them. “For I determined not to know any thing among you, save Jesus Christ, and him crucified” (1 Cor. 2:1–2).

Let us follow the commandments of God and the examples of his servants. Let us focus our teachings on those great and powerful ideas that have eternal significance in promoting righteousness, building up the children of God, and helping each of us toward our destiny of eternal life. That we may do so is my fervent prayer, in the name of Jesus Christ, amen.

\newpage
\settalk{Joseph, the Man and the Prophet}
\setspeaker{Dallin H. Oaks}
\settalkdate{April 1996}
\talkheading{Joseph, the Man and the Prophet}{Dallin H. Oaks}

On this beautiful Easter Sunday, I have chosen to speak about the Prophet Joseph Smith and to emphasize some lesser-known aspects of his life that further affirm his prophetic calling.

During my college studies at BYU, I was introduced to the History of the Church, an edited compilation of the writings of Joseph Smith and others. After I graduated from law school, I carefully read all seven volumes. I also pursued personal research in original records in Illinois, where the Prophet Joseph lived the last five years of his life.

The man I came to know in this way was not the man I had imagined. When I was a boy, growing up in the Church, I imagined the Prophet Joseph to be old and dignified and distant. But the Joseph Smith I met in my reading and personal research was a man of the frontier—young, emotional, dynamic, and so loved and approachable by his people that they often called him “Brother Joseph.” My studies strengthened my testimony of his prophetic calling. What a remarkable man! At the same time, I could see that he was mortal and therefore subject to sin and error, pain and affliction.

Overarching the Prophet Joseph’s entire ministry were his comparative youth, his superficial formal education, and his incredibly rapid acquisition of knowledge and maturity. He was 14 at the First Vision and 17 at the first visit from the angel Moroni. He was 21 when he received the golden plates and just 23 when he finished translating the Book of Mormon (in less than 60 working days). Over half of the revelations in our Doctrine and Covenants were given through the Prophet while he was 25 or younger. He was 26 when the First Presidency was organized and 30 when the Kirtland Temple was dedicated. He was just over 33 when he escaped his imprisonment in Missouri and resumed leadership of the Saints gathering in Nauvoo. He was 38 1/2 when he was martyred.

Joseph Smith had more than his share of mortal afflictions. When he was about seven, he suffered an excruciatingly painful surgery. Pieces of bone in his leg were chipped away without anesthetic. He was on crutches most of the next three years. Because of the poverty of his family, he had little formal education and as a youth was compelled to work long hours to help put food on the family table. The first three children of Joseph and his beloved Emma died at birth. A later child also died at birth and another in early childhood. Only four of Joseph and Emma’s nine birth children survived childhood. They also adopted twins, and one of these died as an infant.

Joseph was attacked physically on many occasions. He was often hounded by false charges. He was almost continually on the edge of financial distress. In the midst of trying to fulfill the staggering responsibilities of his sacred calling, he had to labor as a farmer or merchant to provide a living for his family. He did this without the remarkable spiritual gifts that sustained him in his prophetic calling. The Lord had advised him, “In temporal labors thou shalt not have strength, for this is not thy calling” (D\&C 24:9).

Speaking of his teen years following the First Vision, Joseph wrote, “I frequently fell into many foolish errors, and displayed the weakness of youth, and the foibles of human nature; which, I am sorry to say, led me into divers temptations, offensive in the sight of God.” He hastened to add that this behavior did not include “any great or malignant sins” since “a disposition to commit such was never in my nature” (JS—H 1:28).

Joseph’s candor about his shortcomings is evident in the fact that one of the first revelations he recorded in writing and published to the world was a crushing rebuke he received from the Lord. The first 116 manuscript pages of the Book of Mormon translation were lost because 22-year-old Joseph yielded to entreaties and loaned them to Martin Harris. “Behold,” the Lord declared, “how oft you have transgressed the commandments and the laws of God, and have gone on in the persuasions of men” (D\&C 3:6). The Lord told Joseph to repent or he would be stripped of his prophetic role (see D\&C 3:10–11). Four later revelations, also published by the Prophet, command him to “repent and walk more uprightly” (D\&C 5:21), speak of his having “sinned” (D\&C 64:7; see also D\&C 90:1), and rebuke him for not keeping the commandments (see D\&C 93:47).

The Prophet Joseph had no role models from whom he could learn how to be a prophet and leader of the Lord’s people. He learned from heavenly messengers and from the harvest of his unique spiritual gifts. He had to rely on associates who had no role models either. They struggled and learned together, and the Prophet’s growth was extremely rapid.

When Joseph warned the Saints against mortal imperfections, he did not raise himself above them, and they loved him for it. He cautioned a group of Saints newly arrived in Nauvoo against the tendency to be dissatisfied if everything was not done perfectly right. “He said he was but a man and they must not expect him to be perfect,” an associate recorded. “If they expected perfection from him, he should expect it from them, but if they would bear with his infirmities and the infirmities of the brethren, he would likewise bear with their infirmities” (The Papers of Joseph Smith, ed. Dean C. Jessee, 2 vols. [1992], 2:489).

Joseph had a “native cheery temperament” (JS—H 1:28). He delighted in the society of his friends. “He would play with the people,” one acquaintance recalled, “and he was always cheerful and happy” (Rachel Ridgeway Grant, in “Joseph Smith, the Prophet,” Young Woman’s Journal, Dec. 1905, 551). He loved little children and often frolicked with them in a manner shocking to some accustomed to the exaggerated sobriety of other ministers. These warm, human qualities caused some to deny Joseph’s prophetic role, but they endeared him to many who knew him. Our records contain numerous adoring tributes like that of an acquaintance who said, “The love the saints had for him was inexpressible” (Mary Alice Cannon Lambert, in “Joseph Smith, the Prophet,” Young Woman’s Journal, Dec. 1905, 554).

Despite his familiar and friendly style, the Prophet Joseph Smith was resolute in doing his duty. During a meeting to consider disciplining a man who had rejected the counsel of the Presidency and the Twelve, he declared: “The Saints need not think because I am familiar with them and am playful and cheerful, that I am ignorant of what is going on. Iniquity of any kind cannot be sustained in the Church, and it will not fare well where I am; for I am determined while I do lead the Church, to lead it right” (History of the Church, 5:411). On another occasion he wrote, “I am a lover of the cause of Christ and of virtue chastity and an upright steady course of conduct \& a holy walk, I despise a hypocrite or a covenant breaker” (The Personal Writings of Joseph Smith, ed. Dean C. Jessee [1984], 246).

All of his life, Joseph Smith lived on the frontier, where men had to pit their brute strength against nature and sometimes against one another. He was a large man, strong and physically active. He delighted in competitive sports, including pulling sticks—a test of physical strength (see History of the Church, 5:302). Our archives contain many recollections of his wrestling with friends and acquaintances. On one Sabbath, he and Brigham Young preached to the Saints in Ramus, Illinois, about a day’s ride from Nauvoo. On Monday, before departing Ramus, Joseph matched his wrestling prowess against a man someone described as “the bully of Ramus” (see Joseph Smith Journal, 13 Mar. 1843, recorded by Willard Richards, Joseph Smith Collection, LDS Church Archives). Joseph threw him. I am glad our current conference schedules do not provide the local members opportunity to test the visiting authorities in this manner.

Like most other leaders on the frontier, Joseph Smith did not shrink from physical confrontation, and he had the courage of a lion. Once he was kidnapped by two men who held cocked pistols to his head and repeatedly threatened to shoot him if he moved a muscle. The Prophet endured these threats for a time and then snapped back, “Shoot away; I have endured so much persecution and oppression that I am sick of life; why then don’t you shoot, and have done with it, instead of talking so much about it?” (in Journal of Discourses, 2:167; see also History of the Church, 5:440).

The Prophet Joseph Smith experienced severe opposition and persecution throughout his life, but in the midst of all of this he never wavered from his divine calling. During a public sermon in Nauvoo, he declared:

“The burdens which roll upon me are very great. My persecutors allow me no rest, and I find that in the midst of business and care the spirit is willing, but the flesh is weak. Although I was called of my Heavenly Father to lay the foundation of this great work and kingdom in this dispensation, and testify of His revealed will to scattered Israel, I am subject to like passions as other men, like the prophets of olden times” (History of the Church, 5:516).

In a sermon preached a little over a month before he was martyred, he declared, “I never told you I was perfect—but there is no error in the revelations which I have taught” (The Words of Joseph Smith, ed. Andrew F. Ehat and Lyndon W. Cook [1980], 369).

The event that focused anti-Mormon hostilities and led directly to the Martyrdom was the action of Mayor Joseph Smith and the city council in closing a newly established opposition newspaper in Nauvoo. Mormon historians—including Elder B. H. Roberts—had conceded that this action was illegal, but as a young law professor pursuing original research, I was pleased to find a legal basis for this action in the Illinois law of 1844. The amendment to the United States Constitution that extended the guarantee of freedom of the press to protect against the actions of city and state governments was not adopted until 1868, and it was not enforced as a matter of federal law until 1931. (See Dallin H. Oaks, “The Suppression of the Nauvoo Expositor,” Utah Law Review 9 [1965], 862–902.) We should judge the actions of our predecessors on the basis of the laws and commandments and circumstances of their day, not ours.

As students at the University of Chicago, historian Marvin S. Hill and I were intrigued with the little-known fact that five men went to trial in Illinois for the murders of Joseph and Hyrum Smith. For over 10 years we scoured libraries and archives across the nation to find every scrap of information about this trial and those involved in it. We studied the actions and words of Illinois citizens who knew Joseph Smith personally, some who hated him and plotted to kill him, and others who loved him and risked their lives to witness the trial of his accused assassins. Nothing in our discoveries in the original court records or in the testimony at the lengthy trial disclosed anything that reflected dishonor on the men who were murdered. (See Dallin H. Oaks and Marvin S. Hill, Carthage Conspiracy [1975].)

The accessibility of Illinois court records led to another untouched area of research on Joseph Smith—his financial activities. Joseph I. Bentley, then a law student at Chicago, and I discovered numerous records showing the business activities of Joseph Smith. As we explained in our article, this was a period following a nationwide financial panic and depression. Economic conditions in frontier states like Illinois were ruinous. The biographers of an Illinois contemporary, Abraham Lincoln, have described his financial embarrassments during this decade, when business was precarious, many obligations were in default, and lawsuits were common. The enemies of Joseph charged him with fraud in various property conveyances, mostly in behalf of the Church. A succession of court proceedings that extended for nearly a decade examined these claims in meticulous detail. Finally, in 1852, long after the Saints’ exodus from Illinois (so there was no conceivable political or other cause for anyone to favor the Prophet), a federal judge concluded this litigation with a decree that found no fraud or other moral impropriety by the Prophet. (See Dallin H. Oaks and Joseph I. Bentley, “Joseph Smith and Legal Process: In the Wake of the Steamboat Nauvoo,” Brigham Young University Law Review [1976], 735–82.) Independent of that decree, as one who has examined the hundreds of pages of allegations and evidence in these proceedings, I testify to the Prophet’s innocence of the charges against him.

As one familiar with early Illinois property law and as a lawyer enjoying the benefit of over 100 years of hindsight, I can readily see where Joseph and his fellow Church leaders and members were seriously disadvantaged by poor legal advice in some of the controversies just described. Bad legal advice may have been one of the causes for Brigham Young’s well-publicized negative opinions of lawyers. I have often chuckled at his 1845 declaration that he “would rather have a six-shooter than all the lawyers in Illinois” (History of the Church, 7:386).

Men who knew Joseph best and stood closest to him in Church leadership loved and sustained him as a prophet. His brother Hyrum chose to die at his side. John Taylor, also with him when he was murdered, said, “I testify before God, angels, and men, that he was a good, honorable, virtuous man … —that his private and public character was unimpeachable—and that he lived and died as a man of God” (The Gospel Kingdom, sel. G. Homer Durham [1987], 355; see also D\&C 135:3). Brigham Young declared, “I do not think that a man lives on the earth that knew [Joseph Smith] any better than I did; and I am bold to say that, Jesus Christ excepted, no better man ever lived or does live upon this earth” (in Journal of Discourses, 9:332).

Like other faithful Latter-day Saints, I have built my life on the testimony and mission of the Prophet Joseph Smith. In all of my reading and original research, I have never been dissuaded from my testimony of his prophetic calling and of the gospel and priesthood restoration the Lord initiated through him. I solemnly affirm the testimony Joseph Smith expressed in the famous Wentworth letter of 1842:

“The standard of truth has been erected; no unhallowed hand can stop the work from progressing; persecutions may rage, mobs may combine, armies may assemble, calumny may defame, but the truth of God will go forth boldly, nobly, and independent, till it has penetrated every continent, visited every clime, swept every country, and sounded in every ear, till the purposes of God shall be accomplished and the great Jehovah shall say the work is done” (Times and Seasons, 1 Mar. 1842, 709; in History of the Church, 4:540).

In the name of Jesus Christ, amen.

\newpage
\settalk{“Always Have His Spirit”}
\setspeaker{Dallin H. Oaks}
\settalkdate{October 1996}
\talkheading{“Always Have His Spirit”}{Dallin H. Oaks}

I have chosen to speak about the relationship between our partaking of the sacrament and our enjoying the blessings available from the gift of the Holy Ghost.

In modern revelation the Lord commanded, “That thou mayest more fully keep thyself unspotted from the world, thou shalt go to the house of prayer and offer up thy sacraments upon my holy day” (D\&C 59:9). As we partake of the sacrament each week, we ponder the Atonement of the Lord Jesus Christ and we reaffirm and renew the covenants we made when we were baptized. These acts of worship and commitment are described in the revealed prayer the priest offers upon the bread. As stated in that prayer, we partake of the bread “in remembrance of the body” of our Savior, and by doing so we witness to God, the Eternal Father, “that [we] are willing to take upon [us] the name of [his] Son, and always remember him and keep his commandments which he has given [us]” (D\&C 20:77).

After we were baptized, hands were laid upon our heads and we were given the gift of the Holy Ghost. When we consciously and sincerely renew our baptismal covenants as we partake of the sacrament, we renew our qualification for the promise “that [we] may always have his Spirit to be with [us]” (D\&C 20:77).

We cannot overstate the importance of that promise. President Wilford Woodruff called the gift of the Holy Ghost the greatest gift we can receive in mortality (see The Discourses of Wilford Woodruff, ed. G. Homer Durham [1990], 5). Unfortunately, the great value of that gift and the important conditions for its fulfillment are not well understood. Nephi prophesied that in the last days churches would be built up that would “teach with their learning, and deny the Holy Ghost, which giveth utterance” (2 Ne. 28:4). He also pronounced “wo” upon “him that hearkeneth unto the precepts of men, and denieth the power of God, and the gift of the Holy Ghost!” (2 Ne. 28:26).

The Bible tells us that when the Savior gave His final instructions to His disciples, He promised that He would send them “the Comforter” (John 16:7). Earlier He had taught them the mission of this Comforter, who is otherwise referred to as the Holy Ghost, the Holy Spirit, the Spirit of the Lord, or simply the Spirit. That Comforter dwells in us (see John 14:17). He teaches us all things and brings all things to our remembrance (see John 14:26). He guides us into truth and shows us things to come (see John 16:13). He testifies of the Son (see John 15:26; 1 Cor. 12:3). The Bible also teaches that the Savior and His servants will baptize with the Holy Ghost and with fire (see Matt. 3:11; Mark 1:8; John 1:33; Acts 1:5). I will speak of the meaning of that teaching later.

The Bible’s teachings about the Holy Ghost are reaffirmed and elaborated in the Book of Mormon and in modern revelations. The Holy Ghost is the means by which God inspires and reveals His will to His children (see, for example, D\&C 8:2–3). The Holy Ghost bears record of the Father and of the Son (see 3 Ne. 28:11; D\&C 20:27; 42:17). He enlightens our minds and fills us with joy (see D\&C 11:13). By the power of the Holy Ghost we may know the truth of all things (see Moro. 10:5). By His power we may have the mysteries of God unfolded to us (see 1 Ne. 10:19), all things which are expedient (see D\&C 18:18; 39:6). The Holy Ghost shows us what we should do (see 2 Ne. 32:5). We teach the gospel as we are directed by the Holy Ghost, which carries our words into the hearts of those we teach (see 2 Ne. 33:1).

Latter-day scriptures also teach that the remission of sins, which is made possible by the Atonement, comes “by baptism, and by fire, yea, even the Holy Ghost” (D\&C 19:31; see also 2 Ne. 31:17). Thus, the Risen Lord pleaded with the Nephites to repent and come unto Him and be baptized “that ye may be sanctified by the reception of the Holy Ghost, that ye may stand spotless before me at the last day” (3 Ne. 27:20).

The gift of the Holy Ghost is so important to our faith that a prophet gave it unique emphasis in a conversation with the president of the United States. Joseph Smith had journeyed to Washington to seek help in recovering compensation for injuries and losses the Saints had suffered in the Missouri persecutions. In his meeting with the president, Joseph was asked how this Church differed from the other religions of the day. The Prophet replied that “we differed in mode of baptism, and the gift of the Holy Ghost by the laying on of hands” (History of the Church, 4:42). He later explained that this answer was given because “all other considerations were contained in the gift of the Holy Ghost” (History of the Church, 4:42).

In highlighting the gift of the Holy Ghost as a distinguishing characteristic of our faith, we need to understand the important differences between (1) the Light of Christ, (2) a manifestation of the Holy Ghost, and (3) the gift of the Holy Ghost.

The Light of Christ, which is sometimes called the Spirit of Christ or the Spirit of God, “giveth light to every man that cometh into the world” (D\&C 84:46). This is the light “which is in all things, which giveth life to all things” (D\&C 88:13). The prophet Mormon taught that “the Spirit of Christ is given to every man, that he may know good from evil” (Moro. 7:16; see also Moro. 7:19; 2 Ne. 2:5; Hel. 14:31). Elder Lorenzo Snow spoke of this light when he said, “Everybody has the Spirit of God” (in Journal of Discourses, 14:304). The Light of Christ enlightens and gives understanding to all men (see D\&C 88:11).

In contrast, a manifestation of the Holy Ghost is more focused. This manifestation is given to acquaint sincere seekers with the truth about the Lord and His gospel. For example, the prophet Moroni promises that when we study the Book of Mormon and seek to know whether it is true, sincerely and with real intent, God will “manifest the truth of it unto [us], by the power of the Holy Ghost” (Moro. 10:4). Moroni also records this promise from the Risen Lord: “He that believeth these things which I have spoken, him will I visit with the manifestations of my Spirit, and he shall know and bear record. For because of my Spirit he shall know that these things are true” (Ether 4:11).

These manifestations are available to everyone. The Book of Mormon declares that the Savior “manifesteth himself unto all those who believe in him, by the power of the Holy Ghost; yea, unto every nation, kindred, tongue, and people” (2 Ne. 26:13).

To repeat, the Light of Christ is given to all men and women that they may know good from evil; manifestations of the Holy Ghost are given to lead sincere seekers to gospel truths that will persuade them to repentance and baptism.

The gift of the Holy Ghost is more comprehensive. The Prophet Joseph Smith explained: “There is a difference between the Holy Ghost and the gift of the Holy Ghost. Cornelius received the Holy Ghost before he was baptized, which was the convincing power of God unto him of the truth of the Gospel, but he could not receive the gift of the Holy Ghost until after he was baptized. Had he not taken this sign or ordinance upon him, the Holy Ghost which convinced him of the truth of God, would have left him” (Teachings of the Prophet Joseph Smith, sel. Joseph Fielding Smith [1976], 199; emphasis added).

The gift of the Holy Ghost includes the right to constant companionship, that we may “always have his Spirit to be with [us]” (D\&C 20:77).

A newly baptized member told me what she felt when she received that gift. This was a faithful Christian woman who had spent her life in service to others. She knew and loved the Lord, and she had felt the manifestations of His Spirit. When she received the added light of the restored gospel, she was baptized and the elders placed their hands upon her head and gave her the gift of the Holy Ghost. She recalled, “I felt the influence of the Holy Ghost settle upon me with greater intensity than I had ever felt before. He was like an old friend who had guided me in the past but now had come to stay.”

For faithful members of the Church of Jesus Christ, the companionship of the Holy Spirit should be so familiar that we must use care not to take it for granted. For example, that good feeling you have felt during the messages and music of this conference is a confirming witness of the Spirit, available to faithful members on a continuing basis. A member once asked me why he felt so good about the talks and music in a sacrament meeting, while a guest he had invited that day apparently experienced no such feeling. This is but one illustration of the contrast between one who has the gift of the Holy Ghost and is in tune with His promptings and one who has not, or is not.

If we are practicing our faith and seeking the companionship of the Holy Spirit, His presence can be felt in our hearts and in our homes. A family having daily family prayers and seeking to keep the commandments of God and honor His name and speak lovingly to one another will have a spiritual feeling in their home that will be discernible to all who enter it. I know this because I have felt the presence or absence of that feeling in many LDS homes.

It is important to remember that the illumination and revelation that come to an individual as a result of the gift of the Holy Ghost do not come suddenly or without seeking. President Spencer W. Kimball taught that the Holy Ghost “comes a little at a time as you merit it. And as your life is in harmony, you gradually receive the Holy Ghost in a great measure” (The Teachings of Spencer W. Kimball, ed. Edward L. Kimball [1982], 114).

The blessings available through the gift of the Holy Ghost are conditioned upon worthiness. “The Spirit of the Lord doth not dwell in unholy temples” (Hel. 4:24; see also Mosiah 2:36–37; 1 Cor. 3:16–17). Even though we have a right to His constant companionship, the Spirit of the Lord will dwell with us only when we keep the commandments. He will withdraw when we offend Him by profanity, uncleanliness, disobedience, rebellion, or other serious sins.

Worthy men and women who have the gift of the Holy Ghost can be edified and guided by inspiration and revelation. The Lord has declared that “the mysteries of his kingdom … are only to be seen and understood by the power of the Holy Spirit, which God bestows on those who love him, and purify themselves before him” (D\&C 76:114, 116).

A few years ago I met with a prospective mission president and his wife to discuss their availability for service. I asked whether their responsibilities to aged parents would preclude their service at that time. This sister was the only daughter of a wonderful mother, then about 80, whom she visited and helped each week. Though somewhat dependent physically, this mother was strong spiritually. She had served four missions and 15 years as a temple worker. Because she was in tune with the Spirit, she had a remarkable experience. Several months before this interview she told her daughter that the Spirit had whispered that her daughter’s husband would be called as a mission president. So advised, the mother had prepared herself for the needed separation and assured her daughter, long in advance of my assignment for the exploratory interview, that she would “not be a hindrance” to their service.

The need to keep our personal temple clean in order to have the companionship and guidance of the Holy Ghost explains the importance of the commandment to partake of the sacrament on the Sabbath.

In partaking of the sacrament, we can renew the effects of our baptism. When we desire a remission of our sins through the Atonement of our Savior, we are commanded to repent and come to Him with a broken heart and a contrite spirit (see 3 Ne. 9:20; 12:19; Moro. 6:2; D\&C 20:37). In the waters of baptism we witness to the Lord that we have repented of our sins and are willing to take His name upon us and serve Him to the end (see D\&C 20:37). The effects are described by Nephi: “For the gate by which ye should enter is repentance and baptism by water; and then cometh a remission of your sins by fire and by the Holy Ghost” (2 Ne. 31:17; see also Moro. 6:4). That last promise is fulfilled as a result of our receiving the gift of the Holy Ghost.

The renewal of our covenants by partaking of the sacrament should also be preceded by repentance, so we come to that sacred ordinance with a broken heart and a contrite spirit (see 2 Ne. 2:7; 3 Ne. 12:19; D\&C 59:8). Then, as we renew our baptismal covenants and affirm that we will “always remember him” (D\&C 20:77), the Lord will renew the promised remission of our sins, under the conditions and at the time He chooses. One of the primary purposes and effects of this renewal of covenants and cleansing from sin is “that [we] may always have his Spirit to be with [us]” (D\&C 20:77).

My brothers and sisters, I solemnly witness to you that these doctrines and principles are true. In view of these truths, I plead with all members of the Church, young and old, to attend sacrament meeting each Sabbath day and to partake of the sacrament with the repentant attitude described as “a broken heart and a contrite spirit” (3 Ne. 9:20). I pray that we will do so with the reverence and worship of our Savior that will signify a serious covenant to “always remember him” (D\&C 20:77). The Savior himself has said that we should partake “with an eye single to my glory—remembering unto the Father my body which was laid down for you, and my blood which was shed for the remission of your sins” (D\&C 27:2).

I pray that we will also partake of the sacrament with the submissive manner that will help us accept and serve in Church callings in order to comply with our solemn covenant to take His name and His work upon us. I also plead for us to comply with our solemn covenant to keep His commandments.

To those brothers and sisters who may have allowed themselves to become lax in this vital renewal of the covenants of the sacrament, I plead in words of the First Presidency that you “come back and feast at the table of the Lord, and taste again the sweet and satisfying fruits of fellowship with the saints” (“An Invitation to Come Back,” Church News, 22 Dec. 1985, 3). Let us qualify ourselves for our Savior’s promise that by partaking of the sacrament we will “be filled” (3 Ne. 20:8; see also 3 Ne. 18:9), which means that we will be “filled with the Spirit” (3 Ne. 20:9). That Spirit—the Holy Ghost—is our comforter, our direction finder, our communicator, our interpreter, our witness, and our purifier—our infallible guide and sanctifier for our mortal journey toward eternal life.

Any who may have thought it a small thing to partake of the sacrament should remember the Lord’s declaration that the foundation of a great work is laid by small things, for “out of small things proceedeth that which is great” (D\&C 64:33). Out of the seemingly small act of consciously and reverently renewing our baptismal covenants comes a renewal of the blessings of baptism by water and by the Spirit, that we may always have His Spirit to be with us. In this way all of us will be guided, and in this way all of us can be cleansed. That we may qualify for these precious blessings is my humble prayer, in the name of Jesus Christ, amen.

\newpage
\settalk{“Bishop, Help!”}
\setspeaker{Dallin H. Oaks}
\settalkdate{April 1997}
\talkheading{“Bishop, Help!”}{Dallin H. Oaks}

My brothers and sisters, I begin by sharing an event from a large ward in Provo about 20 years ago. During a sacrament meeting, a little boy made a big disturbance. After several minutes of trying to quiet this noisy three-year-old, the mother desperately handed him to the father, who was seated on the aisle close to the front of the chapel. By this time the noise distracted the speaker and audience, and everyone was very conscious of the parents’ plight. The father’s patience was much shorter than the mother’s. In a few moments he put the little boy over his shoulder, stood up, and started for the back door. Looking back over his father’s shoulder and sensing his determined steps, the little boy became quiet and apprehensive. Just as the father approached the rear door of the chapel, the little fellow reached his arms out toward the stand and shouted, “Bishop, help!”

There are times in the lives of all of us when we must reach out to our bishop or his counselors for help. Perhaps we need inspired counsel and direction to help with our families or our occupations. Perhaps we seek increased understanding of the gospel or the duties of our callings. We may need temporal assistance in a time of stress. We may even reach out for discipline to assist us in getting back on the path of growth. Always we benefit from their stalwart examples. Thank heaven for faithful and inspired bishops and branch presidents and their counselors!

A bishop (or branch president) has many duties. As the president of the Aaronic Priesthood, he personally oversees the programs and activities of the young men and young women in the ward. He and his counselors interview each one each year. They give special attention to teaching correct principles. Always they encourage our youth to prepare for the covenants they will make in the temple.

As the presiding high priest, the bishop gives direction to all quorums, auxiliaries, activities, and programs in the ward. Calls to ward positions are under his direction. So are home teaching and visiting teaching, and the performing of ordinances like baptism. Assisted in all of this by his counselors, he is responsible for sacrament meeting and for the teaching of the gospel in all classes in the ward. The bishopric also directs all of the other meetings of the ward, including the priesthood executive committee and the ward council.

The bishopric is also responsible for monitoring the Church-service time of all ward members serving under their direction. Knowing the circumstances of the ward, they determine the appropriate balance of ward meetings and activities and the time remaining for families. They are also conscious of the purpose of our Sunday consolidated meeting schedule, which was not established to give time for more Sabbath meetings but to allow increased time for families to be together and for individual gospel study and service.

The bishopric (or branch presidency) is also in charge of unit finances. They receive tithes and offerings, oversee the unit budget and expenditures, remit funds, and see that records are properly kept. The bishop is the judge who determines how Church commodities and funds are used to provide for the temporal needs of the members. He is also responsible for seeking out the poor and the needy.

The bishop is the judge and the shepherd who has the power of discernment and the right to revelation and inspiration for the guidance of the flock. He is responsible for holding worthiness interviews in order to authorize attendance at the temple, callings to ward positions, ordinations to priesthood offices, and the callings of missionaries. He administers formal and informal discipline for violation of the laws of the Church, and he counsels and helps members avoid the necessity for discipline.

Although some of their duties cannot be delegated, in most of these tasks the bishop and his counselors need the assistance of many others working under their direction: executive secretary, clerks, presidencies and group leadership of quorums, presidencies of auxiliaries, and officers and teachers. A bishop needs to be a skillful delegator, or he will be crushed under the burden of his responsibilities or frustrated at seeing so many of them unfilled.

I marvel at the work of our bishops and branch presidents. In my lifetime, our family has had many bishops. We have loved each of them and their counselors, and we have felt their love and assistance in our lives. Each of them was different in his personality, but each was a devoted servant of the Lord. I have seen the mantle of responsibility increase their stature, and I have rejoiced in their magnificent service to the people. God bless the bishops and bishoprics of this Church!

There is something else we should mention about bishops. They are not specialists. We do not have bishops whose sole attentions are directed toward the youth, the aged, the married, the abused, or any particular occupational or ethnic group. Under the revelations of the Lord and the directions of His prophets, a bishop is ordained and set apart to preside over a ward whose boundaries are geographic and whose membership includes all who reside there. For this reason, a bishop looks after the old and the young, the married and the single, the rich and the poor, the active and the less active. In this he seeks to unify the flock so that we may be taught and serve in groups of Saints that transcend considerations of age, marital status, ancestry, and economic condition. Our bishops lead us all in our efforts to follow the Savior’s commandment to “be one; and if ye are not one ye are not mine” (D\&C 38:27).

The Lord told the early members of His Church that the voice of His servants is the voice of the Lord, and that the hand of His servants is the hand of the Lord (see D\&C 1:38; 36:2). I testify to the truthfulness of that principle, which imposes a solemn duty upon the members of this Church to be loyal to their leaders and faithful in following their direction. I affirm that the Lord will bless us for doing so. That principle also imposes a great responsibility on the holders of office in this Church. Leaders must assure that they exercise their sacred authority “by persuasion, by long-suffering, by gentleness and meekness, and by love unfeigned” (D\&C 121:41).

We now have over 15,000 bishops and over 8,000 branch presidents in this Church. When we count their counselors, the total serving in bishoprics and branch presidencies is over 65,000. We praise and honor these worthy shepherds of the flock, judges in Israel, leaders and teachers of the people, men who love and are loved by those whom they serve as undershepherds of the Lord Jesus Christ. God bless these good men! And God bless their faithful wives, whose loyalty and support make their service possible.

As I began this talk, I quoted the words of a three-year-old who called, “Bishop, help!” I will now reverse those words and make them a challenge for each of us: “Help the bishop!”

Our current circumstances are different from those experienced by bishops and their counselors and members in earlier times. Today we have local leaders in most parts of the world. Many geographic wards and branches are in large cities and include hundreds of thousands or even millions of people. Some bishops travel during the week or commute long hours and great distances to work, effectively isolated from their families and their members for most of the hours of the week. Nevertheless, we also have communication and transportation resources undreamed of in earlier times. Whatever the physical changes over time, the nature of our local leaders’ callings has not changed, nor has their compensation. They are totally uncompensated by the coin of mortality. For the reward of their labors, all rely on the Lord’s deferred compensation plan.

Unchanged also is the fact that as they struggle with the heavy duties of their callings, bishops and their counselors must also earn a living and fulfill other family responsibilities. They do this not only because of their love for their wife and children, but also because they are responsible for being role models for the members of their flock. The burden is a heavy one that cannot be fulfilled without the supportive efforts of ward officers and members.

How do we help? To lighten the load of the bishopric, auxiliary presidencies and Melchizedek Priesthood quorum presidencies and group leaders need to exercise initiative and fully function in the great responsibilities of their callings. Bishops are responsible to call; they should not be required to beg or push. All of us should accept the callings we are given and serve in all diligence. The most common calling received for men is home teacher and for women is Relief Society visiting teacher. When properly performed, these vital callings can substantially lighten the load of the bishopric. Home teachers and visiting teachers are the eyes and ears and hands of the bishop. Brothers and sisters, help the bishop and his counselors by reliable, faithful performance of your visits and oversight as home teachers and visiting teachers.

Each of us should do all that we can, in the spirit of gospel self-reliance, to provide for ourselves and our families in a temporal and a spiritual way. Then, if it is necessary to reach out for help, we know we have first done all that we can. This includes helping the members of our immediate and extended families to the maximum extent possible so that the bishop is not faced with burdens that should be handled in the first instance by the individual or by the extended family.

Another way to help our busy bishops and their counselors is to be careful not to occupy their time with matters that others can handle. If we need an address or a phone number or help with some other routine task, we should not call a member of the bishopric. Let us reserve their time for the heavy responsibilities that are uniquely theirs. Let us call on others for the things others can handle.

When contacting our local leaders is necessary, we should remember that they have employment responsibilities too. Don’t contact them at their place of work unless there is a true emergency. Let us be careful not to put our leaders’ employment in jeopardy. Members should also be careful not to expect their local leaders to give them the products of those leaders’ occupations. Our leaders are called to give us Church service, not professional services or merchandise inventories.

We should remember that our leaders are also husbands and fathers. They are bishops or counselors for a season, but they will never be released from their family responsibilities, which are for eternity. Our leaders need time to perform their family responsibilities also, and our thoughtful consideration will help.

My heart ached for a young mother who wondered what would necessitate her bishop-husband’s spending six hours counseling a needy member on a Sunday following sacrament meeting. He did not arrive home until 6:00 p.m., which is bad enough, but this particular Sunday happened to be Christmas Day. I am sure the bishop felt he needed to give the help that was requested, but I also wonder whether a member in distress could not have held some of that need in abeyance long enough for a bishop to enjoy this Christmas afternoon with his family. That is admittedly an extreme example, but the problem is not an exceptional one, as many bishops and their wives would affirm.

A more familiar example was mentioned in a ward I recently attended in Salt Lake City. A wife of a member of the bishopric was speaking in sacrament meeting. She thanked the members of the ward for not phoning their home on Monday evening. She said that was the only time in the week when she and her children could plan to have their husband and father all to themselves. That forbearance would be good for all wards and branches.

Brothers and sisters, the offices of bishop and branch president and counselors are sacred in this Church. The men who hold those offices are respected by the Lord, inspired by His Spirit, and given the powers of discernment necessary to their office. We honor and love them, and we show this by our consideration for them.

I testify of the Lord Jesus Christ, whose Church this is and whose servants they are. I ask the blessings of the Lord on the members and leaders of this Church, general and local. In the name of Jesus Christ, amen.

\newpage
\settalk{Following the Pioneers}
\setspeaker{Dallin H. Oaks}
\settalkdate{October 1997}
\talkheading{Following the Pioneers}{Dallin H. Oaks}

A few years ago I showed one of my senior brethren a talk I had prepared for future delivery. He returned it with a stimulating two-word comment: “Therefore, what?” The talk was incomplete because it omitted a vital element: what a listener should do. I had failed to follow the example of King Benjamin, who concluded an important message by saying, “And now, if you believe all these things see that ye do them” (Mosiah 4:10).

For many months we have studied the lives and accomplishments of our pioneers, early and modern. We have thrilled to some modern reenactments, in which many have been blessed to participate. I was humbled to walk in the footsteps and wagon trails of my 31 pioneer ancestors for 13 miles over the Wyoming heights called Rocky Ridge, and for 5 miles on the trail 3 of them later followed down El Cajon Pass to settle what is now San Bernardino, California.

Now after all these studies and activities, it is appropriate to ask ourselves, “Therefore, what?” Are these pioneer celebrations academic, merely increasing our fund of experiences and knowledge? Or will they have a profound impact on how we live our lives?

This question applies to all of us. As President Hinckley reminded us last April, “Whether you are among the posterity of the pioneers or whether you were baptized only yesterday, each is the beneficiary of their great undertaking.” All of us enjoy the blessings of their efforts, and all of us have the responsibilities which go with that heritage.

It is not enough to study or reenact the accomplishments of our pioneers. We need to identify the great, eternal principles they applied to achieve all they achieved for our benefit and then apply those principles to the challenges of our day. In that way we honor their pioneering efforts, and we also reaffirm our heritage and strengthen its capacity to bless our own posterity and “those millions of our Heavenly Father’s children who have yet to hear and accept the gospel of Jesus Christ.” We are all pioneers in doing so.

Many of our challenges are different from those faced by former pioneers but perhaps just as dangerous and surely as significant to our own salvation and the salvation of those who follow us. For example, as for life-threatening obstacles, the wolves that prowled around pioneer settlements were no more dangerous to their children than the drug dealers or pornographers who threaten our children. Similarly, the early pioneers’ physical hunger posed no greater threat to their well-being than the spiritual hunger experienced by many in our day. The children of earlier pioneers were required to do incredibly hard physical work to survive their environment. That was no greater challenge than many of our young people now face from the absence of hard work, which results in spiritually corrosive challenges to discipline, responsibility, and self-worth. Jesus taught, “And fear not them which kill the body, but are not able to kill the soul: but rather fear him which is able to destroy both soul and body in hell” (Matt. 10:28).

The foremost quality of our pioneers was faith. With faith in God, they did what every pioneer does—they stepped forward into the unknown: a new religion, a new land, a new way of doing things. With faith in their leaders and in one another, they stood fast against formidable opposition. When their leader said, “This is the right place,” they trusted and they stayed. When other leaders said, “Do it this way,” they followed in faith.

Two companion qualities evident in the lives of our pioneers, early and modern, are unselfishness and sacrifice. Our Utah pioneers excelled at putting “the general welfare and community goals over individual gain and personal ambition.” That same quality is evident in the conversion stories of modern pioneers. Upon receiving a testimony of the truth of the restored gospel, they have unhesitatingly sacrificed all that was required to assure that its blessings will be available to their children and to generations unborn. Some have sold all their property to travel to a temple. Some have lost employment. Many have lost friends. Some have even lost parents and extended family, as new converts have been disowned for their faith. This must be the greatest sacrifice of all. Here we recall the Savior’s teaching:

“For I am come to set a man at variance against his father, and the daughter against her mother, and the daughter in law against her mother in law. …

“He that loveth father or mother more than me is not worthy of me: and he that loveth son or daughter more than me is not worthy of me.

“And he that taketh not his cross, and followeth after me, is not worthy of me” (Matt. 10:35, 37–38).

We praise what the pioneers’ unselfishness and sacrifice have done for us, but that is not enough. We should also assure that these same qualities are guiding principles for each of us as we have opportunities to sacrifice for our nations, our families, our quorums, our members, and our Church. This is especially important in societies that have exalted personal interest and individual rights to the point where these values seem to erase the principles of individual responsibility and sacrifice.

Other great qualities in our early pioneers were obedience, unity, and cooperation. We have all thrilled at the example of the Saints who responded to President Brigham Young’s call to rescue the stranded handcart companies, or to pull up roots in settled communities and apply their talents and lives to colonizing new areas.

Our people have always been characterized by their loyalty and obedience to the direction of their leaders, by their unity, and by their extraordinary capacity to cooperate in a common venture. We see the modern manifestations of these pioneer qualities in the great contributions our brothers and sisters make in a wide variety of private projects and common efforts that require unity and cooperation. Another modern manifestation of Mormon obedience, unity, and cooperation is our unique missionary program, from the preparation and service of young missionaries to the remarkably diverse activities of mature couples throughout the world.

Our recent Worldwide Pioneer Heritage Service Day, where Church members contributed more than two million hours of service to their local communities, provides visible evidence that the pioneer qualities of obedience, unity, and cooperation live on in our day. In this and the other examples given, I hope we are not too satisfied with an annual demonstration but will practice these pioneer principles all the days of our lives, as individuals, as families, as Church organizations, and as citizens.

In a day when our prophet has challenged us to reach out to welcome and fellowship new members and to reawaken the faith and fellowship of those who have strayed, we can gain strength from the example of the pioneers. The pioneer legacy is a legacy of inclusion. When the Saints were driven out of Missouri, many were so poor that they lacked teams and wagons to move. Their Church leaders were adamant that none of the poor would be left behind. The response was the same in the exodus from Nauvoo. At a conference of the Church in October 1845, the membership entered into a covenant to take all the Saints with them. Thereafter, in the initial epic struggle across Iowa, the companies that arrived first at their stopping place on the Missouri River sent rescue wagons back toward Nauvoo to gather those who had been too poor to leave earlier. The revelation that guided their next exodus on the trip west directed each company to “bear an equal proportion … in taking the poor, the widows, the fatherless, and the families of those who have gone into the army” (D\&C 136:8). When the wagons and handcarts moved west, their movement was always one of inclusion, and no day’s journey ended until every straggler was accounted for.

When the Saints settled in the valleys of the mountains, they promptly established a Perpetual Emigrating Fund to assist the poor to move from Winter Quarters, and later from the nations of Europe. At least half of those who journeyed to join the Saints could not have come without the help of leaders and members who were determined to include everyone who desired to gather to Zion. We need that same spirit of inclusion to accomplish our prophet’s clarion call for retention and reactivation.

Another great pioneer virtue was their commitment to one another, to their leaders, and to their faith. We honor that quality in the words of these favorite hymns:

\begin{verse}
Firm as the mountains around us,\\
Stalwart and brave we stand\\
On the rock our fathers planted\\
For us in this goodly land—\\
The rock of honor and virtue,\\
Of faith in the living God.\\
They raised his banner triumphant—\\
Over the desert sod.\\
And we hear the desert singing:\\
Carry on, carry on, carry on!
\end{verse}

[“Carry On,” Hymns, no. 255]

\begin{verse}
True to the faith that our parents have cherished,\\
True to the truth for which martyrs have perished,\\
To God’s command, Soul, heart, and hand,\\
Faithful and true we will ever stand.
\end{verse}

[“True to the Faith,” Hymns, no. 254]

What does it mean to be true to the faith? That word true implies commitment, integrity, endurance, and courage. It reminds us of the Book of Mormon’s description of the 2,000 young warriors:

“And they were all … exceedingly valiant for courage, and also for strength and activity; but behold, this was not all—they were men who were true at all times in whatsoever thing they were entrusted.

“Yea, they were men of truth and soberness, for they had been taught to keep the commandments of God and to walk uprightly before him” (Alma 53:20–21).

In the spirit of that description I say to our returned missionaries—men and women who have made covenants to serve the Lord and who have already served Him in the great work of proclaiming the gospel and perfecting the Saints—are you being true to the faith? Do you have the faith and continuing commitment to demonstrate the principles of the gospel in your own lives, consistently? You have served well, but do you, like the pioneers, have the courage and the consistency to be true to the faith and to endure to the end?

Here I recall a pioneer example of faith, commitment, and courage by some young men just about the age of our missionaries. A few months before the Prophet Joseph Smith was murdered at Carthage, some of his enemies plotted to kill him. As part of their plan, they sought to enlist others in their conspiracy. Among those they invited to a meeting in Nauvoo were two young men still in their teens, Robert Scott and Dennison L. Harris. Dennison’s father, Emer, was the older brother of Martin Harris, one of the Three Witnesses to the Book of Mormon. Being loyal to the Prophet, these young men immediately reported the invitation to Dennison’s father, who advised the Prophet Joseph and sought his advice. Joseph asked Emer Harris to request that the young men attend the meeting, pay strict attention to what was said, make no commitments, and report the entire matter to the Prophet.

As events proceeded, there were three meetings. They began by denouncing Joseph as a fallen prophet, proceeded to considering how Joseph could be overthrown, and concluded with specific planning to kill him. All of this the two young men reported to the Prophet Joseph after each meeting.

Before the third meeting, the Prophet foresaw what would happen and told the young men this would be the last meeting. He warned them that the conspirators might kill them when they refused the required oath to participate in the murderous scheme. He said he did not think the conspirators would shed their blood because they were so young, but he called upon their loyalty and courage in these words: “Don’t flinch. If you have to die, die like men, you will be martyrs to the cause, and your crowns can be no greater.” He renewed his original caution that they should not make any promises or enter into any covenants with the conspirators. Then he blessed them and expressed his love for their willingness to risk their lives for him.

As Joseph had foreseen, the third and final meeting required all present to unite in a solemn oath to destroy Joseph Smith. When the two boys refused, explaining that Joseph had never harmed them and they were unwilling to participate in his destruction, the leaders declared that since the boys knew the group’s plans, they must agree to join them or they must die on the spot. Knives were drawn.

Some protested killing the boys, especially since their parents knew of their presence, so their failure to return would cast suspicion on some of the conspirators. By the barest margin, the cautious course was chosen, and those who opposed killing prevailed. The boys were threatened with certain death if they ever revealed what had transpired in the meetings or who had participated, and they were then allowed to leave unharmed.

As the boys passed beyond the view of the guards, they were met by the Prophet, who was anxiously watching and praying for their safe return. They reported everything to him. He thanked and praised them, and then, for their safety, counseled them not to speak of this to anyone for 20 years or more.

The faith, commitment, and courage of these young men is an example to all of us. These pioneer qualities and the others I have mentioned—integrity, inclusion, cooperation, unity, unselfishness, sacrifice, and obedience—are as vital today as when they guided the actions of our pioneer forebears, early and modern. To honor those pioneers, we must honor and act upon the eternal principles that guided their actions. As President Hinckley reminded us last April, “We honor best those who have gone before when we serve well in the cause of truth.” That cause of truth is the cause of our Lord and Savior, Jesus Christ, whose servants they were, and whose servants we should strive to be. I testify of this and pray that we, too, may be “true to the faith that our parents have cherished,” in the name of Jesus Christ, amen.

\newpage
\settalk{Have You Been Saved?}
\setspeaker{Dallin H. Oaks}
\settalkdate{April 1998}
\talkheading{Have You Been Saved?}{Dallin H. Oaks}

What do we say when someone asks us, “Have you been saved?” This question, so common in the conversation of some Christians, can be puzzling to members of The Church of Jesus Christ of Latter-day Saints because it is not our usual way of speaking. We tend to speak of saved or salvation as a future event rather than something that has already been realized.

Good Christian people sometimes attach different meanings to some key gospel terms like saved or salvation. If we answer according to what our questioner probably means in asking if we have been “saved,” our answer must be “yes.” If we answer according to the various meanings we attach to the terms saved or salvation, our answer will be either “yes” or “yes, but with conditions.”

\intalkheader{I.}

As I understand what is meant by the good Christians who speak in these terms, we are “saved” when we sincerely declare or confess that we have accepted Jesus Christ as our personal Lord and Savior. This meaning relies on words the Apostle Paul taught the Christians of his day:

“If thou shalt confess with thy mouth the Lord Jesus, and shalt believe in thine heart that God hath raised him from the dead, thou shalt be saved.

“For with the heart man believeth unto righteousness; and with the mouth confession is made unto salvation” (Rom. 10:9–10).

To Latter-day Saints, the words saved and salvation in this teaching signify a present covenant relationship with Jesus Christ in which we are assured salvation from the consequences of sin if we are obedient. Every sincere Latter-day Saint is “saved” according to this meaning. We have been converted to the restored gospel of Jesus Christ, we have experienced repentance and baptism, and we are renewing our covenants of baptism by partaking of the sacrament.

\intalkheader{II.}

As Latter-day Saints use the words saved and salvation, there are at least six different meanings. According to some of these, our salvation is assured—we are already saved. In others, salvation must be spoken of as a future event (e.g., 1 Cor. 5:5) or as conditioned upon a future event (e.g., Mark 13:13). But in all of these meanings, or kinds of salvation, salvation is in and through Jesus Christ.

First, all mortals have been saved from the permanence of death through the Resurrection of Jesus Christ. “For as in Adam all die, even so in Christ shall all be made alive” (1 Cor. 15:22).

As to salvation from sin and the consequences of sin, our answer to the question of whether or not we have been saved is “yes, but with conditions.” Our third article of faith declares our belief:

“We believe that through the Atonement of Christ, all mankind may be saved, by obedience to the laws and ordinances of the Gospel” (A of F 1:3).

Many Bible verses declare that Jesus came to take away the sins of the world (e.g., John 1:29; Matt. 26:28). The New Testament frequently refers to the grace of God and to salvation by grace (e.g., John 1:17; Acts 15:11; Eph. 2:8). But it also has many specific commandments on personal behavior and many references to the importance of works (e.g., Matt. 5:16; Eph. 2:10; James 2:14–17). In addition, the Savior taught that we must endure to the end in order to be saved (see Matt. 10:22; Mark 13:13).

Relying upon the totality of Bible teachings and upon clarifications received through modern revelation, we testify that being cleansed from sin through Christ’s Atonement is conditioned upon the individual sinner’s faith, which must be manifested by obedience to the Lord’s command to repent, be baptized, and receive the Holy Ghost (see Acts 2:37–38). “Verily, verily, I say unto thee,” Jesus taught, “except a man be born of water and of the Spirit, he cannot enter into the kingdom of God” (John 3:5; see also Mark 16:16; Acts 2:37–38). Believers who have had this required rebirth at the hands of those having authority have already been saved from sin conditionally, but they will not be saved finally until they have completed their mortal probation with the required continuing repentance, faithfulness, service, and enduring to the end.

Some Christians accuse Latter-day Saints who give this answer of denying the grace of God through claiming they can earn their own salvation. We answer this accusation with the words of two Book of Mormon prophets. Nephi taught, “For we labor diligently … to persuade our children … to believe in Christ, and to be reconciled to God; for we know that it is by grace that we are saved, after all we can do” (2 Ne. 25:23). And what is “all we can do”? It surely includes repentance (see Alma 24:11) and baptism, keeping the commandments, and enduring to the end. Moroni pleaded, “Yea, come unto Christ, and be perfected in him, and deny yourselves of all ungodliness; and if ye shall deny yourselves of all ungodliness, and love God with all your might, mind and strength, then is his grace sufficient for you, that by his grace ye may be perfect in Christ” (Moro. 10:32).

We are not saved in our sins, as by being unconditionally saved through confessing Christ and then, inevitably, committing sins in our remaining lives (see Alma 11:36–37). We are saved from our sins (see Hel. 5:10) by a weekly renewal of our repentance and cleansing through the grace of God and His blessed plan of salvation (see 3 Ne. 9:20–22).

The question of whether a person has been saved is sometimes phrased in terms of whether that person has been “born again.” Being “born again” is a familiar reference in the Bible and the Book of Mormon. As noted earlier, Jesus taught that except a man was “born again” (John 3:3), of water and of the Spirit, he could not enter into the kingdom of God (see John 3:5). The Book of Mormon has many teachings about the necessity of being “born again” or “born of God” (Mosiah 27:25; see vv. 24–26; Alma 36:24, 26; Moses 6:59). As we understand these scriptures, our answer to whether we have been born again is clearly “yes.” We were born again when we entered into a covenant relationship with our Savior by being born of water and of the Spirit and by taking upon us the name of Jesus Christ. We can renew that rebirth each Sabbath when we partake of the sacrament.

Latter-day Saints affirm that those who have been born again in this way are spiritually begotten sons and daughters of Jesus Christ (see Mosiah 5:7; 15:9–13; 27:25). Nevertheless, in order to realize the intended blessings of this born-again status, we must still keep our covenants and endure to the end. In the meantime, through the grace of God, we have been born again as new creatures with new spiritual parentage and the prospects of a glorious inheritance.

A fourth meaning of being saved is to be saved from the darkness of ignorance of God the Father and His Son, Jesus Christ, and of the purpose of life, and of the destiny of men and women. The gospel made known to us by the teachings of Jesus Christ has given us this salvation. “I am the light of the world,” Jesus taught; “he that followeth me shall not walk in darkness, but shall have the light of life” (John 8:12; see also John 12:46).

For Latter-day Saints, being saved can also mean being saved or delivered from the second death (meaning the final spiritual death) by assurance of a kingdom of glory in the world to come (see 1 Cor. 15:40–42). Just as the Resurrection is universal, we affirm that every person who ever lived upon the face of the earth—except for a very few—is assured of salvation in this sense. As we read in modern revelation:

“And this is the gospel, the glad tidings …

“That he came into the world, even Jesus, to be crucified for the world, and to bear the sins of the world, and to sanctify the world, and to cleanse it from all unrighteousness;

“That through him all might be saved whom the Father had put into his power and made by him;

“Who glorifies the Father, and saves all the works of his hands, except those sons of perdition who deny the Son after the Father has revealed him” (D\&C 76:40–43; emphasis added).

The prophet Brigham Young taught that doctrine when he declared that “every person who does not sin away the day of grace, and become an angel to the Devil, will be brought forth to inherit a kingdom of glory” (Discourses of Brigham Young [1954], 382). This meaning of saved ennobles the whole human race through the grace of our Lord and Savior, Jesus Christ. In this sense of the word, all should answer: “Yes, I have been saved. Glory to God for the gospel and gift and grace of His Son!”

Finally, in another usage familiar and unique to Latter-day Saints, the words saved and salvation are also used to denote exaltation or eternal life (see Abr. 2:11). This is sometimes referred to as the “fulness of salvation” (Bruce R. McConkie, The Mortal Messiah, 4 vols. [1979–81], 1:242). This salvation requires more than repentance and baptism by appropriate priesthood authority. It also requires the making of sacred covenants, including eternal marriage, in the temples of God and faithfulness to those covenants by enduring to the end. If we use the word salvation to mean “exaltation,” it is premature for any of us to say that we have been “saved” in mortality. That glorious status can only follow the final judgment of Him who is the Great Judge of the living and the dead.

I have suggested that the short answer to the question of whether a faithful member of The Church of Jesus Christ of Latter-day Saints has been saved or born again must be a fervent “yes.” Our covenant relationship with our Savior puts us in that “saved” or “born again” condition meant by those who ask this question. Some modern prophets have also used “salvation” or “saved” in that same present sense. President Brigham Young declared:

“It is present salvation and the present influence of the Holy Ghost that we need every day to keep us on saving ground. …

“I want present salvation. … Life is for us, and it is for us to receive it today, and not wait for the Millennium. Let us take a course to be saved today” (Discourses of Brigham Young, sel. John A. Widtsoe [1954], 15–16). President David O. McKay spoke of the revealed gospel of Jesus Christ in that same present sense of “salvation here—here and now” (Gospel Ideals [1953], 6).

\intalkheader{III.}

I will conclude by discussing another important question members and leaders of The Church of Jesus Christ of Latter-day Saints are asked by others: “Why do you send missionaries to preach to other Christians?” Sometimes this is asked with curiosity and sometimes with resentment.

My most memorable experience with that question occurred some years ago in what we then called the Eastern Bloc. After many years of Communist hostility to religion, these countries were suddenly and miraculously given a measure of religious freedom. When that door opened, many Christian faiths sent missionaries. As part of our preparation to do so, the First Presidency sent members of the Quorum of the Twelve Apostles to meet with government and church leaders in these countries. Our assignment was to introduce ourselves and to explain what our missionaries would be doing.

Elder Russell M. Nelson and I called on the leader of the Orthodox Church in one of these countries. Here was a man who had helped keep the light of Christianity burning through the dark decades of Communist repression. I noted in my journal that he was a warm and gracious man who impressed me as a servant of the Lord. I mention this so that you will not think there was any spirit of arrogance or contention in our conversation of nearly an hour. Our visit was pleasant and cordial, filled with the goodwill that should always characterize conversations between men and women who love the Lord and seek to serve Him, each according to his or her own understanding.

Our host told us about the activities of his church during the period of Communist repression. He described the various difficulties his church and its work were experiencing as they emerged from that period and sought to regain their former position in the life of the country and the hearts of the people. We introduced ourselves and our fundamental beliefs. We explained that we would soon be sending missionaries into his country and told him how they would perform their labors.

He asked, “Will your missionaries preach only to unbelievers, or will they also try to preach to believers?” We replied that our message was for everyone, believers as well as unbelievers. We gave two reasons for this answer—one a matter of principle and the other a matter of practicality. We told him that we preached to believers as well as unbelievers because our message, the restored gospel, makes an important addition to the knowledge, happiness, and peace of all mankind. As a matter of practicality, we preach to believers as well as unbelievers because we cannot tell the difference. I remember asking this distinguished leader, “When you stand before a congregation and look into the faces of the people, can you tell the difference between those who are real believers and those who are not?” He smiled wryly, and I sensed an admission that he had understood the point.

Through missionaries and members, the message of the restored gospel is going to all the world. To non-Christians, we witness of Christ and share the truths and ordinances of His restored gospel. To Christians we do the same. Even if a Christian has been “saved” in the familiar single sense discussed earlier, we teach that there remains more to be learned and more to be experienced. As President Hinckley recently said: “[We are] not argumentative. We do not debate. We, in effect, simply say to others, ‘Bring all the good that you have and let us see if we can add to it’” (“The BYU Experience,” BYU devotional address, 4 Nov. 1997).

The Church of Jesus Christ of Latter-day Saints offers all of the children of God the opportunity to learn the fulness of the gospel of Jesus Christ as restored in these latter days. We offer everyone the privilege of receiving all of the ordinances of salvation and exaltation.

We invite all to hear this message, and we invite all who receive the confirming witness of the Spirit to heed it. These things are true, I testify in the name of Jesus Christ, amen.

\newpage
\settalk{The Aaronic Priesthood and the Sacrament}
\setspeaker{Dallin H. Oaks}
\settalkdate{October 1998}
\talkheading{The Aaronic Priesthood and the Sacrament}{Dallin H. Oaks}

My beloved brethren, I appreciate the opportunity to speak to you this evening. I address my remarks to the young men who hold the Aaronic Priesthood and to the bishops and counselors who preside over them. I will speak about the sacred activities of Aaronic Priesthood holders in preparing, administering, and passing the sacrament of the Lord’s Supper to the members of the Church.

\intalkheader{I.}

On May 15, 1829, John the Baptist restored the Aaronic Priesthood to the earth. He did so by laying his hands upon Joseph Smith and Oliver Cowdery and speaking these words: “Upon you my fellow servants, in the name of Messiah I confer the Priesthood of Aaron, which holds the keys of the ministering of angels, and of the gospel of repentance, and of baptism by immersion for the remission of sins; and this shall never be taken again from the earth, until the sons of Levi do offer again an offering unto the Lord in righteousness” (D\&C 13:1).

Later the Lord revealed these further truths: “The lesser priesthood … holdeth the key of the ministering of angels and the preparatory gospel;

“Which gospel is the gospel of repentance and of baptism, and the remission of sins” (D\&C 84:26–27).

What does it mean that the Aaronic Priesthood holds “the key of the ministering of angels” and of the “gospel of repentance and of baptism, and the remission of sins”? The meaning is found in the ordinance of baptism and in the sacrament. Baptism is for the remission of sins, and the sacrament is a renewal of the covenants and blessings of baptism. Both should be preceded by repentance. When we keep the covenants made in these ordinances, we are promised that we will always have His Spirit to be with us. The ministering of angels is one of the manifestations of that Spirit.

\intalkheader{II.}

We begin with the doctrine as taught by the Lord. During His ministry, Jesus taught that baptism is necessary for salvation. “Except a man be born of water and of the Spirit, he cannot enter into the kingdom of God” (John 3:5). Baptism is the first of the saving ordinances. When we are baptized, we covenant that we will take upon us the name of Jesus Christ and serve Him and keep His commandments.

At the conclusion of His ministry, Jesus introduced the sacrament of the Lord’s Supper. He broke bread and blessed it and gave it to His disciples, saying, “Take, eat; this is my body” (Matt. 26:26). “This do in remembrance of me” (Luke 22:19). He took the cup and gave thanks and gave it to them, saying, “This is my blood of the new testament, which is shed for many for the remission of sins” (Matt. 26:28).

When He introduced the sacrament, the Savior also gave teachings and promises about the Holy Ghost. On that sacred occasion known as the Last Supper, Jesus explained the mission of the Comforter, which is the Holy Ghost. The Comforter would testify of Him and reveal other truths. Jesus also explained that He had to leave His disciples in order for the Comforter to come to them. When I depart, He told them, “I will send him unto you” (John 16:7). After His Resurrection, He told His Apostles to tarry in Jerusalem until they were given “power from on high” (Luke 24:49). That power came when “the promise of the Holy Ghost” was “shed forth” upon the Apostles on the day of Pentecost (Acts 2:33).

Similarly, when the Savior introduced the sacrament in the New World, He promised, “He that eateth this bread eateth of my body to his soul; and he that drinketh of this wine drinketh of my blood to his soul; and his soul shall never hunger nor thirst, but shall be filled” (3 Ne. 20:8). The meaning of that promise is evident: “Now, when the multitude had all eaten and drunk, behold, they were filled with the Spirit” (3 Ne. 20:9).

The close relationship between partaking of the sacrament and the companionship of the Holy Ghost is explained in the revealed prayer on the sacrament. In partaking of the bread, we witness that we are willing to take upon us the name of Jesus Christ and always remember Him and keep His commandments. When we do so, we have the promise that we will always have His Spirit to be with us (see D\&C 20:77).

To have the continuous companionship of the Holy Ghost is the most precious possession we can have in mortality. The gift of the Holy Ghost was conferred upon us by the authority of the Melchizedek Priesthood after our baptism. But to realize the blessings of that gift, we must keep ourselves free from sin. When we commit sin, we become unclean and the Spirit of the Lord withdraws from us. The Spirit of the Lord does not dwell in “unholy temples” (see Mosiah 2:36–37; Alma 34:35–36; Hel. 4:24), and no unclean thing can dwell in His presence (see Eph. 5:5; 1 Ne. 10:21; Alma 7:21; Moses 6:57).

A few weeks ago I used a chain saw to cut down a tree in my backyard. It was a dirty job, and when I was done I was splattered with a filthy mixture of sawdust and oil. In that condition I did not want anyone to see me. I just wanted to be cleansed in water so I would again feel comfortable in the presence of other people.

Not one of you young men and not one of your leaders has lived without sin since his baptism. Without some provision for further cleansing after our baptism, each of us is lost to things spiritual. We cannot have the companionship of the Holy Ghost, and at the final judgment we would be bound to be “cast off forever” (1 Ne. 10:21). How grateful we are that the Lord has provided a process for each baptized member of His Church to be periodically cleansed from the soil of sin. The sacrament is an essential part of that process.

We are commanded to repent of our sins and to come to the Lord with a broken heart and a contrite spirit and partake of the sacrament in compliance with its covenants. When we renew our baptismal covenants in this way, the Lord renews the cleansing effect of our baptism. In this way we are made clean and can always have His Spirit to be with us. The importance of this is evident in the Lord’s commandment that we partake of the sacrament each week (see D\&C 59:8–9).

We cannot overstate the importance of the Aaronic Priesthood in this. All of these vital steps pertaining to the remission of sins are performed through the saving ordinance of baptism and the renewing ordinance of the sacrament. Both of these ordinances are officiated by holders of the Aaronic Priesthood under the direction of the bishopric, who exercise the keys of the gospel of repentance and of baptism and the remission of sins.

\intalkheader{III.}

In a closely related way, these ordinances of the Aaronic Priesthood are also vital to the ministering of angels.

“The word ‘angel’ is used in the scriptures for any heavenly being bearing God’s message” (George Q. Cannon, Gospel Truth, sel. Jerreld L. Newquist [1987], 54). The scriptures recite numerous instances where an angel appeared personally. Angelic appearances to Zacharias and Mary (see Luke 1) and to King Benjamin and Nephi (see Mosiah 3:2; 3 Ne. 7:17–18) are only a few examples. When I was young, I thought such personal appearances were the only meaning of the ministering of angels. As a young holder of the Aaronic Priesthood, I did not think I would see an angel, and I wondered what such appearances had to do with the Aaronic Priesthood.

But the ministering of angels can also be unseen. Angelic messages can be delivered by a voice or merely by thoughts or feelings communicated to the mind. President John Taylor described “the action of the angels, or messengers of God, upon our minds, so that the heart can conceive … revelations from the eternal world” (Gospel Kingdom, sel. G. Homer Durham [1943], 31).

Nephi described three manifestations of the ministering of angels when he reminded his rebellious brothers that (1) they had “seen an angel,” (2) they had “heard his voice from time to time,” and (3) also that an angel had “spoken unto [them] in a still small voice” though they were “past feeling” and “could not feel his words” (1 Ne. 17:45). The scriptures contain many other statements that angels are sent to teach the gospel and bring men to Christ (see Heb. 1:14; Alma 39:19; Moro. 7:25, 29, 31–32; D\&C 20:35). Most angelic communications are felt or heard rather than seen.

How does the Aaronic Priesthood hold the key to the ministering of angels? The answer is the same as for the Spirit of the Lord.

In general, the blessings of spiritual companionship and communication are only available to those who are clean. As explained earlier, through the Aaronic Priesthood ordinances of baptism and the sacrament, we are cleansed of our sins and promised that if we keep our covenants we will always have His Spirit to be with us. I believe that promise not only refers to the Holy Ghost but also to the ministering of angels, for “angels speak by the power of the Holy Ghost; wherefore, they speak the words of Christ” (2 Ne. 32:3). So it is that those who hold the Aaronic Priesthood open the door for all Church members who worthily partake of the sacrament to enjoy the companionship of the Spirit of the Lord and the ministering of angels.

\intalkheader{IV.}

The doctrines I have just discussed are contained in the scriptures. From the scriptures we also know that those who officiate in the priesthood act in behalf of the Lord (see D\&C 1:38; 36:2). I will now suggest how teachers and priests and deacons should carry out their sacred responsibilities to act in behalf of the Lord in preparing, administering, and passing the sacrament. I will not suggest detailed rules, since the circumstances in various wards and branches in our worldwide Church are so different that a specific rule that seems required in one setting may be inappropriate in another. Rather, I will suggest a principle based on the doctrines. If all understand this principle and act in harmony with it, there should be little need for rules. If rules or counseling are needed in individual cases, local leaders can provide them, consistent with the doctrines and their related principles.

The principle I suggest to govern those officiating in the sacrament—whether preparing, administering, or passing—is that they should not do anything that would distract any member from his or her worship and renewal of covenants. This principle of non-distraction suggests some companion principles.

Deacons, teachers, and priests should always be clean in appearance and reverent in the manner in which they perform their solemn and sacred responsibilities. Teachers’ special assignments in preparing the sacrament are the least visible but should still be done with dignity, quietly and reverently. Teachers should always remember that the emblems they are preparing represent the body and blood of our Lord.

To avoid distracting from the sacred occasion, priests should speak the sacrament prayers clearly and distinctly. Prayers that are rattled off swiftly or mumbled inaudibly will not do. All present should be helped to understand an ordinance and covenants so important that the Lord prescribed the exact words to be uttered. All should be helped to focus on those sacred words as they renew their covenants by partaking.

On this subject I feel to share a painful experience from my youth. As a 16-year-old priest, I was just beginning a part-time job as a radio announcer at a local station. After I offered a prayer at the sacrament table in our ward, a girl who was present told me I sounded like I was reading a commercial. Can you imagine the shame I felt? After 50 years that rebuke still stings. Brethren, remember the significance of those sacred prayers. You are praying as a servant of the Lord in behalf of the entire congregation. Speak to be heard and understood, and say it like you mean it.

Deacons should pass the sacrament in a reverent and orderly manner, with no needless motions or expressions that call attention to themselves. In all their actions they should avoid distracting any member of the congregation from worship and covenant making.

All who officiate in the sacrament—in preparing, administering, or passing—should be well groomed and modestly dressed, with nothing about their personal appearance that calls special attention to themselves. In appearance as well as actions, they should avoid distracting anyone present from full attention to the worship and covenant making that are the purpose of this sacred ordinance.

This principle of non-distraction applies to things unseen as well as seen. If someone officiating in this sacred ordinance is unworthy to participate, and this is known to anyone present, their participation is a serious distraction to that person. Young men, if any of you is unworthy, talk to your bishop without delay. Obtain his direction on what you should do to qualify yourself to participate in your priesthood duties worthily and appropriately.

I have a final suggestion. With the single exception of those priests occupied breaking the bread, all who hold the Aaronic Priesthood should join in singing the sacrament hymn by which we worship and prepare to partake. No one needs that spiritual preparation more than the priesthood holders who will officiate in it. My young brethren, it is important that you sing the sacrament hymn. Please do so.

The Aaronic Priesthood holds the keys of “the gospel of repentance and of baptism, and the remission of sins” (D\&C 84:27). The cleansing power of our Savior’s Atonement is renewed for us as we partake of the sacrament. The promise that we “may always have his Spirit to be with [us]” (D\&C 20:77) is essential to our spirituality. The ordinances of the Aaronic Priesthood are vital to all of this. I testify that this is true, and I pray that our brethren of the Aaronic Priesthood will understand the importance of their sacred responsibilities and act worthily in them, in the name of Jesus Christ, amen.

\newpage
\settalk{The Witness: Martin Harris}
\setspeaker{Dallin H. Oaks}
\settalkdate{April 1999}
\talkheading{The Witness: Martin Harris}{Dallin H. Oaks}

\intalkheader{The Law of Witnesses}

Witnesses and witnessing are vital in God’s plan for the salvation of His children. In the Godhead, the function of the Holy Ghost is to bear witness of the Father and the Son (see 2 Ne. 31:18). The Father has borne witness of the Son (see Matt. 3:17; 17:5; John 5:31–39), and the Son has borne witness of the Father (see John 17). The Lord has commanded His servants to testify of Him (see Isa. 43:10; Mosiah 18:9; D\&C 84:62), and all of the prophets have borne witness of Jesus Christ (see Acts 10:42–43; Rev. 19:10).

The scriptures state that “in the mouth of two or three witnesses shall every word be established” (2 Cor. 13:1; D\&C 6:28; see also Deut. 19:15). The most important ordinances of salvation—baptism, marriage, and other ordinances of the temple—are required to have witnesses (see D\&C 127:6; 128:3).

The Bible witnesses of Jesus Christ by prophecies of His coming, by accounts of His ministry, and by the testimonies of those who carried His message to the world. The Book of Mormon has the same content: witnesses preceding, during, and following the ministry of the Messiah. Appropriately, it is now subtitled “Another Testament of Jesus Christ.”

\intalkheader{Book of Mormon Witnesses}

There are witnesses of the Book of Mormon itself. I have chosen to speak about the significance of their testimonies and about the life of one of them.

While Joseph Smith was translating the Book of Mormon, the Lord revealed that in addition to the Prophet’s testimony, the world would have “the testimony of three of my servants, whom I shall call and ordain, unto whom I will show these things” (D\&C 5:11; see also 2 Ne. 27:12–13; Ether 5:2–4). “They shall know of a surety that these things are true,” the Lord declared, “for from heaven will I declare it unto them” (D\&C 5:12).

There were also eight witnesses, but their testimony is a subject for another time.

The three men chosen as witnesses of the Book of Mormon were Oliver Cowdery, David Whitmer, and Martin Harris. Their written “Testimony of Three Witnesses” has been included in all of the almost 100 million copies of the Book of Mormon the Church has published since 1830. These witnesses solemnly testify that they “have seen the plates which contain this record” and “the engravings which are upon the plates.” They witness that these writings “have been translated by the gift and power of God, for his voice hath declared it unto us.” They testify, “We declare with words of soberness, that an angel of God came down from heaven, and he brought and laid before our eyes, that we beheld and saw the plates, and the engravings thereon; and we know that it is by the grace of God the Father, and our Lord Jesus Christ, that we beheld and bear record that these things are true.”

Further, “the voice of the Lord commanded us that we should bear record of it; wherefore, to be obedient unto the commandments of God, we bear testimony of these things” (“The Testimony of Three Witnesses,” Book of Mormon).

People who deny the possibility of supernatural beings may reject this remarkable testimony, but people who are open to believe in miraculous experiences should find it compelling. The solemn written testimony of three witnesses to what they saw and heard—two of them simultaneously and the third almost immediately thereafter—is entitled to great weight. Indeed, we know that upon the testimony of one witness great miracles have been claimed and accepted by many religious people, and in the secular world the testimony of one witness has been deemed sufficient for weighty penalties and judgments.

Persons experienced in evaluating testimony commonly consider a witness’s opportunity to observe an event and the possibility of his bias on the subject. Where different witnesses give identical testimony about the same event, skeptics look for evidence of collusion among them or for other witnesses who could contradict them.

Measured against all of these possible objections, the testimony of the Three Witnesses to the Book of Mormon stands forth in great strength. Each of the three had ample reason and opportunity to renounce his testimony if it had been false or to equivocate on details if any had been inaccurate. As is well known, because of disagreements or jealousies involving other leaders of the Church, each one of these three witnesses was excommunicated from The Church of Jesus Christ of Latter-day Saints by about eight years after the publication of their testimony. All three went their separate ways, with no common interest to support a collusive effort. Yet to the end of their lives—periods ranging from 12 to 50 years after their excommunications—not one of these witnesses deviated from his published testimony or said anything that cast any shadow on its truthfulness.

Furthermore, their testimony stands uncontradicted by any other witnesses. Reject it one may, but how does one explain three men of good character uniting and persisting in this published testimony to the end of their lives in the face of great ridicule and other personal disadvantage? Like the Book of Mormon itself, there is no better explanation than is given in the testimony itself, the solemn statement of good and honest men who told what they saw.

\intalkheader{Martin Harris}

Having a special interest in Martin Harris, I have been saddened at how he is remembered by most Church members. He deserves better than to be remembered solely as the man who unrighteously obtained and then lost the initial manuscript pages of the Book of Mormon.

When the Book of Mormon was published, Martin Harris was nearly 47 years of age, more than 20 years older than Joseph Smith and the other two witnesses. He was a prosperous and respected citizen of Palmyra, New York. He owned a farm of over 240 acres, large for the time and place. He was an honored veteran of two battles in the War of 1812. His fellow citizens entrusted him with many elective offices and responsibilities in the community. He was universally respected for his industry and integrity. Assessments by contemporaries described him as “an industrious, hard-working farmer, shrewd in his business calculations, frugal in his habits,” and “strictly upright in his business dealings” (quoted in Richard Lloyd Anderson, Investigating the Book of Mormon Witnesses [1981], 96–97, 98).

This prosperous and upright older man befriended the young and penniless Joseph Smith, giving him the \$50 that permitted him to pay his debts in Palmyra and locate in northeastern Pennsylvania, about 150 miles away. There, in April 1828, Joseph Smith began his first persistent translation of the Book of Mormon. He dictated, and Martin Harris wrote until there were 116 pages of manuscript.

Martin’s persistent requests to show this manuscript to his family wearied Joseph into letting him take it to Palmyra, where its pages were stolen from him, lost, and probably burned. For this the Lord rebuked Martin and Joseph. Joseph had his gift of translation suspended for a season, and Martin was rebuked as “a wicked man” who had “set at naught the counsels of God, and … broken the most sacred promises which were made before God” (D\&C 3:12–13; see also D\&C 10). Fortunately, both Joseph and Martin were later forgiven by the Lord, and the work of translation resumed with other scribes. We obviously honor Joseph for his magnificent ministry, but Martin’s subsequent faithfulness continues under a shadow from which this important man should be rescued.

I will review some of the high points of Martin Harris’s life following the devastating episode of the stolen and lost manuscript.

About nine months after Martin’s rebuke, the Prophet Joseph received a revelation declaring that there would be three witnesses to the plates and if Martin would humble himself he would be privileged to see them (see D\&C 5:11, 15, 24). A few months later, Martin Harris was selected as one of the Three Witnesses and had the experience and bore the testimony described earlier.

One of Martin Harris’s greatest contributions to the Church, for which he should be honored for all time, was his financing the publication of the Book of Mormon. In August 1829 he mortgaged his home and farm to Egbert B. Grandin to secure payment on the printer’s contract. Seven months later the 5,000 copies of the first printing of the Book of Mormon were completed. Later, when the mortgage note fell due, the home and a portion of the farm were sold for \$3,000. In this way Martin Harris was obedient to the Lord’s revelation:

“Thou shalt not covet thine own property, but impart it freely to the printing of the Book of Mormon. …

“Pay the debt thou hast contracted with the printer. Release thyself from bondage” (D\&C 19:26, 35).

Other records and revelations show Martin Harris’s significant involvement in the activities of the restored Church and his standing with God. He was present at the organization of the Church on April 6, 1830, and was baptized that same day. A year later he was called to journey to Missouri with Joseph Smith, Sidney Rigdon, and Edward Partridge (see D\&C 52:24). In Missouri that year—1831—he was commanded to “be an example unto the church, in laying his moneys before the bishop of the church” (D\&C 58:35), thus becoming the first man the Lord called by name to consecrate his property in Zion. Two months later he was named with Joseph Smith, Oliver Cowdery, Sidney Rigdon, and others to be “stewards over the revelations and commandments” (D\&C 70:3; see also D\&C 70:1), a direction to publish and circulate what later became the Doctrine and Covenants.

In 1832 Martin Harris’s older brother, Emer, who is my great-great-grandfather, was called on a mission from Ohio (see D\&C 75:30). Emer spent a year preaching the gospel near his former home in northeastern Pennsylvania. During most of this time Emer’s companion was his brother Martin, whose zeal in preaching even caused him to be jailed for a few days. The Harris brothers baptized about 100 persons. Among those baptized was a family named Oaks, which included my great-great-grandfather. Thus, my middle name and my last name come from the grandfathers who met in that missionary encounter in Susquehanna County in 1832–33.

Back in Kirtland, Ohio, after his mission, in February 1834 Martin Harris was chosen by revelation to serve on the first high council in the Church (see D\&C 102:3). Less than three months later, he left Kirtland with the men of Zion’s Camp, marching 900 miles to Missouri to relieve the oppressed Saints there.

One of the most important events of the Restoration was the calling of a Quorum of Twelve Apostles in February 1835. The Three Witnesses, including Martin Harris, were appointed to “search out the Twelve” (D\&C 18:37), to select them and, under authority granted by the Prophet and his counselors, to ordain them [these ordinations were then confirmed under the hands of the First Presidency] (see B. H. Roberts, A Comprehensive History of the Church, 1:372–75).

From a position of great influence and authority, all three witnesses fell, each in his own way. During 1837 there were intense financial and spiritual conflicts in Kirtland, Ohio. Martin Harris later said that he “lost confidence in Joseph Smith” and “his mind became darkened” (quoted in Investigating the Book of Mormon Witnesses, 110). He was released from the high council in September 1837 and three months later was excommunicated.

Martin’s wife, Lucy, who had been involved in the loss of the manuscript pages, died in Palmyra in 1836. Within a year thereafter, Martin and his family located in Kirtland, and Martin married Caroline Young, a niece of Brigham Young.

When most of the Saints moved on—to Missouri, to Nauvoo, and to the West—Martin Harris remained in Kirtland. There he was rebaptized by a visiting missionary in 1842. In 1856 Caroline and their four children took the long journey to Utah, but Martin, then 73 years of age, remained on his property in Kirtland. In 1860 he told a census taker that he was a “Mormon preacher,” evidence of his continuing loyalty to the restored gospel. Later he would tell a visitor, “I never did leave the Church; the Church left me” (quoted in William H. Homer Jr., “‘… Publish It upon the Mountains’: The Story of Martin Harris,” Improvement Era, July 1955, 505), meaning of course that Brigham Young led the Church west and the aging Martin remained in Kirtland.

During part of his remaining years in Kirtland, Martin Harris acted as a self-appointed guide-caretaker of the deserted Kirtland Temple, which he loved. Visitors reported his alienation from the leaders of the Church in Utah but also his fervent reaffirmation of his published testimony of the Book of Mormon.

Finally, in 1870, Martin’s desire to be reunited with his family in Utah resulted in a warm invitation from Brigham Young, a ticket for his passage, and an official escort from one of the Presidents of Seventy. A Utah interviewer of the 87-year-old man described him as “remarkably vigorous for one of his years, … his memory being very good” (Deseret Evening News, 31 Aug. 1870). He was rebaptized, a common practice at that time, and spoke twice to audiences in this Tabernacle. We have no official report of what he said, but we can be sure of his central message since over 35 persons left similar personal accounts of what he told them during this period. One reported Martin saying, “It is not a mere belief, but is a matter of knowledge. I saw the plates and the inscriptions thereon. I saw the angel, and he showed them unto me” (quoted in Investigating the Book of Mormon Witnesses, 116).

When he reiterated his testimony of the Book of Mormon in the closing days of his life, Martin Harris declared, “I tell you of these things that you may tell others that what I have said is true, and I dare not deny it; I heard the voice of God commanding me to testify to the same” (quoted in Investigating the Book of Mormon Witnesses, 118).

Martin Harris died in Clarkston, Utah, in 1875, at age 92. His life is commemorated in the memorable pageant Martin Harris: The Man Who Knew, produced each summer in Clarkston, Utah.

What do we learn from this example? (1) Witnesses are important, and the testimony of the Three Witnesses to the Book of Mormon is impressive and reliable. (2) Happiness and spiritual progress lie in following the leaders of the Church. (3) There is hope for each of us, even if we have sinned and strayed from a favored position.

The Lord’s invitation is warm and loving: “Come back and feast at the table of the Lord, and taste again the sweet and satisfying fruits of fellowship with the saints” (The First Presidency, “An Invitation to Come Back,” Church News, 22 Dec. 1985, 3). I testify that this is the word of the Lord and the work of the Lord, in the name of Jesus Christ, amen.

\intalkheader{Bibliography}

\newpage
\settalk{Gospel Teaching}
\setspeaker{Dallin H. Oaks}
\settalkdate{October 1999}
\talkheading{Gospel Teaching}{Dallin H. Oaks}

A national author wrote a book about his greatest teacher. At the heart of this college teacher’s powerful impact on his student was the student’s conviction that this teacher really cared for him and wanted him to learn and do what would help him find happiness. The author concluded his tribute with this question: “Have you ever really had a teacher? One who saw you as a raw but precious thing, a jewel that, with wisdom, could be polished to a proud shine? If you are lucky enough to find your way to such teachers, you will always find your way back.”

\intalkheader{I.}

Every member of The Church of Jesus Christ of Latter-day Saints is, or will be, a teacher. Each of us has a vital interest in the content and effectiveness of gospel teaching. We want everyone to have great gospel teachers, and we want those teachers to help all of us find our way back, not just to them but to our Heavenly Father.

Our concern with gospel teaching is not limited to those who are called to teach in the priesthood quorums, in the Primary, Relief Society, Sunday School, Young Women, and in other assignments. In the Lord’s great plan of salvation there are no more important teachers than parents, who teach their children constantly by example and by precept. Each of us teaches those around us by example. Even children teach one another. Every missionary is a teacher. And every leader is a teacher. As President Hinckley taught many years ago, “Effective teaching is the very essence of leadership in the Church.”

Gospel teaching is universal and important. Truly, “no greater responsibility can rest upon any [one of us] than to be a teacher of God’s children.” Our Savior’s occupation was that of a teacher. He was the Master Teacher, and He invites each of us to follow Him in that great service.

Several years ago the First Presidency challenged the Quorum of the Twelve to revitalize teaching in the Church. The Twelve, assisted by the Seventy, accepted that challenge. Now, after years of preparation, engaging the efforts of superb gospel teachers, scholars, writers, and others, the First Presidency has just sent a letter launching a Churchwide effort “to revitalize and improve teaching in the Church.” This letter states, “This renewed emphasis is intended to improve gospel teaching in homes and in Church meetings and help nourish members with the good word of God.”

We have just published a 10-page booklet, Improving Gospel Teaching: A Leader’s Guide. Copies are being distributed to all unit leaders and to every quorum and auxiliary officer in the Church. As it explains, our concern with “gospel teaching in the Church” includes parents’ everyday teachings in the home as well as the work of teachers in the quorums and auxiliaries.

This important effort to “revitalize and improve teaching in the Church” includes three elements. At the outset, it emphasizes leaders’ vital responsibilities to work to improve gospel teaching in their organizations. We want all leaders to encourage and help the teachers and learners over whom they preside.

Next, the effort initiates quarterly teacher improvement meetings for teachers of three different groups—children, youth, and adults—to “instruct and edify each other” (D\&C 43:8) on principles, methods, and skills that will improve gospel teaching and learning.

Finally, a 12-lesson course on “Teaching the Gospel” will be taught at least once each year, generally during Sunday School. Its course material will be drawn from a new abbreviated and improved edition of Teaching, No Greater Call: A Resource Guide for Gospel Teaching. This book is being distributed to all wards and branches in the Church.

We have also reissued the Teaching Guidebook for use in the home and for smaller and developing units that cannot staff the entire Church program.

\intalkheader{II.}

Some may wonder why we are making such an extensive effort to improve gospel teaching. Those who wonder must be blessed with superior teachers, and we have many of those in the Church. Others will understand why such an effort is needed and will pray for its success.

For many years I have sought to learn more about the nature and quality of teaching in the various quorums and auxiliaries of the Church. I have done this by dropping in unannounced on classes in various wards in different parts of the Church. By now I have visited hundreds of classes. I apologize if any of my visits has terrorized a teacher. My impression is that almost all of the teachers I have observed in these surprise visits have appreciated having a visitor who was there to learn and there to show appreciation for their efforts and concern for their students.

For the most part, what I have seen in these visits has been gratifying and reassuring. I have seen inspired teachers whose love for the gospel and their students was so evident that the effect of their teaching was positively electric. I have also seen thoughtful and respectful students, receptive to the message and hungry to learn.

Notwithstanding the great examples I have observed, I am convinced that in the Church as a whole—as with each of us individually—we can always do better. The challenge of progress is inherent in our Father in Heaven’s plan for His children. And in our sacred callings of gospel teaching, no effort is too good for the work of the Lord and the growth of His children.

\intalkheader{III.}

There are many different ways to teach, but all good teaching is based on certain fundamental principles. Without pretending to be exhaustive, I wish to identify and comment on six fundamental principles of gospel teaching.

The first is love. It has two manifestations. When we are called to teach, we should accept our calling and teach because of our love for God the Eternal Father and His Son, Jesus Christ. In addition, a gospel teacher should always teach with love for the students. We are taught that we should pray “with all the energy of heart … [to] be filled with this love” (Moro. 7:48). Love of God and love of His children is the highest reason for service. Those who teach out of love will be magnified as instruments in the hands of Him whom they serve.

Second, a gospel teacher, like the Master we serve, will concentrate entirely on those being taught. His or her total concentration will be on the needs of the sheep—the good of the students. A gospel teacher does not focus on himself or herself. One who understands that principle will not look upon his or her calling as “giving or presenting a lesson,” because that definition views teaching from the standpoint of the teacher, not the student.

Focusing on the needs of the students, a gospel teacher will never obscure their view of the Master by standing in the way or by shadowing the lesson with self-promotion or self-interest. This means that a gospel teacher must never indulge in priestcrafts, which are “that men preach and set themselves up for a light unto the world, that they may get gain and praise of the world” (2 Ne. 26:29). A gospel teacher does not preach “to become popular” (Alma 1:3) or “for the sake of riches and honor” (Alma 1:16). He or she follows the marvelous Book of Mormon example in which “the preacher was no better than the hearer, neither was the teacher any better than the learner” (Alma 1:26). Both will always look to the Master.

Third, a superior teacher of the gospel will teach from the prescribed course material, with greatest emphasis on teaching the doctrine and principles and covenants of the gospel of Jesus Christ. This is commanded in modern revelation, where the Lord said:

“Teachers of this church shall teach the principles of my gospel, which are in the Bible and the Book of Mormon, in the which is the fulness of the gospel.

“And they shall observe the covenants and church articles to do them, and these shall be their teachings, as they shall be directed by the Spirit” (D\&C 42:12–13).

Teachers who are commanded to teach “the principles of [the] gospel” and “the doctrine of the kingdom” (D\&C 88:77) should generally forgo teaching specific rules or applications. For example, they would not teach any rules for determining what is a full tithing, and they would not provide a list of do’s and don’ts for keeping the Sabbath day holy. Once a teacher has taught the doctrine and the associated principles from the scriptures and the living prophets, such specific applications or rules are generally the responsibility of individuals and families.

Well-taught doctrines and principles have a more powerful influence on behavior than rules. When we teach gospel doctrine and principles, we can qualify for the witness and guidance of the Spirit to reinforce our teaching, and we enlist the faith of our students in seeking the guidance of that same Spirit in applying those teachings in their personal lives.

The subject being taught in the Melchizedek Priesthood quorums and Relief Societies of the Church during the second and third Sundays of each month is the Teachings of Presidents of the Church. During the last two years we have studied the teachings of President Brigham Young. For the next two years we will be studying the teachings of President Joseph F. Smith. The books containing these teachings, which are being given to every adult member of the Church as a permanent personal library resource, contain doctrine and principles. They are rich and relevant to the needs of our day, and they are superb for teaching and discussion.

As I have visited in quorums and Relief Societies, I have generally been pleased and impressed at how these Teachings of Presidents of the Church are being presented and received. However, I have sometimes observed teachers who gave the designated chapter no more than a casual mention and then presented a lesson and invited discussion on other materials of the teacher’s choice. That is not acceptable. A gospel teacher is not called to choose the subject of the lesson but to teach and discuss what has been specified. Gospel teachers should also be scrupulous to avoid hobby topics, personal speculations, and controversial subjects. The Lord’s revelations and the directions of His servants are clear on this point. We should all be mindful of President Spencer W. Kimball’s great instruction that a gospel teacher is a “guest”:

“He has been given an authoritative position and a stamp of approval is placed upon him, and those whom he teaches are justified in assuming that, having been chosen and sustained in the proper order, he represents the Church and the things which he teaches are approved by the Church. No matter how brilliant he may be and how many new truths he may think he has found, he has no right to go beyond the program of the Church.”

Fourth, a gospel teacher will prepare diligently and strive to use the most effective means of presenting the prescribed lessons. The new Teaching the Gospel course and the new teacher improvement meetings are obviously intended to assist teachers in this effort.

The fifth fundamental principle of gospel teaching I wish to stress is the Lord’s command, quoted earlier, that gospel teachers should “teach the principles of my gospel … as they shall be directed by the Spirit. … And if ye receive not the Spirit ye shall not teach” (D\&C 42:12–14). It is a gospel teacher’s privilege and duty to seek that level of discipleship where his or her teachings will be directed and endorsed by the Spirit rather than being rigidly selected and prearranged for personal convenience or qualifications. The marvelous principles of “Gospel Teaching and Leadership” in the new Church Handbook of Instructions include the following:

“Teachers and class members should seek the Spirit during the lesson. A person may teach profound truths, and class members may engage in stimulating discussions, but unless the Spirit is present, these things will not be powerfully impressed upon the soul. …

“When the Spirit is present in gospel teaching, ‘the power of the Holy Ghost carrieth [the message] unto the hearts of the children of men’ (2 Ne. 33:1).”

President Hinckley stated an important corollary to the command to teach by the Spirit when he issued this challenge:

“We must … get our teachers to speak out of their hearts rather than out of their books, to communicate their love for the Lord and this precious work, and somehow it will catch fire in the hearts of those they teach.”

That is our objective—to have love of God and commitment to the gospel of Jesus Christ “catch fire” in the hearts of those we teach.

That leads to the sixth and final principle I will discuss. A gospel teacher is concerned with the results of his or her teaching, and such a teacher will measure the success of teaching and testifying by its impact on the lives of the learners. A gospel teacher will never be satisfied with just delivering a message or preaching a sermon. A superior gospel teacher wants to assist in the Lord’s work to bring eternal life to His children.

President Harold B. Lee said: “The calling of the gospel teacher is one of the noblest in the world. The good teacher can make all the difference in inspiring boys and girls and men and women to change their lives and fulfill their highest destiny. The importance of the teacher has been beautifully described by Daniel Webster when he said, ‘If we work upon marble, it will perish; if we work upon brass, time will efface it; but if we work upon immortal minds, if we imbue them with principles and the just fear of God and love of our fellowman, we engrave upon those tablets something that will brighten through all eternity.’”

I testify that this is God’s work, and that we are His servants with the sacred responsibility of teaching the gospel of Jesus Christ, the greatest message of all time. We need more teachers to match that message. I pray that we will all become superior gospel teachers, in the name of Jesus Christ, amen.

\newpage
\settalk{Resurrection}
\setspeaker{Dallin H. Oaks}
\settalkdate{April 2000}
\talkheading{Resurrection}{Dallin H. Oaks}

The book of Job poses the universal question, “If a man die, shall he live again?” (Job 14:14). The question of resurrection from the dead is a central subject of scripture, ancient and modern. The resurrection is a pillar of our faith. It adds meaning to our doctrine, motivation to our behavior, and hope for our future.

\intalkheader{I. The Resurrection of Jesus}

The universal resurrection became a reality with the Resurrection of Jesus Christ (see Matt. 27:52–53). On the third day after His death and burial, Jesus came forth out of the tomb. He appeared to several men and women, and then to the assembled Apostles. Three of the Gospels describe this event. Luke is the most complete:

“Jesus … saith unto them, Peace be unto you.

“But they were terrified and affrighted, and supposed that they had seen a spirit.

“And he said unto them, Why are ye troubled? and why do thoughts arise in your hearts?

“Behold my hands and my feet, that it is I myself: handle me, and see; for a spirit hath not flesh and bones, as ye see me have. …

“Then opened he their understanding, …

“And said unto them, Thus it is written, and thus it behoved Christ to suffer, and to rise from the dead the third day” (Luke 24:36–39, 45–46).

The Savior gave the Apostles a second witness. Thomas, one of the Twelve, had not been with them when Jesus came. He insisted that he would not believe unless he could see and feel for himself. John records:

“And after eight days again his disciples were within, and Thomas with them: then came Jesus, the doors being shut, and stood in the midst, and said, Peace be unto you.

“Then saith he to Thomas, Reach hither thy finger, and behold my hands; and reach hither thy hand, and thrust it into my side: and be not faithless, but believing.

“And Thomas answered and said unto him, My Lord and my God.

“Jesus saith unto him, Thomas, because thou hast seen me, thou hast believed: blessed are they that have not seen, and yet have believed” (John 20:26–29).

Despite these biblical witnesses, many who call themselves Christians reject or confess serious doubts about the reality of the resurrection. As if to anticipate and counter such doubts, the Bible records many appearances of the risen Christ. In some of these He appeared to a single individual, such as to Mary Magdalene at the sepulchre. In others He appeared to large or small groups, such as when “he was seen of [about] five hundred brethren at once” (1 Cor. 15:6).

The Book of Mormon: Another Testament of Jesus Christ records the experience of hundreds who saw the risen Lord in person and touched Him, feeling the prints of the nails in His hands and feet and thrusting their hands into His side. The Savior invited a multitude to have this experience “one by one” (3 Ne. 11:15) so that they could know that He was “the God of Israel, and the God of the whole earth, and [had] been slain for the sins of the world” (3 Ne. 11:14).

During the course of His personal ministry among these faithful people, the resurrected Christ healed the sick and also “took their little children, one by one, and blessed them” (3 Ne. 17:21). This tender episode was witnessed by about 2,500 men, women, and children (see 3 Ne. 17:25).

\intalkheader{II. The Resurrection of Mortals}

The possibility that a mortal who has died will be brought forth and live again in a resurrected body has awakened hope and stirred controversy through much of recorded history. Relying on clear scriptural teachings, Latter-day Saints join in affirming that Christ has “broken the bands of death” (Mosiah 16:7) and that “death is swallowed up in victory” (1 Cor. 15:54; see also Morm. 7:5; Mosiah 15:8; 16:7–8; Alma 22:14). Because we believe the Bible and Book of Mormon descriptions of the literal Resurrection of Jesus Christ, we also readily accept the numerous scriptural teachings that a similar resurrection will come to all mortals who have ever lived upon this earth (see 1 Cor. 15:22; 2 Ne. 9:22; Hel. 14:17; Morm. 9:13; D\&C 29:26; 76:39, 42–44). As Jesus taught, “Because I live, ye shall live also” (John 14:19).

The literal and universal nature of the resurrection is vividly described in the Book of Mormon. The prophet Amulek taught:

“The death of Christ shall loose the bands of this temporal death, that all shall be raised from this temporal death.

“The spirit and the body shall be reunited again in its perfect form; both limb and joint shall be restored to its proper frame, even as we now are at this time; …

“Now, this restoration shall come to all, both old and young, both bond and free, both male and female, both the wicked and the righteous; and even there shall not so much as a hair of their heads be lost; but every thing shall be restored to its perfect frame” (Alma 11:42–44).

Alma also taught that in the resurrection “all things shall be restored to their proper and perfect frame” (Alma 40:23).

Many living witnesses can testify to the literal fulfillment of these scriptural assurances of the resurrection. Many, including some in my own extended family, have seen a departed loved one in vision or personal appearance and have witnessed their restoration in “proper and perfect frame” in the prime of life. Whether these were manifestations of persons already resurrected or of righteous spirits awaiting an assured resurrection, the reality and nature of the resurrection of mortals is evident. What a comfort to know that all who have been disadvantaged in life from birth defects, from mortal injuries, from disease, or from the natural deterioration of old age will be resurrected in “proper and perfect frame.”

\intalkheader{III. The Significance of the Resurrection}

I wonder if we fully appreciate the enormous significance of our belief in a literal, universal resurrection. The assurance of immortality is fundamental to our faith. The Prophet Joseph Smith declared:

“The fundamental principles of our religion are the testimony of the Apostles and Prophets, concerning Jesus Christ, that He died, was buried, and rose again the third day, and ascended into heaven; and all other things which pertain to our religion are only appendages to it” (Teachings of the Prophet Joseph Smith, sel. Joseph Fielding Smith [1976], 121).

Of all things in that glorious ministry, why did the Prophet Joseph Smith use the testimony of the Savior’s death, burial, and Resurrection as the fundamental principle of our religion, saying that “all other things … are only appendages to it”? The answer is found in the fact that the Savior’s Resurrection is central to what the prophets have called “the great and eternal plan of deliverance from death” (2 Ne. 11:5).

In our eternal journey, the resurrection is the mighty milepost that signifies the end of mortality and the beginning of immortality. The Lord described the importance of this vital transition when He declared, “And thus did I, the Lord God, appoint unto man the days of his probation—that by his natural death he might be raised in immortality unto eternal life, even as many as would believe” (D\&C 29:43). Similarly, the Book of Mormon teaches, “For as death hath passed upon all men, to fulfil the merciful plan of the great Creator, there must needs be a power of resurrection” (2 Ne. 9:6). We also know, from modern revelation, that without the reuniting of our spirits and our bodies in the resurrection we could not receive a “fulness of joy” (D\&C 93:33–34).

When we understand the vital position of the resurrection in the “plan of redemption” that governs our eternal journey (Alma 12:25), we see why the Apostle Paul taught, “If there be no resurrection of the dead, then … is our preaching vain, and your faith is also vain” (1 Cor. 15:13–14). We also see why the Apostle Peter referred to the fact that God the Father, in His abundant mercy, “hath begotten us again unto a lively hope by the resurrection of Jesus Christ from the dead” (1 Pet. 1:3; see also 1 Thes. 4:13–18).

\intalkheader{IV. The Resurrection Changes Our View of Mortality}

The “lively hope” we are given by the resurrection is our conviction that death is not the conclusion of our identity but merely a necessary step in the destined transition from mortality to immortality. This hope changes the whole perspective of mortal life. The assurance of resurrection and immortality affects how we look on the physical challenges of mortality, how we live our mortal lives, and how we relate to those around us.

The assurance of resurrection gives us the strength and perspective to endure the mortal challenges faced by each of us and by those we love, such things as the physical, mental, or emotional deficiencies we bring with us at birth or acquire during mortal life. Because of the resurrection, we know that these mortal deficiencies are only temporary!

The assurance of resurrection also gives us a powerful incentive to keep the commandments of God during our mortal lives. Resurrection is much more than merely reuniting a spirit to a body held captive by the grave. We know from the Book of Mormon that the resurrection is a restoration that brings back “carnal for carnal” and “good for that which is good” (Alma 41:13; see also Alma 41:2–4 and Hel. 14:31). The prophet Amulek taught, “That same spirit which doth possess your bodies at the time that ye go out of this life, that same spirit will have power to possess your body in that eternal world” (Alma 34:34). As a result, when persons leave this life and go on to the next, “they who are righteous shall be righteous still” (2 Ne. 9:16), and “whatever principle of intelligence we attain unto in this life … will rise with us in the resurrection” (D\&C 130:18).

The principle of restoration also means that persons who are not righteous in mortal life will not rise up righteous in the resurrection (see 2 Ne. 9:16; 1 Cor. 15:35–44; D\&C 88:27–32). Moreover, unless our mortal sins have been cleansed and blotted out by repentance and forgiveness (see Alma 5:21; 2 Ne. 9:45–46; D\&C 58:42), we will be resurrected with a “bright recollection” (Alma 11:43) and a “perfect knowledge of all our guilt, and our uncleanness” (2 Ne. 9:14; see also Alma 5:18). The seriousness of that reality is emphasized by the many scriptures suggesting that the resurrection is followed immediately by the Final Judgment (see 2 Ne. 9:15, 22; Mosiah 26:25; Alma 11:43–44; 42:23; Morm. 7:6; 9:13–14). Truly, “this life is the time for men to prepare to meet God” (Alma 34:32).

The assurance that the resurrection will include an opportunity to be with our family members—husband, wife, parents, brothers and sisters, children, and grandchildren—is a powerful encouragement for us to fulfill our family responsibilities in mortality. It helps us live together in love in this life in anticipation of joyful reunions and associations in the next.

Our sure knowledge of a resurrection to immortality also gives us the courage to face our own death—even a death that we might call premature. Thus, the people of Ammon in the Book of Mormon “never did look upon death with any degree of terror, for their hope and views of Christ and the resurrection; therefore, death was swallowed up to them by the victory of Christ over it” (Alma 27:28).

The assurance of immortality also helps us bear the mortal separations involved in the death of our loved ones. Every one of us has wept at a death, grieved through a funeral, or stood in pain at a graveside. I am surely one who has. We should all praise God for the assured resurrection that makes our mortal separations temporary and gives us the hope and strength to carry on.

\intalkheader{V. The Resurrection and Temples}

We are living in a glorious season of temple building. This is also a consequence of our faith in the resurrection. Just a few months ago I was privileged to accompany President Hinckley to the dedication of a new temple. In that sacred setting I heard him say:

“Temples stand as a witness of our conviction of immortality. Our temples are concerned with life beyond the grave. For example, there is no need for marriage in the temple if we were only concerned with being married for the period of our mortal lives.”

This prophetic teaching enlarged my understanding. Our temples are living, working testimonies to our faith in the reality of the resurrection. They provide the sacred settings where living proxies can perform all of the necessary ordinances of mortal life in behalf of those who live in the world of the spirits. None of this would be meaningful if we did not have the assurance of universal immortality and the opportunity for eternal life because of the Resurrection of our Lord and Savior, Jesus Christ.

We believe in the literal, universal resurrection of all mankind because of “the resurrection of the Holy One of Israel” (2 Ne. 9:12). We also testify of “the Living Christ,” as was said in the recent apostolic declaration of that same name:

“We solemnly testify that His life, which is central to all human history, neither began in Bethlehem nor concluded on Calvary. …

“We bear testimony, as His duly ordained Apostles—that Jesus is the Living Christ, the immortal Son of God. He is the great King Immanuel, who stands today on the right hand of His Father. He is the light, the life, and the hope of the world. His way is the path that leads to happiness in this life and eternal life in the world to come” (“The Living Christ: The Testimony of the Apostles,” Ensign, Apr. 2000, 2–3).

I testify of that reality and of the reality of His Resurrection and ours, in the name of Jesus Christ, amen.

\newpage
\settalk{The Challenge to Become}
\setspeaker{Dallin H. Oaks}
\settalkdate{October 2000}
\talkheading{The Challenge to Become}{Dallin H. Oaks}

The Apostle Paul taught that the Lord’s teachings and teachers were given that we may all attain “the measure of the stature of the fulness of Christ” (Eph. 4:13). This process requires far more than acquiring knowledge. It is not even enough for us to be convinced of the gospel; we must act and think so that we are converted by it. In contrast to the institutions of the world, which teach us to know something, the gospel of Jesus Christ challenges us to become something.

Many Bible and modern scriptures speak of a final judgment at which all persons will be rewarded according to their deeds or works or the desires of their hearts. But other scriptures enlarge upon this by referring to our being judged by the condition we have achieved.

The prophet Nephi describes the Final Judgment in terms of what we have become: “And if their works have been filthiness they must needs be filthy; and if they be filthy it must needs be that they cannot dwell in the kingdom of God” (1 Ne. 15:33; emphasis added). Moroni declares, “He that is filthy shall be filthy still; and he that is righteous shall be righteous still” (Morm. 9:14; emphasis added; see also Rev. 22:11–12; 2 Ne. 9:16; D\&C 88:35). The same would be true of “selfish” or “disobedient” or any other personal attribute inconsistent with the requirements of God. Referring to the “state” of the wicked in the Final Judgment, Alma explains that if we are condemned by our words, our works, and our thoughts, “we shall not be found spotless; … and in this awful state we shall not dare to look up to our God” (Alma 12:14).

From such teachings we conclude that the Final Judgment is not just an evaluation of a sum total of good and evil acts—what we have done. It is an acknowledgment of the final effect of our acts and thoughts—what we have become. It is not enough for anyone just to go through the motions. The commandments, ordinances, and covenants of the gospel are not a list of deposits required to be made in some heavenly account. The gospel of Jesus Christ is a plan that shows us how to become what our Heavenly Father desires us to become.

A parable illustrates this understanding. A wealthy father knew that if he were to bestow his wealth upon a child who had not yet developed the needed wisdom and stature, the inheritance would probably be wasted. The father said to his child:

“All that I have I desire to give you—not only my wealth, but also my position and standing among men. That which I have I can easily give you, but that which I am you must obtain for yourself. You will qualify for your inheritance by learning what I have learned and by living as I have lived. I will give you the laws and principles by which I have acquired my wisdom and stature. Follow my example, mastering as I have mastered, and you will become as I am, and all that I have will be yours.”

This parable parallels the pattern of heaven. The gospel of Jesus Christ promises the incomparable inheritance of eternal life, the fulness of the Father, and reveals the laws and principles by which it can be obtained.

We qualify for eternal life through a process of conversion. As used here, this word of many meanings signifies not just a convincing but a profound change of nature. Jesus used this meaning when He taught His chief Apostle the difference between a testimony and a conversion. Jesus asked His disciples, “Whom do men say that I the Son of man am?” (Matt. 16:13). Next He asked, “But whom say ye that I am?

“And Simon Peter answered and said, Thou art the Christ, the Son of the living God.

“And Jesus answered and said unto him, Blessed art thou, Simon Bar-jona: for flesh and blood hath not revealed it unto thee, but my Father which is in heaven” (Matt. 16:15–17).

Peter had a testimony. He knew that Jesus was the Christ, the promised Messiah, and he declared it. To testify is to know and to declare.

Later on, Jesus taught these same men about conversion, which is far more than testimony. When the disciples asked who was the greatest in the kingdom of heaven, “Jesus called a little child unto him, and set him in the midst of them,

“And said, Verily I say unto you, Except ye be converted, and become as little children, ye shall not enter into the kingdom of heaven.

“Whosoever therefore shall humble himself as this little child, the same is greatest in the kingdom of heaven” (Matt. 18:2–4; emphasis added).

Later the Savior confirmed the importance of being converted, even for those with a testimony of the truth. In the sublime instructions given at the Last Supper, He told Simon Peter, “I have prayed for thee, that thy faith fail not: and when thou art converted, strengthen thy brethren” (Luke 22:32).

In order to strengthen his brethren—to nourish and lead the flock of God—this man who had followed Jesus for three years, who had been given the authority of the holy apostleship, who had been a valiant teacher and testifier of the Christian gospel, and whose testimony had caused the Master to declare him blessed still had to be “converted.”

Jesus’ challenge shows that the conversion He required for those who would enter the kingdom of heaven (see Matt. 18:3) was far more than just being converted to testify to the truthfulness of the gospel. To testify is to know and to declare. The gospel challenges us to be “converted,” which requires us to do and to become. If any of us relies solely upon our knowledge and testimony of the gospel, we are in the same position as the blessed but still unfinished Apostles whom Jesus challenged to be “converted.” We all know someone who has a strong testimony but does not act upon it so as to be converted. For example, returned missionaries, are you still seeking to be converted, or are you caught up in the ways of the world?

The needed conversion by the gospel begins with the introductory experience the scriptures call being “born again” (e.g., Mosiah 27:25; Alma 5:49; John 3:7; 1 Pet. 1:23). In the waters of baptism and by receiving the gift of the Holy Ghost, we become the spiritual “sons and daughters” of Jesus Christ, “new creatures” who can “inherit the kingdom of God” (Mosiah 27:25–26).

In teaching the Nephites, the Savior referred to what they must become. He challenged them to repent and be baptized and be sanctified by the reception of the Holy Ghost, “that ye may stand spotless before me at the last day” (3 Ne. 27:20). He concluded: “Therefore, what manner of men ought ye to be? Verily I say unto you, even as I am” (3 Ne. 27:27).

The gospel of Jesus Christ is the plan by which we can become what children of God are supposed to become. This spotless and perfected state will result from a steady succession of covenants, ordinances, and actions, an accumulation of right choices, and from continuing repentance. “This life is the time for men to prepare to meet God” (Alma 34:32).

Now is the time for each of us to work toward our personal conversion, toward becoming what our Heavenly Father desires us to become. As we do so, we should remember that our family relationships—even more than our Church callings—are the setting in which the most important part of that development can occur. The conversion we must achieve requires us to be a good husband and father or a good wife and mother. Being a successful Church leader is not enough. Exaltation is an eternal family experience, and it is our mortal family experiences that are best suited to prepare us for it.

The Apostle John spoke of what we are challenged to become when he said, “Beloved, now are we the sons of God, and it doth not yet appear what we shall be: but we know that, when he shall appear, we shall be like him; for we shall see him as he is” (1 Jn. 3:2; see also Moro. 7:48).

I hope the importance of conversion and becoming will cause our local leaders to reduce their concentration on statistical measures of actions and to focus more on what our brothers and sisters are and what they are striving to become.

Our needed conversions are often achieved more readily by suffering and adversity than by comfort and tranquillity, as Elder Hales taught us so beautifully this morning. Father Lehi promised his son Jacob that God would “consecrate [his] afflictions for [his] gain” (2 Ne. 2:2). The Prophet Joseph was promised, “Thine adversity and thine afflictions shall be but a small moment; and then, if thou endure it well, God shall exalt thee on high” (D\&C 121:7–8).

Most of us experience some measure of what the scriptures call “the furnace of affliction” (Isa. 48:10; 1 Ne. 20:10). Some are submerged in service to a disadvantaged family member. Others suffer the death of a loved one or the loss or postponement of a righteous goal like marriage or childbearing. Still others struggle with personal impairments or with feelings of rejection, inadequacy, or depression. Through the justice and mercy of a loving Father in Heaven, the refinement and sanctification possible through such experiences can help us achieve what God desires us to become.

We are challenged to move through a process of conversion toward that status and condition called eternal life. This is achieved not just by doing what is right, but by doing it for the right reason—for the pure love of Christ. The Apostle Paul illustrated this in his famous teaching about the importance of charity (see 1 Cor. 13). The reason charity never fails and the reason charity is greater than even the most significant acts of goodness he cited is that charity, “the pure love of Christ” (Moro. 7:47), is not an act but a condition or state of being. Charity is attained through a succession of acts that result in a conversion. Charity is something one becomes. Thus, as Moroni declared, “except men shall have charity they cannot inherit” the place prepared for them in the mansions of the Father (Ether 12:34; emphasis added).

All of this helps us understand an important meaning of the parable of the laborers in the vineyard, which the Savior gave to explain what the kingdom of heaven is like. As you remember, the owner of the vineyard hired laborers at different times of the day. Some he sent into the vineyard early in the morning, others about the third hour, and others in the sixth and ninth hours. Finally, in the eleventh hour he sent others into the vineyard, promising that he would also pay them “whatsoever is right” (Matt. 20:7).

At the end of the day the owner of the vineyard gave the same wage to every worker, even to those who had come in the eleventh hour. When those who had worked the entire day saw this, “they murmured against the goodman of the house” (Matt. 20:11). The owner did not yield but merely pointed out that he had done no one any wrong, since he had paid each man the agreed amount.

Like other parables, this one can teach several different and valuable principles. For present purposes its lesson is that the Master’s reward in the Final Judgment will not be based on how long we have labored in the vineyard. We do not obtain our heavenly reward by punching a time clock. What is essential is that our labors in the workplace of the Lord have caused us to become something. For some of us, this requires a longer time than for others. What is important in the end is what we have become by our labors. Many who come in the eleventh hour have been refined and prepared by the Lord in ways other than formal employment in the vineyard. These workers are like the prepared dry mix to which it is only necessary to “add water”—the perfecting ordinance of baptism and the gift of the Holy Ghost. With that addition—even in the eleventh hour—these workers are in the same state of development and qualified to receive the same reward as those who have labored long in the vineyard.

This parable teaches us that we should never give up hope and loving associations with family members and friends whose fine qualities (see Moro. 7:5–14) evidence their progress toward what a loving Father would have them become. Similarly, the power of the Atonement and the principle of repentance show that we should never give up on loved ones who now seem to be making many wrong choices.

Instead of being judgmental about others, we should be concerned about ourselves. We must not give up hope. We must not stop striving. We are children of God, and it is possible for us to become what our Heavenly Father would have us become.

How can we measure our progress? The scriptures suggest various ways. I will mention only two.

After King Benjamin’s great sermon, many of his hearers cried out that the Spirit of the Lord “has wrought a mighty change in us, or in our hearts, that we have no more disposition to do evil, but to do good continually” (Mosiah 5:2). If we are losing our desire to do evil, we are progressing toward our heavenly goal.

The Apostle Paul said that persons who have received the Spirit of God “have the mind of Christ” (1 Cor. 2:16). I understand this to mean that persons who are proceeding toward the needed conversion are beginning to see things as our Heavenly Father and His Son, Jesus Christ, see them. They are hearing His voice instead of the voice of the world, and they are doing things in His way instead of by the ways of the world.

I testify of Jesus Christ, our Savior and our Redeemer, whose Church this is. I testify with gratitude of the plan of the Father, under which, through the Resurrection and Atonement of our Savior, we have the assurance of immortality and the opportunity to become what is necessary for eternal life. In the name of Jesus Christ, amen.

\newpage
\settalk{Focus and Priorities}
\setspeaker{Dallin H. Oaks}
\settalkdate{April 2001}
\talkheading{Focus and Priorities}{Dallin H. Oaks}

As we approach the conclusion of this wonderful conference, it is timely to ask ourselves what we are going to strive to become because of what we have heard from the Lord’s servants.

We are accountable and will be judged for how we use what we have received. This eternal principle applies to all we have been given. In the parable of the talents (see Matt. 25:14–30), the Savior taught this principle with reference to the use of property. The principle of accountability also applies to the spiritual resources conferred in the teachings we have been given and to the precious hours and days allotted to each of us during our time in mortality.

I wish to examine how this principle of accountability applies to our use of the enlarged time and information we have been given in our day.

Because of increased life expectancies and modern timesaving devices, most of us have far more discretionary time than our predecessors. We are accountable for how we use that time. “Thou shalt not idle away thy time,” and “Cease to be idle” (D\&C 60:13; 88:124), the Lord commanded the early missionaries and members. “Time flies on wings of lightning,” we sing in a popular hymn; “we cannot call it back. It comes, then passes forward along its onward track. And if we are not mindful, the chance will fade away, for life is quick in passing. ’Tis as a single day” (“Improve the Shining Moments,” Hymns, no. 226).

The significance of our increased discretionary time has been magnified many times by modern data-retrieval technology. For good or for evil, devices like the Internet and the compact disc have put at our fingertips an incredible inventory of information, insights, and images. Along with fast food, we have fast communications and fast facts. The effect of these resources on some of us seems to fulfill the prophet Daniel’s prophecy that in the last days “knowledge shall be increased” and “many shall run to and fro” (Dan. 12:4).

With greatly increased free time and vastly more alternatives for its use, it is prudent to review the fundamental principles that should guide us. Temporal circumstances change, but the eternal laws and principles that should guide our choices never change.

\intalkheader{I.}

A homely story contains a warning. I like this story because it translates easily into different languages and cultures.

Two men formed a partnership. They built a small shed beside a busy road. They obtained a truck and drove it to a farmer’s field, where they purchased a truckload of melons for a dollar a melon. They drove the loaded truck to their shed by the road, where they sold their melons for a dollar a melon. They drove back to the farmer’s field and bought another truckload of melons for a dollar a melon. Transporting them to the roadside, they again sold them for a dollar a melon. As they drove back toward the farmer’s field to get another load, one partner said to the other, “We’re not making much money on this business, are we?” “No, we’re not,” his partner replied. “Do you think we need a bigger truck?”

We don’t need a bigger truckload of information, either. Like the two partners in my story, our biggest need is a clearer focus on how we should value and use what we already have.

Because of modern technology, the contents of huge libraries and other data resources are at the fingertips of many of us. Some choose to spend countless hours in unfocused surfing the Internet, watching trivial television, or scanning other avalanches of information. But to what purpose? Those who engage in such activities are like the two partners in my story, hurrying to and fro, hauling more and more but failing to grasp the essential truth that we cannot make a profit from our efforts until we understand the true value of what is already within our grasp.

A poet described this delusion as an “endless cycle” that brings “knowledge of words, and ignorance of the Word,” in which “wisdom” is “lost in knowledge” and “knowledge” is “lost in information” (T. S. Eliot, “Choruses from ‘The Rock,’” in The Complete Poems and Plays, 1909–1950 [1962], 96).

We have thousands of times more available information than Thomas Jefferson or Abraham Lincoln. Yet which of us would think ourselves a thousand times more educated or more serviceable to our fellowmen than they? The sublime quality of what these two men gave to us—including the Declaration of Independence and the Gettysburg Address—was not attributable to their great resources of information, for their libraries were comparatively small by our standards. Theirs was the wise and inspired use of a limited amount of information.

Available information wisely used is far more valuable than multiplied information allowed to lie fallow. I had to learn this obvious lesson as a law student.

Over 45 years ago, I was introduced to a law library with hundreds of thousands of law books. (Today such a library would include millions of additional pages available by electronic data retrieval.) When I began to prepare an assigned paper, I spent many days searching in hundreds of books for the needed material. I soon learned the obvious truth (already familiar to experienced researchers) that I could never complete my assigned task within the available time unless I focused my research in the beginning and stopped that research soon enough to have time to analyze my findings and compose my conclusions.

Faced with an excess of information in the marvelous resources we have been given, we must begin with focus or we are likely to become like those in the well-known prophecy about people in the last days—“ever learning, and never able to come to the knowledge of the truth” (2 Tim. 3:7). We also need quiet time and prayerful pondering as we seek to develop information into knowledge and mature knowledge into wisdom.

We also need focus to avoid what is harmful. The abundant information and images accessible on the Internet call for sharp focus and control to avoid accessing the pornography that is an increasing scourge in our society. As the Deseret News noted in a recent editorial, “Images that used to be hidden in out-of-the-way store counters now are as close as a mouse click” (“Staying ahead of Pornography,” 21–22 Feb. 2001, A12). The Internet has made pornography accessible almost without effort and often without leaving the privacy of one’s home or room. The Internet has also facilitated the predatory activities of adults who use its anonymity and accessibility to stalk children for evil purposes. Parents and youth, beware!

There are many gospel implications of this easily accessible flood of information. For example, our Church Web site now provides access to all of the general conference addresses and other contents of Church magazines for the past 30 years. Teachers can download bales of information on any subject. When highly focused, a handout can enrich. But a bale of handouts can detract from our attempt to teach gospel principles with clarity and testimony. Stacks of supplementary material can impoverish rather than enrich, because they can blur students’ focus on the assigned principles and draw them away from prayerfully seeking to apply those principles in their own lives.

Nephi taught, “Feast upon the words of Christ; for behold, the words of Christ will tell you all things what ye should do” (2 Ne. 32:3). That is focus. Nephi also said that as he taught from the scriptures, “I did liken all scriptures unto us, that it might be for our profit and learning” (1 Ne. 19:23). That is personal application.

As a further illustration of the need for focus in using and teaching from the great information resources of the past, consider the comparative value today of the advice Brigham Young gave to an audience 140 years ago with what President Hinckley and other servants of the Lord are saying to each of us right now, in this conference. Or compare the value to each of us of some other facts or advice from the distant past with what our stake president said at our last stake conference or what our bishop counseled us last Sunday.

Overarching all of this is the importance of what the Spirit whispered to us last night or this morning about our own specific needs. Each of us should be careful that the current flood of information does not occupy our time so completely that we cannot focus on and hear and heed the still, small voice that is available to guide each of us with our own challenges today.

I hope that these cautions on the need for focus will not be understood as hostile to selective use of the new technology that has put such a wealth of information at our fingertips. In this I echo Brigham Young, who declared:

“Every discovery in science and art, that is really true and useful to mankind, has been given by direct revelation from God. … We should take advantage of all these great discoveries … and give to our children the benefit of every branch of useful knowledge, to prepare them to step forward and efficiently do their part in the great work” (Deseret News, 22 Oct. 1862, 129).

\intalkheader{II.}

We also need priorities. Our priorities determine what we seek in life. Most of what has been taught in this conference concerns priorities. I hope we will heed these teachings.

Jesus taught about priorities when He said, “Seek not the things of this world but seek ye first to build up the kingdom of God, and to establish his righteousness, and all these things shall be added unto you” (JST, Matt. 6:38; in Matt. 6:33, footnote a). “Seek … first to build up the kingdom of God” means to assign first priority to God and to His work. The work of God is to bring to pass the eternal life of His children (see Moses 1:39), and all that this entails in the birth, nurturing, teaching, and sealing of our Heavenly Father’s children. Everything else is lower in priority. Think about that reality as we consider some teachings and some examples on priorities. As someone has said, if we do not choose the kingdom of God first, it will make little difference in the long run what we have chosen instead of it.

As regards knowledge, the highest priority religious knowledge is what we receive in the temple. That knowledge is obtained from the explicit and symbolic teachings of the endowment and from the whisperings of the Spirit that come as we are desirous to seek and receptive to hear the revelation available to us in that sacred place.

As regards property, Jesus taught that “a man’s life consisteth not in the abundance of the things which he possesseth” (Luke 12:15). Consequently, we should not lay up for ourselves “treasures upon earth, where moth and rust doth corrupt, and where thieves break through and steal” (Matt. 6:19). In other words, the treasures of our hearts—our priorities—should not be what the scriptures call “riches [and] the vain things of this world” (Alma 39:14). The “vain things of [the] world” include every combination of that worldly quartet of property, pride, prominence, and power. As to all of these, the scriptures remind us that “you cannot carry them with you” (Alma 39:14). We should be seeking the kind of treasures the scriptures promise the faithful: “great treasures of knowledge, even hidden treasures” (D\&C 89:19).

All around us we have the good examples of those who seek permanent treasures—those who “hunger and thirst after righteousness” (Matt. 5:6) and put the kingdom of God first in their lives. Among the most visible such examples are the men and women who set aside their worldly pursuits and even say good-bye to their families to serve missions for the Lord. Tens of thousands of these are young missionaries. In addition, I pay particular tribute to those who serve missions in their mature years, some as mission leaders and some as what we call couple missionaries. Their remarkable service evidences their priorities, and their impressive example is a guide to their families and to all who know them.

Our priorities are most visible in how we use our time. Someone has said, “Three things never come back—the spent arrow, the spoken word, and the lost opportunity.” We cannot recycle or save the time allotted to us each day. With time, we have only one opportunity for choice, and then it is gone forever.

Good choices are especially important in our family life. For example, how do family members spend their free time together? Time together is necessary but not sufficient. Priorities should govern us in the precious time we give to our family relationships. Compare the impact of time spent merely in the same room as spectators for television viewing with the significance of time spent communicating with one another individually and as a family.

To cite another example, how much time does a family allocate to learning the gospel by scripture study and parental teachings, in contrast to the time family members spend viewing sports contests, talk shows, or soap operas? I believe many of us are overnourished on entertainment junk food and undernourished on the bread of life.

In terms of priorities for each major decision (such as education, occupation, place of residence, marriage, or childbearing), we should ask ourselves, “What will be the eternal impact of this decision?” Some decisions that seem desirable for mortality have unacceptable risks for eternity. In all such choices, we need to have inspired priorities and apply them in ways that will bring eternal blessings to us and to our family members.

Then, after we have done all that we can, we should remember the wise counsel and comforting assurance of King Benjamin, who taught, “And see that all these things are done in wisdom and order; for it is not requisite that a man should run faster than he has strength” (Mosiah 4:27).

The ultimate Latter-day Saint priorities are twofold: First, we seek to understand our relationship to God the Eternal Father and His Son, Jesus Christ, and to secure that relationship by obtaining their saving ordinances and by keeping our personal covenants. Second, we seek to understand our relationship to our family members and to secure those relationships by the ordinances of the temple and by keeping the covenants we make in that holy place. These relationships, secured in the way I have explained, provide eternal blessings available in no other way. No combination of science, success, property, pride, prominence, or power can provide these eternal blessings!

I testify that this is true, and I testify of God the Father, whose plan establishes the way, and of our Savior, Jesus Christ, whose Atonement makes it all possible.

In the name of Jesus Christ, amen.

\newpage
\settalk{Sharing the Gospel}
\setspeaker{Dallin H. Oaks}
\settalkdate{October 2001}
\talkheading{Sharing the Gospel}{Dallin H. Oaks}

Thank you, President Hinckley, for your great message. We are all profoundly grateful for your vigorous and inspired leadership in this difficult time. Under that leadership, we are going forward with the work of the Lord, so urgently needed in this troubled world.

To proclaim the good news of the gospel of Jesus Christ is a fundamental principle of the Christian faith. Three of the gospel writers report this direction by the Savior.

The book of Mark records: “And he said unto them, Go ye into all the world, and preach the gospel to every creature.

“He that believeth and is baptized shall be saved; but he that believeth not shall be damned” (Mark 16:15–16).

Matthew quotes the Savior’s command, “Go ye therefore, and teach all nations, baptizing them in the name of the Father, and of the Son, and of the Holy Ghost” (Matt. 28:19).

Luke states, “Thus it is written … that repentance and remission of sins should be preached in his name among all nations” (Luke 24:46–47).

Applying the Savior’s directions to our day, modern prophets have challenged each of us to share the gospel.

President Gordon B. Hinckley has given the clarion call for our time. In a worldwide satellite address to missionaries and local leaders, he asked for “an infusion of enthusiasm” for missionary work “at every level in the Church” (“Find the Lambs, Feed the Sheep,” Ensign, May 1999, 107). Though missionaries must continue their best efforts to find persons to teach, he declared that the “better way … is through the members of the Church” (105). He asked each of us to give our very best efforts to assisting missionaries in finding persons to teach. He also asked that each stake president and each bishop “accept full responsibility and accountability for the finding and friendshipping of investigators” within their units (107). President Hinckley also invoked the blessings of the Lord upon each of us “in meeting the tremendous challenge that is ours” (108).

Though it has been two and a half years since our president made this plea, most of us have not yet acted effectively upon his challenge.

As I have prayerfully studied President Hinckley’s words and pondered over how we can share the gospel, I have concluded that we need three things to fulfill our prophet’s challenge. First, we need a sincere desire to share the gospel. Second, we need divine assistance. Third, we need to know what to do.

\intalkheader{I. Desire}

As with so many other things, sharing the gospel begins with desire. If we are to become more effective instruments in the hands of the Lord in sharing His gospel, we must sincerely desire to do so. I believe we acquire this desire in two steps.

First, we must have a firm testimony of the truth and importance of the restored gospel of Jesus Christ. This includes the supreme value of God’s plan for His children, the essential position of the Atonement of Jesus Christ in it, and the role of The Church of Jesus Christ in carrying out that plan in mortality.

Second, we must have a love for God and for all of His children. In modern revelation we are told that “love, with an eye single to the glory of God, [qualifies us] for the work” (D\&C 4:5). The early Apostles of this dispensation were told that their love should “abound unto all men” (D\&C 112:11).

From our testimony of the truth and importance of the restored gospel, we understand the value of what we have been given. From our love of God and our fellowmen, we acquire our desire to share that great gift with everyone. The intensity of our desire to share the gospel is a great indicator of the extent of our personal conversion.

The Book of Mormon contains some marvelous examples of the effect of testimony and love. When the sons of Mosiah, who had been “the very vilest of sinners,” acquired their testimony, “they were desirous that salvation should be declared to every creature, for they could not bear that any human soul should perish” (Mosiah 28:3–4). In a later account, their associate, Alma, cried, “O that I were an angel, … that I might go forth and speak with the trump of God, with a voice to shake the earth” and declare “the plan of redemption” to every soul, “that there might not be more sorrow upon all the face of the earth” (Alma 29:1–2).

I like to refer to missionary efforts as sharing the gospel. The word sharing affirms that we have something extraordinarily valuable and desire to give it to others for their benefit and blessing.

The most effective missionaries, member and full-time, always act out of love. I learned this lesson as a young man. I was assigned to visit a less-active member, a successful professional many years older than I. Looking back on my actions, I realize that I had very little loving concern for the man I visited. I acted out of duty, with a desire to report 100 percent on my home teaching. One evening, close to the end of a month, I phoned to ask if my companion and I could come right over and visit him. His chastening reply taught me an unforgettable lesson.

“No, I don’t believe I want you to come over this evening,” he said. “I’m tired. I’ve already dressed for bed. I am reading, and I am just not willing to be interrupted so that you can report 100 percent on your home teaching this month.” That reply still stings me because I knew he had sensed my selfish motivation.

I hope no person we approach with an invitation to hear the message of the restored gospel feels that we are acting out of any reason other than a genuine love for them and an unselfish desire to share something we know to be precious.

If we lack this love for others, we should pray for it. The prophet Mormon’s writings about “the pure love of Christ” teach us to “pray unto the Father with all the energy of heart, that ye may be filled with this love, which he hath bestowed upon all who are true followers of his Son, Jesus Christ” (Moro. 7:47–48).

\intalkheader{II. Divine Assistance/Right Timing}

We also need divine assistance to guide us in sharing the gospel. Just as our desires must be pure and rooted in testimony and love, our actions must be directed by the Lord. It is His work, not ours, and it must be done in His way and on His timing, not ours. Otherwise, our efforts may be fated to frustration and failure.

All of us have family members or friends who need the gospel but are not now interested. To be effective, our efforts with them must be directed by the Lord so that we act in the way and at the time when they will be most receptive. We must pray for the Lord’s help and directions so we can be instruments in His hands for one who is now ready—one He would have us help today. Then, we must be alert to hear and heed the promptings of His Spirit in how we proceed.

Those promptings will come. We know from countless personal testimonies that in His own way and His own time the Lord is preparing persons to accept His gospel. Such persons are searching, and when we are seeking to identify them the Lord will answer their prayers through answering ours. He will prompt and guide those who desire and who sincerely seek guidance in how, where, when, and with whom to share His gospel. In this way, God grants unto us according to our desires (see Alma 29:4; D\&C 6:8).

In modern revelation, the Lord’s prophet has told us that “there are many yet on the earth among all sects, parties, and denominations, who are blinded … and who are only kept from the truth because they know not where to find it” (D\&C 123:12). When we are standing as “witnesses of God at all times and in all things” (Mosiah 18:9), the Lord will open ways for us to find and have appropriate communications with those who are seeking. This will come when we seek direction and when we act out of a sincere and Christlike love for others.

The Lord loves all of His children. He desires that all have the fulness of His truth and the abundance of His blessings. He knows when they are ready, and He wants us to hear and heed His directions on sharing His gospel. When we do so, those who are prepared will respond to the message of Him who said, “My sheep hear my voice … and they follow me” (John 10:27).

\intalkheader{III. How to Do It}

When we have a sincere desire to share the gospel with others, and when we have sought divine assistance in our efforts, what should we do? How do we proceed? We begin by beginning. We should not wait for a further invitation from heaven. Revelation comes most often when we are on the move.

The Lord has given us this instruction as to who and how: “And let your preaching be … every man to his neighbor, in mildness and in meekness” (D\&C 38:41). “Neighbors,” of course, means not only those who live beside us and other friends and associates. When He was asked, “Who is my neighbour?” the Savior told of a Samaritan who recognized a neighbor on the road to Jericho (see Luke 10:25–37). Thus, our neighbors also include those we encounter in our daily travels.

We should pray, as Alma of old, for the Lord to give us “power and wisdom that we may bring” our associates to the Lord (Alma 31:35). We also pray for the welfare of their souls (see Alma 6:6).

We must be sure we act out of love and not in any attempt to gain personal recognition or advantage. The warning against those who use Church position to gratify their pride or vain ambition (see D\&C 121:37) surely applies to our efforts to share the gospel.

The need to act out of love also warns us against manipulation, real or perceived. People who do not share our belief can be repelled when they hear us refer to something as a “missionary tool.” A “tool” is something used to manipulate an inanimate object. If we talk about something as a “missionary tool,” we can convey the impression that we want to manipulate someone. That impression is entirely contrary to the unselfish, sharing spirit of our missionary service.

In his great message President Hinckley declares that “opportunities for sharing the gospel are everywhere” (Ensign, May 1999, 106). He mentions many things we can do. We should live so that what he called “the tremendous power of the example of a member of the Church” (104) will influence those around us. “The most effective tract we will carry,” he said, “will be the goodness of our own lives and example” (107). We must be sincerely friendly to all.

President Hinckley reminded us that we can “leave a piece of Church literature” (106) with those with whom we come in contact. We can offer our homes “to carry on this missionary service” (105). The missionaries “may appropriately ask the members for referrals” (107), and when they do, we should respond.

In summary, President Hinckley said every member of the Church can “work constantly at the task of finding and encouraging investigators” (107).

There are other things we can do, especially as we act upon the prophet Mormon’s great statement, “I fear not what man can do; for perfect love casteth out all fear” (Moro. 8:16; also see 1 Jn. 4:18). We can invite friends to Church meetings or Church-related activities. We can make appreciative references to our Church and the effect of its teachings and ask persons if they would like to know more.

Even easier, we can carry a packet of these attractive pass-along cards and give them to persons—even casual acquaintances—with whom we come in contact in the daily activities of our lives. These cards are an ideal way to invite people to investigate the additional truths we have to share. In a nonintrusive way, they offer something precious, but the gift depends upon the choice and initiative of the potential recipient. In our experience, a significant fraction of those who telephone for the offered gift choose to have it delivered by those who can tell them more.

The Church has just announced another way to share the gospel, worldwide, on the Internet. In its potential, this new initiative is as exciting as the publishing of written tracts in the 19th century and our use of radio, television, and film in the 20th. The Church has activated a new Internet site to which we may refer persons interested in obtaining information about the Church and its doctrine and how they can find a place to worship with us. Its address is www.mormon.org. For missionaries, the value and use of this new resource will emerge with experience. For members of the Church, it will help us answer the questions of friends directly or by referring them to the site. It will also allow us to send our friends electronic greeting cards that include gospel messages and invitations.

\intalkheader{IV. Conclusion}

We have been asked to redouble our efforts and our effectiveness in sharing the gospel to accomplish the Lord’s purposes in this great work. Until we do so, these wonderful full-time missionaries—our sons and daughters and our noble associates in the Lord’s work—will remain underused in their great assignment to teach the restored gospel of Jesus Christ.

We have spoken about loving desire, heavenly guidance, and ways we can proceed with the divine command to share the gospel with our neighbors. The gospel of Jesus Christ is the brightest light and the only hope for this darkened world. “Wherefore,” as Nephi teaches, we “must press forward with a steadfastness in Christ, having a perfect brightness of hope, and a love of God and of all men” (2 Ne. 31:20).

I testify of Jesus Christ, our Savior, and of His desire that we join wholeheartedly in this, His work, in the name of Jesus Christ, amen.

\newpage
\settalk{The Gospel in Our Lives}
\setspeaker{Dallin H. Oaks}
\settalkdate{April 2002}
\talkheading{The Gospel in Our Lives}{Dallin H. Oaks}

Some years ago I enjoyed a newspaper cartoon that showed a clergyman in conversation with a hippie-dressed couple astride a motorcycle. “We are church goers,” one was saying to the clergyman. “We’ve been going for years … we just haven’t got there yet.”

Many of our extended family and friends haven’t yet gotten to church either. They may attend sporadically, but they are not yet enjoying the full blessings of Church participation and service. Others may attend regularly but refrain from commitments and from seeking the personal spiritual rebirth that comes from yielding our hearts to God. Both kinds are missing some unique blessings in this life. And both are in jeopardy of missing the most glorious blessings in the life to come.

Paul taught that the Lord gave prophets and apostles for “the perfecting of the saints, … the work of the ministry, [and] the edifying of the body of Christ” (Eph. 4:12). Persons who are not fully participating in The Church of Jesus Christ of Latter-day Saints and also seeking a personal spiritual conversion are missing out on experiences that are essential under the divinely established great plan of happiness. The teachings and the work of the Church are essential to bring to pass the eternal life of man (see Moses 1:39).

I pray that many within the sound of my voice will have a spiritual witness of the importance of the Church’s mission to edify and exalt the children of God. I pray especially that some who are not yet enjoying the blessings of full Church participation and commitment will seek and obtain that witness and act upon it.

About a decade ago, while I was at a stake conference in the United States, I was introduced to a member who had not participated in the Church for many years. “Why should I return to Church activity?” this member asked me. Considering all the Savior has done for us, I replied it should be easy to offer something in service to Him and our fellowmen. My questioner considered that idea for a moment and then voiced this astonishing reply: “What’s He done for me?”

This astonishing answer moved me to ponder what people expect to receive from Jesus Christ, from His gospel, and from participating in His Church. I thought of some others who said they stopped going to Church because the Church was “not meeting their needs.” Which needs could they be expecting the Church to meet? If persons are simply seeking a satisfying social experience, they might be disappointed in a particular ward or branch and seek other associations. There are satisfying social experiences in many organizations. If they are simply seeking help to learn the gospel, they could pursue that goal through available literature. But are these the principal purposes of the Church? Is this all we are to receive from the gospel of Jesus Christ?

Someone has said that what we get depends on what we seek. Persons who attend church solely in order to get something of a temporal nature may be disappointed. The Apostle Paul wrote disparagingly of persons who “serve not our Lord Jesus Christ, but their own belly” (Rom. 16:18). Persons who attend church in order to give to their fellowmen and serve the Lord will rarely be disappointed. The Savior promised that “he that loseth his life for my sake shall find it” (Matt. 10:39).

The Church gives us opportunities to serve the Lord and our fellowmen. If given in the right way and for the right reasons, that service will reward us beyond anything we have given. Millions serve unselfishly and effectively as officers or teachers in Church organizations, and those who do experience the conversion described by the prophet who pleaded with us to “come unto Christ, and be perfected in him” (Moro. 10:32).

Throughout my life I have been blessed by my membership and participation in The Church of Jesus Christ of Latter-day Saints. It is impossible to describe all the ways the Church has blessed my life and the lives of those I love. But I will give a few examples, in the hope that this will add personal persuasion to the principles described.

Attendance at church each week provides the opportunity to partake of the sacrament, as the Lord has commanded us (see D\&C 59:9). If we act with the right preparation and attitude, partaking of the sacrament renews the cleansing effect of our baptism and qualifies us for the promise that we will always have His Spirit to be with us. A mission of that Spirit, the Holy Ghost, is to testify of the Father and the Son and to lead us into truth (see John 14:26; 2 Ne. 31:18). Testimony and truth, which are essential to our personal conversion, are the choice harvest of this weekly renewing of our covenants. In the day-to-day decisions of my life and in my personal spiritual growth, I have enjoyed the fulfillment of that promise.

I am sorry when any Latter-day Saint does not understand the precious blessing that comes to those who keep the commandment to offer up their sacraments upon each Sabbath day. What is there in life—on the lakes or streams, in places of commercial recreation, or at home reading the Sunday paper—that can provide anything comparable to these blessings? No recreational pleasure can equal the cleansing renewal and the spiritual guidance and growth God has promised those who faithfully partake of the sacrament and honor Him each Sabbath day. I give thanks for the fulfillment of those promises in my life and affirm their availability to all.

As I came to the age of accountability and understood and experienced the effect of personal sin, the teachings of the gospel of Jesus Christ gave me the peace and courage to go forward in the knowledge that my sins could be forgiven and that there is always hope and the possibility of mercy for those who fall short.

As I experienced the death of loved ones, including my father, my mother, and my wife, the comforting revelations of the Holy Ghost gave me the strength to carry on. The Spirit affirms that there is purpose in mortal adversities and gives assurance of the Resurrection and the reality of family relationships sealed for eternity.

Throughout my life I have been blessed by the doctrine and teachings of the gospel of Jesus Christ. As taught in the scriptures and by the leaders and teachers of this Church, the gospel has been a light to my path and the impetus for my temporal and spiritual progress. As Brigham Young taught, the gospel laws “teach men to be truthful, honest, chaste, sober, industrious, frugal and to love and practice every good word and work, … they elevate and ennoble man … [and], if fully obeyed, bring health and strength to the body, clearness to the perceptions, power to the reasoning faculties as well as salvation to the soul.”

Among many blessings I have received from gospel teachings are those promised for keeping the Word of Wisdom. For me these have included health and knowledge and the capacity to “run and not be weary, and … walk and not faint,” and the fulfillment of the promise that “the destroying angel shall pass by them, as the children of Israel, and not slay them” (D\&C 89:21; see also vv. 18–19).

The gospel teaches us to pay our tithes and offerings and assures us of blessings when we do. I testify of the fulfillment of these promises in my life. I have seen the windows of heaven open in my behalf to bestow blessings unnumbered. Among these is the capacity to see the relative unimportance of the property, pride, prominence, and power of this world, in the context of eternity. How grateful I am for the focus and peace that come from a gospel-based understanding of the purpose of life and its relationship to eternity!

From my earliest years and through schooling and marriage to middle age and beyond, the Church has provided me personal associations with the finest people in the world. Teachers and classmates in Sunday School and Primary, in Scouting and other youth activities, in quorum and ward and stake activities have given me the finest possible role models and friends. Of course, our Church does not have a monopoly on good people, but we have a remarkable concentration of them. My associations in the organizations of The Church of Jesus Christ of Latter-day Saints have given me the basis to recognize, appreciate, and enlarge my associations with people of quality in other churches and organizations.

Because my father died before I was eight years old, I had early cause to wonder about the purposes of the Lord in depriving me of a relationship other boys enjoyed and took for granted. As with so many other mortal challenges, the perspective of the gospel of Jesus Christ filled that void. How grateful I am that my brother and sister and I were raised by a widowed mother who used her faith and our parents’ temple marriage to make our departed father a daily presence in our lives. We never had cause to feel that we were without a father. We had a father, but he was away for a season. There are few things more important in this life than knowing your place in mortality and your potential in eternity. Marriages sealed for eternity in a temple of the Lord provide that possibility for every child and for every adult.

Over the years, my active participation in the Church has provided me access to the counsel and inspiration of the leaders of the Church on things I should do as a husband and father and leader in my family. Again and again, in stake and general conferences, in priesthood quorums, and in Sunday School classes, I have been taught and inspired by wonderful and experienced fathers, mothers, and grandparents. I have sought to follow those teachings to improve my participation in those associations that will persist for eternity. To cite only one example, I have been taught the power of a priesthood blessing—not just a blessing of healing, but a blessing of comfort and guidance that a father holding the Melchizedek Priesthood is privileged to give to members of his family. Learning and applying that principle has favored me and my loved ones with the sweetness and closeness that can only come from sensing the significance of the priesthood of God in an eternal family.

I am also grateful for the warnings of the scriptures and Church leaders on things to avoid. By following that counsel I have been able to avoid pitfalls that might otherwise have trapped and enslaved me. Alcohol, tobacco, drugs, pornography, and gambling are but a few examples of dangerous substances and addictive practices we have been warned to avoid. I appeal to all—especially to young people—to hear and heed the words of the men and women God has called as your leaders and teachers. You will be blessed if you refrain from setting your own wisdom or desires ahead of the commandments of your Creator and the warnings of His servants.

The scriptures tell us to take upon us the “whole armor” of God that we “may be able to withstand the evil day.” They promise that the “breastplate of righteousness” and “the shield of faith” will “quench all the fiery darts of the wicked” (D\&C 27:15–17). I urge you to obey those teachings and lay claim on those blessings. They include the personal spiritual conversion—the “mighty change … in our hearts” (Mosiah 5:2)—that helps us become what our Heavenly Father desires us to become.

The leaders of this Church say, as the Savior said: “My doctrine is not mine, but his that sent me. If any man will do his will, he shall know of the doctrine, whether it be of God, or whether I speak of myself” (John 7:16–17).

Your leaders also say, along with King Benjamin: “I would desire that ye should consider on the blessed and happy state of those that keep the commandments of God. For behold, they are blessed in all things, both temporal and spiritual; and if they hold out faithful to the end they are received into heaven, that thereby they may dwell with God in a state of never-ending happiness” (Mosiah 2:41).

In modern revelation the Lord has declared, “I, the Lord, am bound when ye do what I say; but when ye do not what I say, ye have no promise” (D\&C 82:10).

What has our Savior done for us? He has given us His Atonement, His gospel, and His Church, a sacred combination that gives us the assurance of immortality and the opportunity for eternal life. I testify that this is true, and I testify of God the Father, the author of the plan, and of His Son, Jesus Christ, the Atoning One, who has made it all possible, in the name of Jesus Christ, amen.

\newpage
\settalk{I’ll Go Where You Want Me to Go}
\setspeaker{Dallin H. Oaks}
\settalkdate{October 2002}
\talkheading{I’ll Go Where You Want Me to Go}{Dallin H. Oaks}

My text comes from a hymn that has inspired faithful servants of the Lord for many generations:

\begin{verse}
It may not be on the mountain height\\
Or over the stormy sea,\\
It may not be at the battle’s front\\
My Lord will have need of me.\\
But if, by a still, small voice he calls\\
To paths that I do not know,\\
I’ll answer, dear Lord, with my hand in thine:\\
I’ll go where you want me to go.
\end{verse}

[“I’ll Go Where You Want Me to Go,” Hymns, no. 270]

Penned by a poetess who was not a Latter-day Saint, these words express the commitment of the faithful children of God in all ages.

Abraham, who led Isaac on that heartbreaking journey to Mount Moriah, was faithfully going where the Lord wanted him to go (see Genesis 22). So was David when he stepped out before the hosts of Israel to answer the challenge of the giant Goliath (see 1 Samuel 17). Esther, inspired to save her people, walked a life-threatening path to challenge the king in his inner court (see Esther 4–5). “I’ll go where you want me to go, dear Lord” was the motivation for Lehi to leave Jerusalem (see 1 Nephi 2) and for his son Nephi to return for the precious records (see 1 Nephi 3–4). Hundreds of other scriptural examples can be cited.

All of these faithful souls showed their obedience to the Lord’s direction and their faith in His power and goodness. As Nephi explained, “I will go and do the things which the Lord hath commanded, for I know that the Lord giveth no commandments unto the children of men, save he shall prepare a way for them that they may accomplish the thing which he commandeth them” (1 Nephi 3:7).

All about us, and in our memories of earlier times, we have inspiring examples of the submissive, faithful service of Latter-day Saints. One of the best known was that of President J. Reuben Clark. After he had served over 16 years as an extraordinarily influential first counselor, the First Presidency was reorganized and he was called as second counselor. Offering an example of humility and willingness to serve that has influenced generations, he said to the Church: “In the service of the Lord, it is not where you serve but how. In the Church of Jesus Christ of Latter-day Saints, one takes the place to which one is duly called, which place one neither seeks nor declines” (in Conference Report, Apr. 1951, 154).

Just as significant, though less visible, are the millions of members now laboring with similar faith and devotion in the remote corners of the Lord’s vineyard. Our faithful senior missionaries provide the best examples I know.

I recently reviewed the missionary papers of over 50 senior couples. All had already served at least three missions when they submitted their papers for another call. Their homes were everywhere from Australia to Arizona, California to Missouri. Their ages ranged from the 60s and early 70s to the—well, never mind. One couple, who were offering themselves for a seventh mission, had already served on Temple Square, in Alaska, in New Zealand, in Kenya, and in Ghana. They were sent to the Philippines. Scores of similar examples could be cited.

The priesthood leaders’ comments on the papers of these couples are testimonies of service and sacrifice. I quote several:

“Willing to go anyplace, do anything for whatever length of time required.”

“[These] are great examples of Church members who dedicate their lives to the Lord.”

“Will go where the Lord wants [us] to go,” another couple noted. “We pray we will be sent where we are needed.”

Priesthood leader comments on the qualifications of these couples provide a good summary of the work our senior missionaries do so effectively.

“He is great in getting programs running and [in] leadership.”

“Their joy is fullest when they are asked to ‘build’ and develop; therefore an assignment in a developing area of the Church may be appropriate. Willing to serve in whatever capacity called.”

“They will likely be of more value working with [less-actives] and converts rather than in offices.”

“They love the youth and have a gift with them.”

“They feel most effective in and have a fondness for leadership support and fellowshipping work.”

“They have slowed down some physically, but not in spiritual matters or missionary zeal.”

“He is a true missionary. His first name is Nephi, and he follows his namesake. She is a tremendous lady, has always been a great example. Will do great wherever called. This is their fifth mission.” (They had previously served in Guam, Nigeria, Vietnam, Pakistan, Singapore, and Malaysia. Giving them some respite from those arduous paths, the Lord’s servants called that couple to serve in the Nauvoo Temple.)

Another couple spoke for all these heroes and heroines when they wrote: “Will go anywhere and do what is asked. It is not a sacrifice; it is a privilege.”

These senior missionaries offer a special measure of sacrifice and commitment. So do our mission presidents and temple presidents and their loyal companions. All leave their homes and families to serve full-time for a season. The same is true of the army of young missionaries, who put their lives at home on hold and bid good-bye to family and friends and set forth (usually at their own expense) to serve wherever they are assigned by the Lord, speaking through His servants.

\begin{verse}
I’ll go where you want me to go, dear Lord,\\
Over mountain or plain or sea;\\
I’ll say what you want me to say, dear Lord;\\
I’ll be what you want me to be.
\end{verse}

[Hymns, no. 270]

Millions of others serve from their homes on a Church-service basis. So it is with the 26,000 bishoprics and branch presidencies and the faithful presidencies of the quorums and Relief Society, Primary, and Young Women who serve with them and under their direction. So it is with millions of others—faithful teachers in wards, branches, stakes, and districts. And think of the hundreds of thousands of home teachers and visiting teachers who fulfill the Lord’s command to “watch over the church always, and be with and strengthen them” (D\&C 20:53). All of these can join in this inspired verse:

\begin{verse}
Perhaps today there are loving words\\
Which Jesus would have me speak;\\
There may be now in the paths of sin\\
Some wand’rer whom I should seek.\\
O Savior, if thou wilt be my guide,\\
Tho dark and rugged the way,\\
My voice shall echo the message sweet:\\
I’ll say what you want me to say.
\end{verse}

[Hymns, no. 270]

As the prophet-king Benjamin taught, “When [we] are in the service of [our] fellow beings [we] are only in the service of [our] God” (Mosiah 2:17). He also cautioned us to “see that all these things are done in wisdom and order; for it is not requisite that a man should run faster than he has strength” (Mosiah 4:27).

The gospel of Jesus Christ challenges us to become converted. It teaches us what we should do, and it provides us opportunities to become what our Heavenly Father desires us to become. The full measure of this conversion to men and women of God happens best through our labors in His vineyard.

We have a great tradition of unselfish service in The Church of Jesus Christ of Latter-day Saints. Indeed, one of the distinguishing characteristics of this Church is the fact that we have no paid or professional clergy in our thousands of local congregations and in the regional stakes, districts, and missions that oversee them. As an essential part of God’s plan for His children, the leadership and work of His Church is provided by His children who give their time freely for the service of God and their fellowmen. They obey the Lord’s command to love Him and to serve Him (see John 14:15; D\&C 20:19; 42:29; 59:5). This is the way men and women prepare for the ultimate blessing of eternal life.

Still, there is room for improvement in the commitment of some. When I ask stake presidents for suggestions on subjects I should treat at stake conferences, I often hear about members who refuse Church callings or accept callings and fail to fulfill their responsibilities. Some are not committed and faithful. It has always been so. But this is not without consequence.

The Savior spoke of the contrast between the faithful and the unfaithful in three great parables recorded in the 25th chapter of Matthew. Half of the invited guests were excluded from the wedding feast because they were unprepared when the bridegroom came (see Matthew 25:1–13). The unprofitable servants who failed to employ the talents they were given by the Master were not allowed to enter into the joy of the Lord (see Matthew 25:14–30). And when the Lord came in His glory, He separated the sheep, who had served Him and their fellowmen, from the goats, who had not. Only those who had “done it unto one of the least of these my brethren” (Matthew 25:40) were set on His right hand to inherit the kingdom prepared from the foundation of the world (see Matthew 25:31–46).

My brothers and sisters, if you are delinquent in commitment, please consider who it is you are refusing or neglecting to serve when you decline a calling or when you accept, promise, and fail to fulfill. I pray that each of us will follow this inspired declaration:

\begin{verse}
There’s surely somewhere a lowly place\\
In earth’s harvest fields so wide\\
Where I may labor through life’s short day\\
For Jesus, the Crucified.
\end{verse}

[Hymns, no. 270]

Jesus showed the way. Even though He shrank from the bitter path that led through Gethsemane and Calvary (see D\&C 19:18), He submissively said to the Father, “Nevertheless not my will, but thine, be done” (Luke 22:42).

Earlier He taught:

“If any man will come after me, let him deny himself, and take up his cross, and follow me.

“For whosoever will save his life shall lose it: and whosoever will lose his life for my sake shall find it.

“For what is a man profited, if he shall gain the whole world, and lose his own soul? or what shall a man give in exchange for his soul?” (Matthew 16:24–26).

We need to remember the purpose of our service to one another. If it were only to accomplish some part of His work, God could dispatch “legions of angels,” as Jesus taught on another occasion (see Matthew 26:53). But that would not achieve the purpose of the service He has prescribed. We serve God and our fellowmen in order to become the kind of children who can return to live with our heavenly parents.

\begin{verse}
So trusting my all to thy tender care,\\
And knowing thou lovest me,\\
I’ll do thy will with a heart sincere:\\
I’ll be what you want me to be.
\end{verse}

[Hymns, no. 270]

Almost a decade ago, I read a letter from a returned missionary who described this process in his life. He had written to thank those who direct missionary work “for daring to send me where the Lord required rather than where I had deemed appropriate.” He had come, he said, “from a background of proud, competitive intellectualism.” Before his mission he was a student at a prestigious university in the eastern United States. Quote:

“I guess out of a sense of obligation and inertia, I filled out my [missionary] papers and sent them in, extremely careful to mark the column indicating greatest desire to serve abroad and in a foreign language. I was careful to make it apparent that I was an accomplished student of Russian and fully capable of spending two years among the Russian people. Confident that no committee could resist such qualifications, I rested confident that I would enjoy a wonderfully mind-expanding cultural adventure.”

He was shocked to receive a call to serve in a mission in the United States. He didn’t know anything about the state where he would serve, except that it was in his own country speaking English rather than abroad speaking the language he had learned, and, as he said, “The people I would work with would likely be academic incompetents.” He continued, “I almost refused to accept the call, feeling that I would be more fulfilled by enlisting in the Peace Corps or something else.”

Fortunately, this proud young man found the courage and faith to accept the call and to follow the direction and counsel of his fine mission president. Then the miracle of spiritual growth began. He described it thus:

“As I began to serve among the uneducated people of [this state], I struggled mightily for several months, but gradually the sweet workings of the Spirit began to tear down the walls of pride and disbelief that had wrapped themselves so tightly around my soul. The miracle of a conversion to Christ began. The sense of the reality of God and the eternal brotherhood of all men came more and more powerfully to my troubled mind.”

It was not easy, he admitted, but with the influence of his great mission president and with his growing love for the people he served, it was possible, and it occurred.

“My desire to love and serve these people who in the ultimate scale were at least my peers, almost definitely my superiors, waxed stronger and stronger. I learned humility for the first time in my life; I learned what it means to make our valuations of others [without relying on the] irrelevant details of life. I began to feel swelling within my heart a love of the spirits that came here to earth with me” (letter to General Authorities, Feb. 1994).

Such is the miracle of service. As the poetess wrote:

\begin{verse}
But if, by a still, small voice he calls\\
To paths that I do not know,\\
I’ll answer, dear Lord, with my hand in thine:\\
I’ll go where you want me to go.
\end{verse}

[Hymns, no. 270]

I testify of Jesus Christ, who beckons us to His path and His service, and pray that we will have the faith and commitment to follow and the power to be what He wants us to be, in the name of Jesus Christ, amen.

\newpage
\settalk{Give Thanks in All Things}
\setspeaker{Dallin H. Oaks}
\settalkdate{April 2003}
\talkheading{Give Thanks in All Things}{Dallin H. Oaks}

In one of the times of spiritual and temporal adversity recorded in the Book of Mormon, when the people of God were “suffering all manner of afflictions,” the Lord commanded them “to give thanks in all things” (Mosiah 26:38, 39). I wish to apply that teaching to our time.

\intalkheader{I.}

The children of God have always been commanded to give thanks. There are examples throughout the Old and New Testaments. The Apostle Paul wrote, “In every thing give thanks: for this is the will of God in Christ Jesus concerning you” (1 Thessalonians 5:18). The prophet Alma taught, “When thou risest in the morning let thy heart be full of thanks unto God” (Alma 37:37). And in modern revelation the Lord declared that “he who receiveth all things with thankfulness shall be made glorious; and the things of this earth shall be added unto him, even an hundred fold” (D\&C 78:19).

\intalkheader{II.}

We have so much for which to give thanks. First and foremost, we are thankful for our Savior Jesus Christ. Under the plan of the Father, He created the world. Through His prophets, He revealed the plan of salvation with its accompanying commandments and ordinances. He came into mortality to teach and show us the way. He suffered and paid the price for our sins if we would repent. He gave up His life, and He conquered death and rose from the grave that we all will live again. He is the Light and Life of the World. As King Benjamin taught, if we “should render all the thanks and praise which [our] whole soul has power to possess, to that God who has created [us], and has kept and preserved [us], and … should serve him with all [our] whole souls yet [we] would be unprofitable servants” (Mosiah 2:20–21).

We give thanks for the revealed truths that provide a standard against which to measure all things. As the Bible teaches, the Lord gave us apostles and prophets “for the perfecting of the saints” (see Ephesians 4:11–12). We use the revealed truth they give us, “that we henceforth be no more children, tossed to and fro, and carried about with every wind of doctrine, by the sleight of men, and cunning craftiness, whereby they lie in wait to deceive” (Ephesians 4:14). Those who view every calamity and measure every new assertion or discovery against the standard of revealed truth need not be “tossed to and fro” but can be steady and at peace. God is in His heavens, and His promises are sure. “Be not troubled,” He has said to us concerning the destructions that will precede the end of the world, “for, when all these things shall come to pass, ye may know that the promises which have been made unto you shall be fulfilled” (D\&C 45:35). What an anchor to the soul in these troubled times!

We give thanks for commandments. They are directions away from pitfalls, and they are invitations to blessings. Commandments mark the path and show us the way to happiness in this life and eternal life in the world to come.

\intalkheader{III.}

In the past eight months in the Philippines, I have heard many testimonies of the blessings of the gospel. Speaking at the dedication of his ward chapel, a Filipino bishop expressed his gratitude for the gospel message that came into his life about 10 years ago. He described how it rescued him from a life of selfishness, excess, and abusive practices and made him a good husband and father. He testified of the blessings that had come to him from paying his tithing.

Speaking at a leadership meeting, a counselor in a stake presidency who is a lawyer and community leader said: “I can declare to the whole world without mental reservation that the greatest thing that ever happened in my life is my becoming a member of The Church of Jesus Christ of Latter-day Saints. It … made a great difference in my life and that of my family, even if I feel there is more that I should learn and apply in my life. The Church is indeed a marvelous work and a wonder.”

You do not have to travel to the Philippines to experience such testimonies. They are evident wherever the gospel message is received and lived. But Sister Oaks and I are profoundly grateful for our opportunity to live and serve in the Philippines, where we have met thousands of wonderful members in new surroundings and seen the gospel in a new light.

In the developing world we learn the importance of establishing the Church—not just teaching and baptizing, but retaining the new members by loving, by calling and ordaining, and by nourishing with the good word of God. We have learned the importance of challenging members to abandon cultural traditions that are contrary to gospel commandments and covenants and to live so that they and their posterity “are no more strangers and foreigners, but fellowcitizens with the saints, and of the household of God; … built upon the foundation of the apostles and prophets, Jesus Christ himself being the chief corner stone” (Ephesians 2:19–20).

People who do this become part of the worldwide gospel culture of commandments and covenants and ordinances and blessings. Such people experience “a mighty change” in their hearts, “that [they] have no more disposition to do evil, but to do good continually” (Mosiah 5:2). The image of God is “engraven upon [their] countenances” (Alma 5:19). Such followers of Christ are found in every land where the gospel and the Church have been established. We have many of them in the Philippines, and we are working to encourage more of them. We do this by growing from our centers of strength, concentrating our teaching where there are sufficiently large groups of committed members to provide the friendshipping, the teachings, the role models, and the needed assistance to the struggling newly born members who are just learning what the gospel asks of us and gives to us.

\intalkheader{IV.}

The revelations, for which we are grateful, show that we should even give thanks for our afflictions because they turn our hearts to God and give us opportunities to prepare for what God would have us become. The Lord taught the prophet Moroni, “I give unto men weakness that they may be humble,” and then promised that “if they humble themselves … and have faith in me, then will I make weak things become strong unto them” (Ether 12:27). In the midst of the persecutions the Latter-day Saints were suffering in Missouri, the Lord gave a similar teaching and promise: “Verily I say unto you my friends, fear not, let your hearts be comforted; yea, rejoice evermore, and in everything give thanks; … and all things wherewith you have been afflicted shall work together for your good” (D\&C 98:1, 3). And to Joseph Smith in the afflictions of Liberty Jail, the Lord said, “Know thou, my son, that all these things shall give thee experience, and shall be for thy good” (D\&C 122:7). Brigham Young understood. Said he, “There is not a single condition of life [or] one hour’s experience but what is beneficial to all those who make it their study, and aim to improve upon the experience they gain” (Deseret News Weekly, 9 July 1862, 1).

As someone has said, there is a big difference between 20 years’ experience and one year’s experience repeated 20 times. If we understand the Lord’s teachings and promises, we will learn and grow from our adversities.

Many of the inspired teachings of our modern prophets are compiled in Teachings of Presidents of the Church, our course of study for Melchizedek Priesthood and Relief Society. The timeless doctrines and principles included in these books are fountains of divine wisdom and guidance. Wise teachers in wards and branches will not substitute their own subjects and wisdom but focus on these inspired teachings and their application to current circumstances and challenges.

For example, in the current volume we read these words of President John Taylor on the subject of gratitude for suffering: “We have learned many things through suffering. We call it suffering. I call it a school of experience. … I have never looked at these things in any other light than trials for the purpose of purifying the Saints of God that they may be, as the scriptures say, as gold that has been seven times purified by the fire” (Teachings of Presidents of the Church: John Taylor [2001], 203). Pioneers like President John Taylor, who witnessed the murder of their prophet and experienced prolonged persecution and incredible hardships for their faith, praised God and thanked Him. Through their challenges and the courageous and inspired actions they took to meet them, they grew in faith and in spiritual stature. Through their afflictions they became what God desired them to become, and they laid the foundation of the great work that blesses our lives today.

Like the pioneers, we should thank God for our adversities and pray for guidance in meeting them. Through that attitude and through our faith and obedience, we will realize the promises God has given us. It is all part of the plan.

I love the musical and motion picture Fiddler on the Roof. There a wonderful Jewish father sings “If I Were a Rich Man.” His memorable prayer concludes with this pleading question:

\begin{verse}
Lord, who made the lion and the lamb,\\
You decreed I should be what I am;\\
Would it spoil some vast eternal plan,\\
If I were a wealthy man?
\end{verse}

[Lyrics by Sheldon Harnick (1964)]

Yes, Tevye, it might. Let us give thanks for what we are and for the circumstances God has given us for our personal journey through mortality.

In ancient times the prophet Lehi taught this truth to his son Jacob:

“In thy childhood thou hast suffered afflictions and much sorrow, because of the rudeness of thy brethren.

“Nevertheless, Jacob, my first-born in the wilderness, thou knowest the greatness of God; and he shall consecrate thine afflictions for thy gain” (2 Nephi 2:1–2).

My mother loved that scripture and lived its principle. The greatest affliction of her life was the death of her husband, our father, after only 11 years of marriage. This changed her life and imposed great hardships as she proceeded to earn a living and raise her three little children alone. Nevertheless, I often heard her say that the Lord consecrated that affliction for her gain because her husband’s death compelled her to develop her talents and serve and become something that she could never have become without that seeming tragedy. Our mother was a spiritual giant, strong and fully worthy of the loving tribute her three children inscribed on her headstone: “Her Faith Strengthened All.”

The blessings of adversity extend to others. I know it was a blessing to be raised by a widowed mother whose children had to learn how to work, early and hard. I know that relative poverty and hard work are not greater adversities than affluence and abundant free time. I also know that strength is forged in adversity and that faith is developed in a setting where we cannot see ahead.

\intalkheader{V.}

When we give thanks in all things, we see hardships and adversities in the context of the purpose of life. We are sent here to be tested. There must be opposition in all things. We are meant to learn and grow through that opposition, through meeting our challenges, and through teaching others to do the same. Our beloved colleague Elder Neal A. Maxwell has given us a noble example of this. His courage, his submissive attitude in accepting his affliction with cancer, and his stalwart continued service have ministered comfort to thousands and taught eternal principles to millions. His example shows that the Lord will not only consecrate our afflictions for our gain, but He will use them to bless the lives of countless others.

Jesus taught this lesson when He and His disciples met a man who was born blind. “Who did sin, this man, or his parents, that he was born blind?” the disciples asked. “Neither,” Jesus answered. The man was born blind “that the works of God should be made manifest in him” (John 9:2, 3).

If we see life through the lens of spirituality, we can see many examples of the works of God being furthered through the adversities of His children. I often visit the American War Memorial in Manila. To me, that is a sacred place. It is the burial place of over 17,000 soldiers, sailors, and airmen who lost their lives in World War II battles in the Pacific. The memorial also honors over 36,000 other servicemen who also lost their lives but whose bodies were never recovered. As I walk past the beautiful walls where are inscribed their names and the state of their origin, I see many that I suppose were faithful Latter-day Saints.

Reflecting on the wartime deaths of so many worthy and wonderful members and how much suffering this has caused their loved ones, I have thought of President Joseph F. Smith’s great vision recorded in the 138th section of the Doctrine and Covenants. He saw “an innumerable company” of righteous spirits “who had been faithful in the testimony of Jesus while they lived in mortality” (D\&C 138:12). They were organized and appointed as messengers, “clothed with power and authority, and commissioned … to go forth and carry the light of the gospel to them that were in darkness, … and thus was the gospel preached to the dead” (D\&C 138:30). Reflecting on this revelation and remembering the millions who have fallen in war, I rejoice in the Lord’s plan, in which the adversity of the deaths of many righteous individuals is turned into the blessing of righteous messengers to preach the gospel to their countless comrades-in-arms.

When we understand this principle, that God offers us opportunities for blessings and blesses us through our own adversities and the adversities of others, we can understand why He has commanded us again and again to “thank the Lord thy God in all things” (D\&C 59:7).

I pray that we will be blessed to understand the truth and purpose of the doctrines and commandments I have described and that we will be faithful enough and strong enough to give thanks in all things. I testify of Jesus Christ, our Savior and Redeemer and Creator, for whom we give thanks, in the name of Jesus Christ, amen.

\newpage
\settalk{Repentance and Change}
\setspeaker{Dallin H. Oaks}
\settalkdate{October 2003}
\talkheading{Repentance and Change}{Dallin H. Oaks}

I bring you greetings from the Philippines Area, with its 520,000 members in 80 stakes and 80 member districts and its 2,200 missionaries in 13 missions. We are progressing against the challenges the Church encounters where it is not yet fully established.

In these developing areas, we rely heavily on senior missionary couples. I stress this because there are many within the sound of my voice who need to know how much their service is appreciated, and there are others we pray will decide to be available for this vital service.

\intalkheader{I.}

My introduction is something said in my presence by one of these valiant missionaries. “As I look back on my life,” he said, “I can hardly imagine a barefoot surfer from Hawaii completing his third mission. But when I felt the warm embrace of the Savior, I wanted to serve Him, and I changed.” Yes he did! Stanley Y. Q. Ho told me that until he was 30 years old he did nothing but “hang around the beaches at Waikiki.” Then he found the gospel, he married a Latter-day Saint girl, and he changed. Since then he has fulfilled many callings, including bishop and stake president. Now, Elder Ho and his beloved Momi, who is responsible for so many of the changes in his life, have served three full-time missions.

For another example, I turn to the Gospel of Luke:

“And Jesus entered and passed through Jericho.

“And, behold, there was a man named Zacchaeus, which was the chief among the publicans, and he was rich.

“And he sought to see Jesus who he was; and could not for the press, because he was little of stature.

“And he ran before, and climbed up into a sycomore tree to see him: for he was to pass that way.

“And when Jesus came to the place, he looked up, and saw him, and said unto him, Zacchaeus, make haste, and come down; for to day I must abide at thy house.

“And he made haste, and came down, and received him joyfully” (Luke 19:1–6).

Here the Gospel records that Jesus’ followers “murmured” because of His going to the house of a sinner (v. 7). But that did not matter to Jesus. His gospel is for all who will forsake their old ways and make the changes they need to be saved in the kingdom of God.

Now back to the account of the man who opened his house and his heart to the Lord:

“And Zacchaeus stood, and said unto the Lord; Behold, Lord, the half of my goods I give to the poor; and if I have taken any thing from any man by false accusation, I restore him fourfold.

“And Jesus said unto him, This day is salvation come to this house. …

“For the Son of man is come to seek and to save that which was lost” (vv. 8–10).

Zacchaeus of Jericho and Stanley of Hawaii stand for all of us. They are examples of what we pray will be experienced by all of us who decide to receive the Lord “joyfully” and follow where He leads.

\intalkheader{II.}

The gospel of Jesus Christ challenges us to change. “Repent” is its most frequent message, and repenting means giving up all of our practices—personal, family, ethnic, and national—that are contrary to the commandments of God. The purpose of the gospel is to transform common creatures into celestial citizens, and that requires change.

John the Baptist preached repentance. His listeners came from different groups, and he declared the changes each must make to “bring forth … fruits worthy of repentance” (Luke 3:8). Publicans, soldiers, and ordinary people—each had traditions that had to yield to the process of repentance.

The teachings of Jesus also challenged the traditions of different groups. When the scribes and Pharisees complained that His disciples “transgress[ed] the tradition of the elders” by omitting the ritual washings, Jesus replied that the scribes and Pharisees “transgress[ed] the commandment of God by [their] tradition” (Matthew 15:2–3). He described how they had “made the commandment of God of none effect by [their] tradition” (v. 6). “Hypocrites” is what He called those whose adherence to their traditions kept them from keeping the commandments of God (v. 7).

Again, in modern revelation the Lord declares that the “wicked one” takes the innocent children of God away from light and truth “through disobedience … and because of the tradition of their fathers” (D\&C 93:39).

The traditions or culture or way of life of a people inevitably include some practices that must be changed by those who wish to qualify for God’s choicest blessings.

Chastity is an example. “Thou shalt not commit adultery,” the Lord commanded from Sinai (Exodus 20:14) and repeated in modern revelation (D\&C 42:24; see also D\&C 59:6). “Flee fornication,” the New Testament commands (1 Corinthians 6:18; see also Galatians 5:19; 1 Thessalonians 4:3). Always the prophets of God have condemned whoredoms. Yet these eternal commands have frequently been ignored, opposed, or mocked by powerful traditions in many lands. This is especially visible today, when the movies, magazines, and Internet communications of one nation are instantly shared with many others. Sexual relations out of wedlock are tolerated or advocated by many. So is the rapidly expanding culture of pornography. All who have belonged to these cultures of sin must repent and change if they are to become the people of God, for He has warned that “no unclean thing can enter into his kingdom” (3 Nephi 27:19).

Weekly attendance at church is another example of a commandment contrary to popular traditions. The Lord has commanded us to attend church and “offer up [our] sacraments” on His Sabbath day (see D\&C 59:9). This requires more than passive attendance. We are commanded to participate in worship and in service, and that requires a wrenching change for many non-Christians and even for those Christians who have attended church only as irregular spectators.

The Lord’s command that we abstain from alcohol, tobacco, tea, and coffee (see D\&C 89) also runs counter to the traditions of many. Long-standing addictions or habits are not easily broken, but God’s command is clear, and the promised blessings more than compensate for the challenges of change.

Another example is honesty. Some cultures allow lying, stealing, and other dishonest practices. But dishonesty in any form—whether to appease, to save face, or to get gain—is in direct conflict with gospel commandments and culture. God is a God of truth, and God does not change. We are the ones who must change. And that will be a big change for all whose traditions accustom them to thinking that they can lie a little, cheat a little, or engage in deceit whenever it brings personal advantage and is not likely to be detected.

A less serious worldly tradition that conflicts with gospel culture is the idea of upward or downward movement in positions. In the world, we refer to the up or down of promotions or reductions. But there is no up or down in Church positions. We just move around. A bishop released by proper authority and called to teach in Primary does not move down. He moves forward as he accepts his release with gratitude and fulfills the duties of a new calling—even one far less visible.

I saw a memorable example of this a few months ago in the Philippines. I visited a ward in the Pasig stake, near Manila. There I met Augusto Lim, whom I had known in earlier years as a stake president, a mission president, a General Authority, and president of the Manila temple. Now I saw him serving humbly and gratefully in his ward bishopric, second counselor to a man much younger and much less experienced. From temple president to second counselor in a ward bishopric is a beautiful example of the gospel culture in action.

In these examples I am not contrasting the culture or traditions of one part of the world with another. I am contrasting the Lord’s way with the world’s way—the culture of the gospel of Jesus Christ with the culture or traditions of every nation or people. No group has a monopoly on virtue or an immunity from the commandment to change. Jesus and His Apostles did not attempt to make Gentiles into Jews (see Romans 2:11; Galatians 2:11–16; 3:1–29; 5:1–6; 6:15). They taught Gentiles and Jews, attempting to make each of them into followers of Christ.

Similarly, the present-day servants of the Lord do not attempt to make Filipinos or Asians or Africans into Americans. The Savior invites all to come unto Him (see 2 Nephi 26:33; D\&C 43:20), and His servants seek to persuade all—including Americans—to become Latter-day Saints. We say to all, give up your traditions and cultural practices that are contrary to the commandments of God and the culture of His gospel, and join with His people in building the kingdom of God. If we cease to walk in darkness, the Apostle John taught, “we walk in the light, … we have fellowship one with another, and the blood of Jesus Christ his Son cleanseth us from all sin” (1 John 1:7).

\intalkheader{III.}

There is a unique gospel culture, a set of values and expectations and practices common to all members of The Church of Jesus Christ of Latter-day Saints. This gospel way of life comes from the plan of salvation, the commandments of God, and the teachings of the living prophets. It is given expression in the way we raise our families and live our individual lives. The principles stated in the family proclamation are a beautiful expression of our gospel culture.

Those who are baptized in the Church of Jesus Christ make covenants. In modern revelation the Lord declared, “When men are called unto mine everlasting gospel, and covenant with an everlasting covenant, they are accounted as the salt of the earth and the savor of men” (D\&C 101:39). To perform our covenant duty as the salt of the earth, we must be different from those around us.

As Jesus taught: “I give unto you to be the salt of the earth; but if the salt shall lose its savor wherewith shall the earth be salted? The salt shall be thenceforth good for nothing, but to be cast out and to be trodden under foot of men” (3 Nephi 12:13; see also Matthew 5:13; D\&C 101:40).

This requires us to make some changes from our family culture, our ethnic culture, or our national culture. We must change all elements of our behavior that are in conflict with gospel commandments, covenants, and culture.

The gospel plan is based on individual responsibility. Our article of faith states the eternal truth “that men will be punished for their own sins, and not for Adam’s transgression” (Articles of Faith 1:2). This requirement of individual responsibility, which has many expressions in our doctrine, is in sharp contrast to Satan’s plan to “redeem all mankind, that one soul shall not be lost” (Moses 4:1). The plan of the Father and the Savior is based on individual choice and individual effort.

The doctrine and practice of personal responsibility and personal effort collide with individual traditions and local cultures in many lands. We live in a world where there are large differences in income and material possessions and where there are many public and private efforts to narrow these differences. The followers of the Savior are commanded to give to the poor, and many do. But some gifts have promoted a culture of dependency, reducing their recipients’ need for earthly food or shelter but impoverishing them in their eternal need for individual growth. The growth required by the gospel plan only occurs in a culture of individual effort and responsibility. It cannot occur in a culture of dependency. Whatever causes us to be dependent on someone else for decisions or resources we could provide for ourselves weakens us spiritually and retards our growth toward what the gospel plan intends us to be.

The gospel raises people out of poverty and dependency, but only when gospel culture, including the faithful payment of tithing even by the very poor, prevails over the traditions and cultures of dependency. That is the lesson to be learned from the children of Israel, who came out of hundreds of years of slavery in Egypt and followed a prophet into their own land and became a mighty people. That lesson can also be learned from the Mormon pioneers, who never used their persecutions or poverty as an excuse but went forward in faith, knowing that God would bless them when they kept His commandments, which He did.

The changes we must make to become part of the gospel culture require prolonged and sometimes painful effort, and our differences must be visible. As the “salt of the earth,” we are also the “light of the world,” and our light must not be hidden (see Matthew 5:13–16). The Apostle John warned that this will cause the world to hate us (see 1 John 3:13). That is why those who have made the covenant to change have a sacred duty to love and help one another. That encouragement must be extended to every soul who struggles to come out of the culture of the world and into the culture of the gospel of Jesus Christ. The Apostle John concluded, “Let us not love in word, neither in tongue; but in deed and in truth” (1 John 3:18).

No one shows love for their fellowmen more impressively than the noble men and women of this Church who leave comfortable homes and surroundings to serve as couple missionaries. They provide the most authentic and the most valuable assistance to those who are struggling to change. God bless our couple missionaries!

\intalkheader{IV.}

Jesus commanded us to love one another, and we show that love by the way we serve one another. We are also commanded to love God, and we show that love by continually repenting and by keeping His commandments (see John 14:15). And repentance means more than giving up our sins. In its broadest meaning it requires change, giving up all of our traditions that are contrary to the commandments of God. As we become full participants in the culture of the gospel of Jesus Christ, we become “fellowcitizens with the saints, and of the household of God” (Ephesians 2:19).

I testify that this is what our Lord and Savior would have us do so that we may become what His gospel intends us to be, in the name of Jesus Christ, amen.

\newpage
\settalk{Preparation for the Second Coming}
\setspeaker{Dallin H. Oaks}
\settalkdate{April 2004}
\talkheading{Preparation for the Second Coming}{Dallin H. Oaks}

In modern revelation we have the promise that if we are prepared we need not fear (see D\&C 38:30). I was introduced to that principle 60 years ago this summer when I became a Boy Scout and learned the Scout motto: “Be prepared.” Today I have felt prompted to speak of the importance of preparation for a future event of supreme importance to each of us—the Second Coming of the Lord.

The scriptures are rich in references to the Second Coming, an event eagerly awaited by the righteous and dreaded or denied by the wicked. The faithful of all ages have pondered the sequence and meaning of the many events prophesied to precede and follow this hinge point of history.

Four matters are indisputable to Latter-day Saints: (1) The Savior will return to the earth in power and great glory to reign personally during a millennium of righteousness and peace. (2) At the time of His coming there will be a destruction of the wicked and a resurrection of the righteous. (3) No one knows the time of His coming, but (4) the faithful are taught to study the signs of it and to be prepared for it. I wish to speak about the fourth of these great realities: the signs of the Second Coming and what we should do to prepare for it.

\intalkheader{I.}

The Lord has declared, “He that feareth me shall be looking forth for the great day of the Lord to come, even for the signs of the coming of the Son of Man,” signs that will be shown “in the heavens above, and in the earth beneath” (D\&C 45:39–40).

The Savior taught this in the parable of the fig tree whose tender new branches give a sign of the coming of summer. “So likewise,” when the elect shall see the signs of His coming, “they shall know that he is near, even at the doors” (Joseph Smith—Matthew 1:38–39; see also Matthew 24:32–33; D\&C 45:37–38).

Biblical and modern prophecies give many signs of the Second Coming. These include:

(See Matthew 24:5–15; Joseph Smith—Matthew 1:22, 28–32; D\&C 45:26–33.)

In another revelation the Lord declares that some of these signs are His voice calling His people to repentance:

“Hearken, O ye nations of the earth, and hear the words of that God who made you. …

“How oft have I called upon you by the mouth of my servants, and by the ministering of angels, and by mine own voice, and by the voice of thunderings, and by the voice of lightnings, and by the voice of tempests, and by the voice of earthquakes, and great hailstorms, and by the voice of famines and pestilences of every kind, … and would have saved you with an everlasting salvation, but ye would not!” (D\&C 43:23, 25).

These signs of the Second Coming are all around us and seem to be increasing in frequency and intensity. For example, the list of major earthquakes in The World Almanac and Book of Facts, 2004 shows twice as many earthquakes in the decades of the 1980s and 1990s as in the two preceding decades (see pages 189–90). It also shows further sharp increases in the first several years of this century. The list of notable floods and tidal waves and the list of hurricanes, typhoons, and blizzards worldwide show similar increases in recent years (see pages 188–89). Increases by comparison with 50 years ago can be dismissed as changes in reporting criteria, but the accelerating pattern of natural disasters in the last few decades is ominous.

\intalkheader{II.}

Another sign of the times is the gathering of the faithful (see D\&C 133:4). In the early years of this last dispensation, a gathering to Zion involved various locations in the United States: to Kirtland, to Missouri, to Nauvoo, and to the tops of the mountains. Always these were gatherings to prospective temples. With the creation of stakes and the construction of temples in most nations with sizable populations of the faithful, the current commandment is not to gather to one place but to gather in stakes in our own homelands. There the faithful can enjoy the full blessings of eternity in a house of the Lord. There, in their own homelands, they can obey the Lord’s command to enlarge the borders of His people and strengthen her stakes (see D\&C 101:21; 133:9, 14). In this way the stakes of Zion are “for a defense, and for a refuge from the storm, and from wrath when it shall be poured out without mixture upon the whole earth” (D\&C 115:6).

\intalkheader{III.}

While we are powerless to alter the fact of the Second Coming and unable to know its exact time, we can accelerate our own preparation and try to influence the preparation of those around us.

A parable that contains an important and challenging teaching on this subject is the parable of the ten virgins. Of this parable the Lord said, “And at that day, when I shall come in my glory, shall the parable be fulfilled which I spake concerning the ten virgins” (D\&C 45:56).

Given in the 25th chapter of Matthew, this parable contrasts the circumstances of the five foolish and the five wise virgins. All ten were invited to the wedding feast, but only half of them were prepared with oil in their lamps when the bridegroom came. The five who were prepared went into the marriage feast, and the door was shut. The five who had delayed their preparations came late. The door had been closed, and the Lord denied them entrance, saying, “I know you not” (v. 12). “Watch therefore,” the Savior concluded, “for ye know neither the day nor the hour wherein the Son of man cometh” (v. 13).

The arithmetic of this parable is chilling. The ten virgins obviously represent members of Christ’s Church, for all were invited to the wedding feast and all knew what was required to be admitted when the bridegroom came. But only half were ready when he came.

Modern revelation contains this teaching, spoken by the Lord to the early leaders of the Church:

“And after your testimony cometh wrath and indignation upon the people.

“For after your testimony cometh the testimony of earthquakes. …

“And … the testimony of the voice of thunderings, and the voice of lightnings, and the voice of tempests, and the voice of the waves of the sea heaving themselves beyond their bounds.

“And all things shall be in commotion; and surely, men’s hearts shall fail them; for fear shall come upon all people.

“And angels shall fly through the midst of heaven, crying with a loud voice, sounding the trump of God, saying: Prepare ye, prepare ye, O inhabitants of the earth; for the judgment of our God is come. Behold, and lo, the Bridegroom cometh; go ye out to meet him” (D\&C 88:88–92).

\intalkheader{IV.}

Brothers and sisters, as the Book of Mormon teaches, “this life is the time for men to prepare to meet God; … the day of this life is the day for men to perform their labors” (Alma 34:32). Are we preparing?

In His preface to our compilation of modern revelation, the Lord declares, “Prepare ye, prepare ye for that which is to come, for the Lord is nigh” (D\&C 1:12).

The Lord also warned: “Yea, let the cry go forth among all people: Awake and arise and go forth to meet the Bridegroom; behold and lo, the Bridegroom cometh; go ye out to meet him. Prepare yourselves for the great day of the Lord” (D\&C 133:10; see also D\&C 34:6).

Always we are cautioned that we cannot know the day or the hour of His coming. In the 24th chapter of Matthew, Jesus taught:

“Watch therefore: for ye know not what hour your Lord doth come.

“But know this, that if the goodman of the house had known in what watch the thief would come, he would have watched, and would not have suffered his house to be broken up” (Matthew 24:42–43). “But would have been ready” (Joseph Smith—Matthew 1:47).

“Therefore be ye also ready: for in such an hour as ye think not the Son of man cometh” (Matthew 24:44; see also D\&C 51:20).

What if the day of His coming were tomorrow? If we knew that we would meet the Lord tomorrow—through our premature death or through His unexpected coming—what would we do today? What confessions would we make? What practices would we discontinue? What accounts would we settle? What forgivenesses would we extend? What testimonies would we bear?

If we would do those things then, why not now? Why not seek peace while peace can be obtained? If our lamps of preparation are drawn down, let us start immediately to replenish them.

We need to make both temporal and spiritual preparation for the events prophesied at the time of the Second Coming. And the preparation most likely to be neglected is the one less visible and more difficult—the spiritual. A 72-hour kit of temporal supplies may prove valuable for earthly challenges, but, as the foolish virgins learned to their sorrow, a 24-hour kit of spiritual preparation is of greater and more enduring value.

\intalkheader{V.}

We are living in the prophesied time “when peace shall be taken from the earth” (D\&C 1:35), when “all things shall be in commotion” and “men’s hearts shall fail them” (D\&C 88:91). There are many temporal causes of commotion, including wars and natural disasters, but an even greater cause of current commotion is spiritual.

Viewing our surroundings through the lens of faith and with an eternal perspective, we see all around us a fulfillment of the prophecy that “the devil shall have power over his own dominion” (D\&C 1:35). Our hymn describes “the foe in countless numbers, / Marshaled in the ranks of sin” (“Hope of Israel,” Hymns, no. 259), and so it is.

Evil that used to be localized and covered like a boil is now legalized and paraded like a banner. The most fundamental roots and bulwarks of civilization are questioned or attacked. Nations disavow their religious heritage. Marriage and family responsibilities are discarded as impediments to personal indulgence. The movies and magazines and television that shape our attitudes are filled with stories or images that portray the children of God as predatory beasts or, at best, as trivial creations pursuing little more than personal pleasure. And too many of us accept this as entertainment.

The men and women who made epic sacrifices to combat evil regimes in the past were shaped by values that are disappearing from our public teaching. The good, the true, and the beautiful are being replaced by the no-good, the “whatever,” and the valueless fodder of personal whim. Not surprisingly, many of our youth and adults are caught up in pornography, pagan piercing of body parts, self-serving pleasure pursuits, dishonest behavior, revealing attire, foul language, and degrading sexual indulgence.

An increasing number of opinion leaders and followers deny the existence of the God of Abraham, Isaac, and Jacob and revere only the gods of secularism. Many in positions of power and influence deny the right and wrong defined by divine decree. Even among those who profess to believe in right and wrong, there are “them that call evil good, and good evil” (Isaiah 5:20; 2 Nephi 15:20). Many also deny individual responsibility and practice dependence on others, seeking, like the foolish virgins, to live on borrowed substance and borrowed light.

All of this is grievous in the sight of our Heavenly Father, who loves all of His children and forbids every practice that keeps any from returning to His presence.

What is the state of our personal preparation for eternal life? The people of God have always been people of covenant. What is the measure of our compliance with covenants, including the sacred promises we made in the waters of baptism, in receiving the holy priesthood, and in the temples of God? Are we promisers who do not fulfill and believers who do not perform?

Are we following the Lord’s command, “Stand ye in holy places, and be not moved, until the day of the Lord come; for behold, it cometh quickly”? (D\&C 87:8). What are those “holy places”? Surely they include the temple and its covenants faithfully kept. Surely they include a home where children are treasured and parents are respected. Surely the holy places include our posts of duty assigned by priesthood authority, including missions and callings faithfully fulfilled in branches, wards, and stakes.

As the Savior taught in His prophecy of the Second Coming, blessed is the “faithful and wise servant” who is attending to his duty when the Lord comes (see Matthew 24:45–46). As the prophet Nephi taught of that day, “The righteous need not fear” (1 Nephi 22:17; see also 1 Nephi 14:14; D\&C 133:44). And modern revelation promises that “the Lord shall have power over his saints” (D\&C 1:36).

We are surrounded by challenges on all sides (see 2 Corinthians 4:8–9). But with faith in God, we trust the blessings He has promised those who keep His commandments. We have faith in the future, and we are preparing for that future. To borrow a metaphor from the familiar world of athletic competitions, we do not know when this game will end, and we do not know the final score, but we do know that when the game finally ends, our team wins. We will continue to go forward “till the purposes of God shall be accomplished, and the Great Jehovah shall say the work is done” (History of the Church, 4:540).

“Wherefore,” the Savior tells us, “be faithful, praying always, having your lamps trimmed and burning, and oil with you, that you may be ready at the coming of the Bridegroom—For behold, verily, verily, I say unto you, that I come quickly” (D\&C 33:17–18).

I testify of Jesus Christ. I testify that He shall come, as He has promised. And I pray that we will be prepared to meet Him, in the name of Jesus Christ, amen.

\newpage
\settalk{Be Not Deceived}
\setspeaker{Dallin H. Oaks}
\settalkdate{October 2004}
\talkheading{Be Not Deceived}{Dallin H. Oaks}

I am grateful to speak to this worldwide audience of priesthood holders. It is now 8:00 a.m. Sunday morning in the Philippines, my home for the last two years. I send greetings to my beloved associates in that nation and to all of you.

I assume there are no boys in this audience, only young men who are holders of the priesthood. The Apostle Paul wrote that when he was a child he understood as a child, but when he became a man he put away such things (see 1 Corinthians 13:11). You young men are doing the same, so I will speak to you as one man speaks to another.

\intalkheader{I.}

From your position on the road of life, you young men have many miles to go and many choices to make as you seek to return to our Heavenly Father. Along the road there are many signs that beckon. Satan is the author of some of these invitations. He seeks to confuse and deceive us, to get us on a low road that leads away from our eternal destination.

In the beginning, when a powerful spirit was cast down for rebellion, “he became Satan, … the devil, the father of all lies, to deceive and to blind men, and to lead them captive at his will” (Moses 4:4). He and the spirits who follow him are still deceiving the world. Modern revelation declares that “Satan hath sought to deceive you, that he might overthrow you” (see D\&C 50:2–3). Satan’s methods of deception are enticing: music, movies and other media, and the glitter of a good time. When Satan’s lies succeed in deceiving us, we become vulnerable to his power.

Here are some ways the devil will try to deceive us. God’s commandments and the teachings of His prophets warn against each of them.

Surely you have seen and heard these arguments, brethren. They will come at you in classrooms and hallways, in what you read, and in what you see in popular entertainment. Many in the world deny the need for a Savior. Others deny that there is any right or wrong, and they scoff at the idea of sin or a devil. Still others rely on the mercy of God and ignore His justice. The prophet said, “There shall be many which shall teach after this manner, false and vain and foolish doctrines” (2 Nephi 28:9).

The Apostle Paul gave pointed warnings about the “perilous times” that would come in the last days. “For men shall be lovers of their own selves, … disobedient to parents, unthankful, unholy, without natural affection, … despisers of those that are good, … lovers of pleasures more than lovers of God” (2 Timothy 3:1–4). He also said that “evil men and seducers shall wax worse and worse, deceiving, and being deceived” (v. 13). In a moment I will discuss what Paul told young Timothy about how to avoid this wickedness.

The Apostle gave another warning against being deceived by the devil and his pawns:

“Know ye not that the unrighteous shall not inherit the kingdom of God? Be not deceived: neither fornicators, nor idolaters, nor adulterers, nor effeminate, nor abusers of themselves with mankind,

“Nor thieves, nor covetous, nor drunkards, nor revilers, nor extortioners, shall inherit the kingdom of God” (1 Corinthians 6:9–10).

Be not deceived, brethren. Heed the ancient and modern prophetic warnings against thievery, drunkenness, and all forms of sexual sin. The deceiver seeks to destroy your spirituality by all of these means. Paul warns us against those who “lie in wait to deceive” “by the sleight of men, and cunning craftiness” (Ephesians 4:14). Beware of the slick package and the glitz of a good time. What the devil portrays as fun can be spiritually fatal.

\intalkheader{II.}

As we look about us, we see many who are practicing deception. We hear of prominent officials who have lied about their secret acts. We learn of honored sports heroes who have lied about gambling on the outcome of their games or using drugs to enhance their performance. We see less well-known persons engaging in evil acts in secret they would never do in public. Perhaps they think no one will ever know. But God always knows. And He has repeatedly warned that the time will come when “[our] iniquities shall be spoken upon the housetops, and [our] secret acts shall be revealed” (D\&C 1:3; see also Mormon 5:8; D\&C 38:7).

“Be not deceived,” the Apostle Paul taught. “God is not mocked: for whatsoever a man soweth, that shall he also reap. For he that soweth to his flesh shall of the flesh reap corruption; but he that soweth to the Spirit shall of the Spirit reap life everlasting” (Galatians 6:7–8).

In other words, if we indulge in drugs or pornography or other evils that the Apostle called sowing to the flesh, eternal law dictates that we harvest corruption rather than life eternal. That is the justice of God, and mercy cannot rob justice. If an eternal law is broken, the punishment affixed to that law must be suffered. Some of this can be satisfied by the Savior’s Atonement, but the merciful cleansing of a soiled sinner comes only after repentance (see Alma 42:22–25), which for some sins is a prolonged and painful process. Otherwise, “he that exercises no faith unto repentance is exposed to the whole law of the demands of justice; therefore only unto him that has faith unto repentance is brought about the great and eternal plan of redemption” (Alma 34:16).

Fortunately, repentance is possible. For the most serious sins we need to confess to our bishop and seek his loving help. For other sins it may be sufficient for us to confess to the Lord and to whomever we have wronged. Most lying is of this sort. If you have deceived someone, resolve now to stop carrying the burden. Make it right and get on with your life.

\intalkheader{III.}

Now I wish to speak about how each of us can avoid being deceived on matters of eternal importance. I have two texts. The first is what Paul taught Timothy after giving him the warning I quoted earlier. Continue in the things you have learned and been assured of, he wrote, “knowing of whom thou hast learned them” (2 Timothy 3:14). In other words, you have been taught righteousness and assured of its truth, so stay with it. Continuing, Paul reminded his young friend “that from a child thou hast known the holy scriptures, which are able to make thee wise unto salvation” through faith in our Savior (v. 15). Hold fast to the scriptures, whose teachings protect us against evil.

The parable of the ten virgins teaches that when the Lord comes in His glory, of all followers of Christ invited to the wedding feast, only half will be given entrance. The inspired explanation of this parable reveals our second source of protection:

“For they that are wise and have received the truth, and have taken the Holy Spirit for their guide, and have not been deceived—verily I say unto you, they shall not be hewn down and cast into the fire, but shall abide the day” (D\&C 45:57).

The other half will be denied entrance because they are not prepared. It is not enough to have received the truth. We must also “have taken the Holy Spirit for [our] guide” and “not [be] deceived.”

How do we take the Holy Spirit for our guide? We must repent of our sins each week and renew our covenants by partaking of the sacrament with clean hands and a pure heart, as we are commanded to do (see D\&C 59:8–9, 12). Only in this way can we have the divine promise that we will “always have his Spirit to be with [us]” (D\&C 20:77). That Spirit is the Holy Ghost, whose mission is to teach us, to lead us to truth, and to testify of the Father and the Son (see John 14:26; 15:26; 16:13; 3 Nephi 11:32, 36).

To avoid being deceived, we must also follow the promptings of that Spirit. The Lord taught this principle in the 46th section of the Doctrine and Covenants:

“That which the Spirit testifies unto you even so I would that ye should do in all holiness of heart, walking uprightly before me, considering the end of your salvation, doing all things with prayer and thanksgiving, that ye may not be seduced by evil spirits, or doctrines of devils, or the commandments of men. …

“Wherefore, beware lest ye are deceived; and that ye may not be deceived seek ye earnestly the best gifts, always remembering for what they are given” (D\&C 46:7–8).

The Holy Ghost will protect us against being deceived, but to realize that wonderful blessing we must always do the things necessary to retain that Spirit. We must keep the commandments, pray for guidance, and attend church and partake of the sacrament each Sunday. And we must never do anything to drive away that Spirit. Specifically, we should avoid pornography, alcohol, tobacco, and drugs, and always, always avoid violations of the law of chastity. We must never take things into our bodies or do things with our bodies that drive away the Spirit of the Lord and leave us without our spiritual protection against deception.

I will conclude by describing another subtle form of deception—the idea that it is enough to hear and believe without acting on that belief. Many prophets have taught against that deception. The Apostle James wrote, “Be ye doers of the word, and not hearers only, deceiving your own selves” (James 1:22). King Benjamin taught, “And now, if you believe all these things see that ye do them” (Mosiah 4:10). And in modern revelation the Lord declares, “If you will that I give unto you a place in the celestial world, you must prepare yourselves by doing the things which I have commanded you and required of you” (D\&C 78:7).

It is not enough to know that God lives, that Jesus Christ is our Savior, and that the gospel is true. We must take the high road by acting upon that knowledge. It is not enough to know that President Gordon B. Hinckley is God’s prophet. We must put his teachings to work in our lives. It is not enough to have a calling. We must fulfill our responsibilities. The things taught in this conference are not just to fill our minds. They are to motivate and guide our actions.

I testify that these things are true, and I pray that we will do all that is necessary to avoid the deceptions of the devil, in the name of Jesus Christ, amen.

\newpage
\settalk{Pornography}
\setspeaker{Dallin H. Oaks}
\settalkdate{April 2005}
\talkheading{Pornography}{Dallin H. Oaks}

Last summer Sister Oaks and I returned from two years in the Philippines. We loved our service there, and we loved returning home. When we have been away, we see our surroundings in a new light, with increased appreciation and sometimes with new concerns.

We were concerned to see the inroads pornography had made in the United States while we were away. For many years our Church leaders have warned against the dangers of images and words intended to arouse sexual desires. Now the corrupting influence of pornography, produced and disseminated for commercial gain, is sweeping over our society like an avalanche of evil.

At our last conference, President Gordon B. Hinckley devoted an entire talk to this subject, warning in the plainest terms that “this is a very serious problem even among us” (in Conference Report, Oct. 2004, 64; or Ensign, Nov. 2004, 61). Most of the bishops we meet in stake conferences now report major concerns with this problem.

My fellow holders of the Melchizedek Priesthood, and also our young men, I wish to speak to you today about pornography. I know that many of you are exposed to this and that many of you are being stained by it.

In concentrating my talk on this subject, I feel like the prophet Jacob, who told the men of his day that it grieved him to speak so boldly in front of their sensitive wives and children. But notwithstanding the difficulty of the task, he said he had to speak to the men about this subject because God had commanded him (see Jacob 2:7–11). I do so for the same reason.

In the second chapter of the book that bears his name, Jacob condemns men for their “whoredoms” (vv. 23, 28). He told them they had “broken the hearts of [their] tender wives, and lost the confidence of [their] children, because of [their] bad examples before them” (v. 35).

What were these grossly wicked “whoredoms”? No doubt some men were already guilty of evil acts. But the main focus of Jacob’s great sermon was not with evil acts completed, but with evil acts contemplated.

Jacob began his sermon by telling the men that “as yet, [they had] been obedient unto the word of the Lord” (Jacob 2:4). However, he then told them he knew their thoughts, that they were “beginning to labor in sin, which sin appeareth very abominable … unto God” (Jacob 2:5). “I must testify unto you concerning the wickedness of your hearts” (Jacob 2:6), he added. Jacob was speaking as Jesus spoke when He said, “Whosoever looketh on a woman to lust after her hath committed adultery with her already in his heart” (Matthew 5:28; see also 3 Nephi 12:28; D\&C 59:6; 63:16).

More than 30 years ago, I urged BYU students to avoid the “promotional literature of illicit sexual relations” in what they read and viewed. I gave this analogy:

“Pornographic or erotic stories and pictures are worse than filthy or polluted food. The body has defenses to rid itself of unwholesome food. With a few fatal exceptions, bad food will only make you sick but do no permanent harm. In contrast, a person who feasts upon filthy stories or pornographic or erotic pictures and literature records them in this marvelous retrieval system we call a brain. The brain won’t vomit back filth. Once recorded, it will always remain subject to recall, flashing its perverted images across your mind and drawing you away from the wholesome things in life.”

Here, brethren, I must tell you that our bishops and our professional counselors are seeing an increasing number of men involved with pornography, and many of those are active members. Some involved in pornography apparently minimize its seriousness and continue to exercise the priesthood of God because they think no one will know of their involvement. But the user knows, brethren, and so does the Lord.

Some have suggested that pornography should be a separate question in the temple recommend interview. It is already. At least five different questions should elicit a confession and discussion on this subject if the person being interviewed has the spiritual sensitivity and honesty we expect of those who worship in the house of the Lord.

One of the Savior’s most memorable teachings applies to men who are secretly viewing pornography:

“Woe unto you, scribes and Pharisees, hypocrites! for ye make clean the outside of the cup and of the platter, but within they are full of extortion and excess.

“Thou blind Pharisee, cleanse first that which is within the cup and platter, that the outside of them may be clean also” (Matthew 23:25–26; see also Alma 60:23).

The Savior continues His denunciation of those who treat what is visible but neglect to cleanse the inner man:

“Ye are like unto whited sepulchres, which indeed appear beautiful outward, but are within full of dead men’s bones, and of all uncleanness.

“Even so ye also outwardly appear righteous unto men, but within ye are full of hypocrisy and iniquity” (Matthew 23:27–28).

The immediate spiritual consequences of such hypocrisy are devastating. Those who seek out and use pornography forfeit the power of their priesthood. The Lord declares: “When we undertake to cover our sins, … behold, the heavens withdraw themselves; the Spirit of the Lord is grieved; and when it is withdrawn, Amen to the priesthood or the authority of that man” (D\&C 121:37).

Patrons of pornography also lose the companionship of the Spirit. Pornography produces fantasies that destroy spirituality. “To be carnally minded is death”—spiritual death (Romans 8:6; see also 2 Nephi 9:39).

The scriptures repeatedly teach that the Spirit of the Lord will not dwell in an unclean tabernacle. When we worthily partake of the sacrament, we are promised that we will “always have his Spirit to be with [us].” To qualify for that promise we covenant that we will “always remember him” (D\&C 20:77). Those who seek out and use pornography for sexual stimulation obviously violate that covenant. They also violate a sacred covenant to refrain from unholy and impure practices. They cannot have the Spirit of the Lord to be with them. All such need to heed the Apostle Peter’s plea: “Repent therefore of this thy wickedness, and pray God, if perhaps the thought of thine heart may be forgiven thee” (Acts 8:22).

Brethren, you have noticed that I am not discussing the effects of pornography on mental health or criminal behavior. I am discussing its effects on spirituality—on our ability to have the companionship of the Spirit of the Lord and our capacity to exercise the power of the priesthood.

Pornography also inflicts mortal wounds on our most precious personal relationships. In his talk to men of the priesthood last October, President Hinckley quoted the letter of a woman who asked him to warn Church members that pornography “has the effect of damaging hearts and souls to their very depths, strangling the life out of relationships” (in Conference Report, Oct. 2004, 64; or Ensign, Nov. 2004, 60).

At a recent stake conference a woman handed me a similar letter. Her husband had also served in important Church callings for many years while addicted to pornography. She told of great difficulty in getting priesthood leaders to take this problem of pornography seriously: “I got all kinds of responses—like I was overreacting or it was my fault. The bishop we have now has been great. And now after 15 years my husband is trying to deal with his addiction, but now it is 15 years harder to quit for him, and the loss has been incalculable.”

Pornography impairs one’s ability to enjoy a normal emotional, romantic, and spiritual relationship with a person of the opposite sex. It erodes the moral barriers that stand against inappropriate, abnormal, or illegal behavior. As conscience is desensitized, patrons of pornography are led to act out what they have witnessed, regardless of its effects on their life and the lives of others.

Pornography is also addictive. It impairs decision-making capacities and it “hooks” its users, drawing them back obsessively for more and more. A man who had been addicted to pornography and to hard drugs wrote me this comparison: “In my eyes, cocaine doesn’t hold a candle to this. I have done both. … Quitting even the hardest drugs was nothing compared to [trying to quit pornography]” (letter of Mar. 20, 2005).

Some seek to justify their indulgence by arguing that they are only viewing “soft,” not “hard,” porn. A wise bishop called this refusing to see evil as evil. He quoted men seeking to justify their viewing choices by comparisons such as “not as bad as” or “only one bad scene.” But the test of what is evil is not its degree but its effect. When persons entertain evil thoughts long enough for the Spirit to withdraw, they lose their spiritual protection, and they are subject to the power and direction of the evil one. When they use Internet or other pornography for what this bishop described as “arousal on demand” (letter of Mar. 13, 2005), they are deeply soiled by sin.

King Benjamin’s great sermon describes the terrible consequences. When we withdraw from the Spirit of the Lord, we become an enemy to righteousness, we have a lively sense of our guilt, and we “shrink from the presence of the Lord” (see Mosiah 2:36–38). “Mercy hath no claim on that man,” he concluded; “therefore his final doom is to endure a never-ending torment” (v. 39).

Consider the tragic example of King David. Though a spiritual giant in Israel, he allowed himself to look upon something he should not have viewed (see 2 Samuel 11). Tempted by what he saw, he violated two of the Ten Commandments, beginning with “Thou shalt not commit adultery” (Exodus 20:14). In this way a prophet-king fell from his exaltation (see D\&C 132:39).

But the good news is that no one needs to follow the evil, downward descent to torment. Everyone caught on that terrible escalator has the key to reverse his course. He can escape. Through repentance he can be clean.

Alma the Younger described it:

“Yea, I did remember all my sins and iniquities, for which I was tormented with the pains of hell. …

“… The very thought of coming into the presence of my God did rack my soul with inexpressible horror. …

“And it came to pass that as I was thus racked with torment, while I was harrowed up by the memory of my many sins, behold, I remembered also to have heard my father prophesy unto the people concerning the coming of one Jesus Christ, a Son of God, to atone for the sins of the world.

“Now, as my mind caught hold upon this thought, I cried within my heart: O Jesus, thou Son of God, have mercy on me, who am in the gall of bitterness, and am encircled about by the everlasting chains of death.

“And now, behold, when I thought this, I could remember my pains no more; yea, I was harrowed up by the memory of my sins no more.

“And oh, what joy, and what marvelous light I did behold; yea, my soul was filled with joy as exceeding as was my pain!” (Alma 36:13–14, 17–20).

My brethren who are caught in this addiction or troubled by this temptation, there is a way.

First, acknowledge the evil. Don’t defend it or try to justify yourself. For at least a quarter century our leaders have pleaded with men, and also with women and children, to avoid this evil. Our current Church magazines are full of warnings, information, and helps on this subject—with more than a score of articles published or to be published this year and last year alone.

Second, seek the help of the Lord and His servants. Hear and heed President Hinckley’s words:

“Plead with the Lord out of the depths of your soul that He will remove from you the addiction which enslaves you. And may you have the courage to seek the loving guidance of your bishop and, if necessary, the counsel of caring professionals” (in Conference Report, Oct. 2004, 66; or Ensign, Nov. 2004, 62).

Third, do all that you can to avoid pornography. If you ever find yourself in its presence—which can happen to anyone in the world in which we live—follow the example of Joseph of Egypt. When temptation caught him in her grip, he left temptation and “got him out” (Genesis 39:12).

Don’t accommodate any degree of temptation. Prevent sin and avoid having to deal with its inevitable destruction. So turn it off! Look away! Avoid it at all costs. Direct your thoughts in wholesome paths. Remember your covenants and be faithful in temple attendance. The wise bishop I quoted earlier reported that “an endowed priesthood bearer’s fall into pornography never occurs during periods of regular worship in the temple; it happens when he has become casual in his temple worship” (letter of Mar. 13, 2005).

We must also act to protect those we love. Parents install alarms to warn if their household is threatened by smoke or carbon monoxide. We should also install protections against spiritual threats, protections like filters on Internet connections and locating access so others can see what is being viewed. And we should build the spiritual strength of our families by loving relationships, family prayer, and scripture study.

Finally, do not patronize pornography. Do not use your purchasing power to support moral degradation. And young women, please understand that if you dress immodestly, you are magnifying this problem by becoming pornography to some of the men who see you.

Please heed these warnings. Let us all improve our personal behavior and redouble our efforts to protect our loved ones and our environment from the onslaught of pornography that threatens our spirituality, our marriages, and our children.

I testify that this is what we should do to enjoy the blessings of Him whom we worship. I testify of Jesus Christ, the Light and Life of the World, whose Church this is, in the name of Jesus Christ, amen.

\newpage
\settalk{Priesthood Authority in the Family and the Church}
\setspeaker{Dallin H. Oaks}
\settalkdate{October 2005}
\talkheading{Priesthood Authority in the Family and the Church}{Dallin H. Oaks}

My subject is priesthood authority in the family and in the Church.

\intalkheader{I.}

My father died when I was seven. I was the oldest of three small children our widowed mother struggled to raise. When I was ordained a deacon, she said how pleased she was to have a priesthood holder in the home. But Mother continued to direct the family, including calling on which one of us would pray when we knelt together each morning. I was puzzled. I had been taught that the priesthood presided in the family. There must be something I didn’t know about how that principle worked.

About this same time we had a neighbor who dominated and sometimes abused his wife. He roared like a lion, and she cowered like a lamb. When they walked to church, she always walked a few steps behind him. That made my mother mad. She was a strong woman who would not accept such domination, and she was angry to see another woman abused in that way. I think of her reaction whenever I see men misusing their authority to gratify their pride or exercise control or compulsion upon their wives in any degree of unrighteousness (see D\&C 121:37).

I have also seen some faithful women who misunderstand how priesthood authority functions. Mindful of their partnership relationship with their husband in the family, some wives have sought to extend that relationship to their husband’s priesthood calling, such as bishop or mission president. In contrast, some single women who have been abused by men (such as in a divorce) mistakenly confuse the priesthood with male abuse and become suspicious of any priesthood authority. A person who has had a bad experience with a particular electrical appliance should not forego using the power of electricity.

Each of the circumstances I have described results from misunderstanding priesthood authority and the great principle that while this authority presides in both the family and the Church, the priesthood functions in a different way in each of them. This principle is understood and applied by the great Church and family leaders I have known, but it is rarely explained. Even the scriptures, which record various exercises of priesthood authority, seldom state expressly which principles only apply to the exercise of priesthood authority in the family or in the Church or which apply in both of them.

\intalkheader{II.}

In our theology and in our practice, the family and the Church have a mutually reinforcing relationship. The family is dependent upon the Church for doctrine, ordinances, and priesthood keys. The Church provides the teachings, authority, and ordinances necessary to perpetuate family relationships to the eternities.

We have programs and activities in both the family and the Church. Each is so interrelated that service to one is service to the other. When children see their parents faithfully perform Church callings, it strengthens their family relationships. When families are strong, the Church is strong. The two run in parallel. Each is important and necessary, and each must be conducted with careful concern for the other. Church programs and activities should not be so all-encompassing that families cannot have everyone present for family time. And family activities should not be scheduled in conflict with sacrament meeting or other vital Church meetings.

We need both Church activities and family activities. If all families were complete and perfect, the Church could sponsor fewer activities. But in a world where many of our youth grow up in homes where one parent is missing, not a member, or otherwise inactive in gospel leadership, there is a special need for Church activities to fill in the gaps. Our widowed mother wisely saw that Church activities would provide her sons with experiences she could not provide because we had no male role model in the home. I remember her urging me to watch and try to be like the good men in our ward. She pushed me to participate in Scouting and other Church activities that would provide this opportunity.

In a church where there are many single members, who do not presently have the companionship the Lord intends for all of His sons and daughters, the Church and its families should also have special concern for the needs of single adults.

\intalkheader{III.}

Priesthood authority functions in both the family and the Church. The priesthood is the power of God used to bless all of His children, male and female. Some of our abbreviated expressions, like “the women and the priesthood,” convey an erroneous idea. Men are not “the priesthood.” Priesthood meeting is a meeting of those who hold and exercise the priesthood. The blessings of the priesthood, such as baptism, receiving the Holy Ghost, the temple endowment, and eternal marriage, are available to men and women alike. The authority of the priesthood functions in the family and in the Church according to the principles the Lord has established.

When my father died, my mother presided over our family. She had no priesthood office, but as the surviving parent in her marriage she had become the governing officer in her family. At the same time, she was always totally respectful of the priesthood authority of our bishop and other Church leaders. She presided over her family, but they presided over the Church.

\intalkheader{IV.}

There are many similarities and some differences in the way priesthood authority functions in the family and in the Church. If we fail to recognize and honor the differences, we encounter difficulties.

Keys. One important difference between its function in the Church and in the family is the fact that all priesthood authority in the Church functions under the direction of the one who holds the appropriate priesthood keys. In contrast, the authority that presides in the family—whether father or single-parent mother—functions in family matters without the need to get authorization from anyone holding priesthood keys. This family authority includes directing the activities of the family, family meetings like family home evenings, family prayer, teaching the gospel, and counseling and disciplining family members. It also includes ordained fathers giving priesthood blessings.

However, priesthood keys are necessary to authorize the ordaining or setting apart of family members. This is because the organization the Lord has made responsible for the performance and recording of priesthood ordinances is the Church, not the family.

Boundaries. Church organizations like wards, quorums, or auxiliaries always have geographic boundaries that limit the responsibility and authority of the callings associated with them. In contrast, family relationships and responsibilities are not dependent upon where different family members reside.

Duration. Church callings are always temporary, but family relationships are permanent.

Call and release. Another contrast concerns the initiation and termination of positions. In the Church, a priesthood leader who holds the necessary keys has the authority to call or release persons serving under his direction. He can even cause that they lose their membership and have their names “blotted out” (see Mosiah 26:34–38; Alma 5:56–62). In contrast, family relationships are so important that the head of the family lacks the authority to make changes in family membership. That can only be done by someone authorized to adjust family relationships under the laws of man or the laws of God. Thus, while a bishop can release a Relief Society president, he cannot sever his relationship with his wife without a divorce under the laws of man. Again, his sealing for eternity cannot be ended without a cancellation procedure under the laws of God. Similarly, a youth serving in a class or quorum presidency can be released by priesthood authority in the ward, but parents cannot divorce a child whose life choices are offensive to them. Family relationships are more enduring than Church relationships.

Partnership. A most important difference in the functioning of priesthood authority in the family and in the Church results from the fact that the government of the family is patriarchal, whereas the government of the Church is hierarchical. The concept of partnership functions differently in the family than in the Church.

The family proclamation gives this beautiful explanation of the relationship between a husband and a wife: While they have separate responsibilities, “in these sacred responsibilities, fathers and mothers are obligated to help one another as equal partners” (“The Family: A Proclamation to the World,” Liahona, Oct. 2004, 49; Ensign, Nov. 1995, 102; emphasis added).

President Spencer W. Kimball said this: “When we speak of marriage as a partnership, let us speak of marriage as a full partnership. We do not want our LDS women to be silent partners or limited partners in that eternal assignment! Please be a contributing and full partner” (The Teachings of Spencer W. Kimball, ed. Edward L. Kimball [1982], 315).

President Kimball also declared, “We have heard of men who have said to their wives, ‘I hold the priesthood and you’ve got to do what I say.’” He decisively rejected that abuse of priesthood authority in a marriage, declaring that such a man “should not be honored in his priesthood” (The Teachings of Spencer W. Kimball, 316).

There are cultures or traditions in some parts of the world that allow men to oppress women, but those abuses must not be carried into the families of the Church of Jesus Christ. Remember how Jesus taught: “Ye have heard that it was said by them of old time, … but I say unto you …” (Matthew 5:27–28). For example, the Savior contradicted the prevailing culture in His considerate treatment of women. Our guide must be the gospel culture He taught.

If men desire the Lord’s blessings in their family leadership, they must exercise their priesthood authority according to the Lord’s principles for its use:

“No power or influence can or ought to be maintained by virtue of the priesthood, only by persuasion, by long-suffering, by gentleness and meekness, and by love unfeigned; by kindness, and pure knowledge” (D\&C 121:41–42).

When priesthood authority is exercised in that way in the patriarchal family, we achieve the “full partnership” President Kimball taught. As declared in the family proclamation:

“Happiness in family life is most likely to be achieved when founded upon the teachings of the Lord Jesus Christ. Successful marriages and families are established and maintained on principles of faith, prayer, repentance, forgiveness, respect, love, [and] compassion” (Liahona, Oct. 2004, 49; Ensign, Nov. 1995, 102).

Church callings are performed according to the principles that govern all of us in working under priesthood authority in the Church. These principles include the persuasion and gentleness taught in the 121st section, which are especially necessary in the hierarchal organization of the Church.

The principles I have identified for the exercise of priesthood authority are more understandable and more comfortable for a married woman than for a single woman, especially a single woman who has never been married. She does not now experience priesthood authority in the partnership relationship of marriage. Her experiences with priesthood authority are in the hierarchical relationships of the Church, and some single women feel they have no voice in those relationships. It is, therefore, imperative to have an effective ward council, where male and female ward officers sit down together regularly to counsel under the presiding authority of the bishop.

\intalkheader{V.}

I conclude with some general comments and a personal experience.

The theology of The Church of Jesus Christ of Latter-day Saints centers on the family. Our relationship to God and the purpose of earth life are explained in terms of the family. We are the spirit children of heavenly parents. The gospel plan is implemented through earthly families, and our highest aspiration is to perpetuate those family relationships throughout eternity. The ultimate mission of our Savior’s Church is to help us achieve exaltation in the celestial kingdom, and that can only be accomplished in a family relationship.

No wonder our Church is known as a family-centered church. No wonder we are distressed at the current legal and cultural deteriorations in the position of marriage and childbearing. At a time when the world seems to be losing its understanding of the purpose of marriage and the value of childbearing, it is vital that Latter-day Saints have no confusion about these matters.

The faithful widowed mother who raised us had no confusion about the eternal nature of the family. She always honored the position of our deceased father. She made him a presence in our home. She spoke of the eternal duration of their temple marriage. She often reminded us of what our father would like us to do so we could realize the Savior’s promise that we could be a family forever.

I recall an experience that shows the effect of her teachings. Just before Christmas one year, our bishop asked me, as a deacon, to help him deliver Christmas baskets to the widows of the ward. I carried a basket to each door with his greetings. When he drove me home, there was one basket remaining. He handed it to me and said it was for my mother. As he drove away, I stood in the falling snow wondering why there was a basket for my mother. She never referred to herself as a widow, and it had never occurred to me that she was. To a 12-year-old boy, she wasn’t a widow. She had a husband, and we had a father. He was just away for a while.

I anticipate that glorious future day when the separated will be reunited and all of us will be made complete as the Lord has promised. I testify of Jesus Christ, the Only Begotten Son of the Eternal Father, whose priesthood authority and whose Atonement and Resurrection make it all possible, in the name of Jesus Christ, amen.

\newpage
\settalk{All Men Everywhere}
\setspeaker{Dallin H. Oaks}
\settalkdate{April 2006}
\talkheading{All Men Everywhere}{Dallin H. Oaks}

Last year, at the invitation of a prophet, millions read the Book of Mormon. Millions benefited. For each of us there were blessings of obedience, and most of us also grew in knowledge and testimony of the Lord Jesus Christ, of whom this book is a witness.

Many other things were learned, but what was learned depended on the reader. What we get from a book, especially a sacred text, is mostly dependent on what we take to its reading—in desire and readiness to learn, and in attunement to the light communicated by the Spirit of the Lord.

\intalkheader{I.}

One of the things I learned in this most recent reading of the Book of Mormon was how much God loves all of His children in every nation. In the first chapter Father Lehi praises the Lord, whose “power, and goodness, and mercy are over all the inhabitants of the earth” (1 Nephi 1:14). Again and again the Book of Mormon teaches that the gospel of Jesus Christ is universal in its promise and effect, reaching out to all who ever live on the earth. Here are some examples, quoted directly from that book:

We also read that “his blood atoneth for the sins of those … who have died not knowing the will of God concerning them, or who have ignorantly sinned” (Mosiah 3:11). Similarly, “the blood of Christ atoneth for [little children]” (Mosiah 3:16). These teachings that the resurrecting and cleansing power of the Atonement is for all contradict the assertion that the grace of God saves only a chosen few. His grace is for all. These teachings of the Book of Mormon expand our vision and enlarge our understanding of the all-encompassing love of God and the universal effect of His Atonement for all men everywhere.

\intalkheader{II.}

The Book of Mormon teaches that our Savior “inviteth [all the children of men] to come unto him and partake of his goodness; and he denieth none that come unto him, black and white, bond and free, male and female; and he remembereth the heathen; and all are alike unto God, both Jew and Gentile” (2 Nephi 26:33; see also Alma 5:49).

“He inviteth them all.” We understand “male and female.” We also understand “black and white,” which means all races. But what about “bond and free”? Bond—the opposite of free—means more than slavery. It means being bound (in bondage) to anything from which it is difficult to escape. Bond includes those whose freedom is restricted by physical or emotional afflictions. Bond includes those who are addicted to some substance or practice. Bond surely refers to those who are imprisoned by sin—“encircled about” by what another teaching of the Book of Mormon calls “the chains of hell” (Alma 5:7). Bond includes those who are held down by traditions or customs contrary to the commandments of God (see Matthew 15:3–6; Mark 7:7–9; D\&C 74:4–7; 93:39). Finally, bond also includes those who are confined within the boundaries of other erroneous ideas. The Prophet Joseph Smith taught that we preach to “liberate the captives.” Our Savior “inviteth … all to come unto him and partake of his goodness; … he denieth none that come unto him … ; and all are alike unto God.”

\intalkheader{III.}

The children of God in all nations have His promise that He will manifest Himself to them. The Book of Mormon tells us:

“He manifesteth himself unto all those who believe in him, by the power of the Holy Ghost; yea, unto every nation, kindred, tongue, and people, working mighty miracles, signs, and wonders, among the children of men according to their faith” (2 Nephi 26:13).

Note that these promised manifestations of the Lord are to “every nation, kindred, tongue, and people.” Today we are seeing the fulfillment of that promise in every nation where our missionaries are permitted to labor, even among peoples we have not previously associated with Christianity.

For example, we know of many cases where the Lord has been manifesting Himself to men and women in the nation of Russia, so recently released from the long grip of godless communism. While reading critical or mocking articles about Mormons, two different Russian men felt a strong impression to search out our meeting places. Both met missionaries and joined the Church.

A medical doctor in a village in Nigeria had a dream in which he saw his good friend speaking to a congregation. Intrigued, he traveled to his friend’s village on a Sunday and was astonished to find exactly what he had seen in his dream—a congregation called a ward being taught by his friend, who was their bishop. Impressed with what he heard in repeated visits, he and his wife were taught and baptized. Two months later over 30 others in their village had also joined the Church, and their clinic had become the meeting place.

A man I met from northern India had never even heard the name of Jesus Christ until he saw it on a calendar in the shop of a shoemaker. The Spirit led him to conversion in a Protestant church. Later, during a visit to a distant college town, he saw an advertisement for an American group called “The BYU Young Ambassadors.” During their performance, an inner voice told him to go into the lobby after the program and a man in a blue blazer would tell him what to do. In this way he obtained a Book of Mormon, read it, and was converted to the restored gospel. He has since served as a missionary and as a bishop.

A little girl in Thailand felt a memory of a loving Father in Heaven. As she grew older, she would often pray and counsel with Him in her heart. In her early 20s she met our missionaries. Their teachings confirmed the loving personal feelings for God she remembered from her childhood. She was baptized and served a full-time mission in Thailand.

Only 5 percent of the people in Cambodia are Christians. A family in that country was searching for the truth. While their 11-year-old son was riding his bicycle he saw some men in white shirts and ties showing someone a picture and asking who it was. He felt he should stop. As he watched, he was prompted to say, “That is Jesus Christ, the Son of God, and He came to save man.” Then he rode away. It took the missionaries a month to find him and his family. Today, the father is a counselor in the mission presidency.

Last June, a family of five visited the open house for a new chapel in Mongolia. As the father walked through the door a powerful force went through his body, a feeling of peace he had never experienced before. Tears flowed. He asked the missionaries what that amazing feeling was and how he could feel it again. Soon, the entire family was baptized.

These are only a few examples. There are thousands more.

\intalkheader{IV.}

The Book of Mormon also teaches that the great Creator died “for all men, that all men might become subject unto him” (2 Nephi 9:5). Being subject to our Savior means that if our sins are to be forgiven through His Atonement, we must comply with the conditions He has prescribed, including faith, repentance, and baptism. The fulfillment of these conditions depends on our desires, our choices, and our actions. “He cometh into the world that he may save all men if they will hearken unto his voice” (2 Nephi 9:21).

The Lord provides a way for all His children, and He desires that each of us come unto Him. In the closing chapter of the Book of Mormon, Moroni pleads:

“Come unto Christ, and be perfected in him, and deny yourselves of all ungodliness; and if ye shall deny yourselves of all ungodliness, and love God with all your might, mind and strength, then is his grace sufficient for you, that by his grace ye may be perfect in Christ” (Moroni 10:32).

\intalkheader{V.}

The Bible tells us how God made a covenant with Abraham and promised him that through him all “families” or “nations” of the earth would be blessed (see Genesis 12:3; 22:18). What we call the Abrahamic covenant opens the door for God’s choicest blessings to all of His children everywhere. The Bible teaches that “if ye be Christ’s, then are ye Abraham’s seed, and heirs according to the promise” (Galatians 3:29; see also Abraham 2:10). The Book of Mormon promises that all who receive and act upon the Lord’s invitation to “repent and believe in his Son” become “the covenant people of the Lord” (2 Nephi 30:2). This is a potent reminder that neither riches nor lineage nor any other privileges of birth should cause us to believe that we are “better one than another” (Alma 5:54; see also Jacob 3:9). Indeed, the Book of Mormon commands, “Ye shall not esteem one flesh above another, or one man shall not think himself above another” (Mosiah 23:7).

The Bible teaches that some of Abraham’s descendants would be scattered “into all the kingdoms of the earth,” “among all nations,” and from “one end of the earth even unto the other” (Deuteronomy 28:25, 37, 64). The Book of Mormon affirms this teaching, declaring that the descendants of Abraham would be “scattered upon all the face of the earth, and … among all nations” (1 Nephi 22:3).

The Book of Mormon adds to our knowledge of how our Savior’s earthly ministry reached out to all of His scattered flock. In addition to His ministry in what we now call the Middle East, the Book of Mormon records His appearance and teachings to the Nephites on the American continent (see 3 Nephi 11–28). There He repeated that the Father had commanded Him to visit the other sheep which were not of the land of Jerusalem (see 3 Nephi 16:1; John 10:16). He also said that he would visit others “who [had] not as yet heard [His] voice” (see 3 Nephi 16:2–3). As prophesied centuries earlier (see 2 Nephi 29:12), the Savior told His followers in the Americas that He was going “to show [Himself]” to these “lost tribes of Israel, for they are not lost unto the Father, for he knoweth whither he hath taken them” (3 Nephi 17:4).

The Book of Mormon is a great witness that the Lord loves all people everywhere. It declares that “he shall manifest himself unto all nations” (1 Nephi 13:42). “Know ye not that there are more nations than one?” the Lord said through the prophet Nephi.

“Know ye not that I, the Lord your God, have created all men, and that I remember those who are upon the isles of the sea; and that I rule in the heavens above and in the earth beneath; and I bring forth my word unto the children of men, yea, even upon all the nations of the earth?” (2 Nephi 29:7).

Similarly, the prophet Alma taught that “the Lord doth grant unto all nations, of their own nation and tongue, to teach his word, yea, in wisdom, all that he seeth fit that they should have” (Alma 29:8).

\intalkheader{VI.}

The Lord not only manifests Himself to all nations; He also commands that they write His words:

“Know ye not that the testimony of two nations is a witness unto you that I am God, that I remember one nation like unto another? Wherefore, I speak the same words unto one nation like unto another. …

“For I command all men … that they shall write the words which I speak unto them. …

“For behold, I shall speak unto the Jews and they shall write it; and I shall also speak unto the Nephites and they shall write it; and I shall also speak unto the other tribes of the house of Israel, which I have led away, and they shall write it; and I shall also speak unto all nations of the earth and they shall write it” (2 Nephi 29:8, 11–12; see also 1 Nephi 13:38–39).

Furthermore, the Book of Mormon teaches that all of these groups will have the writings of the others (see 2 Nephi 29:13).

We conclude from this that the Lord will eventually cause the inspired teachings He has given to His children in various nations to be brought forth for the benefit of all people. This will include accounts of the visit of the resurrected Lord to what we call the lost tribes of Israel and His revelations to all the seed of Abraham. The finding of the Dead Sea Scrolls shows one way this can occur.

When new writings come forth—and according to prophecy they will—we hope they will not be treated with the rejection some applied to the Book of Mormon because they already had a Bible (see 2 Nephi 29:3–10). As the Lord said through a prophet in that book, “And because that I have spoken one word ye need not suppose that I cannot speak another; for my work is not yet finished; neither shall it be until the end of man” (2 Nephi 29:9).

Truly, the gospel is for all men everywhere—every nation, every people. All are invited.

We live in the day foretold when righteousness is sent down out of heaven and truth out of the earth “to sweep the earth as with a flood,” and to gather out the elect “from the four quarters of the earth” (Moses 7:62). The Book of Mormon has come forth to remind us of the covenants of the Lord, to the convincing of all “that Jesus is the Christ, the Eternal God, manifesting himself unto all nations” (Book of Mormon title page). I add this, my testimony of Him and His mission, in the name of Jesus Christ, amen.

\newpage
\settalk{He Heals the Heavy Laden}
\setspeaker{Dallin H. Oaks}
\settalkdate{October 2006}
\talkheading{He Heals the Heavy Laden}{Dallin H. Oaks}

The Savior said, “Come unto me, all ye that labour and are heavy laden, and I will give you rest” (Matthew 11:28).

Many carry heavy burdens. Some have lost a loved one to death or care for one who is disabled. Some have been wounded by divorce. Others yearn for an eternal marriage. Some are caught in the grip of addictive substances or practices like alcohol, tobacco, drugs, or pornography. Others have crippling physical or mental impairments. Some are challenged by same-gender attraction. Some have terrible feelings of depression or inadequacy. In one way or another, many are heavy laden.

To each of us our Savior gives this loving invitation:

“Come unto me, all ye that labour and are heavy laden, and I will give you rest.

“Take my yoke upon you, and learn of me; for I am meek and lowly in heart: and ye shall find rest unto your souls.

“For my yoke is easy, and my burden is light” (Matthew 11:28–30).

The scriptures contain many accounts of the Savior’s healing the heavy laden. He caused the blind to see; the deaf to hear; the palsied, withered, or maimed to be restored; lepers to be cleansed; and unclean spirits to be cast out. Often we read that the person healed of these physical ailments was “made whole” (see Matthew 14:36; 15:28; Mark 6:56; 10:52; Luke 17:19; John 5:9).

Jesus healed many from physical diseases, but He did not withhold healing from those who sought to be “made whole” from other ailments. Matthew writes that He healed every sickness and every disease among the people (see Matthew 4:23; 9:35). Great multitudes followed Him, and He “healed them all” (Matthew 12:15). Surely these healings included those whose sicknesses were emotional, mental, or spiritual. He healed them all.

In His early sermon in the synagogue, Jesus read aloud from this prophecy of Isaiah: “He hath anointed me to preach the gospel to the poor; he hath sent me to heal the brokenhearted, to preach deliverance to the captives, and recovering of sight to the blind, to set at liberty them that are bruised” (Luke 4:18). As Jesus declared that He was come to fulfill that prophecy, He expressly affirmed that He would heal those with physical ailments and He would also deliver the captive, liberate the bruised, and heal the brokenhearted.

The Gospel of Luke contains many examples of that ministry. It tells of the time when “great multitudes came together to hear [Jesus], and to be healed by him of their infirmities” (Luke 5:15). On other occasions it records that Jesus “cured many of their infirmities” (Luke 7:21) and that He “healed them that had need of healing” (Luke 9:11). It also describes how a great multitude of people out of Judea and Jerusalem and the seacoast of Sidon came down to the plain “to hear him, and to be healed” (Luke 6:17).

When the Savior appeared to the righteous in the New World, He called for persons to come forward who were lame or blind or had other physical ailments. He extended the same invitation to those “that are afflicted in any manner” (3 Nephi 17:7). “Bring them hither,” He said, “and I will heal them” (v. 7). The Book of Mormon tells how the multitude brought forward “all them that were afflicted in any manner” (v. 9). This must have included persons with every kind of physical or emotional or mental affliction, and the scripture testifies that Jesus “did heal them every one” (v. 9).

The Savior teaches that we will have tribulation in the world, but we should “be of good cheer” because He has “overcome the world” (John 16:33). His Atonement reaches and is powerful enough not only to pay the price for sin but also to heal every mortal affliction. The Book of Mormon teaches that “He shall go forth, suffering pains and afflictions and temptations of every kind; and this that the word might be fulfilled which saith he will take upon him the pains and the sicknesses of his people” (Alma 7:11; see also 2 Nephi 9:21).

He knows of our anguish, and He is there for us. Like the good Samaritan in His parable, when He finds us wounded at the wayside, He binds up our wounds and cares for us (see Luke 10:34). Brothers and sisters, the healing power of His Atonement is for you, for us, for all.

His all-encompassing healing power is sought in the prayerful words of our hymn “Master, the Tempest Is Raging”:

\begin{verse}
Master, with anguish of spirit\\
I bow in my grief today.\\
The depths of my sad heart are troubled.\\
Oh, waken and save, I pray!\\
Torrents of sin and of anguish\\
Sweep o’er my sinking soul,\\
And I perish! I perish! dear Master.\\
Oh, hasten and take control!
\end{verse}

(Hymns, no. 105)

We can be healed through the authority of the Melchizedek Priesthood. Jesus gave His Apostles power “to heal all manner of sickness and all manner of disease” (Matthew 10:1; see also Mark 3:15; Luke 9:1–2), and they went forth “preaching the gospel, and healing every where” (Luke 9:6; see also Mark 6:13; Acts 5:16). The Seventy were also sent forth with power and direction to heal the sick (see Luke 10:9; Acts 8:6–7).

Although the Savior could heal all whom He would heal, this is not true of those who hold His priesthood authority. Mortal exercises of that authority are limited by the will of Him whose priesthood it is. Consequently, we are told that some whom the elders bless are not healed because they are “appointed unto death” (D\&C 42:48). Similarly, when the Apostle Paul sought to be healed from the “thorn in the flesh” that buffeted him (2 Corinthians 12:7), the Lord declined to heal him. Paul later wrote that the Lord explained, “My grace is sufficient for thee: for my strength is made perfect in weakness” (v. 9). Paul obediently responded that he would “rather glory in my infirmities, that the power of Christ may rest upon me … for when I am weak, then am I strong” (vv. 9–10).

Healing blessings come in many ways, each suited to our individual needs, as known to Him who loves us best. Sometimes a “healing” cures our illness or lifts our burden. But sometimes we are “healed” by being given strength or understanding or patience to bear the burdens placed upon us.

The people who followed Alma were in bondage to wicked oppressors. When they prayed for relief, the Lord told them He would deliver them eventually, but in the meantime He would ease their burdens “that even you cannot feel them upon your backs, even while you are in bondage; and this will I do that ye may stand as witnesses … that I, the Lord God, do visit my people in their afflictions” (Mosiah 24:14). In that case the people did not have their burdens removed, but the Lord strengthened them so that “they could bear up their burdens with ease, and they did submit cheerfully and with patience to all the will of the Lord” (v. 15).

This same promise and effect applies to you mothers who are widowed or divorced, to you singles who are lonely, to you caregivers who are burdened, to you persons who are addicted, and to all of us whatever our burden. “Come unto Christ,” the prophet says, “and be perfected in him” (Moroni 10:32).

At times we may despair that our burdens are too great. When it seems that a tempest is raging in our lives, we may feel abandoned and cry out like the disciples in the storm, “Master, carest thou not that we perish?” (Mark 4:38). At such times we should remember His reply: “Why are ye so fearful? how is it that ye have no faith?” (v. 40).

The healing power of the Lord Jesus Christ—whether it removes our burdens or strengthens us to endure and live with them like the Apostle Paul—is available for every affliction in mortality.

After I gave a general conference talk on the evils of pornography (see Conference Report, Apr. 2005, 91–95; or Ensign, May 2005, 87–90), I received many letters from persons burdened with this addiction. Some of these letters were from men who had overcome pornography. One man wrote:

“There are several lessons I’ve gleaned from my experience coming out of the darkness of a sin that so thoroughly dominates the lives of the people it ensnares: (1) This is a major problem that is unbelievably difficult to overcome. … (2) The most important source of support and strength in the repentance process is the Savior. … (3) Intense, daily scripture study, regular temple worship, and serious, contemplative participation in the ordinance of the sacrament are all indispensable parts of a true repentance process. This, I assume, is because all of these activities serve to deepen and strengthen one’s relationship with the Savior, one’s understanding of His atoning sacrifice, and one’s faith in His healing power” (letter dated Oct. 24, 2005).

“Come unto me,” the Savior said, “and ye shall find rest unto your souls” (Matthew 11:28–29). That heavy-laden man turned to the Savior, and so can each of us.

A woman whose marriage was threatened by her husband’s addiction to pornography wrote how she stood beside him for five pain-filled years until, as she said, “through the gift of our precious Savior’s glorious Atonement and what He taught me about forgiveness, [my husband] finally is free—and so am I.” As one who needed no cleansing from sin, but only sought a loved one’s deliverance from captivity, she wrote this advice:

“Commune with the Lord. … He is your best friend! He knows your pain because He has felt it for you already. He is ready to carry that burden. Trust Him enough to place it at His feet and allow Him to carry it for you. Then you can have your anguish replaced with His peace, in the very depths of your soul” (letter dated Apr. 18, 2005).

A man wrote a General Authority about how the power of the Atonement helped him with his problem of same-gender attraction. He had been excommunicated for serious transgressions that violated his temple covenants and his responsibilities to his children. He had to choose whether to attempt to live the gospel or whether to continue a course contrary to its teachings.

“I knew it would be difficult,” he wrote, “but I didn’t realize what I would have to go through.” His letter describes the emptiness and loneliness and the incredible pain he experienced from deep within his soul as he sought to return. He prayed mightily for forgiveness, sometimes for hours at a time. He was sustained by reading the scriptures, by the companionship of a loving bishop, and by priesthood blessings. But what finally made the difference was the help of the Savior. He explained:

“It [was] only through Him and His Atonement. … I now feel an overwhelming gratitude. My pains have been almost more than I could bear at times, and yet they were so small compared to what He suffered. Where there once was darkness in my life, there is now love and gratitude.”

He continues: “Some profess that change is possible and therapy is the only answer. They are very learned on the subject and have so much to offer those who struggle … , but I worry that they forget to involve Heavenly Father in the process. If change is to happen, it will happen according to the will of God. I also worry that many people focus on the causes of [same-gender attraction]. … There is no need to determine why I have [this challenge]. I don’t know if I was born with it, or if environmental factors contributed to it. The fact of the matter is that I have this struggle in my life and what I do with it from this point forward is what matters” (letter dated Mar. 25, 2006).

The persons who wrote these letters know that the Atonement of Jesus Christ and the healing it offers do much more than provide the opportunity for repentance from sins. The Atonement also gives us the strength to endure “pains and afflictions and temptations of every kind,” because our Savior also took upon Him “the pains and the sicknesses of his people” (Alma 7:11). Brothers and sisters, if your faith and prayers and the power of the priesthood do not heal you from an affliction, the power of the Atonement will surely give you the strength to bear the burden.

“Come unto me, all ye that labour and are heavy laden,” the Savior said, “and I will give you rest … unto your souls” (Matthew 11:28–29).

As we struggle with the challenges of mortality, I pray for each of us, as the prophet Mormon prayed for his son, Moroni: “May Christ lift thee up, and may his sufferings and death, … and his mercy and long-suffering, and the hope of his glory and of eternal life, rest in your mind forever” (Moroni 9:25).

I testify of Jesus Christ, our Savior, who invites us all to come unto Him and be perfected in Him. He will bind up our wounds and He will heal the heavy laden. In the name of Jesus Christ, amen.

\newpage
\settalk{Divorce}
\setspeaker{Dallin H. Oaks}
\settalkdate{April 2007}
\talkheading{Divorce}{Dallin H. Oaks}

I have felt impressed to speak about divorce. This is a sensitive subject because it evokes such strong emotions from persons it has touched in different ways. Some see themselves or their loved ones as the victims of divorce. Others see themselves as its beneficiaries. Some see divorce as evidence of failure. Others consider it an essential escape hatch from marriage. In one way or another, divorce touches most families in the Church.

Whatever your perspective, please listen as I try to speak plainly about the effects of divorce on the eternal family relationships we seek under the gospel plan. I speak out of concern, but with hope.

\intalkheader{I.}

We live in a world in which the whole concept of marriage is in peril and where divorce is commonplace.

The concept that society has a strong interest in preserving marriages for the common good as well as the good of the couple and their children has been replaced for many by the idea that marriage is only a private relationship between consenting adults, terminable at the will of either.

Nations that had no divorce law have adopted one, and most nations permitting divorces have made them easier to obtain. Unfortunately, under current no-fault divorce laws, it can be easier to sever a marriage relationship with an unwanted spouse than an employment relationship with an unwanted employee. Some even refer to a first marriage as a “starter marriage,” like a small home one uses for a while before moving on.

The weakening of the concept that marriages are permanent and precious has far-reaching consequences. Influenced by their own parents’ divorce or by popular notions that marriage is a ball and chain that prevents personal fulfillment, some young people shun marriage. Many who marry withhold full commitment, poised to flee at the first serious challenge.

In contrast, modern prophets have warned that looking upon marriage “as a mere contract that may be entered into at pleasure … and severed at the first difficulty … is an evil meriting severe condemnation,” especially where children are made to suffer.

In ancient times and even under tribal laws in some countries where we now have members, men have power to divorce their wives for any trivial thing. Such unrighteous oppression of women was rejected by the Savior, who declared:

“Moses because of the hardness of your hearts suffered you to put away your wives: but from the beginning it was not so.

“And I say unto you, Whosoever shall put away his wife, except it be for fornication, and shall marry another, committeth adultery: and whoso marrieth her which is put away doth commit adultery” (Matthew 19:8–9).

The kind of marriage required for exaltation—eternal in duration and godlike in quality—does not contemplate divorce. In the temples of the Lord, couples are married for all eternity. But some marriages do not progress toward that ideal. Because “of the hardness of [our] hearts,” the Lord does not currently enforce the consequences of the celestial standard. He permits divorced persons to marry again without the stain of immorality specified in the higher law. Unless a divorced member has committed serious transgressions, he or she can become eligible for a temple recommend under the same worthiness standards that apply to other members.

\intalkheader{II.}

There are many good Church members who have been divorced. I speak first to them. We know that many of you are innocent victims—members whose former spouses persistently betrayed sacred covenants or abandoned or refused to perform marriage responsibilities for an extended period. Members who have experienced such abuse have firsthand knowledge of circumstances worse than divorce.

When a marriage is dead and beyond hope of resuscitation, it is needful to have a means to end it. I saw examples of this in the Philippines. Two days after their temple marriage, a husband deserted his young wife and has not been heard from for over 10 years. A married woman fled and obtained a divorce in another country, but her husband, who remained behind, is still married in the eyes of the Philippine law. Since there is no provision for divorce in that country, these innocent victims of desertion have no way to end their married status and go forward with their lives.

We know that some look back on their divorces with regret at their own partial or predominant fault in the breakup. All who have been through divorce know the pain and need the healing power and hope that come from the Atonement. That healing power and that hope are there for them and also for their children.

\intalkheader{III.}

Now I speak to married members, especially to any who may be considering divorce.

I strongly urge you and those who advise you to face up to the reality that for most marriage problems, the remedy is not divorce but repentance. Often the cause is not incompatibility but selfishness. The first step is not separation but reformation. Divorce is not an all-purpose solution, and it often creates long-term heartache. A broad-based international study of the levels of happiness before and after “major life events” found that, on average, persons are far more successful in recovering their level of happiness after the death of a spouse than after a divorce. Spouses who hope that divorce will resolve conflicts often find that it aggravates them, since the complexities that follow divorce—especially where there are children—generate new conflicts.

Think first of the children. Because divorce separates the interests of children from the interests of their parents, children are its first victims. Scholars of family life tell us that the most important cause of the current decline in the well-being of children is the current weakening of marriage, because family instability decreases parental investment in children. We know that children raised in a single-parent home after divorce have a much higher risk for drug and alcohol abuse, sexual promiscuity, poor school performance, and various kinds of victimization.

A couple with serious marriage problems should see their bishop. As the Lord’s judge, he will give counsel and perhaps even discipline that will lead toward healing.

Bishops do not counsel members to divorce, but they can help members with the consequences of their decisions. Under the law of the Lord, a marriage, like a human life, is a precious, living thing. If our bodies are sick, we seek to heal them. We do not give up. While there is any prospect of life, we seek healing again and again. The same should be true of our marriages, and if we seek Him, the Lord will help us and heal us.

Latter-day Saint spouses should do all within their power to preserve their marriages. They should follow the marriage enrichment counsel in the First Presidency’s message in the April 2007 Ensign and Liahona. To avoid so-called “incompatibility,” they should be best friends, kind and considerate, sensitive to each other’s needs, always seeking to make each other happy. They should be partners in family finances, working together to regulate their desires for temporal things.

Of course, there can be times when one spouse falls short and the other is wounded and feels pain. When that happens, the one who is wronged should balance current disappointments against the good of the past and the brighter prospects of the future.

Don’t treasure up past wrongs, reprocessing them again and again. In a marriage relationship, festering is destructive; forgiving is divine (see D\&C 64:9–10). Plead for the guidance of the Spirit of the Lord to forgive wrongs (as President Faust has just taught us so beautifully), to overcome faults, and to strengthen relationships.

If you are already descending into the low state of marriage-in-name-only, please join hands, kneel together, and prayerfully plead for help and the healing power of the Atonement. Your humble and united pleadings will bring you closer to the Lord and to each other and will help you in the hard climb back to marital harmony.

Consider these observations of a wise bishop with extensive experience in counseling members with marriage problems. Speaking of those who eventually divorced, he said:

“Universally, every couple or individual said they recognized that divorce was not a good thing, but they all insisted that their situation was different.

“Universally, they focused on the fault of the spouse and attributed little responsibility to their own behavior. Communication had withered.

“Universally, they were looking back, not willing to leave the baggage of past behavior on the roadside and move on.

“Part of the time, serious sin was involved, but more often they had just ‘fallen out of love,’ saying, ‘He doesn’t satisfy my needs anymore,’ or, ‘She has changed.’

“All were worried about the effect on the children, but always the conclusion was ‘it’s worse for them to have us together and fighting.’”

In contrast, the couples who followed this bishop’s counsel and stayed together emerged with their marriages even stronger. That prospect began with their mutual commitment to keep the commandments, stay active in their Church attendance, scripture reading, and prayer, and to work on their own shortcomings. They “recognized the importance and power of the Atonement for their spouse and for themselves,” and “they were patient and would try again and again.” When the couples he counseled did these things, repenting and working to save their marriages, this bishop reported that “healing was achieved 100 percent of the time.”

Even those who think their spouse is entirely to blame should not act hastily. One study found “no evidence that divorce or separation typically made adults happier than staying in an unhappy marriage. Two out of three unhappily married adults who avoided divorce reported being happily married five years later.” A woman who persisted in an intolerable marriage for many years until the children were raised explained: “There were three parties to our marriage—my husband and I and the Lord. I told myself that if two of us could hang in there, we could hold it together.”

The power of hope expressed in these examples is sometimes rewarded with repentance and reformation, but sometimes it is not. Personal circumstances vary greatly. We cannot control and we are not responsible for the choices of others, even when they impact us so painfully. I am sure the Lord loves and blesses husbands and wives who lovingly try to help spouses struggling with such deep problems as pornography or other addictive behavior or with the long-term consequences of childhood abuse.

Whatever the outcome and no matter how difficult your experiences, you have the promise that you will not be denied the blessings of eternal family relationships if you love the Lord, keep His commandments, and just do the best you can. When young Jacob “suffered afflictions and much sorrow” from the actions of other family members, Father Lehi assured him, “Thou knowest the greatness of God; and he shall consecrate thine afflictions for thy gain” (2 Nephi 2:1–2). Similarly, the Apostle Paul assured us that “all things work together for good to them that love God” (Romans 8:28).

\intalkheader{IV.}

In conclusion, I speak briefly to those contemplating marriage. The best way to avoid divorce from an unfaithful, abusive, or unsupportive spouse is to avoid marriage to such a person. If you wish to marry well, inquire well. Associations through “hanging out” or exchanging information on the Internet are not a sufficient basis for marriage. There should be dating, followed by careful and thoughtful and thorough courtship. There should be ample opportunities to experience the prospective spouse’s behavior in a variety of circumstances. Fiancés should learn everything they can about the families with whom they will soon be joined in marriage. In all of this, we should realize that a good marriage does not require a perfect man or a perfect woman. It only requires a man and a woman committed to strive together toward perfection.

President Spencer W. Kimball taught: “Two individuals approaching the marriage altar must realize that to attain the happy marriage which they hope for they must know that marriage … means sacrifice, sharing, and even a reduction of some personal liberties. It means long, hard economizing. It means children who bring with them financial burdens, service burdens, care and worry burdens; but also it means the deepest and sweetest emotions of all.”

From personal experience, I testify to the sweetness of the marriage and family life that the family proclamation describes as founded upon a husband and wife’s “solemn responsibility to love and care for each other and for their children” and “upon the teachings of the Lord Jesus Christ.” I testify of Him as our Savior and pray in His name for all who strive for the supreme blessings of an eternal family, in the name of Jesus Christ, amen.

\newpage
\settalk{Good, Better, Best}
\setspeaker{Dallin H. Oaks}
\settalkdate{October 2007}
\talkheading{Good, Better, Best}{Dallin H. Oaks}

Most of us have more things expected of us than we can possibly do. As breadwinners, as parents, as Church workers and members, we face many choices on what we will do with our time and other resources.

\intalkheader{I.}

We should begin by recognizing the reality that just because something is good is not a sufficient reason for doing it. The number of good things we can do far exceeds the time available to accomplish them. Some things are better than good, and these are the things that should command priority attention in our lives.

Jesus taught this principle in the home of Martha. While she was “cumbered about much serving” (Luke 10:40), her sister, Mary, “sat at Jesus’ feet, and heard his word” (v. 39). When Martha complained that her sister had left her to serve alone, Jesus commended Martha for what she was doing (v. 41) but taught her that “one thing is needful: and Mary hath chosen that good part, which shall not be taken away from her” (v. 42). It was praiseworthy for Martha to be “careful and troubled about many things” (v. 41), but learning the gospel from the Master Teacher was more “needful.” The scriptures contain other teachings that some things are more blessed than others (see Acts 20:35; Alma 32:14–15).

A childhood experience introduced me to the idea that some choices are good but others are better. I lived for two years on a farm. We rarely went to town. Our Christmas shopping was done in the Sears, Roebuck catalog. I spent hours poring over its pages. For the rural families of that day, catalog pages were like the shopping mall or the Internet of our time.

Something about some displays of merchandise in the catalog fixed itself in my mind. There were three degrees of quality: good, better, and best. For example, some men’s shoes were labeled good (\$1.84), some better (\$2.98), and some best (\$3.45).

As we consider various choices, we should remember that it is not enough that something is good. Other choices are better, and still others are best. Even though a particular choice is more costly, its far greater value may make it the best choice of all.

Consider how we use our time in the choices we make in viewing television, playing video games, surfing the Internet, or reading books or magazines. Of course it is good to view wholesome entertainment or to obtain interesting information. But not everything of that sort is worth the portion of our life we give to obtain it. Some things are better, and others are best. When the Lord told us to seek learning, He said, “Seek ye out of the best books words of wisdom” (D\&C 88:118; emphasis added).

\intalkheader{II.}

Some of our most important choices concern family activities. Many breadwinners worry that their occupations leave too little time for their families. There is no easy formula for that contest of priorities. However, I have never known of a man who looked back on his working life and said, “I just didn’t spend enough time with my job.”

In choosing how we spend time as a family, we should be careful not to exhaust our available time on things that are merely good and leave little time for that which is better or best. A friend took his young family on a series of summer vacation trips, including visits to memorable historic sites. At the end of the summer he asked his teenage son which of these good summer activities he enjoyed most. The father learned from the reply, and so did those he told of it. “The thing I liked best this summer,” the boy replied, “was the night you and I laid on the lawn and looked at the stars and talked.” Super family activities may be good for children, but they are not always better than one-on-one time with a loving parent.

The amount of children-and-parent time absorbed in the good activities of private lessons, team sports, and other school and club activities also needs to be carefully regulated. Otherwise, children will be overscheduled, and parents will be frazzled and frustrated. Parents should act to preserve time for family prayer, family scripture study, family home evening, and the other precious togetherness and individual one-on-one time that binds a family together and fixes children’s values on things of eternal worth. Parents should teach gospel priorities through what they do with their children.

Family experts have warned against what they call “the overscheduling of children.” In the last generation, children are far busier and families spend far less time together. Among many measures of this disturbing trend are the reports that structured sports time has doubled, but children’s free time has declined by 12 hours per week, and unstructured outdoor activities have fallen by 50 percent.

The number of those who report that their “whole family usually eats dinner together” has declined 33 percent. This is most concerning because the time a family spends together “eating meals at home [is] the strongest predictor of children’s academic achievement and psychological adjustment.” Family mealtimes have also been shown to be a strong bulwark against children’s smoking, drinking, or using drugs. There is inspired wisdom in this advice to parents: what your children really want for dinner is you.

President Gordon B. Hinckley has pleaded that we “work at our responsibility as parents as if everything in life counted on it, because in fact everything in life does count on it.”

He continued: “I ask you men, particularly, to pause and take stock of yourselves as husbands and fathers and heads of households. Pray for guidance, for help, for direction, and then follow the whisperings of the Spirit to guide you in the most serious of all responsibilities, for the consequences of your leadership in your home will be eternal and everlasting.”

The First Presidency has called on parents “to devote their best efforts to the teaching and rearing of their children in gospel principles. … The home is the basis of a righteous life, and no other instrumentality can take its place … in … this God-given responsibility.” The First Presidency has declared that “however worthy and appropriate other demands or activities may be, they must not be permitted to displace the divinely-appointed duties that only parents and families can adequately perform.”

\intalkheader{III.}

Church leaders should be aware that Church meetings and activities can become too complex and burdensome if a ward or a stake tries to have the membership do everything that is good and possible in our numerous Church programs. Priorities are needed there also.

Members of the Quorum of the Twelve have stressed the importance of exercising inspired judgment in Church programs and activities. Elder L. Tom Perry taught this principle in our first worldwide leadership training meeting in 2003. Counseling the same leaders in 2004, Elder Richard G. Scott said: “Adjust your activities to be consistent with your local conditions and resources. … Make sure that the essential needs are met, but do not go overboard in creating so many good things to do that the essential ones are not accomplished. … Remember, don’t magnify the work to be done—simplify it.”

In general conference last year, Elder M. Russell Ballard warned against the deterioration of family relationships that can result when we spend excess time on ineffective activities that yield little spiritual sustenance. He cautioned against complicating our Church service “with needless frills and embellishments that occupy too much time, cost too much money, and sap too much energy. … The instruction to magnify our callings is not a command to embellish and complicate them. To innovate does not necessarily mean to expand; very often it means to simplify. … What is most important in our Church responsibilities,” he said, “is not the statistics that are reported or the meetings that are held but whether or not individual people—ministered to one at a time just as the Savior did—have been lifted and encouraged and ultimately changed.”

Stake presidencies and bishoprics need to exercise their authority to weed out the excessive and ineffective busyness that is sometimes required of the members of their stakes or wards. Church programs should focus on what is best (most effective) in achieving their assigned purposes without unduly infringing on the time families need for their “divinely appointed duties.”

But here is a caution for families. Suppose Church leaders reduce the time required by Church meetings and activities in order to increase the time available for families to be together. This will not achieve its intended purpose unless individual family members—especially parents—vigorously act to increase family togetherness and one-on-one time. Team sports and technology toys like video games and the Internet are already winning away the time of our children and youth. Surfing the Internet is not better than serving the Lord or strengthening the family. Some young men and women are skipping Church youth activities or cutting family time in order to participate in soccer leagues or to pursue various entertainments. Some young people are amusing themselves to death—spiritual death.

Some uses of individual and family time are better, and others are best. We have to forego some good things in order to choose others that are better or best because they develop faith in the Lord Jesus Christ and strengthen our families.

\intalkheader{IV.}

Here are some other illustrations of good, better, and best:

It is good to belong to our Father in Heaven’s true Church and to keep all of His commandments and fulfill all of our duties. But if this is to qualify as best, it should be done with love and without arrogance. We should, as we sing in a great hymn, “crown [our] good with brotherhood,” showing love and concern for all whom our lives affect.

To our hundreds of thousands of home teachers and visiting teachers, I suggest that it is good to visit our assigned families; it is better to have a brief visit in which we teach doctrine and principle; and it is best of all to make a difference in the lives of some of those we visit. That same challenge applies to the many meetings we hold—it’s good to hold a meeting, better to teach a principle, but best to actually improve lives as a result of the meeting.

As we approach 2008 and a new course of study in our Melchizedek Priesthood quorums and Relief Societies, I renew our caution about how we use the Teachings of Presidents of the Church manuals. Many years of inspired work have produced our 2008 volume of the teachings of Joseph Smith, the founding prophet of this dispensation. This is a landmark among Church books. In the past, some teachers have given a chapter of the Teachings manuals no more than a brief mention and then substituted a lesson of their own choice. It may have been a good lesson, but this is not an acceptable practice. A gospel teacher is called to teach the subject specified from the inspired materials provided. The best thing a teacher can do with Teachings: Joseph Smith is to select and quote from the words of the Prophet on principles that are specially suited to the needs of class members and then direct a class discussion on how to apply those principles in the circumstances of their lives.

I testify of our Heavenly Father, whose children we are and whose plan is designed to qualify us for “eternal life, … the greatest of all the gifts of God” (D\&C 14:7; see also D\&C 76:51–59). I testify of Jesus Christ, whose Atonement makes it possible. And I testify that we are led by prophets—our President, Gordon B. Hinckley, and his counselors—in the name of Jesus Christ, amen.

\newpage
\settalk{Testimony}
\setspeaker{Dallin H. Oaks}
\settalkdate{April 2008}
\talkheading{Testimony}{Dallin H. Oaks}

A testimony of the gospel is a personal witness borne to our souls by the Holy Ghost that certain facts of eternal significance are true and that we know them to be true. Such facts include the nature of the Godhead and our relationship to its three members, the effectiveness of the Atonement, and the reality of the Restoration.

A testimony of the gospel is not a travelogue, a health log, or an expression of love for family members. It is not a sermon. President Kimball taught that the moment we begin preaching to others, our testimony is ended.

\intalkheader{I.}

Various questions arise as we hear others bear testimony or as we consider bearing testimony ourselves.

\intalkheader{II.}

What do we mean when we testify and say that we know the gospel is true? Contrast that kind of knowledge with “I know it is cold outside” or “I know I love my wife.” These are three different kinds of knowledge, each learned in a different way. Knowledge of outside temperature can be verified by scientific proof. Knowledge that we love our spouse is personal and subjective. While not capable of scientific proof, it is still important. The idea that all important knowledge is based on scientific evidence is simply untrue.

While there are some “evidences” for gospel truths (for example, see Psalm 19:1; Helaman 8:24), scientific methods will not yield spiritual knowledge. This is what Jesus taught in response to Simon Peter’s testimony that He was the Christ: “Blessed art thou, Simon Bar-jona: for flesh and blood hath not revealed it unto thee, but my Father which is in heaven” (Matthew 16:17). The Apostle Paul explained this. In a letter to the Corinthian Saints, he said, “The things of God knoweth no man, but [by] the Spirit of God” (1 Corinthians 2:11; see also John 14:17).

In contrast, we know the things of man by the ways of man, but “the natural man receiveth not the things of the Spirit of God: for they are foolishness unto him: neither can he know them, because they are spiritually discerned” (1 Corinthians 2:14).

The Book of Mormon teaches that God will manifest the truth of spiritual things unto us by the power of the Holy Ghost (see Moroni 10:4–5). In modern revelation God promises us that we will receive “knowledge” by His telling us in our mind and in our heart “by the Holy Ghost” (D\&C 8:1–2).

One of the greatest things about our Heavenly Father’s plan for His children is that each of us can know the truth of that plan for ourselves. That revealed knowledge does not come from books, from scientific proof, or from intellectual pondering. As with the Apostle Peter, we can receive that knowledge directly from our Heavenly Father through the witness of the Holy Ghost.

When we know spiritual truths by spiritual means, we can be just as sure of that knowledge as scholars and scientists are of the different kinds of knowledge they have acquired by different methods.

The Prophet Joseph Smith provided a wonderful example of this. When he was persecuted for telling people about his vision, he likened his circumstance to the Apostle Paul, who was ridiculed and reviled as he made his defense before King Agrippa (see Acts 26). “But all this did not destroy the reality of his vision,” Joseph said. “He had seen a vision, he knew he had, and all the persecution under heaven could not make it otherwise. … So it was with me,” Joseph continued. “I had actually seen a light, and in the midst of that light I saw two Personages, and they did in reality speak to me. … I had seen a vision; I knew it, and I knew that God knew it, and I could not deny it, neither dared I” (Joseph Smith—History 1:24–25).

\intalkheader{III.}

That was Joseph Smith’s testimony. What about ours? How can we come to know and testify that what he said was true? How does one gain what we call a testimony?

The first step in gaining any kind of knowledge is to really desire to know. In the case of spiritual knowledge, the next step is to ask God in sincere prayer. As we read in modern revelation, “If thou shalt ask, thou shalt receive revelation upon revelation, knowledge upon knowledge, that thou mayest know the mysteries and peaceable things—that which bringeth joy, that which bringeth life eternal” (D\&C 42:61).

Here is what Alma wrote about what he did: “Behold, I have fasted and prayed many days that I might know these things of myself. And now I do know of myself that they are true; for the Lord God hath made them manifest unto me by his Holy Spirit” (Alma 5:46).

As we desire and seek, we should remember that acquiring a testimony is not a passive thing but a process in which we are expected to do something. Jesus taught, “If any man will do his will, he shall know of the doctrine, whether it be of God, or whether I speak of myself” (John 7:17).

Another way to seek a testimony seems astonishing when compared with the methods of obtaining other knowledge. We gain or strengthen a testimony by bearing it. Someone even suggested that some testimonies are better gained on the feet bearing them than on the knees praying for them.

A personal testimony is fundamental to our faith. Consequently, the things we must do to acquire, strengthen, and retain a testimony are vital to our spiritual life. In addition to those already stated, we need to partake of the sacrament each week (see D\&C 59:9) to qualify for the precious promise that we will “always have his Spirit to be with [us]” (D\&C 20:77). Of course, that Spirit is the source of our testimonies.

\intalkheader{IV.}

Those who have a testimony of the restored gospel also have a duty to share it. The Book of Mormon teaches that we should “stand as witnesses of God at all times and in all things, and in all places that [we] may be in” (Mosiah 18:9).

One of the most impressive teachings on the relationship between the gift of a testimony and the duty to bear it is in the 46th section of the Doctrine and Covenants. In describing different kinds of spiritual gifts, this revelation states:

“To some it is given by the Holy Ghost to know that Jesus Christ is the Son of God, and that he was crucified for the sins of the world.

“To others it is given to believe on their words, that they also might have eternal life if they continue faithful” (vv. 13–14; see also John 20:29).

Those who have the gift to know have an obvious duty to bear their witness so that those who have the gift to believe on their words might also have eternal life.

There has never been a greater need for us to profess our faith, privately and publicly (see D\&C 60:2). Though some profess atheism, there are many who are open to additional truths about God. To these sincere seekers, we need to affirm the existence of God the Eternal Father, the divine mission of our Lord and Savior, Jesus Christ, and the reality of the Restoration. We must be valiant in our testimony of Jesus. Each of us has many opportunities to proclaim our spiritual convictions to friends and neighbors, to fellow workers, and to casual acquaintances. We should use these opportunities to express our love for our Savior, our witness of His divine mission, and our determination to serve Him. Our children should also hear us bear our testimonies frequently. We should also strengthen our children by encouraging them to define themselves by their growing testimonies, not just by their recognitions in scholarship, sports, or other school activities.

\intalkheader{V.}

We live in a time when some misrepresent the beliefs of those they call Mormons and even revile us because of them. When we encounter such misrepresentations, we have a duty to speak out to clarify our doctrine and what we believe. We should be the ones to state our beliefs rather than allowing others the final word in misrepresenting them. This calls for testimony, which can be expressed privately to an acquaintance or publicly in a small or large meeting. As we testify of the truth we know, we should faithfully follow the caution to speak “in mildness and in meekness” (D\&C 38:41). We should never be overbearing, shrill, or reviling. As the Apostle Paul taught, we should speak the truth in love (see Ephesians 4:15). Anyone can disagree with our personal testimony, but no one can refute it.

\intalkheader{VI.}

In closing, I refer to the relationship between obedience and knowledge. Members who have a testimony and who act upon it under the direction of their Church leaders are sometimes accused of blind obedience.

Of course we have leaders, and of course we are subject to their decisions and directions in the operation of the Church and in the performance of needed priesthood ordinances. But when it comes to learning and knowing the truth of the gospel—our personal testimonies—we each have a direct relationship with God, our Eternal Father, and His Son, Jesus Christ, through the powerful witness of the Holy Ghost. This is what our critics fail to understand. It puzzles them that we can be united in following our leaders and yet independent in knowing for ourselves.

Perhaps the puzzle some feel can be explained by the reality that each of us has two different channels to God. We have a channel of governance through our prophet and other leaders. This channel, which has to do with doctrine, ordinances, and commandments, results in obedience. We also have a channel of personal testimony, which is direct to God. This has to do with His existence, our relationship to Him, and the truth of His restored gospel. This channel results in knowledge. These two channels are mutually reinforcing: knowledge encourages obedience (see Deuteronomy 5:27; Moses 5:11), and obedience enhances knowledge (see John 7:17; D\&C 93:1).

We all act upon or give obedience to knowledge. Whether in science or religion, our obedience is not blind when we act upon knowledge suited to the subject of our action. A scientist receives and acts upon a trusted certification of the content or conditions of a particular experiment. In matters of religion, a believer’s source of knowledge is spiritual, but the principle is the same. In the case of Latter-day Saints, when the Holy Ghost gives our souls a witness of the truth of the restored gospel and the calling of a modern prophet, our choice to follow those teachings is not blind obedience.

In all of our testifying we must avoid arrogance and pride. We should remember the Book of Mormon rebuke to a people who had such pride in the greater things God had given them that they afflicted their neighbors (see Jacob 2:20). Jacob said this was “abominable unto him who created all flesh” because “the one being is as precious in his sight as the other” (Jacob 2:21). Later, Alma cautioned that “ye shall not esteem one flesh above another, or one man shall not think himself above another” (Mosiah 23:7).

I close with my testimony. I know that we have a Heavenly Father, whose plan brings us to earth and provides the conditions and destiny of our eternal journey. I know that we have a Savior, Jesus Christ, whose teachings define the plan and whose Atonement gives the assurance of immortality and the opportunity for eternal life. I know that the Father and the Son appeared to the Prophet Joseph Smith to restore the fulness of the gospel in these latter days. And I know that we are led today by a prophet, President Thomas S. Monson, who holds the keys to authorize priesthood holders to perform the ordinances prescribed for our progress toward eternal life. In the name of Jesus Christ, amen.

\newpage
\settalk{Sacrament Meeting and the Sacrament}
\setspeaker{Dallin H. Oaks}
\settalkdate{October 2008}
\talkheading{Sacrament Meeting and the Sacrament}{Dallin H. Oaks}

We live in the perilous times prophesied by the Apostle Paul (see 2 Timothy 3:1). Those who try to walk the straight and narrow path see inviting detours on every hand. We can be distracted, degraded, downhearted, or depressed. How can we have the Spirit of the Lord to guide our choices and keep us on the path?

In modern revelation the Lord gave the answer in this commandment:

“And that thou mayest more fully keep thyself unspotted from the world, thou shalt go to the house of prayer and offer up thy sacraments upon my holy day;

“For verily this is a day appointed unto you to rest from your labors, and to pay thy devotions unto the Most High” (D\&C 59:9–10).

This is a commandment with a promise. By participating weekly and appropriately in the ordinance of the sacrament, we qualify for the promise that we will “always have his Spirit to be with [us]” (D\&C 20:77). That Spirit is the foundation of our testimony. It testifies of the Father and the Son, brings all things to our remembrance, and leads us into truth. It is the compass to guide us on our path. This gift of the Holy Ghost, President Wilford Woodruff taught, “is the greatest gift that can be bestowed upon man” (Deseret Weekly, Apr. 6, 1889, 451).

\intalkheader{I.}

The ordinance of the sacrament makes the sacrament meeting the most sacred and important meeting in the Church. It is the only Sabbath meeting the entire family can attend together. Its content in addition to the sacrament should always be planned and presented to focus our attention on the Atonement and teachings of the Lord Jesus Christ.

My first memories of sacrament meeting are set in the small Utah town where I was ordained a deacon and participated in passing the sacrament. Measured against those memories, the sacrament meetings I now attend in many different wards are greatly improved. Typically the sacrament is administered, passed, and received by the members in an atmosphere of quiet reverence. The conducting of the meeting, including the necessary business, is brief and dignified, and the talks are spiritual in content and delivery. The music is appropriate, and so are the prayers. This is the standard, and it represents great progress since the experiences of my youth.

There are occasional exceptions. I sense that some in the rising generation and even some adults have not yet come to understand the significance of this meeting and the importance of individual reverence and worship in it. The things I feel impressed to teach here are addressed to those who are not yet understanding and practicing these important principles and not yet enjoying the promised spiritual blessings of always having His guiding Spirit to be with them.

\intalkheader{II.}

I begin with how members of the Church should prepare themselves to participate in the ordinance of the sacrament. In a worldwide leadership training meeting five years ago, Elder Russell M. Nelson of the Quorum of the Twelve Apostles taught priesthood leaders of the Church how to plan and conduct sacrament meetings. “We commemorate His Atonement in a very personal way,” Elder Nelson said. “We bring a broken heart and a contrite spirit to our sacrament meeting. It is the highlight of our Sabbath-day observance” (“Worshiping at Sacrament Meeting,” Liahona, Aug. 2004, 12; Ensign, Aug. 2004, 26).

We are seated well before the meeting begins. “During that quiet interval, prelude music is subdued. This is not a time for conversation or transmission of messages but a period of prayerful meditation as leaders and members prepare spiritually for the sacrament” (Liahona, Aug. 2004, 13; Ensign, Aug. 2004, 27).

When the Savior appeared to the Nephites following His Resurrection, He taught them that they should stop the practice of sacrifice by the shedding of blood. Instead, “ye shall offer for a sacrifice unto me a broken heart and a contrite spirit” (3 Nephi 9:20). That commandment, repeated in the modern revelation directing us to partake of the sacrament each week, tells us how we should prepare. As Elder Nelson taught, “Each member of the Church bears responsibility for the spiritual enrichment that can come from a sacrament meeting” (Liahona, Aug. 2004, 14; Ensign, Aug. 2004, 28).

In his writings on the doctrines of salvation, President Joseph Fielding Smith teaches that we partake of the sacrament as our part of commemorating the Savior’s death and sufferings for the redemption of the world. This ordinance was introduced so that we can renew our covenants to serve Him, to obey Him, and to always remember Him. President Smith adds, “We cannot retain the Spirit of the Lord if we do not consistently comply with this commandment” (Doctrines of Salvation, comp. Bruce R. McConkie, 3 vols. [1954–56], 2:341).

\intalkheader{III.}

How we dress is an important indicator of our attitude and preparation for any activity in which we will engage. If we are going swimming or hiking or playing on the beach, our clothing, including our footwear, will indicate this. The same should be true of how we dress when we are to participate in the ordinance of the sacrament. It is like going to the temple. Our manner of dress indicates the degree to which we understand and honor the ordinance in which we will participate.

During sacrament meeting—and especially during the sacrament service—we should concentrate on worship and refrain from all other activities, especially from behavior that could interfere with the worship of others. Even a person who slips into quiet slumber does not interfere with others. Sacrament meeting is not a time for reading books or magazines. Young people, it is not a time for whispered conversations on cell phones or for texting persons at other locations. When we partake of the sacrament, we make a sacred covenant that we will always remember the Savior. How sad to see persons obviously violating that covenant in the very meeting where they are making it.

The music of sacrament meeting is a vital part of our worship. The scriptures teach that the song of the righteous is a prayer unto the Lord (see D\&C 25:12). The First Presidency has declared that “some of the greatest sermons are preached by the singing of hymns” (Hymns, ix). How wonderful when every person in attendance joins in the worship of singing—especially in the hymn that helps us prepare to partake of the sacrament. All sacrament meeting music requires careful planning, always remembering that this music is for worship, not for performance.

President Joseph Fielding Smith taught: “This is an occasion when the gospel should be presented, when we should be called upon to exercise faith, and to reflect on the mission of our Redeemer, and to spend time in the consideration of the saving principles of the gospel, and not for other purposes. Amusement, laughter, light-mindedness, are all out of place in the sacrament meetings of the Latter-day Saints. We should assemble in the spirit of prayer, of meekness, with devotion in our hearts” (Doctrines of Salvation, 2:342).

When we do this—when we join in the solemnity that should always accompany the ordinance of the sacrament and the worship of this meeting—we are qualified for the companionship and revelation of the Spirit. This is the way we get direction for our lives and peace along the way.

\intalkheader{IV.}

The resurrected Lord emphasized the importance of the sacrament when He visited the American continent and instituted this ordinance among the faithful Nephites. He blessed the emblems of the sacrament and gave them to His disciples and the multitude (see 3 Nephi 18:1–10), commanding:

“And this shall ye always do to those who repent and are baptized in my name; and ye shall do it in remembrance of my blood, which I have shed for you, that ye may witness unto the Father that ye do always remember me. And if ye do always remember me ye shall have my Spirit to be with you.

“… And if ye shall always do these things blessed are ye, for ye are built upon my rock.

“But whoso among you shall do more or less than these are not built upon my rock, but are built upon a sandy foundation; and when the rain descends, and the floods come, and the winds blow, and beat upon them, they shall fall” (3 Nephi 18:11–13).

The sacrament is the ordinance that replaced the blood sacrifices and burnt offerings of the Mosaic law, and with it came the Savior’s promise: “And whoso cometh unto me with a broken heart and a contrite spirit, him will I baptize with fire and with the Holy Ghost” (3 Nephi 9:20).

\intalkheader{V.}

Now I speak particularly to the priesthood holders who officiate in the sacrament. This ordinance should always be performed with reverence and dignity. Priests who offer the prayers in behalf of the congregation should speak the words slowly and distinctly, expressing the terms of the covenants and promised blessings. This is a very sacred act.

The teachers who prepare and the deacons who pass the emblems of the sacrament also perform a very sacred act. I love President Thomas S. Monson’s account of how, as a 12-year-old deacon, he was asked by the bishop to take the sacrament to a bedfast brother who longed for this blessing. “His gratitude overwhelmed me,” President Monson said. “The Spirit of the Lord came over me. I stood on sacred ground” (Inspiring Experiences That Build Faith [1994], 188). All who officiate in this sacred ordinance stand on sacred ground.

Young men who officiate in the ordinance of the sacrament should be worthy. The Lord has said, “Be ye clean that bear the vessels of the Lord” (D\&C 38:42). The scriptural warning about partaking of the sacrament unworthily (see 1 Corinthians 11:29; 3 Nephi 18:29) surely applies also to those who officiate in that ordinance. In administering discipline to Church members who have committed serious sins, a bishop can temporarily withdraw the privilege of partaking of the sacrament. That same authority is surely available to withdraw the privilege of officiating in that sacred ordinance.

What I said earlier about the importance of appropriate dress for those who receive the ordinance of the sacrament obviously applies with special force to the young men of the Aaronic Priesthood who officiate in any part of that sacred ordinance. All should be well groomed and modestly dressed. There should be nothing about their personal appearance or actions that would call special attention to themselves or distract anyone present from full attention to the worship and covenant making that are the purpose of this sacred service.

Elder Jeffrey R. Holland gave a valuable teaching on this subject in general conference 13 years ago. Since most of our current deacons were not even born when these words were last spoken here, I repeat them for their benefit and that of their parents and teachers: “May I suggest that wherever possible a white shirt be worn by the deacons, teachers, and priests who handle the sacrament. For sacred ordinances in the Church we often use ceremonial clothing, and a white shirt could be seen as a gentle reminder of the white clothing you wore in the baptismal font and an anticipation of the white shirt you will soon wear into the temple and onto your missions” (in Conference Report, Oct. 1995, 89; or Ensign, Nov. 1995, 68).

Finally, the sacrament is administered only when authorized by the one holding the keys to this priesthood ordinance. This is why the sacrament is not generally served in the home or at family reunions, even where there are sufficient priesthood holders available. Those who officiate at the sacrament table, prepare the sacrament, or pass it to the congregation should be designated by one who holds or exercises the keys of this ordinance. I refer to the bishopric or to the presidencies of the teachers or deacons quorums. “[My] house is a house of order,” the Lord declared (D\&C 132:8).

How can we have the Spirit of the Lord to guide our choices so that we will remain “unspotted from the world” (D\&C 59:9) and on the safe path through mortality? We need to qualify for the cleansing power of the Atonement of Jesus Christ. We do this by keeping His commandment to come to Him with a broken heart and a contrite spirit and in that wonderful weekly meeting partake of the emblems of the sacrament and make the covenants that qualify us for the precious promise that we will always have His Spirit to be with us (see D\&C 20:77). That we may always do so is my humble prayer, which I offer in the name of Him whose Atonement makes it all possible, even Jesus Christ, amen.

\newpage
\settalk{Unselfish Service}
\setspeaker{Dallin H. Oaks}
\settalkdate{April 2009}
\talkheading{Unselfish Service}{Dallin H. Oaks}

Our Savior gave Himself in unselfish service. He taught that each of us should follow Him by denying ourselves of selfish interests in order to serve others.

“If any man will come after me [He said], let him deny himself, and take up his cross, and follow me.

“For whosoever will save his life shall lose it: and whosoever will lose his life for my sake shall find it” (Matthew 16:24–25; see also Matthew 10:39).

\intalkheader{I.}

As a group, Latter-day Saints are unique in following that teaching—unique in the extent of their unselfish service.

Each year tens of thousands of Latter-day Saints submit their papers for full-time missionary service. Seniors put aside the diversions of retirement, the comforts of home, and the loving companionship of children and grandchildren and go forth to serve strangers in unfamiliar places. Young men and women put work and education on hold and make themselves available to serve wherever they are assigned. Hundreds of thousands of faithful members participate in the unselfish service we call “temple work,” which has no motive other than love and service for our fellowmen, living and dead. The same unselfish service is given by legions of officers and teachers in our stakes and wards and branches. All are uncompensated in worldly terms but committed to Christlike service to their fellowmen.

It is not easy to give up our personal priorities and desires. Many years ago a new missionary in England was frustrated and discouraged. He wrote home saying he felt he was wasting his time. His wise father replied, “Forget yourself and go to work.” Young Elder Gordon B. Hinckley went to his knees and covenanted with the Lord that he would try to forget himself and lose himself in the Lord’s service. Years later, as a mature servant of the Lord, Elder Hinckley would say, “He who lives only unto himself withers and dies, while he who forgets himself in the service of others grows and blossoms in this life and in eternity.”

Last January, President Thomas S. Monson taught Brigham Young University students that their student days should include “the matter of spiritual preparation,” including service to others. “An attitude of love characterized the mission of the Master,” President Monson said. “He gave sight to the blind, legs to the lame, and life to the dead. Perhaps when we [face] our Maker, we will not be asked, ‘How many positions did you hold?’ but rather, ‘How many people did you help?’ In reality,” President Monson concluded, “you can never love the Lord until you serve Him by serving His people.”

A familiar example of losing ourselves in the service of others—this one not unique to Latter-day Saints—is the sacrifice parents make for their children. Mothers suffer pain and loss of personal priorities and comforts to bear and rear each child. Fathers adjust their lives and priorities to support a family. The gap between those who are and those who are not willing to do this is widening in today’s world. One of our family members recently overheard a young couple on an airline flight explaining that they chose to have a dog instead of children. “Dogs are less trouble,” they declared. “Dogs don’t talk back, and we never have to ground them.”

We rejoice that so many Latter-day Saint couples are among that unselfish group who are willing to surrender their personal priorities and serve the Lord by bearing and rearing the children our Heavenly Father sends to their care. We also rejoice in those who care for disabled family members and aged parents. None of this service asks, what’s in it for me? All of it requires setting aside personal convenience for unselfish service. All of it stands in contrast to the fame, fortune, and other immediate gratification that are the worldly ways of so many in our day.

Latter-day Saints are uniquely committed to sacrifice. In partaking of the sacrament each week, we witness our commitment to serve the Lord and our fellowmen. In sacred temple ceremonies we covenant to sacrifice and consecrate our time and talents for the welfare of others.

\intalkheader{II.}

Latter-day Saints are also renowned for their ability to unite in cooperative efforts. The Mormon pioneers who colonized the Intermountain West established our honored tradition of unselfish cooperation for the common good. Following in this tradition are our modern “Helping Hands” projects in many nations. In recent elections Latter-day Saints have united with other like-minded persons in efforts to defend marriage. For some, that service has involved great sacrifice and continuing personal pain.

Our members’ religious faith and Church service have taught them how to work in cooperative efforts to benefit the larger community. Because of this, Latter-day Saint volunteers are in great demand in education, local government, charitable causes, and countless other efforts that call for high skills in cooperative efforts and unselfish sacrifice of time and means.

Some attribute our members’ willingness to sacrifice and their skills in cooperative efforts to our effective Church organization or to what skeptics mistakenly call “blind obedience.” Neither explanation is correct. No outside copying of our organization and no application of blind obedience could duplicate the record of this Church or the performance of its members. Our willingness to sacrifice and our skills in cooperative efforts come from our faith in the Lord Jesus Christ, from the inspired teachings of our leaders, and from the commitments and covenants we knowingly make.

\intalkheader{III.}

Unfortunately, some Latter-day Saints seem to forgo unselfish service to others, choosing instead to fix their priorities on the standards and values of the world. Jesus cautioned that Satan desires to sift us like wheat (see Luke 22:31; 3 Nephi 18:18), which means to make us common like all those around us. But Jesus taught that we who follow Him should be precious and unique, “the salt of the earth” (Matthew 5:13), and “the light of the world” to shine forth to all men (Matthew 5:14, 16; see also 3 Nephi 18:24).

We do not serve our Savior well if we fear man more than God. He rebuked some leaders in His restored Church for seeking the praise of the world and for having their minds on the things of the earth more than on the things of the Lord (see D\&C 30:2; 58:39). Those chastisements remind us that we are called to establish the Lord’s standards, not to follow the world’s. Elder John A. Widtsoe declared, “We cannot walk as other men, or talk as other men, or do as other men, for we have a different destiny, obligation, and responsibility placed upon us, and we must fit ourselves [to it].” That reality has current application to every trendy action, including immodest dress. As a wise friend observed, “You can’t be a life saver if you look like all the other swimmers on the beach.”

Those who are caught up in trying to save their lives by seeking the praise of the world are actually rejecting the Savior’s teaching that the only way to save our eternal life is to love one another and lose our lives in service.

C. S. Lewis explained this teaching of the Savior: “The moment you have a self at all, there is a possibility of putting yourself first—wanting to be the centre—wanting to be God, in fact. That was the sin of Satan: and that was the sin he taught the human race. Some people think the fall of man had something to do with sex, but that is a mistake. … What Satan put into the heads of our remote ancestors was the idea that they could ‘be like gods’—could set up on their own as if they had created themselves—be their own masters—invent some sort of happiness for themselves outside God, apart from God. And out of that hopeless attempt has come … the long terrible story of man trying to find something other than God which will make him happy.”

A selfish person is more interested in pleasing man—especially himself—than in pleasing God. He looks only to his own needs and desires. He walks “in his own way, and after the image of his own god, whose image is in the likeness of the world” (D\&C 1:16). Such a person becomes disconnected from the covenant promises of God (see D\&C 1:15) and from the mortal friendship and assistance we all need in these tumultuous times. In contrast, if we love and serve one another as the Savior taught, we remain connected to our covenants and to our associates.

\intalkheader{IV.}

We live in a time when sacrifice is definitely out of fashion, when the outside forces that taught our ancestors the need for unselfish cooperative service have diminished. Someone has called this the “me” generation—a selfish time when everyone seems to be asking, what’s in it for me? Even some who should know better seem to be straining to win the praise of those who mock and scoff from the “great and spacious building” identified in vision as the pride of the world (see 1 Nephi 8:26–28; 11:35–36).

The worldly aspiration of our day is to get something for nothing. The ancient evil of greed shows its face in the assertion of entitlement: I am entitled to this or that because of who I am—a son or a daughter, a citizen, a victim, or a member of some other group. Entitlement is generally selfish. It demands much, and it gives little or nothing. Its very concept causes us to seek to elevate ourselves above those around us. This separates us from the divine, evenhanded standard of reward that when anyone obtains any blessing from God, it is by obedience to the law on which that blessing is predicated (see D\&C 130:21).

The effects of greed and entitlement are evident in the multimillion-dollar bonuses of some corporate executives. But the examples are more widespread than that. Greed and ideas of entitlement have also fueled the careless and widespread borrowing and excessive consumerism behind the financial crises that threaten to engulf the world.

Gambling is another example of greed and selfishness. The gambler ventures a minimum amount in the hope of a huge return that comes by taking it away from others. No matter how it is disguised, getting something for nothing is contrary to the gospel law of the harvest: “Whatsoever a man soweth, that shall he also reap” (Galatians 6:7; see also 2 Corinthians 9:6).

The values of the world wrongly teach that “it’s all about me.” That corrupting attitude produces no change and no growth. It is contrary to eternal progress toward the destiny God has identified in His great plan for His children. The plan of the gospel of Jesus Christ lifts us above our selfish desires and teaches us that this life is all about what we can become.

A great example of unselfish service is the late Mother Teresa of Calcutta, whose vow committed herself and her fellow workers to “wholehearted free service to the poorest of the poor.” She taught that “one thing will always secure heaven for us—the acts of charity and kindness with which we have filled our lives.” “We can do no great things,” Mother Teresa maintained, “only small things with great love.” When this wonderful Catholic servant died, the First Presidency’s message of condolence declared, “Her life of unselfish service is an inspiration to all the world, and her acts of Christian goodness will stand as a memorial for generations to come.” That is what the Savior called losing our lives in service to others.

Each of us should apply that principle to our attitudes in attending church. Some say “I didn’t learn anything today” or “No one was friendly to me” or “I was offended” or “The Church is not filling my needs.” All those answers are self-centered, and all retard spiritual growth.

In contrast, a wise friend wrote:

“Years ago, I changed my attitude about going to church. No longer do I go to church for my sake, but to think of others. I make a point of saying hello to people who sit alone, to welcome visitors, … to volunteer for an assignment. …

“In short, I go to church each week with the intent of being active, not passive, and making a positive difference in people’s lives. Consequently, my attendance at Church meetings is so much more enjoyable and fulfilling.”

All of this illustrates the eternal principle that we are happier and more fulfilled when we act and serve for what we give, not for what we get.

Our Savior teaches us to follow Him by making the sacrifices necessary to lose ourselves in unselfish service to others. If we do, He promises us eternal life, “the greatest of all the gifts of God” (D\&C 14:7), the glory and joy of living in the presence of God the Father and His Son, Jesus Christ. I testify of Them and of Their great plan for the salvation of Their children, in the name of Jesus Christ, amen.

\newpage
\settalk{Love and Law}
\setspeaker{Dallin H. Oaks}
\settalkdate{October 2009}
\talkheading{Love and Law}{Dallin H. Oaks}

I have been impressed to speak about God’s love and God’s commandments. My message is that God’s universal and perfect love is shown in all the blessings of His gospel plan, including the fact that His choicest blessings are reserved for those who obey His laws. These are eternal principles that should guide parents in their love and teaching of their children.

\intalkheader{I.}

I begin with four examples which illustrate some mortal confusion between love and law.

In these examples a person violating commandments asserts that parental love should override the commandments of divine law and the teachings of parents.

The next two examples show mortal confusion about the effect of God’s love.

These persons disbelieve eternal laws which they consider contrary to their concept of the effect of God’s love. Persons who take this position do not understand the nature of God’s love or the purpose of His laws and commandments. The love of God does not supersede His laws and His commandments, and the effect of God’s laws and commandments does not diminish the purpose and effect of His love. The same should be true of parental love and rules.

\intalkheader{II.}

First, consider the love of God, described so meaningfully this morning by President Dieter F. Uchtdorf. “Who shall separate us from the love of Christ?” the Apostle Paul asked. Not tribulation, not persecution, not peril or the sword (see Romans 8:35). “For I am persuaded,” he concluded, “that neither death, nor life, nor angels, nor principalities, nor powers, … nor any other creature, shall be able to separate us from the love of God” (verses 38–39).

There is no greater evidence of the infinite power and perfection of God’s love than is recorded by the Apostle John: “For God so loved the world, that he gave his only begotten Son” (John 3:16). Another Apostle wrote that God “spared not his own Son, but delivered him up for us all” (Romans 8:32). Think how it must have grieved our Heavenly Father to send His Son to endure incomprehensible suffering for our sins. That is the greatest evidence of His love for each of us!

God’s love for His children is an eternal reality, but why does He love us so much, and why do we desire that love? The answer is found in the relationship between God’s love and His laws.

Some seem to value God’s love because of their hope that His love is so great and so unconditional that it will mercifully excuse them from obeying His laws. In contrast, those who understand God’s plan for His children know that God’s laws are invariable, which is another great evidence of His love for His children. Mercy cannot rob justice, and those who obtain mercy are “they who have kept the covenant and observed the commandment” (D\&C 54:6).

We read again and again in the Bible and in modern scriptures of God’s anger with the wicked and of His acting in His wrath against those who violate His laws. How are anger and wrath evidence of His love? Joseph Smith taught that God “institute[d] laws whereby [the spirits that He would send into the world] could have a privilege to advance like himself.” God’s love is so perfect that He lovingly requires us to obey His commandments because He knows that only through obedience to His laws can we become perfect, as He is. For this reason, God’s anger and His wrath are not a contradiction of His love but an evidence of His love. Every parent knows that you can love a child totally and completely while still being creatively angry and disappointed at that child’s self-defeating behavior.

The love of God is so universal that His perfect plan bestows many gifts on all of His children, even those who disobey His laws. Mortality is one such gift, bestowed on all who qualified in the War in Heaven. Another unconditional gift is the universal resurrection: “For as in Adam all die, even so in Christ shall all be made alive” (1 Corinthians 15:22). Many other mortal gifts are not tied to our personal obedience to law. As Jesus taught, our Heavenly Father “maketh his sun to rise on the evil and on the good, and sendeth rain on the just and on the unjust” (Matthew 5:45).

If only we will listen, we can know of God’s love and feel it, even when we are disobedient. A woman recently returned to Church activity gave this description in a sacrament meeting talk: “He has always been there for me, even when I rejected Him. He has always guided me and comforted me with His tender mercies all around me, but I [was] too angry to see and accept incidents and feelings as such.”

\intalkheader{III.}

God’s choicest blessings are clearly contingent upon obedience to God’s laws and commandments. The key teaching is from modern revelation:

“There is a law, irrevocably decreed in heaven before the foundations of this world, upon which all blessings are predicated—

“And when we obtain any blessing from God, it is by obedience to that law upon which it is predicated” (D\&C 130:20–21).

This great principle helps us understand the why of many things, like justice and mercy balanced by the Atonement. It also explains why God will not forestall the exercise of agency by His children. Agency—our power to choose—is fundamental to the gospel plan that brings us to earth. God does not intervene to forestall the consequences of some persons’ choices in order to protect the well-being of other persons—even when they kill, injure, or oppress one another—for this would destroy His plan for our eternal progress. He will bless us to endure the consequences of others’ choices, but He will not prevent those choices.

If a person understands the teachings of Jesus, he or she cannot reasonably conclude that our loving Heavenly Father or His divine Son believes that Their love supersedes Their commandments. Consider these examples.

When Jesus began His ministry, His first message was repentance.

When He exercised loving mercy by not condemning the woman taken in adultery, He nevertheless told her, “Go, and sin no more” (John 8:11).

Jesus taught, “Not every one that saith unto me, Lord, Lord, shall enter into the kingdom of heaven; but he that doeth the will of my Father which is in heaven” (Matthew 7:21).

The effect of God’s commandments and laws is not changed to accommodate popular behavior or desires. If anyone thinks that godly or parental love for an individual grants the loved one license to disobey the law, he or she does not understand either love or law. The Lord declared: “That which breaketh a law, and abideth not by law, but seeketh to become a law unto itself, and willeth to abide in sin, and altogether abideth in sin, cannot be sanctified by law, neither by mercy, justice, nor judgment. Therefore, they must remain filthy still” (D\&C 88:35).

We read in modern revelation, “All kingdoms have a law given” (D\&C 88:36). For example:

“He who is not able to abide the law of a celestial kingdom cannot abide a celestial glory.

“And he who cannot abide the law of a terrestrial kingdom cannot abide a terrestrial glory.

“And he who cannot abide the law of a telestial kingdom cannot abide a telestial glory” (D\&C 88:22–24).

In other words, the kingdom of glory to which the Final Judgment assigns us is not determined by love but by the law that God has invoked in His plan to qualify us for eternal life, “the greatest of all the gifts of God” (D\&C 14:7).

\intalkheader{IV.}

In teaching and reacting to their children, parents have many opportunities to apply these principles. One such opportunity has to do with the gifts parents bestow on their children. Just as God has bestowed some gifts on all of His mortal children without requiring their personal obedience to His laws, parents provide many benefits like housing and food even if their children are not in total harmony with all parental requirements. But, following the example of an all-wise and loving Heavenly Father who has given laws and commandments for the benefit of His children, wise parents condition some parental gifts on obedience.

If parents have a wayward child—such as a teenager indulging in alcohol or drugs—they face a serious question. Does parental love require that these substances or their consumption be allowed in the home, or do the requirements of civil law or the seriousness of the conduct or the interests of other children in the home require that this be forbidden?

To pose an even more serious question, if an adult child is living in cohabitation, does the seriousness of sexual relations outside the bonds of marriage require that this child feel the full weight of family disapproval by being excluded from any family contacts, or does parental love require that the fact of cohabitation be ignored? I have seen both of these extremes, and I believe that both are inappropriate.

Where do parents draw the line? That is a matter for parental wisdom, guided by the inspiration of the Lord. There is no area of parental action that is more needful of heavenly guidance or more likely to receive it than the decisions of parents in raising their children and governing their families. This is the work of eternity.

As parents grapple with these problems, they should remember the Lord’s teaching that we leave the ninety and nine and go out into the wilderness to rescue the lost sheep. President Thomas S. Monson has called for a loving crusade to rescue our brothers and sisters who are wandering in the wilderness of apathy or ignorance. These teachings require continued loving concern, which surely requires continued loving associations.

Parents should also remember the Lord’s frequent teaching that “whom the Lord loveth he chasteneth” (Hebrews 12:6). In his conference talk on tolerance and love, Elder Russell M. Nelson taught that “real love for the sinner may compel courageous confrontation—not acquiescence! Real love does not support self-destructing behavior.”

Wherever the line is drawn between the power of love and the force of law, the breaking of commandments is certain to impact loving family relationships. Jesus taught:

“Suppose ye that I am come to give peace on earth? I tell you, Nay; but rather division:

“For from henceforth there shall be five in one house divided, three against two, and two against three.

“The father shall be divided against the son, and the son against the father; the mother against the daughter, and the daughter against the mother” (Luke 12:51–53).

This sobering teaching reminds us that when family members are not united in striving to keep the commandments of God, there will be divisions. We do all that we can to avoid impairing loving relationships, but sometimes it happens after all we can do.

In the midst of such stress, we must endure the reality that the straying of our loved ones will detract from our happiness, but it should not detract from our love for one another or our patient efforts to be united in understanding God’s love and God’s laws.

I testify of the truth of these things, which are part of the plan of salvation and the doctrine of Christ, of whom I testify in the name of Jesus Christ, amen.

\newpage
\settalk{Healing the Sick}
\setspeaker{Dallin H. Oaks}
\settalkdate{April 2010}
\talkheading{Healing the Sick}{Dallin H. Oaks}

In these times of worldwide turmoil, more and more persons of faith are turning to the Lord for blessings of comfort and healing. I wish to speak to this audience of priesthood holders about healing the sick—by medical science, by prayers of faith, and by priesthood blessings.

\intalkheader{I.}

Latter-day Saints believe in applying the best available scientific knowledge and techniques. We use nutrition, exercise, and other practices to preserve health, and we enlist the help of healing practitioners, such as physicians and surgeons, to restore health.

The use of medical science is not at odds with our prayers of faith and our reliance on priesthood blessings. When a person requested a priesthood blessing, Brigham Young would ask, “Have you used any remedies?” To those who said no because “we wish the Elders to lay hands upon us, and we have faith that we shall be healed,” President Young replied: “That is very inconsistent according to my faith. If we are sick, and ask the Lord to heal us, and to do all for us that is necessary to be done, according to my understanding of the Gospel of salvation, I might as well ask the Lord to cause my wheat and corn to grow, without my plowing the ground and casting in the seed. It appears consistent to me to apply every remedy that comes within the range of my knowledge, and [then] to ask my Father in Heaven … to sanctify that application to the healing of my body.”

Of course we don’t wait until all other methods are exhausted before we pray in faith or give priesthood blessings for healing. In emergencies, prayers and blessings come first. Most often we pursue all efforts simultaneously. This follows the scriptural teachings that we should “pray always” (D\&C 90:24) and that all things should be done in wisdom and order.

\intalkheader{II.}

We know that the prayer of faith, uttered alone or in our homes or places of worship, can be effective to heal the sick. Many scriptures refer to the power of faith in the healing of an individual. The Apostle James taught that we should “pray one for another, that ye may be healed,” adding, “the effectual fervent prayer of a righteous man availeth much” (James 5:16). When the woman who touched Jesus was healed, He told her, “Thy faith hath made thee whole” (Matthew 9:22). Similarly, the Book of Mormon teaches that the Lord “worketh by power, according to the faith of the children of men” (Moroni 10:7).

A recent nationwide survey found that nearly 8 in 10 Americans “believe that miracles still occur today as [they did] in ancient times.” A third of those surveyed said they had “experienced or witnessed a divine healing.” Many Latter-day Saints have experienced the power of faith in healing the sick. We also hear examples of this among people of faith in other churches. A Texas newspaperman described such a miracle. When a five-year-old girl breathed with difficulty and became feverish, her parents rushed her to the hospital. By the time she arrived there, her kidneys and lungs had shut down, her fever was 107 degrees, and her body was bright red and covered with purple lesions. The doctors said she was dying of toxic shock syndrome, cause unknown. As word spread to family and friends, God-fearing people began praying for her, and a special prayer service was held in their Protestant congregation in Waco, Texas. Miraculously, she suddenly returned from the brink of death and was released from the hospital in a little over a week. Her grandfather wrote, “She is living proof that God does answer prayers and work miracles.”

Truly, as the Book of Mormon teaches, God “manifesteth himself unto all those who believe in him, by the power of the Holy Ghost; yea, unto every nation, kindred, tongue, and people, working mighty miracles … among the children of men according to their faith” (2 Nephi 26:13).

\intalkheader{III.}

For this audience—adults who hold the Melchizedek Priesthood and young men who will soon receive this power—I will concentrate my remarks on healing blessings involving the power of the priesthood. We have this priesthood power, and we should all be prepared to use it properly. Current increases in natural disasters and financial challenges show that we will need this power even more in the future than in the past.

Many scriptures teach that the servants of the Lord “shall lay hands on the sick, and they shall recover” (Mark 16:18). Miracles happen when the authority of the priesthood is used to bless the sick. I have experienced these miracles. As a boy and as a man I have seen healings as miraculous as any recorded in the scriptures, and so have many of you.

There are five parts to the use of priesthood authority to bless the sick: (1) the anointing, (2) the sealing of the anointing, (3) faith, (4) the words of the blessing, and (5) the will of the Lord.

The Old Testament frequently mentions anointing with oil as part of a blessing conferred by priesthood authority. Anointings were declared to be for sanctification and perhaps can also be seen as symbolic of the blessings to be poured out from heaven as a result of this sacred act.

In the New Testament we read that Jesus’s Apostles “anointed with oil many that were sick, and healed them” (Mark 6:13). The book of James teaches the role of anointing in connection with the other elements in a healing blessing by priesthood authority:

“Is any sick among you? let him call for the elders of the church; and let them pray over him, anointing him with oil in the name of the Lord:

“And the prayer of faith shall save the sick, and the Lord shall raise him up” (James 5:14–15).

When someone has been anointed by the authority of the Melchizedek Priesthood, the anointing is sealed by that same authority. To seal something means to affirm it, to make it binding for its intended purpose. When elders anoint a sick person and seal the anointing, they open the windows of heaven for the Lord to pour forth the blessing He wills for the person afflicted.

President Brigham Young taught: “When I lay hands on the sick, I expect the healing power and influence of God to pass through me to the patient, and the disease to give way. … When we are prepared, when we are holy vessels before the Lord, a stream of power from the Almighty can pass through the tabernacle of the administrator to the system of the patient, and the sick are made whole.”

Although we know of many cases where persons blessed by priesthood authority have been healed, we rarely refer to these healings in public meetings because modern revelation cautions us not to “boast [ourselves] of these things, neither speak them before the world; for these things are given unto you for your profit and for salvation” (D\&C 84:73).

Faith is essential for healing by the powers of heaven. The Book of Mormon even teaches that “if there be no faith among the children of men God can do no miracle among them” (Ether 12:12). In a notable article on administering to the sick, President Spencer W. Kimball said: “The need of faith is often underestimated. The ill one and the family often seem to depend wholly on the power of the priesthood and the gift of healing that they hope the administering brethren may have, whereas the greater responsibility is with him who is blessed. … The major element is the faith of the individual when that person is conscious and accountable. ‘Thy faith hath made thee whole’ [Matthew 9:22] was repeated so often by the Master that it almost became a chorus.”

President Kimball even suggested that “too frequent administrations may be an indication of lack of faith or of the ill one trying to pass the responsibility for faith development to the elders rather than self.” He told about a faithful sister who received a priesthood blessing. When asked the next day if she wished to be administered to again, she replied: “No, I have been anointed and administered to. The ordinance has been performed. It is up to me now to claim my blessing through my faith.”

Another part of a priesthood blessing is the words of blessing spoken by the elder after he seals the anointing. These words can be very important, but their content is not essential and they are not recorded on the records of the Church. In some priesthood blessings—like a patriarchal blessing—the words spoken are the essence of the blessing. But in a healing blessing it is the other parts of the blessing—the anointing, the sealing, faith, and the will of the Lord—that are the essential elements.

Ideally, the elder who officiates will be so in tune with the Spirit of the Lord that he will know and declare the will of the Lord in the words of the blessing. Brigham Young taught priesthood holders, “It is your privilege and duty to live so that you know when the word of the Lord is spoken to you and when the mind of the Lord is revealed to you.” When that happens, the spoken blessing is fulfilled literally and miraculously. On some choice occasions I have experienced that certainty of inspiration in a healing blessing and have known that what I was saying was the will of the Lord. However, like most who officiate in healing blessings, I have often struggled with uncertainty on the words I should say. For a variety of causes, every elder experiences increases and decreases in his level of sensitivity to the promptings of the Spirit. Every elder who gives a blessing is subject to influence by what he desires for the person afflicted. Each of these and other mortal imperfections can influence the words we speak.

Fortunately, the words spoken in a healing blessing are not essential to its healing effect. If faith is sufficient and if the Lord wills it, the afflicted person will be healed or blessed whether the officiator speaks those words or not. Conversely, if the officiator yields to personal desire or inexperience and gives commands or words of blessing in excess of what the Lord chooses to bestow according to the faith of the individual, those words will not be fulfilled. Consequently, brethren, no elder should ever hesitate to participate in a healing blessing because of fear that he will not know what to say. The words spoken in a healing blessing can edify and energize the faith of those who hear them, but the effect of the blessing is dependent upon faith and the Lord’s will, not upon the words spoken by the elder who officiated.

Young men and older men, please take special note of what I will say now. As we exercise the undoubted power of the priesthood of God and as we treasure His promise that He will hear and answer the prayer of faith, we must always remember that faith and the healing power of the priesthood cannot produce a result contrary to the will of Him whose priesthood it is. This principle is taught in the revelation directing that the elders of the Church shall lay their hands upon the sick. The Lord’s promise is that “he that hath faith in me to be healed, and is not appointed unto death, shall be healed” (D\&C 42:48; emphasis added). Similarly, in another modern revelation the Lord declares that when one “asketh according to the will of God … it is done even as he asketh” (D\&C 46:30).

From all of this we learn that even the servants of the Lord, exercising His divine power in a circumstance where there is sufficient faith to be healed, cannot give a priesthood blessing that will cause a person to be healed if that healing is not the will of the Lord.

As children of God, knowing of His great love and His ultimate knowledge of what is best for our eternal welfare, we trust in Him. The first principle of the gospel is faith in the Lord Jesus Christ, and faith means trust. I felt that trust in a talk my cousin gave at the funeral of a teenage girl who had died of a serious illness. He spoke these words, which first astonished me and then edified me: “I know it was the will of the Lord that she die. She had good medical care. She was given priesthood blessings. Her name was on the prayer roll in the temple. She was the subject of hundreds of prayers for her restoration to health. And I know that there is enough faith in this family that she would have been healed unless it was the will of the Lord to take her home at this time.” I felt that same trust in the words of the father of another choice girl whose life was taken by cancer in her teen years. He declared, “Our family’s faith is in Jesus Christ and is not dependent on outcomes.” Those teachings ring true to me. We do all that we can for the healing of a loved one, and then we trust in the Lord for the outcome.

I testify of the power of the priesthood of God, of the power of the prayer of faith, and of the truth of these principles. Most of all, I testify of the Lord Jesus Christ, whose servants we are, whose Resurrection gives us the assurance of immortality, and whose Atonement gives us the opportunity for eternal life, the greatest of all the gifts of God, in the name of Jesus Christ, amen.

\newpage
\settalk{Two Lines of Communication}
\setspeaker{Dallin H. Oaks}
\settalkdate{October 2010}
\talkheading{Two Lines of Communication}{Dallin H. Oaks}

Our Heavenly Father has given His children two lines of communication with Him—what we may call the personal line and the priesthood line. All should understand and be guided by both of these essential lines of communication.

\intalkheader{I. The Personal Line}

In the personal line we pray directly to our Heavenly Father, and He answers us by the channels He has established, without any mortal intermediary. We pray to our Heavenly Father in the name of Jesus Christ, and He answers us through His Holy Spirit and in other ways. The mission of the Holy Ghost is to testify of the Father and the Son (see John 15:26; 2 Nephi 31:18; 3 Nephi 28:11), to guide us into truth (see John 14:26; 16:13), and to show us all things we should do (see 2 Nephi 32:5). This personal line of communication with our Heavenly Father through His Holy Spirit is the source of our testimony of truth, of our knowledge, and of our personal guidance from a loving Heavenly Father. It is an essential feature of His marvelous gospel plan, which allows each one of His children to receive a personal witness of its truth.

The direct, personal channel of communication to our Heavenly Father through the Holy Ghost is based on worthiness and is so essential that we are commanded to renew our covenants by partaking of the sacrament each Sabbath day. In this way we qualify for the promise that we may always have His Spirit to be with us, to guide us.

On this personal line of communication with the Lord, our belief and practice is similar to that of those Christians who insist that human mediators between God and man are unnecessary because all have direct access to God under the principle Martin Luther espoused that is now known as “the priesthood of all believers.” I will say more of that later.

The personal line is of paramount importance in personal decisions and in the governance of the family. Unfortunately, some members of our Church underestimate the need for this direct, personal line. Responding to the undoubted importance of prophetic leadership—the priesthood line, which operates principally to govern heavenly communications on Church matters—some seek to have their priesthood leaders make personal decisions for them, decisions they should make for themselves by inspiration through their personal line. Personal decisions and family governance are principally a matter for the personal line.

I feel to add two other cautions we should remember in connection with this precious direct, personal line of communication with our Heavenly Father.

First, in its fulness the personal line does not function independent of the priesthood line. The gift of the Holy Ghost—the means of communication from God to man—is conferred by priesthood authority as authorized by those holding priesthood keys. It does not come merely by desire or belief. And the right to the continuous companionship of this Spirit needs to be affirmed each Sabbath as we worthily partake of the sacrament and renew our baptismal covenants of obedience and service.

Similarly, we cannot communicate reliably through the direct, personal line if we are disobedient to or out of harmony with the priesthood line. The Lord has declared that “the powers of heaven cannot be controlled nor handled only upon the principles of righteousness” (D\&C 121:36). Unfortunately, it is common for persons who are violating God’s commandments or disobedient to the counsel of their priesthood leaders to declare that God has revealed to them that they are excused from obeying some commandment or from following some counsel. Such persons may be receiving revelation or inspiration, but it is not from the source they suppose. The devil is the father of lies, and he is ever anxious to frustrate the work of God by his clever imitations.

\intalkheader{II. The Priesthood Line}

Unlike the personal line, in which our Heavenly Father communicates with us directly through the Holy Ghost, the priesthood line of communication has the additional and necessary intermediaries of our Savior, Jesus Christ; His Church; and His appointed leaders.

Because of what He accomplished by His atoning sacrifice, Jesus Christ has the power to prescribe the conditions we must fulfill to qualify for the blessings of His Atonement. That is why we have commandments and ordinances. That is why we make covenants. That is how we qualify for the promised blessings. They all come through the mercy and grace of the Holy One of Israel, “after all we can do” (2 Nephi 25:23).

During His earthly ministry, Jesus Christ conferred the authority of the priesthood that bears His name and He established a church that also bears His name. In this last dispensation, His priesthood authority was restored and His Church was reestablished through heavenly ministrations to the Prophet Joseph Smith. This restored priesthood and this reestablished Church are at the heart of the priesthood line.

The priesthood line is the channel by which God has spoken to His children through the scriptures in times past. And it is this line through which He currently speaks through the teachings and counsel of living prophets and apostles and other inspired leaders. This is the way we receive the required ordinances. This is the way we receive calls to service in His Church. His Church is the way and His priesthood is the power through which we are privileged to participate in those cooperative activities that are essential to accomplishing the Lord’s work. These include preaching the gospel, building temples and chapels, and helping the poor.

In respect to this priesthood line, our belief and practice is similar to the insistence of some Christians that authoritative ordinances (sacraments) are essential and must be performed by one authorized and empowered by Jesus Christ (see John 15:16). We believe the same but of course differ with other Christians on how we trace that authority.

Some members or former members of our Church fail to recognize the importance of the priesthood line. They underestimate the importance of the Church and its leaders and its programs. Relying entirely on the personal line, they go their own way, purporting to define doctrine and to direct competing organizations contrary to the teachings of prophet-leaders. In this they mirror the modern hostility to what is disparagingly called “organized religion.” Those who reject the need for organized religion reject the work of the Master, who established His Church and its officers in the meridian of time and who reestablished them in modern times.

Organized religion, established by divine authority, is essential, as the Apostle Paul taught:

“For the perfecting of the saints, for the work of the ministry, for the edifying of the body of Christ:

“Till we all come in the unity of the faith, and of the knowledge of the Son of God, unto a perfect man, unto the measure of the stature of the fulness of Christ” (Ephesians 4:12–13).

We should all remember the Lord’s declaration in modern revelation that the voice of the Lord’s servants is the voice of the Lord (see D\&C 1:38; 21:5; 68:4).

I feel to add two cautions we should remember in connection with reliance on the vital priesthood line.

First, the priesthood line does not supersede the need for the personal line. We all need a personal testimony of truth. As our faith develops, we necessarily rely on the words and faith of others, like our parents, teachers, or priesthood leaders (see D\&C 46:14). But if we are solely dependent on one particular priesthood leader or teacher for our personal testimony of the truth instead of getting that testimony through the personal line, we will be forever vulnerable to disillusionment by the action of that person. When it comes to a mature knowledge or testimony of the truth, we should not be dependent on a mortal mediator between us and our Heavenly Father.

Second, like the personal line, the priesthood line cannot function fully and properly in our behalf unless we are worthy and obedient. Many scriptures teach that if we persist in serious violations of the commandments of God, we are “cut off from his presence” (Alma 38:1). When that happens, the Lord and His servants are seriously inhibited in giving us spiritual help and we cannot obtain it for ourselves.

History provides us a vivid example of the importance of the Lord’s servants being in tune with the Spirit. The young Prophet Joseph Smith could not translate when he was angry or upset.

David Whitmer recalled: “One morning when he was getting ready to continue the translation, something went wrong about the house and he was put out about it. Something that Emma, his wife, had done. Oliver and I went up stairs, and Joseph came up soon after to continue the translation, but he could not do anything. He could not translate a single syllable. He went down stairs, out into the orchard and made supplication to the Lord; was gone about an hour—came back to the house, asked Emma’s forgiveness and then came up stairs where we were and the translation went on all right. He could do nothing save he was humble and faithful.”

\intalkheader{III. The Need for Both Lines}

I will conclude with further examples of the need for both of the lines our Heavenly Father has established for communication with His children. Both lines are essential to His purpose to bring about the immortality and eternal life of His children. An early scriptural account of this need is in Father Jethro’s counsel that Moses should not try to do so much. The people were waiting upon their priesthood leader from morning till night to “enquire of God” (Exodus 18:15) and also to “judge between one and another” (verse 16). We often note how Jethro counseled Moses to delegate by appointing judges to handle the personal conflicts (see verses 21–22). But Jethro also gave Moses counsel that illustrates the importance of the personal line: “Thou shalt teach them ordinances and laws, and shalt shew them the way wherein they must walk, and the work that they must do” (verse 20; emphasis added).

In other words, the Israelites who followed Moses should be taught not to bring every question to that priesthood leader. They should understand the commandments and seek inspiration to work out most problems for themselves.

Recent events in the nation of Chile illustrate the need for both lines. Chile suffered a devastating earthquake. Many of our members lost homes; some lost family members. Many lost confidence. Quickly—because our Church is prepared to respond to such disasters—food, shelter, and other material aid was provided. The Saints of Chile heard the voice of the Lord through His Church and its leaders responding to their material needs. But however sufficient the priesthood line, it was not enough. Each member needed to seek the Lord in prayer and receive the direct message of comfort and guidance that comes through the Holy Spirit to those who seek and listen.

Our missionary work is another example of the need for both lines. The men and women who are called to be missionaries are worthy and willing because of the teachings they have received through the priesthood line and the testimony they have received through the personal line. They are called through the priesthood line. Then, as representatives of the Lord and under the direction of His priesthood line, they teach investigators. Sincere seekers after truth listen, and the missionaries encourage them to pray to know the truth of the message for themselves through the personal line.

A final example applies these principles to the subject of priesthood authority in the family and the Church. All priesthood authority in the Church functions under the direction of one who holds the appropriate priesthood keys. This is the priesthood line. But the authority that presides in the family—whether father or single-parent mother—functions in family matters without the need to get authorization from anyone holding priesthood keys. That is like the personal line. Both lines must be functioning in our family life and in our personal lives if we are to have the growth and achieve the destiny identified in our Heavenly Father’s plan for His children.

We must use both the personal line and the priesthood line in proper balance to achieve the growth that is the purpose of mortal life. If personal religious practice relies too much on the personal line, individualism erases the importance of divine authority. If personal religious practice relies too much on the priesthood line, individual growth suffers. The children of God need both lines to achieve their eternal destiny. The restored gospel teaches both, and the restored Church provides both.

I testify of the Lord’s prophet, President Thomas S. Monson, who holds the keys that govern the priesthood line. I testify of the Lord Jesus Christ, whose Church this is. And I testify of the restored gospel, whose truth can be known by each of us through the precious personal line to our Heavenly Father. In the name of Jesus Christ, amen.

\newpage
\settalk{Desire}
\setspeaker{Dallin H. Oaks}
\settalkdate{April 2011}
\talkheading{Desire}{Dallin H. Oaks}

I have chosen to talk about the importance of desire. I hope each of us will search our hearts to determine what we really desire and how we rank our most important desires.

Desires dictate our priorities, priorities shape our choices, and choices determine our actions. The desires we act on determine our changing, our achieving, and our becoming.

First I speak of some common desires. As mortal beings we have some basic physical needs. Desires to satisfy these needs compel our choices and determine our actions. Three examples will demonstrate how we sometimes override these desires with other desires that we consider more important.

First, food. We have a basic need for food, but for a time that desire can be overridden by a stronger desire to fast.

Second, shelter. As a 12-year-old boy I resisted a desire for shelter because of my greater desire to fulfill a Boy Scout requirement to spend a night in the woods. I was one of several boys who left comfortable tents and found a way to construct a shelter and make a primitive bed from the natural materials we could find.

Third, sleep. Even this basic desire can be temporarily overridden by an even more important desire. As a young soldier in the Utah National Guard, I learned an example of this from a combat-seasoned officer.

In the early months of the Korean War, a Richfield Utah National Guard field artillery battery was called into active service. This battery, commanded by Captain Ray Cox, consisted of about 40 Mormon men. After additional training and reinforcement by reservists from elsewhere, they were sent to Korea, where they experienced some of the fiercest combat of that war. In one battle they had to repel a direct assault by hundreds of enemy infantry, the kind of attack that overran and destroyed other field artillery batteries.

What does this have to do with overcoming the desire for sleep? During one critical night, when enemy infantry had poured through the front lines and into the rear areas occupied by the artillery, the captain had the field telephone lines wired into his tent and ordered his numerous perimeter guards to phone him personally each hour on the hour all night long. This kept the guards awake, but it also meant that Captain Cox had scores of interruptions to his sleep. “How could you do that?” I asked him. His answer shows the power of an overriding desire.

“I knew that if we ever got home, I would be meeting the parents of those boys on the streets in our small town, and I didn’t want to face any of them if their son didn’t make it home because of anything I failed to do as his commander.”

What an example of the power of an overriding desire on priorities and on actions! What a powerful example for all of us who are responsible for the welfare of others—parents, Church leaders, and teachers!

As a conclusion to that illustration, early in the morning following his nearly sleepless night, Captain Cox led his men in a counterattack on the enemy infantry. They took over 800 prisoners and suffered only two wounded. Cox was decorated for bravery, and his battery received a Presidential Unit Citation for its extraordinary heroism. And, like Helaman’s stripling warriors (see Alma 57:25–26), they all made it home.

The Book of Mormon contains many teachings on the importance of desire.

After many hours of pleading with the Lord, Enos was told that his sins were forgiven. He then “began to feel a desire for the welfare of [his] brethren” (Enos 1:9). He wrote, “And … after I had prayed and labored with all diligence, the Lord said unto me: I will grant unto thee according to thy desires, because of thy faith” (verse 12). Note the three essentials that preceded the promised blessing: desire, labor, and faith.

In his sermon on faith, Alma teaches that faith can begin with “no more than [a] desire to believe” if we will “let this desire work in [us]” (Alma 32:27).

Another great teaching on desire, especially on what should be our ultimate desire, occurs in the experience of the Lamanite king being taught by the missionary Aaron. When Aaron’s teaching caught his interest, the king asked, “What shall I do that I may be born of God” and “have this eternal life?” (Alma 22:15). Aaron replied, “If thou desirest this thing, … if thou wilt repent of all thy sins, and will bow down before God, and call on his name in faith, believing that ye shall receive, then shalt thou receive the hope which thou desirest” (verse 16).

The king did so and in mighty prayer declared, “I will give away all my sins to know thee … and be saved at the last day” (verse 18). With that commitment and that identification of his ultimate desire, his prayer was answered miraculously.

The prophet Alma had a great desire to cry repentance to all people, but he came to understand that he should not desire the compelling power this would require because, he concluded, “a just God … granteth unto men according to their desire, whether it be unto death or unto life” (Alma 29:4). Similarly, in modern revelation the Lord declares that He “will judge all men according to their works, according to the desire of their hearts” (D\&C 137:9).

Are we truly prepared to have our Eternal Judge attach this enormous significance to what we really desire?

Many scriptures speak of what we desire in terms of what we seek. “He that seeketh me early shall find me, and shall not be forsaken” (D\&C 88:83). “Seek ye earnestly the best gifts” (D\&C 46:8). “He that diligently seeketh shall find” (1 Nephi 10:19). “Draw near unto me and I will draw near unto you; seek me diligently and ye shall find me; ask, and ye shall receive; knock, and it shall be opened unto you” (D\&C 88:63).

Readjusting our desires to give highest priority to the things of eternity is not easy. We are all tempted to desire that worldly quartet of property, prominence, pride, and power. We might desire these, but we should not fix them as our highest priorities.

Those whose highest desire is to acquire possessions fall into the trap of materialism. They fail to heed the warning “Seek not after riches nor the vain things of this world” (Alma 39:14; see also Jacob 2:18).

Those who desire prominence or power should follow the example of the valiant Captain Moroni, whose service was not “for power” or for the “honor of the world” (Alma 60:36).

How do we develop desires? Few will have the kind of crisis that motivated Aron Ralston, but his experience provides a valuable lesson about developing desires. While Ralston was hiking in a remote canyon in southern Utah, an 800-pound (360 kg) rock shifted suddenly and trapped his right arm. For five lonely days he struggled to free himself. When he was about to give up and accept death, he had a vision of a three-year-old boy running toward him and being scooped up with his left arm. Understanding this as a vision of his future son and an assurance that he could still live, Ralston summoned the courage and took drastic action to save his life before his strength ran out. He broke the two bones in his trapped right arm and then used the knife in his multitool to cut off that arm. He then summoned the strength to hike five miles (8 km) for help. What an example of the power of an overwhelming desire! When we have a vision of what we can become, our desire and our power to act increase enormously.

Most of us will never face such an extreme crisis, but all of us face potential traps that will prevent progress toward our eternal destiny. If our righteous desires are sufficiently intense, they will motivate us to cut and carve ourselves free from addictions and other sinful pressures and priorities that prevent our eternal progress.

We should remember that righteous desires cannot be superficial, impulsive, or temporary. They must be heartfelt, unwavering, and permanent. So motivated, we will seek for that condition described by the Prophet Joseph Smith, where we have “overcome the evils of [our lives] and lost every desire for sin.” That is a very personal decision. As Elder Neal A. Maxwell said:

“When people are described as ‘having lost their desire for sin,’ it is they, and they only, who deliberately decided to lose those wrong desires by being willing to ‘give away all [their] sins’ in order to know God.”

“Therefore, what we insistently desire, over time, is what we will eventually become and what we will receive in eternity.”

As important as it is to lose every desire for sin, eternal life requires more. To achieve our eternal destiny, we will desire and work for the qualities required to become an eternal being. For example, eternal beings forgive all who have wronged them. They put the welfare of others ahead of themselves. And they love all of God’s children. If this seems too difficult—and surely it is not easy for any of us—then we should begin with a desire for such qualities and call upon our loving Heavenly Father for help with our feelings. The Book of Mormon teaches us that we should “pray unto the Father with all the energy of heart, that [we] may be filled with this love, which he hath bestowed upon all who are true followers of his Son, Jesus Christ” (Moroni 7:48).

I close with a final example of a desire that should be paramount for all men and women—those who are currently married and those who are single. All should desire and seriously work to secure a marriage for eternity. Those who already have a temple marriage should do all they can to preserve it. Those who are single should desire a temple marriage and exert priority efforts to obtain it. Youth and young singles should resist the politically correct but eternally false concept that discredits the importance of marrying and having children.

Single men, please consider the challenge in this letter written by a single sister. She pleaded for “the righteous daughters of God that are sincerely searching for a worthy helpmeet, yet the men seem to be blinded and confused as to whether or not it is their responsibility to seek out these wonderful, choice daughters of our Heavenly Father and court them and be willing to make and keep sacred covenants in the Lord’s house.” She concluded, “There are many single LDS men here that are happy to go out and have fun, and date and hang out, but have absolutely no desire to ever make any kind of commitment to a woman.”

I am sure that some anxiously seeking young men would want me to add that there are some young women whose desires for a worthy marriage and children rank far below their desires for a career or other mortal distinctions. Both men and women need righteous desires that will lead them to eternal life.

Let us remember that desires dictate our priorities, priorities shape our choices, and choices determine our actions. In addition, it is our actions and our desires that cause us to become something, whether a true friend, a gifted teacher, or one who has qualified for eternal life.

I testify of Jesus Christ, whose love, whose teachings, and whose Atonement make it all possible. I pray that above all else we will desire to become like Him so that one day we can return to His presence to receive the fulness of His joy. In the name of Jesus Christ, amen.

\newpage
\settalk{Teachings of Jesus}
\setspeaker{Dallin H. Oaks}
\settalkdate{October 2011}
\talkheading{Teachings of Jesus}{Dallin H. Oaks}

“What think ye of Christ?” (Matthew 22:42). With those words Jesus confounded the Pharisees of His day. With those same words I ask my fellow Latter-day Saints and other Christians what you really believe about Jesus Christ and what you are doing because of that belief.

Most of my scriptural quotations will come from the Bible because it is familiar to most Christians. My interpretations will of course draw on what modern scripture, notably the Book of Mormon, teaches us about the meaning of Bible scriptures so ambiguous that different Christians disagree on their meaning. I address believers but others as well. As Elder Tad R. Callister taught us this morning, some who call themselves Christians praise Jesus as a great teacher but refrain from affirming His divinity. To address them, I have used the words of Jesus Himself. We should all consider what He Himself taught about who He is and what He was sent to earth to do.

\intalkheader{Only Begotten Son}

Jesus taught that He was the Only Begotten Son. Said He:

“For God so loved the world, that he gave his only begotten Son, that whosoever believeth in him should not perish, but have everlasting life.

“For God sent not his Son into the world to condemn the world; but that the world through him might be saved” (John 3:16–17).

God the Father affirmed this. In the culmination of the sacred experience on the Mount of Transfiguration, He declared from heaven, “This is my beloved Son, in whom I am well pleased; hear ye him” (Matthew 17:5).

Jesus also taught that His appearance was the same as His Father’s. To His Apostles, He said:

“If ye had known me, ye should have known my Father also: and from henceforth ye know him, and have seen him.

“Philip saith unto him, Lord, shew us the Father, and it sufficeth us.

“Jesus saith unto him, Have I been so long time with you, and yet hast thou not known me, Philip? he that hath seen me hath seen the Father” (John 14:7–9).

Later the Apostle Paul described the Son as being “the express image of [God the Father’s] person” (Hebrews 1:3; see also 2 Corinthians 4:4).

\intalkheader{Creator}

The Apostle John wrote that Jesus, whom he called “the Word,” “was in the beginning with God. All things were made by him; and without him was not any thing made that was made” (John 1:1–3). Thus, under the plan of the Father, Jesus Christ was the Creator of all things.

\intalkheader{Lord God of Israel}

During His ministry to His people in Palestine, Jesus taught that He was Jehovah, the Lord God of Israel (see John 8:58). Later, as the risen Lord, He ministered to His people on the American continent. There He declared:

“Behold, I am Jesus Christ, whom the prophets testified shall come into the world. …

“… I am the God of Israel, and the God of the whole earth” (3 Nephi 11:10, 14).

\intalkheader{What He Has Done for Us}

At a stake conference many years ago, I met a woman who said she had been asked to come back to church after many years away but could not think of any reason why she should. To encourage her I said, “When you consider all of the things the Savior has done for us, don’t you have many reasons to come back to church to worship and serve Him?” I was astonished at her reply: “What’s He done for me?” For those who do not understand what our Savior has done for us, I will answer that question in His own words and with my own testimony.

\intalkheader{Life of the World}

The Bible records Jesus’s teaching: “I am come that they might have life, and that they might have it more abundantly” (John 10:10). Later, in the New World, He declared, “I am the light and the life of the world” (3 Nephi 11:11). He is the “life of the world” because He is our Creator and because, through His Resurrection, we are all assured that we will live again. And the life He gives us is not merely mortal life. He taught, “I give unto them eternal life; and they shall never perish, neither shall any man pluck them out of my hand” (John 10:28; see also John 17:2).

\intalkheader{Light of the World}

Jesus also taught, “I am the light of the world: he that followeth me shall not walk in darkness” (John 8:12). He also declared, “I am the way, the truth, and the life” (John 14:6). He is the way and He is the light because His teachings light our path in mortal life and show us the way back to the Father.

\intalkheader{Doing the Will of the Father}

Always, Jesus honored the Father and followed Him. Even as a youth He declared to His earthly parents, “Wist ye not that I must be about my Father’s business?” (Luke 2:49). “For I came down from heaven,” He later taught, “not to do mine own will, but the will of him that sent me” (John 6:38; see also John 5:19). And the Savior taught, “No man cometh unto the Father, but by me” (John 14:6; see also Matthew 11:27).

We return to the Father by doing His will. Jesus taught, “Not every one that saith unto me, Lord, Lord, shall enter into the kingdom of heaven; but he that doeth the will of my Father which is in heaven” (Matthew 7:21). He explained:

“Many will say to me in that day, Lord, Lord, have we not prophesied in thy name? and in thy name have cast out devils? and in thy name done many wonderful works?

“And then will I profess unto them, I never knew you: depart from me, ye that work iniquity” (Matthew 7:22–23).

Who then will enter the kingdom of heaven? Not those who merely do wonderful works using the name of the Lord, Jesus taught, but only “he that doeth the will of my Father which is in heaven.”

\intalkheader{The Great Exemplar}

Jesus showed us how to do this. Again and again He invited us to follow Him: “My sheep hear my voice, and I know them, and they follow me” (John 10:27).

\intalkheader{Priesthood Power}

He gave priesthood power to His Apostles (see Matthew 10:1) and to others. To Peter, the senior Apostle, He said, “And I will give unto thee the keys of the kingdom of heaven: and whatsoever thou shalt bind on earth shall be bound in heaven: and whatsoever thou shalt loose on earth shall be loosed in heaven” (Matthew 16:19; see also Matthew 18:18).

Luke records that “the Lord appointed … seventy also, and sent them two and two before his face into every city and place, whither he himself would come” (Luke 10:1). Later these Seventy joyfully told Jesus, “Even the devils are subject unto us through thy name” (Luke 10:17). I am a witness of that priesthood power.

\intalkheader{Guidance by the Holy Ghost}

At the close of His earthly ministry, Jesus taught His Apostles, “The Comforter, which is the Holy Ghost, whom the Father will send in my name, he shall teach you all things, and bring all things to your remembrance, whatsoever I have said unto you” (John 14:26), and “he will guide you into all truth” (John 16:13).

\intalkheader{Guidance by His Commandments}

He also guides us by His commandments. Thus He commanded the Nephites that they should have no more disputes concerning points of doctrine, for, He said:

“He that hath the spirit of contention is not of me, but is of the devil, who is the father of contention, and he stirreth up the hearts of men to contend with anger, one with another.

“Behold, this is not my doctrine, to stir up the hearts of men with anger, one against another; but this is my doctrine, that such things should be done away” (3 Nephi 11:29–30).

\intalkheader{Focus on Eternal Life}

He also challenges us to focus on Him, not on the things of the world. In His great sermon on the bread of life, Jesus explained the contrast between mortal and eternal nourishment. “Labour not for the meat which perisheth,” He said, “but for that meat which endureth unto everlasting life, which the Son of man shall give unto you” (John 6:27). The Savior taught that He was the Bread of Life, the source of eternal nourishment. Speaking of the mortal nourishment the world offered, including the manna Jehovah had sent to feed the children of Israel in the wilderness, Jesus taught that those who relied on this bread were now dead (see John 6:49). In contrast, the nourishment He offered was “the living bread which came down from heaven,” and, Jesus taught, “if any man eat of this bread, he shall live for ever” (John 6:51).

Some of His disciples said this was “an hard saying,” and from that time many of His followers “went back, and walked no more with him” (John 6:60, 66). Apparently they did not accept His earlier teaching that they should “seek … first the kingdom of God” (Matthew 6:33). Even today some who profess Christianity are more attracted to the things of the world—the things that sustain life on earth but give no nourishment toward eternal life. For some, His “hard saying” is still a reason not to follow Christ.

\intalkheader{The Atonement}

The culmination of our Savior’s mortal ministry was His Resurrection and His Atonement for the sins of the world. John the Baptist prophesied this when he said, “Behold the Lamb of God, which taketh away the sin of the world” (John 1:29). Later Jesus taught that “the Son of man came … to minister, and to give his life a ransom for many” (Matthew 20:28). At the Last Supper, Jesus explained, according to the account in Matthew, that the wine He had blessed was “my blood of the new testament, which is shed for many for the remission of sins” (Matthew 26:28).

Appearing to the Nephites, the risen Lord invited them to come forward to feel the wound in His side and the prints of the nails in His hands and His feet. He did this, He explained, “that ye may know that I am the God of Israel, and the God of the whole earth, and have been slain for the sins of the world” (3 Nephi 11:14). And, the account continues, the multitude fell “down at the feet of Jesus, and did worship him” (verse 17). For this, the whole world will ultimately worship Him.

Jesus taught further precious truths about His Atonement. The Book of Mormon, which elaborates the Savior’s teachings and gives the best explanation of His mission, reports this teaching:

“My Father sent me that I might be lifted up upon the cross … , that I might draw all men unto me, …

“… that they may be judged according to their works.

“And … whoso repenteth and is baptized in my name shall be filled; and if he endureth to the end, behold, him will I hold guiltless before my Father at that day when I shall stand to judge the world. …

“And no unclean thing can enter into [the Father’s] kingdom; therefore nothing entereth into his rest save it be those who have washed their garments in my blood, because of their faith, and the repentance of all their sins, and their faithfulness unto the end” (3 Nephi 27:14–16, 19).

And so we understand that the Atonement of Jesus Christ gives us the opportunity to overcome the spiritual death that results from sin and, through making and keeping sacred covenants, to have the blessings of eternal life.

\intalkheader{Challenge and Testimony}

Jesus issued the challenge “What think ye of Christ?” (Matthew 22:42). The Apostle Paul challenged the Corinthians to “examine yourselves, whether ye be in the faith” (2 Corinthians 13:5). All of us should answer these challenges for ourselves. Where is our ultimate loyalty? Are we like the Christians in Elder Neal A. Maxwell’s memorable description who have moved their residence to Zion but still try to keep a second residence in Babylon?

There is no middle ground. We are followers of Jesus Christ. Our citizenship is in His Church and His gospel, and we should not use a visa to visit Babylon or act like one of its citizens. We should honor His name, keep His commandments, and “seek not the things of this world but seek … first to build up the kingdom of God, and to establish his righteousness” (Matthew 6:33, footnote a; from Joseph Smith Translation, Matthew 6:38).

Jesus Christ is the Only Begotten and Beloved Son of God. He is our Creator. He is the Light of the World. He is our Savior from sin and death. This is the most important knowledge on earth, and you can know this for yourself, as I know it for myself. The Holy Ghost, who testifies of the Father and the Son and leads us into truth, has revealed these truths to me, and He will reveal them to you. The way is desire and obedience. As to desire, Jesus taught, “Ask, and it shall be given you; seek, and ye shall find; knock, and it shall be opened unto you” (Matthew 7:7). As to obedience, He taught, “If any man will do his will, he shall know of the doctrine, whether it be of God, or whether I speak of myself” (John 7:17). I testify of the truth of these things in the name of Jesus Christ, amen.

\newpage
\settalk{Sacrifice}
\setspeaker{Dallin H. Oaks}
\settalkdate{April 2012}
\talkheading{Sacrifice}{Dallin H. Oaks}

The atoning sacrifice of Jesus Christ has been called the “most transcendent of all events from creation’s dawn to the endless ages of eternity.” That sacrifice is the central message of all the prophets. It was prefigured by the animal sacrifices prescribed by the law of Moses. A prophet declared that their whole meaning “point[ed] to that great and last sacrifice [of] … the Son of God, yea, infinite and eternal” (Alma 34:14). Jesus Christ endured incomprehensible suffering to make Himself a sacrifice for the sins of all. That sacrifice offered the ultimate good—the pure Lamb without blemish—for the ultimate measure of evil—the sins of the entire world. In the memorable words of Eliza R. Snow:

\begin{verse}
His precious blood he freely spilt;\\
His life he freely gave,\\
A sinless sacrifice for guilt,\\
A dying world to save.
\end{verse}

That sacrifice—the Atonement of Jesus Christ—is at the center of the plan of salvation.

The incomprehensible suffering of Jesus Christ ended sacrifice by the shedding of blood, but it did not end the importance of sacrifice in the gospel plan. Our Savior requires us to continue to offer sacrifices, but the sacrifices He now commands are that we “offer for a sacrifice unto [Him] a broken heart and a contrite spirit” (3 Nephi 9:20). He also commands each of us to love and serve one another—in effect, to offer a small imitation of His own sacrifice by making sacrifices of our own time and selfish priorities. In an inspired hymn, we sing, “Sacrifice brings forth the blessings of heaven.”

I will speak of these mortal sacrifices our Savior asks us to make. This will not include sacrifices we are compelled to make or actions that may be motivated by personal advantage rather than service or sacrifice (see 2 Nephi 26:29).

\intalkheader{I.}

The Christian faith has a history of sacrifice, including the ultimate sacrifice. In the early years of the Christian era, Rome martyred thousands for their faith in Jesus Christ. In later centuries, as doctrinal controversies divided Christians, some groups persecuted and even put to death the members of other groups. Christians killed by other Christians are the most tragic martyrs of the Christian faith.

Many Christians have voluntarily given sacrifices motivated by faith in Christ and the desire to serve Him. Some have chosen to devote their entire adult lives to the service of the Master. This noble group includes those in the religious orders of the Catholic Church and those who have given lifelong service as Christian missionaries in various Protestant faiths. Their examples are challenging and inspiring, but most believers in Christ are neither expected nor able to devote their entire lives to religious service.

\intalkheader{II.}

For most followers of Christ, our sacrifices involve what we can do on a day-to-day basis in our ordinary personal lives. In that experience I know of no group whose members make more sacrifices than Latter-day Saints. Their sacrifices—your sacrifices, my brothers and sisters—stand in contrast to the familiar worldly quests for personal fulfillment.

My first examples are our Mormon pioneers. Their epic sacrifices of lives, family relationships, homes, and comforts are at the foundation of the restored gospel. Sarah Rich spoke for what motivated these pioneers when she described her husband, Charles, being called away on a mission: “This truly was a trying time for me as well as for my husband; but duty called us to part for a season and knowing that we [were] obeying the will of the Lord, we felt to sacrifice our own feelings in order to help establish the work … of helping to build up the Kingdom of God on earth.”

Today the most visible strength of The Church of Jesus Christ of Latter-day Saints is the unselfish service and sacrifice of its members. Prior to the rededication of one of our temples, a Christian minister asked President Gordon B. Hinckley why it did not contain any representation of the cross, the most common symbol of the Christian faith. President Hinckley replied that the symbols of our Christian faith are “the lives of our people.” Truly, our lives of service and sacrifice are the most appropriate expressions of our commitment to serve the Master and our fellowmen.

\intalkheader{III.}

We have no professionally trained and salaried clergy in The Church of Jesus Christ of Latter-day Saints. As a result, the lay members who are called to lead and serve our congregations must carry the whole load of our numerous Church meetings, programs, and activities. They do this in more than 14,000 congregations just in the United States and Canada. Of course, we are not unique in having lay members of our congregations serve as teachers and lay leaders. But the amount of time donated by our members to train and minister to one another is uniquely large. Our efforts to have each family in our congregations visited by home teachers each month and to have each adult woman visited by Relief Society visiting teachers each month are examples of this. We know of no comparable service in any organization in the world.

The best-known examples of unique LDS service and sacrifice are the work of our missionaries. Currently they number more than 50,000 young men and young women and over 5,000 adult men and women. They devote from six months to two years of their lives to teaching the gospel of Jesus Christ and providing humanitarian service in more than 160 countries in the world. Their work always involves sacrifice, including the years they give to the work of the Lord and also the sacrifices made in providing funds for their support.

Those who remain at home—parents and other family members—also sacrifice by forgoing the companionship and service of the missionaries they send forth. For example, a young Brazilian received a missionary call while he was working to support his brothers and sisters after his father and mother died. A General Authority described these children’s meeting in council and remembering that their deceased parents had taught them that they should always be prepared to serve the Lord. The young man accepted his missionary call, and a 16-year-old brother took over the responsibility of working to support the family. Most of us know of many other examples of sacrifice to serve a mission or to support a missionary. We know of no other voluntary service and sacrifice like this in any other organization in the world.

We are frequently asked, “How do you persuade your young people and your older members to leave their schooling or their retirement to sacrifice in this way?” I have heard many give this explanation: “Knowing what my Savior did for me—His grace in suffering for my sins and in overcoming death so I can live again—I feel privileged to make the small sacrifice I am asked to make in His service. I want to share the understanding He has given me.” How do we persuade such followers of Christ to serve? As a prophet explained, “We [just] ask them.”

Other sacrifices resulting from missionary service are the sacrifices of those who act on the teachings of the missionaries and become members of the Church. For many converts, these sacrifices are very significant, including the loss of friends and family associations.

Many years ago this conference heard of a young man who found the restored gospel while he was studying in the United States. As this man was about to return to his native land, President Gordon B. Hinckley asked him what would happen to him when he returned home as a Christian. “My family will be disappointed,” the young man answered. “They may cast me out and regard me as dead. As for my future and my career, all opportunity may be foreclosed against me.”

“Are you willing to pay so great a price for the gospel?” President Hinckley asked.

Tearfully the young man answered, “It’s true, isn’t it?” When that was affirmed, he replied, “Then what else matters?” That is the spirit of sacrifice among many of our new members.

Other examples of service and sacrifice appear in the lives of the faithful members who serve in our temples. Temple service is unique to Latter-day Saints, but the significance of such sacrifice should be understandable to all Christians. Latter-day Saints have no tradition of service in a monastery, but we can still understand and honor the sacrifice of those whose Christian faith motivates them to devote their lives to that religious activity.

In this conference just a year ago, President Thomas S. Monson shared an example of sacrifice in connection with temple service. A faithful Latter-day Saint father on a remote island in the Pacific did heavy physical work in a faraway place for six years to earn the money necessary to take his wife and 10 children for marriage and sealing for eternity in the New Zealand Temple. President Monson explained, “Those who understand the eternal blessings which come from the temple know that no sacrifice is too great, no price too heavy, no struggle too difficult in order to receive those blessings.”

I am grateful for the marvelous examples of Christian love, service, and sacrifice I have seen among the Latter-day Saints. I see you performing your Church callings, often at great sacrifice of time and means. I see you serving missions at your own expense. I see you cheerfully donating your professional skills in service to your fellowmen. I see you caring for the poor through personal efforts and through supporting Church welfare and humanitarian contributions. All of this is affirmed in a nationwide study which concluded that active members of The Church of Jesus Christ of Latter-day Saints “volunteer and donate significantly more than the average American and are even more generous in time and money than the upper [20 percent] of religious people in America.”

Such examples of giving to others strengthen all of us. They remind us of the Savior’s teaching:

“If any man will come after me, let him deny himself. …

“For whosoever will save his life shall lose it: and whosoever will lose his life for my sake shall find it” (Matthew 16:24–25).

\intalkheader{IV.}

Perhaps the most familiar and most important examples of unselfish service and sacrifice are performed in our families. Mothers devote themselves to the bearing and nurturing of their children. Husbands give themselves to supporting their wives and children. The sacrifices involved in the eternally important service to our families are too numerous to mention and too familiar to need mention.

I also see unselfish Latter-day Saints adopting children, including those with special needs, and seeking to provide foster children the hope and opportunities denied them by earlier circumstances. I see you caring for family members and neighbors who suffer from birth defects, mental and physical ailments, and the effects of advancing years. The Lord sees you also, and He has caused His prophets to declare that “as you sacrifice for each other and your children, the Lord will bless you.”

I believe that Latter-day Saints who give unselfish service and sacrifice in worshipful imitation of our Savior adhere to eternal values to a greater extent than any other group of people. Latter-day Saints look on their sacrifices of time and means as a part of their schooling and qualifying for eternity. This is a truth revealed in the Lectures on Faith, which teach that “a religion that does not require the sacrifice of all things never has power sufficient to produce the faith necessary unto life and salvation. … It [is] through this sacrifice, and this only, that God has ordained that men should enjoy eternal life.”

Just as the atoning sacrifice of Jesus Christ is at the center of the plan of salvation, we followers of Christ must make our own sacrifices to prepare for the destiny that plan provides for us.

I know that Jesus Christ is the Only Begotten Son of God the Eternal Father. I know that because of His atoning sacrifice, we have the assurance of immortality and the opportunity for eternal life. He is our Lord, our Savior, and our Redeemer, and I testify of Him in the name of Jesus Christ, amen.

\newpage
\settalk{Protect the Children}
\setspeaker{Dallin H. Oaks}
\settalkdate{October 2012}
\talkheading{Protect the Children}{Dallin H. Oaks}

We can all remember our feelings when a little child cried out and reached up to us for help. A loving Heavenly Father gives us those feelings to impel us to help His children. Please recall those feelings as I speak about our responsibility to protect and act for the well-being of children.

I speak from the perspective of the gospel of Jesus Christ, including His plan of salvation. That is my calling. Local Church leaders have responsibility for a single jurisdiction, like a ward or stake, but an Apostle is responsible to witness to the entire world. In every nation, of every race and creed, all children are children of God.

Although I do not speak in terms of politics or public policy, like other Church leaders, I cannot speak for the welfare of children without implications for the choices being made by citizens, public officials, and workers in private organizations. We are all under the Savior’s command to love and care for each other and especially for the weak and defenseless.

Children are highly vulnerable. They have little or no power to protect or provide for themselves and little influence on so much that is vital to their well-being. Children need others to speak for them, and they need decision makers who put their well-being ahead of selfish adult interests.

\intalkheader{I.}

Worldwide, we are shocked at the millions of children victimized by evil adult crimes and selfishness.

In some war-torn countries, children are abducted to serve as soldiers in contending armies.

A United Nations report estimates that over two million children are victimized each year through prostitution and pornography.

From the perspective of the plan of salvation, one of the most serious abuses of children is to deny them birth. This is a worldwide trend. The national birthrate in the United States is the lowest in 25 years, and the birthrates in most European and Asian countries have been below replacement levels for many years. This is not just a religious issue. As rising generations diminish in numbers, cultures and even nations are hollowed out and eventually disappear.

One cause of the diminishing birthrate is the practice of abortion. Worldwide, there are estimated to be more than 40 million abortions per year. Many laws permit or even promote abortion, but to us this is a great evil. Other abuses of children that occur during pregnancy are the fetal impairments that result from the mother’s inadequate nutrition or drug use.

There is a tragic irony in the multitude of children eliminated or injured before birth while throngs of infertile couples long for and seek babies to adopt.

Childhood abuses or neglect of children that occur after birth are more publicly visible. Worldwide, almost eight million children die before their fifth birthday, mostly from diseases both treatable and preventable. And the World Health Organization reports that one in four children have stunted growth, mentally and physically, because of inadequate nutrition. Living and traveling internationally, we Church leaders see much of this. The general presidency of the Primary report children living in conditions “beyond our imaginations.” A mother in the Philippines said: “Sometimes we do not have enough money for food, but that is all right because it gives me the opportunity to teach my children about faith. We gather and pray for relief, and the children see the Lord bless us.” In South Africa, a Primary worker met a little girl, lonely and sad. In faint responses to loving questions, she said she had no mother, no father, and no grandmother—only a grandfather to care for her. Such tragedies are common on a continent where many caregivers have died of AIDS.

Even in rich nations little children and youth are impaired by neglect. Children growing up in poverty have inferior health care and inadequate educational opportunities. They are also exposed to dangerous environments in their physical and cultural surroundings and even from the neglect of their parents. Elder Jeffrey R. Holland recently shared the experience of an LDS police officer. In an investigation he found five young children huddled together and trying to sleep without bedding on a filthy floor in a dwelling where their mother and others were drinking and partying. The apartment had no food to relieve their hunger. After tucking the children into a makeshift bed, the officer knelt and prayed for their protection. As he walked toward the door, one of them, about six, pursued him, grabbed him by the hand, and pleaded, “Will you please adopt me?”

We remember our Savior’s teaching as He placed a little child before His followers and declared:

“And whoso shall receive one such little child in my name receiveth me.

“But whoso shall offend one of these little ones which believe in me, it were better for him that a millstone were hanged about his neck, and that he were drowned in the depth of the sea” (Matthew 18:5–6).

When we consider the dangers from which children should be protected, we should also include psychological abuse. Parents or other caregivers or teachers or peers who demean, bully, or humiliate children or youth can inflict harm more permanent than physical injury. Making a child or youth feel worthless, unloved, or unwanted can inflict serious and long-lasting injury on his or her emotional well-being and development. Young people struggling with any exceptional condition, including same-gender attraction, are particularly vulnerable and need loving understanding—not bullying or ostracism.

With the help of the Lord, we can repent and change and be more loving and helpful to children—our own and those around us.

\intalkheader{II.}

There are few examples of physical or emotional threats to children as important as those arising out of their relationships with their parents or guardians. President Thomas S. Monson has spoken of what he called the “vile deeds” of child abuse, where a parent has broken or disfigured a child, physically or emotionally. I grieved as I had to study the shocking evidence of such cases during my service on the Utah Supreme Court.

Of utmost importance to the well-being of children is whether their parents were married, the nature and duration of the marriage, and, more broadly, the culture and expectations of marriage and child care where they live. Two scholars of the family explain: “Throughout history, marriage has first and foremost been an institution for procreation and raising children. It has provided the cultural tie that seeks to connect the father to his children by binding him to the mother of his children. Yet in recent times, children have increasingly been pushed from center stage.”

A Harvard law professor describes the current law and attitude toward marriage and divorce: “The [current] American story about marriage, as told in the law and in much popular literature, goes something like this: marriage is a relationship that exists primarily for the fulfillment of the individual spouses. If it ceases to perform this function, no one is to blame and either spouse may terminate it at will. … Children hardly appear in the story; at most they are rather shadowy characters in the background.”

Our Church leaders have taught that looking “upon marriage as a mere contract that may be entered into at pleasure … and severed at the first difficulty … is an evil meriting severe condemnation,” especially where “children are made to suffer.” And children are impacted by divorces. Over half of the divorces in a recent year involved couples with minor children.

Many children would have had the blessing of being raised by both of their parents if only their parents had followed this inspired teaching in the family proclamation: “Husband and wife have a solemn responsibility to love and care for each other and for their children. … Parents have a sacred duty to rear their children in love and righteousness, to provide for their physical and spiritual needs, and to teach them to love and serve one another.” The most powerful teaching of children is by the example of their parents. Divorcing parents inevitably teach a negative lesson.

There are surely cases when a divorce is necessary for the good of the children, but those circumstances are exceptional. In most marital contests the contending parents should give much greater weight to the interests of the children. With the help of the Lord, they can do so. Children need the emotional and personal strength that come from being raised by two parents who are united in their marriage and their goals. As one who was raised by a widowed mother, I know firsthand that this cannot always be achieved, but it is the ideal to be sought whenever possible.

Children are the first victims of current laws permitting so-called “no-fault divorce.” From the standpoint of children, divorce is too easy. Summarizing decades of social science research, a careful scholar concluded that “the family structure that produces the best outcomes for children, on average, are two biological parents who remain married.” A New York Times writer noted “the striking fact that even as traditional marriage has declined in the United States … the evidence has mounted for the institution’s importance to the well-being of children.” That reality should give important guidance to parents and parents-to-be in their decisions involving marriage and divorce. We also need politicians, policy makers, and officials to increase their attention to what is best for children in contrast to the selfish interests of voters and vocal advocates of adult interests.

Children are also victimized by marriages that do not occur. Few measures of the welfare of our rising generation are more disturbing than the recent report that 41 percent of all births in the United States were to women who were not married. Unmarried mothers have massive challenges, and the evidence is clear that their children are at a significant disadvantage when compared with children raised by married parents.

Most of the children born to unmarried mothers—58 percent—were born to couples who were cohabitating. Whatever we may say about these couples’ forgoing marriage, studies show that their children suffer significant comparative disadvantages. For children, the relative stability of marriage matters.

We should assume the same disadvantages for children raised by couples of the same gender. The social science literature is controversial and politically charged on the long-term effect of this on children, principally because, as a New York Times writer observed, “same-sex marriage is a social experiment, and like most experiments it will take time to understand its consequences.”

\intalkheader{III.}

I have spoken for children—children everywhere. Some may reject some of these examples, but none should resist the plea that we unite to increase our concern for the welfare and future of our children—the rising generation.

We are speaking of the children of God, and with His powerful help, we can do more to help them. In this plea I address not only Latter-day Saints but also all persons of religious faith and others who have a value system that causes them to subordinate their own needs to those of others, especially to the welfare of children.

Religious persons are also conscious of the Savior’s New Testament teaching that pure little children are our role models of humility and teachableness:

“Verily I say unto you, Except ye be converted, and become as little children, ye shall not enter into the kingdom of heaven.

“Whosoever therefore shall humble himself as this little child, the same is greatest in the kingdom of heaven” (Matthew 18:3–4).

In the Book of Mormon we read of the risen Lord teaching the Nephites that they must repent and be baptized “and become as a little child” or they could not inherit the kingdom of God (3 Nephi 11:38; see also Moroni 8:10).

I pray that we will humble ourselves as little children and reach out to protect our little children, for they are the future for us, for our Church, and for our nations. In the name of Jesus Christ, amen.

\newpage
\settalk{Followers of Christ}
\setspeaker{Dallin H. Oaks}
\settalkdate{April 2013}
\talkheading{Followers of Christ}{Dallin H. Oaks}

One of our most beloved hymns, performed by the Tabernacle Choir this morning, begins with these words:

\begin{verse}
“Come, follow me,” the Savior said.\\
Then let us in his footsteps tread,\\
For thus alone can we be one\\
With God’s own loved, begotten Son.
\end{verse}

Those words, inspired by the Savior’s earliest invitation to His disciples (see Matthew 4:19), were written by John Nicholson, a Scottish convert. Like many of our early leaders, he had little formal schooling but a profound love for our Savior and the plan of salvation.

All of the messages of this conference help us follow in the footsteps of our Savior, whose example and teachings define the path for every follower of Jesus Christ.

Like all other Christians, members of The Church of Jesus Christ of Latter-day Saints study the life of our Savior as reported in the New Testament books of Matthew, Mark, Luke, and John. I will review examples and teachings contained in these four books of the Holy Bible and invite each of us and all other Christians to consider how this restored Church and each of us qualify as followers of Christ.

Jesus taught that baptism was necessary to enter the kingdom of God (see John 3:5). He began His ministry by being baptized (see Mark 1:9), and He and His followers baptized others (see John 3:22–26). We do likewise.

Jesus began His preaching by inviting His listeners to repent (see Matthew 4:17). That is still His servants’ message to the world.

Throughout His ministry Jesus gave commandments. And He taught, “If ye love me, keep my commandments” (John 14:15; see also verses 21, 23). He affirmed that keeping His commandments would require His followers to leave what He called “that which is highly esteemed among men” (Luke 16:15) and “the tradition of men” (Mark 7:8; see also verse 13). He also warned, “If ye were of the world, the world would love his own: but because ye are not of the world, but I have chosen you out of the world, therefore the world hateth you” (John 15:19). As the Apostle Peter later declared, the followers of Jesus were to be “a peculiar people” (1 Peter 2:9).

Latter-day Saints understand that we should not be “of the world” or bound to “the tradition of men,” but like other followers of Christ, we sometimes find it difficult to separate ourselves from the world and its traditions. Some model themselves after worldly ways because, as Jesus said of some whom He taught, “they loved the praise of men more than the praise of God” (John 12:43). These failures to follow Christ are too numerous and too sensitive to list here. They range all the way from worldly practices like political correctness and extremes in dress and grooming to deviations from basic values like the eternal nature and function of the family.

Jesus’s teachings were not meant to be theoretical. Always they were to be acted upon. Jesus taught, “Whosoever heareth these sayings of mine, and doeth them, I will liken him unto a wise man” (Matthew 7:24; see also Luke 11:28) and “Blessed is that servant, whom his lord when he cometh shall find so doing” (Matthew 24:46). In another beloved hymn we sing:

\begin{verse}
Savior, may I learn to love thee,\\
Walk the path that thou hast shown. …\\
Savior, may I learn to love thee—\\
Lord, I would follow thee.
\end{verse}

As Jesus taught, those who love Him will keep His commandments. They will be obedient, as President Thomas S. Monson taught this morning. Following Christ is not a casual or occasional practice but a continuous commitment and way of life that applies at all times and in all places. The Savior taught this principle and how we should be reminded and strengthened to follow it when He instituted the ordinance of the sacrament (communion, as others call it). We know from modern revelation that He commanded His followers to partake of the emblems in remembrance of Him (see Joseph Smith Translation, Matthew 26:22 [in Matthew 26:26, footnote c], 24 [in the Bible appendix]; Joseph Smith Translation, Mark 14:21–24 [in the Bible appendix]). Members of The Church of Jesus Christ of Latter-day Saints follow that commandment each week by attending a worship service in which we partake of the bread and water and covenant that we will always remember Him and keep His commandments.

Jesus taught that “men ought always to pray” (Luke 18:1). He also set that example, such as when He “continued all night in prayer to God” (Luke 6:12) before He called His Twelve Apostles. Like other Christians, we pray in all our worship services. We also pray for guidance, and we teach that we should have frequent personal prayers and daily kneeling prayers as a family. Like Jesus, we pray to our Father in Heaven, and we do so in the sacred name of Jesus Christ.

The Savior called Twelve Apostles to assist in His Church and gave them the keys and authority to carry on after His death (see Matthew 16:18–19; Mark 3:14–15; 6:7; Luke 6:13). The Church of Jesus Christ of Latter-day Saints, as the restored Church of Jesus Christ, follows this example in its organization and in its conferral of keys and authority on Apostles.

Some whom Jesus called to follow Him did not respond immediately but sought a delay to attend to proper family obligations. Jesus replied, “No man, having put his hand to the plough, and looking back, is fit for the kingdom of God” (Luke 9:62). Many Latter-day Saints practice the priority Jesus taught. This includes the wonderful example of thousands of senior missionaries and others who have left children and grandchildren to perform the missionary duties to which they have been called.

Jesus taught that God created male and female and that a man should leave his parents and cleave to his wife (see Mark 10:6–8). Our commitment to this teaching is well known.

In the familiar parable of the lost sheep, Jesus taught that we should go out of our way to seek after any of the flock who have strayed (see Matthew 18:11–14; Luke 15:3–7). As we know, President Thomas S. Monson has given great emphasis to this direction in his memorable example and teachings about rescuing our fellow men and women.

In our efforts to rescue and serve, we follow our Savior’s unique example and tender teachings about love: “Thou shalt love thy neighbour as thyself” (Matthew 22:39). He even commanded us to love our enemies (see Luke 6:27–28). And in His great teachings at the end of His mortal ministry, He said:

“A new commandment I give unto you, That ye love one another; as I have loved you, that ye also love one another.

“By this shall all men know that ye are my disciples, if ye have love one to another” (John 13:34–35).

As part of loving one another, Jesus taught that when we are wronged by persons, we should forgive them (see Matthew 18:21–35; Mark 11:25–26; Luke 6:37). While many struggle with this difficult commandment, we all know of inspiring examples of Latter-day Saints who have given loving forgiveness, even for the most serious wrongs. For example, Chris Williams drew upon his faith in Jesus Christ to forgive the drunken driver who caused the death of his wife and two of their children. Only two days after the tragedy and still deeply distraught, this forgiving man, then serving as one of our bishops, said, “As a disciple of Christ, I had no other choice.”

Most Christians give to the poor and the needy, as Jesus taught (see Matthew 25:31–46; Mark 14:7). In following this teaching of our Savior, The Church of Jesus Christ of Latter-day Saints and its members excel. Our members make generous contributions to charities and give personal service and other gifts to the poor and needy. In addition, our members fast for two meals each month and donate at least the cost of these meals as a fast offering, which our bishops and branch presidents use to help our needy members. Our fasting to help the hungry is an act of charity and, when done with pure intent, is a spiritual feast.

Less well known is our Church’s global humanitarian service. Using funds donated by generous members, The Church of Jesus Christ of Latter-day Saints sends food, clothing, and other essentials to relieve the suffering of adults and children all over the world. These humanitarian donations, totaling hundreds of millions of dollars in the last decade, are made without any consideration of religion, race, or nationality.

Our massive relief effort following the 2011 Japanese earthquake and tsunami provided \$13 million in cash and relief supplies. In addition, more than 31,000 Church-sponsored volunteers gave more than 600,000 hours of service. Our humanitarian assistance to the victims of Hurricane Sandy in the eastern United States included large donations of various resources, plus almost 300,000 hours of service in cleanup efforts by about 28,000 Church members. Among many other examples last year, we provided 300,000 pounds (136,000 kg) of clothing and shoes for the refugees in the African nation of Chad. During the last quarter century we have assisted nearly 30 million people in 179 countries. Truly, the people called “Mormons” know how to give to the poor and needy.

In His last biblical teaching, our Savior directed His followers to take His teachings to every nation and every creature. From the beginning of the Restoration, The Church of Jesus Christ of Latter-day Saints has sought to follow that teaching. Even when we were a poor and struggling new church with only a few thousand members, our early leaders sent missionaries across the oceans, east and west. As a people, we have continued to teach the Christian message until today our unique missionary program has more than 60,000 full-time missionaries, plus thousands more who serve part-time. We have missionaries in over 150 countries and territories worldwide.

As part of His great Sermon on the Mount, Jesus taught, “Be ye therefore perfect, even as your Father which is in heaven is perfect” (Matthew 5:48). The purpose of this teaching and the purpose of following our Savior is to come to the Father, whom our Savior referred to as “my Father, and your Father; and … my God, and your God” (John 20:17).

From modern revelation, unique to the restored gospel, we know that the commandment to seek perfection is part of God the Father’s plan for the salvation of His children. Under that plan we are all heirs of our heavenly parents. “We are the children of God,” the Apostle Paul taught, “and if children, then heirs; heirs of God, and joint-heirs with Christ” (Romans 8:16–17). This means, as we are told in the New Testament, that we are “heirs … of eternal life” (Titus 3:7) and that if we come to the Father, we are to “inherit all things” (Revelation 21:7)—all that He has—a concept our mortal minds can hardly grasp. But at least we can understand that achieving this ultimate destiny in eternity is possible only if we follow our Savior, Jesus Christ, who taught that “no man cometh unto the Father, but by me” (John 14:6). We seek to follow Him and become more like Him, here and hereafter. So it is that in the final verses of our hymn “Come, Follow Me,” we sing:

\begin{verse}
Is it enough alone to know\\
That we must follow him below,\\
While trav’ling thru this vale of tears?\\
No, this extends to holier spheres. …
\end{verse}

I testify of our Savior, Jesus Christ, whose teachings and example we seek to follow. He invites all of us who are heavy laden to come unto Him, to learn of Him, to follow Him, and thus to find rest to our souls (see Matthew 4:19; 11:28). I testify of the truth of His message and of the divine mission and authority of His restored Church in the name of Jesus Christ, amen.

\newpage
\settalk{No Other Gods}
\setspeaker{Dallin H. Oaks}
\settalkdate{October 2013}
\talkheading{No Other Gods}{Dallin H. Oaks}

The Ten Commandments are fundamental to the Christian and Jewish faiths. Given by God to the children of Israel through the prophet Moses, the first two of these commandments direct our worship and our priorities. In the first, the Lord commanded, “Thou shalt have no other gods before me” (Exodus 20:3). Centuries later, when Jesus was asked, “Which is the great commandment in the law?” He answered, “Thou shalt love the Lord thy God with all thy heart, and with all thy soul, and with all thy mind” (Matthew 22:36–37).

The second of the Ten Commandments elaborates the direction to have no other gods and identifies what should be the ultimate priority in our lives as His children. “Thou shalt not make unto thee any graven image, or any likeness of any thing” in the heavens or the earth (Exodus 20:4). The commandment then adds, “Thou shalt not bow down thyself to them, nor serve them” (Exodus 20:5). More than merely forbidding physical idols, this states a fundamental priority for all time. Jehovah explains, “For I the Lord thy God am a jealous God, … shewing mercy unto … them that love me, and keep my commandments” (Exodus 20:5–6). The meaning of jealous is revealing. Its Hebrew origin means “possessing sensitive and deep feelings” (Exodus 20:5, footnote b). Thus we offend God when we “serve” other gods—when we have other first priorities.

\intalkheader{I.}

What other priorities are being “served” ahead of God by persons—even religious persons—in our day? Consider these possibilities, all common in our world:

If none of these examples seems to apply to any one of us, we can probably suggest others that do. The principle is more important than individual examples. The principle is not whether we have other priorities. The question posed by the second commandment is “What is our ultimate priority?” Are we serving priorities or gods ahead of the God we profess to worship? Have we forgotten to follow the Savior who taught that if we love Him, we will keep His commandments? (see John 14:15). If so, our priorities have been turned upside down by the spiritual apathy and undisciplined appetites so common in our day.

\intalkheader{II.}

For Latter-day Saints, God’s commandments are based on and inseparable from God’s plan for His children—the great plan of salvation. This plan, sometimes called the “great plan of happiness” (Alma 42:8), explains our origin and destiny as children of God—where we came from, why we are here, and where we are going. The plan of salvation explains the purpose of creation and the conditions of mortality, including God’s commandments, the need for a Savior, and the vital role of mortal and eternal families. If we Latter-day Saints, who have been given this knowledge, do not establish our priorities in accord with this plan, we are in danger of serving other gods.

Knowledge of God’s plan for His children gives Latter-day Saints a unique perspective on marriage and family. We are correctly known as a family-centered church. Our theology begins with heavenly parents, and our highest aspiration is to attain the fulness of eternal exaltation. We know this is possible only in a family relationship. We know that the marriage of a man and a woman is necessary for the accomplishment of God’s plan. Only this marriage will provide the approved setting for mortal birth and to prepare family members for eternal life. We look on marriage and the bearing and nurturing of children as part of God’s plan and a sacred duty of those given the opportunity to do so. We believe that the ultimate treasures on earth and in heaven are our children and our posterity.

\intalkheader{III.}

Because of what we understand about the potentially eternal role of the family, we grieve at the sharply declining numbers of births and marriages in many Western countries whose historic cultures are Christian and Jewish. Responsible sources report the following:

In the midst of these concerning trends, we are also conscious that God’s plan is for all of His children and that God loves all of His children, everywhere. The first chapter of the Book of Mormon declares that God’s “power, and goodness, and mercy are over all the inhabitants of the earth” (1 Nephi 1:14). A later chapter declares that “he hath given [his salvation] free for all men” and that “all men are privileged the one like unto the other, and none are forbidden” (2 Nephi 26:27–28). Consequently, the scriptures teach that we are responsible to be compassionate and charitable (loving) toward all men (see 1 Thessalonians 3:12; 1 John 3:17; D\&C 121:45).

\intalkheader{IV.}

We are also respectful of the religious beliefs of all people, even of those increasing numbers who profess no belief in God. We know that through the God-given power of choice, many will hold beliefs contrary to ours, but we are hopeful that others will be equally respectful of our religious beliefs and understand that our beliefs compel us to some different choices and behaviors than theirs. For example, we believe that, as an essential part of His plan of salvation, God has established an eternal standard that sexual relations should occur only between a man and a woman who are married.

The power to create mortal life is the most exalted power God has given to His children. Its use was mandated by God’s first commandment to Adam and Eve (see Genesis 1:28), but other important commandments were given to forbid its misuse (see Exodus 20:14; 1 Thessalonians 4:3). The emphasis we place on the law of chastity is explained by our understanding of the purpose of our procreative powers in the accomplishment of God’s plan. Outside the bonds of marriage between a man and a woman, all uses of our procreative powers are to one degree or another sinful and contrary to God’s plan for the exaltation of His children.

The importance we attach to the law of chastity explains our commitment to the pattern of marriage that originated with Adam and Eve and has continued through the ages as God’s pattern for the procreative relationship between His sons and daughters and for the nurturing of His children. Fortunately, many persons affiliated with other denominations or organizations agree with us on the nature and importance of marriage, some on the basis of religious doctrine and others on the basis of what they deem best for society.

Our knowledge of God’s plan for His children explains why we are distressed that more and more children are born outside of marriage—currently 41 percent of all births in the United States—and that the number of couples living together without marriage has increased dramatically in the past half century. Five decades ago, only a tiny percentage of first marriages were preceded by cohabitation. Now cohabitation precedes 60 percent of marriages. And this is increasingly accepted, especially among teenagers. Recent survey data found about 50 percent of teenagers stating that out-of-wedlock childbearing was a “worthwhile lifestyle.”

\intalkheader{V.}

There are many political and social pressures for legal and policy changes to establish behaviors contrary to God’s decrees about sexual morality and contrary to the eternal nature and purposes of marriage and childbearing. These pressures have already authorized same-gender marriages in various states and nations. Other pressures would confuse gender or homogenize those differences between men and women that are essential to accomplish God’s great plan of happiness.

Our understanding of God’s plan and His doctrine gives us an eternal perspective that does not allow us to condone such behaviors or to find justification in the laws that permit them. And, unlike other organizations that can change their policies and even their doctrines, our policies are determined by the truths God has identified as unchangeable.

Our twelfth article of faith states our belief in being subject to civil authority and “in obeying, honoring, and sustaining the law.” But man’s laws cannot make moral what God has declared immoral. Commitment to our highest priority—to love and serve God—requires that we look to His law for our standard of behavior. For example, we remain under divine command not to commit adultery or fornication even when those acts are no longer crimes under the laws of the states or countries where we reside. Similarly, laws legalizing so-called “same-sex marriage” do not change God’s law of marriage or His commandments and our standards concerning it. We remain under covenant to love God and keep His commandments and to refrain from serving other gods and priorities—even those becoming popular in our particular time and place.

In this determination we may be misunderstood, and we may incur accusations of bigotry, suffer discrimination, or have to withstand invasions of our free exercise of religion. If so, I think we should remember our first priority—to serve God—and, like our pioneer predecessors, push our personal handcarts forward with the same fortitude they exhibited.

A teaching of President Thomas S. Monson applies to this circumstance. At this conference 27 years ago, he boldly declared: “Let us have the courage to defy the consensus, the courage to stand for principle. Courage, not compromise, brings the smile of God’s approval. Courage becomes a living and an attractive virtue when it is regarded not only as a willingness to die manfully, but as the determination to live decently. A moral coward is one who is afraid to do what he thinks is right because others will disapprove or laugh. Remember that all men have their fears, but those who face their fears with dignity have courage as well.”

I pray that we will not let the temporary challenges of mortality cause us to forget the great commandments and priorities we have been given by our Creator and our Savior. We must not set our hearts so much on the things of the world and aspire to the honors of men (see D\&C 121:35) that we stop trying to achieve our eternal destiny. We who know God’s plan for His children—we who have made covenants to participate in it—have a clear responsibility. We must never deviate from our paramount desire, which is to achieve eternal life. We must never dilute our first priority—to have no other gods and to serve no other priorities ahead of God the Father and His Son, our Savior, Jesus Christ.

May God help us to understand this priority and to be understood by others as we seek to pursue it in a wise and loving way, I pray in the name of Jesus Christ, amen.

\newpage
\settalk{The Keys and Authority of the Priesthood}
\setspeaker{Dallin H. Oaks}
\settalkdate{April 2014}
\talkheading{The Keys and Authority of the Priesthood}{Dallin H. Oaks}

\intalkheader{I.}

At this conference we have seen the release of some faithful brothers, and we have sustained the callings of others. In this rotation—so familiar in the Church—we do not “step down” when we are released, and we do not “step up” when we are called. There is no “up or down” in the service of the Lord. There is only “forward or backward,” and that difference depends on how we accept and act upon our releases and our callings. I once presided at the release of a young stake president who had given fine service for nine years and was now rejoicing in his release and in the new calling he and his wife had just received. They were called to be the nursery leaders in their ward. Only in this Church would that be seen as equally honorable!

\intalkheader{II.}

While addressing a women’s conference, Relief Society general president Linda K. Burton said, “We hope to instill within each of us a greater desire to better understand the priesthood.” That need applies to all of us, and I will pursue it by speaking of the keys and authority of the priesthood. Since these subjects are of equal concern to men and to women, I am pleased that these proceedings are broadcast and published for all members of the Church. Priesthood power blesses all of us. Priesthood keys direct women as well as men, and priesthood ordinances and priesthood authority pertain to women as well as men.

\intalkheader{III.}

President Joseph F. Smith described the priesthood as “the power of God delegated to man by which man can act in the earth for the salvation of the human family.” Other leaders have taught us that the priesthood “is the consummate power on this earth. It is the power by which the earth was created.” The scriptures teach that “this same Priesthood, which was in the beginning, shall be in the end of the world also” (Moses 6:7). Thus, the priesthood is the power by which we will be resurrected and proceed to eternal life.

The understanding we seek begins with an understanding of the keys of the priesthood. “Priesthood keys are the authority God has given to priesthood [holders] to direct, control, and govern the use of His priesthood on earth.” Every act or ordinance performed in the Church is done under the direct or indirect authorization of one holding the keys for that function. As Elder M. Russell Ballard has explained, “Those who have priesthood keys … literally make it possible for all who serve faithfully under their direction to exercise priesthood authority and have access to priesthood power.”

In the controlling of the exercise of priesthood authority, the function of priesthood keys both enlarges and limits. It enlarges by making it possible for priesthood authority and blessings to be available for all of God’s children. It limits by directing who will be given the authority of the priesthood, who will hold its offices, and how its rights and powers will be conferred. For example, a person who holds the priesthood is not able to confer his office or authority on another unless authorized by one who holds the keys. Without that authorization, the ordination would be invalid. This explains why a priesthood holder—regardless of office—cannot ordain a member of his family or administer the sacrament in his own home without authorization from the one who holds the appropriate keys.

With the exception of the sacred work that sisters do in the temple under the keys held by the temple president, which I will describe hereafter, only one who holds a priesthood office can officiate in a priesthood ordinance. And all authorized priesthood ordinances are recorded on the records of the Church.

Ultimately, all keys of the priesthood are held by the Lord Jesus Christ, whose priesthood it is. He is the one who determines what keys are delegated to mortals and how those keys will be used. We are accustomed to thinking that all keys of the priesthood were conferred on Joseph Smith in the Kirtland Temple, but the scripture states that all that was conferred there were “the keys of this dispensation” (D\&C 110:16). At general conference many years ago, President Spencer W. Kimball reminded us that there are other priesthood keys that have not been given to man on the earth, including the keys of creation and resurrection.

The divine nature of the limitations put upon the exercise of priesthood keys explains an essential contrast between decisions on matters of Church administration and decisions affecting the priesthood. The First Presidency and the Council of the First Presidency and Quorum of the Twelve, who preside over the Church, are empowered to make many decisions affecting Church policies and procedures—matters such as the location of Church buildings and the ages for missionary service. But even though these presiding authorities hold and exercise all of the keys delegated to men in this dispensation, they are not free to alter the divinely decreed pattern that only men will hold offices in the priesthood.

\intalkheader{IV.}

I come now to the subject of priesthood authority. I begin with the three principles just discussed: (1) priesthood is the power of God delegated to man to act for the salvation of the human family, (2) priesthood authority is governed by priesthood holders who hold priesthood keys, and (3) since the scriptures state that “all other authorities [and] offices in the church are appendages to this [Melchizedek] priesthood” (D\&C 107:5), all that is done under the direction of those priesthood keys is done with priesthood authority.

How does this apply to women? In an address to the Relief Society, President Joseph Fielding Smith, then President of the Quorum of the Twelve Apostles, said this: “While the sisters have not been given the Priesthood, it has not been conferred upon them, that does not mean that the Lord has not given unto them authority. … A person may have authority given to him, or a sister to her, to do certain things in the Church that are binding and absolutely necessary for our salvation, such as the work that our sisters do in the House of the Lord. They have authority given unto them to do some great and wonderful things, sacred unto the Lord, and binding just as thoroughly as are the blessings that are given by the men who hold the Priesthood.”

In that notable address, President Smith said again and again that women have been given authority. To the women he said, “You can speak with authority, because the Lord has placed authority upon you.” He also said that the Relief Society “[has] been given power and authority to do a great many things. The work which they do is done by divine authority.” And, of course, the Church work done by women or men, whether in the temple or in the wards or branches, is done under the direction of those who hold priesthood keys. Thus, speaking of the Relief Society, President Smith explained, “[The Lord] has given to them this great organization where they have authority to serve under the directions of the bishops of the wards … , looking after the interest of our people both spiritually and temporally.”

Thus, it is truly said that Relief Society is not just a class for women but something they belong to—a divinely established appendage to the priesthood.

We are not accustomed to speaking of women having the authority of the priesthood in their Church callings, but what other authority can it be? When a woman—young or old—is set apart to preach the gospel as a full-time missionary, she is given priesthood authority to perform a priesthood function. The same is true when a woman is set apart to function as an officer or teacher in a Church organization under the direction of one who holds the keys of the priesthood. Whoever functions in an office or calling received from one who holds priesthood keys exercises priesthood authority in performing her or his assigned duties.

Whoever exercises priesthood authority should forget about their rights and concentrate on their responsibilities. That is a principle needed in society at large. The famous Russian writer Aleksandr Solzhenitsyn is quoted as saying, “It is time … to defend not so much human rights as human obligations.” Latter-day Saints surely recognize that qualifying for exaltation is not a matter of asserting rights but a matter of fulfilling responsibilities.

\intalkheader{V.}

The Lord has directed that only men will be ordained to offices in the priesthood. But, as various Church leaders have emphasized, men are not “the priesthood.” Men hold the priesthood, with a sacred duty to use it for the blessing of all of the children of God.

The greatest power God has given to His sons cannot be exercised without the companionship of one of His daughters, because only to His daughters has God given the power “to be a creator of bodies … so that God’s design and the Great Plan might meet fruition.” Those are the words of President J. Reuben Clark.

He continued: “This is the place of our wives and of our mothers in the Eternal Plan. They are not bearers of the Priesthood; they are not charged with carrying out the duties and functions of the Priesthood; nor are they laden with its responsibilities; they are builders and organizers under its power, and partakers of its blessings, possessing the complement of the Priesthood powers and possessing a function as divinely called, as eternally important in its place as the Priesthood itself.”

In those inspired words, President Clark was speaking of the family. As stated in the family proclamation, the father presides in the family and he and the mother have separate responsibilities, but they are “obligated to help one another as equal partners.” Some years before the family proclamation, President Spencer W. Kimball gave this inspired explanation: “When we speak of marriage as a partnership, let us speak of marriage as a full partnership. We do not want our LDS women to be silent partners or limited partners in that eternal assignment! Please be a contributing and full partner.”

In the eyes of God, whether in the Church or in the family, women and men are equal, with different responsibilities.

I close with some truths about the blessings of the priesthood. Unlike priesthood keys and priesthood ordinations, the blessings of the priesthood are available to women and to men on the same terms. The gift of the Holy Ghost and the blessings of the temple are familiar illustrations of this truth.

In his insightful talk at BYU Education Week last summer, Elder M. Russell Ballard gave these teachings:

“Our Church doctrine places women equal to and yet different from men. God does not regard either gender as better or more important than the other. …

“When men and women go to the temple, they are both endowed with the same power, which is priesthood power. … Access to the power and the blessings of the priesthood is available to all of God’s children.”

I testify of the power and blessings of the priesthood of God, available for His sons and daughters alike. I testify of the authority of the priesthood, which functions throughout all of the offices and activities of The Church of Jesus Christ of Latter-day Saints. I testify of the divinely directed function of the keys of the priesthood, held and exercised in their fulness by our prophet/president, Thomas S. Monson. Finally and most important, I testify of our Lord and Savior, Jesus Christ, whose priesthood this is and whose servants we are, in the name of Jesus Christ, amen.

\newpage
\settalk{Loving Others and Living with Differences}
\setspeaker{Dallin H. Oaks}
\settalkdate{October 2014}
\talkheading{Loving Others and Living with Differences}{Dallin H. Oaks}

\intalkheader{I.}

In the concluding days of His mortal ministry, Jesus gave His disciples what He called “a new commandment” (John 13:34). Repeated three times, that commandment was simple but difficult: “Love one another, as I have loved you” (John 15:12; see also verse 17). The teaching to love one another had been a central teaching of the Savior’s ministry. The second great commandment was “love thy neighbour as thyself” (Matthew 22:39). Jesus even taught, “Love your enemies” (Matthew 5:44). But the commandment to love others as He had loved His flock was to His disciples—and is to us—a challenge that was unique. “Actually,” President Thomas S. Monson taught us last April, “love is the very essence of the gospel, and Jesus Christ is our Exemplar. His life was a legacy of love.”

Why is it so difficult to have Christlike love for one another? It is difficult because we must live among those who do not share our beliefs and values and covenant obligations. In His great Intercessory Prayer, offered just before His Crucifixion, Jesus prayed for His followers: “I have given them thy word; and the world hath hated them, because they are not of the world, even as I am not of the world” (John 17:14). Then, to the Father He pleaded, “I pray not that thou shouldest take them out of the world, but that thou shouldest keep them from the evil” (verse 15).

We are to live in the world but not be of the world. We must live in the world because, as Jesus taught in a parable, His kingdom is “like leaven,” whose function is to raise the whole mass by its influence (see Luke 13:21; Matthew 13:33; see also 1 Corinthians 5:6–8). His followers cannot do that if they associate only with those who share their beliefs and practices. But the Savior also taught that if we love Him, we will keep His commandments (see John 14:15).

\intalkheader{II.}

The gospel has many teachings about keeping the commandments while living among people with different beliefs and practices. The teachings about contention are central. When the resurrected Christ found the Nephites disputing over the manner of baptism, He gave clear directions on how this ordinance should be performed. Then He taught this great principle:

“There shall be no disputations among you, as there have hitherto been; neither shall there be disputations among you concerning the points of my doctrine, as there have hitherto been.

“For verily, verily I say unto you, he that hath the spirit of contention is not of me, but is of the devil, who is the father of contention, and he stirreth up the hearts of men to contend with anger, one with another.

“Behold, this is … my doctrine, that such things should be done away” (3 Nephi 11:28–30; emphasis added).

The Savior did not limit His warning against contention to those who were not keeping the commandment about baptism. He forbade contention by anyone. Even those who keep the commandments must not stir up the hearts of men to contend with anger. The “father of contention” is the devil; the Savior is the Prince of Peace.

Similarly, the Bible teaches that “wise men turn away wrath” (Proverbs 29:8). The early Apostles taught that we should “follow after the things [that] make for peace” (Romans 14:19) and “[speak] the truth in love” (Ephesians 4:15), “for the wrath of man worketh not the righteousness of God” (James 1:20). In modern revelation the Lord commanded that the glad tidings of the restored gospel should be declared “every man to his neighbor, in mildness and in meekness” (D\&C 38:41), “with all humility, … reviling not against revilers” (D\&C 19:30).

\intalkheader{III.}

Even as we seek to be meek and to avoid contention, we must not compromise or dilute our commitment to the truths we understand. We must not surrender our positions or our values. The gospel of Jesus Christ and the covenants we have made inevitably cast us as combatants in the eternal contest between truth and error. There is no middle ground in that contest.

The Savior showed the way when His adversaries confronted Him with the woman who had been “taken in adultery, in the very act” (John 8:4). When shamed with their own hypocrisy, the accusers withdrew and left Jesus alone with the woman. He treated her with kindness by declining to condemn her at that time. But He also firmly directed her to “sin no more” (John 8:11). Loving-kindness is required, but a follower of Christ—just like the Master—will be firm in the truth.

\intalkheader{IV.}

Like the Savior, His followers are sometimes confronted by sinful behavior, and today when they hold out for right and wrong as they understand it, they are sometimes called “bigots” or “fanatics.” Many worldly values and practices pose such challenges to Latter-day Saints. Prominent among these today is the strong tide that is legalizing same-sex marriage in many states and provinces in the United States and Canada and many other countries in the world. We also live among some who don’t believe in marriage at all. Some don’t believe in having children. Some oppose any restrictions on pornography or dangerous drugs. Another example—familiar to most believers—is the challenge of living with a nonbelieving spouse or family member or associating with nonbelieving fellow workers.

In dedicated spaces, like temples, houses of worship, and our own homes, we should teach the truth and the commandments plainly and thoroughly as we understand them from the plan of salvation revealed in the restored gospel. Our right to do so is protected by constitutional guarantees of freedom of speech and religion, as well as by the privacy that is honored even in countries without formal constitutional guarantees.

In public, what religious persons say and do involves other considerations. The free exercise of religion covers most public actions, but it is subject to qualifications necessary to accommodate the beliefs and practices of others. Laws can prohibit behavior that is generally recognized as wrong or unacceptable, like sexual exploitation, violence, or terrorist behavior, even when done by extremists in the name of religion. Less grievous behaviors, even though unacceptable to some believers, may simply need to be endured if legalized by what a Book of Mormon prophet called “the voice of the people” (Mosiah 29:26).

On the subject of public discourse, we should all follow the gospel teachings to love our neighbor and avoid contention. Followers of Christ should be examples of civility. We should love all people, be good listeners, and show concern for their sincere beliefs. Though we may disagree, we should not be disagreeable. Our stands and communications on controversial topics should not be contentious. We should be wise in explaining and pursuing our positions and in exercising our influence. In doing so, we ask that others not be offended by our sincere religious beliefs and the free exercise of our religion. We encourage all of us to practice the Savior’s Golden Rule: “Whatsoever ye would that men should do to you, do ye even so to them” (Matthew 7:12).

When our positions do not prevail, we should accept unfavorable results graciously and practice civility with our adversaries. In any event, we should be persons of goodwill toward all, rejecting persecution of any kind, including persecution based on race, ethnicity, religious belief or nonbelief, and differences in sexual orientation.

\intalkheader{V.}

I have spoken of general principles. Now I will speak of how those principles should apply in a variety of familiar circumstances in which the Savior’s teachings should be followed more faithfully.

I begin with what our young children learn in their play activities. Too often non-Mormons here in Utah have been offended and alienated by some of our members who will not allow their children to be friends with children of other faiths. Surely we can teach our children values and standards of behavior without having them distance themselves or show disrespect to any who are different.

Many teachers in church and school have grieved at the way some teenagers, including LDS youth, treat one another. The commandment to love one another surely includes love and respect across religious lines and also across racial, cultural, and economic lines. We challenge all youth to avoid bullying, insults, or language and practices that deliberately inflict pain on others. All of these violate the Savior’s command to love one another.

The Savior taught that contention is a tool of the devil. That surely teaches against some of the current language and practices of politics. Living with policy differences is essential to politics, but policy differences need not involve personal attacks that poison the process of government and punish participants. All of us should banish hateful communications and practice civility for differences of opinion.

The most important setting to forgo contention and practice respect for differences is in our homes and family relationships. Differences are inevitable—some minor and some major. As to major differences, suppose a family member is in a cohabitation relationship. That brings two important values into conflict—our love for the family member and our commitment to the commandments. Following the Savior’s example, we can show loving-kindness and still be firm in the truth by forgoing actions that facilitate or seem to condone what we know to be wrong.

I close with another example of a family relationship. At a stake conference in the Midwest about 10 years ago, I met a sister who told me that her nonmember husband had been accompanying her to church for 12 years but had never joined the Church. What should she do? she asked. I counseled her to keep doing all the right things and to be patient and kind with her husband.

About a month later she wrote me as follows: “Well, I thought that the 12 years was a good show of patience, but I didn’t know if I was being very kind about it. So, I practiced real hard for over a month, and he got baptized.”

Kindness is powerful, especially in a family setting. Her letter continued, “I am even trying to be kinder now because we are working on a temple sealing this year!”

Six years later she wrote me another letter: “My husband was [just] called and set apart as the bishop [of our ward].”

\intalkheader{VI.}

In so many relationships and circumstances in life, we must live with differences. Where vital, our side of these differences should not be denied or abandoned, but as followers of Christ we should live peacefully with others who do not share our values or accept the teachings upon which they are based. The Father’s plan of salvation, which we know by prophetic revelation, places us in a mortal circumstance where we are to keep His commandments. That includes loving our neighbors of different cultures and beliefs as He has loved us. As a Book of Mormon prophet taught, we must press forward, having “a love of God and of all men” (2 Nephi 31:20).

As difficult as it is to live in the turmoil surrounding us, our Savior’s command to love one another as He loves us is probably our greatest challenge. I pray that we may understand this and seek to live it in all of our relationships and activities, in the name of Jesus Christ, amen.

\newpage
\settalk{The Parable of the Sower}
\setspeaker{Dallin H. Oaks}
\settalkdate{April 2015}
\talkheading{The Parable of the Sower}{Dallin H. Oaks}

Subjects for general conference talks are assigned—not by mortal authority but by the impressions of the Spirit. Many subjects would address the mortal concerns we all share. But just as Jesus did not teach how to overcome the mortal challenges or political oppression of His day, He usually inspires His modern servants to speak about what we must do to reform our personal lives to prepare us to return to our heavenly home. On this Easter weekend I have felt impressed to talk about the precious and timeless teachings in one of the parables of Jesus.

The parable of the sower is one of a small number of parables reported in all three of the synoptic Gospels. It is also one of an even smaller group of parables Jesus explained to His disciples. The seed that was sown was “the word of the kingdom” (Matthew 13:19), “the word” (Mark 4:14), or “the word of God” (Luke 8:11)—the teachings of the Master and His servants.

The different soils on which the seeds fell represent different ways in which mortals receive and follow these teachings. Thus the seeds that “fell by the way side” (Mark 4:4) have not reached mortal soil where they might possibly grow. They are like teachings that fall upon a heart hardened or unprepared. I will say nothing more of these. My message concerns those of us who have committed to be followers of Christ. What do we do with the Savior’s teachings as we live our lives?

The parable of the sower warns us of circumstances and attitudes that can keep anyone who has received the seed of the gospel message from bringing forth a goodly harvest.

\intalkheader{I. Stony Ground, No Root}

Some seed “fell on stony ground, where it had not much earth; and immediately it sprang up, because it had no depth of earth: but when the sun was up, it was scorched; and because it had no root, it withered away” (Mark 4:5–6).

Jesus explained that this describes those “who, when they have heard the word, immediately receive it with gladness,” but because they “have no root in themselves, … when affliction or persecution ariseth for the word’s sake, immediately they are offended” (Mark 4:16–17).

What causes hearers to “have no root in themselves”? This is the circumstance of new members who are merely converted to the missionaries or to the many attractive characteristics of the Church or to the many great fruits of Church membership. Not being rooted in the word, they can be scorched and wither away when opposition arises. But even those raised in the Church—long-term members—can slip into a condition where they have no root in themselves. I have known some of these—members without firm and lasting conversion to the gospel of Jesus Christ. If we are not rooted in the teachings of the gospel and regular in its practices, any one of us can develop a stony heart, which is stony ground for spiritual seeds.

Spiritual food is necessary for spiritual survival, especially in a world that is moving away from belief in God and the absolutes of right and wrong. In an age dominated by the Internet, which magnifies messages that menace faith, we must increase our exposure to spiritual truth in order to strengthen our faith and stay rooted in the gospel.

Young people, if that teaching seems too general, here is a specific example. If the emblems of the sacrament are being passed and you are texting or whispering or playing video games or doing anything else to deny yourself essential spiritual food, you are severing your spiritual roots and moving yourself toward stony ground. You are making yourself vulnerable to withering away when you encounter tribulation like isolation, intimidation, or ridicule. And that applies to adults also.

Another potential destroyer of spiritual roots—accelerated by current technology but not unique to it—is the keyhole view of the gospel or the Church. This limited view focuses on a particular doctrine or practice or perceived deficiency in a leader and ignores the grand panorama of the gospel plan and the personal and communal fruits of its harvest. President Gordon B. Hinckley gave a vivid description of one aspect of this keyhole view. He told a BYU audience about political commentators “aflame with indignation” at a then-recent news event. “With studied art they poured out the sour vinegar of invective and anger. … Surely,” he concluded, “this is the age and place of the gifted pickle sucker.” In contrast, to be securely rooted in the gospel, we must be moderate and measured in criticism and seek always for the broader view of the majestic work of God.

\intalkheader{II. Thorns: The Cares of This World and the Deceitfulness of Riches}

Jesus taught that “some fell among thorns, and the thorns grew up, and choked it, and it yielded no fruit” (Mark 4:7). He explained that these are “such as hear the word, and the cares of this world, and the deceitfulness of riches, and the lusts of other things entering in, choke the word, and it becometh unfruitful” (Mark 4:18–19). This is surely a warning to be heeded by all of us.

I will speak first of the deceitfulness of riches. Wherever we are in our spiritual journey—whatever our state of conversion—we are all tempted by this. When attitudes or priorities are fixed on the acquisition, use, or possession of property, we call that materialism. So much has been said and written about materialism that little needs to be added here. Those who believe in what has been called the theology of prosperity are suffering from the deceitfulness of riches. The possession of wealth or significant income is not a mark of heavenly favor, and their absence is not evidence of heavenly disfavor. When Jesus told a faithful follower that he could inherit eternal life if he would only give all that he had to the poor (see Mark 10:17–24), He was not identifying an evil in the possession of riches but an evil in that follower’s attitude toward them. As we are all aware, Jesus praised the good Samaritan, who used the same coinage to serve his fellowman that Judas used to betray his Savior. The root of all evil is not money but the love of money (see 1 Timothy 6:10).

The Book of Mormon tells of a time when the Church of God “began to fail in its progress” (Alma 4:10) because “the people of the church began to … set their hearts upon riches and upon the vain things of the world” (Alma 4:8). Whoever has an abundance of material things is in jeopardy of being spiritually “sedated” by riches and other things of the world. That is a suitable introduction to the next of the Savior’s teachings.

The most subtle thorns to choke out the effect of the gospel word in our lives are the worldly forces that Jesus called the “cares and riches and pleasures of this life” (Luke 8:14). These are too numerous to recite. Some examples will suffice.

On one occasion Jesus rebuked His chief Apostle, saying to Peter, “Thou art an offence unto me: for thou savourest not the things that be of God, but those that be of men” (Matthew 16:23; see also D\&C 3:6–7; 58:39). Savoring the things of men means putting the cares of this world ahead of the things of God in our actions, our priorities, and our thinking.

We surrender to the “pleasures of this life” (1) when we are addicted, which impairs God’s precious gift of agency; (2) when we are beguiled by trivial distractions, which draw us away from things of eternal importance; and (3) when we have an entitlement mentality, which impairs the personal growth necessary to qualify us for our eternal destiny.

We are overcome by the “cares … of this life” when we are paralyzed by fear of the future, which hinders our going forward in faith, trusting in God and His promises. Twenty-five years ago my esteemed BYU teacher Hugh W. Nibley spoke of the dangers of surrendering to the cares of the world. He was asked in an interview whether world conditions and our duty to spread the gospel made it desirable to seek some way to “be accommodating of the world in what we do in the Church.”

His reply: “That’s been the whole story of the Church, hasn’t it? You have to be willing to offend here, you have to be willing to take the risk. That’s where the faith comes in. … Our commitment is supposed to be a test, it’s supposed to be hard, it’s supposed to be impractical in the terms of this world.”

This gospel priority was affirmed on the BYU campus just a few months ago by an esteemed Catholic leader, Charles J. Chaput, the archbishop of Philadelphia. Speaking of “concerns that the LDS and Catholic communities share,” such as “about marriage and family, the nature of our sexuality, the sanctity of human life, and the urgency of religious liberty,” he said this:

“I want to stress again the importance of really living what we claim to believe. That needs to be a priority—not just in our personal and family lives but in our churches, our political choices, our business dealings, our treatment of the poor; in other words, in everything we do.”

“Here’s why that’s important,” he continued. “Learn from the Catholic experience. We Catholics believe that our vocation is to be leaven in society. But there’s a fine line between being leaven in society, and being digested by society.”

The Savior’s warning against having the cares of this world choke out the word of God in our lives surely challenges us to keep our priorities fixed—our hearts set—on the commandments of God and the leadership of His Church.

The Savior’s examples could cause us to think of this parable as the parable of the soils. The suitability of the soil depends upon the heart of each one of us who is exposed to the gospel seed. In susceptibility to spiritual teachings, some hearts are hardened and unprepared, some hearts are stony from disuse, and some hearts are set upon the things of the world.

\intalkheader{III. Fell into Good Ground and Brought Forth Fruit}

The parable of the sower ends with the Savior’s description of the seed that “fell into good ground, and brought forth fruit” in various measures (Matthew 13:8). How can we prepare ourselves to be that good ground and to have that good harvest?

Jesus explained that “the good ground are they, which in an honest and good heart, having heard the word, keep it, and bring forth fruit with patience” (Luke 8:15). We have the seed of the gospel word. It is up to each of us to set the priorities and to do the things that make our soil good and our harvest plentiful. We must seek to be firmly rooted and converted to the gospel of Jesus Christ (see Colossians 2:6–7). We achieve this conversion by praying, by scripture reading, by serving, and by regularly partaking of the sacrament to always have His Spirit to be with us. We must also seek that mighty change of heart (see Alma 5:12–14) that replaces evil desires and selfish concerns with the love of God and the desire to serve Him and His children.

I testify of the truth of these things, and I testify of our Savior, Jesus Christ, whose teachings point the way and whose Atonement makes it all possible, in the name of Jesus Christ, amen.

\newpage
\settalk{Strengthened by the Atonement of Jesus Christ}
\setspeaker{Dallin H. Oaks}
\settalkdate{October 2015}
\talkheading{Strengthened by the Atonement of Jesus Christ}{Dallin H. Oaks}

In mortality we have the certainty of death and the burden of sin. The Atonement of Jesus Christ offsets these two certainties of mortal life. But apart from death and sin, we have many other challenges as we struggle through mortality. Because of that same Atonement, our Savior can provide us the strength we need to overcome these mortal challenges. That is my subject today.

\intalkheader{I.}

Most scriptural accounts of the Atonement concern the Savior’s breaking the bands of death and suffering for our sins. In his sermon recorded in the Book of Mormon, Alma taught these fundamentals. But he also provided our clearest scriptural assurances that the Savior also experienced the pains and sicknesses and infirmities of His people.

Alma described this part of the Savior’s Atonement: “And he shall go forth, suffering pains and afflictions and temptations of every kind; and this that the word might be fulfilled which saith he will take upon him the pains and the sicknesses of his people” (Alma 7:11; also see 2 Nephi 9:21).

Think of it! In the Savior’s Atonement, He suffered “pains and afflictions and temptations of every kind.” As President Boyd K. Packer explained: “He had no debt to pay. He had committed no wrong. Nevertheless, an accumulation of all of the guilt, the grief and sorrow, the pain and humiliation, all of the mental, emotional, and physical torments known to man—He experienced them all.”

Why did He suffer these mortal challenges “of every kind”? Alma explained, “And he will take upon him their infirmities, that his bowels may be filled with mercy, according to the flesh, that he may know according to the flesh how to succor his people according to their infirmities” (Alma 7:12).

For example, the Apostle Paul declared that because the Savior “hath suffered being tempted, he is able to succour them that are tempted” (Hebrews 2:18). Similarly, President James E. Faust taught, “Since the Savior has suffered anything and everything that we could ever feel or experience, He can help the weak to become stronger.”

Our Savior experienced and suffered the fulness of all mortal challenges “according to the flesh” so He could know “according to the flesh” how to “succor [which means to give relief or aid to] his people according to their infirmities.” He therefore knows our struggles, our heartaches, our temptations, and our suffering, for He willingly experienced them all as an essential part of His Atonement. And because of this, His Atonement empowers Him to succor us—to give us the strength to bear it all.

\intalkheader{II.}

While Alma’s teaching in the seventh chapter is the single clearest of all the scriptures on this essential power of the Atonement, it is also taught throughout holy writ.

At the beginning of His ministry, Jesus explained that He was sent “to heal the brokenhearted” (Luke 4:18). The Bible often tells us of His healing people “of their infirmities” (Luke 5:15; 7:21). The Book of Mormon records His healing those “that were afflicted in any manner” (3 Nephi 17:9). The Gospel of Matthew explains that Jesus healed the people “that it might be fulfilled which was spoken by Esaias the prophet, saying, Himself took our infirmities, and bare our sicknesses” (Matthew 8:17).

Isaiah taught that the Messiah would bear our “griefs” and our “sorrows” (Isaiah 53:4). Isaiah also taught of His strengthening us: “Fear thou not; for I am with thee: be not dismayed; for I am thy God: I will strengthen thee; yea, I will help thee” (Isaiah 41:10).

Thus, we sing:

\begin{verse}
Fear not, I am with thee; oh, be not dismayed,\\
For I am thy God and will still give thee aid.\\
I’ll strengthen thee, help thee, and cause thee to stand, …\\
Upheld by my righteous, omnipotent hand.
\end{verse}

Speaking of some of his own mortal challenges, the Apostle Paul wrote, “I can do all things through Christ which strengtheneth me” (Philippians 4:13).

And so we see that because of His Atonement, the Savior has the power to succor—to help—every mortal pain and affliction. Sometimes His power heals an infirmity, but the scriptures and our experiences teach that sometimes He succors or helps by giving us the strength or patience to endure our infirmities.

\intalkheader{III.}

What are these mortal pains and afflictions and infirmities that our Savior experienced and suffered?

We all have pains and afflictions and infirmities at one time or another. Apart from what we experience because of our sins, mortality is filled with frequent struggles, heartaches, and suffering.

We and those we love suffer sickness. At some time each of us also experiences pain from traumatic injuries or from other physical or mental difficulties. All of us suffer and grieve in connection with the death of a loved one. We all experience failure in our personal responsibilities, our family relationships, or our occupations.

When a spouse or a child rejects what we know to be true and strays from the path of righteousness, we experience particularly stressful pain, just like the father of the prodigal son in Jesus’s memorable parable (see Luke 15:11–32).

As the Psalmist declared, “Many are the afflictions of the righteous: but the Lord delivereth him out of them all” (Psalm 34:19).

Thus, our hymns contain this true assurance: “Earth has no sorrow that heav’n cannot cure.” What cures us is our Savior and His Atonement.

Particularly poignant for teenagers is the feeling of rejection when peers seem to join in happy relationships and activities and deliberately leave them out. Racial and ethnic prejudices produce other painful rejections, for youth and adults. Life has many other challenges, such as unemployment or other reverses in our plans.

I am still speaking of mortal infirmities not caused by our sins. Some are born with physical or mental disabilities that cause personal suffering for them and struggles for those who love and care for them. For many, the infirmity of depression is painful or permanently disabling. Another painful affliction is the circumstance of singleness. Those who suffer this circumstance should remember that our Savior experienced this kind of pain also and that, through His Atonement, He offers the strength to bear it.

Few disabilities are more crippling to our temporal or spiritual lives than addictions. Some of these, like addictions to pornography or drugs, are likely to have been caused by sinful behavior. Even when that behavior has been repented of, the addiction may remain. That disabling grip can also be relieved by the decisive strength available from the Savior. So can the severe challenge experienced by those sent to prison for crimes. A recent letter testifies to the strength that can come even to one in that circumstance: “I know that our Savior is walking these halls, and I have often felt Christ’s love within these prison walls.”

I love the testimony of our poetess and friend Emma Lou Thayne. In words we now sing as a hymn, she wrote:

\begin{verse}
Where can I turn for peace?\\
Where is my solace\\
When other sources cease to make me whole?\\
When with a wounded heart, anger, or malice,\\
I draw myself apart,\\
Searching my soul?
\end{verse}

\intalkheader{IV.}

Who can be succored and strengthened through the Atonement of Jesus Christ? Alma taught that the Savior would “take upon him the pains and the sicknesses of his people” and “succor his people” (Alma 7:11, 12; emphasis added). Who are “his people” in this promise? Is it all mortals—all who enjoy the reality of resurrection through the Atonement? Or is it only those select servants qualified through ordinances and covenants?

The word people has many meanings in the scriptures. The meaning most appropriate for the teaching that the Savior will succor “his people” is the meaning Ammon employed when he taught that “God is mindful of every people, whatsoever land they may be in” (Alma 26:37). That is also what the angels meant when they announced the birth of the Christ child: “Good tidings of great joy, which shall be to all people” (Luke 2:10).

Because of His atoning experience in mortality, our Savior is able to comfort, heal, and strengthen all men and women everywhere, but I believe He does so only for those who seek Him and ask for His help. The Apostle James taught, “Humble yourselves in the sight of the Lord, and he shall lift you up” (James 4:10). We qualify for that blessing when we believe in Him and pray for His help.

There are millions of God-fearing people who pray to God to be lifted out of their afflictions. Our Savior has revealed that He “descended below all things” (D\&C 88:6). As Elder Neal A. Maxwell taught, “Having ‘descended below all things,’ He comprehends, perfectly and personally, the full range of human suffering.” We might even say that having descended beneath it all, He is perfectly positioned to lift us and give us the strength we need to endure our afflictions. We have only to ask.

Many times in modern revelation, the Lord declares, “Therefore, if you will ask of me you shall receive; if you will knock it shall be opened unto you” (for example, D\&C 6:5; 11:5; see also Matthew 7:7). Indeed, because of Their all-encompassing love, our Heavenly Father and His Beloved Son, Jesus Christ, hear and suitably answer the prayers of all who seek Them in faith. As the Apostle Paul wrote, “We trust in the living God, who is the Saviour of all men, specially of those that believe” (1 Timothy 4:10).

I know these things to be true. Our Savior’s Atonement does more than assure us of immortality by a universal resurrection and give us the opportunity to be cleansed from sin by repentance and baptism. His Atonement also provides the opportunity to call upon Him who has experienced all of our mortal infirmities to give us the strength to bear the burdens of mortality. He knows of our anguish, and He is there for us. Like the good Samaritan, when He finds us wounded at the wayside, He will bind up our wounds and care for us (see Luke 10:34). The healing and strengthening power of Jesus Christ and His Atonement is for all of us who will ask. I testify of that as I also testify of our Savior, who makes it all possible.

One day all of these mortal burdens will pass away and there will be no more pain (see Revelation 21:4). I pray that we will all understand the hope and strength of our Savior’s Atonement: the assurance of immortality, the opportunity for eternal life, and the sustaining strength we can receive if only we will ask, in the name of Jesus Christ, amen.

\newpage
\settalk{Opposition in All Things}
\setspeaker{Dallin H. Oaks}
\settalkdate{April 2016}
\talkheading{Opposition in All Things}{Dallin H. Oaks}

Central to the gospel of Jesus Christ is the Father’s plan of salvation for the eternal progress of His children. That plan, explained in modern revelation, helps us understand many things we face in mortality. My message focuses on the essential role of opposition in that plan.

\intalkheader{I.}

The purpose of mortal life for the children of God is to provide the experiences needed “to progress toward perfection and ultimately realize their divine destiny as heirs of eternal life.” As President Thomas S. Monson taught us so powerfully this morning, we progress by making choices, by which we are tested to show that we will keep God’s commandments (see Abraham 3:25). To be tested, we must have the agency to choose between alternatives. To provide alternatives on which to exercise our agency, we must have opposition.

The rest of the plan is also essential. When we make wrong choices—as we inevitably will—we are soiled by sin and must be cleansed to proceed toward our eternal destiny. The Father’s plan provides the way to do this, the way to satisfy the eternal demands of justice: a Savior pays the price to redeem us from our sins. That Savior is the Lord Jesus Christ, the Only Begotten Son of God the Eternal Father, whose atoning sacrifice—whose suffering—pays the price for our sins if we will repent of them.

One of the best explanations of the planned role of opposition is in the Book of Mormon, in Lehi’s teachings to his son Jacob.

“It must needs be, that there is an opposition in all things. If not so, … righteousness could not be brought to pass, neither wickedness, neither holiness nor misery, neither good nor bad” (2 Nephi 2:11; see also verse 15).

As a result, Lehi continued, “the Lord God gave unto man that he should act for himself. Wherefore, man could not act for himself save it should be that he was enticed by the one or the other” (verse 16). Similarly, in modern revelation the Lord declares, “It must needs be that the devil should tempt the children of men, or they could not be agents unto themselves” (D\&C 29:39).

Opposition was necessary in the Garden of Eden. If Adam and Eve had not made the choice that introduced mortality, Lehi taught, “they would have remained in a state of innocence, … doing no good, for they knew no sin” (2 Nephi 2:23).

From the beginning, agency and opposition were central to the Father’s plan and to Satan’s rebellion against it. As the Lord revealed to Moses, in the council of heaven Satan “sought to destroy the agency of man” (Moses 4:3). That destruction was inherent in the terms of Satan’s offer. He came before the Father and said, “Behold, here am I, send me, I will be thy son, and I will redeem all mankind, that one soul shall not be lost, and surely I will do it; wherefore give me thine honor” (Moses 4:1).

Thus, Satan proposed to carry out the Father’s plan in a way that would prevent the accomplishment of the Father’s purpose and give Satan His glory.

Satan’s proposal would have ensured perfect equality: it would “redeem all mankind,” that not one soul would be lost. There would be no agency or choice by anyone and, therefore, no need for opposition. There would be no test, no failure, and no success. There would be no growth to attain the purpose the Father desired for His children. The scriptures record that Satan’s opposition resulted in a “war in heaven” (Revelation 12:7), in which two-thirds of the children of God earned the right to experience mortal life by choosing the Father’s plan and rejecting Satan’s rebellion.

Satan’s purpose was to gain for himself the Father’s honor and power (see Isaiah 14:12–15; Moses 4:1, 3). “Wherefore,” the Father said, “because that Satan rebelled against me, … I caused that he should be cast down” (Moses 4:3) with all the spirits who had exercised their agency to follow him (see Jude 1:6; Revelation 12:8–9; D\&C 29:36–37). Cast down as unembodied spirits in mortality, Satan and his followers tempt and seek to deceive and captivate the children of God (see Moses 4:4). So it is that the evil one, who opposed and sought to destroy the Father’s plan, actually facilitated it, because it is opposition that enables choice and it is the opportunity of making the right choices that leads to the growth that is the purpose of the Father’s plan.

\intalkheader{II.}

Significantly, the temptation to sin is not the only kind of opposition in mortality. Father Lehi taught that if the Fall had not taken place, Adam and Eve “would have remained in a state of innocence, having no joy, for they knew no misery” (2 Nephi 2:23). Without the experience of opposition in mortality, “all things must needs be a compound in one,” in which there would be no happiness or misery (verse 11). Therefore, Father Lehi continued, after God had created all things, “to bring about his eternal purposes in the end of man, … it must needs be that there was an opposition; even the forbidden fruit in opposition to the tree of life; the one being sweet and the other bitter” (verse 15). His teaching on this part of the plan of salvation concludes with these words:

“Behold, all things have been done in the wisdom of him who knoweth all things.

“Adam fell that men might be; and men are, that they might have joy” (verses 24–25).

Opposition in the form of difficult circumstances we face in mortality is also part of the plan that furthers our growth in mortality.

\intalkheader{III.}

All of us experience various kinds of opposition that test us. Some of these tests are temptations to sin. Some are mortal challenges apart from personal sin. Some are very great. Some are minor. Some are continuous, and some are mere episodes. None of us is exempt. Opposition permits us to grow toward what our Heavenly Father would have us become.

After Joseph Smith had completed translating the Book of Mormon, he still had to find a publisher. This was not easy. The complexity of this lengthy manuscript and the cost of printing and binding thousands of copies were intimidating. Joseph first approached E. B. Grandin, a Palmyra printer, who refused. He then sought another printer in Palmyra, who also turned him down. He traveled to Rochester, 25 miles (40 km) away, and approached the most prominent publisher in western New York, who also turned him down. Another Rochester publisher was willing, but circumstances made this alternative unacceptable.

Weeks had passed, and Joseph must have been bewildered at the opposition to accomplishing his divine mandate. The Lord did not make it easy, but He did make it possible. Joseph’s fifth attempt, a second approach to the Palmyra publisher Grandin, was successful.

Years later, Joseph was painfully imprisoned in Liberty Jail for many months. When he prayed for relief, the Lord told him that “all these things shall give thee experience, and shall be for thy good” (D\&C 122:7).

We are all acquainted with other kinds of mortal opposition not caused by our personal sins, including illness, disability, and death. President Thomas S. Monson explained:

“Some of you may at times have cried out in your suffering, wondering why our Heavenly Father would allow you to go through whatever trials you are facing. …

“Our mortal life, however, was never meant to be easy or consistently pleasant. Our Heavenly Father … knows that we learn and grow and become refined through hard challenges, heartbreaking sorrows, and difficult choices. Each one of us experiences dark days when our loved ones pass away, painful times when our health is lost, feelings of being forsaken when those we love seem to have abandoned us. These and other trials present us with the real test of our ability to endure.”

Our efforts to improve our observance of the Sabbath day pose a less stressful example of opposition. We have the Lord’s commandment to honor the Sabbath. Some of our choices may violate that commandment, but other choices in how to spend time on the Sabbath are simply a question of whether we will do what is merely good or what is better or best.

To illustrate the opposition of temptation, the Book of Mormon describes three methods the devil will use in the last days. First, he will “rage in the hearts of the children of men, and stir them up to anger against that which is good” (2 Nephi 28:20). Second, he will “pacify, and lull [members] away into carnal security,” saying “Zion prospereth, all is well” (verse 21). Third, he will tell us “there is no hell; and … I am no devil, for there is none” (verse 22), and therefore there is no right and wrong. Because of this opposition, we are warned not to be “at ease in Zion!” (verse 24).

The Church in its divine mission and we in our personal lives seem to face increasing opposition today. Perhaps as the Church grows in strength and we members grow in faith and obedience, Satan increases the strength of his opposition so we will continue to have “opposition in all things.”

Some of this opposition even comes from Church members. Some who use personal reasoning or wisdom to resist prophetic direction give themselves a label borrowed from elected bodies—“the loyal opposition.” However appropriate for a democracy, there is no warrant for this concept in the government of God’s kingdom, where questions are honored but opposition is not (see Matthew 26:24).

As another example, there are many things in our early Church history, such as what Joseph Smith did or did not do in every circumstance, that some use as a basis for opposition. To all I say, exercise faith and put reliance on the Savior’s teaching that we should “know them by their fruits” (Matthew 7:16). The Church is making great efforts to be transparent with the records we have, but after all we can publish, our members are sometimes left with basic questions that cannot be resolved by study. That is the Church history version of “opposition in all things.” Some things can be learned only by faith (see D\&C 88:118). Our ultimate reliance must be on faith in the witness we have received from the Holy Ghost.

God rarely infringes on the agency of any of His children by intervening against some for the relief of others. But He does ease the burdens of our afflictions and strengthen us to bear them, as He did for Alma’s people in the land of Helam (see Mosiah 24:13–15). He does not prevent all disasters, but He does answer our prayers to turn them aside, as He did with the uniquely powerful cyclone that threatened to prevent the dedication of the temple in Fiji; or He does blunt their effects, as He did with the terrorist bombing that took so many lives in the Brussels airport but only injured our four missionaries.

Through all mortal opposition, we have God’s assurance that He will “consecrate [our] afflictions for [our] gain” (2 Nephi 2:2). We have also been taught to understand our mortal experiences and His commandments in the context of His great plan of salvation, which tells us the purpose of life and gives us the assurance of a Savior, in whose name I testify of the truth of these things. In the name of Jesus Christ, amen.

\newpage
\settalk{Sharing the Restored Gospel}
\setspeaker{Dallin H. Oaks}
\settalkdate{October 2016}
\talkheading{Sharing the Restored Gospel}{Dallin H. Oaks}

\intalkheader{I.}

Nearing the end of His earthly ministry, our Savior, Jesus Christ, commanded His disciples: “Go ye therefore, and teach all nations” (Matthew 28:19) and “Go ye into all the world, and preach the gospel to every creature” (Mark 16:15). All Christians are under these commands to share the gospel with everyone. Many call this the “great commission.”

As Elder Neil L. Andersen described in the morning session, Latter-day Saints are surely among those most committed to this great responsibility. We should be because we know that God loves all of His children and that in these last days He has restored vital additional knowledge and power to bless all of them. The Savior taught us to love all as our brothers and sisters, and we honor that teaching by sharing the witness and message of the restored gospel “among all nations, kindreds, tongues, and people” (D\&C 112:1). This is a vital part of what it means to be a Latter-day Saint. We look on this as a joyful privilege. What could be more joyful than sharing the truths of eternity with God’s children?

Today we have many resources to share the gospel that were not available in earlier generations. We have TV, the internet, and social media channels. We have many valuable messages to introduce the restored gospel. We have the prominence of the Church in many nations. We have a greatly increased number of missionaries. But are we using all these resources to maximum effect? I believe most of us would say no. We desire to be more effective in fulfilling our divinely appointed responsibility to proclaim the restored gospel in all the world.

There are many good ideas for sharing the gospel that will work in individual stakes or countries. However, because we are a worldwide Church, I wish to speak of ideas that will work everywhere, from the newest units to the most established, from cultures now receptive to the gospel of Jesus Christ to cultures and nations that are increasingly hostile to religion. I want to speak of ideas that you can share with persons who are committed believers in Jesus Christ as well as with persons who have never heard His name, with persons who are satisfied with their current lives as well as with persons who are desperately seeking to improve themselves.

What can I say that will be helpful in your sharing the gospel, whatever your circumstances? We need the help of every member, and every member can help, since there are many tasks to perform as we share the restored gospel with every nation, kindred, tongue, and people.

We all know that member participation in missionary work is vital to achieving both conversion and retention. President Thomas S. Monson has said: “Now is the time for members and missionaries to come together … [and] labor in the Lord’s vineyard to bring souls unto Him. He has prepared the means for us to share the gospel in a multitude of ways, and He will assist us in our labors if we will act in faith to fulfill His work.”

Sharing the restored gospel is our lifelong Christian duty and privilege. Elder Quentin L. Cook reminds us, “Missionary work is not just one of the 88 keys on a piano that is occasionally played; it is a major chord in a compelling melody that needs to be played continuously throughout our lives if we are to remain in harmony with our commitment to Christianity and the gospel of Jesus Christ.”

\intalkheader{II.}

There are three things all members can do to help share the gospel, regardless of the circumstances in which they live and work. All of us should do all of these.

First, we can all pray for desire to help with this vital part of the work of salvation. All efforts begin with desire.

Second, we can keep the commandments ourselves. Faithful, obedient members are the most persuasive witnesses of the truth and value of the restored gospel. Even more important, faithful members will always have the Savior’s Spirit to be with them to guide them as they seek to participate in the great work of sharing the restored gospel of Jesus Christ.

Third, we can pray for inspiration on what we can do in our individual circumstances to share the gospel with others. This is different than praying for the missionaries or praying for what others can do. We should pray for what we can do personally. When we pray, we should remember that prayers for this kind of inspiration will be answered if accompanied by a commitment—something the scriptures call “real intent” or “full purpose of heart.” Pray with a commitment to act upon the inspiration you receive, promising the Lord that if He will inspire you to speak to someone about the gospel, you will do it.

We need the guidance of the Lord because at any particular time some are—and some are not—ready for the additional truths of the restored gospel. We should never set ourselves up as judges of who is ready and who is not. The Lord knows the hearts of all of His children, and if we pray for inspiration, He will help us find persons He knows to be “in a preparation to hear the word” (Alma 32:6).

As an Apostle of the Lord, I urge every member and family in the Church to pray for the Lord to help them find persons prepared to receive the message of the restored gospel of Jesus Christ. Elder M. Russell Ballard has given this important counsel, with which I concur: “Trust the Lord. He is the Good Shepherd. He knows His sheep. … If we are not engaged, many who would hear the message of the Restoration will be passed by. … The principles are pretty simple—pray, personally and in your family, for missionary opportunities.” As we demonstrate our faith, these opportunities will come without any “forced or … contrived response. They will flow as a natural result of our love for our brothers and sisters.”

I know this is true. I add my promise that with faith in the Lord’s help, we will be guided, be inspired, and find great joy in this eternally important work of love. We will come to understand that success in sharing the gospel is inviting people with love and genuine intent to help them, no matter what their response.

\intalkheader{III.}

Here are some other things we can do to share the gospel effectively:

Sharing the gospel is not a burden but a joy. What we call “member missionary work” is not a program but an attitude of love and outreach to help those around us. It is also an opportunity to witness how we feel about the restored gospel of our Savior. As Elder Ballard has taught, “A most significant evidence of our conversion and of how we feel about the gospel in our own lives is our willingness to share it with others.”

I testify of Jesus Christ, who is the Light and Life of the World (see 3 Nephi 11:11). His restored gospel lights our way in mortality. His Atonement gives us the assurance of life after death and the strength to persist toward immortality. And His Atonement gives us the opportunity to be forgiven of our sins and, under God’s glorious plan of salvation, to qualify for eternal life, “the greatest of all the gifts of God” (D\&C 14:7). In the name of Jesus Christ, amen.

\newpage
\settalk{The Godhead and the Plan of Salvation}
\setspeaker{Dallin H. Oaks}
\settalkdate{April 2017}
\talkheading{The Godhead and the Plan of Salvation}{Dallin H. Oaks}

\intalkheader{I.}

Our first article of faith declares, “We believe in God, the Eternal Father, and in His Son, Jesus Christ, and in the Holy Ghost.” We join other Christians in this belief in a Father and a Son and a Holy Ghost, but what we believe about Them is different from the beliefs of others. We do not believe in what the Christian world calls the doctrine of the Holy Trinity. In his First Vision, Joseph Smith saw two distinct personages, two beings, thus clarifying that the then-prevailing beliefs concerning God and the Godhead were not true.

In contrast to the belief that God is an incomprehensible and unknowable mystery is the truth that the nature of God and our relationship to Him is knowable and is the key to everything else in our doctrine. The Bible records Jesus’s great Intercessory Prayer, where He declared that “this is life eternal, that they might know thee the only true God, and Jesus Christ, whom thou hast sent” (John 17:3).

The effort to know God and His work began before mortality and will not be concluded here. The Prophet Joseph Smith taught, “It will be a great while after you have passed through the veil before you will have learned … all the principles of exaltation.” We build on the knowledge we acquired in the premortal spirit world. Thus, in trying to teach Israelites the nature of God and His relationship to His children, the prophet Isaiah declared, as recorded in the Bible:

“To whom then will ye liken God? or what likeness will ye compare unto him? …

“Have ye not known? have ye not heard? hath it not been told you from the beginning? have ye not understood from the foundations of the earth?” (Isaiah 40:18, 21).

We know that the three members of the Godhead are separate and distinct beings. We know this from instruction given by the Prophet Joseph Smith: “The Father has a body of flesh and bones as tangible as man’s; the Son also; but the Holy Ghost has not a body of flesh and bones, but is a personage of Spirit. Were it not so, the Holy Ghost could not dwell in us” (D\&C 130:22).

As to the supreme position of God the Father within the Godhead, as well as the respective roles each personage performs, the Prophet Joseph explained:

“Any person that had seen the heavens opened knows that there are three personages in the heavens who hold the keys of power, and one presides over all. …

“… These personages … are called God the first, the Creator; God the second, the Redeemer; and God the third, the Witness or Testator.

“[It is] the province of the Father to preside as the Chief or President, Jesus as the Mediator, and the Holy Ghost as the Testator or Witness.”

\intalkheader{II. The Plan}

We understand our relationship to the members of the Godhead from what is revealed about the plan of salvation.

Questions like “Where did we come from?” “Why are we here?” and “Where are we going?” are answered in what the scriptures call the “plan of salvation,” the “great plan of happiness,” or the “plan of redemption” (Alma 42:5, 8, 11). The gospel of Jesus Christ is central to this plan.

As spirit children of God, in an existence prior to mortality, we desired a destiny of eternal life but had progressed as far as we could without a mortal experience in a physical body. To provide that opportunity, our Heavenly Father presided over the Creation of this world, where, deprived of our memory of what preceded our mortal birth, we could prove our willingness to keep His commandments and experience and grow through the other challenges of mortal life. But in the course of that mortal experience, and as a result of the Fall of our first parents, we would suffer spiritual death by being cut off from the presence of God, be soiled by sin, and become subject to physical death. The Father’s plan anticipated and provided ways to overcome all of those barriers.

\intalkheader{III. The Godhead}

Knowing the purpose of God’s great plan, we now consider the respective roles of the three members of the Godhead in that plan.

We begin with a teaching from the Bible. In concluding his second letter to the Corinthians, the Apostle Paul makes this almost offhand reference to the Godhead of Father, Son, and Holy Ghost: “The grace of the Lord Jesus Christ, and the love of God, and the communion [or fellowship] of the Holy Ghost, be with you all” (2 Corinthians 13:14).

This biblical scripture represents the Godhead and references the all-defining and motivating love of God the Father, the merciful and saving mission of Jesus Christ, and the fellowship of the Holy Ghost.

It all begins with God the Father. While we know comparatively little about Him, what we know is decisive in understanding His supreme position, our relationship to Him, and His superintending role in the plan of salvation, the Creation, and all else that followed.

As Elder Bruce R. McConkie wrote just before his death: “In the ultimate and final sense of the word, there is only one true and living God. He is the Father, the Almighty Elohim, the Supreme Being, the Creator and Ruler of the universe.” He is the God and Father of Jesus Christ, as well as of all of us. President David O. McKay taught that “the first fundamental truth advocated by Jesus Christ was this, that behind, above and over all there is God the Father, Lord of heaven and earth.”

What we know of the nature of God the Father is mostly what we can learn from the ministry and teachings of His Only Begotten Son, Jesus Christ. As Elder Jeffrey R. Holland has taught, one of the paramount purposes of Jesus’s ministry was to reveal to mortals “what God our Eternal Father is like, … to reveal and make personal to us the true nature of His Father, our Father in Heaven.” The Bible contains an apostolic witness that Jesus was “the express image” of His Father’s person (Hebrews 1:3), which merely elaborates Jesus’s own teaching that “he that hath seen me hath seen the Father” (John 14:9).

God the Father is the Father of our spirits. We are His children. He loves us, and all that He does is for our eternal benefit. He is the author of the plan of salvation, and it is by His power that His plan achieves its purposes for the ultimate glory of His children.

To mortals, the most visible member of the Godhead is Jesus Christ. A great doctrinal statement by the First Presidency in 1909 declares Him to be “the firstborn among all the sons of God—the first begotten in the spirit, and the only begotten in the flesh.” The Son, the greatest of all, was chosen by the Father to carry out the Father’s plan—to exercise the Father’s power to create worlds without number (see Moses 1:33) and to save the children of God from death by His Resurrection and from sin by His Atonement. This supernal sacrifice is truly called “the central act of all human history.”

On those unique and sacred occasions when God the Father personally introduced the Son, He has said, “This is my beloved Son: hear him” (Mark 9:7; Luke 9:35; see also 3 Nephi 11:7; Joseph Smith—History 1:17). Thus, it is Jesus Christ, Jehovah, the Lord God of Israel, who speaks to and through the prophets. So it is that when Jesus appeared to the Nephites after His Resurrection, He introduced Himself as “the God of the whole earth” (3 Nephi 11:14). So it is that Jesus often speaks to the prophets of the Book of Mormon and to the Latter-day Saints as “the Father and the Son,” a title explained in the First Presidency and Quorum of the Twelve’s inspired doctrinal exposition just 100 years ago.

The third member of the Godhead is the Holy Ghost, also referred to as the Holy Spirit, the Spirit of the Lord, and the Comforter. He is the member of the Godhead who is the agent of personal revelation. As a personage of spirit (see D\&C 130:22), He can dwell in us and perform the essential role of communicator between the Father and the Son and the children of God on earth. Many scriptures teach that His mission is to testify of the Father and the Son (see John 15:26; 3 Nephi 28:11; D\&C 42:17). The Savior promised that the Comforter will teach us all things, bring all things to our remembrance, and guide us into all truth (see John 14:26; 16:13). Thus, the Holy Ghost helps us discern between truth and falsehood, guides us in our major decisions, and helps us through the challenges of mortality. He is also the means by which we are sanctified, that is, cleansed and purified from sin (see 2 Nephi 31:17; 3 Nephi 27:20; Moroni 6:4).

\intalkheader{IV.}

So, how does understanding this heavenly revealed doctrine about the Godhead and the plan of salvation help us with our challenges today?

Because we have the truth about the Godhead and our relationship to Them, the purpose of life, and the nature of our eternal destiny, we have the ultimate road map and assurance for our journey through mortality. We know whom we worship and why we worship. We know who we are and what we can become (see D\&C 93:19). We know who makes it all possible, and we know what we must do to enjoy the ultimate blessings that come through God’s plan of salvation. How do we know all of this? We know by the revelations of God to His prophets and to each of us individually.

Attaining what the Apostle Paul described as “the measure of the stature of the fulness of Christ” (Ephesians 4:13) requires far more than acquiring knowledge. It is not even enough for us to be convinced of the gospel; we must act and think so that we are converted by it. In contrast to the institutions of the world, which teach us to know something, the plan of salvation and the gospel of Jesus Christ challenge us to become something.

As President Thomas S. Monson taught us in our last general conference:

“Essential to the plan [of salvation] is our Savior, Jesus Christ. Without His atoning sacrifice, all would be lost. It is not enough, however, merely to believe in Him and His mission. We need to work and learn, search and pray, repent and improve. We need to know God’s laws and live them. We need to receive His saving ordinances. Only by so doing will we obtain true, eternal happiness. …

“From the depths of my soul and in all humility,” President Monson declared, “I testify of the great gift which is our Father’s plan for us. It is the one perfect path to peace and happiness both here and in the world to come.”

I add my testimony to that of our beloved prophet-president. I testify that we have a Heavenly Father, who loves us. I testify that we have a Holy Ghost, who guides us. And I testify of Jesus Christ, our Savior, who makes it all possible, in the name of Jesus Christ, amen.

\newpage
\settalk{The Plan and the Proclamation}
\setspeaker{Dallin H. Oaks}
\settalkdate{October 2017}
\talkheading{The Plan and the Proclamation}{Dallin H. Oaks}

As is evident in our family proclamation, members of The Church of Jesus Christ of Latter-day Saints are blessed with unique doctrine and different ways of viewing the world. We participate and even excel in many worldly activities, but on some subjects we forgo participation as we seek to follow the teachings of Jesus Christ and His Apostles, ancient and modern.

\intalkheader{I.}

In a parable, Jesus described those who “[hear] the word” but become “unfruitful” when that word is “choke[d]” by “the care of this world, and the deceitfulness of riches” (Matthew 13:22). Later, Jesus corrected Peter for not savoring “the things that be of God, but those that be of men,” declaring, “For what is a man profited, if he shall gain the whole world, and lose his own soul?” (Matthew 16:23, 26). In His final teachings in mortality, He told His Apostles, “If ye were of the world, the world would love his own: but because ye are not of the world, … the world hateth you” (John 15:19; see also John 17:14, 16).

Similarly, the writings of Jesus’s early Apostles frequently use the image of “the world” to represent opposition to gospel teachings. “Be not conformed to this world” (Romans 12:2), the Apostle Paul taught. “For the wisdom of this world is foolishness with God” (1 Corinthians 3:19). And, “Beware,” he warned, “lest any man spoil you … after the tradition of men, after the rudiments of the world, and not after Christ” (Colossians 2:8). The Apostle James taught that “the friendship of the world is enmity with God[.] Whosoever therefore will be a friend of the world is the enemy of God” (James 4:4).

The Book of Mormon often uses this image of the opposition of “the world.” Nephi prophesied the ultimate destruction of “those who are built up to become popular in the eyes of the world, and those who seek … the things of the world” (1 Nephi 22:23; see also 2 Nephi 9:30). Alma condemned those who were “puffed up … with the vain things of the world” (Alma 31:27). Lehi’s dream shows that those who seek to follow the iron rod, the word of God, will encounter the opposition of the world. The occupants of the “great and spacious building” Lehi saw were “mocking and pointing” the “finger of scorn” (1 Nephi 8:26–27, 33). In his vision interpreting this dream, Nephi learned that this ridicule and opposition came from “the multitudes of the earth, … the world and the wisdom thereof; … the pride of the world” (1 Nephi 11:34–36).

What is the meaning of these scriptural cautions and commandments not to be “of the world” or the modern commandment to “forsake the world”? (D\&C 53:2). President Thomas S. Monson summarized these teachings: “We must be vigilant in a world which has moved so far from that which is spiritual. It is essential that we reject anything that does not conform to our standards, refusing in the process to surrender that which we desire most: eternal life in the kingdom of God.”

God created this earth according to His plan to provide His spirit children a place to experience mortality as a necessary step toward the glories He desires for all His children. While there are various kingdoms and glories, our Heavenly Father’s ultimate desire for His children is what President Monson called “eternal life in the kingdom of God,” which is exaltation in families. This is more than salvation. President Russell M. Nelson has reminded us, “In God’s eternal plan, salvation is an individual matter; [but] exaltation is a family matter.”

The restored gospel of Jesus Christ and the inspired family proclamation, which I will discuss later, are essential teachings to guide mortal preparation for exaltation. Even as we must live with the marriage laws and other traditions of a declining world, those who strive for exaltation must make personal choices in family life according to the Lord’s way whenever that differs from the world’s way.

In this mortal life, we have no memory of what preceded our birth, and we now experience opposition. We grow and mature spiritually by choosing to obey God’s commandments in a succession of right choices. These include covenants and ordinances and repentance when our choices are wrong. In contrast, if we lack faith in God’s plan and are disobedient to or deliberately refrain from its required actions, we forgo that growth and maturity. The Book of Mormon teaches, “This life is the time for men to prepare to meet God” (Alma 34:32).

\intalkheader{II.}

Latter-day Saints who understand God’s plan of salvation have a unique worldview that helps them see the reason for God’s commandments, the unchangeable nature of His required ordinances, and the fundamental role of our Savior, Jesus Christ. Our Savior’s Atonement reclaims us from death and, subject to our repentance, saves us from sin. With that worldview, Latter-day Saints have distinctive priorities and practices and are blessed with the strength to endure the frustrations and pains of mortal life.

Inevitably, the actions of those who try to follow God’s plan of salvation can cause misunderstanding or even conflict with family members or friends who do not believe its principles. Such conflict is always so. Every generation that has sought to follow God’s plan has had challenges. Anciently, the prophet Isaiah gave strength to the Israelites, whom he called “ye that know righteousness, … in whose heart is my law.” To them he said, “Fear ye not the reproach of men, neither be ye afraid of their revilings” (Isaiah 51:7; see also 2 Nephi 8:7). But whatever the cause of conflict with those who do not understand or believe God’s plan, those who do understand are always commanded to choose the Lord’s way instead of the world’s way.

\intalkheader{III.}

The gospel plan each family should follow to prepare for eternal life and exaltation is outlined in the Church’s 1995 proclamation, “The Family: A Proclamation to the World.” Its declarations are, of course, visibly different from some current laws, practices, and advocacy of the world in which we live. In our day, the differences most evident are cohabitation without marriage, same-sex marriage, and the raising of children in such relationships. Those who do not believe in or aspire to exaltation and are most persuaded by the ways of the world consider this family proclamation as just a statement of policy that should be changed. In contrast, Latter-day Saints affirm that the family proclamation defines the kind of family relationships where the most important part of our eternal development can occur.

We have witnessed a rapid and increasing public acceptance of cohabitation without marriage and of same-sex marriage. The corresponding media advocacy, education, and even occupational requirements pose difficult challenges for Latter-day Saints. We must try to balance the competing demands of following the gospel law in our personal lives and teachings, even as we seek to show love for all. In doing so we sometimes face, but need not fear, what Isaiah called “the reproach of men.”

Converted Latter-day Saints believe that the family proclamation, issued nearly a quarter century ago and now translated into scores of languages, is the Lord’s reemphasis of the gospel truths we need to sustain us through current challenges to the family. Two examples are same-sex marriage and cohabitation without marriage. Just 20 years after the family proclamation, the United States Supreme Court authorized same-sex marriage, overturning thousands of years of marriage being limited to a man and a woman. The shocking percentage of United States children born to a mother not married to the father came more gradually: 5 percent in 1960, 32 percent in 1995, and now 40 percent.

\intalkheader{IV.}

The family proclamation begins by declaring “that marriage between a man and a woman is ordained of God and that the family is central to the Creator’s plan for the eternal destiny of His children.” It also affirms that “gender is an essential characteristic of individual premortal, mortal, and eternal identity and purpose.” It further declares “that God has commanded that the sacred powers of procreation are to be employed only between man and woman, lawfully wedded as husband and wife.”

The proclamation affirms the continuing duty of husband and wife to multiply and replenish the earth and their “solemn responsibility to love and care for each other and for their children”: “Children are entitled to birth within the bonds of matrimony, and to be reared by a father and a mother who honor marital vows with complete fidelity.” It solemnly warns against the abuse of spouse or offspring, and it affirms that “happiness in family life is most likely to be achieved when founded upon the teachings of the Lord Jesus Christ.” Finally, it calls for the promotion of official “measures designed to maintain and strengthen the family as the fundamental unit of society.”

In 1995 a President of the Church and 14 other Apostles of the Lord issued these important doctrinal statements. As one of only seven of those Apostles still living, I feel obliged to share what led to the family proclamation for the information of all who consider it.

The inspiration identifying the need for a proclamation on the family came to the leadership of the Church over 23 years ago. It was a surprise to some who thought the doctrinal truths about marriage and the family were well understood without restatement. Nevertheless, we felt the confirmation and we went to work. Subjects were identified and discussed by members of the Quorum of the Twelve for nearly a year. Language was proposed, reviewed, and revised. Prayerfully we continually pleaded with the Lord for His inspiration on what we should say and how we should say it. We all learned “line upon line, precept upon precept,” as the Lord has promised (D\&C 98:12).

During this revelatory process, a proposed text was presented to the First Presidency, who oversee and promulgate Church teachings and doctrine. After the Presidency made further changes, the proclamation on the family was announced by the President of the Church, Gordon B. Hinckley. In the women’s meeting of September 23, 1995, he introduced the proclamation with these words: “With so much of sophistry that is passed off as truth, with so much of deception concerning standards and values, with so much of allurement and enticement to take on the slow stain of the world, we have felt to warn and forewarn.”

I testify that the proclamation on the family is a statement of eternal truth, the will of the Lord for His children who seek eternal life. It has been the basis of Church teaching and practice for the last 22 years and will continue so for the future. Consider it as such, teach it, live by it, and you will be blessed as you press forward toward eternal life.

Forty years ago, President Ezra Taft Benson taught that “every generation has its tests and its chance to stand and prove itself.” I believe our attitude toward and use of the family proclamation is one of those tests for this generation. I pray for all Latter-day Saints to stand firm in that test.

I close with President Gordon B. Hinckley’s teachings uttered two years after the family proclamation was announced. He said: “I see a wonderful future in a very uncertain world. If we will cling to our values, if we will build on our inheritance, if we will walk in obedience before the Lord, if we will simply live the gospel, we will be blessed in a magnificent and wonderful way. We will be looked upon as a peculiar people who have found the key to a peculiar happiness.”

I testify of the truth and eternal importance of the family proclamation, revealed by the Lord Jesus Christ to His Apostles for the exaltation of the children of God (see Doctrine and Covenants 131:1–4), in the name of Jesus Christ, amen.

\newpage
\settalk{Small and Simple Things}
\setspeaker{Dallin H. Oaks}
\settalkdate{April 2018}
\talkheading{Small and Simple Things}{Dallin H. Oaks}

\intalkheader{I.}

My dear brothers and sisters, like you, I have been profoundly touched and edified and inspired by the messages and music and the feelings of this time together. I’m sure I speak for you in expressing thanks to our brothers and sisters who, as instruments in the hands of the Lord, have given us the strengthening effect of this time together.

I am grateful to speak to this audience on Easter Sunday. Today we join other Christians in celebrating the Resurrection of the Lord Jesus Christ. For members of The Church of Jesus Christ of Latter-day Saints, the literal Resurrection of Jesus Christ is a pillar of our faith.

Because we believe the accounts in both the Bible and the Book of Mormon about the literal Resurrection of Jesus Christ, we also believe the numerous scriptural teachings that a similar resurrection will come to all mortals who have ever lived upon this earth. That resurrection gives us what the Apostle Peter called “a lively hope” (1 Peter 1:3). That lively hope is our conviction that death is not the conclusion of our identity but merely a necessary step in our Heavenly Father’s merciful plan for the salvation of His children. That plan calls for a transition from mortality to immortality. Central to that transition is the sunset of death and the glorious morning made possible by the Resurrection of our Lord and Savior that we celebrate on this Easter Sunday.

\intalkheader{II.}

In a great hymn whose words were written by Eliza R. Snow, we sing:

\begin{verse}
How great, how glorious, how complete\\
Redemption’s grand design,\\
Where justice, love, and mercy meet\\
In harmony divine!
\end{verse}

In furtherance of that divine design and harmony, we assemble in meetings, including this conference, to teach and encourage one another.

This morning I have felt to use as my text Alma’s teaching to his son Helaman, recorded in the Book of Mormon: “By small and simple things are great things brought to pass” (Alma 37:6).

We are taught many small and simple things in the gospel of Jesus Christ. We need to be reminded that in total and over a significant period of time, these seemingly small things bring to pass great things. There have been many talks on this subject by General Authorities and by other respected teachers. The subject is so important that I feel to speak of it again.

I was reminded of the power of small and simple things over time by something I saw on a morning walk. Here is the picture I took. The thick and strong concrete sidewalk is cracking. Is this the result of some large and powerful thrust? No, this cracking is caused by the slow, small growth of one of the roots reaching out from the adjoining tree. Here is a similar example I saw on another street.

The thrusting power that cracked these heavy concrete sidewalks was too small to measure on a daily or even a monthly basis, but its effect over time was incredibly powerful.

So is the powerful effect over time of the small and simple things we are taught in the scriptures and by living prophets. Consider the scripture study we’ve been taught to incorporate into our daily lives. Or consider the personal prayers and the kneeling family prayers that are regular practices for faithful Latter-day Saints. Consider attendance at seminary for youth or institute classes for young adults. Though each of these practices may seem to be small and simple, over time they result in powerful spiritual uplift and growth. This occurs because each of these small and simple things invites the companionship of the Holy Ghost, the Testifier who enlightens us and guides us into truth, as President Eyring has explained.

Another source of spiritual uplift and growth is an ongoing practice of repenting, even of seemingly small transgressions. Our own inspired self-evaluations can help us see how we have fallen short and how we can do better. Such repentance should precede our weekly partaking of the sacrament. Some subjects to consider in this process of repentance are suggested in the hymn “Have I Done Any Good?”

\begin{verse}
Have I done any good in the world today?\\
Have I helped anyone in need?\\
Have I cheered up the sad and made someone feel glad?\\
If not, I have failed indeed.\\
Has anyone’s burden been lighter today\\
Because I was willing to share?\\
Have the sick and the weary been helped on their way?\\
When they needed my help was I there?
\end{verse}

Surely these are small things, but surely they are good examples of what Alma taught his son Helaman: “And the Lord God doth work by means to bring about his great and eternal purposes; and by very small means the Lord … bringeth about the salvation of many souls” (Alma 37:7).

President Steven C. Wheelwright gave an audience at Brigham Young University–Hawaii this inspired description of Alma’s teaching: “Alma confirms for his son that indeed the pattern the Lord follows when we exercise faith in Him and follow His counsel in small and simple things is that He blesses us with small daily miracles, and over time, with marvelous works.”

President Howard W. Hunter taught that “frequently it is the commonplace tasks … that have the greatest positive effect on the lives of others, as compared with the things that the world so often relates to greatness.”

A persuasive secular teaching of this same principle comes from former Senator Dan Coats of Indiana, who wrote: “The only preparation for that one profound decision which can change a life, or even a nation, is those hundreds and thousands of half-conscious, self-defining, seemingly insignificant decisions made in private.”

Those “seemingly insignificant” private decisions include how we use our time, what we view on television and the internet, what we read, the art and music with which we surround ourselves at work and at home, what we seek for entertainment, and how we apply our commitment to be honest and truthful. Another seemingly small and simple thing is being civil and cheerful in our personal interactions.

None of these desirable small and simple things will lift us to great things unless they are practiced consistently and continuously. President Brigham Young was reported as saying: “Our lives are made up of little, simple circumstances that amount to a great deal when they are brought together, and sum up the whole life of the man or woman.”

We are surrounded by media influences and cultural deteriorations that will carry us downstream in our values if we are not continually resisting. To move upstream toward our eternal goal, we must constantly keep paddling. It helps if we are part of a team that is paddling together, like a rowing crew in action. To extend that example even further, the cultural currents are so strong that if we ever stop paddling, we will be carried downstream toward a destination we do not seek but which becomes inevitable if we do not constantly try to move forward.

After reciting a seemingly small event that had great consequences, Nephi wrote, “And thus we see that by small means the Lord can bring about great things” (1 Nephi 16:29). The Old Testament includes a memorable example of this. There we read how the Israelites were plagued by fiery serpents. Many people died from their bites (see Numbers 21:6). When Moses prayed for relief, he was inspired to make “a serpent of brass, and put it upon a pole.” Then, “if a serpent had bitten any man, when he beheld the serpent of brass, he lived” (verse 9). Such a small thing for such a miraculous result! Yet, as Nephi explained when he taught this example to those who were rebelling against the Lord, even when the Lord had prepared a simple way by which they could be healed, “because of the simpleness of the way, or the easiness of it, there were many who perished” (1 Nephi 17:41).

That example and that teaching remind us that the simplicity of the way or the easiness of the commanded task cannot mean that it is unimportant to achieve our righteous desire.

Similarly, even small acts of disobedience or minor failures to follow righteous practices can draw us down toward an outcome we have been warned to avoid. The Word of Wisdom provides an example of this. Likely the effect on the body of one cigarette or one drink of alcohol or one dose of another drug cannot be measured. But over time, the effect is powerful and may be irreversible. Remember the cracking of the sidewalk by the gradual small expansions of the root of the tree. One thing is certain, the terrible consequences of partaking of anything that can become addictive, like drugs that attack our bodies or pornographic material that degrades our thoughts, is totally avoidable if we never partake for the first time—even once.

Many years ago, President M. Russell Ballard described to a general conference audience “how small and simple things can be negative and destructive to a person’s salvation.” He taught: “Like weak fibers that form a yarn, then a strand, and finally a rope, these small things combined together can become too strong to be broken. We must ever be aware of the power that the small and simple things can have in building spirituality,” he said. “At the same time, we must be aware that Satan will use small and simple things to lead us into despair and misery.”

President Wheelwright gave a similar caution to his BYU–Hawaii audience: “It is in failing to do the small and simple things that faith wavers, miracles cease, and progress towards the Lord and His kingdom is first put on hold and then begins to unravel as seeking after the kingdom of God is replaced with more temporal pursuits and worldly ambitions.”

To protect against the cumulative negative effects that are destructive to our spiritual progress, we need to follow the spiritual pattern of small and simple things. Elder David A. Bednar described this principle in a BYU Women’s Conference: “We can learn much about the nature and importance of this spiritual pattern from the technique of … dripping water onto the soil at very low rates,” in contrast to flooding or spraying large quantities of water where it may not be needed.

He explained: “The steady drips of water sink deep into the ground and provide a high moisture level in the soil wherein plants can flourish. In like manner, if you and I are focused and frequent in receiving consistent drops of spiritual nourishment, then gospel roots can sink deep into our soul, can become firmly established and grounded, and can produce extraordinary and delicious fruit.”

Continuing, he said, “The spiritual pattern of small and simple things bringing forth great things produces firmness and steadfastness, deepening devotion, and more complete conversion to the Lord Jesus Christ and His gospel.”

The Prophet Joseph Smith taught this principle in words now included in the Doctrine and Covenants: “Let no man count them as small things; for there is much … pertaining to the saints, which depends upon these things” (D\&C 123:15).

In connection with the earliest attempts to establish the Church in Missouri, the Lord counseled patience for “all things must come to pass in their time” (D\&C 64:32). Then He gave this great teaching: “Wherefore, be not weary in well-doing, for ye are laying the foundation of a great work. And out of small things proceedeth that which is great” (D\&C 64:33).

I believe we all desire to follow President Russell M. Nelson’s challenge to press forward “on the covenant path.” Our commitment to do so is strengthened by consistently following the “small things” we are taught by the gospel of Jesus Christ and the leaders of His Church. I testify of Him and invoke His blessings on all who seek to keep on His covenant path, in the name of Jesus Christ, amen.

\newpage
\settalk{The Powers of the Priesthood}
\setspeaker{Dallin H. Oaks}
\settalkdate{April 2018}
\talkheading{The Powers of the Priesthood}{Dallin H. Oaks}

My beloved brethren, we have heard a revelatory announcement from President Russell M. Nelson. We have heard important elaborations by Elders Christofferson and Rasband and by President Eyring. What will yet be said, including more from President Nelson, will elaborate what you, the Lord’s leaders and priesthood holders, will now do in your responsibilities. To help with that, I will review some fundamental principles governing the priesthood you hold.

\intalkheader{I. The Priesthood}

The Melchizedek Priesthood is the divine authority God has delegated to accomplish His work “to bring to pass the … eternal life of man” (Moses 1:39). In 1829, it was conferred upon Joseph Smith and Oliver Cowdery by the Savior’s Apostles Peter, James, and John (see D\&C 27:12). It is sacred and powerful beyond our powers to describe.

The keys of the priesthood are the powers to direct the exercise of priesthood authority. Thus, when the Apostles conferred the Melchizedek Priesthood upon Joseph and Oliver, they also gave them the keys to direct its exercise (see D\&C 27:12–13). But not all priesthood keys were conferred at that time. The entire keys and knowledge necessary for this “dispensation of the fulness of times” (D\&C 128:18) are given “line upon line” (verse 21). Additional keys were given in the Kirtland Temple seven years later (see D\&C 110:11–16). These keys were given to direct priesthood authority in the additional assignments being given at that time, such as baptism for the dead.

The Melchizedek Priesthood is not a status or a label. It is a divine power held in trust to use for the benefit of God’s work for His children. We should always remember that men who hold the priesthood are not “the priesthood.” It is not appropriate to refer to “the priesthood and the women.” We should refer to “the holders of the priesthood and the women.”

\intalkheader{II. A Ministry of Service}

Now let us consider what the Lord Jesus Christ expects from those who hold His priesthood—how we are to bring souls unto Him.

President Joseph F. Smith taught: “It has truly been said that the Church is perfectly organized. The only trouble is that these organizations are not fully alive to the obligations that rest upon them. When they become thoroughly awakened to the requirements made of them, they will fulfil their duties more faithfully, and the work of the Lord will be all the stronger and more powerful and influential in the world.”

President Smith also cautioned:

“The God-given titles of honor … associated with the several offices in and orders of the Holy Priesthood, are not to be used nor considered as are the titles originated by man; they are not for adornment nor are they expressive of mastership, but rather of appointment to humble service in the work of the one Master whom we profess to serve. …

“… We are laboring for the salvation of souls, and we should feel that this is the greatest duty devolving upon us. Therefore, we should feel willing to sacrifice everything, if need be, for the love of God, the salvation of men, and the triumph of the kingdom of God upon the earth.”

\intalkheader{III. The Offices of the Priesthood}

In the Lord’s Church, the offices in the Melchizedek Priesthood have different functions. The Doctrine and Covenants refers to high priests as “standing presidents or servants over different stakes scattered abroad” (D\&C 124:134). It refers to elders as “standing ministers to [the Lord’s] church” (D\&C 124:137). Here are other teachings on these separate functions.

A high priest officiates and administers in spiritual things (see D\&C 107:10, 12). Also, as President Joseph F. Smith taught, “Inasmuch as he has been ordained a high priest, [he] should feel that he is obliged … to set an example before the old and young worthy of emulation, and to place himself in a position to be a teacher of righteousness, not only by precept but more particularly by example—giving to the younger ones the benefit of the experience of age, and thus becoming individually a power in the midst of the community in which he dwells.”

On the duties of an elder, Elder Bruce R. McConkie of the Quorum of the Twelve taught: “An elder is a minister of the Lord Jesus Christ. … He is commissioned to stand in the place and stead of his Master … in ministering to his fellowmen. He is the Lord’s agent.”

Elder McConkie criticized the idea that one is “only an elder.” “Every elder in the Church holds as much priesthood as the President of the Church … ,” he said. “What is an elder? He is a shepherd, a shepherd serving in the sheepfold of the Good Shepherd.”

In this important function to minister in the sheepfold of the Good Shepherd, there is no distinction between the offices of high priest and elder in the Melchizedek Priesthood. In the great section 107 of the Doctrine and Covenants, the Lord declares, “High priests after the order of the Melchizedek Priesthood have a right to officiate in their own standing, under the direction of the presidency, in administering spiritual things, and also in the office of an elder [or offices in the Aaronic Priesthood]” (D\&C 107:10; see also verse 12).

The most important principle for all priesthood holders is the principle taught by the Book of Mormon prophet Jacob. After he and his brother Joseph were consecrated priests and teachers of the people, he declared, “And we did magnify our office unto the Lord, taking upon us the responsibility, answering the sins of the people upon our own heads if we did not teach them the word of God with all diligence” (Jacob 1:19).

Brethren, our responsibilities as holders of the priesthood are serious matters. Other organizations can be satisfied with worldly standards of performance in delivering their messages and performing their other functions. But we who hold the priesthood of God have the divine power that even governs entrance into the celestial kingdom of God. We have the purpose and the responsibility the Lord defined in the revealed preface to the Doctrine and Covenants. We are to proclaim to the world:

“That every man might speak in the name of God the Lord, even the Savior of the world;

“That faith also might increase in the earth;

“That mine everlasting covenant might be established;

“That the fulness of my gospel might be proclaimed by the weak and the simple unto the ends of the world” (D\&C 1:20–23).

To fulfill this divine charge, we must be faithful in “magnifying” our priesthood callings and responsibilities (see D\&C 84:33). President Harold B. Lee explained what it means to magnify the priesthood: “When one becomes a holder of the priesthood, he becomes an agent of the Lord. He should think of his calling as though he were on the Lord’s errand. That is what it means to magnify the priesthood.”

Therefore, brethren, if the Lord Himself were to ask you to help one of His sons or daughters—which He has done through His servants—would you do it? And if you did, would you act as His agent, “on the Lord’s errand,” relying on His promised help?

President Lee had another teaching about magnifying the priesthood: “When you hold a magnifying glass over something it makes that thing look bigger than you could see it with the naked eye; that’s a magnifying glass. Now, … if anybody magnifies their priesthood—that is, makes it bigger than they first thought it was and more important than anyone else thought it was—that is the way you magnify your priesthood.”

Here is an example of a priesthood holder magnifying his priesthood responsibility. I heard this from Elder Jeffrey D. Erekson, my companion in a stake conference in Idaho. As a young married elder, desperately poor and feeling unable to finish his last year of college, Jeffrey decided to drop out and accept an attractive job offer. A few days later his elders quorum president came to his home. “Do you understand the significance of the priesthood keys I hold?” the elders quorum president asked. When Jeffrey said he did, the president told him that since hearing of his intention to drop out of college, the Lord had tormented him during sleepless nights to give Jeffrey this message: “As your elders quorum president, I counsel you not to drop out of college. That is a message to you from the Lord.” Jeffrey stayed in school. Years later I met him when he was a successful businessman and heard him tell an audience of priesthood holders, “That [counsel] has made all the difference in my life.”

A priesthood holder magnified his priesthood and calling, and that made “all the difference” in the life of another child of God.

\intalkheader{IV. Priesthood in the Family}

Up to now, I have been speaking of the functions of the priesthood in the Church. Now I will speak of priesthood in the family. I begin with keys. The principle that priesthood authority can be exercised only under the direction of the one who holds the keys for that function is fundamental in the Church but does not apply to the exercise of priesthood authority in the family. A father who holds the priesthood presides in his family by the authority of the priesthood he holds. He has no need to have the direction or approval of priesthood keys in order to counsel the members of his family, hold family meetings, give priesthood blessings to his wife and children, or give healing blessings to family members or others.

If fathers would magnify their priesthood in their own family, it would further the mission of the Church as much as anything else they might do. Fathers who hold the Melchizedek Priesthood should keep the commandments so they will have the power of the priesthood to give blessings to their family members. Fathers should also cultivate loving family relationships so that family members will want to ask their fathers for blessings. And parents should encourage more priesthood blessings in the family.

Fathers, function as “equal partners” of your wives, as the family proclamation teaches. And, fathers, when you are privileged to exercise the power and influence of your priesthood authority, do so “by persuasion, by long-suffering, by gentleness and meekness, and by love unfeigned” (D\&C 121:41). That high standard for the exercise of priesthood authority is most important in the family. President Harold B. Lee gave this promise just after he became President of the Church: “Never is the power of the priesthood, which you hold, more wonderful than when there is a crisis in your home, a serious illness, or some great decision that has to be made. … Vested in the power of the priesthood, which is the power of Almighty God, is the power to perform miracles if the Lord wills it so, but in order for us to use that priesthood, we must be worthy to exercise it. A failure to understand this principle is a failure to receive the blessings of holding that great priesthood.”

My beloved brethren, the magnifying of the holy priesthood you hold is vital to the work of the Lord in your families and in your Church callings.

I testify of Him whose priesthood it is. Through His atoning suffering and sacrifice and resurrection, all men and women have the assurance of immortality and the opportunity for eternal life. Each of us should be faithful and diligent in doing our part in this great work of God our Eternal Father, in the name of Jesus Christ, amen.

\newpage
\settalk{Parents and Children}
\setspeaker{Dallin H. Oaks}
\settalkdate{October 2018}
\talkheading{Parents and Children}{Dallin H. Oaks}

My dear sisters, how wonderful to have this new general conference session of women of the Church eight years and older. We have heard inspiring messages from the sister leaders and from President Henry B. Eyring. President Eyring and I love working under the direction of President Russell M. Nelson, and we look forward to his prophetic address.

\intalkheader{I.}

Children are our most precious gift from God—our eternal increase. Yet we live in a time when many women wish to have no part in the bearing and nurturing of children. Many young adults delay marriage until temporal needs are satisfied. The average age of our Church members’ marriages has increased by more than two years, and the number of births to Church members is falling. The United States and some other nations face a future of too few children maturing into adults to support the number of retiring adults. Over 40 percent of births in the United States are to unwed mothers. Those children are vulnerable. Each of these trends works against our Father’s divine plan of salvation.

\intalkheader{II.}

Latter-day Saint women understand that being a mother is their highest priority, their ultimate joy. President Gordon B. Hinckley said: “Women for the most part see their greatest fulfillment, their greatest happiness in home and family. God planted within women something divine that expresses itself in quiet strength, in refinement, in peace, in goodness, in virtue, in truth, in love. And all of these remarkable qualities find their truest and most satisfying expression in motherhood.”

He continued: “The greatest job that any woman will ever do will be in nurturing and teaching and living and encouraging and rearing her children in righteousness and truth. There is no other thing that will compare with that, regardless of what she does.”

Mothers, beloved sisters, we love you for who you are and what you do for all of us.

In his important 2015 address titled “A Plea to My Sisters,” President Russell M. Nelson said:

“The kingdom of God is not and cannot be complete without women who make sacred covenants and then keep them, women who can speak with the power and authority of God!

“Today, … we need women who know how to make important things happen by their faith and who are courageous defenders of morality and families in a sin-sick world. We need women who are devoted to shepherding God’s children along the covenant path toward exaltation; women who know how to receive personal revelation, who understand the power and peace of the temple endowment; women who know how to call upon the powers of heaven to protect and strengthen children and families; women who teach fearlessly.”

These inspired teachings are all based on “The Family: A Proclamation to the World,” in which this restored Church reaffirms doctrine and practices central to the Creator’s plan before He created the earth.

\intalkheader{III.}

Now I address the younger group of this audience. My dear young sisters, because of your knowledge of the restored gospel of Jesus Christ, you are unique. Your knowledge will enable you to endure and overcome the difficulties of growing up. From a young age, you have participated in projects and programs that have developed your talents, such as writing, speaking, and planning. You have learned responsible behavior and how to resist temptations to lie, cheat, steal, or use alcohol or drugs.

Your uniqueness was recognized in a University of North Carolina study of American teens and religion. A Charlotte Observer article had the title “Mormon teens cope best: Study finds they top peers at handling adolescence.” This article concluded that “Mormons fared best at avoiding risky behaviors, doing well in school and having a positive attitude about the future.” One of the researchers in the study, who interviewed most of our youth, said, “Across almost every category we looked at, there was a clear pattern: Mormons were first.”

Why do you cope best with the difficulties of growing up? Young women, it is because you understand our Heavenly Father’s great plan of happiness. This tells you who you are and the purpose of your life. Youth with that understanding are first in problem-solving and first in choosing the right. You know you can have the Lord’s help in overcoming all the difficulties of growing up.

Another reason why you are most effective is that you understand that you are children of a Heavenly Father who loves you. I am sure you are familiar with our great hymn “Dearest Children, God Is Near You.” Here is the first verse we all have sung and believed:

\begin{verse}
Dearest children, God is near you,\\
Watching o’er you day and night,\\
And delights to own and bless you,\\
If you strive to do what’s right.
\end{verse}

There are two teachings in that verse: First, our Heavenly Father is near us and watches over us day and night. Think of it! God loves us, He is near to us, and He watches over us. Second, He delights to bless us as we “strive to do what’s right.” What comfort in the midst of our anxieties and difficulties!

Yes, young women, you are blessed and you are wonderful, but you are like all of Heavenly Father’s children in your need to “strive to do what’s right.”

Here I could give you counsel on many different things, but I have chosen to speak of only two.

My first counsel concerns cell phones. A recent nationwide survey found that over half of teens in the United States said they spend too much time on their cell phones. More than 40 percent said they felt anxious when they were separated from their cell phones. This was more common among girls than boys. My young sisters—and adult women too—it will bless your lives if you limit your use of and dependence on cell phones.

My second counsel is even more important. Be kind to others. Kindness is something many of our youth are doing already. Some groups of youth in some communities have shown the way for all of us. We have been inspired by our young people’s acts of kindness to those in need of love and help. In many ways, you give that help and show that love to one another. We wish all would follow your example.

At the same time, we know that the adversary tempts all of us to be unkind, and there are still many examples of this, even among children and youth. Persistent unkindness is known by many names, such as bullying, ganging up on someone, or joining together to reject others. These examples deliberately inflict pain on classmates or friends. My young sisters, it is not pleasing to the Lord if we are cruel or mean to others.

Here is an example. I know of a young man, a refugee here in Utah, who was teased for being different, including sometimes speaking his native language. He was persecuted by a gang of privileged youth until he retaliated in a way that caused him to be jailed for over 70 days while being considered for deportation. I don’t know what provoked this group of youth, many of them Latter-day Saints like you, but I can see the effect of their meanness, a tragic experience and expense to one of the children of God. Small actions of unkindness can have devastating consequences.

When I heard that story, I compared it with what our prophet, President Nelson, said in his recent worldwide youth devotional. In asking you and all other youth to assist in gathering Israel, he said: “Stand out; be different from the world. You and I know that you are to be a light to the world. Therefore, the Lord needs you to look like, sound like, act like, and dress like a true disciple of Jesus Christ.”

The youth battalion President Nelson invited you to join will not be mean to one another. They will follow the Savior’s teaching to reach out and be loving and considerate of others, even to turn the other cheek when we feel someone has wronged us.

In a general conference address about the time many of you were born, President Gordon B. Hinckley praised “beautiful young women who are striving to live the gospel.” He described them, just as I feel to describe you:

“They are generous toward one another. They seek to strengthen one another. They are a credit to their parents and the homes from which they come. They are approaching womanhood and will carry throughout their lives the ideals which presently motivate them.”

As a servant of the Lord, I say to you young women, our world needs your goodness and love. Be kind to one another. Jesus taught us to love one another and to treat others as we want to be treated. As we strive to be kind, we draw closer to Him and His loving influence.

My dear sisters, if you participate in any meanness or pettiness—individually or with a group—resolve now to change and encourage others to change. That is my counsel, and I give it to you as a servant of the Lord Jesus Christ because His Spirit has prompted me to speak to you about this important subject. I testify of Jesus Christ, our Savior, who taught us to love one another as He loved us. I pray that we will do so, in the name of Jesus Christ, amen.

\newpage
\settalk{Truth and the Plan}
\setspeaker{Dallin H. Oaks}
\settalkdate{October 2018}
\talkheading{Truth and the Plan}{Dallin H. Oaks}

Modern revelation defines truth as a “knowledge of things as they are, and as they were, and as they are to come” (Doctrine and Covenants 93:24). That is a perfect definition for the plan of salvation and “The Family: A Proclamation to the World.”

We live in a time of greatly expanded and disseminated information. But not all of this information is true. We need to be cautious as we seek truth and choose sources for that search. We should not consider secular prominence or authority as qualified sources of truth. We should be cautious about relying on information or advice offered by entertainment stars, prominent athletes, or anonymous internet sources. Expertise in one field should not be taken as expertise on truth in other subjects.

We should also be cautious about the motivation of the one who provides information. That is why the scriptures warn us against priestcraft (see 2 Nephi 26:29). If the source is anonymous or unknown, the information may also be suspect.

Our personal decisions should be based on information from sources that are qualified on the subject and free from selfish motivations.

\intalkheader{I.}

When we seek the truth about religion, we should use spiritual methods appropriate for that search: prayer, the witness of the Holy Ghost, and study of the scriptures and the words of modern prophets. I am always sad when I hear of one who reports a loss of religious faith because of secular teachings. Those who once had spiritual vision can suffer from self-inflicted spiritual blindness. As President Henry B. Eyring said, “Their problem does not lie in what they think they see; it lies in what they cannot yet see.”

The methods of science lead us to what we call scientific truth. But “scientific truth” is not the whole of life. Those who do not learn “by study and also by faith” (Doctrine and Covenants 88:118) limit their understanding of truth to what they can verify by scientific means. That puts artificial limits on their pursuit of truth.

President James E. Faust said: “Those who have been [baptized] put their eternal soul at risk by carelessly pursuing only the secular source of learning. We believe that The Church of Jesus Christ of Latter-day Saints has the fulness of the gospel of Christ, which gospel is the essence of truth and eternal enlightenment.”

We find true and enduring joy by coming to know and acting upon the truth about who we are, the meaning of mortal life, and where we are going when we die. Those truths cannot be learned by scientific or secular methods.

\intalkheader{II.}

I will now speak of restored gospel truths that are fundamental to the doctrine of The Church of Jesus Christ of Latter-day Saints. Please consider these truths carefully. They explain much about our doctrine and practices, perhaps including some things not yet understood.

There is a God, who is the loving Father of the spirits of all who have ever lived or will live.

Gender is eternal. Before we were born on this earth, we all lived as male or female spirits in the presence of God.

We have just heard the Tabernacle Choir at Temple Square sing “I Will Follow God’s Plan.” That is the plan God established so that all of His spirit children could progress eternally. That plan is vital to each of us.

Under that plan, God created this earth as a place where His beloved spirit children could be born into mortality to receive a physical body and to have the opportunity for eternal progress by making righteous choices.

To be meaningful, mortal choices had to be made between contesting forces of good and evil. There had to be opposition and, therefore, an adversary, who was cast out because of rebellion and was allowed to tempt God’s children to act contrary to God’s plan.

The purpose of God’s plan was to give His children the opportunity to choose eternal life. This could be accomplished only by experience in mortality and, after death, by postmortal growth in the spirit world.

In the course of mortal life, we would all be soiled by sin as we yielded to the evil temptations of the adversary, and we would eventually die. We accepted those challenges in reliance upon the plan’s assurance that God our Father would provide a Savior, His Only Begotten Son, who would rescue us by a universal resurrection to an embodied life after death. The Savior would also provide an atonement to pay the price for all to be cleansed from sin on the conditions He prescribed. Those conditions included faith in Christ, repentance, baptism, the gift of the Holy Ghost, and other ordinances performed by priesthood authority.

God’s great plan of happiness provides a perfect balance between eternal justice and the mercy we can obtain through the Atonement of Jesus Christ. It also enables us to be transformed into new creatures in Christ.

A loving God reaches out to each of us. We know that through His love and because of the Atonement of His Only Begotten Son, “all mankind may be saved, by obedience to the laws and ordinances of [His] Gospel” (Articles of Faith 1:3; emphasis added).

The Church of Jesus Christ of Latter-day Saints is properly known as a family-centered Church. But what is not well understood is that our family-centeredness is focused on more than mortal relationships. Eternal relationships are also fundamental to our theology. “The family is ordained of God.” Under the great plan of our loving Creator, the mission of His restored Church is to help the children of God achieve the supernal blessing of exaltation in the celestial kingdom, which can be attained only through an eternal marriage between a man and a woman (see Doctrine and Covenants 131:1–3). We affirm the Lord’s teachings that “gender is an essential characteristic of individual premortal, mortal, and eternal identity and purpose” and that “marriage between man and woman is essential to His eternal plan.”

Finally, God’s love is so great that, except for the few who deliberately become sons of perdition, He has provided a destiny of glory for all of His children. “All of His children” includes all who are dead. We perform ordinances for them by proxy in our temples. The purpose of the Church of Jesus Christ is to qualify His children for the highest degree of glory, which is exaltation or eternal life. For those who do not desire or qualify for that, God has provided other, though lesser, kingdoms of glory.

Anyone who understands these eternal truths can understand why we members of The Church of Jesus Christ of Latter-day Saints think as we do and do as we do.

\intalkheader{III.}

I will now mention some applications of these eternal truths, which can be understood only in light of God’s plan.

First, we honor individual agency. Most are aware of the restored Church’s great efforts to promote religious freedom in the United States and across the world. These efforts do not promote just our own interests but, according to His plan, seek to help all of God’s children enjoy freedom to choose.

Second, we are a missionary people. We are sometimes asked why we send missionaries to so many nations, even among Christian populations. We receive the same question about why we give many millions of dollars of humanitarian aid to persons who are not members of our Church and why we do not link this aid to our missionary efforts. We do this because we esteem all mortals as children of God—our brothers and sisters—and we want to share our spiritual and temporal abundance with everyone.

Third, mortal life is sacred to us. Our commitment to God’s plan requires us to oppose abortion and euthanasia.

Fourth, some are troubled by some of our Church’s positions on marriage and children. Our knowledge of God’s revealed plan of salvation requires us to oppose current social and legal pressures to retreat from traditional marriage and to make changes that confuse or alter gender or homogenize the differences between men and women. We know that the relationships, identities, and functions of men and women are essential to accomplish God’s great plan.

Fifth, we also have a distinctive perspective on children. We look on the bearing and nurturing of children as part of God’s plan and a joyful and sacred duty of those given the power to participate in it. In our view, the ultimate treasures on earth and in heaven are our children and our posterity. Therefore, we must teach and contend for principles and practices that provide the best conditions for the development and happiness of children—all children.

Finally, we are beloved children of a Heavenly Father, who has taught us that maleness and femaleness, marriage between a man and a woman, and the bearing and nurturing of children are all essential to His great plan of happiness. Our positions on these fundamentals frequently provoke opposition to the Church. We consider that inevitable. Opposition is part of the plan, and Satan’s most strenuous opposition is directed at whatever is most important to God’s plan. He seeks to destroy God’s work. His prime methods are to discredit the Savior and His divine authority, to erase the effects of the Atonement of Jesus Christ, to discourage repentance, to counterfeit revelation, and to contradict individual accountability. He also seeks to confuse gender, to distort marriage, and to discourage childbearing—especially by parents who will raise children in truth.

\intalkheader{IV.}

The work of the Lord is going forward despite the organized and constant opposition that confronts us as we strive to practice the teachings of The Church of Jesus Christ of Latter-day Saints. For those who falter under that opposition, I offer these suggestions.

Remember the principle of repentance made possible by the power of the Atonement of Jesus Christ. As Elder Neal A. Maxwell urged, don’t be among those “who would rather try to change the Church than to change themselves.”

As Elder Jeffrey R. Holland urged:

“Hold fast to what you already know and stand strong until additional knowledge comes. …

“… In this Church, what we know will always trump what we do not know.”

Exercise faith in the Lord Jesus Christ, which is the first principle of the gospel.

Finally, seek help. Our Church leaders love you and seek spiritual guidance to help you. We provide many resources such as you will find through LDS.org and other supports for gospel study in the home. We also have ministering brothers and sisters called to give loving assistance.

Our loving Heavenly Father wants His children to have the joy that is the purpose of our creation. That joyful destiny is eternal life, which we can obtain by pressing forward along what our prophet, President Russell M. Nelson, often calls “the covenant path.” Here is what he said in his first message as President of the Church: “Keep on the covenant path. Your commitment to follow the Savior by making covenants with Him and then keeping those covenants will open the door to every spiritual blessing and privilege available to men, women, and children everywhere.”

I solemnly testify that the things I have said are true, and they are made possible by the teachings and the Atonement of Jesus Christ, who makes it all possible under the great plan of God, our Eternal Father. In the name of Jesus Christ, amen.

\newpage
\settalk{Where Will This Lead?}
\setspeaker{Dallin H. Oaks}
\settalkdate{April 2019}
\talkheading{Where Will This Lead?}{Dallin H. Oaks}

The restored gospel of Jesus Christ encourages us to think about the future. It explains the purpose of mortal life and the reality of the life to follow. It teaches great ideas about the future to guide our actions today.

In contrast, we all know persons who are concerned only with the present: spend it today, enjoy it today, and take no thought for the future.

Our present and our future will be happier if we are always conscious of the future. As we make current decisions, we should always be asking, “Where will this lead?”

\intalkheader{I.}

Some decisions are choices between doing something or doing nothing. I heard an example of this kind of choice at a stake conference in the United States many years ago.

The setting was a beautiful college campus. A crowd of young students was seated on the grass. The speaker who described this circumstance said they were watching a handsome tree squirrel with a large, bushy tail playing around the base of a beautiful hardwood tree. Sometimes it was on the ground, sometimes up and down and around the trunk. But why would that familiar sight attract a crowd of students?

Stretched out prone on the grass nearby was an Irish setter. He was the object of the students’ interest, and the squirrel was the object of his. Each time the squirrel was momentarily out of sight circling the tree, the setter would quietly creep forward a few inches and then resume his apparently indifferent posture. This was what held the students’ interest. Silent and immobile, their eyes were riveted on the event whose outcome was increasingly obvious.

Finally, the setter was close enough to bound at the squirrel and catch it in his mouth. A gasp of horror arose, and the crowd of students surged forward and wrested the little animal away from the dog, but it was too late. The squirrel was dead.

Anyone in that crowd could have warned the squirrel at any time by waving his or her arms or crying out, but none did. They just watched while the inevitable outcome got closer and closer. No one asked, “Where will this lead?” When the predictable occurred, all rushed to prevent the outcome, but it was too late. Tearful regret was all they could offer.

That true story is a parable of sorts. It applies to things we see in our own lives and in lives and circumstances around us. As we see threats creeping up on persons or things we love, we have the choice of speaking or acting or remaining silent. It is well to ask ourselves, “Where will this lead?” Where the consequences are immediate and serious, we cannot afford to do nothing. We must sound appropriate warnings or support appropriate preventive efforts while there is still time.

The decisions I have just described involve choices between taking some action or taking no action at all. More common are those choices between one action or another. These include choices between good or evil, but more frequently they are choices between two goods. Here too it is desirable to ask where this will lead. We make many choices between two goods, often involving how we will spend our time. There is nothing bad about playing video games or texting or watching TV or talking on a cell phone. But each of these involves what is called “opportunity cost,” meaning that if we spend time doing one thing, we lose the opportunity to do another. I am sure you can see that we need to measure thoughtfully what we are losing by the time we spend on one activity, even if it is perfectly good in itself.

Some time ago I gave a talk titled “Good, Better, Best.” In that talk I said that “just because something is good is not a sufficient reason for doing it. The number of good things we can do far exceeds the time available to accomplish them. Some things are better than good, and these are the things that should command priority attention in our lives. … We have to forego some good things in order to choose others that are better or best.”

Take the long view. What is the effect on our future of the decisions we make in the present? Remember the importance of getting an education, studying the gospel, renewing our covenants by partaking of the sacrament, and attending the temple.

\intalkheader{II.}

“Where will this lead?” is also important in choosing how we label or think of ourselves. Most important, each of us is a child of God with a potential destiny of eternal life. Every other label, even including occupation, race, physical characteristics, or honors, is temporary or trivial in eternal terms. Don’t choose to label yourselves or think of yourselves in terms that put a limit on a goal for which you might strive.

My brethren, and my sisters who may view or read what I say here, I hope you know why your leaders give the teachings and counsel we give. We love you, and our Heavenly Father and His Son, Jesus Christ, love you. Their plan for us is the “great plan of happiness” (Alma 42:8). That plan and Their commandments and ordinances and covenants lead us to the greatest happiness and joy in this life and in the life to come. As servants of the Father and the Son, we teach and counsel as They have directed us by the Holy Ghost. We have no desire other than to speak what is true and to encourage you to do what They have outlined as the pathway to eternal life, “the greatest of all the gifts of God” (Doctrine and Covenants 14:7).

\intalkheader{III.}

Here is another example of the effect on the future of decisions made in the present. This example concerns the choice to make a present sacrifice to achieve an important future goal.

At a stake conference in Cali, Colombia, a sister told how she and her fiancé desired to be married in the temple, but at that time the closest temple was in faraway Peru. For a long time, they saved their money for the bus fares. Finally they boarded the bus to Bogotá, but when they arrived there, they learned that all seats on the bus to Lima, Peru, were taken. They could go home without being married or be married out of the temple. Fortunately, there was one other alternative. They could ride on the bus to Lima if they were willing to sit on the floor of the bus for the entire five-day and five-night ride. They chose to do this. She said it was difficult, even though some riders sometimes let them sit in their seats so they could stretch out on the floor.

What impressed me in her talk was this sister’s statement that she was grateful she and her husband had been able to go to the temple in this way, because it changed the way they felt about the gospel and the way they felt about marriage in the temple. The Lord had rewarded them with the growth that comes from sacrifice. She also observed that their five-day trip to the temple accomplished a great deal more in building their spirituality than many visits to the temple that were sacrifice-free.

In the years since I heard that testimony, I have wondered how different that young couple’s life would have been if they had made another choice—forgoing the sacrifice necessary to be married in the temple.

Brethren, we make countless choices in life, some large and some seemingly small. Looking back, we can see what a great difference some of our choices made in our lives. We make better choices and decisions if we look at the alternatives and ponder where they will lead. As we do, we will be following President Russell M. Nelson’s counsel to begin with the end in mind. For us, the end is always on the covenant path through the temple to eternal life, the greatest of all the gifts of God.

I testify of Jesus Christ and of the effects of His Atonement and the other truths of His everlasting gospel in the name of Jesus Christ, amen.

\newpage
\settalk{Cleansed by Repentance}
\setspeaker{Dallin H. Oaks}
\settalkdate{April 2019}
\talkheading{Cleansed by Repentance}{Dallin H. Oaks}

In mortality we are subject to the laws of man and the laws of God. I have had the unusual experience of judging serious misbehavior under both of these laws—earlier as a justice of the Utah Supreme Court and now as a member of the First Presidency. The contrast I have experienced between the laws of man and the laws of God has increased my appreciation for the reality and power of the Atonement of Jesus Christ. Under the laws of man, a person guilty of the most serious crimes can be sentenced to life in prison without possibility of parole. But it is different under the merciful plan of a loving Heavenly Father. I have witnessed that these same serious sins can be forgiven in mortality because of our Savior’s atoning sacrifice for the sins of “all those who have a broken heart and a contrite spirit” (2 Nephi 2:7). Christ redeems, and His Atonement is real.

The loving compassion of our Savior is expressed in the great hymn just performed by the choir.

\begin{verse}
Come unto Jesus; He’ll ever heed you,\\
Though in the darkness you’ve gone astray.\\
His love will find you and gently lead you\\
From darkest night into day.
\end{verse}

The atoning sacrifice of Jesus Christ opens the door for “all men [to] repent and come unto him” (Doctrine and Covenants 18:11; see also Mark 3:28; 1 Nephi 10:18; Alma 34:8, 16). The book of Alma reports repentance and forgiveness even of those who had been a wicked and a bloodthirsty people (see Alma 25:16; 27:27, 30). My message today is one of hope for all of us, including those who have lost their membership in the Church by excommunication or name removal. We are all sinners who can be cleansed by repentance. “To repent from sin is not easy,” Elder Russell M. Nelson taught in a prior general conference. “But the prize is worth the price.”

\intalkheader{I. Repentance}

Repentance begins with our Savior, and it is a joy, not a burden. In last December’s Christmas devotional, President Nelson taught: “True repentance is not an event. It is a never-ending privilege. It is fundamental to progression and having peace of mind, comfort, and joy.”

Some of the greatest teachings on repentance are in Alma’s Book of Mormon sermon to members of the Church whom he later described as having been in a state of “much unbelief,” “lifted up in … pride,” and with hearts set “upon riches and the vain things of the world” (Alma 7:6). Each member of this restored Church has much to learn from Alma’s inspired teachings.

We begin with faith in Jesus Christ, because “it is he that cometh to take away the sins of the world” (Alma 5:48). We must repent because, as Alma taught, “except ye repent ye can in nowise inherit the kingdom of heaven” (Alma 5:51). Repentance is an essential part of God’s plan. Because all would sin in our mortal experience and be cut off from God’s presence, man could not “be saved” without repentance (Alma 5:31; see also Helaman 12:22).

This has been taught from the beginning. The Lord commanded Adam, “Teach it unto your children, that all men, everywhere, must repent, or they can in nowise inherit the kingdom of God, for no unclean thing can dwell there, or dwell in his presence” (Moses 6:57). We must repent of all our sins—all of our actions or inactions contrary to the commandments of God. No one is exempt. Just last evening President Nelson challenged us, “Brethren, we all need to repent.”

To be cleansed by repentance, we must forsake our sins and confess them to the Lord and to His mortal judge where required (see Doctrine and Covenants 58:43). Alma taught that we must also “bring forth works of righteousness” (Alma 5:35). All of this is part of the frequent scriptural invitation to come unto Christ.

We need to partake of the sacrament each Sabbath day. In that ordinance we make covenants and receive blessings that help us overcome all acts and desires that block us from the perfection our Savior invites us to achieve (see Matthew 5:48; 3 Nephi 12:48). As we “deny [ourselves] of all ungodliness, and love God with all [our] might, mind and strength,” then we may “be perfect in Christ” and be “sanctified” through the shedding of His blood, to “become holy, without spot” (Moroni 10:32–33). What a promise! What a miracle! What a blessing!

\intalkheader{II. Accountability and Mortal Judgments}

One purpose of God’s plan for this mortal experience is to “prove” us “to see if [we] will do all things whatsoever the Lord [our] God shall command [us]” (Abraham 3:25). As part of this plan, we are accountable to God and to His chosen servants, and that accountability involves both mortal and divine judgments.

In the Lord’s Church, mortal judgments for members or prospective members are administered by leaders who seek divine direction. It is their responsibility to judge persons who are seeking to come unto Christ to receive the power of His Atonement on the covenant path to eternal life. Mortal judgments determine whether a person is ready for baptism. Is a person worthy of a recommend to attend the temple? Has a person whose name was removed from the records of the Church repented sufficiently through the Atonement of Jesus Christ to be readmitted by baptism?

When a mortal judge called of God approves a person for further progress, such as temple privileges, he is not signifying that person as perfect, and he is not forgiving any sins. Elder Spencer W. Kimball taught that after what he called the mortal “waiving [of] penalties,” a person “must also seek and secure from the God of heaven a final repentance, and only he can absolve.” And if sinful acts and desires remain unrepented until the Final Judgment, an unrepentant person will remain unclean. The ultimate accountability, including the final cleansing effect of repentance, is between each of us and God.

\intalkheader{III. Resurrection and the Final Judgment}

The judgment most commonly described in the scriptures is the Final Judgment that follows the Resurrection (see 2 Nephi 9:15). Many scriptures state that “we shall all stand before the judgment seat of Christ” (Romans 14:10; see also 2 Nephi 9:15; Mosiah 27:31) “to be judged according to the deeds [that] have been done in the mortal body” (Alma 5:15; see also Revelation 20:12; Alma 41:3; 3 Nephi 26:4). All will be judged “according to their works” (3 Nephi 27:15) and “according to the desire[s] of their hearts” (Doctrine and Covenants 137:9; see also Alma 41:6).

The purpose of this Final Judgment is to determine whether we have achieved what Alma described as a “mighty change of heart” (see Alma 5:14, 26), where we have become new creatures, with “no more disposition to do evil, but to do good continually” (Mosiah 5:2). The judge of this is our Savior, Jesus Christ (see John 5:22; 2 Nephi 9:41). After His judgment we will all confess “that his judgments are just” (Mosiah 16:1; see also Mosiah 27:31; Alma 12:15), because His omniscience (see 2 Nephi 9:15, 20) has given Him a perfect knowledge of all of our acts and desires, both those righteous or repented and those unrepented or unchanged.

The scriptures describe the process of this Final Judgment. Alma teaches that the justice of our God requires that in the Resurrection “all things should be restored to their proper order” (Alma 41:2). This means that “if their works were good in this life, and the desires of their hearts were good, … at the last day [they will] be restored unto that which is good” (Alma 41:3). Similarly, “if their works [or their desires] are evil they shall be restored [to] them for evil” (Alma 41:4–5; see also Helaman 14:31). Similarly, the prophet Jacob taught that in the Final Judgment “they who are righteous shall be righteous still, and they who are filthy shall be filthy still” (2 Nephi 9:16; see also Mormon 9:14; 1 Nephi 15:33). That is the process preceding our standing before what Moroni calls “the pleasing bar of the great Jehovah, the Eternal Judge of both quick and dead” (Moroni 10:34; see also 3 Nephi 27:16).

To assure that we will be clean before God, we must repent before the Final Judgment (see Mormon 3:22). As Alma told his sinful son, we cannot hide our sins before God, “and except ye repent they will stand as a testimony against you at the last day” (Alma 39:8; emphasis added). The Atonement of Jesus Christ gives us the only way to achieve the needed cleansing through repentance, and this mortal life is the time to do it. Although we are taught that some repentance can occur in the spirit world (see Doctrine and Covenants 138:31, 33, 58), that is not as certain. Elder Melvin J. Ballard taught: “It is much easier to overcome and serve the Lord when both flesh and spirit are combined as one. This is the time when men are more pliable and susceptible. … This life is the time to repent.”

When we repent, we have the Lord’s assurance that our sins, including our acts and desires, will be cleansed and our merciful final judge will “remember them no more” (Doctrine and Covenants 58:42; see also Isaiah 1:18; Jeremiah 31:34; Hebrews 8:12; Alma 41:6; Helaman 14:18–19). Cleansed by repentance, we can qualify for eternal life, which King Benjamin described as “dwell[ing] with God in a state of never-ending happiness” (Mosiah 2:41; see also Doctrine and Covenants 14:7).

As another part of God’s “plan of restoration” (Alma 41:2), the Resurrection will restore “all things … to their proper and perfect frame” (Alma 40:23). This includes the perfection of all of our physical deficiencies and deformities acquired in mortality, including at birth or by trauma or disease.

Does this restoration perfect us of all our unholy or unconquered desires or addictions? That cannot be. We know from modern revelation that we will be judged for our desires as well as our actions (see Alma 41:5; Doctrine and Covenants 137:9) and that even our thoughts will condemn us (see Alma 12:14). We must not “procrastinate the day of [our] repentance” until death, Amulek taught (Alma 34:33), because the same spirit that has possessed our body in this life—whether the Lord’s or the devil’s—“will have power to possess [our] body in that eternal world” (Alma 34:34). Our Savior has the power and stands ready to cleanse us from evil. Now is the time to seek His help to repent of our wicked or unseemly desires and thoughts to be clean and prepared to stand before God at the Final Judgment.

\intalkheader{IV. The Arms of Mercy}

Overarching God’s plan and all of His commandments is His love for each of us, which is “most desirable above all things … and the most joyous to the soul” (1 Nephi 11:22–23). The prophet Isaiah assured even the wicked that when they “return unto the Lord, … he will have mercy … [and] abundantly pardon” (Isaiah 55:7). Alma taught, “Behold, he sendeth an invitation unto all men, for the arms of mercy are extended towards them” (Alma 5:33; see also 2 Nephi 26:25–33). The risen Lord told the Nephites, “Behold, mine arm of mercy is extended towards you, and whosoever will come, him will I receive” (3 Nephi 9:14). From these and many other scriptural teachings, we know that our loving Savior opens His arms to receive all men and women on the loving conditions He has prescribed to enjoy the greatest blessings God has for His children.

Because of God’s plan and the Atonement of Jesus Christ, I testify with a “perfect brightness of hope” that God loves us and we can be cleansed by the process of repentance. We are promised that “if [we] press forward, feasting upon the word of Christ, and endure to the end, behold, thus saith the Father: Ye shall have eternal life” (2 Nephi 31:20). May we all do so, I plead and pray in the name of Jesus Christ, amen.

\newpage
\settalk{Trust in the Lord}
\setspeaker{Dallin H. Oaks}
\settalkdate{October 2019}
\talkheading{Trust in the Lord}{Dallin H. Oaks}

My dear brothers and sisters, a letter I received some time ago introduces the subject of my talk. The writer was contemplating a temple marriage to a man whose eternal companion had died. She would be a second wife. She asked this question: would she be able to have her own house in the next life, or would she have to live with her husband and his first wife? I just told her to trust the Lord.

I continue with an experience I heard from a valued associate, which I share with his permission. After the death of his beloved wife and the mother of his children, a father remarried. Some grown children strongly objected to the remarriage and sought the counsel of a close relative who was a respected Church leader. After hearing the reasons for their objections, which focused on conditions and relationships in the spirit world or in the kingdoms of glory that follow the Final Judgment, this leader said: “You are worried about the wrong things. You should be worried about whether you will get to those places. Concentrate on that. If you get there, all of it will be more wonderful than you can imagine.”

What a comforting teaching! Trust in the Lord!

From letters I have received, I know that others are troubled by questions about the spirit world we will inhabit after we die and before we are resurrected. Some assume that the spirit world will continue many of the temporal circumstances and issues we experience in this mortal life. What do we really know about conditions in the spirit world? I believe a BYU religion professor’s article on this subject had it right: “When we ask ourselves what we know about the spirit world from the standard works, the answer is ‘not as much as we often think.’”

Of course, we know from the scriptures that after our bodies die we continue to live as spirits in the spirit world. The scriptures also teach that this spirit world is divided between those who have been “righteous” or “just” during life and those who have been wicked. They also describe how some faithful spirits teach the gospel to those who have been wicked or rebellious (see 1 Peter 3:19; Doctrine and Covenants 138:19–20, 29, 32, 37). Most important, modern revelation reveals that the work of salvation goes forward in the spirit world (see Doctrine and Covenants 138:30–34, 58), and although we are urged not to procrastinate our repentance during mortality (see Alma 13:27), we are taught that some repentance is possible there (see Doctrine and Covenants 138:58).

The work of salvation in the spirit world consists of freeing spirits from what the scriptures frequently describe as “bondage.” All in the spirit world are under some form of bondage. President Joseph F. Smith’s great revelation, canonized in section 138 of the Doctrine and Covenants, states that the righteous dead, who were in a state of “peace” (Doctrine and Covenants 138:22) as they anticipated the Resurrection (see Doctrine and Covenants 138:16), “had looked upon the long absence of their spirits from their bodies as a bondage” (Doctrine and Covenants 138:50).

The wicked also suffer an additional bondage. Because of unrepented sins, they are in what the Apostle Peter referred to as spirit “prison” (1 Peter 3:19; see also Doctrine and Covenants 138:42). These spirits are described as “bound” or as “captives” (Doctrine and Covenants 138:31, 42) or as “cast out into outer darkness” with “weeping, and wailing, and gnashing of teeth” as they await resurrection and judgment (Alma 40:13–14).

Resurrection for all in the spirit world is assured by the Resurrection of Jesus Christ (see 1 Corinthians 15:22), though it occurs at different times for different groups. Until that appointed time, what the scriptures tell us about activity in the spirit world principally concerns the work of salvation. Little else is revealed. The gospel is preached to the ignorant, the unrepentant, and the rebellious so they can be freed from their bondage and go forward to the blessings a loving Heavenly Father has in store for them.

The spirit-world bondage that applies to righteous converted souls is their need to await—and perhaps even be allowed to prompt—the performance of their proxy ordinances on earth so they can be baptized and enjoy the blessings of the Holy Ghost (see Doctrine and Covenants 138:30–37, 57–58). These mortal proxy ordinances also empower them to go forward under priesthood authority to enlarge the hosts of the righteous who can preach the gospel to the spirits in prison.

Beyond these basics, our canon of scripture contains very little about the spirit world that follows death and precedes the Final Judgment. So what else do we know about the spirit world? Many members of the Church have had visions or other inspirations to inform them about how things operate or are organized in the spirit world, but these personal spiritual experiences are not to be understood or taught as the official doctrine of the Church. And, of course, there is abundant speculation by members and others in published sources like books on near-death experiences.

As to all of these, the wise cautions of Elders D. Todd Christofferson and Neil L. Andersen in earlier general conference messages are important to remember. Elder Christofferson taught: “It should be remembered that not every statement made by a Church leader, past or present, necessarily constitutes doctrine. It is commonly understood in the Church that a statement made by one leader on a single occasion often represents a personal, though well-considered, opinion, not meant to be official or binding for the whole Church.”

In the following conference, Elder Andersen taught this principle: “The doctrine is taught by all 15 members of the First Presidency and Quorum of the Twelve. It is not hidden in an obscure paragraph of one talk.” The family proclamation, signed by all 15 prophets, seers, and revelators, is a wonderful illustration of that principle.

Beyond something as formal as the family proclamation, the prophetic teachings of the Presidents of the Church, affirmed by other prophets and apostles, are also an example of this. As to circumstances in the spirit world, the Prophet Joseph Smith gave two teachings near the close of his ministry that have been frequently taught by his successors. One of these is his teaching in the King Follett sermon that family members who were righteous will be together in the world of spirits. Another is this statement at a funeral in the last year of his life: “The spirits of the just are exalted to a greater and more glorious work … [in] the world of spirits. … They are not far from us, and know and understand our thoughts, feelings, and motions, and are often pained therewith.”

So, what about a question like I mentioned earlier about where spirits live? If that question seems strange or trivial to you, consider many of your own questions, or even those you have been tempted to answer on the basis of something you heard from another person sometime in the past. For all questions about the spirit world, I suggest two answers. First, remember that God loves His children and will surely do what is best for each of us. Second, remember this familiar Bible teaching, which has been most helpful to me on a multitude of unanswered questions:

“Trust in the Lord with all thine heart; and lean not unto thine own understanding.

“In all thy ways acknowledge him, and he shall direct thy paths” (Proverbs 3:5–6).

Similarly, Nephi concluded his great psalm with these words: “O Lord, I have trusted in thee, and I will trust in thee forever. I will not put my trust in the arm of flesh” (2 Nephi 4:34).

We can all wonder privately about circumstances in the spirit world or even discuss these or other unanswered questions in family or other intimate settings. But let us not teach or use as official doctrine what does not meet the standards of official doctrine. To do so does not further the work of the Lord and may even discourage individuals from seeking their own comfort or edification through the personal revelation the Lord’s plan provides for each of us. Excessive reliance on personal teachings or speculations may even draw us aside from concentrating on learning and efforts that will further our understanding and help us go forward on the covenant path.

Trust in the Lord is a familiar and true teaching in The Church of Jesus Christ of Latter-day Saints. That was Joseph Smith’s teaching when the early Saints experienced severe persecutions and seemingly insurmountable obstacles. That is still the best principle we can use when our efforts to learn or our attempts to find comfort encounter obstacles in matters not yet revealed or not adopted as the official doctrine of the Church.

That same principle applies to unanswered questions about sealings in the next life or desired readjustments because of events or transgressions in mortality. There is so much we do not know that our only sure reliance is to trust in the Lord and His love for His children.

In conclusion, what we do know about the spirit world is that the Father’s and the Son’s work of salvation continues there. Our Savior initiated the work of declaring liberty to the captives (see 1 Peter 3:18–19; 4:6; Doctrine and Covenants 138:6–11, 18–21, 28–37), and that work continues as worthy and qualified messengers continue to preach the gospel, including repentance, to those who still need its cleansing effect (see Doctrine and Covenants 138:57). The object of all of that is described in the official doctrine of the Church, given in modern revelation.

“The dead who repent will be redeemed, through obedience to the ordinances of the house of God,

“And after they have paid the penalty of their transgressions, and are washed clean, shall receive a reward according to their works, for they are heirs of salvation” (Doctrine and Covenants 138:58–59).

The duty of each of us is to teach the doctrine of the restored gospel, keep the commandments, love and help one another, and do the work of salvation in the holy temples.

I testify of the truth of what I have said here and of the truths taught and to be taught in this conference. This is all made possible because of the Atonement of Jesus Christ. As we know from modern revelation, He “glorifies the Father, and saves all the works of his hands” (Doctrine and Covenants 76:43; emphasis added). In the name of Jesus Christ, amen.

\newpage
\settalk{TwoGreat Commandments}
\setspeaker{Dallin H. Oaks}
\settalkdate{October 2019}
\talkheading{TwoGreat Commandments}{Dallin H. Oaks}

My dear sisters in the gospel of Jesus Christ, I greet you as divinely assigned guardians of the eternal family. President Russell M. Nelson has taught us, “This Church was restored so that families could be formed, sealed, and exalted eternally.” That teaching has important implications for persons who identify as lesbian, gay, bisexual, or transgender, commonly referred to as LGBT. President Nelson has also reminded us that we don’t “have to [always] agree with each other to love each other.” These prophetic teachings are important for family discussions to answer the questions of children and youth. I have prayerfully sought inspiration to speak to this audience because you are uniquely affected by these questions, which directly or indirectly affect every family in the Church.

\intalkheader{I.}

I begin with what Jesus taught were the two great commandments.

“Thou shalt love the Lord thy God with all thy heart, and with all thy soul, and with all thy mind.

“This is the first and great commandment.

“And the second is like unto it, Thou shalt love thy neighbour as thyself.”

This means we are commanded to love everyone, since Jesus’s parable of the good Samaritan teaches that everyone is our neighbor. But our zeal to keep this second commandment must not cause us to forget the first, to love God with all our heart, soul, and mind. We show that love by “keep[ing] [His] commandments.” God requires us to obey His commandments because only through that obedience, including repentance, can we return to live in His presence and become perfect as He is.

In his recent talk to the young adults of the Church, President Russell M. Nelson spoke of what he called the “strong connection between God’s love and His laws.” The laws that apply most significantly to the issues relating to those identifying as LGBT are God’s law of marriage and its companion law of chastity. Both are essential in our Father in Heaven’s plan of salvation for His children. As President Nelson taught, “God’s laws are motivated entirely by His infinite love for us and His desire for us to become all we can become.”

President Nelson taught: “Many countries … have legalized same-sex marriage. As members of the Church, we respect the laws of the land … , including civil marriage. The truth is, however, that in the beginning … marriage was ordained by God! And to this day it is defined by Him as being between a man and a woman. God has not changed His definition of marriage.”

President Nelson continued: “God has also not changed His law of chastity. Requirements to enter the temple have not changed.”

President Nelson reminded all of us that “our commission as Apostles is to teach nothing but truth. That commission does not give [Apostles] the authority to modify divine law.” Thus, my sisters, the leaders of the Church must always teach the unique importance of marriage between a man and a woman and the related law of chastity.

\intalkheader{II.}

The work of The Church of Jesus Christ of Latter-day Saints is ultimately concerned with preparing the children of God for the celestial kingdom, and most particularly for its highest glory, exaltation or eternal life. That highest destiny is possible only through marriage for eternity. Eternal life includes the creative powers inherent in the combination of male and female—what modern revelation describes as the “continuation of the seeds forever and ever.”

In his talk to young adults, President Nelson taught, “Abiding by God’s laws will keep you safe as you progress toward eventual exaltation”—that is, to become like God, with the exalted life and divine potential of our Heavenly Parents. That is the destiny we desire for all we love. Because of that love, we cannot let our love supersede the commandments and the plan and work of God, which we know will bring those we love their greatest happiness.

But there are many we love, including some who have the restored gospel, who do not believe in or choose not to follow God’s commandments about marriage and the law of chastity. What about them?

God’s doctrine shows that all of us are His children and that He has created us to have joy. Modern revelation teaches that God has provided a plan for a mortal experience in which all can choose obedience to seek His highest blessings or make choices that lead to one of the less glorious kingdoms. Because of God’s great love for all of His children, those lesser kingdoms are still more wonderful than mortals can comprehend. The Atonement of Jesus Christ makes all of this possible, as He “glorifies the Father, and saves all the works of his hands.”

\intalkheader{III.}

I have spoken of the first commandment, but what of the second? How do we keep the commandment to love our neighbors? We seek to persuade our members that those who follow lesbian, gay, bisexual, or transgender teachings and actions should be treated with the love our Savior commands us to show toward all our neighbors. Thus, when same-sex marriage was declared legal in the United States, the First Presidency and Quorum of the Twelve declared: “The gospel of Jesus Christ teaches us to love and treat all people with kindness and civility—even when we disagree. We affirm that those who avail themselves of laws or court rulings authorizing same-sex marriage should not be treated disrespectfully.”

Further, we must never persecute those who do not share our beliefs and commitments. Regretfully, some persons facing these issues continue to feel marginalized and rejected by some members and leaders in our families, wards, and stakes. We must all strive to be kinder and more civil.

\intalkheader{IV.}

For reasons we do not understand, we have different challenges in our mortal experiences. But we do know that God will help each of us overcome these challenges if we sincerely seek His help. After suffering and repenting for violations of laws we have been taught, we are all destined for a kingdom of glory. The ultimate and final judgment will be by the Lord, who alone has the required knowledge, wisdom, and grace to judge each of us.

Meanwhile, we must try to keep both of the great commandments. To do so, we walk a fine line between law and love—keeping the commandments and walking the covenant path, while loving our neighbors along the way. This walk requires us to seek divine inspiration on what to support and what to oppose and how to love and listen respectfully and teach in the process. Our walk demands that we not compromise on commandments but show forth a full measure of understanding and love. Our walk must be considerate of children who are uncertain about their sexual orientation, but it discourages premature labeling because, in most children, such uncertainty decreases significantly over time. Our walk opposes recruitment away from the covenant path, and it denies support to any who lead people away from the Lord. In all of this we remember that God promises hope and ultimate joy and blessings for all who keep His commandments.

\intalkheader{V.}

Mothers and fathers and all of us are responsible to teach both of the two great commandments. For the women of the Church, President Spencer W. Kimball described that duty in this great prophecy: “Much of the major growth that is coming to the Church in the last days will come because many of the good women of the world … will be drawn to the Church in large numbers. This will happen to the degree that the women of the Church reflect righteousness and articulateness in their lives and to the degree that the women of the Church are seen as distinct and different … from the women of the world. … Thus it will be that female exemplars of the Church will be a significant force in both the numerical and the spiritual growth of the Church in the last days.”

Speaking of that prophecy, President Russell M. Nelson declared that “the day that President Kimball foresaw is today. You are the women he foresaw!” Little did we who heard that prophecy 40 years ago realize that among those the women of this Church may save will be their own dear friends and family who are currently influenced by worldly priorities and devilish distortions. My prayer and blessing is that you will teach and act to fulfill that prophecy, in the name of Jesus Christ, amen.

\newpage
\settalk{The Melchizedek Priesthood and the Keys}
\setspeaker{Dallin H. Oaks}
\settalkdate{April 2020}
\talkheading{The Melchizedek Priesthood and the Keys}{Dallin H. Oaks}

I have chosen to speak further about the priesthood of God, the subject already addressed by three earlier speakers who taught us about how the priesthood blesses the lives of women, young women, and young men.

The priesthood is a divine power and authority held in trust to be used for God’s work for the benefit of all of His children. Priesthood is not those who have been ordained to a priesthood office or those who exercise its authority. Men who hold the priesthood are not the priesthood. While we should not refer to ordained men as the priesthood, it is appropriate to refer to them as holders of the priesthood.

The power of the priesthood exists both in the Church and in the family organization. But priesthood power and priesthood authority function differently in the Church than they do in the family. All of this is according to the principles the Lord has established. The purpose of God’s plan is to lead His children to eternal life. Mortal families are essential to that plan. The Church exists to provide the doctrine, the authority, and the ordinances necessary to perpetuate family relationships into the eternities. Thus, the family organization and the Church of Jesus Christ have a mutually reinforcing relationship. The blessings of the priesthood—such as the fulness of the gospel and ordinances like baptism, confirmation and receiving the gift of the Holy Ghost, the temple endowment, and eternal marriage—are available to men and women alike.

The priesthood we speak of here is the Melchizedek Priesthood, restored at the beginning of the Restoration of the gospel. Joseph Smith and Oliver Cowdery were ordained by Peter, James, and John, who declared themselves “as possessing the keys of the kingdom, and of the dispensation of the fulness of times” (Doctrine and Covenants 128:20). These senior Apostles received that authority from the Savior Himself. All other authorities or offices in the priesthood are appendages to the Melchizedek Priesthood (see Doctrine and Covenants 107:5), for it “holds the right of presidency, and has power and authority over all the offices in the church in all ages of the world” (Doctrine and Covenants 107:8).

In the Church the authority of the greater priesthood, the Melchizedek Priesthood, and the lesser or Aaronic Priesthood is exercised under the direction of a priesthood leader, like a bishop or president, who holds the keys of that priesthood. To understand the exercise of priesthood authority in the Church, we must understand the principle of priesthood keys.

The Melchizedek Priesthood keys of the kingdom were conferred by Peter, James, and John, but that did not complete the restoration of priesthood keys. Some keys of the priesthood came later. Following the dedication of the first temple of this dispensation in Kirtland, Ohio, three prophets—Moses, Elias, and Elijah—restored “the keys of this dispensation,” including keys pertaining to the gathering of Israel and the work of the temples of the Lord (see Doctrine and Covenants 110), as President Eyring has just described so persuasively.

The most familiar example of the function of keys is in the performance of priesthood ordinances. An ordinance is a solemn act signifying the making of covenants and the promising of blessings. In the Church all ordinances are performed under the authorization of the priesthood leader who holds the keys for that ordinance.

An ordinance is most commonly officiated by persons who have been ordained to an office in the priesthood acting under the direction of one who holds priesthood keys. For example, the holders of the various offices of the Aaronic Priesthood officiate in the ordinance of the sacrament under the keys and direction of the bishop, who holds the keys of the Aaronic Priesthood. The same principle applies to the priesthood ordinances in which women officiate in the temple. Though women do not hold an office in the priesthood, they perform sacred temple ordinances under the authorization of the president of the temple, who holds the keys for the ordinances of the temple.

Another example of priesthood authority under the direction of one who holds the keys are the teachings of men and women called to teach the gospel, whether in classes in their home wards or in the mission field. Other examples are those who hold leadership positions in the ward and exercise priesthood authority in their leadership by reason of their callings and under the setting apart and direction of the priesthood leader who holds the keys in the ward or the stake. This is how the authority and power of the priesthood is exercised and enjoyed in The Church of Jesus Christ of Latter-day Saints.

Priesthood authority is also exercised and its blessings realized in the families of Latter-day Saints. By families I mean a priesthood-holding man and a woman who are married and their children. I also include the variations from the ideal relationships such as caused by death or divorce.

The principle that priesthood authority can be exercised only under the direction of the one who holds the keys for that function is fundamental in the Church, but this does not apply in the family. For example, a father presides and exercises the priesthood in his family by the authority of the priesthood he holds. He has no need to have the direction or approval of one holding priesthood keys in order to perform his various family functions. These include counseling the members of his family, holding family meetings, giving priesthood blessings to his wife and children, or giving healing blessings to family members or others. Church authorities teach family members but do not direct the exercise of priesthood authority in the family.

The same principle applies when a father is absent and a mother is the family leader. She presides in her home and is instrumental in bringing the power and blessings of the priesthood into her family through her endowment and sealing in the temple. While she is not authorized to give the priesthood blessings that can be given only by a person holding a certain office in the priesthood, she can perform all of the other functions of family leadership. In doing so, she exercises the power of the priesthood for the benefit of the children over whom she presides in her position of leadership in the family.

If fathers would magnify their priesthood in their own family, it would further the mission of the Church as much as anything else they might do. Fathers who hold the Melchizedek Priesthood should exercise their authority “by persuasion, by long-suffering, by gentleness and meekness, and by love unfeigned” (Doctrine and Covenants 121:41). That high standard for the exercise of all priesthood authority is most important in the family. Holders of the priesthood should also keep the commandments so they will have the power of the priesthood to give blessings to their family members. They should cultivate loving family relationships so that family members will want to ask them for blessings. And parents should encourage more priesthood blessings in the family.

In these conference meetings, as we seek brief shelter from our mortal concerns with a devastating pandemic, we have been taught great principles of eternity. I encourage each of us to have our eye “single” to receive these truths of eternity so that our bodies “shall be full of light” (3 Nephi 13:22).

In His sermon to multitudes recorded in the Bible and in the Book of Mormon, the Savior taught that mortal bodies can be full of light or full of darkness. We, of course, want to be filled with light, and our Savior taught us how we can make this happen. We should listen to messages about the truths of eternity. He used the example of our eye, through which we take light into our bodies. If our “eye be single”—in other words, if we are concentrating on receiving eternal light and understanding—He explained, “thy whole body shall be full of light” (Matthew 6:22; 3 Nephi 13:22). But if our “eye be evil”—that is, if we look for evil and take that into our bodies—He warned, “thy whole body shall be full of darkness” (verse 23). In other words, the light or darkness in our bodies depends on how we see—or receive—the eternal truths we are taught.

We should follow the Savior’s invitation to seek and ask to understand the truths of eternity. He promises that our Father in Heaven is willing to teach everyone the truths they seek (see 3 Nephi 14:8). If we desire this and have our eye single to receive them, the Savior promises that the truths of eternity “shall be opened” unto us (see 3 Nephi 14:7–8).

In contrast, Satan is anxious to confuse our thinking or to lead us astray on important matters like the operations of the priesthood of God. The Savior warned of such “false prophets, who come to you in sheep’s clothing, but inwardly they are ravening wolves” (3 Nephi 14:15). He gave us this test to help us choose the truth from among different teachings that might confuse us: “Ye shall know them by their fruits,” He taught (3 Nephi 14:16). “A good tree cannot bring forth evil fruit, neither [can] a corrupt tree bring forth good fruit” (verse 18). Therefore, we should look to the results—“the fruits”—of principles that are taught and the persons who teach them. That is the best answer to many of the objections we hear against the Church and its doctrines and policies and leadership. Follow the test the Savior taught. Look to the fruits—the results.

When we think of the fruits of the gospel and the restored Church of Jesus Christ, we rejoice in how the Church, in the lifetimes of its living members, has expanded from local congregations in the Intermountain West to where a majority of its more than 16 million members reside in nations other than the United States. With that growth, we have felt increases in the Church’s capacity to assist its members. We assist in keeping the commandments, in fulfilling responsibilities to preach the restored gospel, in gathering Israel, and in building temples throughout the world.

We are led by a prophet, President Russell M. Nelson, whose leadership the Lord has used to achieve the progress we have felt during all of the more than two years of his leadership. Now we will be blessed to hear from President Nelson, who will teach us how to further our progress in this restored Church of Jesus Christ in these challenging times.

I testify of the truth of these things and join you in praying for our prophet, from whom we will next hear, in the name of Jesus Christ, amen.

\newpage
\settalk{The Great Plan}
\setspeaker{Dallin H. Oaks}
\settalkdate{April 2020}
\talkheading{The Great Plan}{Dallin H. Oaks}

Even in the midst of unique trials and challenges, we are truly blessed! This general conference has given us an outpouring of the riches and joy of the Restoration of the gospel of Jesus Christ. We have rejoiced in the vision of the Father and the Son that commenced the Restoration. We have been reminded of the miraculous coming forth of the Book of Mormon, whose central purpose is to testify of Jesus Christ and His doctrine. We have been renewed with the joyful reality of revelation—to prophets and to us personally. We have heard precious testimonies of the infinite Atonement of Jesus Christ and of His literal Resurrection. And we have been taught other truths of the fulness of His gospel revealed to Joseph Smith after God the Father declared to that newly called prophet: “This is My Beloved Son. Hear Him!” (Joseph Smith—History 1:17).

We have been affirmed in our knowledge of the restoration of the priesthood and its keys. We have been renewed in our determination to have the Lord’s restored Church known by its proper name, The Church of Jesus Christ of Latter-day Saints. And we have been invited to join in fasting and prayer to minimize the present and future effects of a devastating worldwide pandemic. This morning we were inspired by the Lord’s living prophet presenting an historic proclamation of the Restoration. We affirm its declaration that “those who prayerfully study the message of the Restoration and act in faith will be blessed to gain their own witness of its divinity and of its purpose to prepare the world for the promised Second Coming of our Lord and Savior, Jesus Christ.”

\intalkheader{The Plan}

All of this is part of a divine plan whose purpose is to enable the children of God to be exalted and become like Him. Referred to in the scriptures as the “great plan of happiness,” “the plan of redemption,” and the “plan of salvation” (Alma 42:8, 11, 5), that plan—revealed in the Restoration—began with a Council in Heaven. As spirits, we desired to achieve the eternal life enjoyed by our heavenly parents. At that point we had progressed as far as we could without a mortal experience in a physical body. To provide that experience, God the Father planned to create this earth. In the planned mortal life, we would be soiled by sin as we faced the opposition necessary for our spiritual growth. We would also become subject to physical death. To reclaim us from death and sin, our Heavenly Father’s plan would provide a Savior. His Resurrection would redeem all from death, and His atoning sacrifice would pay the price necessary for all to be cleansed from sin on the conditions prescribed to promote our growth. This Atonement of Jesus Christ is central to the Father’s plan.

In the Council in Heaven, all the spirit children of God were introduced to the Father’s plan, including its mortal consequences and trials, its heavenly helps, and its glorious destiny. We saw the end from the beginning. All of the myriads of mortals who have been born on this earth chose the Father’s plan and fought for it in the heavenly contest that followed. Many also made covenants with the Father concerning what they would do in mortality. In ways that have not been revealed, our actions in the spirit world have influenced our circumstances in mortality.

\intalkheader{Mortality and Spirit World}

I will now summarize some of the principal elements of the Father’s plan as they affect us during our mortal journeys and in the spirit world that follows.

The purpose of mortal life and the postmortal growth that can follow it is for the offspring of God to become like He is. This is Heavenly Father’s desire for all His children. To achieve this joyful destiny, eternal laws require that we must become purified beings through the Atonement of Jesus Christ so we can dwell in the presence of the Father and the Son and enjoy the blessings of exaltation. As the Book of Mormon teaches, He invites “all to come unto him and partake of his goodness; and he denieth none that come unto him, black and white, bond and free, male and female; and he remembereth the heathen; and all are alike unto God” (2 Nephi 26:33; see also Alma 5:49).

The divine plan for us to become what we are destined to become requires us to make choices to reject the evil opposition that tempts mortals to act contrary to God’s commandments and His plan. It also requires that we be subject to other mortal opposition, such as from the sins of others or from some defects of birth. Sometimes our needed growth is achieved better by suffering and adversity than by comfort and tranquility. And none of this mortal opposition could achieve its eternal purpose if divine intervention relieved us from all the adverse consequences of mortality.

The plan reveals our destiny in eternity, the purpose and conditions of our journey in mortality, and the heavenly helps we will receive. The commandments of God warn us against straying into dangerous circumstances. The teachings of inspired leaders guide our path and give assurances that promote our eternal journey.

God’s plan gives us four great assurances to assist our journey through mortality. All are given to us through the Atonement of Jesus Christ, the centerpiece of the plan. The first assures us that through His suffering for the sins of which we repent, we can be cleansed of those sins. Then the merciful final judge will “remember them no more” (Doctrine and Covenants 58:42).

Second, as part of our Savior’s Atonement, He took upon Him all other mortal infirmities. This allows us to receive divine help and strength to bear the inevitable burdens of mortality, personal and general, such as war and pestilence. The Book of Mormon provides our clearest scriptural description of this essential power of the Atonement. The Savior took upon Him “the pains and the sicknesses [and infirmities] of his people. … He will take upon him their infirmities, that his bowels may be filled with mercy, according to the flesh, that he may know according to the flesh how to succor his people according to their infirmities” (Alma 7:11–12).

Third, the Savior, through His infinite Atonement, revokes the finality of death and gives us the joyful assurance that all of us will be resurrected. The Book of Mormon teaches, “This restoration shall come to all, both old and young, both bond and free, both male and female, both the wicked and the righteous; and even there shall not so much as a hair of their heads be lost; but every thing shall be restored to its perfect frame” (Alma 11:44).

We celebrate the reality of the Resurrection in this Easter season. This gives us the perspective and strength to endure the mortal challenges faced by each of us and those we love, such things as the physical, mental, or emotional deficiencies we acquire at birth or experience during our mortal lives. Because of the Resurrection, we know that these mortal deficiencies are only temporary!

The restored gospel assures us that the Resurrection can include the opportunity to be with our family members—husband, wife, children, and parents. This is a powerful encouragement for us to fulfill our family responsibilities in mortality. It helps us live together in love in this life in anticipation of joyful reunions and associations in the next.

Fourth and finally, modern revelation teaches us that our progress need not conclude with the end of mortality. Little has been revealed about this important assurance. We are told that this life is the time to prepare to meet God and that we should not procrastinate our repentance (see Alma 34:32–33). Still, we are taught that in the spirit world the gospel is preached even to “the wicked and the disobedient who had rejected the truth” (Doctrine and Covenants 138:29) and that those taught there are capable of repentance in advance of the Final Judgment (see verses 31–34, 57–59).

Here are some other fundamentals of our Heavenly Father’s plan:

The restored gospel of Jesus Christ gives us a unique perspective on the subjects of chastity, marriage, and the bearing of children. It teaches that marriage according to God’s plan is necessary for accomplishing the purpose of God’s plan, to provide the divinely appointed setting for mortal birth, and to prepare family members for eternal life. “Marriage is ordained of God unto man,” the Lord said, “… that the earth might answer the end of its creation” (Doctrine and Covenants 49:15–16). In this, His plan, of course, runs counter to some strong worldly forces in law and custom.

The power to create mortal life is the most exalted power God has given His children. Its use was mandated in the first commandment to Adam and Eve, but another important commandment was given to forbid its misuse. Outside the bonds of marriage, all uses of the procreative power are to one degree or another a sinful degrading and perversion of the most divine attribute of men and women. The emphasis the restored gospel places on this law of chastity is because of the purpose of our procreative powers in the accomplishment of God’s plan.

\intalkheader{What Next?}

During this 200th anniversary of the First Vision, which initiated the Restoration, we know the Lord’s plan and we are encouraged by two centuries of its blessings through His restored Church. In this year of 2020, we have what is popularly called 20/20 vision for the events of the past.

As we look to the future, however, our vision is far less sure. We know that two centuries after the Restoration, the spirit world now includes many mortally experienced workers to accomplish the preaching that occurs there. We also know that we now have many more temples to perform the ordinances of eternity for those who repent and embrace the Lord’s gospel on either side of the veil of death. All of this furthers our Heavenly Father’s plan. God’s love is so great that, except for the few who deliberately become sons of perdition, He has provided a destiny of glory for all of His children (see Doctrine and Covenants 76:43).

We know that the Savior will return and that there will be a millennium of peaceful reign to wrap up the mortal part of God’s plan. We also know that there will be different resurrections, of the just and the unjust, with the final judgment of each person always following his or her resurrection.

We will be judged according to our actions, the desires of our hearts, and the kind of person we have become. This judgment will cause all of the children of God to proceed to a kingdom of glory for which their obedience has qualified them and where they will be comfortable. The judge of all this is our Savior, Jesus Christ (see John 5:22; 2 Nephi 9:41). His omniscience gives Him a perfect knowledge of all of our acts and desires, both those unrepented or unchanged and those repented or righteous. Therefore, after His judgment we will all confess “that his judgments are just” (Mosiah 16:1).

In conclusion, I share the conviction that has come to me from many letters and by reviewing many requests to return to the Church after name removal or apostasy. Many of our members do not fully understand this plan of salvation, which answers most questions about the doctrine and inspired policies of the restored Church. We who know God’s plan and who have covenanted to participate have a clear responsibility to teach these truths and do all that we can to further them for others and in our own circumstances in mortality. I testify of Jesus Christ, our Savior and Redeemer, who makes it all possible, in the name of Jesus Christ, amen.

\newpage
\settalk{Love Your Enemies}
\setspeaker{Dallin H. Oaks}
\settalkdate{October 2020}
\talkheading{Love Your Enemies}{Dallin H. Oaks}

The Lord’s teachings are for eternity and for all of God’s children. In this message I will give some examples from the United States, but the principles I teach are applicable everywhere.

We live in a time of anger and hatred in political relationships and policies. We felt it this summer when some went beyond peaceful protests and engaged in destructive behavior. We feel it in some current campaigns for public offices. Unfortunately, some of this has even spilled over into political statements and unkind references in our Church meetings.

In a democratic government we will always have differences over proposed candidates and policies. However, as followers of Christ we must forgo the anger and hatred with which political choices are debated or denounced in many settings.

Here is one of our Savior’s teachings, probably well known but rarely practiced:

“Ye have heard that it hath been said, Thou shalt love thy neighbour, and hate thine enemy.

“But I say unto you, Love your enemies, bless them that curse you, do good to them that hate you, and pray for them which despitefully use you, and persecute you” (Matthew 5:43–44).

For generations, Jews had been taught to hate their enemies, and they were then suffering under the domination and cruelties of Roman occupation. Yet Jesus taught them, “Love your enemies” and “do good to them that … despitefully use you.”

What revolutionary teachings for personal and political relationships! But that is still what our Savior commands. In the Book of Mormon we read, “For verily, verily I say unto you, he that hath the spirit of contention is not of me, but is of the devil, who is the father of contention, and he stirreth up the hearts of men to contend with anger, one with another” (3 Nephi 11:29).

Loving our enemies and our adversaries is not easy. “Most of us have not reached that stage of … love and forgiveness,” President Gordon B. Hinckley observed, adding, “It requires a self-discipline almost greater than we are capable of.” But it must be essential, for it is part of the Savior’s two great commandments to “love the Lord thy God” and to “love thy neighbour as thyself” (Matthew 22:37, 39). And it must be possible, for He also taught, “Ask, and it shall be given you; seek, and ye shall find” (Matthew 7:7).

How do we keep these divine commandments in a world where we are also subject to the laws of man? Fortunately, we have the Savior’s own example of how to balance His eternal laws with the practicalities of man-made laws. When adversaries sought to trap Him with a question about whether Jews should pay taxes to Rome, He pointed to the image of Caesar on their coins and declared, “Render therefore unto Caesar the things which be Caesar’s, and unto God the things which be God’s” (Luke 20:25).

So, we are to follow the laws of men (render unto Caesar) to live peacefully under civil authority, and we follow the laws of God toward our eternal destination. But how do we do this—especially how do we learn to love our adversaries and our enemies?

The Savior’s teaching not to “contend with anger” is a good first step. The devil is the father of contention, and it is he who tempts men to contend with anger. He promotes enmity and hateful relationships among individuals and within groups. President Thomas S. Monson taught that anger is “Satan’s tool,” for “to be angry is to yield to the influence of Satan. No one can make us angry. It is our choice.” Anger is the way to division and enmity. We move toward loving our adversaries when we avoid anger and hostility toward those with whom we disagree. It also helps if we are even willing to learn from them.

Among other ways to develop the power to love others is the simple method described in a long-ago musical. When we are trying to understand and relate to people of a different culture, we should try getting to know them. In countless circumstances, strangers’ suspicion or even hostility give way to friendship or even love when personal contacts produce understanding and mutual respect.

An even greater help in learning to love our adversaries and our enemies is to seek to understand the power of love. Here are three of many prophetic teachings about this.

The Prophet Joseph Smith taught that “it is a time-honored adage that love begets love. Let us pour forth love—show forth our kindness unto all mankind.”

President Howard W. Hunter taught: “The world in which we live would benefit greatly if men and women everywhere would exercise the pure love of Christ, which is kind, meek, and lowly. It is without envy or pride. … It seeks nothing in return. … It has no place for bigotry, hatred, or violence. … It encourages diverse people to live together in Christian love regardless of religious belief, race, nationality, financial standing, education, or culture.”

And President Russell M. Nelson has urged us to “expand our circle of love to embrace the whole human family.”

An essential part of loving our enemies is to render unto Caesar by keeping the laws of our various countries. Though Jesus’s teachings were revolutionary, He did not teach revolution or lawbreaking. He taught a better way. Modern revelation teaches the same:

“Let no man break the laws of the land, for he that keepeth the laws of God hath no need to break the laws of the land.

“Wherefore, be subject to the powers that be” (Doctrine and Covenants 58:21–22).

And our article of faith, written by the Prophet Joseph Smith after the early Saints had suffered severe persecution from Missouri officials, declares, “We believe in being subject to kings, presidents, rulers, and magistrates, in obeying, honoring, and sustaining the law” (Articles of Faith 1:12).

This does not mean that we agree with all that is done with the force of law. It means that we obey the current law and use peaceful means to change it. It also means that we peacefully accept the results of elections. We will not participate in the violence threatened by those disappointed with the outcome. In a democratic society we always have the opportunity and the duty to persist peacefully until the next election.

The Savior’s teaching to love our enemies is based on the reality that all mortals are beloved children of God. That eternal principle and some basic principles of law were tested in the recent protests in many American cities.

At one extreme, some seem to have forgotten that the First Amendment to the United States Constitution guarantees the “right of the people peaceably to assemble, and to petition the Government for a redress of grievances.” That is the authorized way to raise public awareness and to focus on injustices in the content or administration of the laws. And there have been injustices. In public actions and in our personal attitudes, we have had racism and related grievances. In a persuasive personal essay, the Reverend Theresa A. Dear of the National Association for the Advancement of Colored People (NAACP) has reminded us that “racism thrives on hatred, oppression, collusion, passivity, indifference and silence.” As citizens and as members of The Church of Jesus Christ of Latter-day Saints, we must do better to help root out racism.

At the other extreme, a minority of participants and supporters of these protests and the illegal acts that followed them seem to have forgotten that the protests protected by the Constitution are peaceful protests. Protesters have no right to destroy, deface, or steal property or to undermine the government’s legitimate police powers. The Constitution and laws contain no invitation to revolution or anarchy. All of us—police, protesters, supporters, and spectators—should understand the limits of our rights and the importance of our duties to stay within the boundaries of existing law. Abraham Lincoln was right when he said, “There is no grievance that is a fit object of redress by mob law.” Redress of grievances by mobs is redress by illegal means. That is anarchy, a condition that has no effective governance and no formal police, which undermines rather than protects individual rights.

One reason the recent protests in the United States were shocking to so many was that the hostilities and illegalities felt among different ethnicities in other nations should not be felt in the United States. This country should be better in eliminating racism not only against Black Americans, who were most visible in the recent protests, but also against Latinos, Asians, and other groups. This nation’s history of racism is not a happy one, and we must do better.

The United States was founded by immigrants of different nationalities and different ethnicities. Its unifying purpose was not to establish a particular religion or to perpetuate any of the diverse cultures or tribal loyalties of the old countries. Our founding generation sought to be unified by a new constitution and laws. That is not to say that our unifying documents or the then-current understanding of their meanings were perfect. The history of the first two centuries of the United States showed the need for many refinements, such as voting rights for women and, particularly, the abolition of slavery, including laws to ensure that those who had been enslaved would have all the conditions of freedom.

Two Yale University scholars recently reminded us:

“For all its flaws, the United States is uniquely equipped to unite a diverse and divided society. …

“… Its citizens don’t have to choose between a national identity and multiculturalism. Americans can have both. But the key is constitutional patriotism. We have to remain united by and through the Constitution, regardless of our ideological disagreements.”

Many years ago, a British foreign secretary gave this great counsel in a debate in the House of Commons: “We have no eternal allies and we have no perpetual enemies. Our interests are eternal and perpetual, and these interests it is our duty to follow.”

That is a good secular reason for following “eternal and perpetual” interests in political matters. In addition, the doctrine of the Lord’s Church teaches us another eternal interest to guide us: the teachings of our Savior, who inspired the Constitution of the United States and the basic laws of many of our countries. Loyalty to established law instead of temporary “allies” is the best way to love our adversaries and our enemies as we seek unity in diversity.

Knowing that we are all children of God gives us a divine vision of the worth of all others and the will and ability to rise above prejudice and racism. As I have lived for many years in different places in this nation, the Lord has taught me that it is possible to obey and seek to improve our nation’s laws and also to love our adversaries and our enemies. While not easy, it is possible with the help of our Lord, Jesus Christ. He gave this command to love, and He promises His help as we seek to obey it. I testify that we are loved and will be helped by our Heavenly Father and His Son, Jesus Christ. In the name of Jesus Christ, amen.

\newpage
\settalk{Be of Good Cheer}
\setspeaker{Dallin H. Oaks}
\settalkdate{October 2020}
\talkheading{Be of Good Cheer}{Dallin H. Oaks}

In the final days of His mortal life, Jesus Christ told His Apostles of the persecutions and hardships they would suffer. He concluded with this great assurance: “In the world ye shall have tribulation: but be of good cheer; I have overcome the world” (John 16:33). That is the Savior’s message to all of our Heavenly Father’s children. That is the ultimate good news for each of us in our mortal lives.

“Be of good cheer” was also a needed assurance in the world into which the resurrected Christ sent His Apostles. “We are troubled on every side,” the Apostle Paul later told the Corinthians, “yet not distressed; we are perplexed, but not in despair; persecuted, but not forsaken; cast down, but not destroyed” (2 Corinthians 4:8–9).

Two thousand years later we are also “troubled on every side,” and we also need that same message not to despair but to be of good cheer. The Lord has special love and concern for His precious daughters. He knows of your wants, your needs, and your fears. The Lord is all powerful. Trust Him.

The Prophet Joseph Smith was taught that “the works, and the designs, and the purposes of God cannot be frustrated, neither can they come to naught” (Doctrine and Covenants 3:1). To His struggling children, the Lord gave these great assurances:

“Behold, this is the promise of the Lord unto you, O ye my servants.

“Wherefore, be of good cheer, and do not fear, for I the Lord am with you, and will stand by you; and ye shall bear record of me, even Jesus Christ, that I am the Son of the living God” (Doctrine and Covenants 68:5–6).

The Lord stands near us, and He has said:

“What I say unto one I say unto all, be of good cheer, little children; for I am in your midst, and I have not forsaken you” (Doctrine and Covenants 61:36).

“For after much tribulation come the blessings” (Doctrine and Covenants 58:4).

Sisters, I testify that these promises, given in the midst of persecutions and personal tragedies, apply to each of you in your troubling circumstances today. They are precious and remind each of us to be of good cheer and to have joy in the fulness of the gospel as we press forward through the challenges of mortality.

Tribulation and challenges are the common experiences of mortality. Opposition is an essential part of the divine plan for helping us grow, and in the midst of that process, we have God’s assurance that, in the long view of eternity, opposition will not be allowed to overcome us. With His help and our faithfulness and endurance, we will prevail. Like the mortal life of which they are a part, all tribulations are temporary. In the controversies that preceded a disastrous war, United States president Abraham Lincoln wisely reminded his audience of the ancient wisdom that “this, too, shall pass away.”

As you know, the mortal adversities of which I speak—which make it difficult to be of good cheer—sometimes come to us in common with many others, like the millions now struggling through some of the many devastating effects of the COVID-19 pandemic. Similarly, in the United States millions are suffering through a season of enmity and contention that always seems to accompany presidential elections but this time is the most severe many of the oldest of us can ever remember.

On a personal basis, each of us struggles individually with some of the many adversities of mortality, such as poverty, racism, ill health, job losses or disappointments, wayward children, bad marriages or no marriages, and the effects of sin—our own or others’.

Yet, in the midst of all of this, we have that heavenly counsel to be of good cheer and to find joy in the principles and promises of the gospel and the fruits of our labors. That counsel has always been so, for prophets and for all of us. We know this from the experiences of our predecessors and what the Lord said to them.

Remember the circumstances of the Prophet Joseph Smith. Looked at through the lens of adversities, his life was one of poverty, persecution, frustration, family sorrows, and ultimate martyrdom. As he suffered imprisonment, his wife and children and the other Saints suffered incredible hardships as they were driven out of Missouri.

When Joseph pleaded for relief, the Lord answered:

“My son, peace be unto thy soul; thine adversity and thine afflictions shall be but a small moment;

“And then, if thou endure it well, God shall exalt thee on high; thou shalt triumph over all thy foes” (Doctrine and Covenants 121:7–8).

This was the personal, eternal counsel that helped the Prophet Joseph to maintain his native cheery temperament and the love and loyalty of his people. These same qualities strengthened the leaders and pioneers who followed and can strengthen you as well.

Think of those early members! Again and again, they were driven from place to place. Finally they faced the challenges of establishing their homes and the Church in a wilderness. Two years after the initial band of pioneers arrived in the valley of the Great Salt Lake, the pioneers’ grip on survival in that hostile area was still precarious. Most members were still on the trail across the plains or struggling to get resources to do so. Yet leaders and members were still of hope and good cheer.

Even though the Saints were not settled in their new homes, at October 1849 general conference a new wave of missionaries was sent out to Scandinavia, France, Germany, Italy, and the South Pacific. At what could have been thought their lowest level, the pioneers rose to new heights. And just three years later, another 98 were also called to begin to gather scattered Israel. One of the Church leaders explained that these missions “are generally, not to be very long ones; probably from 3 to 7 years will be as long as any man will be absent from his family.”

Sisters, the First Presidency is concerned about your challenges. We love you and pray for you. At the same time, we often give thanks that our physical challenges—apart from earthquakes, fires, floods, and hurricanes—are usually less than our predecessors faced.

In the midst of hardships, the divine assurance is always “be of good cheer, for I will lead you along. The kingdom is yours and the blessings thereof are yours, and the riches of eternity are yours” (Doctrine and Covenants 78:18). How does this happen? How did it happen for the pioneers? How will it happen to women of God today? By our following prophetic guidance, “the gates of hell shall not prevail against [us],” the Lord said by revelation in April 1830. “Yea,” He said, “… the Lord God will disperse the powers of darkness from before you, and cause the heavens to shake for your good, and his name’s glory” (Doctrine and Covenants 21:6). “Fear not, little flock; do good; let earth and hell combine against you, for if ye are built upon my rock, they cannot prevail” (Doctrine and Covenants 6:34).

With the Lord’s promises, we “lift up [our] heart[s] and rejoice” (Doctrine and Covenants 25:13), and “with a glad heart and a cheerful countenance” (Doctrine and Covenants 59:15), we go forward on the covenant path. Most of us do not face decisions of giant proportions, like leaving our homes to pioneer an unknown land. Our decisions are mostly in the daily routines of life, but as the Lord has told us, “Be not weary in well-doing, for ye are laying the foundation of a great work. And out of small things proceedeth that which is great” (Doctrine and Covenants 64:33).

There is boundless power in the doctrine of the restored gospel of Jesus Christ. Our unshakable faith in that doctrine guides our steps and gives us joy. It enlightens our minds and gives strength and confidence to our actions. This guidance and enlightenment and power are promised gifts we have received from our Heavenly Father. By understanding and conforming our lives to that doctrine, including the divine gift of repentance, we can be of good cheer as we keep ourselves on the path toward our eternal destiny—reunion and exaltation with our loving heavenly parents.

“You may be facing overwhelming challenges,” Elder Richard G. Scott taught. “Sometimes they are so concentrated, so unrelenting, that you may feel they are beyond your capacity to control. Don’t face the world alone. ‘Trust in the Lord with all thine heart; and lean not unto thine own understanding’ [Proverbs 3:5]. … It was intended that life be a challenge, not so that you would fail, but that you might succeed through overcoming.”

It is all part of the plan of God the Father and His Son, Jesus Christ, of which I testify, as I pray that we will all persist to our heavenly destination, in the name of Jesus Christ, amen.

\newpage
\settalk{Sustaining of General Authorities, Area Seventies, and General Officers}
\setspeaker{Presented by President Dallin H. Oaks}
\settalkdate{April 2021}
\talkheading{Sustaining of General Authorities, Area Seventies, and General Officers}{Presented by President Dallin H. Oaks}

Brothers and sisters, I will now present the General Authorities, Area Seventies, and General Officers of the Church for your sustaining vote.

Please express your vote in the usual way wherever you may be. If there are those who oppose any of the proposals, we ask that you contact your stake president.

It is proposed that we sustain Russell Marion Nelson as prophet, seer, and revelator and President of The Church of Jesus Christ of Latter-day Saints; Dallin Harris Oaks as First Counselor in the First Presidency; and Henry Bennion Eyring as Second Counselor in the First Presidency.

Those in favor may manifest it.

Those opposed, if any, may manifest it.

It is proposed that we sustain Dallin H. Oaks as President of the Quorum of the Twelve Apostles and M. Russell Ballard as Acting President of the Quorum of the Twelve Apostles.

Those in favor, please signify.

Any opposed may manifest it.

It is proposed that we sustain the following as members of the Quorum of the Twelve Apostles: M. Russell Ballard, Jeffrey R. Holland, Dieter F. Uchtdorf, David A. Bednar, Quentin L. Cook, D. Todd Christofferson, Neil L. Andersen, Ronald A. Rasband, Gary E. Stevenson, Dale G. Renlund, Gerrit W. Gong, and Ulisses Soares.

Those in favor, please manifest it.

Any opposed may so indicate.

It is proposed that we sustain the counselors in the First Presidency and the Quorum of the Twelve Apostles as prophets, seers, and revelators.

All in favor, please manifest it.

Contrary, if there be any, by the same sign.

As a matter of information, Elders Robert C. Gay and Terence M. Vinson will be released from their service as members of the Presidency of the Seventy, effective on August 1, 2021.

Those who wish to express gratitude to these Brethren for their dedicated service may do so by the uplifted hand.

The following Area Seventies have been released: Elders Sean Douglas, Michael A. Dunn, Clark G. Gilbert, Alfred Kyungu, Carlos G. Revillo Jr., and Vaiangina Sikahema.

Those who wish to join us in expressing appreciation for their excellent service, please manifest it.

We have released the Primary General Presidency as follows: Joy D. Jones as President, Lisa L. Harkness as First Counselor, and Cristina B. Franco as Second Counselor.

All who wish to join us in expressing appreciation to these sisters for their devoted service, please manifest it.

It is proposed that we sustain Elders Paul V. Johnson and S. Mark Palmer to serve as members of the Presidency of the Seventy, effective on August 1, 2021.

Those in favor may manifest it.

Any opposed, by the same sign.

It is proposed that we sustain the following as General Authority Seventies: Sean Douglas, Michael A. Dunn, Clark G. Gilbert, Patricio M. Giuffra, Alfred Kyungu, Alvin F. Meredith III, Carlos G. Revillo Jr., and Vaiangina Sikahema.

All in favor, please manifest it.

Those opposed, by the same sign.

It is proposed that we sustain the new Area Seventies as announced by the Church earlier this week.

Those in favor, please manifest it.

Any opposed, by the same sign.

It is proposed that we sustain as the new Primary General Presidency Camille N. Johnson as President, Susan H. Porter as First Counselor, and Amy A. Wright as Second Counselor.

Those in favor, manifest it.

Any opposed may so signify.

It is proposed that we sustain the other General Authorities, Area Seventies, and General Officers as presently constituted.

All in favor, please manifest it.

Those opposed, if any.

Again, those who oppose any of the proposals are invited to contact their stake presidents.

Thank you for your continued faith and prayers in behalf of the leaders of the Church.

\intalkheader{Changes to Area Seventies}

The following Area Seventies were sustained during a leadership session held as part of general conference:

Daniel P. Amato, Rodney A. Ames, Marcelo Andrezzo, Samuel Annan-Simons, Patrick Appianti-Sarpong, Eduardo M. Argana, Steven C. Barlow, Erik Bernskov, Mark E. Bonham, K. Bruce Boucher, Jonathan G. Cannon, Juan P. Casco, Gregorio E. Casillas, Fernando R. Castro, Ranulfo Cervantes, Thomas K. Checketts, Ross A. Chiles, Benjamin Cinco, David C. Clark, Félix Conde, Jorge A. Contreras, Corbin E. Coombs, Moroni Costa, Leandro J. Curaba, B. Corey Cuvelier, Ernesto A. Deyro Jr., J. Kimo Esplin, Tomás Familia, Michael D. Groll, John Gutty, Oleksiy H. Hakalenko, Tommy D. Haws, Levi W. Heath, Brian J. Holmes, Hal C. Hunsaker, Yuichi Imai, Bruce H. Ixcot, Paul H. Jean Baptiste, Dong Hwan Jeong, Frederick M. Kamya, Gaëtan Kelounou, David S. Kinard, Julio E. Lee, R. Darío Lorenzana, Odair José Castro de Lira, Enrique M. Loo, Hernán D. Lucero, Bartolome Madriaga, Douglas P. Maxfield, Héctor Méndez, Steven C. Merrell, Quinn S. Millington, Siegfried A. Naumann, Ricardo J. Nieves, Lorenzo E. Norambuena, Enefiok Ntem, Charles O. Oide, Juan L. Orquera, Roberto C. Pacheco, Damon Page, Franck A. Poznanski, T. Michael Price, Alexandre Ret, Frédéric T. Riemer, Russell A. Robinson, Leonardo S. Rojas, Douglas A. Rozsa, Lee M. Shumway, Robert H. Simpson, Vance K. Smith, Martiniano S. Soquila Jr., Victor H. Suazo, Raul Tapia, Carlos Torres, Bruno E. Vásquez, M. Travis Wolsey, and Richard G. Youngblood.

\newpage
\settalk{What Has Our Savior Done for Us?}
\setspeaker{Dallin H. Oaks}
\settalkdate{April 2021}
\talkheading{What Has Our Savior Done for Us?}{Dallin H. Oaks}

In a Saturday evening meeting at a stake conference many years ago, I met a woman who said her friends had asked her to come back to church after many years of inactivity, but she could not think of any reason why she should. To encourage her, I said, “When you consider all of the things the Savior has done for you, you have many reasons to come back to worship and serve Him.” I was astonished when she replied, “What’s He done for me?”

What has Jesus Christ done for each of us? He has done everything that is essential for our journey through mortality toward the destiny outlined in the plan of our Heavenly Father. I will speak of four of the principal features of that plan. In each of these, His Only Begotten Son, Jesus Christ, is the central figure. Motivating all of this is “the love of God, which sheddeth itself abroad in the hearts of the children of men; wherefore, it is the most desirable above all things” (1 Nephi 11:22).

\intalkheader{I.}

Just before Easter Sunday, it is timely to speak first of the Resurrection of Jesus Christ. The Resurrection from the dead is the reassuring personal pillar of our faith. It adds meaning to our doctrine, motivation to our behavior, and hope for our future.

Because we believe the Bible and Book of Mormon descriptions of the literal Resurrection of Jesus Christ, we also accept the numerous scriptural teachings that a similar resurrection will come to all mortals who have ever lived upon this earth. As Jesus taught, “Because I live, ye shall live also” (John 14:19). And His Apostle taught that “the dead shall be raised incorruptible” and “this mortal shall have put on immortality” (1 Corinthians 15:52, 54).

But the Resurrection gives us more than this assurance of immortality. It changes the way we view mortal life.

The Resurrection gives us the perspective and the strength to endure the mortal challenges faced by each of us and those we love. It gives us a new way to view the physical, mental, or emotional deficiencies we have at birth or acquire during mortal life. It gives us the strength to endure sorrows, failures, and frustrations. Because each of us has an assured resurrection, we know that these mortal deficiencies and oppositions are only temporary.

The Resurrection also gives us a powerful incentive to keep the commandments of God during our mortal lives. When we rise from the dead and proceed to our prophesied Final Judgment, we want to have qualified for the choicest blessings promised to resurrected beings.

In addition, the promise that the Resurrection can include an opportunity to be with our family members—husband, wife, children, parents, and posterity—is a powerful encouragement to fulfill our family responsibilities in mortality. It also helps us live together in love in this life, and it comforts us in the death of our loved ones. We know that these mortal separations are only temporary, and we anticipate future joyful reunions and associations. The Resurrection provides us hope and the strength to be patient as we wait. It also prepares us with the courage and dignity to face our own death—even a death that might be called premature.

All of these effects of the Resurrection are part of the first answer to the question “What has Jesus Christ done for me?”

\intalkheader{II.}

For most of us, the opportunity to be forgiven of our sins is the major meaning of the Atonement of Jesus Christ. In worship, we reverently sing:

\begin{verse}
His precious blood he freely spilt;\\
His life he freely gave,\\
A sinless sacrifice for guilt,\\
A dying world to save.
\end{verse}

Our Savior and Redeemer endured incomprehensible suffering to become a sacrifice for the sins of all mortals who would repent. This atoning sacrifice offered the ultimate good, the pure lamb without blemish, for the ultimate measure of evil, the sins of the entire world. It opened the door for each of us to be cleansed of our personal sins so we can be readmitted to the presence of God, our Eternal Father. This open door is available to all of the children of God. In worship, we sing:

\begin{verse}
I marvel that he would descend from his throne divine\\
To rescue a soul so rebellious and proud as mine,\\
That he should extend his great love unto such as I.
\end{verse}

The magnificent and incomprehensible effect of the Atonement of Jesus Christ is based on God’s love for each of us. It affirms His declaration that “the worth of souls”—every one—“is great in the sight of God” (Doctrine and Covenants 18:10). In the Bible, Jesus Christ explained this in terms of our Heavenly Father’s love: “For God so loved the world, that he gave his only begotten Son, that whosoever believeth in him should not perish, but have everlasting life” (John 3:16). In modern revelation, our Redeemer, Jesus Christ, declared that He “so loved the world that he gave his own life, that as many as would believe might become the sons of God” (Doctrine and Covenants 34:3).

Is it any wonder, then, that the Book of Mormon: Another Testament of Jesus Christ concludes with the teaching that to become “perfect” and “sanctified in Christ,” we must “love God with all [our] might, mind and strength”? (Moroni 10:32–33). His plan motivated by love must be received with love.

\intalkheader{III.}

What else has our Savior, Jesus Christ, done for us? Through the teachings of His prophets and through His personal ministry, Jesus taught us the plan of salvation. This plan includes the Creation, the purpose of life, the necessity of opposition, and the gift of agency. He also taught us the commandments and covenants we must obey and the ordinances we must experience to take us back to our heavenly parents.

In the Bible, we read His teaching: “I am the light of the world: he that followeth me shall not walk in darkness, but shall have the light of life” (John 8:12). And in modern revelation, we read, “Behold, I am Jesus Christ, … a light which cannot be hid in darkness” (Doctrine and Covenants 14:9). If we follow His teachings, He lights our path in this life and assures our destiny in the next.

Because He loves us, He challenges us to focus on Him instead of the things of this mortal world. In His great sermon on the bread of life, Jesus taught that we should not be among those who are most attracted to the things of the world—the things that sustain life on earth but give no nourishment toward eternal life. As Jesus invited us again and again and again, “Follow me.”

\intalkheader{IV.}

Finally, the Book of Mormon teaches that as part of His Atonement, Jesus Christ “suffer[ed] pains and afflictions and temptations of every kind; and this that the word might be fulfilled which saith he will take upon him the pains and the sicknesses of his people” (Alma 7:11).

Why did our Savior suffer these mortal challenges “of every kind”? Alma explained, “And he will take upon him their infirmities, that his bowels may be filled with mercy, according to the flesh, that he may know according to the flesh how to succor [which means to give relief or aid to] his people according to their infirmities” (Alma 7:12).

Our Savior feels and knows our temptations, our struggles, our heartaches, and our sufferings, for He willingly experienced them all as part of His Atonement. Other scriptures affirm this. The New Testament declares, “In that he himself hath suffered being tempted, he is able to succour them that are tempted” (Hebrews 2:18). Isaiah teaches, “Fear thou not; for I am with thee: … I will strengthen thee; yea, I will help thee” (Isaiah 41:10). All who suffer any kind of mortal infirmities should remember that our Savior experienced that kind of pain also, and that through His Atonement, He offers each of us the strength to bear it.

The Prophet Joseph Smith summarized all of this in our third article of faith: “We believe that through the Atonement of Christ, all mankind may be saved, by obedience to the laws and ordinances of the Gospel.”

“What has Jesus Christ done for me?” that sister asked. Under the plan of our Heavenly Father, He “created the heavens and the earth” (Doctrine and Covenants 14:9) so that each of us could have the mortal experience necessary to seek our divine destiny. As part of the Father’s plan, the Resurrection of Jesus Christ overcame death to assure each of us immortality. Jesus Christ’s atoning sacrifice gives each of us the opportunity to repent of our sins and return clean to our heavenly home. His commandments and covenants show us the way, and His priesthood gives the authority to perform the ordinances that are essential to reach that destiny. And our Savior willingly experienced all mortal pains and infirmities that He would know how to strengthen us in our afflictions.

Jesus Christ did all of this because He loves all of the children of God. Love is the motivation for it all, and it was so from the very beginning. God has told us in modern revelation that “he created … male and female, after his own image … ; and gave unto them commandments that they should love and serve him” (Doctrine and Covenants 20:18–19).

I testify of all of this and pray that we all will remember what our Savior has done for each of us and that we all will love and serve Him, in the name of Jesus Christ, amen.

\newpage
\settalk{Defending Our Divinely Inspired Constitution}
\setspeaker{Dallin H. Oaks}
\settalkdate{April 2021}
\talkheading{Defending Our Divinely Inspired Constitution}{Dallin H. Oaks}

In this troubled time, I have felt to speak about the inspired Constitution of the United States. This Constitution is of special importance to our members in the United States, but it is also a common heritage of constitutions around the world.

\intalkheader{I.}

A constitution is the foundation of government. It provides structure and limits for the exercise of government powers. The United States Constitution is the oldest written constitution still in force today. Though originally adopted by only a small number of colonies, it soon became a model worldwide. Today, every nation except three have adopted written constitutions.

In these remarks I do not speak for any political party or other group. I speak for the United States Constitution, which I have studied for more than 60 years. I speak from my experience as a law clerk to the chief justice of the United States Supreme Court. I speak from my 15 years as a professor of law and my 3½ years as a justice on the Utah Supreme Court. Most important, I speak from 37 years as an Apostle of Jesus Christ, responsible to study the meaning of the divinely inspired United States Constitution to the work of His restored Church.

The United States Constitution is unique because God revealed that He “established” it “for the rights and protection of all flesh” (Doctrine and Covenants 101:77; see also verse 80). That is why this constitution is of special concern for The Church of Jesus Christ of Latter-day Saints throughout the world. Whether or how its principles should be applied in other nations of the world is for them to decide.

What was God’s purpose in establishing the United States Constitution? We see it in the doctrine of moral agency. In the first decade of the restored Church, its members on the western frontier were suffering private and public persecution. Partly this was because of their opposition to the human slavery then existing in the United States. In these unfortunate circumstances, God revealed through the Prophet Joseph Smith eternal truths about His doctrine.

God has given His children moral agency—the power to decide and to act. The most desirable condition for the exercise of that agency is maximum freedom for men and women to act according to their individual choices. Then, the revelation explains, “every man may be accountable for his own sins in the day of judgment” (Doctrine and Covenants 101:78). “Therefore,” the Lord revealed, “it is not right that any man should be in bondage one to another” (Doctrine and Covenants 101:79). This obviously means that human slavery is wrong. And according to the same principle, it is wrong for citizens to have no voice in the selection of their rulers or the making of their laws.

\intalkheader{II.}

Our belief that the United States Constitution was divinely inspired does not mean that divine revelation dictated every word and phrase, such as the provisions allocating the number of representatives from each state or the minimum age of each. The Constitution was not “a fully grown document,” said President J. Reuben Clark. “On the contrary,” he explained, “we believe it must grow and develop to meet the changing needs of an advancing world.” For example, inspired amendments abolished slavery and gave women the right to vote. However, we do not see inspiration in every Supreme Court decision interpreting the Constitution.

I believe the United States Constitution contains at least five divinely inspired principles.

First is the principle that the source of government power is the people. In a time when sovereign power was universally assumed to come from the divine right of kings or from military power, attributing sovereign power to the people was revolutionary. Philosophers had advocated this, but the United States Constitution was the first to apply it. Sovereign power in the people does not mean that mobs or other groups of people can intervene to intimidate or force government action. The Constitution established a constitutional democratic republic, where the people exercise their power through their elected representatives.

A second inspired principle is the division of delegated power between the nation and its subsidiary states. In our federal system, this unprecedented principle has sometimes been altered by inspired amendments, such as those abolishing slavery and extending voting rights to women, mentioned earlier. Significantly, the United States Constitution limits the national government to the exercise of powers granted expressly or by implication, and it reserves all other government powers “to the States respectively, or to the people.”

Another inspired principle is the separation of powers. Well over a century before our 1787 Constitutional Convention, the English Parliament pioneered the separation of legislative and executive authority when they wrested certain powers from the king. The inspiration in the American convention was to delegate independent executive, legislative, and judicial powers so these three branches could exercise checks upon one another.

A fourth inspired principle is in the cluster of vital guarantees of individual rights and specific limits on government authority in the Bill of Rights, adopted by amendment just three years after the Constitution went into force. A Bill of Rights was not new. Here the inspiration was in the practical implementation of principles pioneered in England, beginning with the Magna Carta. The writers of the Constitution were familiar with these because some of the colonial charters had such guarantees.

Without a Bill of Rights, America could not have served as the host nation for the Restoration of the gospel, which began just three decades later. There was divine inspiration in the original provision that there should be no religious test for public office, but the addition of the religious freedom and antiestablishment guarantees in the First Amendment was vital. We also see divine inspiration in the First Amendment’s freedoms of speech and press and in the personal protections in other amendments, such as for criminal prosecutions.

Fifth and finally, I see divine inspiration in the vital purpose of the entire Constitution. We are to be governed by law and not by individuals, and our loyalty is to the Constitution and its principles and processes, not to any office holder. In this way, all persons are to be equal before the law. These principles block the autocratic ambitions that have corrupted democracy in some countries. They also mean that none of the three branches of government should be dominant over the others or prevent the others from performing their proper constitutional functions to check one another.

\intalkheader{III.}

Despite the divinely inspired principles of the United States Constitution, when exercised by imperfect mortals their intended effects have not always been achieved. Important subjects of lawmaking, such as some laws governing family relationships, have been taken from the states by the federal government. The First Amendment guarantee of free speech has sometimes been diluted by suppression of unpopular speech. The principle of separation of powers has always been under pressure with the ebb and flow of one branch of government exercising or inhibiting the powers delegated to another.

There are other threats that undermine the inspired principles of the United States Constitution. The stature of the Constitution is diminished by efforts to substitute current societal trends as the reason for its founding, instead of liberty and self-government. The authority of the Constitution is trivialized when candidates or officials ignore its principles. The dignity and force of the Constitution is reduced by those who refer to it like a loyalty test or a political slogan, instead of its lofty status as a source of authorization for and limits on government authority.

\intalkheader{IV.}

Our belief in divine inspiration gives Latter-day Saints a unique responsibility to uphold and defend the United States Constitution and principles of constitutionalism wherever we live. We should trust in the Lord and be positive about this nation’s future.

What else are faithful Latter-day Saints to do? We must pray for the Lord to guide and bless all nations and their leaders. This is part of our article of faith. Being subject to presidents or rulers of course poses no obstacle to our opposing individual laws or policies. It does require that we exercise our influence civilly and peacefully within the framework of our constitutions and applicable laws. On contested issues, we should seek to moderate and unify.

There are other duties that are part of upholding the inspired Constitution. We should learn and advocate the inspired principles of the Constitution. We should seek out and support wise and good persons who will support those principles in their public actions. We should be knowledgeable citizens who are active in making our influence felt in civic affairs.

In the United States and in other democracies, political influence is exercised by running for office (which we encourage), by voting, by financial support, by membership and service in political parties, and by ongoing communications to officials, parties, and candidates. To function well, a democracy needs all of these, but a conscientious citizen does not need to provide all of them.

There are many political issues, and no party, platform, or individual candidate can satisfy all personal preferences. Each citizen must therefore decide which issues are most important to him or her at any particular time. Then members should seek inspiration on how to exercise their influence according to their individual priorities. This process will not be easy. It may require changing party support or candidate choices, even from election to election.

Such independent actions will sometimes require voters to support candidates or political parties or platforms whose other positions they cannot approve. That is one reason we encourage our members to refrain from judging one another in political matters. We should never assert that a faithful Latter-day Saint cannot belong to a particular party or vote for a particular candidate. We teach correct principles and leave our members to choose how to prioritize and apply those principles on the issues presented from time to time. We also insist, and we ask our local leaders to insist, that political choices and affiliations not be the subject of teachings or advocacy in any of our Church meetings.

The Church of Jesus Christ of Latter-day Saints will, of course, exercise its right to endorse or oppose specific legislative proposals that we believe will impact the free exercise of religion or the essential interests of Church organizations.

I testify of the divinely inspired Constitution of the United States and pray that we who recognize the Divine Being who inspired it will always uphold and defend its great principles. In the name of Jesus Christ, amen.

\newpage
\settalk{The Need for a Church}
\setspeaker{Dallin H. Oaks}
\settalkdate{October 2021}
\talkheading{The Need for a Church}{Dallin H. Oaks}

Many years ago, Elder Mark E. Petersen, a member of the Quorum of the Twelve Apostles, began a talk with this example:

“Kenneth and his wife, Lucille, are good people, honest and upright. They don’t go to church, though, and they feel they can be good enough without it. They teach their children honesty and virtue and they tell themselves that is about all the Church would do for them.

“And, anyway, they insist that they need their weekends for family recreation … [and] church-going would really get in their way.”

Today, my message concerns such good and religious-minded people who have stopped attending or participating in their churches. When I say “churches,” I include synagogues, mosques, or other religious organizations. We are concerned that attendance in all of these is down significantly, nationwide. If we cease valuing our churches for any reason, we threaten our personal spiritual life, and significant numbers separating themselves from God reduce His blessings to our nations.

Attendance and activity in a church help us become better people and better influences on the lives of others. In church we are taught how to apply religious principles. We learn from one another. A persuasive example is more powerful than a sermon. We are strengthened by associating with others of like minds. In church attendance and participation, our hearts are, as the Bible says, “knit together in love.”

\intalkheader{I.}

The scriptures God has given Christians in the Bible and in modern revelation clearly teach the need for a church. Both show that Jesus Christ organized a church and contemplated that a church would carry on His work after Him. He called Twelve Apostles and gave them authority and keys to direct it. The Bible teaches that Christ is “the head of the church” and that its officers were given “for the perfecting of the saints, for the work of the ministry, for the edifying of the body of Christ.” Surely the Bible is clear on the origin of a church and the need for it now.

Some say that attending church meetings is not helping them. Some say, “I didn’t learn anything today” or “No one was friendly to me” or “I was offended.” Personal disappointments should never keep us from the doctrine of Christ, who taught us to serve, not to be served. With this in mind, another member described the focus of his Church attendance:

“Years ago, I changed my attitude about going to church. No longer do I go to church for my sake, but to think of others. I make a point of saying hello to people who sit alone, to welcome visitors, … to volunteer for an assignment. …

“In short, I go to church each week with the intent of being active, not passive, and making a positive difference in people’s lives.”

President Spencer W. Kimball taught that “we do not go to Sabbath meetings to be entertained or even solely to be instructed. We go to worship the Lord. It is an individual responsibility. … If the service is a failure to you, you have failed. No one can worship for you; you must do your own waiting upon the Lord.”

Church attendance can open our hearts and sanctify our souls.

In a church we don’t just serve alone or by our own choice or at our convenience. We usually serve in a team. In service we find heaven-sent opportunities to rise above the individualism of our age. Church-directed service helps us overcome the personal selfishness that can retard our spiritual growth.

There are other important advantages to mention, even briefly. In church we associate with wonderful people striving to serve God. This reminds us that we are not alone in our religious activities. We all need associations with others, and church associations are some of the best we can experience, for us and our companions and children. Without those associations, especially between children and faithful parents, research shows increasing difficulty for parents to raise children in their faith.

\intalkheader{II.}

So far, I have spoken about churches generally. Now I address the special reasons for membership, attendance, and participation in the Savior’s restored Church of Jesus Christ of Latter-day Saints.

We, of course, affirm that the scriptures, ancient and modern, clearly teach the origin and need for a church directed by and with the authority of our Lord, Jesus Christ. We also testify that the restored Church of Jesus Christ has been established to teach the fulness of His doctrine and to officiate with His priesthood authority to perform the ordinances necessary to enter the kingdom of God. Members who forgo Church attendance and rely only on individual spirituality separate themselves from these gospel essentials: the power and blessings of the priesthood, the fulness of restored doctrine, and the motivations and opportunities to apply that doctrine. They forfeit their opportunity to qualify to perpetuate their family for eternity.

Another great advantage of the restored Church is that it helps us grow spiritually. Growth means change. In spiritual terms this means repenting and seeking to draw nearer to the Lord. In the restored Church we have doctrine, procedures, and inspired helpers that assist us to repent. Their purpose, even in membership councils, is not punishment, like the outcome of a criminal court. Church membership councils lovingly seek to help us qualify for the mercy of forgiveness made possible through the Atonement of Jesus Christ.

Individual spirituality can seldom provide the motivation and structure for unselfish service provided by the restored Church. Great examples of this are the young men and women and seniors who put aside their schooling or retirement activities to accept missionary callings. They work as missionaries to strangers in unfamiliar places they have not chosen. The same is true of faithful members who participate in the unselfish service we call “temple work.” None of such service would be possible without the Church that sponsors it, organizes it, and directs it.

Our members’ religious faith and Church service have taught them how to work in cooperative efforts to benefit the larger community. That kind of experience and development does not happen in the individualism so prevalent in the practices of our current society. In the geographic organization of our local wards, we associate and work with persons we might not otherwise have chosen, persons who teach us and test us.

In addition to helping us learn spiritual qualities like love, compassion, forgiveness, and patience, this gives us the opportunities to learn how to work with persons of very different backgrounds and preferences. This advantage has helped many of our members, and many organizations are blessed by their participation. Latter-day Saints are renowned for their ability to lead and unite in cooperative efforts. That tradition originated with our courageous pioneers who colonized the Intermountain West and established our valued tradition of unselfish cooperation for the common good.

Most humanitarian and charitable efforts need to be accomplished by pooling and managing individual resources on a large scale. The restored Church does this with its enormous humanitarian efforts worldwide. These include distributing educational and medical supplies, feeding the hungry, caring for refugees, helping to reverse the effects of addictions, and a host of others. Our Church members are renowned for their Helping Hands projects in natural disasters. Church membership allows us to be part of such large-scale efforts. Members also pay fast offerings to help the poor in their own midst.

In addition to feeling peace and joy through the companionship of the Spirit, our Church-attending members enjoy the fruits of gospel living, such as the blessings of living the Word of Wisdom and the material and spiritual prosperity promised for living the law of tithing. We also have the blessing of counsel from inspired leaders.

Crowning all of this are the authoritative priesthood ordinances necessary for eternity, including the sacrament we receive each Sabbath day. The culminating ordinance in the restored Church is the everlasting covenant of marriage, which makes possible the perpetuation of glorious family relationships. President Russell M. Nelson taught this principle in a memorable way. He said: “We cannot wish our way into the presence of God. We are to obey the laws upon which [that blessing is] predicated.”

One of those laws is to worship in church each Sabbath day. Our worship and application of eternal principles draw us closer to God and magnify our capacity to love. Parley P. Pratt, one of the original Apostles of this dispensation, described how he felt when the Prophet Joseph Smith explained these principles: “I felt that God was my heavenly Father indeed; that Jesus was my brother, and that the wife of my bosom was an immortal, eternal companion: a kind, ministering angel, given to me as a comfort, and a crown of glory for ever and ever. In short, I could now love with the spirit and with the understanding also.”

In closing, I remind all that we do not believe that good can be accomplished only through a church. Independent of a church, we see millions of people supporting and carrying out innumerable good works. Individually, Latter-day Saints participate in many of them. We see these works as a manifestation of the eternal truth that “the Spirit giveth light to every man that cometh into the world.”

Despite the good works that can be accomplished without a church, the fulness of doctrine and its saving and exalting ordinances are available only in the restored Church. In addition, Church attendance gives us the strength and enhancement of faith that come from associating with other believers and worshipping together with those who are also striving to stay on the covenant path and be better disciples of Christ. I pray that we will all be steadfast in these Church experiences as we seek eternal life, the greatest of all the gifts of God, in the name of Jesus Christ, amen.

\newpage
\settalk{Sustaining of General Authorities, Area Seventies, and General Officers}
\setspeaker{Presented by President Dallin H. Oaks}
\settalkdate{April 2022}
\talkheading{Sustaining of General Authorities, Area Seventies, and General Officers}{Presented by President Dallin H. Oaks}

Brothers and sisters, I will now present the General Authorities, Area Seventies, and General Officers of the Church to be sustained.

Please express your support in the usual way wherever you may be. If there are those who oppose any of the proposals, we ask that you contact your stake president.

It is proposed that we sustain Russell Marion Nelson as prophet, seer, and revelator and President of The Church of Jesus Christ of Latter-day Saints; Dallin Harris Oaks as First Counselor in the First Presidency; and Henry Bennion Eyring as Second Counselor in the First Presidency.

Those in favor may manifest it.

Those opposed, if any, may manifest it.

It is proposed that we sustain Dallin H. Oaks as President of the Quorum of the Twelve Apostles and M. Russell Ballard as Acting President of the Quorum of the Twelve Apostles.

Those in favor, please signify.

Any opposed may manifest it.

It is proposed that we sustain the following as members of the Quorum of the Twelve Apostles: M. Russell Ballard, Jeffrey R. Holland, Dieter F. Uchtdorf, David A. Bednar, Quentin L. Cook, D. Todd Christofferson, Neil L. Andersen, Ronald A. Rasband, Gary E. Stevenson, Dale G. Renlund, Gerrit W. Gong, and Ulisses Soares.

Those in favor, please manifest it.

Any opposed may so indicate.

It is proposed that we sustain the counselors in the First Presidency and the Quorum of the Twelve Apostles as prophets, seers, and revelators.

All in favor, please manifest it.

Contrary, if there be any, by the same sign.

The following Area Seventies have been released from their assignments: Mark D. Eddy, Ryan K. Olsen, Jonathan S. Schmitt, and Denelson Silva.

Those who wish to join us in expressing appreciation for their excellent service, please manifest it.

We have extended releases to the Relief Society General Presidency to be effective on August 1, 2022, as follows: Jean B. Bingham as President, Sharon Eubank as First Counselor, and Reyna I. Aburto as Second Counselor.

We have also extended releases to the Primary General Presidency, which will be effective on August 1, 2022, as follows: Camille N. Johnson as President, Susan H. Porter as First Counselor, and Amy A. Wright as Second Counselor.

All who wish to join us in expressing appreciation to these sisters for their devoted service, please manifest it.

It is proposed that we sustain the following as General Authority Seventies: Mark D. Eddy, James W. McConkie III, Isaac K. Morrison, Ryan K. Olsen, Jonathan S. Schmitt, and Denelson Silva.

All in favor, please manifest it.

Those opposed, by the same sign.

We note that 45 new Area Seventies were sustained during the general conference leadership meeting on Thursday, March 31, and then announced on newsroom.ChurchofJesusChrist.org earlier this week.

We invite you to sustain them in their new assignments.

Those in favor, please manifest it.

Any opposed, by the same sign.

It is proposed that we sustain the following as the new Relief Society General Presidency, to be effective on August 1, 2022: Camille N. Johnson as President, Jeannie Anette Dennis as First Counselor, and Kristin Mae Yee as Second Counselor.

Those in favor may manifest it.

Any opposed may so signify.

It is proposed that we sustain the following as the new Primary General Presidency, also to be effective on August 1, 2022: Susan H. Porter as President, Amy A. Wright as First Counselor, and Tracy Y. Browning as Second Counselor.

All in favor may manifest it.

Any opposed may so manifest.

It is proposed that we sustain the other General Authorities, Area Seventies, and General Officers as presently constituted.

All in favor, please manifest it.

Those opposed, if any.

Thank you, brothers and sisters, for your continued faith and prayers on behalf of the leaders of the Church.

\intalkheader{Changes to Area Seventies}

The following Area Seventies were sustained during a leadership session held as part of general conference:

Elzimar Gouvêa de Albuquerque, Roland J. Bäck, Raúl Barrón, Bruno V. Barros, Eric Baxter, Oscar Bedregal, Joep Boom, Michael P. Brady, Randall A. Brown, Kennedy F. Canuto, Stephen K. Christensen, Nathan A. Craig, Mark Anthony Dundon, Favio M. Durán, Amândio A. Feijó, Claude R. Gamiette, Scott L. Hymas, Jason C. Jensen, Roberto G. F. Leite, John W. Lewis, Paulo Renato Marinho, Blaine R. Maxfield, Eduardo R. Mora, David Ngabizele, João Luis Oppe, Justice N. Otuonye, Emanuel Petrignani, Daniel Piros, Craig W. J. Raeside, Nelson Ramírez, Alexey V. Samaykin, Jose Antonio San Gabriel, Jose Estuardo Sazo, Steven D. Shumway, Oswaldo J. Soto, Mark G. Stewart, Scott N. Taylor, Roseveltt de Pina Teixeira, Gordon L. Treadway, Harold Truque, Nikolai Ustyuzhaninov, Carlos Ernesto Velasco, Kyle A. Vest, Sergio Villa, and Min Zu Wang.

\newpage
\settalk{Introductory Message}
\setspeaker{Dallin H. Oaks}
\settalkdate{April 2022}
\talkheading{Introductory Message}{Dallin H. Oaks}

My dear sisters, as we begin this unique women’s session of general conference, I am pleased to deliver this introductory message from the First Presidency.

Our Saturday sessions have a history of different purposes and different audiences. This evening we add to that history as we embark upon a new purpose and procedure for the foreseeable future. The gospel of Jesus Christ does not change. Gospel doctrine does not change. Our personal covenants do not change. But over the years, the meetings we hold to communicate our messages do change and very likely will continue to change over the years.

For now, this Saturday evening meeting is a session of general conference, not a session of any organization. Like all sessions of general conference, the planning, speakers, and music are designated by the First Presidency.

We have asked President Jean B. Bingham, General President of the Relief Society, to conduct this session. Future Saturday evening sessions may be conducted by one of the other General Officers of the Church, such as members of the General Presidencies of Relief Society, Young Women, and Primary, designated by the First Presidency.

Tonight this Saturday evening session of general conference will concentrate on the concerns of Latter-day Saint women. This will include the doctrine of The Church of Jesus Christ of Latter-day Saints, the policies of the Church that relate specially to women, and the general responsibilities and work of the organizations that include the women and girls of the Church. Though this session is deliberately broadcast to a worldwide audience like all sessions of general conference, the audience invited to be present in the Conference Center for this session is women and girls age 12 and older. We have included some priesthood leaders who preside over the participating organizations.

What we are initiating here is responsive to the communication resources currently available to the Lord’s worldwide Church leadership and membership. The doctrine of the gospel of Jesus Christ is for everyone, so that is our principal motive and extent of dissemination. We honor the daughters of God in this special session by concentrating on their concerns and those of their organizations.

We are grateful that broadcast technology now gives Church leaders the capacity to conduct detailed training by addressing specific audiences in the field. We also welcome the fact that current travel opportunities are increasing. That allows us to send Church leaders to conduct needed regular face-to-face leadership training in the field.

This is the work of the Lord Jesus Christ. We are His servants, directed by His Holy Spirit. We invoke the blessings of our Lord upon the leaders of these organizations and upon the faithful women and girls who serve the Lord in these organizations and in their individual lives. In the name of Jesus Christ, amen.

\newpage
\settalk{Divine Love in the Father’s Plan}
\setspeaker{Dallin H. Oaks}
\settalkdate{April 2022}
\talkheading{Divine Love in the Father’s Plan}{Dallin H. Oaks}

The gospel plan shows our Heavenly Father’s love for all His children. To understand this, we must seek to understand His plan and His commandments. He loves His children so much that He gave His Only Begotten Son, Jesus Christ, to be our Savior and Redeemer, to suffer and die for us. In the restored Church of Jesus Christ of Latter-day Saints, we have a unique understanding of our Heavenly Father’s plan. This gives us a different way of viewing the purpose of mortal life, the divine judgment that follows it, and the ultimate glorious destiny of all of God’s children.

I love you, my brothers and sisters. I love all of God’s children. When Jesus was asked, “Which is the great commandment in the law?” He taught that to love God and to love our neighbors are the first of God’s great commandments. Those commands are first because they invite us to grow spiritually by seeking to imitate God’s love for us. I wish we all had a better understanding of the loving doctrine and policies that our Heavenly Father and His Son, Jesus Christ, have established in The Church of Jesus Christ of Latter-day Saints. What I say here seeks to clarify how God’s love explains that doctrine and the Church’s inspired policies.

\intalkheader{I.}

A common misunderstanding of the judgment that ultimately follows mortal life is that good people go to a place called heaven and bad people go to an everlasting place called hell. This erroneous assumption of only two ultimate destinations implies that those who cannot keep all the commandments required for heaven will necessarily be forever destined for hell.

A loving Heavenly Father has a better plan for His children. The revealed doctrine of the restored Church of Jesus Christ teaches that all the children of God—with exceptions too limited to consider here—will finally wind up in a kingdom of glory. “In my Father’s house are many mansions,” Jesus taught. From modern revelation we know that those mansions are in three different kingdoms of glory. In the Final Judgment each of us will be judged according to our deeds and the desires of our hearts. Before that, we will need to suffer for our unrepented sins. The scriptures are clear on that. Then our righteous Judge will grant us residence in one of those kingdoms of glory. Thus, as we know from modern revelation, all “shall be judged … , and every man shall receive according to his own works, his own dominion, in the mansions which are prepared.”

The Lord has chosen to reveal comparatively little about two of these kingdoms of glory. In contrast, the Lord has revealed much about the highest kingdom of glory, which the Bible describes as the “glory of the sun.”

In the “celestial” glory there are three degrees, or levels. The highest of these is exaltation in the celestial kingdom, wherein we may become like our Father and His Son, Jesus Christ. To help us develop the godly attributes and the change in nature necessary to realize our divine potential, the Lord has revealed doctrine and established commandments based on eternal law. This is what we teach in The Church of Jesus Christ of Latter-day Saints because the purpose of the doctrine and policies of this restored Church is to prepare God’s children for salvation in the celestial glory and, more particularly, for exaltation in its highest degree.

The covenants made and the blessings promised to the faithful in the temples of God are the key. This explains our worldwide building of temples, about which the choir has sung so beautifully. Some are puzzled at this emphasis, not understanding that the covenants and ordinances of the temple guide us toward achieving exaltation. This can be understood only in the context of the revealed truth of three degrees of glory. Because of our Heavenly Father’s great love for all of His children, He has provided other kingdoms of glory—as Elder Quentin L. Cook explained yesterday—all of which are more wonderful than we can comprehend.

The Atonement of Jesus Christ makes all of this possible. He has revealed that He “glorifies the Father, and saves all the works of his hands.” That salvation is granted in different kingdoms of glory. We know from modern revelation that “all kingdoms have a law given.” Significantly:

“He who is not able to abide the law of a celestial kingdom cannot abide a celestial glory.

“And he who cannot abide the law of a terrestrial kingdom cannot abide a terrestrial glory.

“And he who cannot abide the law of a telestial kingdom cannot abide a telestial glory.”

In other words, the kingdom of glory we receive in the Final Judgment is determined by the laws we choose to abide by in our Heavenly Father’s loving plan. Under that plan there are multiple kingdoms so that all of His children can be assigned to a kingdom where they can “abide.”

\intalkheader{II.}

The teachings and policies of the Lord’s restored Church apply these eternal truths in a way that can be fully understood only in the context of our Heavenly Father’s loving plan for all of His children.

Thus, we honor individual agency. Most are aware of this Church’s great efforts to promote religious freedom. These efforts are in furtherance of our Heavenly Father’s plan. We seek to help all of His children—not just our own members—enjoy the precious freedom to choose.

Similarly, we are sometimes asked why we send missionaries to so many nations, even among Christian populations. We are also asked why we give enormous humanitarian aid to persons who are not members of our Church without linking this to our missionary efforts. We do this because the Lord has taught us to esteem all of His children as our brothers and sisters, and we want to share our spiritual and temporal abundance with everyone.

Eternal doctrine also provides a distinctive perspective on children. Through this perspective we see the bearing and nurturing of children as part of the divine plan. It is a joyful and sacred duty of those given the power to participate in it. Therefore, we are commanded to teach and contend for principles and practices that provide the best conditions for the development and happiness of children under God’s plan.

\intalkheader{III.}

Finally, The Church of Jesus Christ of Latter-day Saints is properly known as a family-centered Church. But not well understood is the reality that our family-centeredness is not limited to mortal relationships. Eternal relationships are also fundamental to our theology. The mission of the restored Church is to help all the children of God qualify for what God desires as their ultimate destiny. By the redemption provided through the Atonement of Christ, all may attain eternal life (exaltation in the celestial kingdom), which Mother Eve declared “God giveth unto all the obedient.” This is more than salvation. President Russell M. Nelson has reminded us that “in God’s eternal plan, salvation is an individual matter; [but] exaltation is a family matter.”

Fundamental to us is God’s revelation that exaltation can be attained only through faithfulness to the covenants of an eternal marriage between a man and a woman. That divine doctrine is why we teach that “gender is an essential characteristic of individual premortal, mortal, and eternal identity and purpose.”

That is also why the Lord has required His restored Church to oppose social and legal pressures to retreat from His doctrine of marriage between a man and a woman, to oppose changes that homogenize the differences between men and women or confuse or alter gender.

The restored Church’s positions on these fundamentals frequently provoke opposition. We understand that. Our Heavenly Father’s plan allows for “opposition in all things,” and Satan’s most strenuous opposition is directed at whatever is most important to that plan. Consequently, he seeks to oppose progress toward exaltation by distorting marriage, discouraging childbearing, or confusing gender. However, we know that in the long run, the divine purpose and plan of our loving Heavenly Father will not be changed. Personal circumstances may change, and God’s plan assures that in the long run, the faithful who keep their covenants will have the opportunity to qualify for every promised blessing.

A uniquely valuable teaching to help us prepare for eternal life, “the greatest of all the gifts of God,” is the 1995 proclamation on the family. Its declarations are, of course, different from some current laws, practices, and advocacy, such as cohabitation and same-sex marriage. Those who do not fully understand the Father’s loving plan for His children may consider this family proclamation no more than a changeable statement of policy. In contrast, we affirm that the family proclamation, founded on irrevocable doctrine, defines the kind of family relationships where the most important part of our eternal development can occur.

That is the context for the unique doctrine and policies of the restored Church of Jesus Christ of Latter-day Saints.

\intalkheader{IV.}

In many relationships and circumstances in mortal life, each of us must live with differences. As followers of Christ who should love our fellow men, we should live peacefully with those who do not believe as we do. We are all children of a loving Heavenly Father. For all of us, He has destined life after death and, ultimately, a kingdom of glory. God desires all of us to strive for His highest possible blessings by keeping His highest commandments, covenants, and ordinances, all of which culminate in His holy temples being built throughout the world. We must seek to share these truths of eternity with others. But with the love we owe to all of our neighbors, we always accept their decisions. As a Book of Mormon prophet taught, we must press forward, having “a love of God and of all men.”

As President Russell M. Nelson declared in our last conference: “There has never been a time in the history of the world when knowledge of our Savior is more personally vital and relevant to every human soul. … The pure doctrine of Christ is powerful. It changes the life of everyone who understands it and seeks to implement it in his or her life.”

May we all implement that sacred doctrine in our own lives, I pray in the name of Jesus Christ, amen.

\newpage
\settalk{Helping the Poor and Distressed}
\setspeaker{Dallin H. Oaks}
\settalkdate{October 2022}
\talkheading{Helping the Poor and Distressed}{Dallin H. Oaks}

Brothers and sisters, our beloved President Russell M. Nelson will address us later in this session. He has asked me to be the first speaker.

My subject today concerns what The Church of Jesus Christ of Latter-day Saints and its members give and do for the poor and distressed. I will also speak of similar giving by other good people. Giving to those in need is a principle in all Abrahamic religions and in others as well.

A few months ago, The Church of Jesus Christ of Latter-day Saints reported for the first time the extent of our humanitarian work worldwide. Our 2021 expenditures for those in need in 188 countries worldwide totaled \$906 million—almost a billion dollars. In addition, our members volunteered over 6 million hours of labor in the same cause.

Those figures are, of course, an incomplete report of our giving and helping. They do not include the personal services our members give individually as they minister to one another in called positions and voluntary member-to-member service. And our 2021 report makes no mention of what our members do individually through innumerable charitable organizations not formally connected with our Church. I begin with these.

In 1831, less than two years after the restored Church was organized, the Lord gave this revelation to guide its members and, I believe, all of His children worldwide:

“Behold, it is not meet that I should command in all things; for he that is compelled in all things, the same is a slothful and not a wise servant. …

“Verily I say, men should be anxiously engaged in a good cause, and do many things of their own free will, and bring to pass much righteousness;

“For the power is in them, wherein they are agents unto themselves. And inasmuch as men do good they shall in nowise lose their reward.”

In more than 38 years as an Apostle and over 30 years of professional employment, I have seen many generous efforts by organizations and persons of the kind this revelation describes as “a good cause” and “bring[ing] to pass much righteousness.” There are uncounted examples of such humanitarian service throughout the world, beyond our own borders and beyond our common knowledge. Contemplating this, I think of the Book of Mormon prophet-king Benjamin, whose sermon included this eternal truth: “When ye are in the service of your fellow beings ye are only in the service of your God.”

Much welfare and humanitarian service to our fellow beings is taught and practiced by The Church of Jesus Christ of Latter-day Saints and by us as its members. For example, we fast at the first of each month and contribute at least the equivalent of the uneaten meals to help those in need in our own congregations. The Church also makes enormous contributions for humanitarian and other services throughout the world.

Despite all that our Church does directly, most humanitarian service to the children of God worldwide is carried out by persons and organizations having no formal connection with our Church. As one of our Apostles observed: “God is using more than one people for the accomplishment of his great and marvelous work. … It is too vast, too arduous, for any one people.” As members of the restored Church, we need to be more aware and more appreciative of the service of others.

The Church of Jesus Christ is committed to serving those in need, and it is also committed to cooperating with others in that effort. We recently made a large gift to the United Nations World Food Programme. Over the many decades of our humanitarian work, two organizations stand out as key collaborators: projects with the Red Cross and Red Crescent agencies in dozens of countries have provided the children of God crucial relief during natural disasters and conflicts. Likewise, we have a long record of assistance with Catholic Relief Services. These organizations have taught us much about world-class relief.

We have also had fruitful collaborations with other organizations, including Muslim Aid, Water for People, and IsraAID, to name just a few. While each humanitarian organization has its own areas of specialization, we share the common goal of relieving suffering among God’s children. All of this is part of God’s work for His children.

Modern revelation teaches that our Savior, Jesus Christ, is “the true light that lighteth every man that cometh into the world.” By this, all the children of God are enlightened to serve Him and one another to the best of their knowledge and ability.

The Book of Mormon teaches that “every thing which inviteth and enticeth to do good, and to love God, and to serve him, is inspired of God.”

Continuing:

“For behold, the Spirit of Christ is given to every man, that he may know good from evil; wherefore, I show unto you the way to judge; for every thing which inviteth to do good, and to persuade to believe in Christ, is sent forth by the power and gift of Christ. …

“And now, my brethren, … ye know the light by which ye may judge, which light is the light of Christ.”

Here are some examples of children of God helping other children of God with their vital needs for food, medical care, and teaching:

Ten years ago, the Kandharis, a Sikh husband and wife in the United Arab Emirates, personally launched a remarkable effort to feed the hungry. Through the Guru Nanak Darbar Sikh temple, they are currently serving over 30,000 vegetarian meals every weekend to anyone who enters their doors, regardless of religion or race. Dr. Kandhari explains, “We believe that all are one, we are children of one God, and we are here to serve humanity.”

The provision of medical and dental care to those in need is another example. In Chicago, I met a Syrian-American critical care physician, Dr. Zaher Sahloul. He is one of the founders of MedGlobal, which organizes medical professionals to volunteer their time, skills, knowledge, and leadership to help others in crises, such as in the Syrian war, where Dr. Sahloul risked his life in giving medical care to civilians. MedGlobal and similar organizations (including many Latter-day Saint professionals) demonstrate that God is moving professionals of faith to bring the poor needed relief worldwide.

Many unselfish children of God are involved in teaching efforts, also worldwide. A good example, known to us through our humanitarian efforts, is the activity of a man known as Mr. Gabriel, who has been a refugee from various conflicts on several occasions. He recently observed that hundreds of thousands of refugee children in East Africa needed help to keep their hopes alive and their minds active. He organized other teachers in the refugee population into what they called “tree schools,” where children were gathered for lessons under the shade of a tree. He did not wait for others to organize or direct but personally led efforts that have provided learning opportunities for thousands of primary school children during stressful years of displacement.

Of course, these three examples do not mean that everything said or done by organizations or individuals purporting to be good or of God is truly that. These examples do show that God inspires many organizations and individuals to do much good. It also shows that more of us should be recognizing the good done by others and supporting it as we have the time and means to do so.

Here are some examples of service the Church supports and which our members and other good people and organizations also support with individual donations of time and money:

I begin with religious freedom. In supporting that, we serve our own interests but also the interests of other religions. As our first President, Joseph Smith, taught, “We claim the privilege of worshiping Almighty God according to the dictates of our own conscience, and allow all men the same privilege, let them worship how, where, or what they may.”

Other examples of the restored Church’s humanitarian and other assistance that are also supported voluntarily by our members are our well-known schools, colleges, and universities and our less-known but now published large donations for the relief of those suffering from the destructions and dislocations of natural disasters like tornadoes and earthquakes.

Other charitable activities our members support by their voluntary donations and efforts are too numerous to list, but just mentioning these few will suggest their variety and importance: combating racism and other prejudices; researching how to prevent and cure diseases; helping the disabled; supporting music organizations; and improving the moral and physical environment for all.

All of the humanitarian efforts of The Church of Jesus Christ of Latter-day Saints seek to follow the example of a righteous people described in the Book of Mormon: “And thus, in their prosperous circumstances, they did not send away any who were naked, or that were hungry, or that were athirst, or that were sick, … and they … were liberal to all, both old and young, both bond and free, both male and female, whether out of the church or in the church.”

I testify of Jesus Christ, whose light and Spirit guide all of the children of God in helping the poor and distressed throughout the world. In the name of Jesus Christ, amen.

\newpage
\settalk{Sustaining of General Authorities, Area Seventies, and General Officers}
\setspeaker{Presented by President Dallin H. Oaks}
\settalkdate{April 2023}
\talkheading{Sustaining of General Authorities, Area Seventies, and General Officers}{Presented by President Dallin H. Oaks}

Brothers and sisters, it is my privilege to present the General Authorities, Area Seventies, and General Officers of the Church for your sustaining vote.

Please express your support in the usual way wherever you may be. If there are those who oppose any of the proposals, we ask that you contact your stake president.

It is proposed that we sustain Russell Marion Nelson as prophet, seer, and revelator and President of The Church of Jesus Christ of Latter-day Saints; Dallin Harris Oaks as First Counselor in the First Presidency; and Henry Bennion Eyring as Second Counselor in the First Presidency.

Those in favor may manifest it.

Those opposed, if any, may manifest it.

It is proposed that we sustain Dallin H. Oaks as President of the Quorum of the Twelve Apostles and M. Russell Ballard as Acting President of the Quorum of the Twelve Apostles.

Those in favor, please signify.

Any opposed may manifest it.

It is proposed that we sustain the following as members of the Quorum of the Twelve Apostles: M. Russell Ballard, Jeffrey R. Holland, Dieter F. Uchtdorf, David A. Bednar, Quentin L. Cook, D. Todd Christofferson, Neil L. Andersen, Ronald A. Rasband, Gary E. Stevenson, Dale G. Renlund, Gerrit W. Gong, and Ulisses Soares.

Those in favor, please manifest it.

Any opposed may so indicate.

It is proposed that we sustain the counselors in the First Presidency and the Quorum of the Twelve Apostles as prophets, seers, and revelators.

All in favor, please manifest it.

Contrary, if there be any, by the same sign.

The following General Authority Seventies will be released from their assignments and given emeritus status, effective on August 1, 2023: Elders Benjamín De Hoyos, Juan A. Uceda, and Kazuhiko Yamashita.

Those who wish to express gratitude to these brethren and to their wives and families for their years of dedicated service throughout the world may do so by the uplifted hand.

The following Area Seventies have been released from their assignments, effective immediately: J. Kimo Esplin and Alan T. Phillips.

Those who wish to join us in expressing appreciation for their excellent service, please manifest it.

We note with appreciation the other Area Seventies who will complete their service this year and whose names can be found on the Church’s website.

Those who wish to join in expressing gratitude to these brethren for their selfless service may manifest it.

Effective immediately, we hereby release Brothers Ahmad S. Corbitt and Bradley Ray Wilcox from serving as First and Second Counselors in the Young Men General Presidency.

Those who wish to show appreciation to these brethren for their service may so manifest.

We extend releases to the Young Women General Presidency, effective on August 1, 2023, as follows: Bonnie H. Cordon as President, Michelle D. Craig as First Counselor, and Rebecca L. Craven as Second Counselor.

All who wish to join us in expressing appreciation to these sisters for their devoted service, please manifest it.

It is proposed that we sustain the following as General Authority Seventies: Ahmad S. Corbitt, Robert M. Daines, J. Kimo Esplin, Christophe G. Giraud-Carrier, and Alan T. Phillips.

All in favor, please manifest it.

Those opposed, by the same sign.

We note that 61 new Area Seventies were sustained during the general conference leadership meetings on Thursday, March 30, and then announced on the Church’s website.

We invite you to sustain these brethren in their new assignments.

Those in favor, please manifest it.

Any opposed, by the same sign.

It is proposed that we sustain the following as the new Young Women General Presidency, to be effective on August 1, 2023: Emily Belle Freeman as President, Tamara Wood Runia as First Counselor, and Andrea Muñoz Spannaus as Second Counselor.

Those in favor may manifest it.

Any opposed may so signify.

It is proposed that we sustain the following as counselors in the Young Men General Presidency, effective immediately: Bradley Ray Wilcox as First Counselor and Michael T. Nelson as Second Counselor.

All in favor may manifest it.

Any opposed may so manifest.

It is proposed that we sustain the other General Authorities, Area Seventies, and General Officers as presently constituted.

All in favor may do so by the uplifted hand.

Those opposed, if any.

Thank you, brothers and sisters, for your continued faith and prayers on behalf of the leadership of the Church.

Changes to Area Seventies

The following Area Seventies were sustained during a leadership session held as part of general conference:

Isaías Alcala, John D. Amos, Johnny O. Baddoo, Victor O. Bassey, Adrian Bettridge, A. Kaulle Bezerra, Carlos G. Cantero, Emerson B. Carnavale, Orlando A. Castaños, Bun Huoch Eng, Hutch U. Fale, Fernando R. García, Tomás García, C. Alan Gauldin, Aaron T. Hall, Darwin W. Halvorson, Jed J Hancock, Henry Herrera, Ndalamba Ilunga, Samuel M. T. Koivisto, Carlos J. Lantigua, Esau Lara, Stephen J. Larson, Thabo Lebethoa, G. Kenneth Lee, Israel Marin, Wayne E. Maurer, Lee G. McCann II, Robert Mendenhall, Adrian Mendez, Siyabonga Mkhize, Javier F. Monestel, Thomas B. Morgan, Jared V. Ormsby, Z. Rudy Palhua, Arturo D. Palmieri, Kenneth Pambu, Hugo O. Panameño, Kevin J. Parks, Paul Picard, David J. Pickett, Martin Pilka, Irineu E. Prado, Christopher R. Price, Miguel Ribeiro, James N. Robinson, Edward B. Rowe, Robert Schwartz, Gregory A. Scott, Dominic R. Sénéchal, Kofi G. Sosu, Michael B. Strong, Nithya Kumar Sunderraj, Thomas A. Thomas, Alejandro H. Treviño, Nefi M. Trujillo, Chimaroke G. Udeichi, Fernando Valdes, Helton C. Vecchi, Brent B. Ward, and Tomasito S. Zapanta.

The following Area Seventies will be released on or before August 1, 2023:

Richard K. Ahadjie, Duane D. Bell, Hubermann Bien Aimé, Víctor R. Calderón, Michel J. Carter, Daniel Córdova, John N. Craig, William H. K. Davis, Fernando P. Del Carpio, Richard J. DeVries, Kylar G. Dominguez, Torben Engbjerg, Kenneth J. Firmage, Edgar Flores, Silvio Flores, Carlos A. Genaro, Mark A. Gilmour, Sergio A. Gómez, Roberto Gonzalez, Virgilio Gonzalez, Spencer R. Griffin, Marcel Guei, Oleksiy H. Hakalenko, Matthew S. Harding, David J. Harris, Kevin J. Hathaway, José Hernández, Glenn M. Holmes, Richard Neitzel Holzapfel, Okechukwu I. Imo, Michael D. Jones, Pungwe S. Kongolo, Ricardo C. Leite, Aretemio C. Maligon, Edgar A. Mantilla, Lincoln P. Martins, Carl R. Maurer, Daniel S. Mehr II, Glen D. Mella, Tomas S. Merdegia Jr., Allistair B. Odgers, R. Jeffrey Parker, Victor P. Patrick, Denis E. Pineda, Henrique S. Simplicio, Jeffrey H. Singer, Michael L. Staheli, Jeffrey K. Wetzel, Michael S. Wilstead, and David L. Wright.

\newpage
\settalk{The Teachings of Jesus Christ}
\setspeaker{Dallin H. Oaks}
\settalkdate{April 2023}
\talkheading{The Teachings of Jesus Christ}{Dallin H. Oaks}

We believe in Christ. As members of The Church of Jesus Christ of Latter-day Saints, we worship Him and follow His teachings in the scriptures.

Before the Fall, our Heavenly Father spoke directly to Adam and Eve. Thereafter, the Father introduced His Only Begotten Son, Jesus Christ, as our Savior and Redeemer and gave us the command to “hear Him.” From this direction we conclude that the scriptural records of words spoken by “God” or the “Lord” are almost always the words of Jehovah, our risen Lord, Jesus Christ.

We are given the scriptures to direct our lives. As the prophet Nephi taught us, we should “feast upon the words of Christ; for behold, the words of Christ will tell you all things what ye should do.” Most of the scriptures reporting Jesus’s mortal ministries are descriptions of what He did. My message today consists of a selection of the words of our Savior—what He said. These are words recorded in the New Testament (including the inspired additions of Joseph Smith) and in the Book of Mormon. Most of these selections are in the sequence in which our Savior spoke them.

“Verily, verily, I say unto thee, Except a man be born of water and of the Spirit, he cannot enter into the kingdom of God.”

“Blessed are … they [which] do hunger and thirst after righteousness, for they shall be filled with the Holy Ghost.”

“Blessed are the peacemakers: for they shall be called the children of God.”

“Ye have heard that it was said by them of old time, Thou shalt not commit adultery:

“But I say unto you, That whosoever looketh on a woman to lust after her hath committed adultery with her already in his heart.”

“Ye have heard that it hath been said, Thou shalt love thy neighbour, and hate thine enemy.

“But I say unto you, Love your enemies, bless them that curse you, do good to them that hate you, and pray for them which despitefully use you, and persecute you;

“That ye may be the children of your Father which is in heaven: for he maketh his sun to rise on the evil and on the good, and sendeth rain on the just and on the unjust.”

“If ye forgive men their trespasses, your heavenly Father will also forgive you:

“But if ye forgive not men their trespasses, neither will your Father forgive your trespasses.”

“If ye were of the world, the world would love his own: but because ye are not of the world, but I have chosen you out of the world, therefore the world hateth you.”

“Wherefore, seek not the things of this world but seek ye first to build up the kingdom of God, and to establish his righteousness, and all these things shall be added unto you.”

“Therefore all things whatsoever ye would that men should do to you, do ye even so to them: for this is the law and the prophets.”

“Beware of false prophets, which come to you in sheep’s clothing, but inwardly they are ravening wolves.

“Ye shall know them by their fruits. Do men gather grapes of thorns, or figs of thistles?

“Even so every good tree bringeth forth good fruit; but a corrupt tree bringeth forth evil fruit.”

“Not every one that saith unto me, Lord, Lord, shall enter into the kingdom of heaven; but he that doeth the will of my Father which is in heaven.”

“Come unto me, all ye that labour and are heavy laden, and I will give you rest.

“Take my yoke upon you, and learn of me; for I am meek and lowly in heart: and ye shall find rest unto your souls.

“For my yoke is easy, and my burden is light.”

“If any man will come after me, let him deny himself, and take up his cross and follow me.

“And now for a man to take up his cross, is to deny himself all ungodliness, and every worldly lust, and keep my commandments.”

“Therefore, forsake the world, and save your souls; for what is a man profited, if he shall gain the whole world, and lose his own soul? Or what shall a man give in exchange for his soul?”

“If any man will do his will, he shall know of the doctrine, whether it be of God, or whether I speak of myself.”

“And I say unto you, Ask, and it shall be given you; seek, and ye shall find; knock, and it shall be opened unto you.

“For every one that asketh receiveth; and he that seeketh findeth; and to him that knocketh it shall be opened.”

“Other sheep I have, which are not of this fold: them also I must bring, and they shall hear my voice; and there shall be one fold, and one shepherd.”

“Jesus said unto her, I am the resurrection, and the life: he that believeth in me, though he were dead, yet shall he live:

“And whosoever liveth and believeth in me shall never die.”

“[The great commandment in the law is this:] Thou shalt love the Lord thy God with all thy heart, and with all thy soul, and with all thy mind.

“This is the first and great commandment.

“And the second is like unto it, Thou shalt love thy neighbour as thyself.

“On these two commandments hang all the law and the prophets.”

“He that hath my commandments, and keepeth them, he it is that loveth me: and he that loveth me shall be loved of my Father, and I will love him, and will manifest myself to him.”

“Peace I leave with you, my peace I give unto you: not as the world giveth, give I unto you. Let not your heart be troubled, neither let it be afraid.”

“This is my commandment, That ye love one another, as I have loved you.”

“Behold my hands and my feet, that it is I myself: handle me, and see; for a spirit hath not flesh and bones, as ye see me have.”

“Go ye therefore, and teach all nations, baptizing them in the name of the Father, and of the Son, and of the Holy Ghost:

“Teaching them to observe all things whatsoever I have commanded you: and, lo, I am with you alway, even unto the end of the world.”

After His ministry in the Holy Land, Jesus Christ appeared to the righteous on the American continent. These are some of the words He spoke there:

“Behold, I am Jesus Christ the Son of God. I created the heavens and the earth, and all things that in them are. I was with the Father from the beginning. I am in the Father, and the Father in me; and in me hath the Father glorified his name.”

“I am the light and the life of the world. I am Alpha and Omega, the beginning and the end.

“And ye shall offer up unto me no more the shedding of blood; yea, your sacrifices and your burnt offerings shall be done away, for I will accept none of your sacrifices and your burnt offerings.

“And ye shall offer for a sacrifice unto me a broken heart and a contrite spirit. And whoso cometh unto me with a broken heart and a contrite spirit, him will I baptize with fire and with the Holy Ghost. …

“Behold, I have come unto the world to bring redemption unto the world, to save the world from sin.”

“And again I say unto you, ye must repent, and be baptized in my name, and become as a little child, or ye can in nowise inherit the kingdom of God.”

“Therefore I would that ye should be perfect even as I, or your Father who is in heaven is perfect.”

“Verily, verily, I say unto you, ye must watch and pray always, lest ye be tempted by the devil, and ye be led away captive by him.”

“Therefore ye must always pray unto the Father in my name.”

“Therefore, whatsoever ye shall do, ye shall do it in my name; therefore ye shall call the church in my name.”

“Behold I have given unto you my gospel, and this is the gospel which I have given unto you—that I came into the world to do the will of my Father, because my Father sent me.

“And my Father sent me that I might be lifted up upon the cross; and after that I had been lifted up upon the cross, that I might draw all men unto me … to be judged of their works, whether they be good or whether they be evil.”

“Now this is the commandment: Repent, all ye ends of the earth, and come unto me and be baptized in my name, that ye may be sanctified by the reception of the Holy Ghost, that ye may stand spotless before me at the last day.”

We believe in Christ. I conclude with what He said about how we should know and follow His teachings:

“But the Comforter, which is the Holy Ghost, whom the Father will send in my name, he shall teach you all things, and bring all things to your remembrance, whatsoever I have said unto you.”

I affirm the truth of these teachings in the name of Jesus Christ, amen.

\newpage
\settalk{Kingdoms of Glory}
\setspeaker{Dallin H. Oaks}
\settalkdate{October 2023}
\talkheading{Kingdoms of Glory}{Dallin H. Oaks}

Members of The Church of Jesus Christ of Latter-day Saints are frequently asked, “How is your church different from other Christian churches?” Among the answers we give is the fulness of the doctrine of Jesus Christ. Foremost among that doctrine is the fact that our Heavenly Father loves all His children so much that He wants us all to live in a kingdom of glory forever. Moreover, He wants us to live with Him and His Son, Jesus Christ, eternally. His plan gives us the teachings and the opportunity to make the choices that will assure us the destiny and the life we choose.

\intalkheader{I.}

From modern revelation we know that the ultimate destiny of all who live on the earth is not the inadequate idea of heaven for the righteous and the eternal sufferings of hell for the rest. God’s loving plan for His children includes this reality taught by our Savior, Jesus Christ: “In my Father’s house are many mansions.”

The revealed doctrine of the restored Church of Jesus Christ of Latter-day Saints teaches that all the children of God—with exceptions too limited to consider here—will ultimately inherit one of three kingdoms of glory, even the least of which “surpasses all understanding.” After a period in which the disobedient suffer for their sins, which suffering prepares them for what is to follow, all will be resurrected and proceed to the Final Judgment of the Lord Jesus Christ. There, our loving Savior, who, we are taught, “glorifies the Father, and saves all the works of his hands,” will send all the children of God to one of these kingdoms of glory according to the desires manifested through their choices.

Another unique doctrine and practice of the restored Church is the revealed commandments and covenants that offer all the children of God the sacred privilege of qualifying for the highest degree of glory in the celestial kingdom. That highest destination—exaltation in the celestial kingdom—is the focus of The Church of Jesus Christ of Latter-day Saints.

From modern revelation, Latter-day Saints have this unique understanding of God’s plan of happiness for His children. That plan begins with our life as spirits before we were born, and it reveals the purpose and conditions of our chosen journey in mortality and our desired destination thereafter.

\intalkheader{II.}

We know from modern revelation that “all kingdoms have a law given” and that the kingdom of glory we receive in the Final Judgment is determined by the laws we choose to follow in our mortal journey. Under that loving plan, there are multiple kingdoms—many mansions—so that all of God’s children will inherit a kingdom of glory whose laws they can comfortably “abide.”

As we describe the nature and requirements of each of the three kingdoms in the Father’s plan, we begin with the highest, which is the focus of the divine commandments and ordinances God has revealed through The Church of Jesus Christ of Latter-day Saints. In the “celestial” glory there are three levels, of which the highest is exaltation in the celestial kingdom. This is the dwelling of those “who have received of his fulness, and of his glory,” wherefore, “they are gods, even the sons [and daughters] of God” and “dwell in the presence of God and his Christ forever and ever.” Through revelation, God has revealed the eternal laws, ordinances, and covenants that must be observed to develop the godly attributes necessary to realize this divine potential. The Church of Jesus Christ of Latter-day Saints focuses on these because the purpose of this restored Church is to prepare God’s children for salvation in the celestial glory and, more particularly, for exaltation in its highest degree.

God’s plan, founded on eternal truth, requires that exaltation can be attained only through faithfulness to the covenants of an eternal marriage between a man and a woman in the holy temple, which marriage will ultimately be available to all the faithful. That is why we teach that “gender is an essential characteristic of individual premortal, mortal, and eternal identity and purpose.”

A uniquely valuable teaching to help us prepare for exaltation is the 1995 proclamation on the family. Its declarations clarify the requirements that prepare us to live with God the Father and His Son, Jesus Christ. Those who do not fully understand the Father’s loving plan for His children may consider this family proclamation no more than a changeable statement of policy. In contrast, we affirm that the family proclamation, founded on irrevocable doctrine, defines the mortal family relationship where the most important part of our eternal development can occur.

The Apostle Paul describes the three degrees of glory, likening them to the glories of the sun, moon, and stars. He names the highest “celestial” and the second “terrestrial.” He does not name the lowest, but a revelation to Joseph Smith added its name: “telestial.” Another revelation also describes the nature of the persons to be assigned to each of these kingdoms of glory. Those who do not choose “to abide the law of a celestial kingdom” will inherit another kingdom of glory, lesser than the celestial but suited to the laws they have chosen and can comfortably “abide.” That word abide, so common in the scriptures, means a secure placement. For example, those in the terrestrial kingdom—comparable to the popular concept of heaven—“are they who receive of the presence of the Son, but not of the fulness of the Father.” They were “honorable men of the earth, who were blinded by the craftiness of men,” but “not valiant in the testimony of Jesus.”

The revealing description of those assigned to the lowest of the kingdoms of glory, the telestial, is “he who cannot abide … a terrestrial glory.” That describes those who reject the Savior and have observed no divine limits on their behavior. This is the kingdom where the wicked abide, after they have suffered for their sins. These are described in modern revelation as “they who received not the gospel of Christ, neither the testimony of Jesus. …

“These are they who are liars, and sorcerers, and adulterers, and whoremongers, and whosoever loves and makes a lie.”

Speaking of the three kingdoms of glory with his prophetic vision, President Russell M. Nelson recently wrote: “Mortal lifetime is barely a nanosecond compared with eternity. But what a crucial nanosecond it is! Consider carefully how it works: During this mortal life you get to choose which laws you are willing to obey—those of the celestial kingdom, or the terrestrial, or the telestial—and, therefore, in which kingdom of glory you will live forever. What a plan! It is a plan that completely honors your agency.”

\intalkheader{III.}

The Apostle Paul taught that the Lord’s teachings and commandments were given that we may all attain “the measure of the stature of the fulness of Christ.” That process requires far more than acquiring knowledge. It is not even enough to be convinced of the gospel; we must act so that we are converted by it. In contrast to other preaching, which teaches us to know something, the gospel of Jesus Christ challenges us to become something.

From such teachings we conclude that the Final Judgment is not just an evaluation of a sum total of good and evil acts—what we have done. It is based on the final effect of our acts and thoughts—what we have become. We qualify for eternal life through a process of conversion. As used here, this word of many meanings signifies a profound change of nature. It is not enough for anyone just to go through the motions. The commandments, ordinances, and covenants of the gospel are not a list of deposits required to be made in some heavenly account. The gospel of Jesus Christ is a plan that shows us how to become what our Heavenly Father desires us to become.

\intalkheader{IV.}

Because of Jesus Christ and His Atonement, when we fall short in this life, we can repent and rejoin the covenant path that leads to what our Heavenly Father desires for us.

The Book of Mormon teaches that “this life is the time for [us] to prepare to meet God.” But that challenging limitation to “this life” was given a hopeful context (at least to some extent for some persons) by what the Lord revealed to President Joseph F. Smith, now recorded in Doctrine and Covenants section 138. “I beheld,” the prophet wrote, “that the faithful elders of this dispensation, when they depart from mortal life, continue their labors in the preaching of the gospel of repentance and redemption, through the sacrifice of the Only Begotten Son of God, among those who are in darkness and under the bondage of sin in the great world of the spirits of the dead.

“The dead who repent will be redeemed, through obedience to the ordinances of the house of God,

“And after they have paid the penalty of their transgressions, and are washed clean, shall receive a reward according to their works, for they are heirs of salvation.”

In addition, we know that the Millennium, the thousand years that follow the Second Coming of the Savior, will be a time to perform the required ordinances for those who have not received them in their mortal lives.

There is much we do not know about the three major periods in the plan of salvation and their relationship to one another: (1) the premortal spirit world, (2) mortality, and (3) the next life. But we do know these eternal truths: “Salvation is an individual matter, but exaltation is a family matter.” We have a loving Heavenly Father who will see that we receive every blessing and every advantage that our own desires and choices allow. We also know that He will force no one into a sealing relationship against his or her will. The blessings of a sealed relationship are assured for all who keep their covenants but never by forcing a sealed relationship on another person who is unworthy or unwilling.

My dear brothers and sisters, I testify of the truth of these things. I testify of our Lord Jesus Christ, “the author and finisher of our faith,” whose Atonement, under the plan of our Father in Heaven, makes it all possible, in the name of Jesus Christ, amen.

\newpage
\settalk{Sustaining of General Authorities, Area Seventies, and General Officers}
\setspeaker{Presented by President Dallin H. Oaks}
\settalkdate{April 2024}
\talkheading{Sustaining of General Authorities, Area Seventies, and General Officers}{Presented by President Dallin H. Oaks}

Brothers and sisters, it will now be my privilege to present the General Authorities, Area Seventies, and General Officers of the Church for your sustaining vote.

Please express your support in the usual way. If there are those who oppose any of the proposals, we ask that you contact your stake president.

It is proposed that we sustain Russell Marion Nelson as prophet, seer, and revelator and President of The Church of Jesus Christ of Latter-day Saints; Dallin Harris Oaks as First Counselor in the First Presidency; and Henry Bennion Eyring as Second Counselor in the First Presidency.

Those in favor may manifest it.

Those opposed, if any, may manifest it.

It is proposed that we sustain Dallin H. Oaks as President of the Quorum of the Twelve Apostles and Jeffrey R. Holland as Acting President of the Quorum of the Twelve Apostles.

Those in favor, please signify.

Any opposed may manifest it.

It is proposed that we sustain the following as members of the Quorum of the Twelve Apostles: Jeffrey R. Holland, Dieter F. Uchtdorf, David A. Bednar, Quentin L. Cook, D. Todd Christofferson, Neil L. Andersen, Ronald A. Rasband, Gary E. Stevenson, Dale G. Renlund, Gerrit W. Gong, Ulisses Soares, and Patrick Kearon.

Those in favor, please manifest it.

Any opposed may so indicate.

It is proposed that we sustain the counselors in the First Presidency and the Quorum of the Twelve Apostles as prophets, seers, and revelators.

All in favor, please manifest it.

Contrary, if there be any, by the same sign.

The following General Authorities will be released from their assignments and given emeritus status, effective on August 1, 2024: Elders Ian S. Ardern, Shayne M. Bowen, Paul V. Johnson, S. Gifford Nielsen, Brent H. Nielson, Adrián Ochoa, Gary B. Sabin, and Evan A. Schmutz.

Those who wish to express gratitude to these brethren and their wives for their years of dedicated service throughout the Church may do so by the uplifted hand.

We also release Elder Carlos A. Godoy from serving as a member of the Presidency of the Seventy, effective on August 1, 2024.

Those who wish to express appreciation to Elder Godoy for his service in this capacity may do so.

We note with gratitude the Area Seventies who will complete their service and whose names can be found on the Church’s website.

Those who wish to join in expressing appreciation to these brethren and to their families for their years of selfless service may manifest it.

We extend releases to the Sunday School General Presidency, effective on August 1, 2024, as follows: Mark L. Pace as President, Milton Camargo as First Counselor, and Jan E. Newman as Second Counselor.

All who wish to join in expressing appreciation to these brothers for their devoted service may manifest it.

It is proposed that we sustain the following as members of the Presidency of the Seventy: Elder Marcus B. Nash, who began his service in January 2024, and Elders Michael T. Ringwood, Arnulfo Valenzuela, and Edward Dube, who will begin their service on August 1, 2024.

Those in favor may manifest it.

Any opposed may manifest it.

It is proposed that we sustain the following as General Authority Seventies: David L. Buckner, Gregorio E. Casillas, Aroldo B. Cavalcante, I. Raymond Egbo, D. Martin Goury, Karl D. Hirst, Christopher H. Kim, Sandino Roman, Steven D. Shumway, Michael B. Strong, and Sergio R. Vargas.

All in favor, please manifest it.

Those opposed, by the same sign.

We note that 64 new Area Seventies were sustained during the general conference leadership meetings on Thursday, April 4, and then announced on the Church’s website. We invite you to sustain these brethren in their new assignments.

Those in favor, please manifest it.

Any opposed, by the same sign.

It is proposed that we sustain the following as the new Sunday School General Presidency, to be effective on August 1, 2024: Paul V. Johnson as President, Chad H Webb as First Counselor, and Gabriel W. Reid as Second Counselor.

Those in favor may manifest it.

Any opposed may so signify.

We note that Brother Reid is currently serving as president of the Australia Sydney Mission and is therefore not in Salt Lake City for the conference.

It is proposed that we sustain the other General Authorities, Area Seventies, and General Officers as presently constituted.

All in favor may do so by the uplifted hand.

Those opposed, if any.

Thank you, brothers and sisters, for your continued faith and prayers on behalf of the leadership of the Church.

\intalkheader{Changes to Area Seventies}

The following Area Seventies were sustained during a leadership session held as part of general conference:

Daniel A. Abeo, Mauricio A. Araújo, Randy T. Austin, Michel D. Avegnon, Philip J. Barton, Bradley S. Bateman, Eber Antônio Beck, Eric D. Bednar, Jared Black, Bryan G. Borela, Jaime A. Bravo, Juan G. Cardenas, Sancho N. Chukwu, Mark J Cluff, Danilo F. Costales, Daniel A. Cruzado, Gregorio Davalos, Julio N. Del Sero, Ryan E. Dobbs, Stephen W. Dyer, Brik V. Eyre, Denny Fa‘alogo, Timothy L. Farnes, Martín P. Fernández, Luis A. Ferrizo, Ángel J. Gómez, Georgie E. Guidi, Shinjiro Hara, Daniel L. Harris, Todd D. Haynie, Thomas Hengst, John R. Higgins, Niels O. Jensen, Fritzner A. Joseph, Kyoni Kasongo, John S. K. Kauwe III, Dan Kawashima, J. Joseph Kiehl, Carl F. Krauss, Yew Mun Kwan, Woo Cheol Lee, Wai Hung Mak, David R. Marriott, Ignatius Maziofa, Derek B. Miller, Albert Mutariswa, Marvin I. Palomo, Kyung Yeol Park, Domingo J. Perez, Oscar A. Perez, Raul Perez, Gayle L. Pollock, Pierre Portes, Marco A. Quezada, Stephen T. Rockwood, Guillermo Rojas, Kgomotso T. Sehloho, Sandro Alex Silva, Juswan Tandiman, Asuquo E. Udobong, Dwayne J. Van Heerden, Shih Ning (Steve) Yang, Juan F. Zorrilla, Leopoldo Zuñiga.

The following Area Seventies will be released on or before August 1, 2024:

Solomon I. Aliche, Guillermo A. Alvarez, Daren R. Barney, Julius F. Barrientos, James H. Bekker, David L. Buckner, Glenn Burgess, Marcos Cabral, Gregorio E. Casillas, Dunstan G. B. T. Chadambuka, Alan C. K. Cheung, Paul N. Clayton, Michael Cziesla, Hiroyuki Domon, Mernard P. Donato, I. Raymond Egbo, Zachary F. Evans, Sapele Fa‘alogo Jr., Saulo G. Franco, David Frischknecht, John J. Gallego, Efraín R. García, Robert Gordon, Mark A. Gottfredson, D. Martin Goury, Michael J. Hess, Bhanu K. Hiranandani, Richard S. Hutchins, Tito Ibañez, Eustache Ilunga, Akinori Ito, Anthony M. Kaku, Christopher H. Kim, H. Moroni Klein, Stephen Chee Kong Lai, V. Daniel Lattaro, Thabo Lebethoa, Tarmo Lepp, Itzcoatl Lozano, Kevin Lythgoe, Clement M. Matswagothata, Edgar P. Montes, Luiz C. D. Queiroz, Ifano Rasolondraibe, Eduardo D. Resek, Tomás G. Román, Ramon E. Sarmiento, Steven D. Shumway, Luis Spina, Jared W. Stone, Michael B. Strong, Djarot Subiantoro, Carlos G. Süffert, Voi R. Taeoalii, Karim Del Valle, Sergio R. Vargas, Helmut Wondra.

\newpage
\settalk{Covenants and Responsibilities}
\setspeaker{Dallin H. Oaks}
\settalkdate{April 2024}
\talkheading{Covenants and Responsibilities}{Dallin H. Oaks}

“How does your Church differ from others?” My answer to this important question has varied as I have matured and as the Church has grown. When I was born in Utah in 1932, our Church membership was only about 700,000, clustered mostly in Utah and nearby states. At that time, we had only 7 temples. Today the membership of The Church of Jesus Christ of Latter-day Saints numbers more than 17 million in about 170 nations. As of this April 1, we have 189 dedicated temples in many nations and 146 more in planning and construction. I have felt to speak about the purpose of these temples and the history and role of covenants in our worship. This will supplement the inspired teachings of earlier speakers.

\intalkheader{I.}

A covenant is a commitment to fulfill certain responsibilities. Personal commitments are essential to the regulation of our individual lives and to the functioning of society. This idea is currently being challenged. A vocal minority oppose institutional authority and insist that persons should be free from any restrictions that limit their individual freedom. Yet we know from millennia of experience that persons give up some individual freedoms to gain the advantages of living in organized communities. Such relinquishments of individual freedoms are principally based on commitments or covenants, expressed or implied.

Here are some examples of covenant responsibilities in our society: (1) judges, (2) military, (3) medical personnel, and (4) firefighters. All of those involved in these familiar occupations make a commitment—often formalized by oath or covenant—to perform their assigned duties. The same is true of our full-time missionaries. Distinctive clothing or name tags are intended to signify that the wearer is under covenant and therefore has a duty to teach and serve and should be supported in that service. A related purpose is to remind the wearers of their covenant responsibilities. There is no magic in their distinctive clothing or symbols, only a needed reminder of the special responsibilities the wearers have assumed. This is also true of the symbols of the engagement and wedding rings and their role in giving notice to observers or reminding wearers of covenant responsibilities.

\intalkheader{II.}

What I have said about covenants being a foundation for the regulation of individual lives applies particularly to religious covenants. The foundation and history of many religious affiliations and requirements are based on covenants. For example, the Abrahamic covenant is fundamental to several great religious traditions. It introduces the holy idea of God’s covenant promises with His children. The Old Testament frequently refers to God’s covenant with Abraham and his seed.

The first part of the Book of Mormon, which was written during the Old Testament period, clearly demonstrates the role of covenants in the Israelite history and worship. Nephi was told that the Israelite writings of that period were “a record of the Jews, which contains the covenants of the Lord, which he hath made unto the house of Israel.” The books of Nephi make frequent reference to the Abrahamic covenant and to Israel as “the covenant people of the Lord.” The practice of covenanting with God or religious leaders is also recorded in the Book of Mormon writings about Nephi, Joseph in Egypt, King Benjamin, Alma, and Captain Moroni.

\intalkheader{III.}

When the time came for the Restoration of the fulness of the gospel of Jesus Christ, God called a prophet, Joseph Smith. We do not know the full content of the angel Moroni’s early instructions to this maturing young prophet. We do know he told Joseph that “God had a work for [him] to do” and that “the fulness of the everlasting Gospel” must be brought forth, including “the promises made to the fathers.” We also know that the scriptures young Joseph read most intensively—even before he was directed to organize a church—were the many teachings about covenants he was translating in the Book of Mormon. That book is the Restoration’s major source for the fulness of the gospel, including God’s plan for His children, and the Book of Mormon is filled with references to covenants.

Being well read in the Bible, Joseph must have known of the book of Hebrews’ reference to the Savior’s intent to “make a new covenant with the house of Israel and with the house of Judah.” Hebrews also refers to Jesus as “the mediator of the new covenant.” Significantly, the biblical account of the Savior’s mortal ministry is titled “The New Testament,” a virtual synonym for “The New Covenant.”

Covenants were foundational in the Restoration of the gospel. This is evident in the earliest steps the Lord directed the Prophet to take in organizing His Church. As soon as the Book of Mormon was published, the Lord directed the organization of His restored Church, soon to be named The Church of Jesus Christ of Latter-day Saints. Revelation recorded in April 1830 directs that persons “shall be received by baptism into his church” after they “witness” (which means solemnly testify) “that they have truly repented of all their sins, and are willing to take upon them the name of Jesus Christ, having a determination to serve him to the end.”

This same revelation directs that the Church “meet together often to partake of bread and wine [water] in the remembrance of the Lord Jesus.” The importance of this ordinance is evident in the words of covenants specified for the elder or priest who officiates. He blesses the emblems of the bread for “the souls of all those who partake of it … , that they … witness unto thee, O God, the Eternal Father, that they are willing to take upon them the name of thy Son, and always remember him and keep his commandments which he has given them.”

The central role of covenants in the newly restored Church was reaffirmed in the preface the Lord gave for the first publication of His revelations. There the Lord declares that He has called Joseph Smith because the inhabitants of the earth “have strayed from mine ordinances, and have broken mine everlasting covenant.” This revelation further explains that His commandments are being given “that mine everlasting covenant might be established.”

Today we understand the role of covenants in the restored Church and the worship of its members. President Gordon B. Hinckley gave this summary of the effect of our baptism and our weekly partaking of the sacrament: “Every member of this church who has entered the waters of baptism has become a party to a sacred covenant. Each time we partake of the sacrament of the Lord’s supper, we renew that covenant.”

We have been reminded by many speakers at this conference that President Russell M. Nelson often refers to the plan of salvation as the “covenant path” that “leads us back to [God]” and “is all about our relationship with God.” He teaches about the significance of covenants in our temple ceremonies and urges us to see the end from the beginning and to “think celestial.”

\intalkheader{IV.}

Now I speak more of temple covenants. In fulfillment of his responsibility to restore the fulness of the gospel of Jesus Christ, the Prophet Joseph Smith spent much of his final years directing the construction of a temple in Nauvoo, Illinois. Through him the Lord revealed sacred teachings, doctrine, and covenants for his successors to administer in temples. There persons who were endowed were to be taught God’s plan of salvation and invited to make sacred covenants. Those who lived faithful to those covenants were promised eternal life, wherein “all things are theirs” and they “shall dwell in the presence of God and his Christ forever and ever.”

The endowment ceremonies in the Nauvoo Temple were administered just before our early pioneers were expelled to begin their historic trek to the mountains in the West. We have the testimonies of many pioneers that the power they received from being bound to Christ in their endowments in the Nauvoo Temple gave them the strength to make their epic journey and establish themselves in the West.

Persons who have been endowed in a temple are responsible to wear a temple garment, an article of clothing not visible because it is worn beneath outer clothing. It reminds endowed members of the sacred covenants they have made and the blessings they have been promised in the holy temple. To achieve those holy purposes, we are instructed to wear temple garments continuously, with the only exceptions being those obviously necessary. Because covenants do not “take a day off,” to remove one’s garments can be understood as a disclaimer of the covenant responsibilities and blessings to which they relate. In contrast, persons who wear their garments faithfully and keep their temple covenants continually affirm their role as disciples of the Lord Jesus Christ.

The Church of Jesus Christ of Latter-day Saints is constructing temples all over the world. Their purpose is to bless the covenant children of God with temple worship and with the sacred responsibilities and powers and unique blessings of being bound to Christ they receive by covenant.

The Church of Jesus Christ is known as a church that emphasizes making covenants with God. Covenants are inherent in each of the ordinances of salvation and exaltation this restored Church administers. The ordinance of baptism and its associated covenants are requirements for entrance into the celestial kingdom. The ordinances and associated covenants of the temple are requirements for exaltation in the celestial kingdom, which is eternal life, “the greatest of all the gifts of God.” That is the focus of The Church of Jesus Christ of Latter-day Saints.

I testify of Jesus Christ, who is the head of that Church, and invoke His blessings on all who seek to keep their sacred covenants. In the name of Jesus Christ, amen.

\newpage
\settalk{Following Christ}
\setspeaker{Dallin H. Oaks}
\settalkdate{October 2024}
\talkheading{Following Christ}{Dallin H. Oaks}

This year millions have been inspired by the gospel study plan known by the Savior’s invitation “Come, follow me.” Following Christ is not a casual or occasional practice. It is a continuous commitment and way of life that should guide us at all times and in all places. His teachings and His example define the path for every disciple of Jesus Christ. And all are invited to this path, for He invites all to come unto Him, “black and white, bond and free, male and female; … and all are alike unto God.”

\intalkheader{I.}

The first step in following Christ is to obey what He defined as “the great commandment in the law”:

“Thou shalt love the Lord thy God with all thy heart, and with all thy soul, and with all thy mind.

“This is the first and great commandment.

“And the second is like unto it, Thou shalt love thy neighbour as thyself.

“On these two commandments hang all the law and the prophets.”

The commandments of God provide the guiding and steadying force in our lives. Our experiences in mortality are like the little boy and his father flying a kite on a windy day. As the kite rose higher, the winds caused it to tug on the connecting string in the little boy’s hand. Inexperienced with the force of mortal winds, he proposed to cut the string so the kite could rise higher. His wise father counseled no, explaining that the string is what holds the kite in place against mortal winds. If we lose our hold on the string, the kite will not rise higher. It will be carried about by these winds and inevitably crash to the earth.

That essential string represents the covenants that connect us to God, our Heavenly Father, and His Son, Jesus Christ. As we honor those covenants by keeping Their commandments and following Their plan of redemption, Their promised blessings enable us to soar to celestial heights.

The Book of Mormon frequently declares that Christ is “the light of the world.” During His appearance to the Nephites, the risen Lord explained that teaching by telling them: “I have set an example for you.” “I am the light which ye shall hold up—that which ye have seen me do.” He is our role model. We learn what He has said and done by studying the scriptures and following prophetic teachings, as President Russell M. Nelson has urged us to do. In the ordinance of the sacrament, we covenant each Sabbath day that we will “always remember him and keep his commandments.”

\intalkheader{II.}

In the Book of Mormon, the Lord gave us the fundamentals in what He called “the doctrine of Christ.” These are faith in the Lord Jesus Christ, repentance, baptism, receiving the gift of the Holy Ghost, enduring to the end, and becoming as a little child, which means to trust the Lord and submit to all He requires of us.

The Lord’s commandments are of two types: permanent, like the doctrine of Christ, and temporary. Temporary commandments are those necessary for the needs of the Lord’s Church or the faithful in temporary circumstances, but to be set aside when the need has passed. An example of temporary commandments are the Lord’s directions to the early leadership of the Church to move the Saints from New York to Ohio, to Missouri, and to Illinois and finally to lead the pioneer exodus to the Intermountain West. Though only temporary, when still in force these commandments were given to be obeyed.

Some permanent commandments have taken considerable time to be generally observed. For example, President Lorenzo Snow’s famous sermon on the law of tithing emphasized a commandment given earlier but not yet generally observed by Church members. It needed reemphasis in the circumstances then faced by the Church and its members. Recent examples of reemphases have also been needed because of current circumstances faced by Latter-day Saints or the Church. These include the proclamation on the family, issued by President Gordon B. Hinckley a generation ago, and President Russell M. Nelson’s recent call for the Church to be known by its revealed name, The Church of Jesus Christ of Latter-day Saints.

\intalkheader{III.}

Another of our Savior’s teachings seems to require reemphasis in the circumstances of our day.

This is a time of many harsh and hurtful words in public communications and sometimes even in our families. Sharp differences on issues of public policy often result in actions of hostility—even hatred—in public and personal relationships. This atmosphere of enmity sometimes even paralyzes capacities for lawmaking on matters of importance where most citizens see an urgent need for some action in the public interest.

What should followers of Christ teach and do in this time of toxic communications? What were His teachings and examples?

It is significant that among the first principles Jesus taught when He appeared to the Nephites was to avoid contention. While He taught this in the context of disputes over religious doctrine, the reasons He gave clearly apply to communications and relationships in politics, public policy, and family relationships. Jesus taught:

“He that hath the spirit of contention is not of me, but is of the devil, who is the father of contention, and he stirreth up the hearts of men to contend with anger, one with another.

“Behold, this is not my doctrine, to stir up the hearts of men with anger, one against another; but this is my doctrine, that such things should be done away.”

In His remaining ministry among the Nephites, Jesus taught other commandments closely related to His prohibition of contention. We know from the Bible that He had previously taught each of these in His great Sermon on the Mount, usually in precisely the same language He later used with the Nephites. I will quote the familiar Bible language:

“Love your enemies, bless them that curse you, do good to them that hate you, and pray for them which despitefully use you, and persecute you.”

This is one of Christ’s best-known commandments—most revolutionary and most difficult to follow. Yet it is a most fundamental part of His invitation for all to follow Him. As President David O. McKay taught, “There is no better way to manifest love for God than to show an unselfish love for one’s fellowmen.”

Here is another fundamental teaching by Him who is our role model: “Blessed are the peacemakers: for they shall be called the children of God.”

Peacemakers! How it would change personal relationships if followers of Christ would forgo harsh and hurtful words in all their communications.

In general conference last year, President Russell M. Nelson gave us these challenges:

“One of the easiest ways to identify a true follower of Jesus Christ is how compassionately that person treats other people. …

“… True disciples of Jesus Christ are peacemakers.

“… One of the best ways we can honor the Savior is to become a peacemaker.”

Concluding his teachings: “Contention is a choice. Peacemaking is a choice. You have your agency to choose contention or reconciliation. I urge you to choose to be a peacemaker, now and always.”

Potential adversaries should begin their discussions by identifying common ground on which all agree.

To follow our Perfect Role Model and His prophet, we need to practice what is popularly known as the Golden Rule: “All things whatsoever ye would that men should do to you, do ye even so to them: for this is the law and the prophets.” We need to love and do good to all. We need to avoid contention and be peacemakers in all our communications. This does not mean to compromise our principles and priorities but to cease harshly attacking others for theirs. That is what our Perfect Role Model did in His ministry. That is the example He set for us as He invited us to follow Him.

In this conference four years ago, President Nelson gave us a prophetic challenge for our own day:

“Are you willing to let God prevail in your life? Are you willing to let God be the most important influence in your life? Will you allow His words, His commandments, and His covenants to influence what you do each day? Will you allow His voice to take priority over any other?”

As followers of Christ, we teach and testify of Jesus Christ, our Perfect Role Model. So let us follow Him by forgoing contention. As we pursue our preferred policies in public actions, let us qualify for His blessings by using the language and methods of peacemakers. In our families and other personal relationships, let us avoid what is harsh and hateful. Let us seek to be holy, like our Savior, in whose holy name I testify and invoke His blessing to help us be Saints. In the name of Jesus Christ, amen.

\newpage
\settalk{Divine Helps for Mortality}
\setspeaker{Dallin H. Oaks}
\settalkdate{April 2025}
\talkheading{Divine Helps for Mortality}{Dallin H. Oaks}

\intalkheader{I.}

Through the Prophet Joseph Smith, the Lord revealed a few things about our pre-earth life. There we existed as spirit children of God. Because God desired to help His children progress, He decided to create an earth on which we could receive a body, learn through experience, develop divine attributes, and be proven to see if we would keep God’s commandments. Those who qualified would “have glory added upon their heads for ever and ever” (Abraham 3:26).

To establish the conditions of this divine plan, God chose His Only Begotten Son to be our Savior. Lucifer, whose suggested alternative would destroy the agency of man, became Satan and was “cast down.” Banished to the earth and denied the privilege of mortal life, Satan was permitted to attempt “to deceive and to blind men, and to lead them captive at his will, even as many as would not hearken unto [God’s] voice” (Moses 4:4).

Essential to God’s great plan for the mortal growth of His children was for them to experience “opposition in all things” (2 Nephi 2:11). Just as our physical muscles cannot be developed or maintained without straining against the law of gravity, so mortal growth requires us to strain against Satan’s temptations and other mortal opposition. Most important for spiritual growth is the requirement to choose between good and evil. Those who choose good would progress toward their eternal destiny. Those who choose evil—as all would do in the various temptations of mortality—would need saving help, which a loving God designed to provide.

\intalkheader{II.}

By far, God’s strongest mortal help was His provision of a Savior, Jesus Christ, who would suffer to pay the price and provide forgiveness for repented sins. That merciful and glorious Atonement explains why faith in the Lord Jesus Christ is the first principle of the gospel. His Atonement “bringeth to pass the resurrection of the dead” (Alma 42:23), and it “atone[s] for the sins of the world” (Alma 34:8), erasing all of our repented sins and giving our Savior power to succor us in our mortal infirmities.

Beyond that glorious erasing of sins committed and being forgiven, the plan of a loving Father in Heaven provides many other gifts to protect us, including protecting us from sinning in the first place. Our mortal life always begins with a father and a mother. Ideally, both are present, with different gifts to guide our growth. If not, their absence is part of the opposition we must overcome.

\intalkheader{III.}

Our Heavenly Father’s plan provides other helps to guide us through our mortal journeys. I will speak of four of these. Please don’t hold me to my number of four, because these helps are overlapping. Moreover, there are other merciful protections in addition to these.

First, I speak of the Light or Spirit of Christ. In his great teaching in the book of Moroni, Moroni quotes his father that “the Spirit of Christ is given to every man, that he may know good from evil” (Moroni 7:16). We read this same teaching in modern revelations:

“And the Spirit giveth light to every man that cometh into the world; and the Spirit enlighteneth every man through the world, that hearkeneth to the voice of the Spirit” (Doctrine and Covenants 84:46).

Again: “For my Spirit is sent forth into the world to enlighten the humble and contrite, and to the condemnation of the ungodly” (Doctrine and Covenants 136:33).

President Joseph Fielding Smith explained these scriptures: “The Lord has not left men (when they are born into this world) helpless, groping to find the light and truth, but every man … is born with the right to receive the guidance, the instruction, the counsel of the Spirit of Christ, or Light of Truth.”

The second of the great assistances provided by the Lord to help us choose what is right is a cluster of divine directions in the scriptures as part of the plan of salvation (plan of happiness). These directions are commandments, ordinances, and covenants.

Commandments define the path our Heavenly Father has marked out for us to progress toward eternal life. People who imagine commandments as the way God decides who to punish fail to understand this purpose of God’s loving plan of happiness. On that path, we can gradually achieve the needed relationship with our Savior and qualify for an increase of His power to help us on our way to the destination He desires for all of us. Our Heavenly Father desires all of His children to return to the celestial kingdom, where God and our Savior reside, and to have the kind of life of those who reside in that celestial glory.

Ordinances and covenants are part of the law that defines the path to eternal life. Ordinances, and the sacred covenants we make with God through them, are required steps and essential guardrails along that path. I like to think of the role of covenants as demonstrating that under God’s plan, His highest blessings are given to those who promise in advance to keep certain commandments and who keep those promises.

Other God-given helps for making right choices are the manifestations of the Holy Ghost. The Holy Ghost is the third member of the Godhead. His function, defined in scripture, is to testify of the Father and the Son, to teach us, to bring all things to our remembrance, and to guide us into all truth. The scriptures include many descriptions of the manifestations of the Holy Ghost, such as a spiritual witness in response to an inquiry about the truth of the Book of Mormon. A manifestation is not to be confused with the gift of the Holy Ghost, which is conferred following baptism.

One of the most significant of God’s helps for His faithful children is the gift of the Holy Ghost. The importance of this gift is evident in the fact that it is formally conferred after repentance and baptism by water, “and then [the scriptures explain] cometh a remission of your sins by fire and by the Holy Ghost” (2 Nephi 31:17). Persons who have this remission of sins—and then regularly renew their cleansing by daily repentance and living according to the covenants they make through the ordinance of the sacrament—qualify for the promise that the Holy Ghost, the Spirit of the Lord, “may always … be with them” (Doctrine and Covenants 20:77).

Thus, President Joseph F. Smith taught that the Holy Ghost will “enlighten the minds of the people with regard to the things of God, to convince them at the time of their conversion of their having done the will of the Father, and to be in them an abiding testimony as a companion through life, acting as the sure and safe guide into all truth and filling them day by day with joy and gladness, with a disposition to do good to all men, to suffer wrong rather than to do wrong, to be kind and merciful, long suffering and charitable. All who possess this inestimable gift, this pearl of great price, have a continual thirst after righteousness. Without the aid of the Holy Spirit,” President Smith concluded, “no mortal can walk in the straight and narrow way.”

\intalkheader{IV.}

With so many powerful helps to guide us in our mortal journeys, it is disappointing that so many remain unprepared for their appointed meeting with our Savior and Redeemer, Jesus Christ. His parable of the ten virgins, spoken of so frequently in this conference, suggests that of those invited to meet Him, only half will be prepared.

We all know examples of the unprepared: returned missionaries who have interrupted their spiritual growth by periods of inactivity, youth who have jeopardized their spiritual growth by separating themselves from Church teaching and activities, men who have postponed their ordination to the Melchizedek Priesthood, men and women—sometimes the posterity of noble pioneers or worthy parents—who have departed the covenant path short of making and keeping covenants in the holy temple.

Many of such deviations occur when members fail to follow the fundamental spiritual maintenance plan of personal prayer, regular scripture study, and frequent repentance. In contrast, some neglect weekly renewal of covenants by not partaking of the sacrament. Some say the Church is not meeting their needs; those substitute what they perceive as their future needs ahead of what the Lord has provided in His many teachings and opportunities for our essential service to others.

Humility and trust in the Lord are the remedies for such deviations. As the Book of Mormon teaches, the Lord “doth bless and prosper those who put their trust in him” (Helaman 12:1). Trusting in the Lord is a particular need for all who wrongly measure the commandments of God and the teachings of His prophets against the latest findings and wisdom of man.

I have spoken of the many mortal helps our loving Father in Heaven has given to help His children return to Him. Our part in this divine plan is to trust in God and seek and use these divine helps, most notably the Atonement of His Beloved Son, our Savior and Redeemer, Jesus Christ. I pray that we will teach and live these principles, in the name of Jesus Christ, amen.

\newpage
\settalk{Introduction}
\setspeaker{Dallin H. Oaks}
\settalkdate{October 2025}
\talkheading{Introduction}{Dallin H. Oaks}

My dear brothers and sisters, we meet as we mourn the death of our beloved President Russell M. Nelson. I speak to you as the President of the Quorum of the Twelve, in which position I have already been sustained in each of our ward, stake, and general conferences for the past seven and a half years.

We go forward with this semiannual general conference, proceeding as planned and with speakers and music assigned by President Nelson after months of preparation. The only additions are this, my introductory explanation, and my speaking in place of President Nelson in the closing session on Sunday afternoon.

This is the first time in about 75 years that a Church President has died a few days before a general conference. That has called for us to plan how we could hold the vital leadership meetings in connection with general conference as well as the general sessions of conference, at the same time scheduling his funeral as soon as possible. President Nelson understood the value of general conference to provide direction to the Saints in the coming months. We honor him by following the planned conference schedule he approved.

In this time of grief, we have already held a meeting to pay tribute to President Russell M. Nelson before our general conference meetings and then scheduled his funeral as soon as possible after conference. All of this has been done with the approval of the family of President Nelson and the Quorum of the Twelve.

Since all of us who have been assigned to speak in general conference would like to devote our assigned time to paying a personal tribute to President Nelson, we have therefore asked that all of our conference speakers hold any such tributes to a minimum, deferring elaborate tributes for the funeral, which we have already begun to plan.

My own brief tribute—suitable to this conference and to the tributes we have already paid—is this: I loved Russell M. Nelson and have learned more about the gospel and gospel leadership from my long friendship and association with him than from any other leader I have personally known. He is our model as a servant and follower of the Lord Jesus Christ. Of these things I testify, in the name of Jesus Christ, amen.

\newpage
\settalk{The Family-Centered Gospel of Jesus Christ}
\setspeaker{Dallin H. Oaks}
\settalkdate{October 2025}
\talkheading{The Family-Centered Gospel of Jesus Christ}{Dallin H. Oaks}

My loving brothers and sisters, thank you for your prayers on my behalf. I have felt them.

\intalkheader{I.}

The doctrine of The Church of Jesus Christ of Latter-day Saints centers on the family. Essential to our doctrine on the family is the temple. The ordinances received there enable us to return as eternal families to the presence of our Heavenly Father.

As of the April 2025 general conference, President Russell M. Nelson had announced the construction of 200 new temples. He loved to announce new temples at the conclusion of each general conference, and we all rejoiced with him. However, with the large number of temples now in the very earliest phases of planning and construction, it is appropriate that we slow down the announcement of new temples. Therefore, with the approval of the Quorum of the Twelve Apostles, we will not announce any new temples at this conference. We will now move forward in providing the ordinances of the temple to members of the Church throughout the world, including when and where to announce the construction of new temples.

The portion of my talk that I have just delivered was written after the death of our beloved President Russell M. Nelson. What now follows was written and approved weeks before, but it still represents my teachings, inspired by the Lord.

\intalkheader{II.}

The family proclamation, announced 30 years ago, declares that “the family is ordained of God” and “is central to the Creator’s plan for the eternal destiny of His children.” It also declares “that God’s commandment for His children to multiply and replenish the earth remains in force.” And “we further declare that God has commanded that the sacred powers of procreation are to be employed only between man and woman, lawfully wedded as husband and wife.” As then-Elder Russell M. Nelson taught a Brigham Young University audience, the family is “pivotal to God’s plan. … In fact, a purpose of the plan is to exalt the family.”

The Church of Jesus Christ is sometimes known as a family-centered church. It is! Our relationship to God and the purpose of our mortal life are explained in terms of the family. The gospel of Jesus Christ is the plan of our Heavenly Father for the benefit of His spirit children. We can truly say that the gospel plan was first taught to us in the council of an eternal family, it is implemented through our mortal families, and its intended destiny is to exalt the children of God in eternal families.

\intalkheader{III.}

Despite that doctrinal context, there is opposition. In the United States we are suffering from a deterioration in marriage and childbearing. For nearly a hundred years the proportion of households headed by married couples has declined, and so has the birthrate. The marriages and birthrates of our Church members are much more positive, but they have also declined significantly. It is vital that Latter-day Saints do not lose their understanding of the purpose of marriage and the value of children. That is the future for which we strive. “Exaltation is a family affair,” President Nelson has taught us. “Only through the saving ordinances of the gospel of Jesus Christ can families be exalted.”

The national declines in marriage and childbearing are understandable for historic reasons, but Latter-day Saint values and practices should improve—not follow—those trends.

In my boyhood 80 years ago, I lived on my grandparents’ farm in a setting where almost all that happened during the day was under the direction of the family. There was no television or other electronics to distract from family activities. In contrast, in today’s urban society, few members experience consistent family-centered activities. Urban living and modern transportation, organized entertainment, and high-speed communication have made it easy for youth to treat their homes as boardinghouses, where they sleep and take an occasional meal but where there is far less parental direction of their activities.

Parental influences have also been diluted by the way in which most current members of the Church earn a living. In times past, one of the great influences that unified families was the experience of struggling together in pursuit of a common goal—such as taming the wilderness or earning a living. The family was an organized and conducted unit of economic production. Today, most families are units of economic consumption, which do not require a high degree of family organization and cooperation.

\intalkheader{IV.}

As parental influences diminish, Latter-day Saints still have a God-given responsibility to teach their children to prepare for our family destiny in eternity (see Doctrine and Covenants 68:25). Many of us must do this when not all of our families are traditional. Divorce, death, and separation are realities. I experienced that in the family in which I was raised.

My father died when I was seven years old, so my younger brother and sister and I were raised by a widowed mother. In the most difficult of situations, she pressed on. She was alone and broken, but with the Lord’s help, her powerful teaching of the doctrine of the restored Church guided us. How she prayed for heavenly assistance in raising her children, and she was blessed! We were raised in a happy home in which our deceased father was always a reality. She taught us that we had a father and she had a husband and we would always be a family because of their temple marriage. Our father was just away temporarily because the Lord had called him to a different work.

I know that many other families are not so happy, but every single mother can teach of the love of a Heavenly Father and the eventual blessings of a temple marriage. You too can do this! Heavenly Father’s plan assures this possibility for everyone. We are all grateful for temple marriage and for the prospective blessings of being sealed as an eternal family. Like my mother, we love to quote Lehi’s promise to his son Jacob that God “shall consecrate thine afflictions for thy gain” (2 Nephi 2:2). That applies to every Latter-day Saint family, complete or currently incomplete. We are a family church.

Our doctrine and our belief in eternal families strengthen and bond us. I will never forget the promise of my maternal Grandfather Harris, when we children were living on his farm near Payson, Utah. He gave me the tragic news that my father had died in faraway Denver, Colorado. I ran into the bedroom and knelt beside the bed, crying my heart out. Grandpa followed me and went to his knees beside me and said, “I will be your father.” That tender promise is a powerful example of what grandparents can do to fill in the gaps when families lose or are missing a member.

Parents, single or married—and others, like grandparents, who fill that role for children—are the master teachers. Their most effective teaching is by example. The family circle is the ideal place to demonstrate and learn eternal values, such as the importance of marriage and children, the purpose of life, and the true source of joy. It is also the best place to learn other essential lessons of life, such as kindness, forgiveness, self-control, and the value of education and honest work.

Of course, many Church members have beloved family members who do not embrace gospel values and expectations. Such members need our love and patience. In relating to one another, we should remember that the perfection we seek is not limited to the stressful circumstances of mortality. The great teaching in Doctrine and Covenants 138:57–59 assures us that repentance and spiritual growth can continue in the spirit world that follows mortality. More important, as families unite to strengthen one another, we should all remember that the sins and inevitable shortcomings all of us experience in mortality can be forgiven through repentance because of the glorious and saving Atonement of Jesus Christ.

\intalkheader{V.}

Our Savior, Jesus Christ, is our ultimate role model. We will be blessed if we model our lives after His teachings and self-sacrifice. Following Christ and giving ourselves in service to one another is the best remedy for the selfishness and individualism that now seem to be so common.

Parents also have a duty to teach their children practical knowledge apart from gospel principles. Families unite when they do meaningful things together. Family gardens build family relationships. Happy family experiences strengthen family ties. Camping, sports activities, and other recreation are especially valuable to bond families. Families should organize family reunions to remember ancestors, which lead to the temple.

Parents should educate children in the basic skills of living, including working in the yard and home. Learning languages is a useful preparation for missionary service and modern life. The teachers of these subjects can be parents or grandparents or members of the extended family. Families flourish when they learn as a group and counsel together on all matters of concern to the family and its members.

Some may say, “But we have no time for any of that.” To find time to do what is truly worthwhile, many parents will find that they can turn their family on if they all turn their technologies off. And parents, remember, what those children really want for dinner is time with you.

Great blessings come to families if they pray together, kneeling night and morning to offer thanks for blessings and to pray over common concerns. Families are also blessed as they worship together in Church services and in other devotional settings. Family bonds are also strengthened by family stories, creating family traditions, and sharing sacred experiences. President Spencer W. Kimball reminded us that “stories of inspiration from our own lives and those of our forebears … are powerful teaching tools.” They are often the best sources of inspiration for us and our posterity.

I testify of the Lord Jesus Christ, who is the Only Begotten Son of God, our Eternal Father. He invites us to follow the covenant path that leads to a heavenly family reunion. The sealing powers of the priesthood, directed by the keys restored in the Kirtland Temple, bring families together for eternity (see Doctrine and Covenants 110:13–16). They are currently being exercised in a growing number of temples of the Lord throughout the world. This is real. Let us be part of it, I pray, in the name of Jesus Christ, amen.

\end{document}
